\documentclass[10pt,letterpaper,twoside]{book}
\usepackage{amsmath}
\usepackage{amssymb}
\usepackage{color}
% Latin Modern replaces the unhinted cm-super T1 fonts, which render thin and washed out on
% screen. It is metric-compatible with Computer Modern, so pagination is unchanged, and it also
% supplies the bold typewriter face that cmtt lacks.
\usepackage{lmodern}
\usepackage[T1]{fontenc}
\usepackage{geometry}
\usepackage{graphicx}
\usepackage{hyperref}
\usepackage{listings}
\usepackage{makeidx}
\usepackage{needspace}
\usepackage{paralist}
\usepackage{sectsty}
\usepackage{tcolorbox}
\usepackage{textcomp}
\usepackage{upquote}
\tcbuselibrary{breakable,listings}
\geometry{left=2cm,right=2cm,top=2.5cm,bottom=2.5cm}

% Keep natural vertical spacing instead of stretching glue to fill each page,
% which otherwise opens large gaps (e.g. below headings) around big listings.
\raggedbottom

% Never strand the first or last line of a paragraph (or code listing, which
% is set as one paragraph) alone at a page boundary.
\clubpenalty=10000
\widowpenalty=10000
\displaywidowpenalty=10000

% Lowercase section headers
\let\MakeUppercase\relax

% A cross-reference to a section, written as its number in a source docstring. Restyle here to
% change every such reference at once.
\newcommand{\secref}[1]{\hyperref[sec:#1]{#1}}

% No colored border on hyperlinks
\hypersetup{colorlinks=false,pdfborder={0 0 0}}

% For the horizontal rule under sections
\sectionfont{\sectionrule{0pt}{0pt}{-10pt}{0.8pt}}

% For displaying code. Hue carries the signal in syntax highlighting, not darkness, so these stay
% saturated rather than dark. The first three are the palette the book has always used; the other
% two extend it in the same key.
\lstloadaspects{directives}  % so directivestyle is a known key at preamble time
\definecolor{codekw}{HTML}{0000FF}       % language keywords
\definecolor{codecomment}{HTML}{009900}  % comments
\definecolor{codestring}{HTML}{9400D1}   % string and character literals
\definecolor{codetype}{HTML}{008B8B}     % types, std, and standard library headers
\definecolor{codelit}{HTML}{CC6600}      % true/false/nullptr and preprocessor directives
\lstset{
  aboveskip=5mm,
  basicstyle={\footnotesize\ttfamily},
  belowskip=5mm,
  breaklines=true,
  breakatwhitespace=true,
  frame=none,
  columns=flexible,
  keepspaces=true,
  rulecolor=\color{black},
  showstringspaces=false,
  tabsize=2,
  upquote=true
}
\lstdefinestyle{cpp}{
  language=C++,
  commentstyle=\color{codecomment},
  directivestyle=\color{codelit},
  identifierstyle=\color{black},
  stringstyle=\color{codestring},
  deletekeywords={auto,bool,char,double,float,int,long,short,signed,unsigned,void,wchar_t,
    true,false,register,export,bad_cast,bad_typeid,type_info},
  keywordstyle=\color{codekw},
  morekeywords={alignas,alignof,constexpr,decltype,final,noexcept,override,static_assert,
    thread_local},
  morekeywords=[2]{std,auto,bool,char,char16_t,char32_t,double,float,int,long,short,signed,
    unsigned,void,wchar_t,size_t,ptrdiff_t,int8_t,int16_t,int32_t,int64_t,intmax_t,uint8_t,
    uint16_t,uint32_t,uint64_t,uintmax_t,string,vector,pair,tuple,array,deque,list,map,set,
    multimap,multiset,unordered_map,unordered_set,queue,stack,priority_queue,bitset,optional,
    complex,function,mt19937,mt19937_64,stringstream,istringstream,ostringstream,unique_ptr,
    uniform_int_distribution,uniform_real_distribution,
    algorithm,cassert,cctype,chrono,climits,cmath,cstddef,cstdint,cstdio,cstdlib,cstring,ctime,
    functional,initializer_list,iomanip,iostream,istream,iterator,limits,memory,numeric,ostream,
    random,regex,sstream,stdexcept,thread,type_traits,utility},
  keywordstyle=[2]\color{codetype},
  morekeywords=[3]{true,false,nullptr,NULL},
  keywordstyle=[3]\color{codelit}
}

\newtcblisting{exampleoutput}{
  listing only,
  title={Example Output},
  listing options={
    language={},
    numbers=none,
    basicstyle={\footnotesize\ttfamily},
    aboveskip=0pt,
    belowskip=0pt,
    breaklines=true,
    columns=flexible,
    keepspaces=true,
    showstringspaces=false,
    upquote=true
  },
  breakable,
  colback=black!4,
  colframe=black!30,
  boxrule=0.4pt,
  arc=2mm,
  boxsep=1mm,
  left=2mm,
  right=2mm,
  top=0mm,
  bottom=0mm,
  before skip=0.5em,
  after skip=\baselineskip,
  fonttitle=\small\bfseries,
  colbacktitle=black!7,
  coltitle=black
}

\newtcbox{\inlinecode}{
  on line,
  colback=black!5,
  colframe=black!20,
  boxrule=0.3pt,
  arc=0.8mm,
  boxsep=0.2mm,
  left=0.5mm,
  right=0.5mm,
  top=0mm,
  bottom=0mm,
  fontupper=\footnotesize\ttfamily
}

% Same grey rounded box as \inlinecode, but a step smaller, for code (e.g. variable
% names) that appears inside math mode. \mbox keeps the tcolorbox safe within a formula.
\newtcbox{\inlinecodemathbox}{
  on line,
  colback=black!5,
  colframe=black!20,
  boxrule=0.3pt,
  arc=0.8mm,
  boxsep=0.2mm,
  left=0.5mm,
  right=0.5mm,
  top=0mm,
  bottom=0mm,
  fontupper=\scriptsize\ttfamily
}
\newcommand{\inlinecodemath}[1]{\mbox{\inlinecodemathbox{#1}}}

\begin{document} 

\providecommand{\BookVersion}{v1.0}
\providecommand{\BookDate}{August 19, 2026}
 
\frontmatter
\pagenumbering{gobble}

\begin{titlepage}
	\centering
	\newgeometry{left=2.5cm,right=2.5cm,top=2.5cm,bottom=2.5cm}

	\vspace*{6 cm}

	\textbf{\LARGE Alex's Anthology of Algorithms}
	\\[1.0\baselineskip]
	\textbf{\Large Common Code for Contests in Concise C++}
	\\[1.0\baselineskip]
	{\large\textit{\BookVersion, \BookDate}}

	\vspace*{11 cm}

	{\Large Alex Li}
	\\[0.5\baselineskip]
	{\small \textcircled{c} 2026. All rights reserved.}
\end{titlepage}

\clearpage
\thispagestyle{empty}
\vspace*{8 cm}
\centering

{\large This page is intentionally left blank.}

\vspace{\fill}
\clearpage
\makeatletter

\let\cleardoublepage\clearpage
\pagenumbering{roman}
\chapter{Preface}

\raggedright
\setlength{\parskip}{0.5em}

Browse the latest version online at: 
\href{https://algorithms.alexli.ca}{\texttt{algorithms.alexli.ca}}

Contributions are always welcome on GitHub: 
\href{https://github.com/alxli/algorithm-anthology}{\texttt{github.com/alxli/algorithm-anthology}}

\section*{Introduction}

Welcome to a practical collection of concise C++ implementations of algorithms, data structures, and
contest utilities. This is not a first textbook; it is a codebook about implementation details,
interfaces, edge cases, and the choices that make familiar algorithms usable in real code. Use it:
\begin{itemize}
	\item as a reference when you know an algorithm conceptually but want a clean implementation;
	\item as a printed codebook for contests that allow prepared material, including ICPC-style
	events;
	\item as a checklist for edge cases, complexity, overflow, indexing, and API choices;
	\item as a map of related techniques: slower versions, faster versions, variants, and common
	reductions sit near one another.
\end{itemize}

For learning, do not start by copying. Study the idea, prove the invariant or recurrence, implement
it yourself, then compare; the differences reveal the practical choices. Each section gives only
enough theory to orient the implementation: problem, interface, assumptions, and complexity. Source
files are organized as documentation, reusable implementation, then example usage. The reusable part
ends at the \texttt{Example Usage} marker; everything after exists only so the file compiles and
demonstrates itself. When adapting a section, drop the example block and keep the implementation
above it.

Guiding principles:
\begin{itemize}
	\item \textit{Clarity:} A reader familiar with the algorithm should follow the code without
	reverse-engineering it. Tricks, non-obvious invariants, overflow risks, and language assumptions
	are documented.
	\item \textit{Concision:} Contest code is found, read, adapted, and typed under pressure. Short
	code helps when it removes noise, not when it hides the idea. So the code here is often more
	verbose than compact ICPC snippets; the extra lines buy clearer names, configurable types, safer
	defaults, reset functions, explicit handling of variants, and more.
	\item \textit{Efficiency:} Expected asymptotic behavior and reasonable constants. Both simpler
	and more robust versions are often included, but educational-only sections are labeled so.
	\item \textit{Adaptability:} Interfaces make common changes easy, such as swapping a numeric type,
	changing a combine function, exposing an edge ID, or choosing whether boundaries are included.
	\item \textit{Consistency:} Similar structures expose similar operations, with repeated naming and
	example styles, so moving between chapters is easy.
	\item \textit{Portability:} Targets modern C++17 and avoids unnecessary dependencies; the few
	contest-useful exceptions are noted at each use.
\end{itemize}

This began as a personal contest codebook, so it favors single-file snippets, predictable APIs,
compact examples, and low constants. Use it actively: annotate it, delete what you dislike, rename
functions to match your habits, and stress-test anything you intend to trust. Cheers!

\section*{Portability Notes}

Examples target a modern GCC or Clang environment with C++17. A typical compile command:

\begin{verbatim}
   g++ -std=c++17 -O2 -Wall -pedantic
\end{verbatim}

Most code assumes common contest-platform type sizes: 8-bit \texttt{char}, 32-bit \texttt{int} and
\texttt{float}, 64-bit \texttt{double} and \texttt{long long}. \texttt{long double} width and
precision are implementation-dependent; sections relying on it say so.

The code targets ISO C++17 unless noted. The most common exceptions:

\begin{itemize}
	\item GCC/Clang integer extensions like \texttt{\_\_int128} or \texttt{\_\_uint128\_t}, usually
	guarded with a portable fallback when overflow safety matters.
	\item Compiler builtins like \texttt{\_\_builtin\_popcount()} and \texttt{\_\_builtin\_clz()} ---
	convenient and widely available, though C++20's \texttt{\textless{}bit\textgreater{}} header
	standardizes many of them. See \url{https://gcc.gnu.org/onlinedocs/gcc/Other-Builtins.html} for
	GCC's reference.
	\item GNU policy-based data structures like \texttt{\_\_gnu\_pbds::tree}, used in a few sections
	to show useful non-standard components.
	\item A few documented representation assumptions, such as \texttt{signbit\_()} in the math
	utilities assuming an IEEE-like floating-point sign bit layout.
\end{itemize}

\pdfbookmark[0]{Contents}{toc}
\tableofcontents

\mainmatter
\chapter{\texorpdfstring{Elementary Algorithms}{1. Elementary Algorithms}}

\section{\texorpdfstring{Array Transformations}{1.1 Array Transformations}}
\setcounter{section}{1}
\setcounter{subsection}{0}
\subsection{Sorting Algorithms}
The following functions are equivalent to \inlinecode{std::sort()}, taking random-access iterators as a half-open range $[{\inlinecodemath{lo}}, {\inlinecodemath{hi}})$ to be sorted. The range is sorted into ascending order after the function call. Optionally, a comparison function object specifying a strict weak ordering may be specified to replace the default \inlinecode{operator\textless{}}.

\begin{compactitem}
  \item \inlinecode{quicksort(lo, hi, comp = std::less\textless{}\textgreater{})} sorts the range using quicksort.
  \item \inlinecode{mergesort(lo, hi, comp = std::less\textless{}\textgreater{})} sorts the range using merge sort, which is stable.
  \item \inlinecode{heapsort(lo, hi, comp = std::less\textless{}\textgreater{})} sorts the range using heapsort.
  \item \inlinecode{combsort(lo, hi, comp = std::less\textless{}\textgreater{})} sorts the range using comb sort.
  \item \inlinecode{insertion\_sort(lo, hi, comp = std::less\textless{}\textgreater{})} sorts the range using insertion sort, which is stable.
  \item \inlinecode{shellsort(lo, hi, comp = std::less\textless{}\textgreater{})} sorts the range using shell sort.
  \item \inlinecode{counting\_sort(lo, hi, keys, key)} sorts the range by the index that \inlinecode{key(element)} returns in $[0, {\inlinecodemath{keys}})$, which is stable; unlike the shared interface above, it takes no comparator.
  \item \inlinecode{radix\_sort(lo, hi)} sorts an integer range using least-significant-byte radix sort; unlike the shared interface above, it takes no comparator.
  \item \inlinecode{bitonic\_sort(lo, hi, comp = std::less\textless{}\textgreater{})} sorts a range of any length using a compare-exchange sequence generated from that length.
\end{compactitem}

These functions are not meant to compete with standard library implementations in terms of speed. Instead, they demonstrate how common sorting algorithms can be concisely implemented in C++.

\begin{lstlisting}[language=C++]
#include <algorithm>
#include <climits>
#include <functional>
#include <iterator>
#include <type_traits>
#include <vector>
\end{lstlisting}
Quicksort repeatedly selects a pivot and partitions the range so that elements comparing less than the pivot precede it, elements comparing equal stay in the middle, and elements comparing greater follow it. Divide and conquer is then applied to the two outer ranges until the original range is sorted. The swaps used during partitioning make quicksort unstable. Despite having a worst case of $O(n^{2})$, quicksort is often faster in practice than merge sort and heapsort, which both have a worst case time complexity of $O(n \log n)$.

The pivot chosen in this implementation is always a middle element of the range to be sorted. Shuffling the input first or choosing each pivot randomly makes consistently unbalanced partitions unlikely. Median-of-three is another practical choice, while median-of-medians can guarantee a balanced pivot at greater constant-factor cost.

\textbf{Time Complexity}: $O(n)$ best (all equal), $O(n \log n)$ average, and $O(n^{2})$ worst.

\textbf{Space Complexity}: $O(\log n)$ auxiliary stack space.

\begin{lstlisting}[language=C++]
template<typename It, typename Compare = std::less<>>
void quicksort(It lo, It hi, Compare comp = Compare{}) {
  using T = typename std::iterator_traits<It>::value_type;
  while (hi - lo >= 2) {
    T pivot = *(lo + (hi - lo) / 2);
    It lt = lo, mid = lo, gt = hi;
    while (mid != gt) {
      if (comp(*mid, pivot)) {
        std::iter_swap(lt++, mid++);
      } else if (comp(pivot, *mid)) {
        std::iter_swap(mid, --gt);
      } else {
        ++mid;
      }
    }
    if (lt - lo < hi - gt) {
      quicksort(lo, lt, comp);
      lo = gt;
    } else {
      quicksort(gt, hi, comp);
      hi = lt;
    }
  }
}
\end{lstlisting}
Merge sort first divides a list into $n$ sublists of one element each, then recursively merges the sublists into sorted order until only a single sorted sublist remains. Merge sort is a stable sort, meaning that it preserves the relative order of elements which compare equal by \inlinecode{operator\textless{}} or the custom comparator given.

An analogous function in the C++ standard library is \inlinecode{std::stable\_sort()}, except that the implementation here requires sufficient memory to be available. When $O(n)$ auxiliary memory is not available, \inlinecode{std::stable\_sort()} falls back to a time complexity of $O(n \log^{2} n)$ whereas the implementation here will simply fail.

\textbf{Time Complexity}: $O(n \log n)$ in all cases.

\textbf{Space Complexity}: $O(\log n)$ auxiliary stack space and $O(n)$ auxiliary heap space.

\begin{lstlisting}[language=C++]
template<typename It, typename Compare = std::less<>>
void mergesort(It lo, It hi, Compare comp = Compare{}) {
  if (hi - lo < 2) {
    return;
  }
  It mid = lo + (hi - lo - 1) / 2, a = lo, c = mid + 1;
  mergesort(lo, mid + 1, comp);
  mergesort(mid + 1, hi, comp);
  using T = typename std::iterator_traits<It>::value_type;
  std::vector<T> merged;
  merged.reserve(hi - lo);
  while (a <= mid && c < hi) {
    merged.push_back(comp(*c, *a) ? *c++ : *a++);
  }
  merged.insert(merged.end(), a, mid + 1);
  merged.insert(merged.end(), c, hi);
  std::copy(merged.begin(), merged.end(), lo);
}
\end{lstlisting}
Heapsort first rearranges an array to satisfy the max-heap property. Then, it repeatedly pops the max element of the heap (the left, unsorted subrange), moving it to the beginning of the right, sorted subrange until the entire range is sorted. Heapsort has a better worst case time complexity than quicksort and also a better space complexity than merge sort, but its swaps make it unstable.

The C++ standard library equivalent is calling \inlinecode{std::make\_heap(lo, hi)}, followed by \inlinecode{std::sort\_heap(lo, hi)}.

\inlinecode{sift\_down()} restores the heap property below one node. The first loop applies it to each non-leaf from right to left, which builds the heap in $O(n)$ time. The second loop repeatedly moves the maximum to the back of the unsorted range and restores the heap among the remaining elements.

\textbf{Time Complexity}: $O(n \log n)$ in all cases.

\textbf{Space Complexity}: $O(1)$ auxiliary.

\begin{lstlisting}[language=C++]
template<typename It, typename Compare = std::less<>>
void heapsort(It lo, It hi, Compare comp = Compare{}) {
  if (hi - lo < 2) {
    return;
  }
  using T = typename std::iterator_traits<It>::value_type;
  auto sift_down = [&](It root, It end) {
    T value = *root;
    auto parent = root - lo, n = end - lo;
    while (2 * parent + 1 < n) {
      auto child = 2 * parent + 1;
      if (child + 1 < n && comp(*(lo + child), *(lo + child + 1))) {
        child++;
      }
      if (!comp(value, *(lo + child))) {
        break;
      }
      *(lo + parent) = *(lo + child);
      parent = child;
    }
    *(lo + parent) = value;
  };
  for (It root = lo + (hi - lo) / 2; root != lo;) {
    sift_down(--root, hi);
  }
  for (It end = hi; --end != lo;) {
    std::iter_swap(lo, end);
    sift_down(lo, end);
  }
}
\end{lstlisting}
Comb sort is an improved bubble sort. While bubble sort compares only adjacent elements, comb sort compares elements separated by a fixed gap, decreasing that gap after every pass. Large gaps move small values near the end toward the front quickly, avoiding the slow movement of such values in bubble sort. These nonadjacent swaps make comb sort unstable. The shrink factor $1.3$ is a common empirical choice.

\textbf{Time Complexity}: $O(n \log n)$ best and $O(n^{2})$ worst.

\textbf{Space Complexity}: $O(1)$ auxiliary.

\begin{lstlisting}[language=C++]
template<typename It, typename Compare = std::less<>>
void combsort(It lo, It hi, Compare comp = Compare{}) {
  int gap = static_cast<int>(hi - lo);
  bool swapped = true;
  while (gap > 1 || swapped) {
    if (gap > 1) {
      gap = gap * 10 / 13;
    }
    swapped = false;
    for (It it = lo; it + gap < hi; ++it) {
      if (comp(*(it + gap), *it)) {
        std::iter_swap(it, it + gap);
        swapped = true;
      }
    }
  }
}
\end{lstlisting}
Insertion sort builds the sorted range one element at a time. It scans left to right, and for each element shifts the larger elements of the already-sorted prefix one position to the right to open a slot where the element belongs. Although its average and worst cases are quadratic, it is simple, stable, in-place, and online (it can sort a stream as elements arrive).

Its strength is being adaptive: on an already-sorted or nearly-sorted range, each element travels only a short distance, giving $O(n)$ best-case time and $O(n + d)$ in general for $d$ total inversions. This is why it outperforms the asymptotically faster sorts on small or nearly-sorted ranges, and why hybrid sorts such as introsort and Timsort fall back to it once a subrange is small enough.

The comparison \inlinecode{comp(key, *(j - 1))} is strict, so an element never moves past an earlier element it compares equal to, which keeps the sort stable.

\textbf{Time Complexity}: $O(n + d)$, i.e. $O(n)$ best and $O(n^{2})$ average/worst.

\textbf{Space Complexity}: $O(1)$ auxiliary.

\begin{lstlisting}[language=C++]
template<typename It, typename Compare = std::less<>>
void insertion_sort(It lo, It hi, Compare comp = Compare{}) {
  using T = typename std::iterator_traits<It>::value_type;
  for (It i = lo; i != hi; ++i) {
    T key = *i;
    It j = i;
    while (j != lo && comp(key, *(j - 1))) {
      *j = *(j - 1);
      --j;
    }
    *j = key;
  }
}
\end{lstlisting}
Shell sort is insertion sort performed on subsequences of elements a fixed gap apart, repeating with smaller gaps until the final pass uses a gap of one and is an ordinary insertion sort. The earlier passes are what make that final pass cheap: each one leaves the range closer to sorted, and insertion sort runs in time proportional to the number of remaining inversions. Comparing elements a gap apart also moves a value many positions in one step, which is what insertion sort alone cannot do. Like comb sort above, the nonadjacent swaps make it unstable.

The gap sequence determines the bound, and choosing it well is the whole difficulty. This implementation uses Ciura's empirical sequence, extended upward by a factor of $2.25$, which is the fastest known for small and medium ranges though it has no proven bound. The theoretical choices are Sedgewick's sequence at $O(n^{4/3})$ and Pratt's at $O(n \log^{2} n)$, the latter being slower in practice despite the better bound.

\textbf{Time Complexity}: $O(n \log n)$ best, and $O(n^{4/3})$ or better in practice for the gaps used here.

\textbf{Space Complexity}: $O(1)$ auxiliary.

\begin{lstlisting}[language=C++]
template<typename It, typename Compare = std::less<>>
void shellsort(It lo, It hi, Compare comp = Compare{}) {
  static const int GAPS[] = {1, 4, 10, 23, 57, 132, 301, 701, 1577, 3548, 7983, 17962};
  int n = static_cast<int>(hi - lo), count = sizeof(GAPS) / sizeof(GAPS[0]);
  for (int g = count - 1; g >= 0; g--) {
    int gap = GAPS[g];
    if (gap >= n) {
      continue;
    }
    for (It it = lo + gap; it < hi; ++it) {
      typename std::iterator_traits<It>::value_type key = *it;
      It j = it;
      while (j - lo >= gap && comp(key, *(j - gap))) {
        *j = *(j - gap);
        j -= gap;
      }
      *j = key;
    }
  }
}
\end{lstlisting}
Counting sort sorts values from a small integer range by counting how many times each value occurs, converting those counts into starting offsets with a prefix sum, and then placing each element at the offset for its key. It never compares two elements, so the $O(n \log n)$ lower bound for comparison sorts does not apply, and it runs in linear time whenever the range of keys is proportional to the number of elements.

Placing elements in a second pass over the input from right to left keeps the sort stable, which is what makes it usable as the inner pass of the radix sort below. The \inlinecode{key} function maps an element to an index in $[0, {\inlinecodemath{keys}})$, so records can be sorted by one field while carrying the rest.

\textbf{Time Complexity}: $O(n + k)$ for $n$ elements with $k$ possible keys.

\textbf{Space Complexity}: $O(n + k)$ auxiliary.

\begin{lstlisting}[language=C++]
template<typename It, typename KeyFn>
void counting_sort(It lo, It hi, int keys, KeyFn key) {
  if (hi - lo < 2) {
    return;
  }
  std::vector<int> count(keys);
  for (It it = lo; it != hi; ++it) {
    count[key(*it)]++;
  }
  for (int i = 1; i < keys; ++i) {
    count[i] += count[i - 1];
  }
  using T = typename std::iterator_traits<It>::value_type;
  std::vector<T> res(hi - lo);
  for (It it = hi; it != lo;) {
    --it;
    res[--count[key(*it)]] = *it;
  }
  std::copy(res.begin(), res.end(), lo);
}
\end{lstlisting}
Radix sort is used to sort integer elements with a constant number of bits in linear time. This implementation works on ranges pointing to any signed or unsigned integer primitive. Signed values are handled by sorting on a key that flips the sign bit, which maps the two's-complement order onto the unsigned order so the most negative value sorts first.

Digits are processed from least significant to most significant, and each counting-sort pass is stable. After one pass, the values are sorted by the processed digit; stability ensures that a later pass preserves the ordering already established by all less-significant digits. Inductively, after the final pass the values are sorted by the entire key.

This implementation uses one byte per digit, so digits can be extracted with shifts and masks and the counting table has $2^8 = 256$ entries.

\textbf{Time Complexity}: $O(n \cdot w)$ for $n$ integers of $w$ bits each.

\textbf{Space Complexity}: $O(n + 2^{b})$ auxiliary for a radix of $b$ bits, i.e. $O(n)$ for constant $b$.

\begin{lstlisting}[language=C++]
template<typename It>
void radix_sort(It lo, It hi) {
  if (hi - lo < 2) {
    return;
  }
  const int radix_bits = 8;
  const int radix_base = 1 << radix_bits;  // e.g. 2^8 = 256
  const int radix_mask = radix_base - 1;   // e.g. 2^8 - 1 = 0xFF
  using T = typename std::iterator_traits<It>::value_type;
  static_assert(std::is_integral_v<T> && !std::is_same_v<T, bool>);
  using U = typename std::make_unsigned_t<T>;
  const int num_bits = sizeof(T) * CHAR_BIT;
  // Sort on an unsigned key. For signed types, flipping the sign bit sends the most negative value
  // to 0, mapping the signed order onto the unsigned order; logical shifts then extract each digit.
  auto key = [](T x) -> U {
    U u = static_cast<U>(x);
    return std::is_signed<T>::value ? (u ^ (U{1} << (num_bits - 1))) : u;
  };
  std::vector<T> buf(hi - lo);
  for (int pos = 0; pos < num_bits; pos += radix_bits) {
    std::vector<int> count(radix_base);
    for (It it = lo; it != hi; ++it) {
      count[(key(*it) >> pos) & radix_mask]++;
    }
    std::vector<T *> bucket(radix_base);
    T *curr = buf.data();
    for (int i = 0; i < radix_base; curr += count[i++]) {
      bucket[i] = curr;
    }
    for (It it = lo; it != hi; ++it) {
      *bucket[(key(*it) >> pos) & radix_mask]++ = *it;
    }
    std::copy(buf.begin(), buf.end(), lo);
  }
}
\end{lstlisting}
Bitonic sort is a sorting network: a fixed sequence of compare-exchanges on fixed positions, each swapping its pair when out of order. The sequence depends only on how many elements there are, never on what they are, so no comparison branches on data. It sorts the first half ascending and the second half descending, leaving a bitonic sequence that rises and then falls, which a fixed pattern of halving compare-exchanges merges. Splitting each merge at the largest power of two below its length is what admits lengths that are not powers of two, which the textbook formulation requires.

Its $O(n \log^{2} n)$ comparators exceed the $O(n \log n)$ comparisons of a comparison sort, so run one at a time it loses to every sort above, as the benchmark below shows. What it buys instead is that each group of comparators is independent and may run at once, which is why networks are the standard choice across SIMD lanes and on a GPU. The same independence pays off at a small size fixed at compile time: unrolled into branchless minimum and maximum, a network sorts eight integers about an order of magnitude faster than a general sort. Driving that same comparator sequence from a lookup table forfeits it, since each comparator then costs an indexed load and an unpredictable branch.

\textbf{Time Complexity}: $O(n \log^{2} n)$ in all cases.

\textbf{Space Complexity}: $O(\log n)$ auxiliary stack space.

\begin{lstlisting}[language=C++]
template<typename It, typename Compare = std::less<>>
void bitonic_sort(It lo, It hi, Compare comp = Compare{}) {
  auto merge = [&](auto &&merge, It lo, int n, bool up) -> void {
    if (n < 2) {
      return;
    }
    int m = 1;
    while (2 * m < n) {
      m *= 2;
    }
    for (int i = 0; i < n - m; i++) {
      if (comp(*(lo + i + m), *(lo + i)) == up) {
        std::iter_swap(lo + i, lo + i + m);
      }
    }
    merge(merge, lo, m, up);
    merge(merge, lo + m, n - m, up);
  };
  auto rec = [&](auto &&rec, It lo, int n, bool up) -> void {
    if (n < 2) {
      return;
    }
    rec(rec, lo, n / 2, !up);
    rec(rec, lo + n / 2, n - n / 2, up);
    merge(merge, lo, n, up);
  };
  rec(rec, lo, static_cast<int>(hi - lo), true);
}
\end{lstlisting}
\Needspace*{4\baselineskip}
\textbf{Example Usage}\par
\begin{lstlisting}[language=C++]
#include <cassert>
#include <ctime>
#include <iomanip>
#include <iostream>
#include <random>
#include <utility>
#include <vector>
using namespace std;

template<typename It>
void print_range(It lo, It hi) {
  while (lo != hi) {
    cout << *lo++ << " ";
  }
  cout << endl;
}

int main() {
  {  // Can be used to sort arrays like std::sort().
    int a[]{32, 71, 12, 45, 26, 80, 53, 33};
    quicksort(begin(a), end(a));
    assert(is_sorted(begin(a), end(a)));
  }
  {  // STL containers work too.
    vector<int> a{32, 71, 12, 45, 26, 80, 53, 33};
    quicksort(a.begin(), a.end());
    assert(is_sorted(a.begin(), a.end()));
  }
  {  // Reverse iterators work as expected.
    vector<int> a{32, 71, 12, 45, 26, 80, 53, 33};
    heapsort(a.rbegin(), a.rend());
    assert(is_sorted(a.rbegin(), a.rend()));
  }
  {  // Insertion sort is adaptive: an already-sorted range is left untouched in O(n).
    vector<int> a{12, 26, 32, 33, 45, 53, 71, 80};
    insertion_sort(a.begin(), a.end());
    assert(is_sorted(a.begin(), a.end()));
  }
  {  // We can sort doubles just as well.
    vector<double> a{1.1, -5.0, 6.23, 4.123, 155.2};
    combsort(a.begin(), a.end());
    assert(is_sorted(a.begin(), a.end()));
  }
  {  // Shell sort moves values far in one step, so it is not adaptive but is nearly in-place.
    vector<int> a{32, 71, 12, 45, 26, 80, 53, 33};
    shellsort(a.begin(), a.end());
    assert(is_sorted(a.begin(), a.end()));
  }
  {  // Counting sort keys records by one field, and equal keys keep their input order.
    vector<pair<int, char>> a{{2, 'a'}, {0, 'b'}, {2, 'c'}, {1, 'd'}, {0, 'e'}};
    counting_sort(a.begin(), a.end(), 3, [](const pair<int, char> &x) { return x.first; });
    vector<pair<int, char>> expected{{0, 'b'}, {0, 'e'}, {1, 'd'}, {2, 'a'}, {2, 'c'}};
    assert(a == expected);
  }
  {  // radix_sort() handles signed integers (including negatives), unlike a plain counting sort.
    vector<int> a{32, -71, 12, -45, 26, -80, 53, 33};
    radix_sort(a.begin(), a.end());
    assert(is_sorted(a.begin(), a.end()));
  }
  {  // Bitonic sort is oblivious too, and takes any length rather than only a power of two.
    vector<int> a{32, 71, 12, 45, 26, 80, 53, 33, 19};
    bitonic_sort(a.begin(), a.end());
    assert(is_sorted(a.begin(), a.end()));
  }
  {  // Empty and singleton ranges are valid no-ops.
    vector<int> empty, single{42};
    quicksort(empty.begin(), empty.end());
    mergesort(empty.begin(), empty.end());
    heapsort(empty.begin(), empty.end());
    insertion_sort(empty.begin(), empty.end());
    combsort(empty.begin(), empty.end());
    radix_sort(empty.begin(), empty.end());
    mergesort(single.begin(), single.end());
    assert(empty.empty() && single[0] == 42);
  }

  // Stable sort.
  const vector<double> a{3.14, 1.41, 2.72, 4.67, 1.73, 1.32, 1.62, 2.58};
  {
    vector<double> v(a);
    cout << "mergesort() with default comparisons: ";
    mergesort(v.begin(), v.end());
    print_range(v.begin(), v.end());
  }
  {
    vector<double> v(a);
    cout << "mergesort() with integer comparisons: ";
    mergesort(v.begin(), v.end(), [](double i, double j) {
      return static_cast<int>(i) < static_cast<int>(j);
    });
    print_range(v.begin(), v.end());
  }
  cout << "------" << endl;

  mt19937 rng(1234567);  // Fixed seed for reproducibility.
  vector<int> data(5000000);
  const int maxval = 10000000;
  for (int &x : data) {
    x = static_cast<int>(rng() % maxval);  // limit magnitude for counting sort
  }
  cout << "Sorting five million integers..." << endl;
  cout.precision(3);
  auto benchmark = [&](const string &name, auto sort) {
    vector<int> v = data;
    clock_t start = clock();
    sort(v.begin(), v.end());
    double t = static_cast<double>(clock() - start) / CLOCKS_PER_SEC;
    cout << setw(18) << left << name + "(): " << fixed << t << "s" << endl;
    assert(is_sorted(v.begin(), v.end()));
  };
  benchmark("std::sort", [](auto lo, auto hi) { sort(lo, hi); });
  benchmark("quicksort", [](auto lo, auto hi) { quicksort(lo, hi); });
  benchmark("mergesort", [](auto lo, auto hi) { mergesort(lo, hi); });
  benchmark("heapsort", [](auto lo, auto hi) { heapsort(lo, hi); });
  benchmark("combsort", [](auto lo, auto hi) { combsort(lo, hi); });
  benchmark("shellsort", [](auto lo, auto hi) { shellsort(lo, hi); });
  benchmark("counting_sort", [](auto lo, auto hi) {
    counting_sort(lo, hi, maxval, [](int x) { return x; });
  });
  benchmark("radix_sort", [](auto lo, auto hi) { radix_sort(lo, hi); });
  benchmark("bitonic_sort", [](auto lo, auto hi) { bitonic_sort(lo, hi); });
  return 0;
}
\end{lstlisting}
\begin{exampleoutput}
mergesort() with default comparisons: 1.32 1.41 1.62 1.73 2.58 2.72 3.14 4.67
mergesort() with integer comparisons: 1.41 1.73 1.32 1.62 2.72 2.58 3.14 4.67
------
Sorting five million integers...
std::sort():      0.285s
quicksort():      0.325s
mergesort():      0.517s
heapsort():       0.509s
combsort():       0.469s
shellsort():      0.574s
counting_sort():  0.044s
radix_sort():     0.024s
bitonic_sort():   1.040s
\end{exampleoutput}
\subsection{Partitioning and Selection}
Partition a range around a pivot, or select the element with a given sorted rank without fully sorting the range. A comparator \inlinecode{comp} defines the ordering (defaults to \inlinecode{std::less\textless{}\textgreater{}}). Three-way partitioning rearranges the elements into consecutive groups that precede, are equivalent to, and follow the pivot. \inlinecode{std::partition()} can be used to place elements satisfying a unary predicate before those that do not, but forms only two groups and does not preserve relative order. Forming three groups can be done by first partitioning by \inlinecode{comp(x, pivot)}, then partitioning the remaining suffix by \inlinecode{!comp(pivot, x)}; using \inlinecode{std::stable\_partition()} for both calls preserves relative order within each group.

\inlinecode{partition\_three\_way()} implements the Dutch National Flag algorithm, forming the three groups in one pass with $O(1)$ auxiliary space. At every iteration, $[{\inlinecodemath{lo}}, {\inlinecodemath{lt}})$ contains elements preceding the pivot, $[{\inlinecodemath{lt}}, {\inlinecodemath{cur}})$ contains equivalent elements, $[{\inlinecodemath{cur}}, {\inlinecodemath{gt}})$ remains unclassified, and $[{\inlinecodemath{gt}}, {\inlinecodemath{hi}})$ contains elements following the pivot. The current element \inlinecode{*cur} is compared with the pivot. If it precedes the pivot, it is swapped with \inlinecode{*lt} and both \inlinecode{lt} and \inlinecode{cur} advance. If it is equivalent, only \inlinecode{cur} advances. Otherwise, it is swapped with \inlinecode{*--gt} without advancing \inlinecode{cur}, since the incoming element remains unclassified.

Quickselect repeatedly applies this partition and continues only into the group containing the desired rank. This is the same task as \inlinecode{std::nth\_element()}: after the call, \inlinecode{*nth} is the value that would appear there in comparator order, no earlier value follows it, and no later value precedes it. Choosing pivots uniformly at random gives expected linear time, while the three-way split avoids unnecessary work on duplicate-heavy inputs.

\begin{compactitem}
  \item \inlinecode{partition\_three\_way(lo, hi, pivot, comp = std::less\textless{}\textgreater{}())} rearranges $[{\inlinecodemath{lo}}, {\inlinecodemath{hi}})$ in-place and returns a pair of iterators (\inlinecode{mid1}, \inlinecode{mid2}). The three resulting ranges contain elements that precede, are equivalent to, and follow \inlinecode{pivot} according to \inlinecode{comp}, respectively. The resulting partition is not stable.
  \item \inlinecode{quickselect(lo, nth, hi, comp = std::less\textless{}\textgreater{}())} rearranges $[{\inlinecodemath{lo}}, {\inlinecodemath{hi}})$ in-place around the 0-based rank represented by iterator \inlinecode{nth} according to \inlinecode{comp}. This requires random-access iterators.
\end{compactitem}

Random pivots leave an adversary free to force $O(n^{2})$, which the median-of-medians rule removes. Split the range into groups of five, sort each, and pivot on the median of the group medians: at least three elements in half of the groups precede it, so every partition discards at least three tenths of the range. Finding that median is itself a selection, solved by recursing on the $n/5$ medians gathered to the front, and $T(n) = T(n/5) + T(7n/10) + O(n)$ is linear because the two fractions sum to less than one, which five is the smallest group size to achieve. Each group is sorted directly, since at five elements \inlinecode{std::sort} is an inlined insertion sort. This guarantee is not free, but on random input the cost over a random pivot is modest.

\begin{compactitem}
  \item \inlinecode{quickselect\_linear(lo, nth, hi, comp = std::less\textless{}\textgreater{}())} does the same in linear time even in the worst case, choosing each pivot by the median-of-medians rule instead of at random.
\end{compactitem}

\Needspace*{3\baselineskip}
\textbf{Time Complexity}:
\begin{compactitem}
  \item $O(n)$ per call to \inlinecode{partition\_three\_way()}, where $n$ is the range length.
  \item $O(n)$ expected per call to \inlinecode{quickselect()}, and $O(n^{2})$ in the worst case.
  \item $O(n)$ per call to \inlinecode{quickselect\_linear()} in all cases.
\end{compactitem}

\Needspace*{3\baselineskip}
\textbf{Space Complexity}:
\begin{compactitem}
  \item $O(1)$ auxiliary for \inlinecode{partition\_three\_way()} and \inlinecode{quickselect()}.
  \item $O(\log n)$ auxiliary stack space for \inlinecode{quickselect\_linear()}.
\end{compactitem}

\begin{lstlisting}[language=C++]
#include <algorithm>
#include <chrono>
#include <functional>
#include <iterator>
#include <random>
#include <utility>

template<typename It, typename T, typename Compare = std::less<>>
std::pair<It, It> partition_three_way(It lo, It hi, const T &pivot, Compare comp = Compare{}) {
  It lt = lo, cur = lo, gt = hi;
  while (cur != gt) {
    if (comp(*cur, pivot)) {
      std::iter_swap(lt++, cur++);
    } else if (comp(pivot, *cur)) {
      std::iter_swap(cur, --gt);
    } else {
      ++cur;
    }
  }
  return {lt, gt};
}

template<typename It, typename Compare = std::less<>>
void quickselect(It lo, It nth, It hi, Compare comp = Compare{}) {
  using T = typename std::iterator_traits<It>::value_type;
  static std::mt19937 rng(std::chrono::steady_clock::now().time_since_epoch().count());
  while (hi - lo > 1) {
    std::uniform_int_distribution<int> dist(0, hi - lo - 1);
    T pivot = *(lo + dist(rng));
    auto [lt, gt] = partition_three_way(lo, hi, pivot, comp);
    if (nth < lt) {
      hi = lt;
    } else if (nth >= gt) {
      lo = gt;
    } else {
      return;
    }
  }
}

template<typename It, typename Compare = std::less<>>
void quickselect_linear(It lo, It nth, It hi, Compare comp = Compare{}) {
  while (hi - lo > 5) {
    int groups = 0;
    for (It g = lo; g < hi; g += 5, groups++) {
      It ge = hi - g < 5 ? hi : g + 5;
      std::sort(g, ge, comp);  // At most 5 elements, so this is an O(1) step.
      std::iter_swap(lo + groups, g + (ge - g) / 2);
    }
    It mid = lo + groups / 2;
    quickselect_linear(lo, mid, lo + groups, comp);
    using T = typename std::iterator_traits<It>::value_type;
    T pivot = *mid;
    std::pair<It, It> parts = partition_three_way(lo, hi, pivot, comp);
    if (nth < parts.first) {
      hi = parts.first;
    } else if (nth < parts.second) {
      return;
    } else {
      lo = parts.second;
    }
  }
  std::sort(lo, hi, comp);
}
\end{lstlisting}
\Needspace*{4\baselineskip}
\textbf{Example Usage}\par
\begin{lstlisting}[language=C++]
#include <cassert>
#include <vector>
using namespace std;

int main() {
  vector<int> values{2, 0, 1, 2, 1, 0, 1};
  // On a range containing only 0, 1, and 2, partitioning around 1 sorts the range.
  partition_three_way(values.begin(), values.end(), 1);
  assert(is_sorted(values.begin(), values.end()));

  vector<int> b{4, 2, 5, 3, 3, 1};
  auto [mid1, mid2] = partition_three_way(b.begin(), b.end(), 3);
  for (auto it = b.begin(); it != mid1; ++it) {
    assert(*it < 3);
  }
  for (auto it = mid1; it != mid2; ++it) {
    assert(*it == 3);
  }
  for (auto it = mid2; it != b.end(); ++it) {
    assert(*it > 3);
  }

  vector<int> a{5, 6, 4, 3, 2, 6, 7, 9, 3};
  int n = static_cast<int>(a.size());
  quickselect(a.begin(), a.begin() + n / 2, a.end());
  assert(a[n / 2] == 5);
  // Values left of the median are <=, and values right are >= (the exact order within each side is
  // randomized, since the pivot is chosen at random).
  for (int i = 0; i < n / 2; i++) {
    assert(a[i] <= a[n / 2]);
  }
  for (int i = n / 2 + 1; i < n; i++) {
    assert(a[i] >= a[n / 2]);
  }

  // The median-of-medians variant gives the same answer with a worst case that is still linear.
  vector<int> c{5, 6, 4, 3, 2, 6, 7, 9, 3};
  quickselect_linear(c.begin(), c.begin() + n / 2, c.end());
  assert(c[n / 2] == 5);

  vector<int> desc{1, 4, 2, 3};
  quickselect(desc.begin(), desc.begin() + 1, desc.end(), greater<int>());
  assert(desc[1] == 3);
  return 0;
}
\end{lstlisting}
\subsection{Array Rotation}
Rotates the half-open iterator range $[{\inlinecodemath{lo}}, {\inlinecodemath{hi}})$ left around the split point \inlinecode{mid}, where ${\inlinecodemath{lo}} \leq {\inlinecodemath{mid}} \leq {\inlinecodemath{hi}}$. Writing the adjacent subranges $[{\inlinecodemath{lo}}, {\inlinecodemath{mid}})$ and $[{\inlinecodemath{mid}}, {\inlinecodemath{hi}})$ as \inlinecode{A} and \inlinecode{B}, respectively, the operation transforms \inlinecode{A B} into \inlinecode{B A} while preserving the relative order within each subrange. The standard library function \inlinecode{std::rotate(lo, mid, hi)} performs the same rearrangement.

All three versions below operate in place, using different algorithms and iterator capabilities.

\begin{compactitem}
  \item \inlinecode{rotate1(lo, mid, hi)} requires ForwardIterators and repeatedly swaps the next elements of the two unfinished subranges. When the right iterator reaches \inlinecode{hi}, it wraps to the current \inlinecode{mid}; when the left iterator reaches \inlinecode{mid}, that boundary advances to the right iterator.
  \item \inlinecode{rotate2(lo, mid, hi)} requires BidirectionalIterators, applying a trick with three reversals. Writing the input subranges as \inlinecode{A B}, two initial reversals yield \inlinecode{reverse(A) reverse(B)}; then the whole range is reversed to yield \inlinecode{B A}, restoring the order within each subrange.
  \item \inlinecode{rotate3(lo, mid, hi)} requires random-access iterators, applying a juggling algorithm which first divides the range into \inlinecode{gcd(hi - lo, mid - lo)} sets and then rotates the corresponding elements in each set. The method follows the permutation that sends each position to its rotated destination. These form that many disjoint cycles, so walking each cycle once moves every element to its final position.
\end{compactitem}

\textbf{Time Complexity}: $O(n)$ per call to all versions, where $n$ is the distance between \inlinecode{lo} and \inlinecode{hi}.

\textbf{Space Complexity}: $O(1)$ auxiliary for all versions.

\begin{lstlisting}[language=C++]
#include <algorithm>
#include <numeric>

template<typename It>
void rotate1(It lo, It mid, It hi) {
  if (lo == mid || mid == hi) {
    return;
  }
  It next = mid;
  while (lo != next) {
    std::iter_swap(lo++, next++);
    if (next == hi) {
      next = mid;
    } else if (lo == mid) {
      mid = next;
    }
  }
}

template<typename It>
void rotate2(It lo, It mid, It hi) {
  if (lo == mid || mid == hi) {
    return;
  }
  std::reverse(lo, mid);
  std::reverse(mid, hi);
  std::reverse(lo, hi);
}

template<typename It>
void rotate3(It lo, It mid, It hi) {
  if (lo == mid || mid == hi) {
    return;
  }
  int n = static_cast<int>(hi - lo);
  int jump = static_cast<int>(mid - lo);
  int g = std::gcd(jump, n), cycle = n / g;
  for (int i = 0; i < g; i++) {
    int curr = i, next;
    for (int j = 0; j < cycle - 1; j++) {
      next = curr + jump;
      if (next >= n) {
        next -= n;
      }
      std::iter_swap(lo + curr, lo + next);
      curr = next;
    }
  }
}
\end{lstlisting}
\Needspace*{4\baselineskip}
\textbf{Example Usage}\par
\begin{lstlisting}[language=C++]
#include <algorithm>
#include <cassert>
#include <vector>
using namespace std;

int main() {
  vector<int> a0(10000), a1, a2, a3;
  iota(a0.begin(), a0.end(), 0);
  a1 = a2 = a3 = a0;
  int mid = 5678;
  std::rotate(a0.begin(), a0.begin() + mid, a0.end());
  rotate1(a1.begin(), a1.begin() + mid, a1.end());
  rotate2(a2.begin(), a2.begin() + mid, a2.end());
  rotate3(a3.begin(), a3.begin() + mid, a3.end());
  assert(a0 == a1 && a0 == a2 && a0 == a3);

  vector<int> no_op{1, 2, 3};
  rotate1(no_op.begin(), no_op.begin(), no_op.end());
  assert((no_op == vector<int>{1, 2, 3}));
  rotate2(no_op.begin(), no_op.end(), no_op.end());
  assert((no_op == vector<int>{1, 2, 3}));
  rotate3(no_op.begin(), no_op.begin(), no_op.end());
  assert((no_op == vector<int>{1, 2, 3}));

  vector<int> a{2, 4, 2, 0, 5, 10, 7, 3, 7, 1};

  // Insertion sort.
  for (auto i = a.begin(); i != a.end(); ++i) {
    rotate1(upper_bound(a.begin(), i, *i), i, i + 1);
  }
  assert((a == vector<int>{0, 1, 2, 2, 3, 4, 5, 7, 7, 10}));

  // Simple rotation to the left.
  rotate2(a.begin(), a.begin() + 1, a.end());
  assert((a == vector<int>{1, 2, 2, 3, 4, 5, 7, 7, 10, 0}));

  // Simple rotation to the right.
  rotate3(a.rbegin(), a.rbegin() + 1, a.rend());
  assert((a == vector<int>{0, 1, 2, 2, 3, 4, 5, 7, 7, 10}));
  return 0;
}
\end{lstlisting}
\subsection{Coordinate Compression}
Given a range of $n$ numerical elements, reassign each element to an integer in the domain $[0, k)$, where $k$ is the number of distinct elements in the original range, while preserving the initial relative ordering of elements. That is, if $a$ is the array of original values and $b$ is the array of compressed values, then every pair of indices $i, j$ in $[0, n)$ shall satisfy $a[i] < a[j]$ if and only if $b[i] < b[j]$.

The two functions below take a ForwardIterator range, rewrite it in place, and discard the mapping. The comparator \inlinecode{comp} defines the value ordering: \inlinecode{comp(a, b)} is true when \inlinecode{a} precedes \inlinecode{b}.

\begin{compactitem}
  \item \inlinecode{compress1(lo, hi, comp = std::less\textless{}\textgreater{}())} performs the compression by sorting the array, removing duplicates, and binary searching for the position of each original value.
  \item \inlinecode{compress2(lo, hi, comp = std::less\textless{}\textgreater{}())} achieves the same result by inserting all values into an \inlinecode{std::map}, which automatically removes duplicate values and supports efficient lookups of the compressed values.
\end{compactitem}

\textbf{Time Complexity}: $O(n \log n)$ per call to \inlinecode{compress1(lo, hi)} and \inlinecode{compress2(lo, hi)}, where $n$ is the distance between \inlinecode{lo} and \inlinecode{hi}.

\textbf{Space Complexity}: $O(n)$ auxiliary for \inlinecode{compress1()} and \inlinecode{compress2()}.

\begin{lstlisting}[language=C++]
#include <algorithm>
#include <cassert>
#include <functional>
#include <iterator>
#include <map>
#include <vector>

template<typename It, typename Compare = std::less<>>
void compress1(It lo, It hi, Compare comp = Compare{}) {
  using T = typename std::iterator_traits<It>::value_type;
  std::vector<T> v(lo, hi);
  std::sort(v.begin(), v.end(), comp);
  v.erase(
      std::unique(
          v.begin(), v.end(), [&](const T &a, const T &b) { return !comp(a, b) && !comp(b, a); }
      ),
      v.end()
  );
  for (It it = lo; it != hi; ++it) {
    *it = static_cast<int>(std::lower_bound(v.begin(), v.end(), *it, comp) - v.begin());
  }
}

template<typename It, typename Compare = std::less<>>
void compress2(It lo, It hi, Compare comp = Compare{}) {
  using T = typename std::iterator_traits<It>::value_type;
  std::map<T, int, Compare> m(comp);
  for (It it = lo; it != hi; ++it) {
    m[*it] = 0;
  }
  int id = 0;
  for (auto &[key, val] : m) {
    val = id++;
  }
  for (It it = lo; it != hi; ++it) {
    *it = m[*it];
  }
}
\end{lstlisting}
\inlinecode{Compressor} retains the sorted table of distinct values so that arbitrary values can be mapped to and from compressed ranks long after construction, such as for offline queries arriving separately from the array being compressed. It uses \inlinecode{std::less\textless{}T\textgreater{}} by default; to customize the ordering, instantiate \inlinecode{Compressor\textless{}T, Compare\textgreater{}} and pass the comparator to either constructor.

\begin{compactitem}
  \item \inlinecode{Compressor\textless{}T\textgreater{}()} constructs an empty compressor. Register values with \inlinecode{add(x)}, then call \inlinecode{build()} once before the first query.
  \item \inlinecode{Compressor\textless{}T\textgreater{}(lo, hi)} constructs a compressor from the half-open iterator range $[{\inlinecodemath{lo}}, {\inlinecodemath{hi}})$.
  \item \inlinecode{size()} returns the number of distinct registered values $k$.
  \item \inlinecode{value(r)} returns the original value with rank \inlinecode{r}, inverting \inlinecode{rank()}.
  \item \inlinecode{contains(x)} returns whether \inlinecode{x} is a registered value.
  \item \inlinecode{rank(x)} returns the compressed value (rank) of \inlinecode{x} in $[0, k)$. \inlinecode{x} must have been registered.
\end{compactitem}

\Needspace*{3\baselineskip}
\textbf{Time Complexity}:
\begin{compactitem}
  \item $O(n \log n)$ per call to \inlinecode{Compressor(lo, hi)}, where $n$ is the range length.
  \item $O(1)$ amortized per call to \inlinecode{add()}, and $O(m \log m)$ per call to \inlinecode{build()}, where $m$ is the total number of values registered.
  \item $O(\log k)$ per call to \inlinecode{rank()} and \inlinecode{contains()}, where $k$ is the number of distinct values.
  \item $O(1)$ per call to \inlinecode{size()} and \inlinecode{value()}.
\end{compactitem}

\textbf{Space Complexity}: $O(k)$ storage, where $k$ is the number of distinct registered values.

\begin{lstlisting}[language=C++]
template<typename T, typename Compare = std::less<T>>
class Compressor {
  Compare comp;
  std::vector<T> v;

 public:
  explicit Compressor(Compare comp = Compare{}) : comp(std::move(comp)) {}

  template<typename It>
  Compressor(It lo, It hi, Compare comp = Compare{}) : comp(std::move(comp)), v(lo, hi) {
    build();
  }

  void build() {
    std::sort(v.begin(), v.end(), comp);
    v.erase(
        std::unique(
            v.begin(), v.end(), [&](const T &a, const T &b) { return !comp(a, b) && !comp(b, a); }
        ),
        v.end()
    );
  }

  void add(const T &x) { v.push_back(x); }
  int size() const { return static_cast<int>(v.size()); }
  const T &value(int r) const { return v[r]; }
  bool contains(const T &x) const { return std::binary_search(v.begin(), v.end(), x, comp); }

  int rank(const T &x) const {
    int r = static_cast<int>(std::lower_bound(v.begin(), v.end(), x, comp) - v.begin());
    assert(r < size() && !comp(x, v[r]));  // x must be registered.
    return r;
  }
};
\end{lstlisting}
\Needspace*{4\baselineskip}
\textbf{Example Usage}\par
\begin{lstlisting}[language=C++]
#include <cassert>
using namespace std;

int main() {
  {
    vector<int> a{1, 30, 30, 7, 9, 8, 99, 99};
    compress1(a.begin(), a.end());
    assert((a == vector<int>{0, 4, 4, 1, 3, 2, 5, 5}));
  }
  {
    vector<int> a{1, 30, 30, 7, 9, 8, 99, 99};
    compress2(a.begin(), a.end());
    assert((a == vector<int>{0, 4, 4, 1, 3, 2, 5, 5}));
  }
  {  // Non-integral types work too, as long as ints can be assigned to them.
    vector<double> a{0.5, -1.0, 3, -1.0, 20, 0.5};
    compress1(a.begin(), a.end());
    assert((a == vector<double>{1, 0, 2, 0, 3, 1}));
  }
  {
    vector<int> a{10, 5, 30, 5, 20};
    Compressor<int> cc(a.begin(), a.end());
    assert(cc.size() == 4);
    assert(cc.rank(5) == 0);
    assert(cc.rank(30) == 3);
    assert(cc.value(2) == 20);
    assert(cc.contains(10) && !cc.contains(7));
  }
  {  // Register values incrementally, then build once before querying.
    Compressor<double> cc;
    cc.add(0.5);
    cc.add(-1.0);
    cc.add(0.5);
    cc.build();
    assert(cc.size() == 2);
    assert(cc.rank(0.5) == 1);
    assert(cc.value(0) == -1.0);
  }
  {
    vector<int> a{10, 30, 20, 10};
    Compressor<int, greater<int>> cc(a.begin(), a.end());
    assert(cc.rank(30) == 0 && cc.rank(10) == 2);
    assert(cc.value(1) == 20);
    compress1(a.begin(), a.end(), greater<int>());
    assert((a == vector<int>{2, 0, 1, 2}));
  }
  return 0;
}
\end{lstlisting}
\subsection{Counting Inversions}
The number of inversions for a sequence $a$ is the number of ordered pairs $(i, j)$ such that $i < j$ and $a[i] > a[j]$. This is roughly how ``close'' an array is to being sorted, but is \textit{not} the minimum number of swaps required to sort the array. If the array is sorted, then the inversion count is $0$. If the array is sorted in decreasing order, then the inversion count is maximal. The following two functions are each techniques to efficiently count inversions. In the merge sort approach, whenever the merge step emits an element from the right half, that element jumps ahead of every unmerged left-half element, and exactly that many inversions are added to the count.

\begin{compactitem}
  \item \inlinecode{inversions(lo, hi, comp = std::less\textless{}\textgreater{})} uses merge sort to return the number of inversions given two random-access iterators as a range $[{\inlinecodemath{lo}}, {\inlinecodemath{hi}})$. The input range will be sorted after the function call. Optionally, a comparison function object specifying a strict weak ordering may be specified to replace the default \inlinecode{operator\textless{}}.
  \item \inlinecode{inversions(a)} uses coordinate compression and a Fenwick tree to return the number of inversions in an integer vector without modifying it.
\end{compactitem}

\Needspace*{3\baselineskip}
\textbf{Time Complexity}:
\begin{compactitem}
  \item $O(n \log n)$ per call to \inlinecode{inversions(lo, hi)}, where $n$ is the distance between \inlinecode{lo} and \inlinecode{hi}.
  \item $O(n \log n)$ per call to \inlinecode{inversions(a)}.
\end{compactitem}

\Needspace*{3\baselineskip}
\textbf{Space Complexity}:
\begin{compactitem}
  \item $O(\log n)$ auxiliary stack space and $O(n)$ auxiliary heap space for \inlinecode{inversions(lo, hi)}.
  \item $O(n)$ auxiliary heap space for \inlinecode{inversions(a)}.
\end{compactitem}

\begin{lstlisting}[language=C++]
#include <algorithm>
#include <cstdint>
#include <functional>
#include <iterator>
#include <vector>

template<typename It, typename Compare = std::less<>>
int64_t inversions(It lo, It hi, Compare comp = Compare{}) {
  int n = static_cast<int>(hi - lo);
  if (n < 2) {
    return 0;
  }
  It mid = lo + (n - 1) / 2, a = lo, c = mid + 1;
  int64_t res = inversions(lo, mid + 1, comp) + inversions(mid + 1, hi, comp);
  using T = typename std::iterator_traits<It>::value_type;
  std::vector<T> merged;
  merged.reserve(n);
  while (a <= mid && c < hi) {
    if (comp(*c, *a)) {
      merged.push_back(*(c++));
      res += (mid - a) + 1;
    } else {
      merged.push_back(*(a++));
    }
  }
  if (a > mid) {
    for (It it = c; it != hi; ++it) {
      merged.push_back(*it);
    }
  } else {
    for (It it = a; it <= mid; ++it) {
      merged.push_back(*it);
    }
  }
  for (It it = lo; it != hi; ++it) {
    *it = merged[it - lo];
  }
  return res;
}

int64_t inversions(const std::vector<int> &a) {
  int n = static_cast<int>(a.size());
  std::vector<int> values(a);
  std::sort(values.begin(), values.end());
  values.resize(std::unique(values.begin(), values.end()) - values.begin());
  std::vector<int> bit(values.size() + 1);
  int64_t res = 0;
  for (int i = n - 1; i >= 0; i--) {
    int id =
        static_cast<int>(std::lower_bound(values.begin(), values.end(), a[i]) - values.begin()) + 1;
    for (int j = id - 1; j > 0; j -= j & -j) {
      res += bit[j];
    }
    for (int j = id; j < static_cast<int>(bit.size()); j += j & -j) {
      bit[j]++;
    }
  }
  return res;
}
\end{lstlisting}
\Needspace*{4\baselineskip}
\textbf{Example Usage}\par
\begin{lstlisting}[language=C++]
#include <cassert>
using namespace std;

int main() {
  {
    vector<int> a{6, 9, 1, 14, 8, 12, 3, 2};
    assert(inversions(a.begin(), a.end()) == 16);
    assert(is_sorted(a.begin(), a.end()));
  }
  {
    vector<int> a{6, 9, 1, 14, 8, 12, 3, 2};
    assert(inversions(a) == 16);
    assert((a == vector<int>{6, 9, 1, 14, 8, 12, 3, 2}));
  }
  {
    vector<int> a{2, 2, 1};
    assert(inversions(a) == 2);  // Equal elements are not inversions.
  }
  return 0;
}
\end{lstlisting}

\section{\texorpdfstring{1D Scan and Window Techniques}{1.2 1D Scan and Window Techniques}}
\setcounter{section}{2}
\setcounter{subsection}{0}
\subsection{Prefix Sums}
Precomputes prefix sums so range-sum queries can be answered in constant time. This is one of the most common array preprocessing tools, and also appears as a building block for subarray sums, difference arrays, and two-dimensional grids. Each table entry stores the sum of all elements before it, so any range sum is the difference of two entries; in two dimensions, rectangle sums combine four entries by inclusion-exclusion.

Prefix sums can also be combined with a frequency table to count subarrays having a target sum. When the current prefix sum is $s$, every earlier prefix sum equal to $s - t$ identifies a subarray with sum $t$. Recording frequencies rather than only whether a prefix exists counts duplicate sums correctly.

\begin{compactitem}
  \item \inlinecode{prefix\_sums(a)} returns the prefix sum array \inlinecode{pref} of length $n + 1$, with \inlinecode{pref[0] = 0} and \inlinecode{pref[i + 1]} equal to the sum of \inlinecode{a[0]} through \inlinecode{a[i]}.
  \item \inlinecode{range\_sum(pref, lo, hi)} returns the sum of range $[{\inlinecodemath{lo}}, {\inlinecodemath{hi}}]$.
  \item \inlinecode{prefix\_sums\_2d(a)} returns a two-dimensional prefix sum table for matrix \inlinecode{a}.
  \item \inlinecode{rect\_sum(pref, r1, c1, r2, c2)} returns the sum of the rectangle with rows $[{\inlinecodemath{r1}}, {\inlinecodemath{r2}}]$ and columns $[{\inlinecodemath{c1}}, {\inlinecodemath{c2}}]$.
  \item \inlinecode{count\_subarrays\_with\_sum(a, target)} returns the number of contiguous subarrays whose sum is \inlinecode{target}.
\end{compactitem}

Overflow warning: All prefix sums, target differences, inclusion-exclusion results, and returned counts must fit in \inlinecode{int64\_t}.

\Needspace*{3\baselineskip}
\textbf{Time Complexity}:
\begin{compactitem}
  \item $O(n)$ per call to \inlinecode{prefix\_sums(a)}, where $n$ is the array size.
  \item $O(R \cdot C)$ per call to \inlinecode{prefix\_sums\_2d(a)}, where $R$ and $C$ are the number of rows and columns of \inlinecode{a}, respectively.
  \item $O(1)$ per range or rectangle query.
  \item $O(n)$ expected per call to \inlinecode{count\_subarrays\_with\_sum(a, target)}.
\end{compactitem}

\Needspace*{3\baselineskip}
\textbf{Space Complexity}:
\begin{compactitem}
  \item $O(n)$ for the array returned by \inlinecode{prefix\_sums(a)}.
  \item $O(R \cdot C)$ for the table returned by \inlinecode{prefix\_sums\_2d(a)}.
  \item $O(1)$ auxiliary for \inlinecode{range\_sum()} and \inlinecode{rect\_sum()}.
  \item $O(n)$ auxiliary for \inlinecode{count\_subarrays\_with\_sum(a, target)}.
\end{compactitem}

\begin{lstlisting}[language=C++]
#include <cassert>
#include <cstdint>
#include <unordered_map>
#include <vector>

std::vector<int64_t> prefix_sums(const std::vector<int> &a) {
  std::vector<int64_t> pref(a.size() + 1);
  for (int i = 0; i < static_cast<int>(a.size()); i++) {
    pref[i + 1] = pref[i] + a[i];
  }
  return pref;
}

int64_t range_sum(const std::vector<int64_t> &pref, int lo, int hi) {
  assert(0 <= lo && lo <= hi && hi < static_cast<int>(pref.size()) - 1);
  return pref[hi + 1] - pref[lo];
}

std::vector<std::vector<int64_t>> prefix_sums_2d(const std::vector<std::vector<int>> &a) {
  int rows = static_cast<int>(a.size());
  int cols = a.empty() ? 0 : static_cast<int>(a[0].size());
  std::vector<std::vector<int64_t>> pref(rows + 1, std::vector<int64_t>(cols + 1));
  for (int i = 0; i < rows; i++) {
    for (int j = 0; j < cols; j++) {
      pref[i + 1][j + 1] = a[i][j] + pref[i][j + 1] + pref[i + 1][j] - pref[i][j];
    }
  }
  return pref;
}

int64_t rect_sum(const std::vector<std::vector<int64_t>> &pref, int r1, int c1, int r2, int c2) {
  assert(
      !pref.empty() && !pref[0].empty() && 0 <= r1 && r1 <= r2 &&
      r2 < static_cast<int>(pref.size()) - 1 && 0 <= c1 && c1 <= c2 &&
      c2 < static_cast<int>(pref[0].size()) - 1
  );
  return pref[r2 + 1][c2 + 1] - pref[r1][c2 + 1] - pref[r2 + 1][c1] + pref[r1][c1];
}

int64_t count_subarrays_with_sum(const std::vector<int> &a, int64_t target) {
  std::unordered_map<int64_t, int64_t> count{{0, 1}};
  int64_t sum = 0, result = 0;
  for (int x : a) {
    sum += x;
    int64_t needed = sum - target;
    if (auto it = count.find(needed); it != count.end()) {
      result += it->second;
    }
    count[sum]++;
  }
  return result;
}
\end{lstlisting}
\Needspace*{4\baselineskip}
\textbf{Example Usage}\par
\begin{lstlisting}[language=C++]
using namespace std;

int main() {
  vector<int> a{3, -1, 4, 1, 5};
  auto pref = prefix_sums(a);
  assert(range_sum(pref, 0, 4) == 12);  // Whole array.
  assert(range_sum(pref, 1, 3) == 4);   // -1 + 4 + 1.
  assert(range_sum(pref, 2, 2) == 4);   // Single element.
  assert(count_subarrays_with_sum(a, 4) == 2);
  assert(count_subarrays_with_sum(vector<int>{0, 0, 0}, 0) == 6);

  vector<vector<int>> grid{
      {1, 2, 3},
      {4, 5, 6},
  };
  auto pre2 = prefix_sums_2d(grid);
  assert(rect_sum(pre2, 0, 1, 1, 2) == 16);  // Rows 0-1 and columns 1-2.
  assert(rect_sum(pre2, 0, 0, 1, 2) == 21);  // Whole grid.
  assert(rect_sum(pre2, 1, 1, 1, 1) == 5);   // Single cell.
  return 0;
}
\end{lstlisting}
\subsection{Difference Array}
Applies many range-add updates to an array in linear total time using a difference array. Instead of immediately adding \inlinecode{delta} to every element of range $[{\inlinecodemath{lo}}, {\inlinecodemath{hi}}]$, we can add \inlinecode{delta} at \inlinecode{diff[lo]} and subtract it at \inlinecode{diff[hi + 1]}; a final prefix sum reconstructs the updated values. This idea extends to two dimensions: a rectangle update touches just the four corner cells of a difference grid by inclusion-exclusion, and a final two-dimensional prefix sum reconstructs the grid.

\begin{compactitem}
  \item \inlinecode{DifferenceArray(n)} constructs an initially zero array of size \inlinecode{n}.
  \item \inlinecode{add(lo, hi, delta)} adds \inlinecode{delta} to every position in range $[{\inlinecodemath{lo}}, {\inlinecodemath{hi}}]$.
  \item \inlinecode{build()} returns the final array after all range updates.
  \item \inlinecode{DifferenceArray2D(rows, cols)} constructs an initially zero grid with \inlinecode{rows} rows and \inlinecode{cols} columns.
  \item \inlinecode{add(r1, c1, r2, c2, delta)} adds \inlinecode{delta} to every cell of the rectangle with rows $[{\inlinecodemath{r1}}, {\inlinecodemath{r2}}]$ and columns $[{\inlinecodemath{c1}}, {\inlinecodemath{c2}}]$.
  \item \inlinecode{build()} returns the final grid after all rectangle updates.
\end{compactitem}

Overflow warning: All stored differences and reconstructed values must fit in \inlinecode{int64\_t}.

\Needspace*{3\baselineskip}
\textbf{Time Complexity}:
\begin{compactitem}
  \item $O(1)$ per call to \inlinecode{add()} of either version.
  \item $O(n)$ per call to \inlinecode{build()} of \inlinecode{DifferenceArray}, where $n$ is the array size.
  \item $O(R \cdot C)$ per call to \inlinecode{build()} of \inlinecode{DifferenceArray2D}, where $R$ and $C$ are the number of rows and columns, respectively.
\end{compactitem}

\Needspace*{3\baselineskip}
\textbf{Space Complexity}:
\begin{compactitem}
  \item $O(n)$ for storage and $O(n)$ for the array returned by \inlinecode{DifferenceArray::build()}.
  \item $O(R \cdot C)$ for storage and $O(R \cdot C)$ for the grid returned by \inlinecode{DifferenceArray2D::build()}.
\end{compactitem}

\begin{lstlisting}[language=C++]
#include <cassert>
#include <cstdint>
#include <vector>

class DifferenceArray {
  std::vector<int64_t> diff;

 public:
  explicit DifferenceArray(int n) {
    assert(n >= 0);
    diff.assign(n + 1, 0);
  }

  void add(int lo, int hi, int64_t delta) {
    assert(0 <= lo && lo <= hi && hi + 1 < static_cast<int>(diff.size()));
    diff[lo] += delta;  // Overflow warning.
    diff[hi + 1] -= delta;
  }

  std::vector<int64_t> build() const {
    std::vector<int64_t> res(static_cast<int>(diff.size()) - 1);
    int64_t cur = 0;
    for (int i = 0; i < static_cast<int>(res.size()); i++) {
      cur += diff[i];  // Overflow warning.
      res[i] = cur;
    }
    return res;
  }
};

class DifferenceArray2D {
  std::vector<std::vector<int64_t>> diff;

 public:
  DifferenceArray2D(int rows, int cols) {
    assert(rows >= 0 && cols >= 0);
    diff.assign(rows + 1, std::vector<int64_t>(cols + 1));
  }

  void add(int r1, int c1, int r2, int c2, int64_t delta) {
    int rows = static_cast<int>(diff.size()) - 1;
    int cols = static_cast<int>(diff[0].size()) - 1;
    assert(0 <= r1 && r1 <= r2 && r2 < rows);
    assert(0 <= c1 && c1 <= c2 && c2 < cols);
    diff[r1][c1] += delta;  // Overflow warning.
    diff[r1][c2 + 1] -= delta;
    diff[r2 + 1][c1] -= delta;
    diff[r2 + 1][c2 + 1] += delta;
  }

  std::vector<std::vector<int64_t>> build() const {
    int rows = static_cast<int>(diff.size()) - 1;
    int cols = static_cast<int>(diff[0].size()) - 1;
    std::vector<std::vector<int64_t>> res(rows, std::vector<int64_t>(cols));
    for (int i = 0; i < rows; i++) {
      for (int j = 0; j < cols; j++) {
        int64_t up = i > 0 ? res[i - 1][j] : 0;
        int64_t left = j > 0 ? res[i][j - 1] : 0;
        int64_t diag = i > 0 && j > 0 ? res[i - 1][j - 1] : 0;
        res[i][j] = diff[i][j] + up + left - diag;  // Overflow warning.
      }
    }
    return res;
  }
};
\end{lstlisting}
\Needspace*{4\baselineskip}
\textbf{Example Usage}\par
\begin{lstlisting}[language=C++]
#include <cassert>
using namespace std;

int main() {
  DifferenceArray d(5);
  d.add(0, 2, 2);   // Add 2 to indices 0-2.
  d.add(1, 4, 4);   // Add 4 to indices 1-4.
  d.add(2, 3, -1);  // Subtract 1 from indices 2-3.
  d.add(4, 4, 3);   // Single index update.
  assert((d.build() == vector<int64_t>{2, 6, 5, 3, 7}));

  DifferenceArray2D d2(3, 4);
  d2.add(0, 0, 1, 2, 1);  // Add 1 to rows 0-1 and columns 0-2.
  d2.add(1, 1, 2, 3, 2);  // Add 2 to rows 1-2 and columns 1-3.
  d2.add(2, 0, 2, 0, 5);  // Single cell update.
  vector<vector<int64_t>> expected2{
      {1, 1, 1, 0},
      {1, 3, 3, 2},
      {5, 2, 2, 2},
  };
  assert(d2.build() == expected2);
  return 0;
}
\end{lstlisting}
\subsection{Two Pointers and Sliding Window}
Uses two moving indices to avoid trying all pairs of positions. This technique has two common forms: opposite-end pointers over sorted arrays, and a sliding window over contiguous subarrays. For sorted 2-sum, if the current sum is too small, pairing the left value with any smaller right value also fails, so the left pointer can safely advance; if the sum is too large, the symmetric argument lets the right pointer retreat. Each step therefore discards an entire row or column of candidate pairs.

A sliding window applies the same idea when extending and shrinking a range changes its state monotonically. With nonnegative values, extending a window cannot decrease its sum and removing its leftmost value cannot increase it, so no shorter feasible window is skipped. For a bound such as at most $k$ distinct values, the left pointer advances only until the window becomes valid again. Since neither endpoint moves backward, each element enters and leaves the window at most once.

For fixed-length windows, a monotonic deque of indices can report every minimum or maximum in linear total time. It stores only candidates that could still become the answer, ordered from the current extreme at the front to the least useful value at the back. A new element removes every older value that it dominates, while expired indices are removed from the front. Each index is again pushed and popped at most once.

\begin{lstlisting}[language=C++]
#include <algorithm>
#include <cassert>
#include <cstdint>
#include <deque>
#include <functional>
#include <tuple>
#include <unordered_map>
#include <utility>
#include <vector>
\end{lstlisting}
The examples below cover sorted 2-sum/3-sum variants, finding the shortest subarray with sum at least a target when all values are nonnegative, and finding the longest subarray with at most $k$ distinct values. They also cover the minimum or maximum of every fixed-length window and an online dynamic programming recurrence.

\begin{compactitem}
  \item \inlinecode{two\_sum\_sorted(a, target)} returns two indices whose values sum to \inlinecode{target}, or $(-1, -1)$ if no such pair exists. The input array must be sorted.
  \item \inlinecode{three\_sum(a, target)} returns a tuple of three original indices whose values sum to \inlinecode{target}, or $(-1, -1, -1)$ if none exist.
  \item \inlinecode{min\_subarray\_at\_least(a, target)} returns a tuple (\inlinecode{length}, \inlinecode{lo}, \inlinecode{hi}), the minimum length and inclusive endpoints of a contiguous subarray with sum at least \inlinecode{target}, or length $-1$ if none exists. Values in \inlinecode{a} must be nonnegative. If \inlinecode{target} is nonpositive, it returns the empty subarray.
  \item \inlinecode{max\_subarray\_at\_most\_k\_distinct(a, k)} returns a tuple (\inlinecode{length}, \inlinecode{lo}, \inlinecode{hi}), the maximum length and inclusive endpoints of a contiguous subarray containing at most \inlinecode{k} distinct values. If \inlinecode{k} is nonpositive, it returns the empty subarray.
  \item \inlinecode{sliding\_window\_extrema(a, k, comp = std::less\textless{}\textgreater{})} returns the extreme value in each window of length \inlinecode{k}. With the default \inlinecode{less\textless{}\textgreater{}} comparator it returns minimums; passing \inlinecode{greater\textless{}\textgreater{}} returns maximums.
\end{compactitem}

\Needspace*{3\baselineskip}
\textbf{Time Complexity}:
\begin{compactitem}
  \item $O(n)$ per call to \inlinecode{two\_sum\_sorted()} and \inlinecode{min\_subarray\_at\_least()}, where $n$ is the array size.
  \item $O(n^{2})$ per call to \inlinecode{three\_sum()}.
  \item $O(n)$ expected per call to \inlinecode{max\_subarray\_at\_most\_k\_distinct()}.
  \item $O(n)$ per call to \inlinecode{sliding\_window\_extrema()}.
\end{compactitem}

\Needspace*{3\baselineskip}
\textbf{Space Complexity}:
\begin{compactitem}
  \item $O(1)$ auxiliary for \inlinecode{two\_sum\_sorted(a, target)}.
  \item $O(n)$ auxiliary for \inlinecode{three\_sum(a, target)}.
  \item $O(1)$ auxiliary for \inlinecode{min\_subarray\_at\_least(a, target)}.
  \item $O(k)$ auxiliary for \inlinecode{max\_subarray\_at\_most\_k\_distinct(a, k)}.
  \item $O(k)$ auxiliary and $O(n)$ for the result of \inlinecode{sliding\_window\_extrema()}.
\end{compactitem}

\begin{lstlisting}[language=C++]
std::pair<int, int> two_sum_sorted(const std::vector<int> &a, int64_t target) {
  for (int lo = 0, hi = static_cast<int>(a.size()) - 1; lo < hi;) {
    int64_t sum = static_cast<int64_t>(a[lo]) + a[hi];
    if (sum == target) {
      return {lo, hi};
    } else if (sum < target) {
      lo++;
    } else {
      hi--;
    }
  }
  return {-1, -1};
}

std::tuple<int, int, int> three_sum(const std::vector<int> &a, int64_t target) {
  int n = static_cast<int>(a.size());
  std::vector<std::pair<int, int>> sorted;
  sorted.reserve(a.size());
  for (int i = 0; i < n; i++) {
    sorted.emplace_back(a[i], i);
  }
  std::sort(sorted.begin(), sorted.end());
  for (int i = 0; i < n; i++) {
    int lo = i + 1, hi = n - 1;
    while (lo < hi) {
      int64_t sum = static_cast<int64_t>(sorted[i].first) + sorted[lo].first + sorted[hi].first;
      if (sum == target) {
        return {sorted[i].second, sorted[lo].second, sorted[hi].second};
      } else if (sum < target) {
        lo++;
      } else {
        hi--;
      }
    }
  }
  return {-1, -1, -1};
}

std::tuple<int, int, int> min_subarray_at_least(const std::vector<int> &a, int64_t target) {
  if (target <= 0) {
    return {0, 0, -1};
  }
  int best = static_cast<int>(a.size()) + 1, best_lo = 0, best_hi = -1;
  int64_t sum = 0;
  for (int lo = 0, hi = 0; hi < static_cast<int>(a.size()); hi++) {
    sum += a[hi];  // Overflow warning.
    while (sum >= target) {
      if (hi - lo + 1 < best) {
        best = hi - lo + 1;
        best_lo = lo;
        best_hi = hi;
      }
      sum -= a[lo++];
    }
  }
  if (best == static_cast<int>(a.size()) + 1) {
    return {-1, 0, -1};
  }
  return {best, best_lo, best_hi};
}

std::tuple<int, int, int> max_subarray_at_most_k_distinct(const std::vector<int> &a, int k) {
  if (k <= 0) {
    return {0, 0, -1};
  }
  std::unordered_map<int, int> count;
  int best = 0, best_lo = 0, best_hi = -1;
  for (int lo = 0, hi = 0; hi < static_cast<int>(a.size()); hi++) {
    count[a[hi]]++;
    while (static_cast<int>(count.size()) > k) {
      if (--count[a[lo]] == 0) {
        count.erase(a[lo]);
      }
      lo++;
    }
    if (best < hi - lo + 1) {
      best = hi - lo + 1;
      best_lo = lo;
      best_hi = hi;
    }
  }
  return {best, best_lo, best_hi};
}

template<typename T, typename Compare = std::less<>>
std::vector<T> sliding_window_extrema(const std::vector<T> &a, int k, Compare comp = Compare{}) {
  int n = static_cast<int>(a.size());
  assert(1 <= k && k <= n);
  std::deque<int> window;
  std::vector<T> res;
  res.reserve(n - k + 1);
  for (int i = 0; i < n; i++) {
    while (!window.empty() && !comp(a[window.back()], a[i])) {
      window.pop_back();
    }
    window.push_back(i);
    if (window.front() <= i - k) {
      window.pop_front();
    }
    if (i >= k - 1) {
      res.push_back(a[window.front()]);
    }
  }
  return res;
}
\end{lstlisting}
A monotone queue can be used to maintain the minimum or maximum value in an online sliding window. This is useful for dynamic programming recurrences where each transition may only come from one of the last $w$ states, such as $dp(i) = a_i + \min_{j \in [i - w, i - 1]} dp(j)$.

The queue stores candidate (\inlinecode{index}, \inlinecode{value}) pairs in monotone order. Expired indices are removed from the front, and dominated values are removed from the back before inserting a new candidate. An index is an ordered \inlinecode{int} timestamp used only for expiration: indices need not be contiguous or nonnegative, but they must be pushed in strictly increasing order, and expiration thresholds must not decrease.

\begin{compactitem}
  \item \inlinecode{MonotoneQueue\textless{}T\textgreater{}()} constructs an empty queue for minimum queries. Instantiate \inlinecode{MonotoneQueue\textless{}T, std::greater\textless{}T\textgreater{}\textgreater{}} for maximum queries.
  \item \inlinecode{empty()} returns whether the queue has no active candidates.
  \item \inlinecode{push(index, value)} inserts candidate \inlinecode{value} with timestamp \inlinecode{index}, removing dominated candidates from the back.
  \item \inlinecode{expire(first\_valid)} removes candidates whose indices are less than \inlinecode{first\_valid}.
  \item \inlinecode{top()} returns the best active (\inlinecode{index}, \inlinecode{value}) pair. The queue must be nonempty.
\end{compactitem}

\Needspace*{3\baselineskip}
\textbf{Time Complexity}:
\begin{compactitem}
  \item $O(n)$ for any sequence of $n$ calls to \inlinecode{push(index, value)} and \inlinecode{expire(first\_valid)}.
  \item $O(1)$ amortized per call to \inlinecode{push(index, value)} and \inlinecode{expire(first\_valid)}.
  \item $O(1)$ worst-case per call to \inlinecode{top()} and \inlinecode{empty()}.
\end{compactitem}

\textbf{Space Complexity}: $O(w)$ for storage of a sliding window of width $w$.

\begin{lstlisting}[language=C++]
template<typename T, typename Compare = std::less<T>>
class MonotoneQueue {
  std::deque<std::pair<int, T>> q;
  Compare better;

 public:
  bool empty() const { return q.empty(); }

  void push(int index, const T &value) {
    while (!q.empty() && !better(q.back().second, value)) {
      q.pop_back();
    }
    q.emplace_back(index, value);
  }

  void expire(int first_valid) {
    while (!q.empty() && q.front().first < first_valid) {
      q.pop_front();
    }
  }

  std::pair<int, T> top() const {
    assert(!q.empty());
    return q.front();
  }
};
\end{lstlisting}
\Needspace*{4\baselineskip}
\textbf{Example Usage}\par
\begin{lstlisting}[language=C++]
using namespace std;

int main() {
  vector<int> c{-3, -1, 2, 4, 8, 11};
  assert(two_sum_sorted(c, 10) == make_pair(1, 5));  // -1 + 11.
  assert(two_sum_sorted(c, 7) == make_pair(1, 4));   // -1 + 8.
  assert(two_sum_sorted(c, 100) == make_pair(-1, -1));

  vector<int> d{4, -1, 8, 2, -3, 11};
  auto [i, j, k] = three_sum(d, 9);
  assert(d[i] + d[j] + d[k] == 9);
  assert((three_sum(d, 30) == make_tuple(-1, -1, -1)));

  vector<int> a{2, 3, 1, 2, 4, 3};
  auto [length, lo, hi] = min_subarray_at_least(a, 7);
  assert(length == 2 && lo == 4 && hi == 5);  // [4, 3].
  auto [empty_length, empty_lo, empty_hi] = min_subarray_at_least(a, 0);
  assert(empty_length == 0 && empty_lo == 0 && empty_hi == -1);

  vector<int> b{1, 2, 1, 3, 4, 3, 5};
  auto [longest, longest_lo, longest_hi] = max_subarray_at_most_k_distinct(b, 2);
  assert(longest == 3 && longest_lo == 0 && longest_hi == 2);  // [1, 2, 1].

  vector<int> e{1, 3, -1, -3, 5, 3, 6, 7};
  assert((sliding_window_extrema(e, 3) == vector<int>{-1, -3, -3, -3, 3, 3}));
  assert((sliding_window_extrema(e, 3, greater<>()) == vector<int>{3, 3, 5, 5, 6, 7}));

  // Toy DP: dp[i] = f[i] + min(dp[j]) for j in [i - w, i - 1].
  // Scanning the window costs O(w) per state; the monotone queue makes it O(1) amortized.
  vector<int> f{4, 2, 7, 1, 3, 6};
  int w = 3;
  vector<int> dp(f.size());
  MonotoneQueue<int> mq;
  dp[0] = f[0];
  mq.push(0, dp[0]);
  for (int i = 1; i < static_cast<int>(f.size()); i++) {
    mq.expire(i - w);
    dp[i] = f[i] + mq.top().second;
    mq.push(i, dp[i]);
  }
  assert((dp == vector<int>{4, 6, 11, 5, 8, 11}));
  return 0;
}
\end{lstlisting}
\subsection{Monotonic Stack}
A monotonic stack maintains its elements in sorted order as values are pushed, popping away any entries that would violate the ordering. Scanning an array once while keeping such a stack of indices answers, for every position, where the nearest smaller or larger neighbor lies. These ``nearest smaller/greater'' relationships underlie many array problems, the most famous being the largest rectangle in a histogram. Each index is pushed and popped at most once, so a full scan runs in linear time. A largest all-zero submatrix can then be found by sweeping rows, turning each row into a histogram of consecutive zeros ending at that row.

For any histogram bar, the nearest strictly lower bars to its left and right are the first positions that prevent a rectangle at the current bar's height from extending farther. The widest such rectangle therefore spans the positions strictly between those boundaries. For a binary matrix, the height at each column counts consecutive zeros ending in the current row; every all-zero rectangle has some bottom row, so solving one histogram per row considers every possible rectangle.

All functions below take a vector \inlinecode{a} of $n$ comparable values and return a vector of $n$ indices. A ``less'' query uses a strictly smaller neighbor and a ``greater'' query uses a strictly larger one; changing the comparison from strict to non-strict (e.g. \inlinecode{\textgreater{}=} to \inlinecode{\textgreater{}}) toggles how ties are handled.

\begin{compactitem}
  \item \inlinecode{prev\_less(a)} returns, for each $i$, the largest index $j < i$ with $a[j] < a[i]$, or $-1$ if there's no such index.
  \item \inlinecode{next\_less(a)} returns, for each $i$, the smallest index $j > i$ with $a[j] < a[i]$, or $n$ if there's no such index.
  \item \inlinecode{prev\_greater(a)} returns, for each $i$, the largest index $j < i$ with $a[j] > a[i]$, or $-1$ if there's no such index.
  \item \inlinecode{next\_greater(a)} returns, for each $i$, the smallest index $j > i$ with $a[j] > a[i]$, or $n$ if there's no such index.
  \item \inlinecode{largest\_histogram\_rectangle(heights)} returns the maximum area and its inclusive left and right endpoints in a histogram where bars have width $1$ and nonnegative heights given by \inlinecode{heights}.
  \item \inlinecode{largest\_zero\_submatrix(a)} returns the maximum area and its inclusive top, left, bottom, and right boundaries among all-zero rectangular submatrices of the 0/1 matrix \inlinecode{a}.
\end{compactitem}

Overflow warning: Histogram areas must fit in the height value type.

\Needspace*{3\baselineskip}
\textbf{Time Complexity}:
\begin{compactitem}
  \item $O(n)$ per call to all one-dimensional functions, where $n$ is the size of the input.
  \item $O(R \cdot C)$ per call to \inlinecode{largest\_zero\_submatrix()}, where $R$ and $C$ are the matrix dimensions.
\end{compactitem}

\Needspace*{3\baselineskip}
\textbf{Space Complexity}:
\begin{compactitem}
  \item $O(n)$ auxiliary and $O(n)$ for the returned indices from each neighbor query.
  \item $O(n)$ auxiliary for \inlinecode{largest\_histogram\_rectangle()}.
  \item $O(C)$ auxiliary for \inlinecode{largest\_zero\_submatrix()}.
\end{compactitem}

\begin{lstlisting}[language=C++]
#include <cassert>
#include <stack>
#include <tuple>
#include <vector>

template<typename T>
std::vector<int> prev_less(const std::vector<T> &a) {
  int n = static_cast<int>(a.size());
  std::vector<int> res(n);
  std::stack<int> s;
  for (int i = 0; i < n; i++) {
    while (!s.empty() && a[s.top()] >= a[i]) {
      s.pop();
    }
    res[i] = s.empty() ? -1 : s.top();
    s.push(i);
  }
  return res;
}

template<typename T>
std::vector<int> next_less(const std::vector<T> &a) {
  int n = static_cast<int>(a.size());
  std::vector<int> res(n);
  std::stack<int> s;
  for (int i = n - 1; i >= 0; i--) {
    while (!s.empty() && a[s.top()] >= a[i]) {
      s.pop();
    }
    res[i] = s.empty() ? n : s.top();
    s.push(i);
  }
  return res;
}

template<typename T>
std::vector<int> prev_greater(const std::vector<T> &a) {
  int n = static_cast<int>(a.size());
  std::vector<int> res(n);
  std::stack<int> s;
  for (int i = 0; i < n; i++) {
    while (!s.empty() && a[s.top()] <= a[i]) {
      s.pop();
    }
    res[i] = s.empty() ? -1 : s.top();
    s.push(i);
  }
  return res;
}

template<typename T>
std::vector<int> next_greater(const std::vector<T> &a) {
  int n = static_cast<int>(a.size());
  std::vector<int> res(n);
  std::stack<int> s;
  for (int i = n - 1; i >= 0; i--) {
    while (!s.empty() && a[s.top()] <= a[i]) {
      s.pop();
    }
    res[i] = s.empty() ? n : s.top();
    s.push(i);
  }
  return res;
}

template<typename T>
std::tuple<T, int, int> largest_histogram_rectangle(const std::vector<T> &heights) {
  int n = static_cast<int>(heights.size());
  std::vector<int> left = prev_less(heights), right = next_less(heights);
  T best{};
  int lo = 0, hi = -1;
  for (int i = 0; i < n; i++) {
    T area = heights[i] * (right[i] - left[i] - 1);  // Overflow warning.
    if (area > best) {
      best = area;
      lo = left[i] + 1;
      hi = right[i] - 1;
    }
  }
  return {best, lo, hi};
}

std::tuple<int, int, int, int, int> largest_zero_submatrix(
    const std::vector<std::vector<char>> &a
) {
  if (a.empty()) {
    return {0, 0, 0, -1, -1};
  }
  int rows = static_cast<int>(a.size()), cols = static_cast<int>(a[0].size());
  int best = 0, rlo = 0, clo = 0, rhi = -1, chi = -1;
  std::vector<int> height(cols);
  for (int i = 0; i < rows; i++) {
    for (int j = 0; j < cols; j++) {
      height[j] = a[i][j] ? 0 : height[j] + 1;
    }
    auto [area, lo, hi] = largest_histogram_rectangle(height);
    if (area > best) {
      best = area;
      rlo = i - area / (hi - lo + 1) + 1;
      clo = lo;
      rhi = i;
      chi = hi;
    }
  }
  return {best, rlo, clo, rhi, chi};
}
\end{lstlisting}
\Needspace*{4\baselineskip}
\textbf{Example Usage}\par
\begin{lstlisting}[language=C++]
using namespace std;

int main() {
  vector<int> a{2, 1, 4, 3};
  assert((prev_less(a) == vector<int>{-1, -1, 1, 1}));
  assert((next_less(a) == vector<int>{1, 4, 3, 4}));
  assert((prev_greater(a) == vector<int>{-1, 0, -1, 2}));
  assert((next_greater(a) == vector<int>{2, 2, 4, 4}));

  vector<int> hist{2, 1, 5, 6, 2, 3};
  // The best histogram rectangle uses bars 2 and 3: min height 5 times width 2.
  assert((largest_histogram_rectangle(hist) == make_tuple(10, 2, 3)));
  assert((largest_histogram_rectangle(vector<int>{}) == make_tuple(0, 0, -1)));

  vector<vector<char>> grid{
      {1, 0, 1, 1, 0, 0},
      {1, 0, 0, 1, 0, 0},
      {0, 0, 0, 0, 0, 1},
      {1, 0, 0, 1, 0, 0},
      {1, 0, 1, 0, 0, 1}
  };
  // The largest zero rectangle has 3 rows and 2 columns.
  assert((largest_zero_submatrix(grid) == make_tuple(6, 1, 1, 3, 2)));
  assert((largest_zero_submatrix({}) == make_tuple(0, 0, 0, -1, -1)));
  return 0;
}
\end{lstlisting}
\subsection{Aggregate Queue}
Maintain a first-in, first-out queue that additionally answers, in $O(1)$ amortized time, the fold of an associative operation over all elements currently in the queue. This is the ``sliding window aggregation'' (SWAG) technique: by pushing new elements and popping old ones as a window advances, it reports an aggregate of every window in linear total time.

Unlike the monotonic deque used for sliding-window extrema (see 1.2.3), the operation need only be associative, not order-based, so the same structure handles min, max, sum, gcd, matrix products, and other monoids. It is sometimes called a ``monotonic queue,'' but that term usually denotes the monotonic deque of 1.2.3 (which handles only min/max); the two are different structures. The queue is built from two stacks: a back stack receives pushes while a front stack serves pops, and each stack stores running folds so the two halves combine in $O(1)$. Whenever the front stack empties, the back stack is flipped onto it; this costs $O(1)$ amortized per element, since each element is moved between stacks at most once.

The fold is defined by an associative \inlinecode{combine(a, b)} function. The default code below returns the ``min'' of the elements; for ``sum'', \inlinecode{combine(a, b)} should return \inlinecode{a + b}. The fold is taken in queue order (front to back), so non-commutative operations such as matrix products are also supported.

\begin{compactitem}
  \item \inlinecode{AggregateQueue\textless{}T\textgreater{}()} constructs an empty queue.
  \item \inlinecode{empty()} returns whether the queue is empty.
  \item \inlinecode{size()} returns the number of elements in the queue.
  \item \inlinecode{push(x)} appends \inlinecode{x} to the back of the queue.
  \item \inlinecode{pop()} removes the element at the front of the queue, which must be non-empty.
  \item \inlinecode{front()} returns the element at the front of the queue, which must be non-empty.
  \item \inlinecode{fold()} returns \inlinecode{combine()} applied to all elements in queue order; the queue must be non-empty.
\end{compactitem}

\Needspace*{3\baselineskip}
\textbf{Time Complexity}:
\begin{compactitem}
  \item $O(1)$ per call to the constructor, \inlinecode{empty()}, \inlinecode{size()}, \inlinecode{front()}, and \inlinecode{fold()}.
  \item $O(1)$ amortized per call to \inlinecode{push()} and \inlinecode{pop()}.
\end{compactitem}

\Needspace*{3\baselineskip}
\textbf{Space Complexity}:
\begin{compactitem}
  \item $O(n)$ for storage of the queue elements, where $n$ is the number of elements.
  \item $O(1)$ auxiliary per operation, amortized (a \inlinecode{pop()} may flip the back stack onto the front stack).
\end{compactitem}

\begin{lstlisting}[language=C++]
#include <algorithm>
#include <cassert>
#include <utility>
#include <vector>

template<typename T>
class AggregateQueue {
  static T combine(const T &a, const T &b) { return std::min(a, b); }

  // Each stack entry is (value, fold), where fold combines the value with all entries between it
  // and its end of the queue (the front for front_stack, the back for back_stack).
  std::vector<std::pair<T, T>> front_stack, back_stack;

 public:
  bool empty() const { return front_stack.empty() && back_stack.empty(); }
  int size() const { return static_cast<int>(front_stack.size() + back_stack.size()); }

  void push(const T &x) {
    T acc = back_stack.empty() ? x : combine(back_stack.back().second, x);
    back_stack.emplace_back(x, acc);
  }

  void pop() {
    assert(!empty());
    if (front_stack.empty()) {
      // Flip the back stack onto the front stack, reversing order so the oldest element ends up on
      // top, and recompute the front folds in queue order.
      while (!back_stack.empty()) {
        T x = back_stack.back().first;
        back_stack.pop_back();
        T acc = front_stack.empty() ? x : combine(x, front_stack.back().second);
        front_stack.emplace_back(x, acc);
      }
    }
    front_stack.pop_back();
  }

  T front() const {
    assert(!empty());
    return front_stack.empty() ? back_stack.front().first : front_stack.back().first;
  }

  T fold() const {
    assert(!empty());
    if (front_stack.empty()) {
      return back_stack.back().second;
    }
    if (back_stack.empty()) {
      return front_stack.back().second;
    }
    return combine(front_stack.back().second, back_stack.back().second);
  }
};
\end{lstlisting}
\Needspace*{4\baselineskip}
\textbf{Example Usage}\par
\begin{lstlisting}[language=C++]
using namespace std;

int main() {
  // Used directly as a FIFO queue with an O(1) min query.
  AggregateQueue<int> q;
  assert(q.empty());
  q.push(5);
  q.push(3);
  q.push(8);
  assert(!q.empty());
  assert(q.front() == 5 && q.size() == 3);
  assert(q.fold() == 3);  // min(5, 3, 8)
  q.pop();                // Removes 5.
  assert(q.fold() == 3);  // min(3, 8)
  q.pop();                // Removes 3.
  assert(q.fold() == 8);
  q.pop();  // Removes 8.
  assert(q.empty());

  // Canonical use: the minimum of every width-3 window, advancing the window by push/pop.
  vector<int> a{4, 2, 5, 1, 3, 6};
  AggregateQueue<int> w;
  vector<int> mins;
  for (int i = 0; i < static_cast<int>(a.size()); i++) {
    w.push(a[i]);
    if (w.size() > 3) {
      w.pop();
    }
    if (w.size() == 3) {
      mins.push_back(w.fold());
    }
  }
  assert((mins == vector<int>{2, 1, 1, 1}));  // min of [4,2,5], [2,5,1], [5,1,3], [1,3,6]
  return 0;
}
\end{lstlisting}

\section{\texorpdfstring{Dynamic Programming}{1.3 Dynamic Programming}}
\setcounter{section}{3}
\setcounter{subsection}{0}
\subsection{Maximum Subarray Sum (Kadane)}
Given an array of numbers, determine the maximum possible sum of any contiguous subarray. Kadane's algorithm scans the array while maintaining a nonnegative running sum. At each position, it adds the current value, records the largest sum seen, and resets the running sum to zero if it becomes negative. Discarding a negative sum is safe because it could only reduce every subarray extending it. This can be adapted to compute the maximal submatrix sum as well.

Two related scans maintain slightly different states. For a circular array, the best wrapped subarray is the whole-array sum minus the minimum-sum subarray. For a maximum-product subarray, a negative value swaps the roles of the minimum and maximum products ending at the previous position, so both products must be tracked.

\begin{compactitem}
  \item \inlinecode{max\_subarray\_sum(a)} returns the sum and inclusive endpoints (\inlinecode{sum}, \inlinecode{begin}, \inlinecode{end}) for the maximal-sum subarray of vector \inlinecode{a}. This implementation requires operators \inlinecode{+} and \inlinecode{\textless{}} on the value type. By convention, the empty subarray is allowed, so an input containing only negative values returns sum $0$ and endpoints $[0, -1]$.
  \item \inlinecode{max\_circular\_subarray\_sum(a)} returns a tuple (\inlinecode{sum}, \inlinecode{begin}, \inlinecode{end}) containing the maximal circular-subarray sum and its inclusive endpoints. If \inlinecode{begin} $\leq$ \inlinecode{end}, the result is an ordinary range; if \inlinecode{begin} \textgreater{} \inlinecode{end}, it wraps from \inlinecode{begin} through the end of \inlinecode{a} and continues through \inlinecode{end}. The empty subarray is allowed under the same convention as \inlinecode{max\_subarray\_sum()}.
  \item \inlinecode{max\_product\_subarray(a)} returns a tuple (\inlinecode{product}, \inlinecode{begin}, \inlinecode{end}) containing the product and inclusive endpoints of a nonempty maximal-product subarray, or product $0$ and endpoints $[0, -1]$ for an empty input.
  \item \inlinecode{max\_submatrix\_sum(a)} returns the sum and inclusive boundaries (\inlinecode{sum}, \inlinecode{r1}, \inlinecode{c1}, \inlinecode{r2}, \inlinecode{c2}) for the largest rectangular submatrix of a matrix \inlinecode{a}. This implementation requires operators \inlinecode{+} and \inlinecode{\textless{}} on the matrix value type. By convention, the empty submatrix is allowed, so a matrix containing only negative values returns sum $0$ with empty row and column intervals.
\end{compactitem}

Overflow warning: All resulting sums and products must fit in the value type.

\Needspace*{3\baselineskip}
\textbf{Time Complexity}:
\begin{compactitem}
  \item $O(n)$ per call to \inlinecode{max\_subarray\_sum()}, \inlinecode{max\_circular\_subarray\_sum()}, and \inlinecode{max\_product\_subarray()}, where $n$ is the size of \inlinecode{a}.
  \item $O(R \cdot C^{2})$ per call to \inlinecode{max\_submatrix\_sum()}, where $R$ is the number of rows and $C$ is the number of columns in the matrix.
\end{compactitem}

\Needspace*{3\baselineskip}
\textbf{Space Complexity}:
\begin{compactitem}
  \item $O(1)$ auxiliary for each subarray function.
  \item $O(R)$ auxiliary for \inlinecode{max\_submatrix\_sum()}.
\end{compactitem}

\begin{lstlisting}[language=C++]
#include <tuple>
#include <utility>
#include <vector>

template<typename T>
std::tuple<T, int, int> max_subarray_sum(const std::vector<T> &a) {
  int n = static_cast<int>(a.size());
  if (n == 0) {
    return {T{}, 0, -1};
  }
  int curr_begin = 0, begin = 0, end = -1;
  T sum = 0, max_sum = 0;
  for (int i = 0; i < n; i++) {
    sum += a[i];  // Overflow warning.
    if (sum < 0) {
      sum = 0;
      curr_begin = i + 1;
    } else if (max_sum < sum) {
      max_sum = sum;
      begin = curr_begin;
      end = i;
    }
  }
  return {max_sum, begin, end};
}

template<typename T>
std::tuple<T, int, int> max_circular_subarray_sum(const std::vector<T> &a) {
  auto [max_sum, begin, end] = max_subarray_sum(a);
  int n = static_cast<int>(a.size());
  if (n == 0) {
    return {max_sum, begin, end};
  }
  int curr_begin = 0, min_begin = 0, min_end = -1;
  T total = 0, min_sum = 0, curr_sum = 0;
  for (int i = 0; i < n; i++) {
    total += a[i];     // Overflow warning.
    curr_sum += a[i];  // Overflow warning.
    if (0 < curr_sum) {
      curr_sum = 0;
      curr_begin = i + 1;
    } else if (curr_sum < min_sum) {
      min_sum = curr_sum;
      min_begin = curr_begin;
      min_end = i;
    }
  }
  T wrapped_sum = total - min_sum;
  if (max_sum < wrapped_sum) {
    max_sum = wrapped_sum;
    begin = (min_end + 1) % n;
    end = (min_begin + n - 1) % n;
  }
  return {max_sum, begin, end};
}

template<typename T>
std::tuple<T, int, int> max_product_subarray(const std::vector<T> &a) {
  int n = static_cast<int>(a.size());
  if (n == 0) {
    return {T{}, 0, -1};
  }
  T max_end = a[0], min_end = a[0], best = a[0];
  int max_begin = 0, min_begin = 0, begin = 0, end = 0;
  for (int i = 1; i < n; i++) {
    if (a[i] < 0) {
      std::swap(max_end, min_end);
      std::swap(max_begin, min_begin);
    }
    T next_max = max_end * a[i];  // Overflow warning.
    T next_min = min_end * a[i];  // Overflow warning.
    if (next_max < a[i]) {
      max_end = a[i];
      max_begin = i;
    } else {
      max_end = next_max;
    }
    if (a[i] < next_min) {
      min_end = a[i];
      min_begin = i;
    } else {
      min_end = next_min;
    }
    if (best < max_end) {
      best = max_end;
      begin = max_begin;
      end = i;
    }
  }
  return {best, begin, end};
}

template<typename T>
std::tuple<T, int, int, int, int> max_submatrix_sum(const std::vector<std::vector<T>> &a) {
  if (a.empty() || a[0].empty()) {
    return {T{}, 0, 0, -1, -1};
  }
  int rows = static_cast<int>(a.size()), cols = static_cast<int>(a[0].size());
  int r1 = 0, c1 = 0, r2 = -1, c2 = -1;
  T max_sum = 0;
  for (int clo = 0; clo < cols; clo++) {
    std::vector<T> sums(rows);
    for (int chi = clo; chi < cols; chi++) {
      for (int i = 0; i < rows; i++) {
        sums[i] += a[i][chi];  // Overflow warning.
      }
      auto [sum, rlo, rhi] = max_subarray_sum(sums);
      if (max_sum < sum) {
        max_sum = sum;
        r1 = rlo;
        c1 = clo;
        r2 = rhi;
        c2 = chi;
      }
    }
  }
  return {max_sum, r1, c1, r2, c2};
}
\end{lstlisting}
\Needspace*{4\baselineskip}
\textbf{Example Usage}\par
\begin{lstlisting}[language=C++]
#include <cassert>
#include <iostream>
using namespace std;

int main() {
  {
    // All negative values, so the empty subarray is maximal.
    auto [sum, begin, end] = max_subarray_sum(vector<int>{-2, -1, -3});
    assert(sum == 0 && begin == 0 && end == -1);
  }
  {
    vector<int> a{-2, -1, -3, 4, -1, 2, 1, -5, 4};
    auto [sum, begin, end] = max_subarray_sum(a);
    assert(sum == 6);
    cout << "Maximal sum subarray:" << endl;
    for (int i = begin; i <= end; i++) {
      cout << a[i] << " ";
    }
    cout << endl;
  }
  {
    vector<int> a{5, -3, 5};
    auto [sum, begin, end] = max_circular_subarray_sum(a);
    assert(sum == 10 && begin == 2 && end == 0);
  }
  {
    vector<int> a{2, 3, -2, 4};
    auto [product, begin, end] = max_product_subarray(a);
    assert(product == 6 && begin == 0 && end == 1);
  }
  {
    vector<vector<int>> a{
        {0, -2, -7, 0, 5},
        {9, 2, -6, 2, -4},
        {-4, 1, -4, 1, 0},
        {-1, 8, 0, -2, 3},
    };
    auto [sum, r1, c1, r2, c2] = max_submatrix_sum(a);
    assert(sum == 15);
    cout << "\nMaximal sum submatrix:" << endl;
    for (int i = r1; i <= r2; i++) {
      for (int j = c1; j <= c2; j++) {
        cout << a[i][j] << " ";
      }
      cout << endl;
    }
  }
  return 0;
}
\end{lstlisting}
\begin{exampleoutput}
Maximal sum subarray:
4 -1 2 1

Maximal sum submatrix:
9 2
-4 1
-1 8
\end{exampleoutput}
\subsection{Longest Increasing Subsequence}
Given an array of elements, determine a longest subsequence where all elements are in strictly increasing order. The subsequence is not necessarily contiguous or unique, so only one such answer will be found. The answer is computed by scanning left to right while maintaining, for each length, the smallest possible tail value of an increasing subsequence of that length. Each element extends or improves one of these tails, located in $O(\log n)$ by binary search; predecessor links then reconstruct the subsequence. The comparator \inlinecode{comp} defines the ordering and defaults to \inlinecode{std::less\textless{}\textgreater{}}.

\begin{compactitem}
  \item \inlinecode{longest\_increasing\_subsequence(a, comp = std::less\textless{}\textgreater{}())} returns the indices of one longest strictly increasing subsequence of \inlinecode{a} according to \inlinecode{comp}, in order.
\end{compactitem}

\textbf{Time Complexity}: $O(n \log n)$ per call, where $n$ is the size of \inlinecode{a}.

\textbf{Space Complexity}: $O(n)$ auxiliary and $O(n)$ for the returned indices.

\begin{lstlisting}[language=C++]
#include <algorithm>
#include <functional>
#include <vector>

template<typename T, typename Compare = std::less<>>
std::vector<int> longest_increasing_subsequence(const std::vector<T> &a, Compare comp = Compare{}) {
  int n = static_cast<int>(a.size());
  if (n == 0) {
    return {};
  }
  int len = 0;
  // prev[i] = index of the predecessor of element i in its LIS (or -1 for the first LIS element).
  // tail[i] = index of the last element of the best known LIS of length i + 1.
  std::vector<int> prev(n), tail(n);
  for (int i = 0; i < n; i++) {
    // Find the tail to extend or improve. Comparing by value through the stored indices keeps tail
    // index-based for reconstruction. Switch to an upper-bound search for a non-decreasing result.
    int pos = static_cast<int>(
        std::lower_bound(
            tail.begin(), tail.begin() + len, i, [&](int t, int x) { return comp(a[t], a[x]); }
        ) -
        tail.begin()
    );
    if (pos == len) {
      len++;
    }
    prev[i] = pos > 0 ? tail[pos - 1] : -1;
    tail[pos] = i;
  }
  // Optional: reconstruct one longest increasing subsequence.
  std::vector<int> res(len);
  for (int i = tail[len - 1]; i != -1; i = prev[i]) {
    res[--len] = i;
  }
  return res;
}
\end{lstlisting}
\Needspace*{4\baselineskip}
\textbf{Example Usage}\par
\begin{lstlisting}[language=C++]
#include <cassert>
using namespace std;

int main() {
  vector<int> a{-2, -5, 1, 9, 10, 8, 11, 10, 13, 11};
  assert((longest_increasing_subsequence(a) == vector<int>{1, 2, 3, 4, 6, 8}));

  vector<int> b{1, 4, 3, 2};
  assert((longest_increasing_subsequence(b, greater<int>()) == vector<int>{1, 2, 3}));
  return 0;
}
\end{lstlisting}
\subsection{0-1 Knapsack}
Solves the 0-1 knapsack problem: given items with nonnegative integer weights and values, choose each item at most once to maximize total value without exceeding a capacity limit.

The one-dimensional dynamic programming array \inlinecode{dp[w]} stores the best value that can be achieved with total weight at most \inlinecode{w} after processing some prefix of the items. Iterating weights in decreasing order is what enforces the 0-1 constraint; it prevents the same item from being reused during the current item update.

\begin{compactitem}
  \item \inlinecode{knapsack\_01(weight, value, capacity)} returns a pair (\inlinecode{best\_value}, \inlinecode{items}) containing the maximum value and the selected item indices in increasing order, where item \inlinecode{i} has weight \inlinecode{weight[i]} and value \inlinecode{value[i]}. All weights must be nonnegative integers. Items with weight $0$ are allowed and can each be taken once. The take/skip decision for every item and capacity is stored to reconstruct the optimal subset.
\end{compactitem}

\textbf{Time Complexity}: $O(n \cdot W)$ per call, where $n$ is the number of items and $W$ is the capacity.

\textbf{Space Complexity}: $O(n \cdot W)$ auxiliary.

\begin{lstlisting}[language=C++]
#include <algorithm>
#include <cassert>
#include <cstdint>
#include <utility>
#include <vector>

std::pair<int64_t, std::vector<int>> knapsack_01(
    const std::vector<int> &weight, const std::vector<int64_t> &value, int capacity
) {
  int n = static_cast<int>(weight.size());
  assert(static_cast<int>(value.size()) == n && capacity >= 0);
  for (int i = 0; i < n; i++) {
    assert(weight[i] >= 0 && value[i] >= 0);
  }
  std::vector<int64_t> dp(capacity + 1);
  std::vector<std::vector<char>> take(n, std::vector<char>(capacity + 1));
  for (int i = 0; i < n; i++) {
    for (int w = capacity; w >= weight[i]; w--) {
      int64_t candidate = dp[w - weight[i]] + value[i];  // Overflow warning.
      if (dp[w] < candidate) {
        dp[w] = candidate;
        take[i][w] = true;
      }
    }
  }
  // Optional: reconstruct one optimal subset.
  std::vector<int> items;
  for (int i = n - 1, w = capacity; i >= 0; i--) {
    if (take[i][w]) {
      items.push_back(i);
      w -= weight[i];
    }
  }
  std::reverse(items.begin(), items.end());
  return {dp[capacity], items};
}
\end{lstlisting}
\Needspace*{4\baselineskip}
\textbf{Example Usage}\par
\begin{lstlisting}[language=C++]
#include <cassert>
using namespace std;

int main() {
  vector<int> weight{3, 4, 5, 9};
  vector<int64_t> value{4, 5, 7, 10};
  auto [best_value, items] = knapsack_01(weight, value, 8);
  assert(best_value == 11 && (items == vector<int>{0, 2}));
  return 0;
}
\end{lstlisting}
\subsection{Unbounded Knapsack and Coin Change}
Solves unbounded knapsack and two related coin change problems, where each item or coin may be used any number of times. In the maximum-value version, each item has a weight and value, and the goal is to maximize total value subject to a capacity limit. In coin change variants, the goal is either to count unordered ways to form a target sum or to minimize the number of coins used.

For unbounded maximum-value knapsack, capacities are processed in increasing order. Thus, when computing \inlinecode{dp[w]}, every smaller-capacity state \inlinecode{dp[w - weight[i]]} is already final and may itself contain item \inlinecode{i}, allowing unlimited reuse. For unordered coin change, coins are the outer loop so each multiset of coins is counted once, independent of ordering. For minimum coins, each target sum chooses the best previous sum \inlinecode{sum - coin}, and the predecessor coin reconstructs one optimal multiset.

\begin{compactitem}
  \item \inlinecode{unbounded\_knapsack(weight, value, capacity)} returns a pair (\inlinecode{best\_value}, \inlinecode{count}) containing the maximum value achievable with total weight at most \inlinecode{capacity} and \inlinecode{count[i]}, the number of copies selected of each item \inlinecode{i}.
  \item \inlinecode{count\_coin\_change(coins, target)} returns the number of unordered ways to make sum \inlinecode{target} using the given coin denominations.
  \item \inlinecode{min\_coin\_change(coins, target)} returns a pair (\inlinecode{count}, \inlinecode{used}), where \inlinecode{count} is the minimum number of coins needed to make sum \inlinecode{target}, and \inlinecode{used[i]} is the number of copies of \inlinecode{coins[i]} in one optimal multiset. If \inlinecode{target} cannot be formed, \inlinecode{count} is $-1$ and \inlinecode{used} is empty.
\end{compactitem}

\inlinecode{capacity} and \inlinecode{target} must be nonnegative integers. Weights and coin denominations must be positive. Coin denominations must also be distinct when counting ways; duplicate denominations are treated as different coin types and overcount identical multisets.

\Needspace*{3\baselineskip}
\textbf{Time Complexity}:
\begin{compactitem}
  \item $O(n \cdot W)$ per call to \inlinecode{unbounded\_knapsack()}, where $n$ is the number of items and $W$ is \inlinecode{capacity}.
  \item $O(c \cdot t)$ per call to \inlinecode{count\_coin\_change()} and \inlinecode{min\_coin\_change()}, where $c$ is the number of coin denominations and $t$ is the target.
\end{compactitem}

\Needspace*{3\baselineskip}
\textbf{Space Complexity}:
\begin{compactitem}
  \item $O(W)$ auxiliary and $O(n)$ for the returned \inlinecode{count} from \inlinecode{unbounded\_knapsack()}.
  \item $O(t)$ auxiliary for \inlinecode{count\_coin\_change()}.
  \item $O(t)$ auxiliary and $O(c)$ for the returned \inlinecode{used} from \inlinecode{min\_coin\_change()}.
\end{compactitem}

\begin{lstlisting}[language=C++]
#include <cassert>
#include <cstdint>
#include <utility>
#include <vector>

std::pair<int64_t, std::vector<int>> unbounded_knapsack(
    const std::vector<int> &weight, const std::vector<int64_t> &value, int capacity
) {
  int n = static_cast<int>(weight.size());
  assert(static_cast<int>(value.size()) == n && capacity >= 0);
  for (int w : weight) {
    assert(w > 0);
  }
  std::vector<int64_t> dp(capacity + 1);
  std::vector<int> choice(capacity + 1, -1);
  for (int w = 1; w <= capacity; w++) {
    for (int i = 0; i < n; i++) {
      if (weight[i] <= w) {
        int64_t candidate = dp[w - weight[i]] + value[i];  // Overflow warning.
        if (dp[w] < candidate) {
          dp[w] = candidate;
          choice[w] = i;
        }
      }
    }
  }
  // Optional: reconstruct one optimal multiset of items.
  std::vector<int> count(n);
  for (int w = capacity; choice[w] != -1; w -= weight[choice[w]]) {
    count[choice[w]]++;
  }
  return {dp[capacity], count};
}

int64_t count_coin_change(const std::vector<int> &coins, int target) {
  assert(target >= 0);
  std::vector<int64_t> dp(target + 1);
  dp[0] = 1;
  for (int coin : coins) {
    assert(coin > 0);
    for (int sum = coin; sum <= target; sum++) {
      dp[sum] += dp[sum - coin];  // Overflow warning.
    }
  }
  return dp[target];
}

std::pair<int, std::vector<int>> min_coin_change(const std::vector<int> &coins, int target) {
  assert(target >= 0);
  for (int coin : coins) {
    assert(coin > 0);
  }
  int inf = target + 1;
  std::vector<int> dp(target + 1, inf), choice(target + 1, -1);
  dp[0] = 0;
  for (int sum = 1; sum <= target; sum++) {
    for (int i = 0; i < static_cast<int>(coins.size()); i++) {
      if (coins[i] <= sum && dp[sum] > dp[sum - coins[i]] + 1) {
        dp[sum] = dp[sum - coins[i]] + 1;
        choice[sum] = i;
      }
    }
  }
  if (dp[target] == inf) {
    return {-1, {}};
  }
  // Optional: reconstruct one minimum-size multiset of coins.
  std::vector<int> used(coins.size());
  for (int sum = target; sum > 0; sum -= coins[choice[sum]]) {
    used[choice[sum]]++;
  }
  return {dp[target], used};
}
\end{lstlisting}
\Needspace*{4\baselineskip}
\textbf{Example Usage}\par
\begin{lstlisting}[language=C++]
#include <cassert>
using namespace std;

int main() {
  vector<int> weight{3, 4, 5};
  vector<int64_t> value{4, 5, 7};
  auto [best_value, count] = unbounded_knapsack(weight, value, 10);
  assert(best_value == 14 && (count == vector<int>{0, 0, 2}));

  vector<int> coins{1, 2, 5};
  assert(count_coin_change(coins, 5) == 4);  // 5, 2+2+1, 2+1+1+1, 1*5.

  auto [coin_count, used] = min_coin_change(coins, 11);
  assert(coin_count == 3 && (used == vector<int>{1, 0, 2}));

  auto [impossible, empty_used] = min_coin_change(vector<int>{4, 7}, 6);
  assert(impossible == -1 && empty_used.empty());
  return 0;
}
\end{lstlisting}
\subsection{Matrix Chain Parenthesization}
Given a fixed-order chain of compatible matrices, finds a parenthesization that minimizes the number of scalar multiplications needed to compute their product. The algorithm examines only the matrix dimensions; it does not multiply matrix entries. Matrix multiplication is associative, but multiplying an $a$-by-$b$ matrix with a $b$-by-$c$ matrix costs $abc$ scalar multiplications, so different parenthesizations can have very different costs. This is a canonical example of interval dynamic programming.

Write $d_i$ for \inlinecode{dimensions[i]}, so matrix $i$ has $d_i$ rows and $d_{i+1}$ columns. For a chain of matrices, $dp(l, r)$ is the minimum cost of multiplying matrices $l$ through $r$. Trying every final split $k$ gives the recurrence $dp(l, r) = \min_{l \le k < r}(dp(l, k) + dp(k + 1, r) + d_l d_{k+1} d_{r+1})$. The final product multiplies a $d_l$-by-$d_{k+1}$ matrix with a $d_{k+1}$-by-$d_{r+1}$ matrix. Processing intervals in increasing order of length ensures that both subintervals have already been solved. The same pattern applies whenever a contiguous interval is formed by combining two smaller intervals.

\begin{compactitem}
  \item \inlinecode{matrix\_chain\_order(dimensions)} returns a pair (\inlinecode{operations}, \inlinecode{parenthesization}) containing the minimum number of scalar multiplications and one optimal parenthesization. Matrix \inlinecode{i} has dimensions \inlinecode{dimensions[i]} by \inlinecode{dimensions[i + 1]}; all dimensions must be positive.
\end{compactitem}

Overflow warning: The minimum cost and all intermediate products must fit in \inlinecode{int64\_t}.

\textbf{Time Complexity}: $O(n^{3})$ per call, where $n$ is the number of matrices. The dimension values affect the returned cost but not the running time.

\textbf{Space Complexity}: $O(n^{2})$ auxiliary for the dynamic programming and split tables.

\begin{lstlisting}[language=C++]
#include <cassert>
#include <cstdint>
#include <string>
#include <utility>
#include <vector>

std::pair<int64_t, std::string> matrix_chain_order(const std::vector<int> &dimensions) {
  assert(dimensions.size() >= 2);
  for (int d : dimensions) {
    assert(d > 0);
  }
  int n = static_cast<int>(dimensions.size()) - 1;
  std::vector<std::vector<int64_t>> dp(n, std::vector<int64_t>(n));
  std::vector<std::vector<int>> split(n, std::vector<int>(n));
  for (int length = 2; length <= n; length++) {
    for (int lo = 0; lo + length <= n; lo++) {
      int hi = lo + length - 1;
      for (int k = lo; k < hi; k++) {
        int64_t candidate = dp[lo][k] + dp[k + 1][hi] +
                            static_cast<int64_t>(dimensions[lo]) * dimensions[k + 1] *
                                dimensions[hi + 1];  // Overflow warning.
        if (k == lo || candidate < dp[lo][hi]) {
          dp[lo][hi] = candidate;
          split[lo][hi] = k;
        }
      }
    }
  }
  // Optional: reconstruct one optimal parenthesization.
  auto rec = [&](auto &&rec, int lo, int hi) -> std::string {
    if (lo == hi) {
      return "A" + std::to_string(lo);
    }
    int k = split[lo][hi];
    return "(" + rec(rec, lo, k) + " * " + rec(rec, k + 1, hi) + ")";
  };
  return {dp[0][n - 1], rec(rec, 0, n - 1)};
}
\end{lstlisting}
\Needspace*{4\baselineskip}
\textbf{Example Usage}\par
\begin{lstlisting}[language=C++]
using namespace std;

int main() {
  auto [ops, parenthesization] = matrix_chain_order({40, 20, 30, 10, 30});
  assert(ops == 26000);
  assert(parenthesization == "((A0 * (A1 * A2)) * A3)");

  auto [single_ops, single_parenthesization] = matrix_chain_order({3, 5});
  assert(single_ops == 0 && single_parenthesization == "A0");
  return 0;
}
\end{lstlisting}
\subsection{Digit DP}
Digit dynamic programming counts integers in a large bounded range whose decimal representations satisfy a digit-based property. The concrete routine below counts integers with a specified digit sum and provides the standard template to adapt for other finite-state properties.

To count valid integers in $[0, x]$, process the digits of $x$ from left to right. The state records the current position, accumulated digit sum, and whether the chosen prefix still equals the prefix of $x$. Once that ``tight'' condition is false, the remaining choices depend only on the position and sum, so those states can be memoized. Subtracting the prefix counts for \inlinecode{hi} and \inlinecode{lo - 1} gives the answer for the requested range.

Leading zeros do not affect a digit sum, so every integer can be treated as a zero-padded string of the same length as the bound. For properties where leading zeros matter, add a \inlinecode{started} flag to the state. More generally, the sum can be replaced by any small state with a transition for appending a digit.

\begin{compactitem}
  \item \inlinecode{count\_digit\_sum(lo, hi, target)} returns the number of integers in the inclusive range $[{\inlinecodemath{lo}}, {\inlinecodemath{hi}}]$ whose decimal digits sum to \inlinecode{target}. The bounds and \inlinecode{target} must be nonnegative.
\end{compactitem}

\textbf{Time Complexity}: $O(d \cdot s)$ per call, where $d$ is the number of digits in \inlinecode{hi} and $s$ is \inlinecode{target}.

\textbf{Space Complexity}: $O(d \cdot s)$ auxiliary heap space and $O(d)$ call stack space.

\begin{lstlisting}[language=C++]
#include <cassert>
#include <cstdint>
#include <string>
#include <vector>

int64_t count_digit_sum(int64_t lo, int64_t hi, int target) {
  assert(0 <= lo && lo <= hi && target >= 0);
  auto count_up_to = [&](int64_t bound) {
    if (bound < 0) {
      return 0LL;
    }
    std::string digits = std::to_string(bound);
    int n = static_cast<int>(digits.size());
    if (target > 9 * n) {
      return 0LL;
    }
    std::vector<std::vector<int64_t>> memo(n, std::vector<int64_t>(target + 1, -1));
    auto rec = [&](auto &&rec, int pos, int sum, bool tight) -> int64_t {
      if (sum > target || sum + 9 * (n - pos) < target) {
        return 0;
      }
      if (pos == n) {
        return sum == target;
      }
      if (!tight && memo[pos][sum] != -1) {
        return memo[pos][sum];
      }
      int bound_digit = digits[pos] - '0';
      int limit = tight ? bound_digit : 9;
      int64_t res = 0;
      for (int digit = 0; digit <= limit; digit++) {
        res += rec(rec, pos + 1, sum + digit, tight && digit == bound_digit);
      }
      if (!tight) {
        memo[pos][sum] = res;
      }
      return res;
    };
    return rec(rec, 0, 0, true);
  };
  return count_up_to(hi) - count_up_to(lo - 1);
}
\end{lstlisting}
\Needspace*{4\baselineskip}
\textbf{Example Usage}\par
\begin{lstlisting}[language=C++]
#include <cassert>
using namespace std;

int main() {
  assert(count_digit_sum(0, 20, 2) == 3);   // 2, 11, and 20.
  assert(count_digit_sum(10, 30, 3) == 3);  // 12, 21, and 30.
  assert(count_digit_sum(0, 99, 9) == 10);
  assert(count_digit_sum(0, 0, 0) == 1);
  assert(count_digit_sum(0, 999, 28) == 0);
  return 0;
}
\end{lstlisting}
\subsection{Egg Dropping}
Given $e$ identical eggs and a building of $f$ floors, find the highest floor from which an egg survives a drop, using as few drops as possible in the worst case. An egg that survives can be reused; an egg that breaks is gone. With one egg, the only safe strategy is to climb one floor at a time, needing $f$ drops. With unlimited eggs, binary search needs $\lceil \log_2(f + 1) \rceil$ drops. The problem is what happens in between.

The natural recurrence minimizes over the drop floor and maximizes over the two outcomes, costing $O(e \cdot f^{2})$. Removing that minimization is a technique worth knowing beyond this puzzle: when the quantity being minimized is small and bounded, promote it to a state dimension and solve for the old one. So ask how many floors $d$ drops with $e$ eggs can cover, not how many drops $f$ floors need. A drop splits the searchable floors into those below it, the drop floor, and those above, giving $dp(d, e) = dp(d-1, e-1) + dp(d-1, e) + 1$ with nothing left to minimize. Each entry is then $O(1)$, and the answer is the smallest $d$ whose reach is at least $f$.

That table is $dp(d, e) = \sum_{i=1}^{e} \binom{d}{i}$ in the binomial coefficients of section 6.2.1, which is what keeps $d$ small: enough eggs sum to $2^d - 1$, while one egg leaves only $\binom{d}{1} = d$. Two eggs reach $91$ floors in 13 drops and $105$ in 14, the classic answer. The same inversion fits any question of how many tests distinguish $n$ cases, such as finding one defective item with a limited number of destructive tests.

\begin{compactitem}
  \item \inlinecode{min\_drops(eggs, floors)} returns the number of drops needed in the worst case, where \inlinecode{eggs} is positive and \inlinecode{floors} is nonnegative.
  \item \inlinecode{drop\_floor(eggs, floors)} returns the floor to drop from first under an optimal strategy, which is $dp(d-1, e-1) + 1$ floors above the bottom of the current range, or $0$ when no drop is needed.
\end{compactitem}

\textbf{Time Complexity}: $O(e \cdot d)$ per call, where $e$ is \inlinecode{eggs} and $d$ is the returned number of drops, which is at most $f$ and is $O(\log f)$ once $e \geq \log_2(f)$.

\textbf{Space Complexity}: $O(e)$ auxiliary.

\begin{lstlisting}[language=C++]
#include <cassert>
#include <cstdint>
#include <vector>

int min_drops(int eggs, int floors) {
  assert(eggs > 0 && floors >= 0);
  // reach[e] is the number of floors distinguishable by the current number of drops with e eggs.
  std::vector<int64_t> reach(eggs + 1);
  int drops = 0;
  while (reach[eggs] < floors) {
    drops++;
    for (int e = eggs; e > 0; e--) {  // Descending, so reach[e - 1] is still the previous row.
      reach[e] += reach[e - 1] + 1;
    }
  }
  return drops;
}

int drop_floor(int eggs, int floors) {
  assert(eggs > 0 && floors >= 0);
  int drops = min_drops(eggs, floors);
  if (drops == 0) {
    return 0;
  }
  std::vector<int64_t> reach(eggs + 1);
  for (int i = 0; i < drops - 1; i++) {  // Rebuild the table one drop short of the answer.
    for (int e = eggs; e > 0; e--) {
      reach[e] += reach[e - 1] + 1;
    }
  }
  return static_cast<int>(reach[eggs - 1]) + 1;
}
\end{lstlisting}
\Needspace*{4\baselineskip}
\textbf{Example Usage}\par
\begin{lstlisting}[language=C++]
#include <cassert>

int main() {
  assert(min_drops(1, 0) == 0);
  assert(min_drops(1, 10) == 10);  // One egg forces a linear scan.
  assert(min_drops(2, 10) == 4);
  assert(min_drops(2, 100) == 14);  // The classic two-egg, hundred-floor puzzle.
  assert(min_drops(3, 100) == 9);
  assert(min_drops(50, 1000) == 10);  // Enough eggs to binary search.

  // With two eggs and a hundred floors, the first drop is from floor 14.
  assert(drop_floor(2, 100) == 14);
  assert(drop_floor(1, 10) == 1);
  assert(drop_floor(2, 0) == 0);
  return 0;
}
\end{lstlisting}

\section{\texorpdfstring{Grid Algorithms}{1.4 Grid Algorithms}}
\setcounter{section}{4}
\setcounter{subsection}{0}
\subsection{Grid Transformations}
Transforms a rectangular grid by transposing, rotating, or reflecting its cells. Transposition exchanges rows and columns. Rotation is clockwise by default, while a negative angle rotates counter-clockwise; the angle must be a multiple of $90$ degrees. Together, rotations and reflections generate all eight symmetries of a square grid.

\begin{compactitem}
  \item \inlinecode{transpose(a)} returns the transpose of grid \inlinecode{a}, so each output cell $(c, r)$ equals input cell $(r, c)$.
  \item \inlinecode{transpose\_in\_place(a)} transposes the square grid \inlinecode{a} in place and returns a reference to it.
  \item \inlinecode{rotate(a, degrees = 90)} returns \inlinecode{a} rotated clockwise by \inlinecode{degrees}.
  \item \inlinecode{rotate\_in\_place(a, degrees = 90)} rotates \inlinecode{a} in place and returns a reference to it. The grid must be square for rotations by $90$ or $270$ degrees.
  \item \inlinecode{reflect\_horizontal(a)} returns \inlinecode{a} reflected across its horizontal axis, reversing the row order.
  \item \inlinecode{reflect\_vertical(a)} returns \inlinecode{a} reflected across its vertical axis, reversing every row.
\end{compactitem}

\textbf{Time Complexity}: $O(R \cdot C)$ per call, where $R$ and $C$ are the grid dimensions.

\Needspace*{3\baselineskip}
\textbf{Space Complexity}:
\begin{compactitem}
  \item $O(R \cdot C)$ for the returned grid from the non-in-place operations.
  \item $O(1)$ auxiliary for the in-place operations.
\end{compactitem}

\begin{lstlisting}[language=C++]
#include <algorithm>
#include <cassert>
#include <utility>
#include <vector>

using Grid = std::vector<std::vector<int>>;

Grid transpose(const Grid &a) {
  int rows = static_cast<int>(a.size());
  int cols = a.empty() ? 0 : static_cast<int>(a[0].size());
  Grid res(cols, std::vector<int>(rows));
  for (int r = 0; r < rows; r++) {
    for (int c = 0; c < cols; c++) {
      res[c][r] = a[r][c];
    }
  }
  return res;
}

Grid &transpose_in_place(Grid &a) {
  assert(a.empty() || a.size() == a[0].size());
  for (int r = 0; r < static_cast<int>(a.size()); r++) {
    for (int c = r + 1; c < static_cast<int>(a.size()); c++) {
      std::swap(a[r][c], a[c][r]);
    }
  }
  return a;
}

Grid rotate(const Grid &a, int degrees = 90) {
  assert(degrees % 90 == 0);
  int rows = static_cast<int>(a.size());
  int cols = a.empty() ? 0 : static_cast<int>(a[0].size());
  degrees = (degrees % 360 + 360) % 360;
  Grid res(degrees % 180 == 0 ? rows : cols, std::vector<int>(degrees % 180 == 0 ? cols : rows));
  for (int r = 0; r < rows; r++) {
    for (int c = 0; c < cols; c++) {
      if (degrees == 90) {
        res[c][rows - r - 1] = a[r][c];
      } else if (degrees == 180) {
        res[rows - r - 1][cols - c - 1] = a[r][c];
      } else if (degrees == 270) {
        res[cols - c - 1][r] = a[r][c];
      } else {
        res[r][c] = a[r][c];
      }
    }
  }
  return res;
}

Grid &rotate_in_place(Grid &a, int degrees = 90) {
  assert(degrees % 90 == 0);
  degrees = (degrees % 360 + 360) % 360;
  assert(degrees % 180 == 0 || a.empty() || a.size() == a[0].size());
  if (degrees == 90) {
    transpose_in_place(a);
    for (auto &row : a) {
      std::reverse(row.begin(), row.end());
    }
  } else if (degrees == 180) {
    std::reverse(a.begin(), a.end());
    for (auto &row : a) {
      std::reverse(row.begin(), row.end());
    }
  } else if (degrees == 270) {
    transpose_in_place(a);
    std::reverse(a.begin(), a.end());
  }
  return a;
}

Grid reflect_horizontal(Grid a) {
  std::reverse(a.begin(), a.end());
  return a;
}

Grid reflect_vertical(Grid a) {
  for (auto &row : a) {
    std::reverse(row.begin(), row.end());
  }
  return a;
}
\end{lstlisting}
\Needspace*{4\baselineskip}
\textbf{Example Usage}\par
\begin{lstlisting}[language=C++]
#include <cassert>
using namespace std;

int main() {
  Grid a{{1, 2, 3}, {4, 5, 6}};
  assert((transpose(a) == Grid{{1, 4}, {2, 5}, {3, 6}}));
  assert((rotate(a) == Grid{{4, 1}, {5, 2}, {6, 3}}));
  assert((rotate(a, -90) == Grid{{3, 6}, {2, 5}, {1, 4}}));
  assert((rotate(a, 180) == Grid{{6, 5, 4}, {3, 2, 1}}));
  assert((reflect_horizontal(a) == Grid{{4, 5, 6}, {1, 2, 3}}));
  assert((reflect_vertical(a) == Grid{{3, 2, 1}, {6, 5, 4}}));

  Grid square{{1, 2, 3}, {4, 5, 6}, {7, 8, 9}};
  for (int degrees = -360; degrees <= 360; degrees += 90) {
    Grid rotated = square;
    assert(rotate_in_place(rotated, degrees) == rotate(square, degrees));
  }
  Grid transposed = square;
  assert(transpose_in_place(transposed) == transpose(square));

  Grid rectangular = a;
  assert(rotate_in_place(rectangular, 180) == rotate(a, 180));
  return 0;
}
\end{lstlisting}
\subsection{Grid Traversal}
Traverses a rectangular grid as an implicit unweighted graph: each unblocked cell is a node joined to its unblocked (four-way by default) orthogonal neighbors. Depth-first flood fill finds connected regions, while breadth-first search finds shortest distances because every move has unit length. Seeding the BFS with several cells computes the distance to the nearest source in one pass.

\begin{compactitem}
  \item \inlinecode{inside(r, c, rows, cols)} returns whether $0 \leq {\inlinecodemath{r}} < {\inlinecodemath{rows}}$ and $0 \leq {\inlinecodemath{c}} < {\inlinecodemath{cols}}$.
  \item \inlinecode{adj4(r, c)} returns the four orthogonally adjacent coordinates without filtering those outside the grid. For eight-way movement, also include the four diagonal neighbors.
  \item \inlinecode{component\_area\_perimeter(blocked, r, c)} uses flood fill to return the area and perimeter of the component containing (\inlinecode{r}, \inlinecode{c}). If the starting cell is blocked, both values are $0$.
  \item \inlinecode{grid\_components(blocked)} uses repeated flood fills to return the connected components of cells marked zero, with each cell represented by a pair (\inlinecode{r}, \inlinecode{c}). Movement is allowed in the four orthogonal directions.
  \item \inlinecode{grid\_bfs(blocked, sources)} returns the pair (\inlinecode{dist}, \inlinecode{pred}) from every cell in \inlinecode{sources}. Each source must be an unblocked cell represented by a pair (\inlinecode{r}, \inlinecode{c}). Passing one source gives ordinary single-source BFS. Blocked and unreachable cells have distance $-1$.
  \item \inlinecode{get\_path(pred, tr, tc)} uses the predecessor grid returned by \inlinecode{grid\_bfs()} to return the path from a nearest source to (\inlinecode{tr}, \inlinecode{tc}), including both endpoints, or an empty vector if the target is blocked or unreachable.
  \item \inlinecode{spiral\_order(a)} returns the cells of grid \inlinecode{a} in clockwise spiral order, starting at the top-left corner and walking inward.
\end{compactitem}

\Needspace*{3\baselineskip}
\textbf{Time Complexity}:
\begin{compactitem}
  \item $O(R \cdot C)$ per call to \inlinecode{spiral\_order()}.
  \item $O(R \cdot C)$ per call to \inlinecode{component\_area\_perimeter()}, \inlinecode{grid\_components()}, or \inlinecode{grid\_bfs()}, where $R$ and $C$ are the grid dimensions.
  \item $O(p)$ per call to \inlinecode{get\_path()}, where $p$ is the number of cells in the returned path.
\end{compactitem}

\Needspace*{3\baselineskip}
\textbf{Space Complexity}:
\begin{compactitem}
  \item $O(R \cdot C)$ auxiliary, including the recursion stack used by \inlinecode{component\_area\_perimeter()} and \inlinecode{grid\_components()}.
  \item $O(R \cdot C)$ for the returned \inlinecode{dist} and \inlinecode{pred}, and $O(p)$ for the path returned by \inlinecode{get\_path()}.
\end{compactitem}

\begin{lstlisting}[language=C++]
#include <algorithm>
#include <array>
#include <cassert>
#include <cstddef>
#include <queue>
#include <utility>
#include <vector>

bool inside(int r, int c, int rows, int cols) {
  return 0 <= r && r < rows && 0 <= c && c < cols;
}

std::array<std::pair<int, int>, 4> adj4(int r, int c) {
  return {{{r - 1, c}, {r + 1, c}, {r, c - 1}, {r, c + 1}}};
}

template<typename Fn>
void grid_dfs(
    const std::vector<std::vector<char>> &blocked, int r, int c,
    std::vector<std::vector<char>> &visit, const Fn &f
) {
  int rows = static_cast<int>(blocked.size());
  int cols = static_cast<int>(blocked[0].size());
  visit[r][c] = true;
  f(r, c);
  for (auto [r2, c2] : adj4(r, c)) {
    if (inside(r2, c2, rows, cols) && !blocked[r2][c2] && !visit[r2][c2]) {
      grid_dfs(blocked, r2, c2, visit, f);
    }
  }
}

std::pair<int, int> component_area_perimeter(
    const std::vector<std::vector<char>> &blocked, int r, int c
) {
  int rows = static_cast<int>(blocked.size());
  int cols = blocked.empty() ? 0 : static_cast<int>(blocked[0].size());
  assert(inside(r, c, rows, cols));
  if (blocked[r][c]) {
    return {0, 0};
  }
  int area = 0, perimeter = 0;
  std::vector<std::vector<char>> visit(rows, std::vector<char>(cols));
  grid_dfs(blocked, r, c, visit, [&](int r, int c) {
    area++;
    for (auto [r2, c2] : adj4(r, c)) {
      perimeter += !inside(r2, c2, rows, cols) || blocked[r2][c2];
    }
  });
  return {area, perimeter};
}

std::vector<std::vector<std::pair<int, int>>> grid_components(
    const std::vector<std::vector<char>> &blocked
) {
  int rows = static_cast<int>(blocked.size());
  int cols = blocked.empty() ? 0 : static_cast<int>(blocked[0].size());
  std::vector<std::vector<std::pair<int, int>>> components;
  std::vector<std::vector<char>> visit(rows, std::vector<char>(cols));
  for (int r = 0; r < rows; r++) {
    for (int c = 0; c < cols; c++) {
      if (blocked[r][c] || visit[r][c]) {
        continue;
      }
      components.emplace_back();
      grid_dfs(blocked, r, c, visit, [&](int r, int c) { components.back().emplace_back(r, c); });
    }
  }
  return components;
}

auto grid_bfs(
    const std::vector<std::vector<char>> &blocked, const std::vector<std::pair<int, int>> &sources
) {
  int rows = static_cast<int>(blocked.size());
  int cols = blocked.empty() ? 0 : static_cast<int>(blocked[0].size());
  std::vector<std::vector<int>> dist(rows, std::vector<int>(cols, -1));
  std::vector<std::vector<std::pair<int, int>>> pred(
      rows, std::vector<std::pair<int, int>>(cols, {-1, -1})
  );
  std::queue<std::pair<int, int>> q;
  for (auto [r, c] : sources) {
    assert(inside(r, c, rows, cols) && !blocked[r][c]);
    if (dist[r][c] == -1) {
      dist[r][c] = 0;
      pred[r][c] = {r, c};
      q.emplace(r, c);
    }
  }
  while (!q.empty()) {
    auto [r, c] = q.front();
    q.pop();
    for (auto [r2, c2] : adj4(r, c)) {
      if (inside(r2, c2, rows, cols) && !blocked[r2][c2] && dist[r2][c2] == -1) {
        dist[r2][c2] = dist[r][c] + 1;
        pred[r2][c2] = {r, c};
        q.emplace(r2, c2);
      }
    }
  }
  return std::pair{std::move(dist), std::move(pred)};
}

std::vector<std::pair<int, int>> get_path(
    const std::vector<std::vector<std::pair<int, int>>> &pred, int tr, int tc
) {
  int rows = static_cast<int>(pred.size());
  int cols = pred.empty() ? 0 : static_cast<int>(pred[0].size());
  assert(inside(tr, tc, rows, cols));
  if (pred[tr][tc].first == -1) {
    return {};
  }
  std::vector<std::pair<int, int>> path;
  std::pair cur{tr, tc};
  while (pred[cur.first][cur.second] != cur) {
    path.push_back(cur);
    cur = pred[cur.first][cur.second];
  }
  path.push_back(cur);
  std::reverse(path.begin(), path.end());
  return path;
}

template<typename T>
std::vector<T> spiral_order(const std::vector<std::vector<T>> &a) {
  int rows = static_cast<int>(a.size());
  int cols = a.empty() ? 0 : static_cast<int>(a[0].size());
  std::vector<T> res;
  res.reserve(static_cast<size_t>(rows) * cols);
  // Peel one ring at a time. Guard the last 2 passes so single remaining rows/cols don't repeat.
  for (int top = 0, bottom = rows - 1, left = 0, right = cols - 1; top <= bottom && left <= right;
       top++, bottom--, left++, right--) {
    for (int c = left; c <= right; c++) {
      res.push_back(a[top][c]);
    }
    for (int r = top + 1; r <= bottom; r++) {
      res.push_back(a[r][right]);
    }
    if (top < bottom) {
      for (int c = right - 1; c >= left; c--) {
        res.push_back(a[bottom][c]);
      }
    }
    if (left < right) {
      for (int r = bottom - 1; r > top; r--) {
        res.push_back(a[r][left]);
      }
    }
  }
  return res;
}
\end{lstlisting}
\Needspace*{4\baselineskip}
\textbf{Example Usage}\par
\begin{lstlisting}[language=C++]
#include <cassert>
using namespace std;

int main() {
  vector<vector<char>> blocked{
      {0, 0, 1, 0},
      {1, 0, 1, 0},
      {0, 0, 1, 0},
  };
  auto components = grid_components(blocked);
  assert(components.size() == 2);
  assert(components[0].size() == 5 && components[1].size() == 3);
  assert((component_area_perimeter(blocked, 0, 0) == pair{5, 12}));
  assert((component_area_perimeter(blocked, 0, 2) == pair{0, 0}));

  auto [dist, pred] = grid_bfs(blocked, {{0, 0}, {2, 3}});
  assert((dist == vector<vector<int>>{{0, 1, -1, 2}, {-1, 2, -1, 1}, {4, 3, -1, 0}}));
  assert((get_path(pred, 2, 1) == vector<pair<int, int>>{{0, 0}, {0, 1}, {1, 1}, {2, 1}}));
  assert((get_path(pred, 0, 3) == vector<pair<int, int>>{{2, 3}, {1, 3}, {0, 3}}));
  assert(get_path(pred, 0, 2).empty());

  vector<vector<int>> g{
      {1, 2, 3, 4},
      {5, 6, 7, 8},
      {9, 10, 11, 12},
  };
  assert((spiral_order(g) == vector<int>{1, 2, 3, 4, 8, 12, 11, 10, 9, 5, 6, 7}));
  assert((spiral_order(vector<vector<int>>{{1, 2, 3}}) == vector<int>{1, 2, 3}));
  assert((spiral_order(vector<vector<int>>{{1}, {2}, {3}}) == vector<int>{1, 2, 3}));
  return 0;
}
\end{lstlisting}
\subsection{Grid DP}
Solves basic dynamic programming problems on a rectangular grid. In grid DP, each cell represents a state and transitions usually come from already-computed neighboring cells. Processing rows from top to bottom and columns from left to right makes predecessors above and to the left available when a state is computed.

\begin{compactitem}
  \item \inlinecode{count\_grid\_paths(blocked)} returns the number of paths from the upper-left cell to the lower-right cell, moving only down or right and avoiding blocked cells marked nonzero. The state $dp(r, c)$ counts paths to cell $(r, c)$ from the counts above and left. Only the current row is retained.
  \item \inlinecode{min\_cost\_grid\_path(cost)} returns a pair (\inlinecode{min\_cost}, \inlinecode{path}) containing the minimum cost of a down/right path from the upper-left cell to the lower-right cell and the cells of one optimal path in order. The state $dp(r, c)$ stores the cheapest cost to cell $(r, c)$. Only the current row of costs is retained, while predecessor directions reconstruct one optimal path. If the grid is empty, \inlinecode{min\_cost} is $0$ and \inlinecode{path} is empty.
  \item \inlinecode{largest\_one\_square(a)} returns a tuple (\inlinecode{side}, \inlinecode{top}, \inlinecode{left}) describing a largest all-nonzero square in \inlinecode{a}. The state $dp(c)$ is the largest square ending at the current cell in column $c$; a nonzero cell extends the minimum of the states above, left, and diagonally above-left. If no such square exists, the function returns $(0, -1, -1)$.
\end{compactitem}

Overflow warning: Path counts and path costs must fit in \inlinecode{int64\_t}.

\textbf{Time Complexity}: $O(R \cdot C)$ per call, where $R$ and $C$ are the grid dimensions.

\Needspace*{3\baselineskip}
\textbf{Space Complexity}:
\begin{compactitem}
  \item $O(C)$ auxiliary for \inlinecode{count\_grid\_paths(blocked)}.
  \item $O(R \cdot C)$ auxiliary for predecessor directions, $O(C)$ for path costs, and $O(R + C)$ for the returned path from \inlinecode{min\_cost\_grid\_path(cost)}.
  \item $O(C)$ auxiliary for \inlinecode{largest\_one\_square(a)}.
\end{compactitem}

\begin{lstlisting}[language=C++]
#include <algorithm>
#include <cstdint>
#include <tuple>
#include <utility>
#include <vector>

int64_t count_grid_paths(const std::vector<std::vector<char>> &blocked) {
  int rows = static_cast<int>(blocked.size());
  int cols = blocked.empty() ? 0 : static_cast<int>(blocked[0].size());
  if (rows == 0 || cols == 0 || blocked[0][0] || blocked[rows - 1][cols - 1]) {
    return 0;
  }
  std::vector<int64_t> dp(cols);
  dp[0] = 1;
  for (int r = 0; r < rows; r++) {
    for (int c = 0; c < cols; c++) {
      if (blocked[r][c]) {
        dp[c] = 0;
      } else if (c > 0) {
        dp[c] += dp[c - 1];  // Overflow warning.
      }
    }
  }
  return dp[cols - 1];
}

std::pair<int64_t, std::vector<std::pair<int, int>>> min_cost_grid_path(
    const std::vector<std::vector<int>> &cost
) {
  int rows = static_cast<int>(cost.size());
  int cols = cost.empty() ? 0 : static_cast<int>(cost[0].size());
  if (rows == 0 || cols == 0) {
    return {0, {}};
  }
  std::vector<int64_t> dp(cols);
  std::vector<std::vector<char>> pred(rows, std::vector<char>(cols));
  dp[0] = cost[0][0];
  for (int r = 0; r < rows; r++) {
    for (int c = 0; c < cols; c++) {
      if (r == 0 && c == 0) {
        continue;
      }
      if (r > 0 && (c == 0 || dp[c] <= dp[c - 1])) {
        dp[c] += cost[r][c];  // Overflow warning.
        pred[r][c] = 'U';
      } else {
        dp[c] = dp[c - 1] + cost[r][c];
        pred[r][c] = 'L';
      }
    }
  }
  // Optional: reconstruct one optimal path.
  std::vector<std::pair<int, int>> path;
  for (int r = rows - 1, c = cols - 1;;) {
    path.emplace_back(r, c);
    if (r == 0 && c == 0) {
      break;
    }
    if (pred[r][c] == 'U') {
      r--;
    } else {
      c--;
    }
  }
  std::reverse(path.begin(), path.end());
  return {dp.back(), path};
}

std::tuple<int, int, int> largest_one_square(const std::vector<std::vector<char>> &a) {
  int rows = static_cast<int>(a.size());
  int cols = a.empty() ? 0 : static_cast<int>(a[0].size());
  std::vector<int> dp(cols);
  int best = 0, top = -1, left = -1;
  for (int r = 0; r < rows; r++) {
    int diagonal = 0;
    for (int c = 0; c < cols; c++) {
      int above = dp[c];
      dp[c] = a[r][c] ? 1 + std::min({above, c == 0 ? 0 : dp[c - 1], diagonal}) : 0;
      diagonal = above;
      if (dp[c] > best) {
        best = dp[c];
        top = r - best + 1;
        left = c - best + 1;
      }
    }
  }
  return {best, top, left};
}
\end{lstlisting}
\Needspace*{4\baselineskip}
\textbf{Example Usage}\par
\begin{lstlisting}[language=C++]
#include <cassert>
using namespace std;

int main() {
  vector<vector<char>> blocked{
      {0, 0, 0},
      {0, 1, 0},
      {0, 0, 0},
  };
  assert(count_grid_paths(blocked) == 2);

  vector<vector<int>> cost{
      {1, 3, 1},
      {1, 5, 1},
      {4, 2, 1},
  };
  auto [min_cost, path] = min_cost_grid_path(cost);
  assert(min_cost == 7 && (path == vector<pair<int, int>>{{0, 0}, {0, 1}, {0, 2}, {1, 2}, {2, 2}}));

  vector<vector<char>> binary{
      {1, 0, 1, 0, 0},
      {1, 0, 1, 1, 1},
      {1, 1, 1, 1, 1},
      {1, 0, 1, 1, 1},
  };
  assert((largest_one_square(binary) == tuple{3, 1, 2}));
  assert((largest_one_square({{0, 0}}) == tuple{0, -1, -1}));
  return 0;
}
\end{lstlisting}

\section{\texorpdfstring{Greedy and Scheduling}{1.5 Greedy and Scheduling}}
\setcounter{section}{5}
\setcounter{subsection}{0}
\subsection{Fractional Knapsack}
Solves the fractional knapsack problem: given items with positive integer weights and nonnegative integer values, choose any fraction of each item to maximize total value without exceeding a capacity limit.

The greedy algorithm processes items in decreasing order of value per unit weight, taking each item in full until only part of the next item fits. This is optimal because replacing any weight taken from a lower-density item with the same weight from a higher-density item cannot decrease the total value. Unlike 0-1 knapsack, allowing fractions makes each such exchange feasible.

\begin{compactitem}
  \item \inlinecode{fractional\_knapsack(weight, value, capacity)} returns a pair (\inlinecode{best\_value}, \inlinecode{fraction}) containing the maximum value and the fraction of each input item taken. Item \inlinecode{i} has weight \inlinecode{weight[i]} and value \inlinecode{value[i]}, \inlinecode{capacity} is the maximum total weight, and \inlinecode{fraction[i]} is in $[0, 1]$.
\end{compactitem}

\textbf{Time Complexity}: $O(n \log n)$ per call due to sorting, where $n$ is the number of items.

\textbf{Space Complexity}: $O(n)$ auxiliary and $O(n)$ for the returned fractions.

\begin{lstlisting}[language=C++]
#include <algorithm>
#include <cassert>
#include <cstdint>
#include <numeric>
#include <utility>
#include <vector>

std::pair<double, std::vector<double>> fractional_knapsack(
    const std::vector<int> &weight, const std::vector<int> &value, int capacity
) {
  int n = static_cast<int>(weight.size());
  assert(static_cast<int>(value.size()) == n && capacity >= 0);
  for (int i = 0; i < n; i++) {
    assert(weight[i] > 0 && value[i] >= 0);
  }
  std::vector<int> order(n);
  std::iota(order.begin(), order.end(), 0);
  std::sort(order.begin(), order.end(), [&](int i, int j) {
    int64_t lhs = static_cast<int64_t>(value[i]) * weight[j];
    int64_t rhs = static_cast<int64_t>(value[j]) * weight[i];
    return lhs != rhs ? lhs > rhs : i < j;
  });
  double best_value = 0;
  std::vector<double> fraction(n);
  for (int i : order) {
    int taken = std::min(capacity, weight[i]);
    fraction[i] = static_cast<double>(taken) / weight[i];
    best_value += fraction[i] * value[i];
    capacity -= taken;
    if (capacity == 0) {
      break;
    }
  }
  return {best_value, fraction};
}
\end{lstlisting}
\Needspace*{4\baselineskip}
\textbf{Example Usage}\par
\begin{lstlisting}[language=C++]
#include <cmath>
using namespace std;

int main() {
  vector<int> weight{10, 20, 30};
  vector<int> value{60, 100, 120};
  auto [best_value, fraction] = fractional_knapsack(weight, value, 50);
  assert(fabs(best_value - 240) < 1e-9 && (fraction == vector<double>{1, 1, 2.0 / 3}));
  return 0;
}
\end{lstlisting}
\subsection{Interval Scheduling}
Selects the maximum number of non-overlapping intervals using the classic earliest-finish-time greedy algorithm. This is the unweighted interval scheduling problem, also known as activity selection: every interval has the same value, so the objective is to choose as many compatible intervals as possible.

The greedy choice is safe because among all intervals that could be chosen first, taking one with the earliest finish time leaves the most room for the remaining intervals. In any optimal solution, the first chosen interval can be exchanged for an earliest-finishing compatible interval without reducing the number of intervals selected. Repeating this argument after each choice proves the greedy algorithm optimal. The weighted version in 1.5.3 needs dynamic programming instead.

Intervals are represented as half-open ranges $[{\inlinecodemath{start}}, {\inlinecodemath{finish}})$, so two intervals are compatible if the next interval's \inlinecode{start} is at least the previous interval's \inlinecode{finish}.

\begin{compactitem}
  \item \inlinecode{select\_intervals(intervals)} returns one maximum-size compatible subset of intervals, in the order they are selected, as original indices into an input vector of \inlinecode{Interval} with fields \inlinecode{start} and \inlinecode{finish}. Every interval must satisfy \inlinecode{start} \textless{} \inlinecode{finish}.
\end{compactitem}

\textbf{Time Complexity}: $O(n \log n)$ per call due to sorting, where $n$ is the number of intervals.

\textbf{Space Complexity}: $O(n)$ auxiliary and $O(n)$ for the returned indices.

\begin{lstlisting}[language=C++]
#include <algorithm>
#include <cassert>
#include <climits>
#include <numeric>
#include <vector>

struct Interval {
  int start, finish;
};

std::vector<int> select_intervals(const std::vector<Interval> &intervals) {
  for (const auto &iv : intervals) {
    assert(iv.start < iv.finish);
  }
  std::vector<int> order(intervals.size());
  std::iota(order.begin(), order.end(), 0);
  std::sort(order.begin(), order.end(), [&](int i, int j) {
    return intervals[i].finish != intervals[j].finish ? intervals[i].finish < intervals[j].finish
                                                      : intervals[i].start < intervals[j].start;
  });
  std::vector<int> selected;
  int last_finish = INT_MIN;
  for (int i : order) {
    const auto &iv = intervals[i];
    if (iv.start >= last_finish) {
      selected.push_back(i);
      last_finish = iv.finish;
    }
  }
  return selected;
}
\end{lstlisting}
\Needspace*{4\baselineskip}
\textbf{Example Usage}\par
\begin{lstlisting}[language=C++]
#include <cassert>
using namespace std;

int main() {
  vector<Interval> intervals{{1, 4}, {3, 5}, {0, 6}, {5, 7}, {8, 9}};
  // Earliest-finish greedy chooses original indices 0, 3, 4.
  assert((select_intervals(intervals) == vector<int>{0, 3, 4}));

  vector<Interval> touching{{0, 2}, {2, 4}, {4, 5}};
  assert((select_intervals(touching) == vector<int>{0, 1, 2}));
  return 0;
}
\end{lstlisting}
\subsection{Weighted Interval Scheduling}
Selects a maximum-weight subset of non-overlapping intervals using dynamic programming and binary search. Unlike unweighted interval scheduling, the earliest finish-time greedy choice is not sufficient when intervals have weights. With intervals sorted by finish time, the best total for the first $i$ intervals either skips interval $i$ or adds its weight to the best total over intervals finishing no later than its start, with that predecessor located by binary search.

Intervals are represented as half-open ranges $[{\inlinecodemath{start}}, {\inlinecodemath{finish}})$, so two intervals are compatible if the next interval's \inlinecode{start} is at least the previous interval's \inlinecode{finish}.

\begin{compactitem}
  \item \inlinecode{select\_weighted\_intervals(intervals)} returns a pair (\inlinecode{weight}, \inlinecode{selected}) containing that maximum weight and the selected intervals as original input indices in execution order, from an input vector of \inlinecode{WeightedInterval} with fields \inlinecode{start}, \inlinecode{finish}, and \inlinecode{weight}. \inlinecode{dp[i]} stores the best answer using the first \inlinecode{i} intervals after sorting by finish time. Each interval must satisfy \inlinecode{start} \textless{} \inlinecode{finish}. The empty subset is allowed, so nonpositive weights need not be selected.
\end{compactitem}

Overflow warning: Accumulated weights must fit in \inlinecode{int64\_t}.

\textbf{Time Complexity}: $O(n \log n)$ per call due to sorting and binary searching compatible intervals.

\textbf{Space Complexity}: $O(n)$ auxiliary and $O(n)$ for the returned indices.

\begin{lstlisting}[language=C++]
#include <algorithm>
#include <cassert>
#include <cstdint>
#include <numeric>
#include <utility>
#include <vector>

struct WeightedInterval {
  int start, finish;
  int64_t weight;
};

std::pair<int64_t, std::vector<int>> select_weighted_intervals(
    const std::vector<WeightedInterval> &intervals
) {
  for (const auto &iv : intervals) {
    assert(iv.start < iv.finish);
  }
  int n = static_cast<int>(intervals.size());
  std::vector<int> order(n), finish(n), prev(n);
  std::iota(order.begin(), order.end(), 0);
  std::sort(order.begin(), order.end(), [&](int i, int j) {
    return intervals[i].finish != intervals[j].finish ? intervals[i].finish < intervals[j].finish
                                                      : intervals[i].start < intervals[j].start;
  });
  for (int i = 0; i < n; i++) {
    finish[i] = intervals[order[i]].finish;
  }
  std::vector<int64_t> dp(n + 1);
  std::vector<char> take(n + 1);
  for (int i = 1; i <= n; i++) {
    int j =
        std::upper_bound(finish.begin(), finish.begin() + i - 1, intervals[order[i - 1]].start) -
        finish.begin();
    prev[i - 1] = j;
    int64_t candidate = dp[j] + intervals[order[i - 1]].weight;  // Overflow warning.
    if (dp[i - 1] < candidate) {
      dp[i] = candidate;
      take[i] = true;
    } else {
      dp[i] = dp[i - 1];
    }
  }
  // Optional: reconstruct one optimal subset of intervals.
  std::vector<int> selected;
  for (int i = n; i > 0;) {
    if (take[i]) {
      selected.push_back(order[i - 1]);
      i = prev[i - 1];
    } else {
      i--;
    }
  }
  std::reverse(selected.begin(), selected.end());
  return {dp[n], selected};
}
\end{lstlisting}
\Needspace*{4\baselineskip}
\textbf{Example Usage}\par
\begin{lstlisting}[language=C++]
#include <cassert>
using namespace std;

int main() {
  vector<WeightedInterval> intervals{{1, 3, 5}, {2, 5, 6}, {4, 6, 5}, {6, 7, 4}, {5, 8, 11}};
  auto [weight, selected] = select_weighted_intervals(intervals);
  // Taking ids 1 and 4 beats the earliest-finish unweighted-looking choices.
  assert(weight == 17 && (selected == vector<int>{1, 4}));

  vector<WeightedInterval> touching{{0, 2, 5}, {2, 4, 6}, {1, 3, 7}};
  auto [touching_weight, touching_selected] = select_weighted_intervals(touching);
  assert(touching_weight == 11 && (touching_selected == vector<int>{0, 1}));
  return 0;
}
\end{lstlisting}
\subsection{Interval Merging and Covering}
Two fundamental interval problems ask either for the union of a collection of intervals or for a minimum-size subcollection that covers a target interval. For the union, sorting by start time makes all intervals that overlap the current merged interval consecutive, so one scan can extend its finish or begin the next disjoint interval.

For minimum covering, maintain the rightmost point covered so far. Among every interval starting at or before that frontier, greedily select one with the farthest finish. Any feasible cover must next use one of those intervals, and replacing that choice with the farthest-reaching one cannot make the remaining target harder to cover. This is different from interval scheduling, which selects the compatible interval with the earliest finish.

Intervals are represented as half-open ranges $[{\inlinecodemath{start}}, {\inlinecodemath{finish}})$. Touching intervals may be merged because their union is another half-open interval with no gap.

\begin{compactitem}
  \item \inlinecode{merge\_intervals(intervals)} returns the union as disjoint intervals sorted by start. Every input interval must satisfy \inlinecode{start} \textless{} \inlinecode{finish}.
  \item \inlinecode{cover\_interval(intervals, lo, hi)} returns a minimum-size list of original input indices whose intervals cover $[{\inlinecodemath{lo}}, {\inlinecodemath{hi}})$, in the order selected, or \inlinecode{std::nullopt} if coverage is impossible. Every input interval must satisfy \inlinecode{start} \textless{} \inlinecode{finish}, and \inlinecode{lo} must not exceed \inlinecode{hi}.
\end{compactitem}

\textbf{Time Complexity}: $O(n \log n)$ per call due to sorting, where $n$ is the number of intervals.

\textbf{Space Complexity}: $O(n)$ auxiliary and $O(n)$ for the returned intervals or indices.

\begin{lstlisting}[language=C++]
#include <algorithm>
#include <cassert>
#include <numeric>
#include <optional>
#include <vector>

struct Interval {
  int start, finish;
};

std::vector<Interval> merge_intervals(std::vector<Interval> intervals) {
  for (const auto &iv : intervals) {
    assert(iv.start < iv.finish);
  }
  std::sort(intervals.begin(), intervals.end(), [](const Interval &a, const Interval &b) {
    return a.start != b.start ? a.start < b.start : a.finish < b.finish;
  });
  std::vector<Interval> merged;
  for (const auto &iv : intervals) {
    if (merged.empty() || merged.back().finish < iv.start) {
      merged.push_back(iv);
    } else {
      merged.back().finish = std::max(merged.back().finish, iv.finish);
    }
  }
  return merged;
}

std::optional<std::vector<int>> cover_interval(
    const std::vector<Interval> &intervals, int lo, int hi
) {
  assert(lo <= hi);
  for (const auto &iv : intervals) {
    assert(iv.start < iv.finish);
  }
  int n = static_cast<int>(intervals.size());
  std::vector<int> order(n);
  std::iota(order.begin(), order.end(), 0);
  std::sort(order.begin(), order.end(), [&](int i, int j) {
    if (intervals[i].start != intervals[j].start) {
      return intervals[i].start < intervals[j].start;
    }
    return intervals[i].finish != intervals[j].finish ? intervals[i].finish > intervals[j].finish
                                                      : i < j;
  });
  std::vector<int> selected;
  int pos = 0, frontier = lo;
  while (frontier < hi) {
    int best = -1, best_finish = frontier;
    while (pos < n && intervals[order[pos]].start <= frontier) {
      int i = order[pos++];
      if (best_finish < intervals[i].finish ||
          (best_finish == intervals[i].finish && best != -1 && i < best)) {
        best = i;
        best_finish = intervals[i].finish;
      }
    }
    if (best == -1) {
      return std::nullopt;
    }
    selected.push_back(best);
    frontier = best_finish;
  }
  return selected;
}
\end{lstlisting}
\Needspace*{4\baselineskip}
\textbf{Example Usage}\par
\begin{lstlisting}[language=C++]
#include <cassert>
#include <vector>
using namespace std;

int main() {
  vector<Interval> intervals{{8, 10}, {1, 3}, {2, 6}, {15, 18}, {10, 12}};
  auto merged = merge_intervals(intervals);
  assert(merged.size() == 3);
  assert(merged[0].start == 1 && merged[0].finish == 6);
  assert(merged[1].start == 8 && merged[1].finish == 12);
  assert(merged[2].start == 15 && merged[2].finish == 18);

  vector<Interval> cover{{0, 2}, {1, 5}, {4, 7}, {6, 10}, {2, 6}, {8, 11}};
  assert(cover_interval(cover, 0, 10) == vector<int>({0, 4, 3}));
  assert(!cover_interval({{0, 2}, {3, 5}}, 0, 5));
  assert(cover_interval({}, 4, 4) == vector<int>());
  return 0;
}
\end{lstlisting}
\subsection{Interval Partitioning}
Assigns intervals to the minimum number of groups so that no two intervals in the same group overlap. This is the classic interval partitioning problem, also known as finding the minimum number of rooms needed for meetings.

The greedy algorithm sorts intervals by start time and reuses the group whose current finish time is earliest whenever possible. Otherwise, it creates a new group. This is optimal because the number of groups must be at least the maximum number of intervals overlapping at any time. When the earliest-finishing active group still overlaps the next interval, every active group overlaps it, so creating a new group is unavoidable. If that group is free, reusing it preserves as many later options as possible.

Intervals are treated as half-open ranges $[{\inlinecodemath{start}}, {\inlinecodemath{finish}})$, so one interval may reuse a room that another interval vacates at the same time.

\begin{compactitem}
  \item \inlinecode{partition\_intervals(intervals)} returns a pair (\inlinecode{rooms}, \inlinecode{room}), where \inlinecode{rooms} is the minimum number of rooms and \inlinecode{room[i]} is the assigned room for input interval \inlinecode{i}. Each \inlinecode{Interval} has fields \inlinecode{start} and \inlinecode{finish}, which must satisfy \inlinecode{start} \textless{} \inlinecode{finish}.
\end{compactitem}

\textbf{Time Complexity}: $O(n \log n)$ per call due to sorting and priority queue operations.

\textbf{Space Complexity}: $O(n)$ auxiliary and $O(n)$ for the returned assignment.

\begin{lstlisting}[language=C++]
#include <algorithm>
#include <cassert>
#include <functional>
#include <numeric>
#include <queue>
#include <utility>
#include <vector>

struct Interval {
  int start, finish;
};

std::pair<int, std::vector<int>> partition_intervals(const std::vector<Interval> &intervals) {
  for (const auto &iv : intervals) {
    assert(iv.start < iv.finish);
  }
  int n = static_cast<int>(intervals.size());
  std::vector<int> order(n);
  std::iota(order.begin(), order.end(), 0);
  std::sort(order.begin(), order.end(), [&](int i, int j) {
    return intervals[i].start != intervals[j].start ? intervals[i].start < intervals[j].start
                                                    : intervals[i].finish < intervals[j].finish;
  });
  std::vector<int> room(n);
  std::priority_queue<std::pair<int, int>, std::vector<std::pair<int, int>>, std::greater<>> pq;
  int rooms = 0;
  for (int i : order) {
    if (!pq.empty() && pq.top().first <= intervals[i].start) {
      int r = pq.top().second;
      pq.pop();
      room[i] = r;
    } else {
      room[i] = rooms++;
    }
    pq.emplace(intervals[i].finish, room[i]);
  }
  return {rooms, room};
}
\end{lstlisting}
\Needspace*{4\baselineskip}
\textbf{Example Usage}\par
\begin{lstlisting}[language=C++]
#include <algorithm>
#include <cassert>
using namespace std;

int main() {
  auto [rooms, room] = partition_intervals(vector<Interval>{{0, 30}, {5, 10}, {15, 20}});
  // The long interval overlaps both short intervals, but the short intervals can share a room.
  assert(rooms == 2);
  assert(room[1] == room[2]);
  assert(room[0] != room[1]);

  vector<Interval> touching{{0, 2}, {2, 4}, {4, 5}};
  auto [touching_rooms, touching_room] = partition_intervals(touching);
  assert(touching_rooms == 1 && *max_element(touching_room.begin(), touching_room.end()) == 0);
  assert(partition_intervals({}).first == 0);
  return 0;
}
\end{lstlisting}
\subsection{Deadline Scheduling}
Schedules unit-time jobs with deadlines and profits to maximize total profit. Each job takes one time slot and must be scheduled at or before its deadline. The greedy algorithm processes jobs in decreasing profit order and places each chosen job in the latest available slot before its deadline.

The latest-slot choice is safe because it uses the current high-profit job while leaving earlier slots available for jobs with tighter deadlines. A disjoint-set forest over time slots makes this efficient: \inlinecode{find(t)} returns the latest still-free slot at or before \inlinecode{t}, and occupying slot \inlinecode{t} links it to the next candidate slot \inlinecode{t - 1}.

\begin{compactitem}
  \item \inlinecode{select\_deadline\_jobs(jobs)} returns a pair (\inlinecode{profit}, \inlinecode{slot}) containing that maximum profit and \inlinecode{slot[i]}, the assigned time slot of input job \inlinecode{i}, or $-1$ if that job is not selected, for a vector of \inlinecode{Job} with fields \inlinecode{deadline} and \inlinecode{profit}. Deadlines are positive integers, and assigned slots are numbered starting from $1$. Jobs with nonpositive profit are not selected.
\end{compactitem}

Overflow warning: The total profit must fit in \inlinecode{int64\_t}.

\textbf{Time Complexity}: $O(n \log n)$ per call due to sorting, with near-constant disjoint-set operations.

\textbf{Space Complexity}: $O(n)$ auxiliary and $O(n)$ for the returned slot assignment.

\begin{lstlisting}[language=C++]
#include <algorithm>
#include <cassert>
#include <cstdint>
#include <numeric>
#include <utility>
#include <vector>

struct Job {
  int deadline;
  int64_t profit;
};

class SlotDSU {
  std::vector<int> root;

 public:
  explicit SlotDSU(int n) : root(n + 1) { std::iota(root.begin(), root.end(), 0); }

  int find(int u) { return root[u] == u ? u : root[u] = find(root[u]); }
  void occupy(int u) { root[u] = find(u - 1); }
};

std::pair<int64_t, std::vector<int>> select_deadline_jobs(const std::vector<Job> &jobs) {
  for (const auto &job : jobs) {
    assert(job.deadline > 0);
  }
  int n = static_cast<int>(jobs.size());
  std::vector<int> order(n);
  std::iota(order.begin(), order.end(), 0);
  std::sort(order.begin(), order.end(), [&](int i, int j) {
    return jobs[i].profit != jobs[j].profit ? jobs[i].profit > jobs[j].profit
                                            : jobs[i].deadline < jobs[j].deadline;
  });
  int max_deadline = 0;
  for (const auto &j : jobs) {
    max_deadline = std::max(max_deadline, j.deadline);
  }
  max_deadline = std::min(max_deadline, n);
  SlotDSU slots(max_deadline);
  std::vector<int> assigned(n, -1);
  int64_t res = 0;
  for (int i : order) {
    if (jobs[i].profit <= 0) {
      continue;
    }
    int slot = slots.find(std::min(jobs[i].deadline, max_deadline));
    if (slot > 0) {
      res += jobs[i].profit;  // Overflow warning.
      assigned[i] = slot;
      slots.occupy(slot);
    }
  }
  return {res, assigned};
}
\end{lstlisting}
\Needspace*{4\baselineskip}
\textbf{Example Usage}\par
\begin{lstlisting}[language=C++]
#include <cassert>
using namespace std;

int main() {
  vector<Job> jobs{{2, 100}, {1, 19}, {2, 27}, {1, 25}, {3, 15}};
  auto [profit, slot] = select_deadline_jobs(jobs);
  // Jobs 0, 2, 4 fill slots 2, 1, 3 for total profit 100 + 27 + 15.
  assert(profit == 142 && (slot == vector<int>{2, -1, 1, -1, 3}));
  return 0;
}
\end{lstlisting}
\subsection{Minimum Late Jobs (Moore-Hodgson)}
Selects a maximum-size subset of jobs that can all be completed by their deadlines. Each job has a processing time and deadline, and all jobs are available at time $0$. The Moore-Hodgson algorithm sorts by deadline and keeps the accepted jobs in a max-heap by processing time; whenever the schedule becomes late, it removes the longest accepted job.

After considering jobs up to some deadline, if the accepted set no longer fits, at least one accepted job must be removed. Removing the longest one leaves the most remaining time while keeping the same number of removed jobs, so it is never worse than removing a shorter accepted job. Repeating this repair after each deadline leaves a largest feasible accepted set.

\begin{compactitem}
  \item \inlinecode{max\_on\_time\_jobs(jobs)} returns the selected jobs as original input indices in a feasible earliest-deadline-first execution order, for an input vector of \inlinecode{TimedJob} with fields \inlinecode{duration} and \inlinecode{deadline}. Durations and deadlines must be nonnegative integers.
\end{compactitem}

Overflow warning: The total duration must fit in \inlinecode{int64\_t}.

\textbf{Time Complexity}: $O(n \log n)$ per call due to sorting and heap operations.

\textbf{Space Complexity}: $O(n)$ auxiliary and $O(n)$ for the returned indices.

\begin{lstlisting}[language=C++]
#include <algorithm>
#include <cassert>
#include <cstdint>
#include <numeric>
#include <queue>
#include <utility>
#include <vector>

struct TimedJob {
  int duration, deadline;
};

std::vector<int> max_on_time_jobs(const std::vector<TimedJob> &jobs) {
  for (const auto &job : jobs) {
    assert(job.duration >= 0 && job.deadline >= 0);
  }
  int n = static_cast<int>(jobs.size());
  std::vector<int> order(n);
  std::iota(order.begin(), order.end(), 0);
  std::sort(order.begin(), order.end(), [&](int i, int j) {
    return jobs[i].deadline != jobs[j].deadline ? jobs[i].deadline < jobs[j].deadline
                                                : jobs[i].duration < jobs[j].duration;
  });
  std::priority_queue<std::pair<int, int>> accepted;
  std::vector<char> selected(n);
  int64_t time = 0;
  for (int i : order) {
    time += jobs[i].duration;  // Overflow warning.
    accepted.emplace(jobs[i].duration, i);
    selected[i] = true;
    if (time > jobs[i].deadline) {
      time -= accepted.top().first;
      selected[accepted.top().second] = false;
      accepted.pop();
    }
  }
  std::vector<int> res;
  for (int i : order) {
    if (selected[i]) {
      res.push_back(i);
    }
  }
  return res;
}
\end{lstlisting}
\Needspace*{4\baselineskip}
\textbf{Example Usage}\par
\begin{lstlisting}[language=C++]
#include <cassert>
using namespace std;

int main() {
  vector<TimedJob> jobs{{3, 4}, {2, 3}, {1, 2}, {2, 7}};
  assert((max_on_time_jobs(jobs) == vector<int>{2, 1, 3}));
  return 0;
}
\end{lstlisting}
\subsection{Weighted Completion Time}
Given jobs that must run one at a time on a single machine, each with a processing time and a weight, order them to minimize the weighted sum of completion times $\sum w_j C_j$, where $C_j$ is the moment job $j$ finishes. Unlike the deadline problems of sections 1.5.6 and 1.5.7, nothing is rejected and no deadline is missed; every order is feasible, and only the total cost differs.

Smith's rule sorts by the ratio of processing time to weight, running the job with the smallest $p_j / w_j$ first. An exchange argument proves it optimal. Consider any schedule that runs job $b$ immediately before job $a$ with $p_a / w_a < p_b / w_b$, and swap the two. Every other job finishes at the same moment, since the pair occupies the same total time, so the change in cost is $p_a w_b - p_b w_a$, which the ratio inequality makes negative. Any schedule violating the ratio order therefore improves under a swap, so an optimal schedule must be sorted by it. With equal weights the rule reduces to shortest processing time first, the classic way to minimize the average completion time.

\begin{compactitem}
  \item \inlinecode{min\_weighted\_completion\_time(jobs)} returns the pair (\inlinecode{cost}, \inlinecode{order}), where \inlinecode{cost} is the minimum weighted sum of completion times and \inlinecode{order} lists the job indices in an optimal sequence. Each job is a pair (\inlinecode{time}, \inlinecode{weight}) with a positive weight and a nonnegative time.
\end{compactitem}

Ties break by index, and the comparison cross-multiplies rather than divides so that it stays exact on integers. Parallel machines are a different problem, NP-hard for general weights.

Overflow warning: the products in the comparator and the running completion time are accumulated in \inlinecode{int64\_t}, so the total processing time multiplied by the total weight must fit in that type.

\textbf{Time Complexity}: $O(n \log n)$ per call, where $n$ is the number of jobs.

\textbf{Space Complexity}: $O(n)$ auxiliary and $O(n)$ for the returned order.

\begin{lstlisting}[language=C++]
#include <algorithm>
#include <cassert>
#include <cstdint>
#include <numeric>
#include <utility>
#include <vector>

std::pair<int64_t, std::vector<int>> min_weighted_completion_time(
    const std::vector<std::pair<int, int>> &jobs
) {
  for (auto [time, weight] : jobs) {
    assert(time >= 0 && weight > 0);
  }
  int n = static_cast<int>(jobs.size());
  std::vector<int> order(n);
  std::iota(order.begin(), order.end(), 0);
  std::sort(order.begin(), order.end(), [&](int a, int b) {
    // Sort by time / weight ascending, cross-multiplied to stay exact.
    int64_t lhs = static_cast<int64_t>(jobs[a].first) * jobs[b].second;
    int64_t rhs = static_cast<int64_t>(jobs[b].first) * jobs[a].second;
    return lhs != rhs ? lhs < rhs : a < b;
  });
  int64_t clock = 0, cost = 0;
  for (int id : order) {
    clock += jobs[id].first;
    cost += clock * jobs[id].second;  // Overflow warning.
  }
  return {cost, order};
}
\end{lstlisting}
\Needspace*{4\baselineskip}
\textbf{Example Usage}\par
\begin{lstlisting}[language=C++]
#include <cassert>
using namespace std;

int main() {
  // (time, weight) pairs. Ratios are 3, 1, and 2, so job 1 runs first, then job 2, then job 0.
  auto [cost, order] = min_weighted_completion_time({{3, 1}, {2, 2}, {4, 2}});
  assert((order == vector<int>{1, 2, 0}));
  assert(cost == 2 * 2 + 6 * 2 + 9 * 1);  // Completion times 2, 6, and 9.

  // Equal weights reduce the rule to shortest processing time first.
  auto [spt_cost, spt_order] = min_weighted_completion_time({{5, 1}, {1, 1}, {3, 1}});
  assert((spt_order == vector<int>{1, 2, 0}));
  assert(spt_cost == 1 + 4 + 9);

  assert(min_weighted_completion_time({}).first == 0);
  return 0;
}
\end{lstlisting}

\section{\texorpdfstring{Streaming and Randomized Algorithms}{1.6 Streaming and Randomized Algorithms}}
\setcounter{section}{6}
\setcounter{subsection}{0}
\subsection{Fisher-Yates Shuffle}
Randomly permutes a range in-place using the Fisher-Yates shuffle. Each possible permutation is produced with equal probability, assuming the random number source is uniform.

The algorithm scans from right to left. At each position $i$, it chooses a random position $j$ from $[0, i]$ and swaps the two elements. This fixes one uniformly random remaining element at a time.

\begin{compactitem}
  \item \inlinecode{fisher\_yates\_shuffle(lo, hi)} randomly shuffles the range $[{\inlinecodemath{lo}}, {\inlinecodemath{hi}})$ using an internal random engine. The range must provide random-access iterators.
\end{compactitem}

\textbf{Time Complexity}: $O(n)$ per call, where $n$ is the distance between \inlinecode{lo} and \inlinecode{hi}.

\textbf{Space Complexity}: $O(1)$ auxiliary.

\begin{lstlisting}[language=C++]
#include <algorithm>
#include <chrono>
#include <random>

template<typename It>
void fisher_yates_shuffle(It lo, It hi) {
  static std::mt19937 rng(std::chrono::steady_clock::now().time_since_epoch().count());
  int n = static_cast<int>(hi - lo);
  for (int i = n - 1; i > 0; i--) {
    std::uniform_int_distribution<int> dist(0, i);
    int j = dist(rng);
    std::iter_swap(lo + i, lo + j);
  }
}
\end{lstlisting}
\Needspace*{4\baselineskip}
\textbf{Example Usage}\par
\begin{lstlisting}[language=C++]
#include <algorithm>
#include <cassert>
#include <vector>
using namespace std;

int main() {
  vector<int> a{1, 2, 3, 4, 5}, original(a);
  fisher_yates_shuffle(a.begin(), a.end());
  sort(a.begin(), a.end());
  assert(a == original);
  return 0;
}
\end{lstlisting}
\subsection{Reservoir Sampling}
Selects uniformly random samples from a stream whose length may be unknown in advance. This is useful when values arrive online or when storing the full dataset is unnecessary.

Reservoir sampling keeps only the chosen sample in memory while processing each element once. The first $k$ arrivals fill the reservoir, and the $i$-th element thereafter is accepted with probability $k/i$, evicting a uniformly chosen occupant. One draw performs both choices: an integer uniform on $[0, i)$ is below $k$ with probability $k/i$, and is then uniform on the $k$ occupants.

The invariant is that after $n \geq k$ arrivals the reservoir is a uniformly random $k$-subset, so each of the $\binom{n}{k}$ possibilities is equally likely. This is certain at $n = k$. Assuming it after $n - 1$, fix a $k$-subset $A$ of the first $n$. If $n \notin A$, the reservoir must already hold $A$ and reject the arrival, with probability $\binom{n-1}{k}^{-1}(n-k)/n$. If $n \in A$, it must hold $A$ with $n$ replaced by one of the $n - k$ outsiders, each then evicted with probability $(k/n)(1/k) = 1/n$, giving $\binom{n-1}{k}^{-1}(n-k)/n$ again. Both equal $\binom{n}{k}^{-1}$, since $\binom{n}{k} = \binom{n-1}{k} \cdot n/(n-k)$. Taking marginals, every element seen is in the sample with probability $k/n$, and \inlinecode{ReservoirSampleOne} is the case $k = 1$.

Each class maintains its reservoir incrementally; call \inlinecode{add(x)} once per stream element in any order, then call \inlinecode{sample()} to retrieve the result.

\begin{compactitem}
  \item \inlinecode{ReservoirSampleOne\textless{}T\textgreater{}()} constructs a single-element sampler.
  \item \inlinecode{ReservoirSampleK\textless{}T\textgreater{}(k)} constructs a \inlinecode{k}-element sampler, allocating the reservoir up front so that no later \inlinecode{add()} reallocates.
  \item \inlinecode{add(x)} incorporates one more stream element.
  \item \inlinecode{sample()} returns the current sample. For \inlinecode{ReservoirSampleOne}, \inlinecode{sample()} requires at least one \inlinecode{add()} call; for \inlinecode{ReservoirSampleK}, it returns the full reservoir, which may contain fewer than \inlinecode{k} elements if the stream was shorter.
  \item \inlinecode{count()} returns the number of stream elements seen so far.
\end{compactitem}

\Needspace*{3\baselineskip}
\textbf{Time Complexity}:
\begin{compactitem}
  \item $O(1)$ per call to \inlinecode{ReservoirSampleOne::add()} and \inlinecode{ReservoirSampleK::add()}.
  \item $O(k)$ per call to the \inlinecode{ReservoirSampleK} constructor.
  \item $O(n)$ to process a range of $n$ elements.
\end{compactitem}

\Needspace*{3\baselineskip}
\textbf{Space Complexity}:
\begin{compactitem}
  \item $O(1)$ storage for \inlinecode{ReservoirSampleOne}.
  \item $O(k)$ storage for \inlinecode{ReservoirSampleK}, allocated on construction.
\end{compactitem}

\begin{lstlisting}[language=C++]
#include <cassert>
#include <chrono>
#include <optional>
#include <random>
#include <vector>

int rand_int(int lo, int hi) {
  static std::mt19937 rng(std::chrono::steady_clock::now().time_since_epoch().count());
  return std::uniform_int_distribution<int>(lo, hi)(rng);
}

template<typename T>
class ReservoirSampleOne {
  std::optional<T> value;
  int seen = 0;

 public:
  void add(const T &x) {
    seen++;
    if (rand_int(1, seen) == 1) {
      value = x;
    }
  }

  const T &sample() const {
    assert(value.has_value());
    return *value;
  }

  int count() const { return seen; }
};

template<typename T>
class ReservoirSampleK {
  int k, seen = 0;
  std::vector<T> reservoir;

 public:
  explicit ReservoirSampleK(int k) : k(k) {
    assert(k >= 0);
    reservoir.reserve(k);
  }

  void add(const T &x) {
    ++seen;
    if (k == 0) {
      return;
    }
    if (static_cast<int>(reservoir.size()) < k) {
      reservoir.push_back(x);
    } else {
      int j = rand_int(0, seen - 1);
      if (j < k) {
        reservoir[j] = x;
      }
    }
  }

  const std::vector<T> &sample() const { return reservoir; }
  int count() const { return seen; }
};
\end{lstlisting}
\Needspace*{4\baselineskip}
\textbf{Example Usage}\par
\begin{lstlisting}[language=C++]
#include <cassert>
using namespace std;

int main() {
  vector<int> a{10, 20, 30, 40, 50};

  // Streaming interface: feed elements one at a time.
  ReservoirSampleOne<int> s1;
  for (int x : a) {
    s1.add(x);
  }
  assert(s1.count() == 5);
  int one = s1.sample();
  assert(one == 10 || one == 20 || one == 30 || one == 40 || one == 50);

  ReservoirSampleK<int> sk(3);
  for (int x : a) {
    sk.add(x);
  }
  assert(sk.count() == 5);
  assert(sk.sample().size() == 3);

  ReservoirSampleK<int> large(10);
  for (int x : a) {
    large.add(x);
  }
  assert(large.count() == 5);
  assert(large.sample().size() == 5);
  return 0;
}
\end{lstlisting}
\subsection{Frequent Elements}
Finds frequent elements in a range using cancellation algorithms that retain only a small set of candidates. Boyer-Moore finds a possible majority element, while Misra-Gries generalizes the same idea to frequencies above any threshold $n/k$.

The Boyer-Moore voting algorithm maintains one candidate and counter, incrementing on a match and decrementing on a mismatch. Each decrement cancels two different values. Removing such pairs cannot eliminate an element occurring more than $\lfloor n/2 \rfloor$ times, so a true majority survives as the final candidate. A second pass verifies that candidate.

Misra-Gries keeps at most $k - 1$ candidates. A tracked value increments its counter, an untracked value claims a free counter, and otherwise every counter is decremented. Each full decrement cancels $k$ distinct occurrences, so any value occurring more than $\lfloor n/k \rfloor$ times must remain a candidate. These candidates require a second pass for exact verification because the algorithm does not retain the stream.

\begin{compactitem}
  \item \inlinecode{majority\_element(lo, hi)} returns an iterator to the first occurrence of the majority element of $[{\inlinecodemath{lo}}, {\inlinecodemath{hi}})$, or \inlinecode{hi} if no majority element exists. The range must provide ForwardIterators, and its value type must support equality.
  \item \inlinecode{frequent\_candidates(lo, hi, k)} returns a hash table of candidate values to their residual counters. \inlinecode{k} must be at least $2$, and the value type must support equality and \inlinecode{std::hash}. Use \inlinecode{k = 2} for the unverified Boyer-Moore candidate phase.
\end{compactitem}

\Needspace*{3\baselineskip}
\textbf{Time Complexity}:
\begin{compactitem}
  \item $O(n)$ per call to \inlinecode{majority\_element(lo, hi)}, where $n$ is the range length.
  \item $O(n)$ expected per call to \inlinecode{frequent\_candidates(lo, hi, k)}: a full-table decrement cancels $k$ occurrences, so all decrement sweeps take $O(n)$ total.
  \item $O(n \cdot k)$ per call to \inlinecode{frequent\_candidates()} in the collision-heavy worst case.
\end{compactitem}

\Needspace*{3\baselineskip}
\textbf{Space Complexity}:
\begin{compactitem}
  \item $O(1)$ auxiliary for \inlinecode{majority\_element()}.
  \item $O(k)$ for the candidates returned by \inlinecode{frequent\_candidates()}.
\end{compactitem}

\begin{lstlisting}[language=C++]
#include <cassert>
#include <iterator>
#include <unordered_map>

template<typename It>
It majority_element(It lo, It hi) {
  int count = 0;
  It candidate = lo;
  for (It it = lo; it != hi; ++it) {
    if (count == 0) {
      candidate = it;
      count = 1;
    } else if (*it == *candidate) {
      count++;
    } else {
      count--;
    }
  }
  // Second pass to verify. If a majority is guaranteed, skip and return candidate directly instead.
  int n = 0;
  count = 0;
  It first = hi;
  for (It it = lo; it != hi; ++it) {
    n++;
    if (*it == *candidate) {
      if (first == hi) {
        first = it;
      }
      count++;
    }
  }
  if (count <= n / 2) {
    return hi;
  }
  return first;
}

template<typename It>
auto frequent_candidates(It lo, It hi, int k) {
  using T = typename std::iterator_traits<It>::value_type;
  assert(k >= 2);
  std::unordered_map<T, int> count;
  for (It it = lo; it != hi; ++it) {
    if (auto found = count.find(*it); found != count.end()) {
      found->second++;
    } else if (static_cast<int>(count.size()) < k - 1) {
      count[*it] = 1;
    } else {
      for (auto jt = count.begin(); jt != count.end();) {
        if (--jt->second == 0) {
          jt = count.erase(jt);
        } else {
          ++jt;
        }
      }
    }
  }
  return count;
}
\end{lstlisting}
\Needspace*{4\baselineskip}
\textbf{Example Usage}\par
\begin{lstlisting}[language=C++]
#include <vector>
using namespace std;

int main() {
  vector<int> a{3, 2, 3, 1, 3};
  assert(majority_element(a.begin(), a.end()) == a.begin());
  vector<int> b{2, 3, 3, 3, 2, 1};
  assert(majority_element(b.begin(), b.end()) == b.end());

  vector<int> c{1, 2, 1, 3, 1, 2, 1, 4, 2, 2, 2};
  auto candidates = frequent_candidates(c.begin(), c.end(), 3);
  assert(candidates.count(1));
  assert(candidates.count(2));
  assert(!candidates.count(3));
  return 0;
}
\end{lstlisting}
\subsection{Online Statistics}
Online stream summaries update useful statistics as values arrive without storing the full stream. This section uses Welford's algorithm to maintain the mean and variance of a stream in constant space, then merges Welford summaries inside a queue to maintain the same statistics over a sliding window. For the median of a stream or of a trailing window, see section 1.6.5.

The naive variance formula $\operatorname{Var}(X) = E[X^2] - E[X]^2$ can lose precision when the two terms are large and close together. Welford instead maintains the current mean $\bar{x}_n$ (stored as \inlinecode{avg}) and the unnormalized second central moment $M_{2,n} = \sum_{i=1}^{n}(x_i - \bar{x}_n)^2$ (stored as \inlinecode{m2}). Each new value shifts the running mean by its deviation divided by the new count, then adds the product of its deviations from the old and new means to \inlinecode{m2}.

\begin{compactitem}
  \item \inlinecode{OnlineStatistics()} constructs an empty summary.
  \item \inlinecode{add(x)} incorporates one more value \inlinecode{x} into the summary.
  \item \inlinecode{count()} returns the number of values seen so far.
  \item \inlinecode{mean()} returns the current arithmetic mean, or $0$ if the summary is empty.
  \item \inlinecode{var\_pop()} returns population variance, dividing by $n$, or $0$ if the summary is empty.
  \item \inlinecode{var\_samp()} returns sample variance, dividing by $n - 1$, or $0$ if fewer than two values have been added.
\end{compactitem}

\textbf{Time Complexity}: $O(1)$ per call to \inlinecode{add()} and each query.

\textbf{Space Complexity}: $O(1)$ auxiliary for all operations.

\begin{lstlisting}[language=C++]
#include <cassert>
#include <utility>
#include <vector>

class OnlineStatistics {
  int n;
  double avg, m2;

 public:
  OnlineStatistics() : n(0), avg(0), m2(0) {}

  void add(double x) {
    n++;
    double delta = x - avg;
    avg += delta / n;
    double delta2 = x - avg;
    m2 += delta * delta2;
  }

  int count() const { return n; }
  double mean() const { return avg; }
  double var_pop() const { return n == 0 ? 0 : m2 / n; }
  double var_samp() const { return n <= 1 ? 0 : m2 / (n - 1); }
};
\end{lstlisting}
To additionally support FIFO removals, Welford's update cannot simply be run backwards to drop the oldest value: reversing it turns \inlinecode{m2} into a difference, which loses the very property that made the forward pass stable, and the resulting error has no way to correct itself as the window advances.

Summaries are mergeable instead. Two summaries of disjoint groups combine by shifting the mean toward the larger group and adding the spread between the two group means, so that $M_2 = M_{2,A} + M_{2,B} + \delta^2 n_A n_B / n$, where $\delta$ is the gap between the group means. That merge is associative, which makes summaries a monoid and puts them within reach of the sliding window aggregation (SWAG) technique of section 1.2.5: a queue of two stacks, each entry storing the fold of everything between it and its end of the queue, reports the fold of the entire queue in $O(1)$ amortized time. The class below is that structure specialized to this summary, so section 1.2.5 is the place to look for the general version over any associative operation, and for why flipping the back stack onto the front stack costs only $O(1)$ amortized per value.

Every value still enters through a forward Welford update of a one-element summary, and the folds only ever add nonnegative terms, so the statistics of a window are as stable as those of a whole stream.

\begin{compactitem}
  \item \inlinecode{StatisticsQueue()} constructs an empty summary.
  \item \inlinecode{add(x)} incorporates \inlinecode{x} as the newest value of the queue.
  \item \inlinecode{remove()} drops the oldest value of the queue, which must be nonempty.
  \item \inlinecode{count()} returns the number of values currently in the queue.
  \item \inlinecode{mean()}, \inlinecode{var\_pop()}, and \inlinecode{var\_samp()} return the same statistics as above, over the values currently in the queue.
\end{compactitem}

The window is left to the caller rather than fixed by the constructor. A trailing window of width $w$ is one call to \inlinecode{add()} followed by a call to \inlinecode{remove()} whenever \inlinecode{count()} exceeds $w$, while a variable-width window from the two-pointer patterns of section 1.2.3 instead removes values until its own condition holds again.

\textbf{Time Complexity}: $O(1)$ amortized per call to \inlinecode{add()} and \inlinecode{remove()}, and $O(1)$ per query.

\textbf{Space Complexity}: $O(n)$ storage, where $n$ is the number of values currently in the queue.

\begin{lstlisting}[language=C++]
class StatisticsQueue {
  struct Summary {
    int n;
    double mean, m2;
  };

  // Chan's parallel merge of two disjoint groups, which never subtracts a value out of m2.
  static Summary merge(const Summary &a, const Summary &b) {
    int n = a.n + b.n;
    double delta = b.mean - a.mean;
    return {n, a.mean + delta * b.n / n, a.m2 + b.m2 + delta * delta * a.n * b.n / n};
  }

  // Each stack entry is (summary of one value, fold), where the fold merges that value with all
  // entries between it and its end of the queue (the front for front_stack, back for back_stack).
  std::vector<std::pair<Summary, Summary>> front_stack, back_stack;

  Summary fold() const {
    if (front_stack.empty()) {
      return back_stack.empty() ? Summary{} : back_stack.back().second;
    }
    if (back_stack.empty()) {
      return front_stack.back().second;
    }
    return merge(front_stack.back().second, back_stack.back().second);
  }

 public:
  void add(double x) {
    Summary value{1, x, 0};
    back_stack.emplace_back(
        value, back_stack.empty() ? value : merge(back_stack.back().second, value)
    );
  }

  void remove() {
    assert(count() > 0);
    if (front_stack.empty()) {
      // Flip the back stack onto the front stack, reversing order so the oldest value ends up on
      // top, and recompute the front folds in queue order.
      while (!back_stack.empty()) {
        Summary oldest = back_stack.back().first;
        back_stack.pop_back();
        front_stack.emplace_back(
            oldest, front_stack.empty() ? oldest : merge(oldest, front_stack.back().second)
        );
      }
    }
    front_stack.pop_back();
  }

  int count() const { return static_cast<int>(front_stack.size() + back_stack.size()); }
  double mean() const { return fold().mean; }

  double var_pop() const {
    Summary s = fold();
    return s.n == 0 ? 0 : s.m2 / s.n;
  }

  double var_samp() const {
    Summary s = fold();
    return s.n <= 1 ? 0 : s.m2 / (s.n - 1);
  }
};
\end{lstlisting}
\Needspace*{4\baselineskip}
\textbf{Example Usage}\par
\begin{lstlisting}[language=C++]
#include <cassert>
#include <cmath>
#include <vector>
using namespace std;

bool EQ(double a, double b) {
  return fabs(a - b) < 1e-9;
}

int main() {
  OnlineStatistics empty;
  assert(empty.count() == 0 && EQ(empty.mean(), 0.0));
  assert(EQ(empty.var_pop(), 0.0) && EQ(empty.var_samp(), 0.0));

  OnlineStatistics singleton;
  singleton.add(7);
  assert(singleton.count() == 1 && EQ(singleton.mean(), 7.0));
  assert(EQ(singleton.var_pop(), 0.0) && EQ(singleton.var_samp(), 0.0));

  OnlineStatistics stats;
  for (int x = 1; x <= 5; x++) {
    stats.add(x);
  }
  assert(stats.count() == 5 && EQ(stats.mean(), 3.0));
  assert(EQ(stats.var_pop(), 2.0) && EQ(stats.var_samp(), 2.5));

  OnlineStatistics shifted;
  shifted.add(1e12 + 1);
  shifted.add(1e12 + 2);
  shifted.add(1e12 + 3);
  assert(EQ(shifted.mean(), 1e12 + 2));
  assert(EQ(shifted.var_pop(), 2.0 / 3) && EQ(shifted.var_samp(), 1.0));

  StatisticsQueue empty_queue;
  assert(empty_queue.count() == 0 && EQ(empty_queue.mean(), 0.0));
  assert(EQ(empty_queue.var_pop(), 0.0) && EQ(empty_queue.var_samp(), 0.0));

  // A trailing window of width 3, maintained by the caller.
  StatisticsQueue window;
  for (int x = 1; x <= 5; x++) {
    window.add(x);
    if (window.count() > 3) {
      window.remove();
    }
  }
  assert(window.count() == 3);  // The window holds 3, 4, and 5.
  assert(EQ(window.mean(), 4.0));
  assert(EQ(window.var_pop(), 2.0 / 3) && EQ(window.var_samp(), 1.0));

  // The same shifted values as above, now with an extra value removed from the front first.
  StatisticsQueue shifted_window;
  shifted_window.add(1e12);
  shifted_window.add(1e12 + 1);
  shifted_window.add(1e12 + 2);
  shifted_window.add(1e12 + 3);
  shifted_window.remove();
  assert(EQ(shifted_window.mean(), 1e12 + 2));
  assert(EQ(shifted_window.var_pop(), 2.0 / 3) && EQ(shifted_window.var_samp(), 1.0));

  // A variable-width window, which a fixed size cannot express: extend while the population
  // variance stays at most 1, shrinking from the front whenever it does not.
  StatisticsQueue spread;
  int longest = 0;
  for (double x : {1.0, 1.0, 2.0, 8.0, 8.0, 9.0}) {
    spread.add(x);
    while (spread.var_pop() > 1) {
      spread.remove();
    }
    if (spread.count() > longest) {
      longest = spread.count();
    }
  }
  assert(longest == 3);  // The values 1, 1, and 2 have population variance 2/9.

  return 0;
}
\end{lstlisting}
\subsection{Online Median}
The median of a stream can be maintained as values arrive, without ever sorting the stream. This section keeps the median of a growing stream in two heaps, and the median of a changing collection in a multiset, the difference between them being whether values ever have to leave again. For the mean and variance of a stream, see section 1.6.4.

Two heaps hold the values in halves, using linear storage. A max-heap stores the lower half and a min-heap stores the upper half, balanced so that the lower heap has either the same number of values as the upper heap or one extra value. Their roots are therefore the middle value or values of the stream.

\begin{compactitem}
  \item \inlinecode{OnlineMedian\textless{}T\textgreater{}()} constructs an empty median tracker for an ordered numeric type \inlinecode{T}.
  \item \inlinecode{empty()} returns whether the tracker has no values.
  \item \inlinecode{count()} returns the number of values in the tracker.
  \item \inlinecode{add(x)} inserts \inlinecode{x} into the tracker.
  \item \inlinecode{lower\_median()} and \inlinecode{upper\_median()} return the lower and the upper of the two middle values, or the middle value itself for an odd count. The tracker must be nonempty.
  \item \inlinecode{median()} returns the middle value for an odd count or the arithmetic mean of the two middle values for an even count. The tracker must be nonempty.
\end{compactitem}

Since \inlinecode{median()} averages in \inlinecode{double}, an even count over a type whose values exceed the exactly representable range of \inlinecode{double} should be averaged by the caller from \inlinecode{lower\_median()} and \inlinecode{upper\_median()} instead.

\textbf{Time Complexity}: $O(\log n)$ per call to \inlinecode{add()} and $O(1)$ per query, where $n$ is the number of values.

\textbf{Space Complexity}: $O(n)$ storage.

\begin{lstlisting}[language=C++]
#include <cassert>
#include <functional>
#include <iterator>
#include <queue>
#include <set>
#include <vector>

template<typename T>
class OnlineMedian {
  std::priority_queue<T> lower;
  std::priority_queue<T, std::vector<T>, std::greater<T>> upper;

 public:
  bool empty() const { return lower.empty(); }
  int count() const { return static_cast<int>(lower.size() + upper.size()); }

  void add(const T &x) {
    if (lower.empty() || x <= lower.top()) {
      lower.push(x);
    } else {
      upper.push(x);
    }
    if (lower.size() < upper.size()) {
      lower.push(upper.top());
      upper.pop();
    } else if (lower.size() > upper.size() + 1) {
      upper.push(lower.top());
      lower.pop();
    }
  }

  const T &lower_median() const {
    assert(!empty());
    return lower.top();
  }

  const T &upper_median() const {
    assert(!empty());
    return lower.size() != upper.size() ? lower.top() : upper.top();
  }

  double median() const {
    return static_cast<double>(lower_median()) / 2 + static_cast<double>(upper_median()) / 2;
  }
};
\end{lstlisting}
To additionally support arbitrary removals, a multiset (balanced binary search tree) holds the values in sorted order while an iterator tracks the lower median. Each insertion or erasure shifts the middle of the multiset by at most one position, so the iterator only ever steps one element in either direction: whether it moves depends on which side of the median the value lies and on the parity of the resulting size. Unlike the two heaps above, and unlike the statistics queue of section 1.6.4 that can only drop its oldest value, a multiset erases any value it holds, so removals may come in any order.

\begin{compactitem}
  \item \inlinecode{DeletableMedian\textless{}T\textgreater{}()} constructs an empty median tracker for an ordered numeric type \inlinecode{T}.
  \item \inlinecode{add(x)} inserts \inlinecode{x} into the median tracker.
  \item \inlinecode{remove(x)} erases one copy of \inlinecode{x}, which must currently be present.
  \item \inlinecode{count()} and \inlinecode{empty()} return the number of values and whether the tracker is empty.
  \item \inlinecode{lower\_median()}, \inlinecode{upper\_median()}, and \inlinecode{median()} answer as above, over the values currently held. The tracker must be nonempty.
\end{compactitem}

A trailing window of width $w$ over an array is one call to \inlinecode{add()} followed by a call to \inlinecode{remove()} on the value leaving the window, and a two-pointer scan from section 1.2.3 removes by its own condition instead. Because removals are unordered, the same structure also serves the add and remove operations of Mo's algorithm in section 2.5.4, which answers median queries over arbitrary ranges offline. Replacing the parked iterator with an order-statistic tree, such as the \inlinecode{\_\_gnu\_pbds} tree of section 8.6, generalizes further to any rank rather than just the middle.

\textbf{Time Complexity}: $O(\log n)$ per call to \inlinecode{add()} and \inlinecode{remove()}, and $O(1)$ per query, where $n$ is the number of values currently held.

\textbf{Space Complexity}: $O(n)$ storage.

\begin{lstlisting}[language=C++]
template<typename T>
class DeletableMedian {
  std::multiset<T> sorted;
  typename std::multiset<T>::iterator mid;

 public:
  bool empty() const { return sorted.empty(); }
  int count() const { return static_cast<int>(sorted.size()); }

  void add(const T &x) {
    if (sorted.empty()) {
      mid = sorted.insert(x);
      return;
    }
    bool below = x < *mid;
    sorted.insert(x);
    if (below) {
      if (sorted.size() % 2 == 0) {
        --mid;
      }
    } else if (sorted.size() % 2 == 1) {
      ++mid;
    }
  }

  void remove(const T &x) {
    assert(!empty());
    if (sorted.size() == 1) {
      assert(x == *mid);
      sorted.clear();  // The iterator is meaningless while empty, and add() resets it.
      return;
    }
    // Erasing the median's own node requires stepping off it first. Any other node holding x
    // lies strictly on one side of the median, so erasing it can never invalidate the iterator.
    auto target = (x == *mid) ? mid : sorted.find(x);
    assert(target != sorted.end());
    if (x <= *mid && sorted.size() % 2 == 0) {
      ++mid;
    } else if (x >= *mid && sorted.size() % 2 == 1) {
      --mid;
    }
    sorted.erase(target);
  }

  const T &lower_median() const {
    assert(!empty());
    return *mid;
  }

  const T &upper_median() const {
    assert(!empty());
    return sorted.size() % 2 == 1 ? *mid : *std::next(mid);
  }

  double median() const {
    return static_cast<double>(lower_median()) / 2 + static_cast<double>(upper_median()) / 2;
  }
};
\end{lstlisting}
\Needspace*{4\baselineskip}
\textbf{Example Usage}\par
\begin{lstlisting}[language=C++]
#include <cmath>
using namespace std;

bool EQ(double a, double b) {
  return fabs(a - b) < 1e-9;
}

int main() {
  OnlineMedian<int> medians;
  assert(medians.empty() && medians.count() == 0);
  medians.add(5);
  assert(EQ(medians.median(), 5));
  medians.add(1);
  assert(medians.lower_median() == 1 && medians.upper_median() == 5);
  assert(EQ(medians.median(), 3));
  medians.add(9);
  assert(medians.lower_median() == 5 && medians.upper_median() == 5);
  medians.add(3);
  assert(medians.lower_median() == 3 && medians.upper_median() == 5);
  assert(EQ(medians.median(), 4));

  // A trailing window of width 3 over an array, maintained by the caller.
  vector<int> a{5, 1, 9, 3, 3};
  DeletableMedian<int> window;
  assert(window.empty() && window.count() == 0);
  for (int hi = 0, lo = 0; hi < static_cast<int>(a.size()); hi++) {
    window.add(a[hi]);
    if (window.count() > 3) {
      window.remove(a[lo++]);
    }
  }
  assert(window.count() == 3);  // The window holds 9, 3, and 3.
  assert(window.lower_median() == 3 && EQ(window.median(), 3));

  // Removals need not follow arrival order, and may target the median's own value.
  DeletableMedian<int> tracker;
  for (int x : {1, 2, 3, 4, 5}) {
    tracker.add(x);
  }
  assert(EQ(tracker.median(), 3));
  tracker.remove(3);  // Erasing the median itself steps the iterator off its own node first.
  assert(tracker.lower_median() == 2 && tracker.upper_median() == 4);
  tracker.remove(5);
  assert(EQ(tracker.median(), 2));

  // Duplicates of the median are erased one copy at a time.
  DeletableMedian<int> repeats;
  for (int x : {7, 7, 7}) {
    repeats.add(x);
  }
  repeats.remove(7);
  assert(repeats.count() == 2 && EQ(repeats.median(), 7));
  repeats.remove(7);
  assert(repeats.count() == 1 && repeats.lower_median() == 7);
  return 0;
}
\end{lstlisting}
\subsection{Probabilistic Sketches}
A sketch answers questions about a stream using far less memory than the stream itself, in exchange for a bounded, one-sided error. This section implements a Bloom filter, which tests set membership, and a count-min sketch, which estimates how often each value occurred. Both spread every value over several independent positions of a fixed-size table, and both derive those positions from one pair of base hashes: the Kirsch-Mitzenmacher technique uses $h_i(x) = h_1(x) + i \cdot h_2(x)$ instead of computing an unrelated hash for every position, which matches the false-positive behavior of truly independent hashes while paying for only two. Both base hashes come from the SplitMix64 mixer of section 3.2.1, and the second is forced odd so that repeated steps never stand still.

A Bloom filter represents a set as a bit array. Inserting a value sets the $k$ bits it hashes to, and a membership test reports that a value is present only if all $k$ of its bits are set. Bits are never cleared, so a value that was inserted always tests positive: the error is one-sided, and only a false positive is possible, produced when unrelated insertions happen to have set all $k$ bits. For $n$ values in $m$ bits, the false-positive probability is minimized by $k = (m/n) \ln 2$, which yields $m = -n \ln p / (\ln 2)^2$ bits for a target rate $p$.

\begin{compactitem}
  \item \inlinecode{BloomFilter\textless{}T\textgreater{}(capacity, false\_positive\_rate)} constructs an empty filter sized for \inlinecode{capacity} insertions at the given target rate, where \inlinecode{capacity} is positive and the rate lies in $(0, 1)$. \inlinecode{BloomFilter\textless{}T, Hash\textgreater{}(capacity, false\_positive\_rate, hash)} instead stores the supplied hasher.
  \item \inlinecode{insert(x)} adds \inlinecode{x} to the filter.
  \item \inlinecode{contains(x)} returns whether \inlinecode{x} may be in the filter. It never returns \inlinecode{false} for a value that was inserted, but may return \inlinecode{true} for one that was not.
  \item \inlinecode{count()} returns the number of insertions performed, counting repeated values separately.
  \item \inlinecode{bit\_count()} and \inlinecode{hash\_count()} return the derived table size $m$ and number of positions $k$.
  \item \inlinecode{false\_positive\_rate()} returns the current estimated rate $(1 - e^{-kn/m})^k$, which rises as the filter fills past its planned capacity.
\end{compactitem}

Deletion is not supported, since clearing a bit could hide an unrelated value that also set it. Use a counting Bloom filter, which replaces each bit with a small counter, when values must leave.

\Needspace*{3\baselineskip}
\textbf{Time Complexity}:
\begin{compactitem}
  \item $O(1)$ per call to \inlinecode{count()}, \inlinecode{bit\_count()}, \inlinecode{hash\_count()}, and \inlinecode{false\_positive\_rate()}.
  \item $O(k)$ per call to \inlinecode{insert()} and \inlinecode{contains()}, where $k$ is \inlinecode{hash\_count()}, plus the cost of one call to the hasher.
  \item $O(m)$ per call to the constructor, where $m$ is \inlinecode{bit\_count()}.
\end{compactitem}

\Needspace*{3\baselineskip}
\textbf{Space Complexity}:
\begin{compactitem}
  \item $O(m)$ for storage of the filter, where $m$ is \inlinecode{bit\_count()}.
  \item $O(1)$ auxiliary for all operations.
\end{compactitem}

\begin{lstlisting}[language=C++]
#include <algorithm>
#include <cassert>
#include <cmath>
#include <cstdint>
#include <functional>
#include <utility>
#include <vector>

uint64_t splitmix64(uint64_t x) {
  x += 0x9e3779b97f4a7c15ULL;
  x = (x ^ (x >> 30)) * 0xbf58476d1ce4e5b9ULL;
  x = (x ^ (x >> 27)) * 0x94d049bb133111ebULL;
  return x ^ (x >> 31);
}

template<typename T, typename Hash>
std::pair<uint64_t, uint64_t> hash_pair(const Hash &hash, const T &x) {
  uint64_t h1 = splitmix64(static_cast<uint64_t>(hash(x)));
  return {h1, splitmix64(h1 + 1) | 1};
}

template<typename T, typename Hash = std::hash<T>>
class BloomFilter {
  std::vector<uint64_t> words;
  uint64_t nbits;
  int hashes, inserted;
  Hash hash;

 public:
  BloomFilter(int capacity, double false_positive_rate, Hash hash = Hash{})
      : nbits(0), hashes(0), inserted(0), hash(hash) {
    assert(capacity > 0 && false_positive_rate > 0 && false_positive_rate < 1);
    double ln2 = std::log(2.0);
    nbits =
        static_cast<uint64_t>(std::ceil(-capacity * std::log(false_positive_rate) / (ln2 * ln2)));
    hashes = std::max(1, static_cast<int>(std::lround(nbits * ln2 / capacity)));
    words.resize((nbits + 63) / 64);
  }

  void insert(const T &x) {
    auto [h1, h2] = hash_pair(hash, x);
    for (int i = 0; i < hashes; i++) {
      uint64_t bit = (h1 + static_cast<uint64_t>(i) * h2) % nbits;
      words[bit >> 6] |= 1ULL << (bit & 63);
    }
    inserted++;
  }

  bool contains(const T &x) const {
    auto [h1, h2] = hash_pair(hash, x);
    for (int i = 0; i < hashes; i++) {
      uint64_t bit = (h1 + static_cast<uint64_t>(i) * h2) % nbits;
      if (((words[bit >> 6] >> (bit & 63)) & 1) == 0) {
        return false;
      }
    }
    return true;
  }

  int count() const { return inserted; }
  uint64_t bit_count() const { return nbits; }
  int hash_count() const { return hashes; }

  double false_positive_rate() const {
    double filled = 1 - std::exp(-1.0 * hashes * inserted / nbits);
    return std::pow(filled, hashes);
  }
};
\end{lstlisting}
A count-min sketch estimates the total weight added under each value. It stores an $R$ by $C$ table of counters and, for each row, adds the weight to the one counter that the value hashes to in that row. Every counter that a value touches therefore holds that value's true total plus the totals of whatever else collided there, so with nonnegative weights each row is an overestimate and the minimum over the rows is the sharpest one available. Choosing $C = \lceil e/\epsilon \rceil$ and $R = \lceil \ln(1/\delta) \rceil$ bounds the overestimate by $\epsilon N$ with probability at least $1 - \delta$, where $N$ is the total weight added.

\begin{compactitem}
  \item \inlinecode{CountMinSketch\textless{}T\textgreater{}(eps, delta)} constructs an empty sketch whose estimates exceed the true count by at most \inlinecode{eps} times the total weight, with probability at least $1 - {\inlinecodemath{delta}}$. Both parameters must lie in $(0, 1)$. \inlinecode{CountMinSketch\textless{}T, Hash\textgreater{}(eps, delta, hash)} instead stores the supplied hasher.
  \item \inlinecode{add(x, c = 1)} adds weight \inlinecode{c} to the count of \inlinecode{x}.
  \item \inlinecode{count(x)} returns the estimated total weight of \inlinecode{x}, which is never an underestimate.
  \item \inlinecode{total()} returns the total weight added over all values.
  \item \inlinecode{num\_rows()} and \inlinecode{num\_cols()} return the derived table dimensions.
\end{compactitem}

The error bound is relative to the total weight of the whole stream, so a sketch estimates heavy values well and light ones poorly. That makes it a natural companion to the deterministic heavy-hitter candidates of section 1.6.3: run both, then use the sketch to rank the candidates. Negative weights break the guarantee, since a row is then no longer an overestimate; deletions require the count-mean-min variant or a conservative sketch.

\Needspace*{3\baselineskip}
\textbf{Time Complexity}:
\begin{compactitem}
  \item $O(R)$ per call to \inlinecode{add()} and \inlinecode{count()}, plus the cost of one call to the hasher.
  \item $O(R \cdot C)$ per call to the constructor.
  \item $O(1)$ per call to \inlinecode{total()}, \inlinecode{num\_rows()}, and \inlinecode{num\_cols()}.
\end{compactitem}

\Needspace*{3\baselineskip}
\textbf{Space Complexity}:
\begin{compactitem}
  \item $O(w \cdot d)$ for storage of the sketch.
  \item $O(1)$ auxiliary for all operations.
\end{compactitem}

\begin{lstlisting}[language=C++]
template<typename T, typename Hash = std::hash<T>>
class CountMinSketch {
  std::vector<std::vector<int64_t>> table;
  int rows, cols;
  int64_t weight;
  Hash hash;

 public:
  CountMinSketch(double eps, double delta, Hash hash = Hash{}) : weight(0), hash(hash) {
    assert(eps > 0 && eps < 1 && delta > 0 && delta < 1);
    rows = std::max(1, static_cast<int>(std::ceil(std::log(1 / delta))));
    cols = static_cast<int>(std::ceil(std::exp(1.0) / eps));
    table.assign(rows, std::vector<int64_t>(cols));
  }

  void add(const T &x, int64_t c = 1) {
    auto [h1, h2] = hash_pair(hash, x);
    for (int r = 0; r < rows; r++) {
      table[r][(h1 + static_cast<uint64_t>(r) * h2) % cols] += c;
    }
    weight += c;  // Overflow warning.
  }

  int64_t count(const T &x) const {
    auto [h1, h2] = hash_pair(hash, x);
    int64_t res = INT64_MAX;
    for (int r = 0; r < rows; r++) {
      res = std::min(res, table[r][(h1 + static_cast<uint64_t>(r) * h2) % cols]);
    }
    return res;
  }

  int64_t total() const { return weight; }
  int num_rows() const { return rows; }
  int num_cols() const { return cols; }
};
\end{lstlisting}
\Needspace*{4\baselineskip}
\textbf{Example Usage}\par
\begin{lstlisting}[language=C++]
#include <string>
using namespace std;

int main() {
  BloomFilter<string> filter(1000, 0.01);
  assert(filter.bit_count() == 9586 && filter.hash_count() == 7);
  for (int i = 0; i < 1000; i++) {
    filter.insert("key" + to_string(i));
  }
  assert(filter.count() == 1000);
  for (int i = 0; i < 1000; i++) {
    assert(filter.contains("key" + to_string(i)));  // Inserted values never test negative.
  }
  int false_positives = 0;
  for (int i = 1000; i < 2000; i++) {
    false_positives += filter.contains("key" + to_string(i));
  }
  assert(false_positives < 40);  // Absent values may test positive, about 10 of these 1000.
  assert(filter.false_positive_rate() < 0.02);

  CountMinSketch<string> sketch(0.01, 0.01);
  sketch.add("apple", 3);
  sketch.add("banana");
  assert(sketch.total() == 4);
  assert(sketch.count("apple") == 3 && sketch.count("banana") == 1);
  // A value that was never added reads 0 unless it collides in every row, and an estimate is never
  // lower than the truth.
  assert(sketch.count("cherry") <= sketch.total());
  return 0;
}
\end{lstlisting}

\section{\texorpdfstring{Bit Manipulation}{1.7 Bit Manipulation}}
\setcounter{section}{7}
\setcounter{subsection}{0}
\subsection{Bit Operations and Masks}
Integers are all fixed-size sets of bits, where the $i$-th bit (counting from the least significant bit where $i = 0$) is ``present'' when set to $1$. Treating an \inlinecode{int} as a bitmask makes set membership, insertion, deletion, and iteration into single machine instructions, which is why bitmasks are the backbone of subset dynamic programming and many low-level tricks. The helpers below operate on \inlinecode{Mask}, which is \inlinecode{uint32\_t} by default; change its alias to \inlinecode{uint64\_t} to use 64-bit masks. The loop implementations below exist for education only; production code should prefer the matching GCC built-ins when available, or on C++20 systems the generic \inlinecode{constexpr} equivalents from the \inlinecode{\textless{}bit\textgreater{}} header, which are additionally well-defined at $0$.

\begin{compactitem}
  \item \inlinecode{test\_bit(x, i)} returns whether the \inlinecode{i}-th bit of \inlinecode{x} is set, where $0 \leq {\inlinecodemath{i}} < {\inlinecodemath{MASK\_BITS}}$.
  \item \inlinecode{set\_bit(x, i)}, \inlinecode{clear\_bit(x, i)}, and \inlinecode{toggle\_bit(x, i)} return \inlinecode{x} with the \inlinecode{i}-th respectively forced to $1$, forced to $0$, or flipped, where $0 \leq {\inlinecodemath{i}} < {\inlinecodemath{MASK\_BITS}}$.
  \item \inlinecode{lowest\_set\_bit(x)} returns the value of the lowest set bit of \inlinecode{x} (a power of two), or $0$ if \inlinecode{x} is $0$. \inlinecode{clear\_lowest\_set\_bit(x)} returns \inlinecode{x} with its lowest set bit removed.
  \item \inlinecode{popcount(x)} returns the number of set bits, analogous to C++20's \inlinecode{std::popcount()}.
  \item \inlinecode{parity(x)} returns $1$ if the number of set bits is odd and $0$ otherwise.
  \item \inlinecode{ctz(x)} returns the number of trailing 0-bits for \inlinecode{x \textgreater{} 0}, analogous to C++20's \inlinecode{std::countr\_zero()}.
  \item \inlinecode{ffs(x)} returns one plus the position of the lowest 1-bit of \inlinecode{x}, or $0$ if \inlinecode{x} is $0$.
  \item \inlinecode{clz(x)} returns the number of leading 0-bits for \inlinecode{x \textgreater{} 0}, analogous to C++20's \inlinecode{std::countl\_zero()}.
  \item \inlinecode{is\_pow2(x)} returns whether \inlinecode{x} has exactly one set bit, analogous to C++20's \inlinecode{std::has\_single\_bit()}.
  \item \inlinecode{floor\_pow2(x)} returns the largest power of two that is $\leq$ \inlinecode{x} (for \inlinecode{x} $> 0$), analogous to C++20's \inlinecode{std::bit\_floor()}.
  \item \inlinecode{ceil\_pow2(x)} returns the smallest power of two that is $\geq$ \inlinecode{x}, for $0 < {\inlinecodemath{x}} \leq 2^{b - 1}$ where $b$ is \inlinecode{MASK\_BITS}, analogous to C++20's \inlinecode{std::bit\_ceil()}.
  \item \inlinecode{for\_each\_set\_bit(x, f)} calls \inlinecode{f(i)} once for each set bit position \inlinecode{i} of \inlinecode{x}, in increasing order.
\end{compactitem}

\Needspace*{3\baselineskip}
\textbf{Time Complexity}:
\begin{compactitem}
  \item $O(b)$ worst case per call to \inlinecode{popcount()}, \inlinecode{parity()}, \inlinecode{ctz()}, \inlinecode{ffs()}, \inlinecode{clz()}, and \inlinecode{for\_each\_set\_bit()}, where $b$ is \inlinecode{MASK\_BITS}.
  \item $O(1)$ per call to every other function.
\end{compactitem}

\textbf{Space Complexity}: $O(1)$ auxiliary for all operations.

\begin{lstlisting}[language=C++]
#include <cassert>
#include <climits>
#include <cstdint>

using Mask = uint32_t;
const int MASK_BITS = sizeof(Mask) * CHAR_BIT;

bool test_bit(Mask x, int i) { return (x >> i) & 1; }
Mask set_bit(Mask x, int i) { return x | (Mask{1} << i); }
Mask clear_bit(Mask x, int i) { return x & ~(Mask{1} << i); }
Mask toggle_bit(Mask x, int i) { return x ^ (Mask{1} << i); }
Mask lowest_set_bit(Mask x) { return x & -x; }
Mask clear_lowest_set_bit(Mask x) { return x & (x - 1); }

// popcount(), parity(), ctz(), ffs(), and clz() are spelled out here for educational purposes.

int popcount(Mask x) {  // std::popcount(x) in C++20.
  int count = 0;
  for (; x != 0; x = clear_lowest_set_bit(x)) {
    count++;
  }
  return count;
}

int parity(Mask x) {
  return popcount(x) & 1;
}

int ctz(Mask x) {  // std::countr_zero(x) in C++20.
  assert(x != 0);
  int count = 0;
  for (; (x & 1) == 0; x >>= 1) {
    count++;
  }
  return count;
}

int ffs(Mask x) {
  return x == 0 ? 0 : ctz(x) + 1;
}

int clz(Mask x) {  // std::countl_zero(x) in C++20.
  assert(x != 0);
  int count = 0;
  for (Mask mask = Mask{1} << (MASK_BITS - 1); (x & mask) == 0; mask >>= 1) {
    count++;
  }
  return count;
}

bool is_pow2(Mask x) {  // std::has_single_bit(x) in C++20.
  return x != 0 && (x & (x - 1)) == 0;
}

Mask floor_pow2(Mask x) {  // std::bit_floor(x) in C++20.
  assert(x != 0);
  return Mask{1} << (MASK_BITS - 1 - clz(x));
}

Mask ceil_pow2(Mask x) {  // std::bit_ceil(x) in C++20.
  assert(0 < x && x <= (Mask{1} << (MASK_BITS - 1)));
  return is_pow2(x) ? x : Mask{1} << (MASK_BITS - clz(x));
}

template<typename Fn>
void for_each_set_bit(Mask x, Fn f) {
  while (x != 0) {
    f(ctz(x));
    x = clear_lowest_set_bit(x);
  }
}
\end{lstlisting}
\Needspace*{4\baselineskip}
\textbf{Example Usage}\par
\begin{lstlisting}[language=C++]
#include <cassert>
#include <vector>
using namespace std;

int main() {
  Mask x = 0b101100;
  assert(test_bit(x, 2) == true);
  assert(test_bit(x, 0) == false);
  assert(set_bit(x, 0) == 0b101101U);
  assert(clear_bit(x, 2) == 0b101000U);
  assert(toggle_bit(x, 3) == 0b100100U);

  assert(lowest_set_bit(x) == 0b100U);
  assert(clear_lowest_set_bit(x) == 0b101000U);
  assert(lowest_set_bit(0) == 0U);
  assert(clear_lowest_set_bit(0) == 0U);
  assert(popcount(x) == 3);
  assert(popcount(0) == 0);
  assert(parity(x) == 1);
  assert(parity(0b101101U) == 0);
  assert(parity(0) == 0);
  assert(ctz(x) == 2);
  assert(ffs(x) == 3);
  assert(ffs(0) == 0);
  assert(clz(x) == MASK_BITS - 6);

  assert(is_pow2(16) == true);
  assert(is_pow2(24) == false);
  assert(floor_pow2(20) == 16U);
  assert(ceil_pow2(20) == 32U);
  assert(ceil_pow2(16) == 16U);
  assert(ceil_pow2((Mask{1} << (MASK_BITS - 1)) - 1) == (Mask{1} << (MASK_BITS - 1)));

  vector<int> bits;
  for_each_set_bit(x, [&](int b) { bits.push_back(b); });
  assert((bits == vector<int>{2, 3, 5}));
  return 0;
}
\end{lstlisting}
\subsection{Subset and Submask Iteration}
Many bitmask algorithms must visit related masks in a structured order: every submask of a fixed mask, every mask with a fixed number of set bits, or every mask while accumulating contributions from its submasks. Each pattern has a short, well-known idiom.

The classic submask loop \inlinecode{for (int s = m; s \textgreater{} 0; s = (s - 1) \& m)} walks every nonzero submask of \inlinecode{m} in decreasing order; subtracting 1 borrows through the zero bits of \inlinecode{m}, and the \inlinecode{\& m} masks them back off. Summed over all masks \inlinecode{m} of \inlinecode{n} bits, the total work is $\sum_k \binom{n}{k} 2^k = 3^n$, since each of the \inlinecode{n} bit positions is independently absent, present in \inlinecode{m} only, or present in both \inlinecode{m} and \inlinecode{s}.

Gosper's hack advances to the next mask with the same popcount. It isolates the lowest set bit and adds it to move the lowest movable 1 one position left, clearing the block of 1-bits beneath it. The remaining expression counts those cleared bits and packs them into the least-significant positions, producing the smallest larger integer with the same number of set bits.

\begin{compactitem}
  \item \inlinecode{submasks(m)} returns all submasks of \inlinecode{m}, including \inlinecode{m} itself and 0, in decreasing order.
  \item \inlinecode{next\_same\_popcount(x)} returns the smallest integer greater than \inlinecode{x} with the same number of set bits (Gosper's hack), for \inlinecode{x \textgreater{} 0} when that next mask is representable in \inlinecode{Mask}.
  \item \inlinecode{masks\_with\_popcount(n, k)} returns all \inlinecode{k}-bit subsets of an \inlinecode{n}-bit universe as masks, in increasing order, where $0 \leq k \leq n < {\inlinecodemath{MASK\_BITS}}$.
  \item \inlinecode{subset\_sum\_transform(f)} overwrites \inlinecode{f} (indexed by mask over \inlinecode{n} bits) so that \inlinecode{f[m]} becomes the sum of the original \inlinecode{f[s]} over all submasks \inlinecode{s} of \inlinecode{m}. This ``sum over subsets'' (SOS) DP is the bitmask analog of a prefix sum, accumulating one bit dimension at a time. The value type must support \inlinecode{+=}.
\end{compactitem}

Overflow warning: All intermediate sums must be representable in the value type.

\Needspace*{3\baselineskip}
\textbf{Time Complexity}:
\begin{compactitem}
  \item $O(2^{p})$ per call to \inlinecode{submasks(m)}, where $p$ is \inlinecode{popcount(m)}.
  \item $O(1)$ per call to \inlinecode{next\_same\_popcount()}.
  \item $O(\binom{n}{k})$ per call to \inlinecode{masks\_with\_popcount(n, k)}.
  \item $O(n \cdot 2^{n})$ per call to \inlinecode{subset\_sum\_transform(f)}, where \inlinecode{f} has $2^n$ entries.
\end{compactitem}

\Needspace*{3\baselineskip}
\textbf{Space Complexity}:
\begin{compactitem}
  \item $O(1)$ auxiliary for \inlinecode{next\_same\_popcount()} and \inlinecode{subset\_sum\_transform()}.
  \item $O(1)$ auxiliary and $O(s)$ for the returned masks from \inlinecode{submasks(m)} and \inlinecode{masks\_with\_popcount(n, k)}, where $s$ is the result size.
\end{compactitem}

\begin{lstlisting}[language=C++]
#include <cassert>
#include <climits>
#include <cstdint>
#include <vector>

using Mask = uint32_t;
const int MASK_BITS = sizeof(Mask) * CHAR_BIT;

std::vector<Mask> submasks(Mask m) {
  std::vector<Mask> res;
  Mask s = m;
  while (true) {
    res.push_back(s);
    if (s == 0) {
      break;
    }
    s = (s - 1) & m;
  }
  return res;
}

Mask next_same_popcount(Mask x) {
  assert(x != 0);
  Mask c = x & (Mask{0} - x), r = x + c;
  assert(r != 0);
  return r | (((x ^ r) >> 2) / c);
}

std::vector<Mask> masks_with_popcount(int n, int k) {
  assert(0 <= k && k <= n && n < MASK_BITS);
  std::vector<Mask> res;
  if (k == 0) {
    return {0};
  }
  Mask limit = Mask{1} << n;
  for (Mask x = (Mask{1} << k) - 1; x < limit; x = next_same_popcount(x)) {
    res.push_back(x);
  }
  return res;
}

template<typename T>
void subset_sum_transform(std::vector<T> &f) {
  int size = static_cast<int>(f.size());
  assert(size > 0 && (size & (size - 1)) == 0);
  int n = __builtin_ctz(static_cast<unsigned>(size));  // size = 2^n.
  for (int b = 0; b < n; b++) {
    for (int m = 0; m < size; m++) {
      if (m & (1U << b)) {
        f[m] += f[m ^ (1U << b)];
      }
    }
  }
}
\end{lstlisting}
\Needspace*{4\baselineskip}
\textbf{Example Usage}\par
\begin{lstlisting}[language=C++]
using namespace std;

int main() {
  assert((submasks(0b101) == vector<Mask>{0b101, 0b100, 0b001, 0b000}));

  assert(next_same_popcount(0b0110) == 0b1001U);
  assert(
      (masks_with_popcount(4, 2) == vector<Mask>{0b0011, 0b0101, 0b0110, 0b1001, 0b1010, 0b1100})
  );

  // f[m] = 1 for every mask; after the transform f[m] counts submasks of m = 2^popcount(m).
  vector<int> f(1 << 3, 1);
  subset_sum_transform(f);
  assert(f[0b000] == 1);
  assert(f[0b101] == 4);
  assert(f[0b111] == 8);
  return 0;
}
\end{lstlisting}
\subsection{Bitmask DP}
Bitmask dynamic programming uses a bitmask to represent a small set or frontier, then treats that mask as the DP state. This is useful when one dimension is small enough for $2^n$ states.

\begin{lstlisting}[language=C++]
#include <algorithm>
#include <cassert>
#include <cstdint>
#include <utility>
#include <vector>
\end{lstlisting}
The assignment problem matches workers one-to-one with jobs while minimizing the total cost. The workers are processed in a fixed order: \inlinecode{dp[mask]} is the minimum cost after assigning the first \inlinecode{popcount(mask)} workers to the selected jobs, and the next worker tries every unused job.

\begin{compactitem}
  \item \inlinecode{min\_cost\_assignment(cost)} returns a pair (\inlinecode{sum}, \inlinecode{job}), where \inlinecode{sum} is the minimum cost of assigning each worker to a distinct job, and \inlinecode{job[i]} is the job assigned to worker \inlinecode{i}. The input is a square matrix with one row per worker and one column per job.
\end{compactitem}

Overflow warning: The accumulated cost must fit in \inlinecode{int64\_t}.

\textbf{Time Complexity}: $O(n \cdot 2^{n})$ per call to \inlinecode{min\_cost\_assignment()}, where $n$ is the number of workers and jobs.

\textbf{Space Complexity}: $O(2^{n})$ auxiliary and $O(n)$ for the returned assignment from \inlinecode{min\_cost\_assignment()}.

\begin{lstlisting}[language=C++]
std::pair<int64_t, std::vector<int>> min_cost_assignment(
    const std::vector<std::vector<int64_t>> &cost
) {
  int n = static_cast<int>(cost.size());
  assert(n < 31);
  int states = 1 << n;
  std::vector<int64_t> dp(states);
  std::vector<int> parent(states, -1);
  for (int mask = 0; mask < states; mask++) {
    int worker = __builtin_popcount(static_cast<unsigned>(mask));
    if (worker == n) {
      continue;
    }
    for (int job = 0; job < n; job++) {
      if ((mask & (1 << job)) == 0) {
        int next = mask | (1 << job);
        int64_t candidate = dp[mask] + cost[worker][job];  // Overflow warning.
        if (parent[next] == -1 || dp[next] > candidate) {
          dp[next] = candidate;
          parent[next] = job;
        }
      }
    }
  }
  // Optional: reconstruct one optimal assignment.
  std::vector<int> job(n);
  for (int mask = states - 1, worker = n - 1; worker >= 0; worker--) {
    job[worker] = parent[mask];
    mask ^= 1 << job[worker];
  }
  return {dp[states - 1], job};
}
\end{lstlisting}
The minimum set cover problem asks for the fewest input sets whose union contains every element of a fixed universe. Both input sets and covered subsets are represented as masks: \inlinecode{dp[mask]} is the minimum number of chosen sets needed to cover exactly the elements in \inlinecode{mask}, and adding one set moves to the union mask.

\begin{compactitem}
  \item \inlinecode{set\_cover(sets, universe\_size)} returns a pair (\inlinecode{count}, \inlinecode{chosen}), where \inlinecode{count} is the minimum number of sets needed to cover all elements in $[0, u)$, where $u$ is \inlinecode{universe\_size}, and \inlinecode{chosen} contains one optimal list of set indices. If no cover exists, \inlinecode{count} is $-1$ and \inlinecode{chosen} is empty. Each input set is represented as a bitmask.
\end{compactitem}

\textbf{Time Complexity}: $O(s \cdot 2^{u})$ per call, where $s$ is the number of sets and $u$ is \inlinecode{universe\_size}.

\textbf{Space Complexity}: $O(2^{u})$ auxiliary and $O(s)$ for the returned indices.

\begin{lstlisting}[language=C++]
std::pair<int, std::vector<int>> set_cover(const std::vector<int> &sets, int universe_size) {
  assert(0 <= universe_size && universe_size < 31);
  int states = 1 << universe_size;
  int full = states - 1;
  for (int subset : sets) {
    assert((subset & ~full) == 0);
  }
  int n = static_cast<int>(sets.size());
  int inf = n + 1;
  std::vector<int> dp(states, inf), parent(states, -1), prev(states, -1);
  dp[0] = 0;
  for (int mask = 0; mask < states; mask++) {
    if (dp[mask] == inf) {
      continue;
    }
    for (int i = 0; i < n; i++) {
      int next = mask | sets[i];
      if (dp[next] > dp[mask] + 1) {
        dp[next] = dp[mask] + 1;
        parent[next] = i;
        prev[next] = mask;
      }
    }
  }
  if (dp[full] == inf) {
    return {-1, {}};
  }
  // Optional: reconstruct one minimum set cover.
  std::vector<int> chosen;
  for (int mask = full; mask != 0; mask = prev[mask]) {
    int i = parent[mask];
    chosen.push_back(i);
  }
  std::reverse(chosen.begin(), chosen.end());
  return {dp[full], chosen};
}
\end{lstlisting}
The minimum-cost set partition problem divides all elements into disjoint nonempty groups when the cost of every possible group is known. For each set \inlinecode{mask}, choosing a submask as one group leaves \inlinecode{mask \textasciicircum{} sub} to be partitioned, giving the recurrence \inlinecode{dp[mask] = min(dp[mask \textasciicircum{} sub] + cost[sub])}. Over all masks, this has $O(3^{n})$ transitions.

\begin{compactitem}
  \item \inlinecode{min\_cost\_partition(group\_cost)} returns the minimum total cost to partition all elements into disjoint nonempty groups, where \inlinecode{group\_cost[mask]} is the cost of taking \inlinecode{mask} as one group.
\end{compactitem}

Overflow warning: The accumulated cost must fit in \inlinecode{int64\_t}.

\textbf{Time Complexity}: $O(3^{n})$ per call, where $n$ is the number of elements.

\textbf{Space Complexity}: $O(2^{n})$ auxiliary.

\begin{lstlisting}[language=C++]
int64_t min_cost_partition(const std::vector<int64_t> &group_cost) {
  int states = static_cast<int>(group_cost.size());
  assert(states > 0 && (states & (states - 1)) == 0);
  std::vector<int64_t> dp(states);
  for (int mask = 1; mask < states; mask++) {
    dp[mask] = group_cost[mask];
    for (int sub = (mask - 1) & mask; sub > 0; sub = (sub - 1) & mask) {
      int64_t candidate = dp[mask ^ sub] + group_cost[sub];  // Overflow warning.
      dp[mask] = std::min(dp[mask], candidate);
    }
  }
  return dp[states - 1];
}
\end{lstlisting}
Domino tiling asks for the number of ways to cover a rectangle with horizontal ($1 \times 2$) and vertical ($2 \times 1$) dominoes. Processing cells in row-major order, a profile mask records which cells on the current frontier are already occupied. An occupied current cell is cleared; otherwise, placing a horizontal domino marks the next column, while placing a vertical one carries the current column into the next row. This technique is often called DP on a broken profile or plug DP.

\begin{compactitem}
  \item \inlinecode{count\_domino\_tilings(rows, cols)} returns the number of ways to tile a \inlinecode{rows} by \inlinecode{cols} rectangle with $1 \times 2$ and $2 \times 1$ dominoes.
\end{compactitem}

Overflow warning: The tiling count must fit in \inlinecode{int64\_t}.

\textbf{Time Complexity}: $O(R \cdot C \cdot 2^{C})$ per call, where $R$ and $C$ are the number of rows and columns, respectively.

\textbf{Space Complexity}: $O(2^{C})$ auxiliary.

\begin{lstlisting}[language=C++]
int64_t count_domino_tilings(int rows, int cols) {
  assert(rows >= 0 && cols >= 0);
  if (rows == 0 || cols == 0) {
    return 1;
  }
  if (rows < cols) {
    std::swap(rows, cols);
  }
  assert(cols < 31);
  int states = 1 << cols;
  std::vector<int64_t> dp(states), next(states);
  dp[0] = 1;
  for (int r = 0; r < rows; r++) {
    for (int c = 0; c < cols; c++) {
      next.assign(states, 0);
      for (int mask = 0; mask < states; mask++) {
        if (dp[mask] == 0) {
          continue;
        }
        if (mask & (1 << c)) {
          next[mask ^ (1 << c)] += dp[mask];  // Overflow warning.
        } else {
          if (c + 1 < cols && (mask & (1 << (c + 1))) == 0) {
            next[mask | (1 << (c + 1))] += dp[mask];
          }
          if (r + 1 < rows) {
            next[mask | (1 << c)] += dp[mask];
          }
        }
      }
      dp.swap(next);
    }
  }
  return dp[0];
}
\end{lstlisting}
\Needspace*{4\baselineskip}
\textbf{Example Usage}\par
\begin{lstlisting}[language=C++]
#include <cassert>
using namespace std;

int main() {
  vector<vector<int64_t>> cost{{9, 2, 7}, {6, 4, 3}, {5, 8, 1}};
  auto [assignment_cost, job] = min_cost_assignment(cost);
  assert(assignment_cost == 9 && (job == vector<int>{1, 0, 2}));

  auto [cover_count, chosen] = set_cover(vector<int>{0b0011, 0b0110, 0b1100, 0b1000}, 4);
  assert(cover_count == 2 && (chosen == vector<int>{0, 2}));

  auto [impossible_count, impossible_chosen] = set_cover(vector<int>{0b001, 0b010}, 3);
  assert(impossible_count == -1 && impossible_chosen.empty());

  vector<int64_t> group_cost{
      0,        // Empty group is unused.
      4, 6, 7,  // Singletons.
      6, 2, 3,  // Pairs.
      9,        // All three together.
  };
  assert(min_cost_partition(group_cost) == 7);  // Groups {0} and {1, 2}.

  assert(count_domino_tilings(0, 1000000000) == 1);
  assert(count_domino_tilings(2, 3) == 3);
  assert(count_domino_tilings(3, 3) == 0);
  assert(count_domino_tilings(4, 4) == 36);
  return 0;
}
\end{lstlisting}
\subsection{Gray Code}
The binary reflected Gray code is an ordering of all $2^n$ bitmasks in which consecutive masks differ in exactly one bit. It is useful whenever the cost of moving between adjacent states depends on a single bit flip: enumerating subsets while incrementally maintaining a value, hardware where only one bit may toggle at a time, and constructions like Hamiltonian paths on the hypercube.

The Gray code at 0-based rank \inlinecode{k} is obtained from the ordinary binary value \inlinecode{k} by \inlinecode{k \textasciicircum{} (k \textgreater{}\textgreater{} 1)}. Incrementing \inlinecode{k} flips one binary bit and every lower bit, but XORing each bit with its higher neighbor cancels all but the highest of those changes, so consecutive Gray codes differ in exactly one bit. The map is a bijection: each binary bit is the prefix XOR of the Gray bits at and above it. The inverse loop computes those prefix XORs in doubling steps to recover the original rank.

\begin{compactitem}
  \item \inlinecode{gray\_code(k)} returns the mask at 0-based rank \inlinecode{k} in Gray code order.
  \item \inlinecode{inverse\_gray\_code(g)} returns the 0-based rank \inlinecode{k} such that \inlinecode{gray\_code(k)} equals \inlinecode{g}.
  \item \inlinecode{gray\_sequence(n)} returns all $2^n$ masks in Gray code order; consecutive entries (and the last with the first when \inlinecode{n \textgreater{} 0}) differ in exactly one bit, where $0 \leq {\inlinecodemath{n}} < {\inlinecodemath{MASK\_BITS}}$.
\end{compactitem}

\Needspace*{3\baselineskip}
\textbf{Time Complexity}:
\begin{compactitem}
  \item $O(1)$ per call to \inlinecode{gray\_code()}.
  \item $O(\log b)$ per call to \inlinecode{inverse\_gray\_code()}, where $b$ is \inlinecode{MASK\_BITS}.
  \item $O(2^{n})$ per call to \inlinecode{gray\_sequence(n)}.
\end{compactitem}

\Needspace*{3\baselineskip}
\textbf{Space Complexity}:
\begin{compactitem}
  \item $O(1)$ auxiliary for \inlinecode{gray\_code()} and \inlinecode{inverse\_gray\_code()}.
  \item $O(1)$ auxiliary and $O(2^{n})$ for the returned sequence from \inlinecode{gray\_sequence(n)}.
\end{compactitem}

\begin{lstlisting}[language=C++]
#include <cassert>
#include <climits>
#include <cstdint>
#include <vector>

using Mask = uint32_t;
const int MASK_BITS = sizeof(Mask) * CHAR_BIT;

Mask gray_code(Mask k) {
  return k ^ (k >> 1);
}

Mask inverse_gray_code(Mask g) {
  for (int shift = 1; shift < MASK_BITS && (g >> shift) != 0; shift <<= 1) {
    g ^= g >> shift;
  }
  return g;
}

std::vector<Mask> gray_sequence(int n) {
  assert(0 <= n && n < MASK_BITS);
  std::vector<Mask> res(Mask{1} << n);
  for (Mask k = 0; k < static_cast<Mask>(res.size()); k++) {
    res[k] = gray_code(k);
  }
  return res;
}
\end{lstlisting}
\Needspace*{4\baselineskip}
\textbf{Example Usage}\par
\begin{lstlisting}[language=C++]
#include <cassert>
using namespace std;

int main() {
  assert(gray_code(0) == 0b000U);
  assert(gray_code(1) == 0b001U);
  assert(gray_code(2) == 0b011U);
  assert(gray_code(3) == 0b010U);

  for (Mask k = 0; k < 256; k++) {
    assert(inverse_gray_code(gray_code(k)) == k);
  }
  assert(inverse_gray_code(gray_code(0xdeadbeefU)) == 0xdeadbeefU);

  vector<Mask> seq = gray_sequence(3);
  assert((seq == vector<Mask>{0b000, 0b001, 0b011, 0b010, 0b110, 0b111, 0b101, 0b100}));
  for (int i = 0; i < static_cast<int>(seq.size()); i++) {
    Mask diff = seq[i] ^ seq[(i + 1) % seq.size()];
    assert((diff & (diff - 1)) == 0);  // Cyclically, one bit changes per step.
  }
  return 0;
}
\end{lstlisting}
\subsection{XOR Basis}
An XOR basis (also called a linear basis over the field GF(2)) is the analog of a vector space basis for integers under the bitwise XOR operation. Each integer is treated as a vector of bits, where XOR plays the role of vector addition. Inserting a set of integers one at a time and keeping only those that are linearly independent yields a basis whose XOR-combinations (subset XORs) span exactly the same set of values as the original integers. From a basis of size $r$, exactly $2^r$ distinct values are reachable, and queries such as the maximum attainable XOR become a simple greedy scan.

The basis is kept in row-echelon form: the $i$-th basis element is either $0$ or a value whose highest set bit is bit $i$, and no two basis elements share the same highest bit. Reducing a value means XORing it against the basis elements from the highest bit down, eliminating each set bit that a basis element covers.

\begin{compactitem}
  \item \inlinecode{XorBasis()} constructs an empty basis over \inlinecode{Mask}, which is \inlinecode{uint64\_t} by default.
  \item \inlinecode{insert(x)} adds \inlinecode{x} to the basis, returning \inlinecode{true} if \inlinecode{x} was linearly independent of the current basis (increasing its size), or \inlinecode{false} if \inlinecode{x} was already representable as a subset XOR.
  \item \inlinecode{contains(x)} returns whether \inlinecode{x} is representable as the XOR of some subset of the inserted values, i.e. whether it reduces to $0$ against the basis.
  \item \inlinecode{max\_xor(base = 0)} returns the maximum value of \inlinecode{base} XORed with some subset XOR of the basis. With the default \inlinecode{base} of $0$, this is the maximum subset XOR.
  \item \inlinecode{size()} returns the number of elements in the basis (its rank).
\end{compactitem}

\Needspace*{3\baselineskip}
\textbf{Time Complexity}:
\begin{compactitem}
  \item $O(b)$ per call to \inlinecode{insert()}, \inlinecode{contains()}, and \inlinecode{max\_xor()}, where $b$ is the number of bits.
  \item $O(1)$ per call to \inlinecode{size()}.
\end{compactitem}

\Needspace*{3\baselineskip}
\textbf{Space Complexity}:
\begin{compactitem}
  \item $O(b)$ for storage of the basis.
  \item $O(1)$ auxiliary for all operations.
\end{compactitem}

\begin{lstlisting}[language=C++]
#include <array>
#include <climits>
#include <cstdint>

using Mask = uint64_t;
const int MASK_BITS = sizeof(Mask) * CHAR_BIT;

class XorBasis {
  std::array<Mask, MASK_BITS> basis{};  // basis[i] != 0 has its highest set bit at bit i.
  int basis_size = 0;

 public:
  bool insert(Mask x) {
    for (int i = MASK_BITS - 1; i >= 0; i--) {
      if (!((x >> i) & 1)) {
        continue;
      }
      if (basis[i] == 0) {
        basis[i] = x;
        basis_size++;
        return true;
      }
      x ^= basis[i];
    }
    return false;
  }

  bool contains(Mask x) const {
    for (int i = MASK_BITS - 1; i >= 0; i--) {
      if (((x >> i) & 1) && basis[i] != 0) {
        x ^= basis[i];
      }
    }
    return x == 0;
  }

  Mask max_xor(Mask base = 0) const {
    for (int i = MASK_BITS - 1; i >= 0; i--) {
      if ((base ^ basis[i]) > base) {
        base ^= basis[i];
      }
    }
    return base;
  }

  int size() const { return basis_size; }
};
\end{lstlisting}
\Needspace*{4\baselineskip}
\textbf{Example Usage}\par
\begin{lstlisting}[language=C++]
#include <cassert>

int main() {
  XorBasis b;
  assert(b.insert(1));   // 001
  assert(b.insert(2));   // 010
  assert(!b.insert(3));  // 011 = 1 ^ 2, already representable.
  assert(b.size() == 2);

  assert(b.contains(3));
  assert(b.contains(0));
  assert(!b.contains(4));

  assert(b.max_xor() == 3);   // Best subset XOR over {0, 1, 2, 3}.
  assert(b.max_xor(4) == 7);  // 4 ^ 3.
  return 0;
}
\end{lstlisting}
\subsection{Binary Trie}
Maintain a multiset of nonnegative $b$-bit integers as a binary trie, storing each value as a root-to-leaf path of its bits from the most significant (bit $b - 1$) down to the least significant. Every node keeps a count of the stored values passing through it, which supports deletion and counting queries, and lets the trie answer XOR-extremal queries by a greedy walk down the bits.

The classic application is the maximum-XOR query: to maximize \inlinecode{x \textasciicircum{} y} over all stored \inlinecode{y}, walk from the most significant bit. At each step, descend toward the child whose bit differs from that of \inlinecode{x} whenever such a branch exists, since setting a higher bit of the result always dominates any combination of lower bits. The minimum-XOR query is the mirror image, preferring the matching bit.

This is a multiset (a value may be inserted more than once), and is distinct from the XOR basis: the basis spans subset XORs of a fixed set, whereas this answers XOR queries against the stored elements themselves with support for insertion and deletion.

Values must lie in $[0, 2^b)$ for bit width $b$; the template parameters select the unsigned word type \inlinecode{U} and the bit width \inlinecode{BITS}, which must be less than the width of \inlinecode{U} (for example, $30$ with \inlinecode{uint32\_t} or $62$ with \inlinecode{uint64\_t}). Nodes are kept in a pool indexed by integer, where index $0$ is the root and also serves as the null child.

\begin{compactitem}
  \item \inlinecode{BinaryTrie\textless{}U, BITS\textgreater{}()} constructs an empty multiset.
  \item \inlinecode{size()} returns the number of stored values, and \inlinecode{empty()} returns whether there are none.
  \item \inlinecode{insert(x)} adds one occurrence of \inlinecode{x}.
  \item \inlinecode{count(x)} returns the number of occurrences of \inlinecode{x}.
  \item \inlinecode{erase(x)} removes one occurrence of \inlinecode{x}, returning whether one was present.
  \item \inlinecode{max\_xor(x)} and \inlinecode{min\_xor(x)} return the maximum and minimum of \inlinecode{x \textasciicircum{} y} over all stored \inlinecode{y}. The trie must be non-empty.
  \item \inlinecode{count\_xor\_less(x, k)} returns the number of stored \inlinecode{y} (with multiplicity) such that \inlinecode{x \textasciicircum{} y} is strictly less than \inlinecode{k}.
\end{compactitem}

\Needspace*{3\baselineskip}
\textbf{Time Complexity}:
\begin{compactitem}
  \item $O(b)$ per call to \inlinecode{insert()}, \inlinecode{erase()}, \inlinecode{count()}, \inlinecode{max\_xor()}, \inlinecode{min\_xor()}, and \inlinecode{count\_xor\_less()}, where $b$ is the bit width \inlinecode{BITS}.
  \item $O(1)$ per call to \inlinecode{size()} and \inlinecode{empty()}.
\end{compactitem}

\Needspace*{3\baselineskip}
\textbf{Space Complexity}:
\begin{compactitem}
  \item $O(n \cdot b)$ for storage, where $n$ is the number of distinct values inserted so far.
  \item $O(1)$ auxiliary per operation.
\end{compactitem}

\begin{lstlisting}[language=C++]
#include <array>
#include <cassert>
#include <climits>
#include <cstdint>
#include <vector>

template<typename U = uint32_t, int BITS = 30>
class BinaryTrie {
  static_assert(0 < BITS && BITS < static_cast<int>(sizeof(U) * CHAR_BIT));

  std::vector<std::array<int, 2>> child;  // child[node][bit], where 0 means no child.
  std::vector<int> cnt;                   // Number of stored values passing through each node.

  static void assert_fits(U x) { assert((x >> BITS) == 0); }

  int new_node() {
    child.push_back({0, 0});
    cnt.push_back(0);
    return static_cast<int>(cnt.size()) - 1;
  }

 public:
  BinaryTrie() { new_node(); }  // Node 0 is the root (and doubles as the null child).

  int size() const { return cnt[0]; }
  bool empty() const { return cnt[0] == 0; }

  void insert(U x) {
    assert_fits(x);
    int node = 0;
    cnt[node]++;
    for (int b = BITS - 1; b >= 0; b--) {
      int bit = (x >> b) & 1;
      if (child[node][bit] == 0) {
        int next = new_node();
        child[node][bit] = next;
      }
      node = child[node][bit];
      cnt[node]++;
    }
  }

  int count(U x) const {
    assert_fits(x);
    int node = 0;
    for (int b = BITS - 1; b >= 0; b--) {
      int next = child[node][(x >> b) & 1];
      if (next == 0) {
        return 0;
      }
      node = next;
    }
    return cnt[node];
  }

  bool erase(U x) {
    if (count(x) == 0) {
      return false;
    }
    int node = 0;
    cnt[node]--;
    for (int b = BITS - 1; b >= 0; b--) {
      node = child[node][(x >> b) & 1];
      cnt[node]--;
    }
    return true;
  }

  U max_xor(U x) const {
    assert(!empty());
    assert_fits(x);
    int node = 0;
    U res = 0;
    for (int b = BITS - 1; b >= 0; b--) {
      int bit = (x >> b) & 1;
      int opp = child[node][bit ^ 1];
      if (opp != 0 && cnt[opp] > 0) {
        res |= U{1} << b;
        node = opp;
      } else {
        node = child[node][bit];
      }
    }
    return res;
  }

  U min_xor(U x) const {
    assert(!empty());
    assert_fits(x);
    int node = 0;
    U res = 0;
    for (int b = BITS - 1; b >= 0; b--) {
      int bit = (x >> b) & 1;
      int same = child[node][bit];
      if (same != 0 && cnt[same] > 0) {
        node = same;
      } else {
        res |= U{1} << b;
        node = child[node][bit ^ 1];
      }
    }
    return res;
  }

  int count_xor_less(U x, U k) const {
    assert_fits(x);
    if ((k >> BITS) != 0) {
      return cnt[0];  // Every XOR result is below 2^BITS, which is at most k.
    }
    int node = 0, res = 0;
    for (int b = BITS - 1; b >= 0; b--) {
      int xb = (x >> b) & 1, kb = (k >> b) & 1;
      if (kb == 1) {
        // Stored values matching x at bit b give XOR bit 0 here, hence are already below k.
        int below = child[node][xb];
        if (below != 0) {
          res += cnt[below];
        }
        node = child[node][xb ^ 1];  // Stay tied with k by taking XOR bit 1.
      } else {
        node = child[node][xb];  // XOR bit must be 0 to remain below k.
      }
      if (node == 0) {
        break;
      }
    }
    return res;
  }
};
\end{lstlisting}
\Needspace*{4\baselineskip}
\textbf{Example Usage}\par
\begin{lstlisting}[language=C++]
#include <cassert>
using namespace std;

int main() {
  BinaryTrie<> trie;
  assert(trie.empty());
  for (int x : {3, 10, 5, 25, 2}) {
    trie.insert(x);
  }
  assert(!trie.empty());
  assert(trie.size() == 5);

  // x XOR y over the stored y: {26, 19, 28, 0, 27} for x = 25.
  assert(trie.max_xor(25) == 28);                  // 25 XOR 5.
  assert(trie.min_xor(25) == 0);                   // 25 XOR 25.
  assert(trie.count_xor_less(25, 20) == 2);        // XOR values 0 and 19 are below 20.
  assert(trie.count_xor_less(25, 1U << 30) == 5);  // Bound at/above 2^BITS counts all.

  assert(trie.count(10) == 1);
  trie.insert(10);
  assert(trie.count(10) == 2);  // Multiset: two copies of 10.
  assert(trie.erase(10) && trie.count(10) == 1);
  assert(trie.erase(25) && trie.count(25) == 0);
  assert(!trie.erase(25));        // Already gone.
  assert(trie.max_xor(0) == 10);  // Largest remaining value is 10.
  return 0;
}
\end{lstlisting}
\subsection{Bitwise Convolution (FWHT)}
Bitwise convolution combines two arrays indexed by masks, grouping pairs of masks by a bitwise operation instead of by ordinary addition. For arrays \inlinecode{a} and \inlinecode{b} over an $n$-bit universe, the XOR convolution \inlinecode{c[m]} is the sum of \inlinecode{a[x]*b[y]} over all \inlinecode{x \textasciicircum{} y = m}. OR and AND convolution replace \inlinecode{x \textasciicircum{} y} with \inlinecode{x | y} or \inlinecode{x \& y}. These appear in subset DP, counting pairs by bitwise result, and algebra over the Boolean hypercube.

The fast Walsh-Hadamard transform (FWHT) diagonalizes XOR convolution in the same way that the FFT diagonalizes polynomial convolution: transform both arrays, multiply pointwise, and transform back. The OR and AND versions are Zeta/Mobius transforms over the subset lattice, but they have the same butterfly-shaped implementation and the same $O(n \log n)$ runtime for arrays of length $n = 2^k$.

\begin{compactitem}
  \item \inlinecode{xor\_transform(f, invert)} applies the XOR FWHT to \inlinecode{f} in place. The inverse transform uses \inlinecode{invert = true} and divides every coefficient by \inlinecode{f.size()}.
  \item \inlinecode{or\_transform(f, invert)} applies the OR Zeta transform in place; the inverse is Mobius inversion.
  \item \inlinecode{and\_transform(f, invert)} applies the AND Zeta transform in place; the inverse is Mobius inversion.
  \item \inlinecode{xor\_convolve(a, b)} returns the XOR convolution of \inlinecode{a} and \inlinecode{b}.
  \item \inlinecode{or\_convolve(a, b)} returns the OR convolution of \inlinecode{a} and \inlinecode{b}.
  \item \inlinecode{and\_convolve(a, b)} returns the AND convolution of \inlinecode{a} and \inlinecode{b}.
\end{compactitem}

Subset convolution restricts OR convolution to pairs of disjoint masks, so \inlinecode{c[m]} must count only \inlinecode{x | y = m} with \inlinecode{x \& y = 0}. Rank each entry by its number of set bits, transform every rank layer, and multiply the layers as polynomials in the rank. Mobius inversion then recovers pairs whose union is exactly \inlinecode{m}; among these pairs, the ranks sum to the number of set bits in \inlinecode{m} exactly when the pair is disjoint.

\begin{compactitem}
  \item \inlinecode{subset\_convolve(a, b)} returns the subset convolution of \inlinecode{a} and \inlinecode{b}, where \inlinecode{c[m]} sums \inlinecode{a[x]*b[y]} over the pairs of disjoint masks satisfying \inlinecode{x | y = m}.
\end{compactitem}

The in-place transforms require a nonempty array whose length is a power of two. A convolution returns empty if either input is empty; otherwise, both inputs are padded with zeros to the smallest power of two at least their maximum length, which is also the output length. For exact integer XOR convolution, the inverse divisions are exact; for modular arithmetic, replace the division by multiplication with the modular inverse of the transform length.

Overflow warning: All intermediate sums, differences, and products must be representable in the value type.

\Needspace*{3\baselineskip}
\textbf{Time Complexity}:
\begin{compactitem}
  \item $O(n \log n)$ per transform or bitwise convolution, where $n$ is the padded power-of-two length.
  \item $O(n \log^{2} n)$ per call to \inlinecode{subset\_convolve()}.
\end{compactitem}

\Needspace*{3\baselineskip}
\textbf{Space Complexity}:
\begin{compactitem}
  \item $O(1)$ auxiliary for the in-place transforms.
  \item $O(n)$ auxiliary and $O(n)$ for the returned array from each bitwise convolution.
  \item $O(n \log n)$ auxiliary for \inlinecode{subset\_convolve()}, from the rank layers.
\end{compactitem}

\begin{lstlisting}[language=C++]
#include <algorithm>
#include <cassert>
#include <cstdint>
#include <utility>
#include <vector>

template<typename T>
void xor_transform(std::vector<T> &f, bool invert) {
  int n = static_cast<int>(f.size());
  assert(n > 0 && (n & (n - 1)) == 0);
  for (int len = 1; len < n; len <<= 1) {
    for (int i = 0; i < n; i += len << 1) {
      for (int j = 0; j < len; j++) {
        T u = f[i + j], v = f[i + j + len];
        f[i + j] = u + v;  // Overflow warning.
        f[i + j + len] = u - v;
      }
    }
  }
  if (invert) {
    for (T &x : f) {
      x /= n;
    }
  }
}

template<typename T>
void or_transform(std::vector<T> &f, bool invert) {
  int n = static_cast<int>(f.size());
  assert(n > 0 && (n & (n - 1)) == 0);
  for (int bit = 1; bit < n; bit <<= 1) {
    for (int mask = 0; mask < n; mask++) {
      if (mask & bit) {
        if (invert) {
          f[mask] -= f[mask ^ bit];  // Overflow warning.
        } else {
          f[mask] += f[mask ^ bit];
        }
      }
    }
  }
}

template<typename T>
void and_transform(std::vector<T> &f, bool invert) {
  int n = static_cast<int>(f.size());
  assert(n > 0 && (n & (n - 1)) == 0);
  for (int bit = 1; bit < n; bit <<= 1) {
    for (int mask = 0; mask < n; mask++) {
      if ((mask & bit) == 0) {
        if (invert) {
          f[mask] -= f[mask ^ bit];  // Overflow warning.
        } else {
          f[mask] += f[mask ^ bit];
        }
      }
    }
  }
}

template<typename T, typename Transform>
std::vector<T> bitwise_convolve(std::vector<T> a, std::vector<T> b, Transform transform) {
  if (a.empty() || b.empty()) {
    return {};
  }
  int needed = static_cast<int>(std::max(a.size(), b.size()));
  int n = 1;
  while (n < needed) {
    n <<= 1;
  }
  a.resize(n);
  b.resize(n);
  transform(a, false);
  transform(b, false);
  for (int i = 0; i < n; i++) {
    a[i] *= b[i];  // Overflow warning.
  }
  transform(a, true);
  return a;
}

template<typename T>
std::vector<T> xor_convolve(std::vector<T> a, std::vector<T> b) {
  return bitwise_convolve(std::move(a), std::move(b), xor_transform<T>);
}

template<typename T>
std::vector<T> or_convolve(std::vector<T> a, std::vector<T> b) {
  return bitwise_convolve(std::move(a), std::move(b), or_transform<T>);
}

template<typename T>
std::vector<T> and_convolve(std::vector<T> a, std::vector<T> b) {
  return bitwise_convolve(std::move(a), std::move(b), and_transform<T>);
}

template<typename T>
std::vector<T> subset_convolve(std::vector<T> a, std::vector<T> b) {
  if (a.empty() || b.empty()) {
    return {};
  }
  int needed = static_cast<int>(std::max(a.size(), b.size()));
  int n = 1, bits = 0;
  while (n < needed) {
    n <<= 1;
    bits++;
  }
  a.resize(n);
  b.resize(n);
  std::vector<int> ranks(n);
  for (int mask = 1; mask < n; mask++) {
    ranks[mask] = ranks[mask >> 1] + (mask & 1);
  }
  // Split each input into layers by rank, so that layer k holds only the masks with k set bits.
  std::vector<std::vector<T>> fa(bits + 1, std::vector<T>(n)), fb = fa, fc = fa;
  for (int mask = 0; mask < n; mask++) {
    fa[ranks[mask]][mask] = a[mask];
    fb[ranks[mask]][mask] = b[mask];
  }
  for (int k = 0; k <= bits; k++) {
    or_transform(fa[k], false);
    or_transform(fb[k], false);
  }
  // Multiplying layers groups pairs by total rank. Once Mobius inversion recovers each pair's exact
  // union, total rank ranks[mask] selects precisely the disjoint pairs covering mask.
  for (int i = 0; i <= bits; i++) {
    for (int j = 0; i + j <= bits; j++) {
      for (int mask = 0; mask < n; mask++) {
        fc[i + j][mask] += fa[i][mask] * fb[j][mask];  // Overflow warning.
      }
    }
  }
  std::vector<T> res(n);
  for (int k = 0; k <= bits; k++) {
    or_transform(fc[k], true);
  }
  for (int mask = 0; mask < n; mask++) {
    res[mask] = fc[ranks[mask]][mask];
  }
  return res;
}
\end{lstlisting}
\Needspace*{4\baselineskip}
\textbf{Example Usage}\par
\begin{lstlisting}[language=C++]
#include <cassert>
using namespace std;

int main() {
  {
    vector<int64_t> f{1, 2, 3, 4}, original(f);
    xor_transform(f, false);
    xor_transform(f, true);
    assert(f == original);
  }
  {
    vector<int64_t> f{1, 2, 3, 4}, original(f);
    or_transform(f, false);
    or_transform(f, true);
    assert(f == original);
  }
  {
    vector<int64_t> f{1, 2, 3, 4}, original(f);
    and_transform(f, false);
    and_transform(f, true);
    assert(f == original);
  }

  vector<int64_t> a{1, 2, 3, 4}, b{5, 6, 7, 8};

  // XOR: c[m] sums a[x] * b[y] over pairs with x ^ y == m.
  assert((xor_convolve(a, b) == vector<int64_t>{70, 68, 62, 60}));

  // OR: c[m] sums a[x] * b[y] over pairs with x | y == m.
  assert((or_convolve(a, b) == vector<int64_t>{5, 28, 43, 184}));

  // AND: c[m] sums a[x] * b[y] over pairs with x & y == m.
  assert((and_convolve(a, b) == vector<int64_t>{103, 52, 73, 32}));

  // Subset: c[m] sums a[x] * b[y] over disjoint pairs with x | y == m. Every entry is at most the
  // corresponding OR convolution, which also counts the overlapping pairs.
  assert((subset_convolve(a, b) == vector<int64_t>{5, 16, 22, 60}));

  vector<int64_t> one_mask{0, 1, 0, 0};
  vector<int64_t> two_mask{0, 0, 1};
  assert((xor_convolve(one_mask, two_mask) == vector<int64_t>{0, 0, 0, 1}));
  assert((or_convolve(one_mask, two_mask) == vector<int64_t>{0, 0, 0, 1}));
  assert((and_convolve(one_mask, two_mask) == vector<int64_t>{1, 0, 0, 0}));
  assert((subset_convolve(one_mask, two_mask) == vector<int64_t>{0, 0, 0, 1}));

  // Overlapping masks contribute nothing at all to a subset convolution.
  assert((subset_convolve(one_mask, one_mask) == vector<int64_t>{0, 0, 0, 0}));
  return 0;
}
\end{lstlisting}

\section{\texorpdfstring{Classic Problems}{1.8 Classic Problems}}
\setcounter{section}{8}
\setcounter{subsection}{0}
\subsection{Minimum Excluded Value}
Given a multiset of integers, find the minimum excluded nonnegative value (MEX), i.e. the smallest integer $x \geq 0$ not present in the multiset. MEX often appears in array problems, game theory with Grundy numbers, and dynamic programming states. Since the MEX of $n$ values never exceeds $n$, the one-shot version tracks seen values in $[0, n]$ before scanning to find the first missing.

\begin{compactitem}
  \item \inlinecode{mex(lo, hi)} returns the MEX of the values in $[{\inlinecodemath{lo}}, {\inlinecodemath{hi}})$.
\end{compactitem}

The MEX of a dynamic multiset can be computed using counts of present values and maintaining an ordered set of currently missing candidates.

\begin{compactitem}
  \item \inlinecode{DynamicMex()} maintains the MEX of a multiset of nonnegative integers.
  \item \inlinecode{add(x)} inserts one copy of value \inlinecode{x} if \inlinecode{x} $\geq 0$.
  \item \inlinecode{remove(x)} removes one copy of value \inlinecode{x} if present.
  \item \inlinecode{mex()} returns the current MEX.
\end{compactitem}

\Needspace*{3\baselineskip}
\textbf{Time Complexity}:
\begin{compactitem}
  \item $O(n)$ per call to \inlinecode{mex(lo, hi)}, where $n$ is the number of values.
  \item $O(\log n)$ expected per call to \inlinecode{add(x)} and \inlinecode{remove(x)} for \inlinecode{DynamicMex}, and $O(n)$ in the collision-heavy worst case for the hash table.
  \item $O(1)$ per call to \inlinecode{mex()} for \inlinecode{DynamicMex}.
\end{compactitem}

\Needspace*{3\baselineskip}
\textbf{Space Complexity}:
\begin{compactitem}
  \item $O(n)$ auxiliary for \inlinecode{mex(lo, hi)}.
  \item $O(n)$ storage for \inlinecode{DynamicMex}.
\end{compactitem}

\begin{lstlisting}[language=C++]
#include <iterator>
#include <set>
#include <unordered_map>
#include <vector>

template<typename It>
int mex(It lo, It hi) {
  int n = std::distance(lo, hi);
  std::vector<char> seen(n + 1);
  for (It it = lo; it != hi; ++it) {
    if (0 <= *it && *it <= n) {
      seen[*it] = true;
    }
  }
  for (int x = 0; x <= n; x++) {
    if (!seen[x]) {
      return x;
    }
  }
  return n + 1;
}

class DynamicMex {
  std::unordered_map<int, int> count;
  std::set<int> missing;  // Missing candidates in [0, size].
  int size = 0;

 public:
  DynamicMex() { missing.insert(0); }

  void add(int x) {
    if (x < 0) {
      return;
    }
    if (!count.count(size + 1)) {
      missing.insert(size + 1);
    }
    if (count[x]++ == 0) {
      missing.erase(x);
    }
    size++;
  }

  void remove(int x) {
    auto it = count.find(x);
    if (it == count.end()) {
      return;
    }
    bool removed_last = --it->second == 0;
    if (removed_last) {
      count.erase(it);
    }
    size--;
    missing.erase(size + 1);
    if (removed_last && x <= size) {
      missing.insert(x);
    }
  }

  int mex() const { return *missing.begin(); }
};
\end{lstlisting}
\Needspace*{4\baselineskip}
\textbf{Example Usage}\par
\begin{lstlisting}[language=C++]
#include <cassert>
using namespace std;

int main() {
  vector<int> a{0, 1, 4, 2, 1};
  assert(mex(a.begin(), a.end()) == 3);
  vector<int> b{-1, 0, 1, 2};
  assert(mex(b.begin(), b.end()) == 3);  // Negative values are ignored.

  DynamicMex m;
  m.add(0);
  m.add(1);
  m.add(1);
  m.add(3);
  assert(m.mex() == 2);
  m.add(2);
  assert(m.mex() == 4);
  m.remove(1);
  assert(m.mex() == 4);  // One copy of 1 remains.
  m.remove(1);
  assert(m.mex() == 1);
  m.remove(100);
  assert(m.mex() == 1);
  return 0;
}
\end{lstlisting}
\subsection{Josephus Problem}
The Josephus problem arranges $n$ people, numbered $[0, n)$, in a circle and repeatedly removes every $k$-th remaining person until one survives. Counting starts at person $0$, which counts as the first person. After the first removal, relabeling the smaller circle from the next person gives the same problem on $n - 1$ people. Mapping its survivor back to the original labels yields the recurrence $J(1, k) = 0$ and $J(n, k) = (J(n - 1, k) + k) \bmod n$.

When $k$ is small, several removals can be processed at once. One complete pass removes $\lfloor n/k \rfloor$ people; recursively solving the compressed circle and undoing the resulting index shift reduces the running time to $O(k \log n)$. To recover the full elimination order, a Fenwick tree stores which labels remain and selects each next victim by rank.

\begin{compactitem}
  \item \inlinecode{josephus(n, k)} returns the 0-based label of the survivor, where \inlinecode{n} and \inlinecode{k} must be positive.
  \item \inlinecode{josephus\_small\_k(n, k)} returns the same survivor using the batched recurrence. It is preferable when $k$ is small relative to $n$.
  \item \inlinecode{josephus\_order(n, k)} returns all labels in elimination order, ending with the survivor.
\end{compactitem}

\Needspace*{3\baselineskip}
\textbf{Time Complexity}:
\begin{compactitem}
  \item $O(n)$ per call to \inlinecode{josephus()}.
  \item $O(\min\left(n, k \log n\right))$ per call to \inlinecode{josephus\_small\_k()}.
  \item $O(n \log n)$ per call to \inlinecode{josephus\_order()}.
\end{compactitem}

\Needspace*{3\baselineskip}
\textbf{Space Complexity}:
\begin{compactitem}
  \item $O(1)$ auxiliary for \inlinecode{josephus()}.
  \item $O(\min\left(n, k \log n\right))$ call stack space for \inlinecode{josephus\_small\_k()}.
  \item $O(n)$ auxiliary and $O(n)$ for the returned elimination order from \inlinecode{josephus\_order()}.
\end{compactitem}

\begin{lstlisting}[language=C++]
#include <cassert>
#include <cstdint>
#include <vector>

int josephus(int n, int k) {
  assert(n > 0 && k > 0);
  int survivor = 0;
  for (int size = 2; size <= n; size++) {
    survivor = static_cast<int>((static_cast<int64_t>(survivor) + k) % size);
  }
  return survivor;
}

int josephus_small_k(int n, int k) {
  assert(n > 0 && k > 0);
  if (n == 1) {
    return 0;
  }
  if (k == 1) {
    return n - 1;
  }
  if (k > n) {
    return josephus(n, k);
  }
  int removed = n / k;
  int survivor = josephus_small_k(n - removed, k);
  survivor -= n % k;
  if (survivor < 0) {
    survivor += n;
  } else {
    survivor += survivor / (k - 1);
  }
  return survivor;
}

std::vector<int> josephus_order(int n, int k) {
  assert(n > 0 && k > 0);
  std::vector<int> bit(n + 1), order;
  order.reserve(n);
  for (int i = 1; i <= n; i++) {
    bit[i] = i & -i;  // Initially every label is alive.
  }
  int max_step = 1;
  while (max_step <= n / 2) {
    max_step *= 2;
  }
  auto find_by_rank = [&](int rank) {
    int pos = 0;
    for (int step = max_step; step > 0; step /= 2) {
      int next = pos + step;
      if (next <= n && bit[next] <= rank) {
        pos = next;
        rank -= bit[next];
      }
    }
    return pos;
  };
  int rank = 0;
  for (int remaining = n; remaining > 0; remaining--) {
    rank = static_cast<int>((static_cast<int64_t>(rank) + k - 1) % remaining);
    int removed = find_by_rank(rank);
    order.push_back(removed);
    for (int i = removed + 1; i <= n; i += i & -i) {
      bit[i]--;
    }
    if (remaining > 1) {
      rank %= remaining - 1;
    }
  }
  return order;
}
\end{lstlisting}
\Needspace*{4\baselineskip}
\textbf{Example Usage}\par
\begin{lstlisting}[language=C++]
#include <cassert>
#include <vector>
using namespace std;

int main() {
  assert(josephus(1, 5) == 0);
  assert(josephus(7, 3) == 3);
  assert(josephus(10, 1) == 9);
  for (int n = 1; n <= 30; n++) {
    for (int k = 1; k <= 30; k++) {
      assert(josephus_small_k(n, k) == josephus(n, k));
      assert(josephus_order(n, k).back() == josephus(n, k));
    }
  }
  assert(josephus_order(7, 3) == vector<int>({2, 5, 1, 6, 4, 0, 3}));
  return 0;
}
\end{lstlisting}
\subsection{Tower of Hanoi}
Move a stack of $n$ disks of distinct sizes from one peg to another using a third as scratch space, moving one disk at a time and never placing a larger disk on a smaller one. The recursive solution is forced by the largest disk: it can only move once the other $n - 1$ disks are stacked on the spare peg, so the puzzle is solved by moving $n - 1$ disks aside, moving the largest disk across, and moving the $n - 1$ disks back on top of it. That gives $T(n) = 2T(n - 1) + 1$ with $T(0) = 0$, so exactly $2^n - 1$ moves are needed, and no shorter solution exists because the largest disk cannot move any earlier.

The sequence also has a closed form, for when one move is wanted rather than all of them: numbering moves from $1$, move $k$ transfers the disk given by the number of trailing zero bits of $k$, and that disk always steps in one rotational direction fixed by the parity of $n$. The smallest disk therefore moves every other turn, which is the rule an iterative solver follows.

\begin{compactitem}
  \item \inlinecode{hanoi\_moves(n, from = 0, to = 2)} returns the moves solving the puzzle for \inlinecode{n} disks, as pairs (\inlinecode{from\_peg}, \inlinecode{to\_peg}). Pegs are numbered $[0, 3)$ and the spare peg is whichever of the three is neither \inlinecode{from} nor \inlinecode{to}. The number of disks must be nonnegative.
  \item \inlinecode{hanoi\_move\_count(n)} returns $2^n - 1$, the number of moves, without generating them. It requires $n \leq 63$, the largest count that fits in \inlinecode{int64\_t}.
\end{compactitem}

With four or more pegs the problem becomes Frame-Stewart: move a prefix aside using every peg, transfer the rest with one peg fewer, bring the prefix back, and minimize over the split. That this family of strategies is optimal was only proved in 2014.

\Needspace*{3\baselineskip}
\textbf{Time Complexity}:
\begin{compactitem}
  \item $O(2^{n})$ per call to \inlinecode{hanoi\_moves()}, which is optimal since that is the size of the output.
  \item $O(1)$ per call to \inlinecode{hanoi\_move\_count()}.
\end{compactitem}

\Needspace*{3\baselineskip}
\textbf{Space Complexity}:
\begin{compactitem}
  \item $O(n)$ auxiliary stack space and $O(2^{n})$ for the returned moves from \inlinecode{hanoi\_moves()}.
  \item $O(1)$ auxiliary for \inlinecode{hanoi\_move\_count()}.
\end{compactitem}

\begin{lstlisting}[language=C++]
#include <cassert>
#include <cstdint>
#include <utility>
#include <vector>

std::vector<std::pair<int, int>> hanoi_moves(int n, int from = 0, int to = 2) {
  assert(n >= 0 && from >= 0 && from < 3 && to >= 0 && to < 3 && from != to);
  std::vector<std::pair<int, int>> moves;
  auto rec = [&](auto &&rec, int disks, int src, int dest) -> void {
    if (disks == 0) {
      return;
    }
    int spare = 3 - src - dest;
    rec(rec, disks - 1, src, spare);
    moves.emplace_back(src, dest);
    rec(rec, disks - 1, spare, dest);
  };
  rec(rec, n, from, to);
  return moves;
}

int64_t hanoi_move_count(int n) {
  assert(n >= 0 && n <= 63);
  return static_cast<int64_t>((uint64_t{1} << n) - 1);  // Unsigned, so n = 63 does not overflow.
}
\end{lstlisting}
\Needspace*{4\baselineskip}
\textbf{Example Usage}\par
\begin{lstlisting}[language=C++]
#include <cassert>
using namespace std;

int main() {
  assert(hanoi_moves(0).empty());
  assert((hanoi_moves(1) == vector<pair<int, int>>{{0, 2}}));
  assert((hanoi_moves(2) == vector<pair<int, int>>{{0, 1}, {0, 2}, {1, 2}}));
  assert(hanoi_move_count(10) == 1023);

  // Replay the moves of a larger puzzle to confirm no disk ever lands on a smaller one.
  const int n = 12;
  vector<vector<int>> pegs(3);
  for (int disk = n; disk >= 1; disk--) {
    pegs[0].push_back(disk);
  }
  auto moves = hanoi_moves(n);
  assert(static_cast<int64_t>(moves.size()) == hanoi_move_count(n));
  for (auto [src, dest] : moves) {
    assert(!pegs[src].empty());
    assert(pegs[dest].empty() || pegs[dest].back() > pegs[src].back());
    pegs[dest].push_back(pegs[src].back());
    pegs[src].pop_back();
  }
  assert(pegs[0].empty() && pegs[1].empty() && static_cast<int>(pegs[2].size()) == n);
  return 0;
}
\end{lstlisting}
\subsection{Cycle Detection (Floyd)}
Given a function $f$ mapping a finite set to itself and a starting value $x_0$, return the entry value, position, and length of the cycle reached by repeatedly applying $f$. Formally, the sequence $x_0$, $x_1 = f(x_0)$, $x_2 = f(x_1)$, $\ldots$, $x_n = f(x_{n - 1})$, $\ldots$ must eventually repeat some value. If $x_i = x_j$ for $i < j$, then the sequence from $x_i$ to $x_{j - 1}$ repeats forever. This is useful for detecting cycles in functional graphs (see 4.2.9), pseudo-random generators, Pollard rho style algorithms, degenerate linked lists, and arrays where each index points to the next index. For example, if an array $a_0, \ldots, a_{n+1}$ contains only values in $[1, n]$, with exactly one value occurring more than once, then $f(i) = a_i$ forms a functional graph; the duplicate value is the entry point of the cycle reached from $0$.

Floyd's cycle-finding algorithm, a.k.a. the ``tortoise and the hare algorithm'', is a space-efficient default choice: it moves two pointers through the sequence at different speeds, finds a meeting point inside the cycle, then locates the cycle start. Brent's algorithm is a compact variant that usually calls $f$ fewer times by only resetting the tortoise at powers of two; prefer it when evaluating $f$ is expensive and constant factors matter. In the detection phase, Brent advances one pointer with one call to $f$ per step, while Floyd advances the tortoise once and the hare twice, using three calls to $f$ per step.

For Floyd's algorithm, suppose the tail has length $m$ and the cycle has length $n$. At a meeting, the hare has traveled twice as far as the tortoise, so the tortoise's distance is a multiple of $n$. Consequently, the distance from the meeting point around the cycle to its entry is congruent to $m$ modulo $n$. Resetting one pointer to $x_0$ and advancing both one step at a time makes them meet at the cycle entry after $m$ steps. Brent first determines $n$, then advances one pointer $n$ steps; moving both together from that fixed separation likewise makes them meet at the entry.

\begin{compactitem}
  \item \inlinecode{find\_cycle\_floyd(f, x0)} repeatedly applies \inlinecode{f} starting from \inlinecode{x0} and returns the reached cycle as a tuple (\inlinecode{entry}, \inlinecode{start}, \inlinecode{length}), where \inlinecode{entry} is the first cycle value, \inlinecode{start} is the number of applications of \inlinecode{f} needed to reach it from \inlinecode{x0}, and \inlinecode{length} is the number of distinct values in the cycle.
  \item \inlinecode{find\_cycle\_brent(f, x0)} does the same with fewer calls to \inlinecode{f} in many cases.
\end{compactitem}

\textbf{Time Complexity}: $O(m + n)$ per call, where $m$ is the first index in the cycle and $n$ is the cycle length.

\textbf{Space Complexity}: $O(1)$ auxiliary.

\begin{lstlisting}[language=C++]
#include <tuple>

template<typename Fn, typename T>
std::tuple<T, int, int> find_cycle_floyd(Fn f, T x0) {
  // Find a meeting point inside the cycle.
  T tortoise = f(x0), hare = f(f(x0));
  while (tortoise != hare) {
    tortoise = f(tortoise);
    hare = f(f(hare));
  }
  // Find the cycle entry.
  int start = 0;
  tortoise = x0;
  while (tortoise != hare) {
    tortoise = f(tortoise);
    hare = f(hare);
    start++;
  }
  // Measure the cycle length.
  int length = 1;
  hare = f(tortoise);
  while (tortoise != hare) {
    hare = f(hare);
    length++;
  }
  return {tortoise, start, length};
}

template<typename Fn, typename T>
std::tuple<T, int, int> find_cycle_brent(Fn f, T x0) {
  // Find the cycle length by doubling the search block.
  int power = 1, length = 1;
  T tortoise = x0, hare = f(x0);
  while (tortoise != hare) {
    if (power == length) {
      tortoise = hare;
      power *= 2;
      length = 0;
    }
    hare = f(hare);
    length++;
  }
  // Separate the pointers by one cycle length.
  hare = x0;
  for (int i = 0; i < length; i++) {
    hare = f(hare);
  }
  // Find the cycle entry.
  int start = 0;
  tortoise = x0;
  while (tortoise != hare) {
    tortoise = f(tortoise);
    hare = f(hare);
    start++;
  }
  return {tortoise, start, length};
}
\end{lstlisting}
\Needspace*{4\baselineskip}
\textbf{Example Usage}\par
\begin{lstlisting}[language=C++]
#include <cassert>
#include <set>
#include <vector>
using namespace std;

int f(int x) {
  return (123 * x * x + 4567890) % 1337;
}

void verify(int x0, int start, int length) {
  set<int> s;
  int x = x0;
  for (int i = 0; i < start; i++) {
    assert(!s.count(x));
    s.insert(x);
    x = f(x);
  }
  int startx = x;
  s.clear();
  for (int i = 0; i < length; i++) {
    assert(!s.count(x));
    s.insert(x);
    x = f(x);
  }
  assert(startx == x);
}

// Given n + 1 integers within [1, n], where one value in [1, n] is present 2 or more times (but
// other values in [1, n] are present at most once), find the duplicate without modifying the array.
// Interpret each value as a pointer to the next index; the duplicate is the cycle entry.
int find_duplicate_integer(const vector<int> &a) {
  return get<0>(find_cycle_floyd([&](int i) { return a[i]; }, 0));
}

int main() {
  int x0 = 0;
  auto [floyd_entry, floyd_start, floyd_len] = find_cycle_floyd(f, x0);
  assert(floyd_start == 4 && floyd_len == 2);
  verify(x0, floyd_start, floyd_len);

  auto [brent_entry, brent_start, brent_len] = find_cycle_brent(f, x0);
  assert(brent_entry == floyd_entry && brent_start == floyd_start && brent_len == floyd_len);

  assert(find_duplicate_integer({1, 3, 4, 2, 2}) == 2);
  assert(find_duplicate_integer({3, 1, 3, 4, 2}) == 3);
  return 0;
}
\end{lstlisting}
\subsection{Subset Sum}
Given a collection of integers and a target, the subset sum problem asks for the maximum sum of any subset that does not exceed the target. Two algorithms below solve this problem in complementary regimes, and which one wins depends on whether the bottleneck is the number of items or the size of the weights.

The meet-in-the-middle method splits the collection in half, enumerates every subset sum for each half, sorts one side, and binary searches for the best compatible partner. It makes no assumption about the sign or size of the values, so it handles negative inputs, and it is the method of choice when the number of items $n$ is small (up to roughly $40$).

Pisinger's algorithm targets the opposite regime: nonnegative weights whose maximum value is small, with possibly many items. Let $W$ be the largest weight. It greedily takes a maximal prefix with sum $s$ that fits the target; because the next weight is at most $W$ and does not fit, $s$ is greater than \inlinecode{target} $- W$. The prefix itself is feasible, so the optimum lies between $s$ and \inlinecode{target} and only that width-$W$ window matters. A balanced sweep explores exchanges that add later items to sums below the target and remove earlier items from sums above it. For each sum, it retains the farthest prefix boundary reached, which dominates smaller boundaries because it permits at least the same future exchanges. This gives a running time linear in the number of items times $W$, independent of the target beyond that window. This routine is sometimes called ``fast knapsack,'' a slight misnomer: it maximizes a subset sum rather than packing items with distinct values, so it is really the bounded-weight case of subset sum rather than the 0-1 value knapsack.

\begin{compactitem}
  \item \inlinecode{max\_subset\_sum\_at\_most(values, target)} returns a pair (\inlinecode{sum}, \inlinecode{items}) containing that maximum sum and the selected item indices. Values may be negative, and \inlinecode{target} must be nonnegative.
  \item \inlinecode{max\_subset\_sum\_bounded(weights, target)} returns the maximum sum of any subset of \inlinecode{weights} that does not exceed \inlinecode{target}. All weights and \inlinecode{target} must be nonnegative integers. This is faster than meet-in-the-middle when the largest weight is small relative to $2^{n/2}$.
\end{compactitem}

Overflow warning: All subset sums and intermediate differences in \inlinecode{max\_subset\_sum\_at\_most()} must fit in \inlinecode{int64\_t}.

\Needspace*{3\baselineskip}
\textbf{Time Complexity}:
\begin{compactitem}
  \item $O(n \cdot 2^{n/2})$ per call to \inlinecode{max\_subset\_sum\_at\_most()}, where $n$ is the number of values.
  \item $O(n \cdot W)$ per call to \inlinecode{max\_subset\_sum\_bounded()}, where $n$ is the number of items and $W$ is the largest weight.
\end{compactitem}

\Needspace*{3\baselineskip}
\textbf{Space Complexity}:
\begin{compactitem}
  \item $O(2^{n/2})$ auxiliary for \inlinecode{max\_subset\_sum\_at\_most()}.
  \item $O(W)$ auxiliary for \inlinecode{max\_subset\_sum\_bounded()}.
\end{compactitem}

\begin{lstlisting}[language=C++]
#include <algorithm>
#include <cassert>
#include <cstdint>
#include <utility>
#include <vector>

std::pair<int64_t, std::vector<int>> max_subset_sum_at_most(
    const std::vector<int> &values, int64_t target
) {
  assert(target >= 0);
  int n = static_cast<int>(values.size());
  assert(n / 2 < 63 && n - n / 2 < 63);
  int64_t llen = 1LL << (n / 2), rlen = 1LL << (n - n / 2);
  std::vector<int64_t> lsum(llen);
  std::vector<std::pair<int64_t, int64_t>> rsum(rlen);
  for (int64_t mask = 1; mask < llen; mask++) {
    int bit = __builtin_ctzll(mask);
    lsum[mask] = lsum[mask ^ (1LL << bit)] + values[bit];  // Overflow warning.
  }
  for (int64_t mask = 1; mask < rlen; mask++) {
    int bit = __builtin_ctzll(mask);
    rsum[mask].first = rsum[mask ^ (1LL << bit)].first + values[n / 2 + bit];
    rsum[mask].second = mask;
  }
  std::sort(rsum.begin(), rsum.end());
  int64_t best = INT64_MIN, lmask = 0, rmask = 0;
  for (int64_t mask = 0; mask < llen; mask++) {
    int64_t limit = target - lsum[mask];  // Overflow warning.
    auto it = std::upper_bound(rsum.begin(), rsum.end(), std::make_pair(limit, INT64_MAX));
    if (it != rsum.begin()) {
      --it;
      int64_t candidate = lsum[mask] + it->first;
      if (best < candidate) {
        best = candidate;
        lmask = mask;
        rmask = it->second;
      }
    }
  }
  // Optional: reconstruct one optimal subset.
  std::vector<int> items;
  for (int i = 0; i < n / 2; i++) {
    if ((lmask >> i) & 1) {
      items.push_back(i);
    }
  }
  for (int i = 0; i < n - n / 2; i++) {
    if ((rmask >> i) & 1) {
      items.push_back(n / 2 + i);
    }
  }
  return {best, items};
}

int max_subset_sum_bounded(const std::vector<int> &weights, int target) {
  assert(target >= 0);
  for (int w : weights) {
    assert(w >= 0);
  }
  int n = static_cast<int>(weights.size());
  // Greedily take a prefix while it still fits; the optimum then lies within one weight of target.
  int sum = 0, b = 0;
  while (b < n && weights[b] <= target - sum) {
    sum += weights[b++];
  }
  if (b == n) {
    return sum;  // Every item fits, so the whole collection is the best subset.
  }
  int m = 0;
  for (int w : weights) {
    m = std::max(m, w);
  }
  // reach[s - target + m] = largest prefix boundary that can realize a balanced subset of sum s.
  // Only sums within m of target are tracked, which is where the optimum must lie.
  std::vector<int> reach(2 * m, -1);
  reach[sum - target + m] = b;
  for (int i = b; i < n; i++) {
    std::vector<int> prev = reach;
    for (int x = 0; x < m; x++) {  // Add item i to sums currently below target.
      reach[x + weights[i]] = std::max(reach[x + weights[i]], prev[x]);
    }
    for (int x = 2 * m - 1; x > m; x--) {  // Drop an earlier item from sums above target.
      for (int j = std::max(0, prev[x]); j < reach[x]; j++) {
        reach[x - weights[j]] = std::max(reach[x - weights[j]], j);
      }
    }
  }
  int s = target;
  while (reach[s - target + m] < 0) {
    s--;
  }
  return s;
}
\end{lstlisting}
\Needspace*{4\baselineskip}
\textbf{Example Usage}\par
\begin{lstlisting}[language=C++]
#include <cassert>
using namespace std;

int main() {
  vector<int> a{9, 1, 5, 0, 1, 11, 5};
  auto [sum, items] = max_subset_sum_at_most(a, 8);
  int64_t selected_sum = 0;
  for (int i : items) {
    selected_sum += a[i];
  }
  assert(sum == 7 && selected_sum == sum);
  vector<int> b{-7, -3, -2, 5, 8};
  assert(max_subset_sum_at_most(b, 0).first == 0);

  // The bounded-weight method agrees with meet-in-the-middle on nonnegative inputs.
  assert(max_subset_sum_bounded(a, 8) == 7);
  vector<int> c{3, 34, 4, 12, 5, 2};
  assert(max_subset_sum_bounded(c, 9) == 9);     // 4 + 5.
  assert(max_subset_sum_bounded(c, 10) == 10);   // 3 + 5 + 2.
  assert(max_subset_sum_bounded(c, 100) == 60);  // All items fit.
  return 0;
}
\end{lstlisting}
\subsection{N-Queens}
The N-queens problem asks for a placement of $n$ queens on an $n$-by-$n$ chessboard such that no two share a row, column, or diagonal. Backtracking places one queen in each successive row and tracks which columns and diagonals are occupied. A conflicting placement is rejected immediately; after a failed recursive branch, its column and diagonals are released before trying the next column.

\begin{compactitem}
  \item \inlinecode{n\_queens(n)} returns a placement as a vector \inlinecode{cols} of length \inlinecode{n}, where \inlinecode{cols[row]} is the queen's column in that row. It returns an empty vector if no placement exists.
\end{compactitem}

\textbf{Time Complexity}: $O(n!)$ per call.

\textbf{Space Complexity}: $O(n)$ auxiliary and $O(n)$ for the returned placement.

\begin{lstlisting}[language=C++]
#include <cassert>
#include <vector>

std::vector<int> n_queens(int n) {
  assert(n > 0);
  std::vector<int> cols(n);
  std::vector<char> used_col(n), used_diag1(2 * n - 1), used_diag2(2 * n - 1);
  auto rec = [&](auto &&rec, int row) -> bool {
    if (row == n) {
      return true;
    }
    for (int col = 0; col < n; col++) {
      int diag1 = row - col + n - 1, diag2 = row + col;
      if (used_col[col] || used_diag1[diag1] || used_diag2[diag2]) {
        continue;
      }
      cols[row] = col;
      used_col[col] = used_diag1[diag1] = used_diag2[diag2] = true;
      if (rec(rec, row + 1)) {
        return true;
      }
      used_col[col] = used_diag1[diag1] = used_diag2[diag2] = false;
      cols[row] = -1;
    }
    return false;
  };
  return rec(rec, 0) ? cols : std::vector<int>{};
}
\end{lstlisting}
\Needspace*{4\baselineskip}
\textbf{Example Usage}\par
\begin{lstlisting}[language=C++]
#include <cassert>
#include <vector>
using namespace std;

int main() {
  assert(n_queens(4) == vector<int>({1, 3, 0, 2}));
  assert(n_queens(2).empty());
  assert(n_queens(1) == vector<int>({0}));
  return 0;
}
\end{lstlisting}
\subsection{Knights Tour}
A knight's tour visits every square of an $n$-by-$n$ board exactly once, moving as a chess knight. Plain backtracking over the eight moves is hopeless past a small board, since the search tree has branching factor up to eight and depth $n^2$.

Warnsdorff's rule fixes the move order: from each square, try the onward squares in increasing order of how many unvisited moves they themselves have. The intuition is that squares with few remaining exits, such as corners, become unreachable if they are not used early, so the rule visits the most constrained squares first and leaves the flexible ones for later. This is the same minimum-remaining-values idea that drives the exact cover search of section 1.8.8.

Warnsdorff's rule alone is a heuristic and does fail on some boards, so the search below keeps backtracking as a fallback and merely orders the moves by the rule. That ordering is strong enough that a tour is normally found along the first path tried, with backtracking almost never invoked.

\begin{compactitem}
  \item \inlinecode{knight\_tour(n, sr = 0, sc = 0)} returns an \inlinecode{n}-by-\inlinecode{n} grid whose entry is the step at which the knight visits that square, numbered $[0, {\inlinecodemath{n}}^2)$, starting from square (\inlinecode{sr}, \inlinecode{sc}). It returns an empty grid when no tour exists, which is the case for ${\inlinecodemath{n}} \in \{2, 3, 4\}$. The board size must be positive and the starting square must lie on the board.
\end{compactitem}

A tour is closed when the knight can leap from its final square back to its start, which is possible only for even $n$ at least $6$. Requiring one means rejecting a completed tour whose last square is not a knight's move from the first, so the same search finds it by adding that test at the end.

\textbf{Time Complexity}: $O(n^{2})$ per call in practice, since the move ordering finds a tour without backtracking on the boards where one exists. The worst case remains exponential.

\textbf{Space Complexity}: $O(n^{2})$ auxiliary stack space and $O(n^{2})$ for the returned grid.

\begin{lstlisting}[language=C++]
#include <algorithm>
#include <array>
#include <cassert>
#include <utility>
#include <vector>

const std::array<std::pair<int, int>, 8> KNIGHT_MOVES = {
    {{-2, -1}, {-2, 1}, {-1, -2}, {-1, 2}, {1, -2}, {1, 2}, {2, -1}, {2, 1}}
};

std::vector<std::vector<int>> knight_tour(int n, int sr = 0, int sc = 0) {
  assert(n > 0 && sr >= 0 && sr < n && sc >= 0 && sc < n);
  std::vector<std::vector<int>> visit(n, std::vector<int>(n, -1));
  auto onward_moves = [&](int r, int c) {
    int count = 0;
    for (auto [dr, dc] : KNIGHT_MOVES) {
      int nr = r + dr, nc = c + dc;
      if (nr >= 0 && nr < n && nc >= 0 && nc < n && visit[nr][nc] == -1) {
        count++;
      }
    }
    return count;
  };
  auto rec = [&](auto &&rec, int r, int c, int step) -> bool {
    visit[r][c] = step;
    if (step == n * n - 1) {
      return true;
    }
    // Order the onward squares by how few exits they have left, which is Warnsdorff's rule.
    std::vector<std::pair<int, std::pair<int, int>>> candidates;
    for (auto [dr, dc] : KNIGHT_MOVES) {
      int r2 = r + dr, c2 = c + dc;
      if (r2 >= 0 && r2 < n && c2 >= 0 && c2 < n && visit[r2][c2] == -1) {
        candidates.emplace_back(onward_moves(r2, c2), std::make_pair(r2, c2));
      }
    }
    std::sort(candidates.begin(), candidates.end());
    for (auto [degree, square] : candidates) {
      if (rec(rec, square.first, square.second, step + 1)) {
        return true;
      }
    }
    visit[r][c] = -1;
    return false;
  };
  return rec(rec, sr, sc, 0) ? visit : std::vector<std::vector<int>>{};
}
\end{lstlisting}
\Needspace*{4\baselineskip}
\textbf{Example Usage}\par
\begin{lstlisting}[language=C++]
#include <cassert>
using namespace std;

int main() {
  assert(knight_tour(1) == vector<vector<int>>{{0}});
  assert(knight_tour(3).empty());  // No tour exists on boards of size 2, 3, or 4.
  assert(knight_tour(4).empty());

  for (int n : {5, 6, 8}) {
    auto visit = knight_tour(n, n / 2, n / 2);
    assert(!visit.empty());
    vector<pair<int, int>> at(n * n);
    for (int r = 0; r < n; r++) {
      for (int c = 0; c < n; c++) {
        assert(visit[r][c] >= 0 && visit[r][c] < n * n);
        at[visit[r][c]] = {r, c};  // Every step number is used exactly once.
      }
    }
    assert(at[0] == make_pair(n / 2, n / 2));
    for (int step = 1; step < n * n; step++) {
      int dr = abs(at[step].first - at[step - 1].first);
      int dc = abs(at[step].second - at[step - 1].second);
      assert(dr * dc == 2);  // Each consecutive pair differs by a knight's move.
    }
  }
  return 0;
}
\end{lstlisting}
\subsection{Exact Cover (Dancing Links)}
Given a binary matrix, the exact cover problem asks for a set of rows in which every column contains exactly one $1$. Many search problems reduce to it: a sudoku grid becomes a matrix whose columns demand that each cell hold one digit and that each digit appear once per row, column, and box, while a tiling becomes a matrix whose columns demand that each cell be covered by exactly one piece placement.

Knuth's Algorithm X solves it by repeatedly choosing an uncovered column, trying each row that covers it, and deleting from the matrix every column that row satisfies along with every other row conflicting with it. Dancing links is the representation that makes the deletions cheap. Each $1$ of the matrix is a node in two circular doubly linked lists, one along its row and one along its column, and every column has a header node holding the count of its remaining nodes. Removing a node \inlinecode{x} from a list is \inlinecode{left[right[x]] = left[x]} and \inlinecode{right[left[x]] = right[x]}, which leaves \inlinecode{x} itself still pointing at its old neighbors; restoring it on backtrack is therefore just \inlinecode{left[right[x]] = x} and \inlinecode{right[left[x]] = x}, with no allocation and no search. Undoing the covering operations in the exact reverse of the order that performed them restores the matrix exactly, which is what makes the whole search run in place on one static set of nodes.

The column to branch on is the one with the fewest remaining rows. This is the minimum-remaining- values heuristic, and it fails as early as possible: a column with no remaining rows ends the branch immediately, and a column with one forces the choice.

Secondary columns spare a problem with optional requirements from padding its matrix with one slack row per optional column. In the N-queens reduction below, every board row and column must hold exactly one queen and so is primary, while a diagonal may hold none and so is secondary. Asking for two solutions is the usual test of uniqueness, as a sudoku generator does when validating a puzzle.

\begin{compactitem}
  \item \inlinecode{ExactCover(primary, secondary = 0)} constructs an empty matrix with \inlinecode{primary + secondary} columns, both counts nonnegative. Columns in $[0, {\inlinecodemath{primary}})$ must be covered exactly once, while the optional columns in $[{\inlinecodemath{primary}}, {\inlinecodemath{primary}} + {\inlinecodemath{secondary}})$ may be covered at most once.
  \item \inlinecode{add\_row(cols)} appends a row containing a $1$ in each column of \inlinecode{cols}, which must all be in $[0, {\inlinecodemath{primary}} + {\inlinecodemath{secondary}})$, and returns the row's ID. The IDs of the first $r$ rows added are $[0, r)$, in the order they were added.
  \item \inlinecode{solve(limit = 1)} returns up to \inlinecode{limit} solutions, where \inlinecode{limit} is positive. Each solution is the sorted list of the IDs of its rows, and an empty result means no exact cover exists. The matrix is left unchanged, so it may be queried again.
\end{compactitem}

Solutions come in the order the search finds them, which follows the branching heuristic rather than row IDs.

\Needspace*{3\baselineskip}
\textbf{Time Complexity}:
\begin{compactitem}
  \item $O(k)$ per call to \inlinecode{add\_row(cols)}, where $k$ is the size of \inlinecode{cols}.
  \item Exponential per call to \inlinecode{solve()}. Exact cover is NP-complete, and dancing links prunes the search rather than avoiding it: the cost is proportional to the number of nodes the search unlinks and relinks, which the branching heuristic reduces but does not bound polynomially.
\end{compactitem}

\Needspace*{3\baselineskip}
\textbf{Space Complexity}:
\begin{compactitem}
  \item $O(n + c)$ for storage, where $n$ is the number of $1$s in the matrix and $c$ is the number of columns.
  \item $O(d)$ auxiliary stack space per call to \inlinecode{solve()}, where $d$ is the number of rows in the deepest partial solution reached, plus $O(s \cdot d)$ for the returned solutions, where $s$ is the number found.
\end{compactitem}

\begin{lstlisting}[language=C++]
#include <algorithm>
#include <cassert>
#include <vector>

class ExactCover {
  int columns, rows;
  std::vector<int> left, right, up, down, col_of, row_of, col_size;
  std::vector<int> partial;

  int add_node(int c, int r) {
    left.push_back(0);
    right.push_back(0);
    up.push_back(0);
    down.push_back(0);
    col_of.push_back(c);
    row_of.push_back(r);
    return static_cast<int>(col_of.size()) - 1;
  }

  // Unlinks column c and every row meeting it, so both vanish from all remaining traversals.
  void cover(int c) {
    right[left[c]] = right[c];
    left[right[c]] = left[c];
    for (int i = down[c]; i != c; i = down[i]) {
      for (int j = right[i]; j != i; j = right[j]) {
        up[down[j]] = up[j];
        down[up[j]] = down[j];
        col_size[col_of[j]]--;
      }
    }
  }

  // Relinks in the mirror order of cover(), which is what restores the matrix exactly.
  void uncover(int c) {
    for (int i = up[c]; i != c; i = up[i]) {
      for (int j = left[i]; j != i; j = left[j]) {
        col_size[col_of[j]]++;
        up[down[j]] = j;
        down[up[j]] = j;
      }
    }
    right[left[c]] = c;
    left[right[c]] = c;
  }

  bool search(int limit, std::vector<std::vector<int>> &found) {
    if (right[0] == 0) {  // Every primary column is covered.
      found.push_back(partial);
      std::sort(found.back().begin(), found.back().end());
      return static_cast<int>(found.size()) >= limit;
    }
    int best = right[0];
    for (int c = right[0]; c != 0; c = right[c]) {
      if (col_size[c] < col_size[best]) {
        best = c;
      }
    }
    if (col_size[best] == 0) {
      return false;  // No row can cover this column, so this branch is dead.
    }
    cover(best);
    for (int r = down[best]; r != best; r = down[r]) {
      partial.push_back(row_of[r]);
      for (int j = right[r]; j != r; j = right[j]) {
        cover(col_of[j]);
      }
      bool done = search(limit, found);
      for (int j = left[r]; j != r; j = left[j]) {
        uncover(col_of[j]);
      }
      partial.pop_back();
      if (done) {
        break;
      }
    }
    uncover(best);
    return static_cast<int>(found.size()) >= limit;
  }

 public:
  explicit ExactCover(int primary, int secondary = 0) : columns(primary + secondary), rows(0) {
    assert(primary >= 0 && secondary >= 0);
    col_size.assign(columns + 1, 0);
    for (int c = 0; c <= columns; c++) {  // Node 0 is the root, and node c + 1 heads column c.
      int header = add_node(c, -1);
      up[header] = down[header] = header;
      left[header] = right[header] = header;
    }
    int prev = 0;  // Secondary headers stay outside the root ring, so they are never branched on.
    for (int c = 1; c <= primary; c++) {
      right[prev] = c;
      left[c] = prev;
      prev = c;
    }
    right[prev] = 0;
    left[0] = prev;
  }

  int add_row(const std::vector<int> &cols) {
    for (int c : cols) {
      assert(c >= 0 && c < columns);
    }
    int row = rows++, first = -1;
    for (int c : cols) {
      int header = c + 1, node = add_node(header, row);
      up[node] = up[header];
      down[node] = header;
      down[up[header]] = node;
      up[header] = node;
      col_size[header]++;
      if (first == -1) {
        first = node;
        left[node] = right[node] = node;
      } else {
        left[node] = left[first];
        right[node] = first;
        right[left[first]] = node;
        left[first] = node;
      }
    }
    return row;
  }

  std::vector<std::vector<int>> solve(int limit = 1) {
    assert(limit > 0);
    std::vector<std::vector<int>> found;
    partial.clear();
    search(limit, found);
    return found;
  }
};
\end{lstlisting}
\Needspace*{4\baselineskip}
\textbf{Example Usage}\par
\begin{lstlisting}[language=C++]
#include <cassert>
using namespace std;

int main() {
  // Knuth's example matrix over 7 columns, whose unique cover is rows 1, 3, and 5.
  //     col: 0123456
  //   row 0: 1001001
  //   row 1: 1001000
  //   row 2: 0001101
  //   row 3: 0010110
  //   row 4: 0110011
  //   row 5: 0100001
  ExactCover matrix(7);
  matrix.add_row(vector<int>{0, 3, 6});
  matrix.add_row(vector<int>{0, 3});
  matrix.add_row(vector<int>{3, 4, 6});
  matrix.add_row(vector<int>{2, 4, 5});
  matrix.add_row(vector<int>{1, 2, 5, 6});
  matrix.add_row(vector<int>{1, 6});
  assert((matrix.solve() == vector<vector<int>>{{1, 3, 5}}));
  assert(matrix.solve(2).size() == 1);  // Asking for a second solution proves uniqueness.

  ExactCover impossible(2);
  impossible.add_row(vector<int>{0});
  assert(impossible.solve().empty());

  // The 8-queens problem. Each board row and each board column must hold exactly one queen, while
  // each diagonal may hold at most one, so the 2*n board lines are primary and the 2*(2*n - 1)
  // diagonals are secondary. One matrix row per square places a queen there.
  const int n = 8;
  ExactCover queens(2 * n, 2 * (2 * n - 1));
  for (int r = 0; r < n; r++) {
    for (int c = 0; c < n; c++) {
      queens.add_row(vector<int>{r, n + c, 2 * n + r + c, 2 * n + (2 * n - 1) + r - c + n - 1});
    }
  }
  assert(queens.solve().size() == 1);
  assert(queens.solve(1000).size() == 92);
  return 0;
}
\end{lstlisting}

\chapter{Data Structures}

\section{Linked Lists}
\setcounter{section}{1}
\setcounter{subsection}{0}
\subsection{Singly Linked List Operations}
Implements common singly linked list operations using raw nodes. Linked lists are rarely the best contest choice compared with arrays, vectors, and \inlinecode{std::list}, but the pointer patterns are useful for interview-style problems and for understanding constant-time splicing.

For production C++, prefer containers such as \inlinecode{std::list}, \inlinecode{std::forward\_list}, or \inlinecode{std::vector} when they fit the problem. Use the raw-node style below when the problem statement already gives nodes, or when manual pointer manipulation is the point of the exercise.

\begin{compactitem}
  \item \inlinecode{reverse\_list(head)} reverses a singly linked list and returns the new head.
  \item \inlinecode{merge\_sorted\_lists(a, b)} merges two sorted lists using a dummy node and returns the merged head.
  \item \inlinecode{split\_half(head, \&second)} cuts a list into two halves, returning the first half and storing the second half in \inlinecode{second}.
\end{compactitem}

For linked-list cycle detection, use Floyd or Brent cycle detection from section 1.5.

\Needspace*{3\baselineskip}
\textbf{Time Complexity}:
\begin{compactitem}
  \item $O(n)$ per call to \inlinecode{reverse\_list(head)} and \inlinecode{split\_half(head, \&second)}.
  \item $O(n + m)$ per call to \inlinecode{merge\_sorted\_lists(a, b)}.
\end{compactitem}

\textbf{Space Complexity}: $O(1)$ auxiliary for all operations.

\begin{lstlisting}[language=C++]
#include <cstddef>

struct ListNode {
  int value;
  ListNode *next;

  explicit ListNode(int value = 0) : value(value), next(nullptr) {}
};

ListNode *reverse_list(ListNode *head) {
  ListNode *prev = nullptr;
  while (head != nullptr) {
    ListNode *next = head->next;
    head->next = prev;
    prev = head;
    head = next;
  }
  return prev;
}

ListNode *merge_sorted_lists(ListNode *a, ListNode *b) {
  ListNode dummy;
  ListNode *tail = &dummy;
  while (a != nullptr && b != nullptr) {
    if (b->value < a->value) {
      tail->next = b;
      b = b->next;
    } else {
      tail->next = a;
      a = a->next;
    }
    tail = tail->next;
  }
  tail->next = a != nullptr ? a : b;
  return dummy.next;
}

ListNode *split_half(ListNode *head, ListNode **second) {
  if (head == nullptr || head->next == nullptr) {
    *second = nullptr;
    return head;
  }
  ListNode *slow = head, *fast = head->next;
  while (fast != nullptr && fast->next != nullptr) {
    slow = slow->next;
    fast = fast->next->next;
  }
  *second = slow->next;
  slow->next = nullptr;
  return head;
}
\end{lstlisting}
\Needspace*{4\baselineskip}
\textbf{Example Usage}\par
\begin{lstlisting}[language=C++]
#include <cassert>
#include <vector>
using namespace std;

vector<int> values(ListNode *head) {
  vector<int> res;
  for (; head != nullptr; head = head->next) {
    res.push_back(head->value);
  }
  return res;
}

int main() {
  vector<ListNode> a{ListNode(1), ListNode(3), ListNode(5)};
  vector<ListNode> b{ListNode(2), ListNode(4), ListNode(6)};
  for (int i = 0; i + 1 < static_cast<int>(a.size()); i++) {
    a[i].next = &a[i + 1];
    b[i].next = &b[i + 1];
  }
  ListNode *merged = merge_sorted_lists(&a[0], &b[0]);
  vector<int> merged_values = values(merged);
  for (int i = 0; i < 6; i++) {
    assert(merged_values[i] == i + 1);
  }

  ListNode *second = nullptr;
  ListNode *first = split_half(merged, &second);
  assert(values(first).size() == 3);
  assert(values(second).size() == 3);
  assert(values(reverse_list(first))[0] == 3);
  return 0;
}
\end{lstlisting}
\subsection{Linked List Splicing}
Moves nodes or ranges between singly linked lists in $O(1)$ pointer changes once the neighboring nodes are known. Splicing is the main operation that makes linked lists useful: it moves existing nodes without copying values or invalidating pointers to the moved nodes.

The functions below use a dummy head node for each list. This makes insertions and removals at the beginning of a list match all other positions. In the C++ standard library, \inlinecode{std::list::splice()} provides analogous operations for doubly linked lists, while \inlinecode{std::forward\_list::splice\_after()} provides analogous operations for singly linked lists.

\begin{compactitem}
  \item \inlinecode{splice\_after(pos, before)} moves the node after \inlinecode{before} so it appears after \inlinecode{pos}.
  \item \inlinecode{splice\_range\_after(pos, before\_first, last)} moves the half-open range \inlinecode{(before\_first, last)} so it appears after \inlinecode{pos}.
\end{compactitem}

\Needspace*{3\baselineskip}
\textbf{Time Complexity}:
\begin{compactitem}
  \item $O(1)$ per call to \inlinecode{splice\_after(pos, before)}.
  \item $O(k)$ per call to \inlinecode{splice\_range\_after(pos, before\_first, last)}, where $k$ is the number of moved nodes. The actual pointer splicing is $O(1)$, but this compact version walks to the end of the moved range.
\end{compactitem}

\textbf{Space Complexity}: $O(1)$ auxiliary.

\begin{lstlisting}[language=C++]
#include <cstddef>

struct ListNode {
  int value;
  ListNode *next;

  explicit ListNode(int value = 0) : value(value), next(nullptr) {}
};

void splice_after(ListNode *pos, ListNode *before) {
  ListNode *node = before->next;
  if (node == nullptr || pos == before || pos == node) {
    return;
  }
  before->next = node->next;
  node->next = pos->next;
  pos->next = node;
}

void splice_range_after(ListNode *pos, ListNode *before_first, ListNode *last) {
  ListNode *first = before_first->next;
  if (first == last) {
    return;
  }
  ListNode *tail = first;
  while (tail->next != last) {
    tail = tail->next;
  }
  before_first->next = last;
  tail->next = pos->next;
  pos->next = first;
}
\end{lstlisting}
\Needspace*{4\baselineskip}
\textbf{Example Usage}\par
\begin{lstlisting}[language=C++]
#include <cassert>
#include <vector>
using namespace std;

vector<int> values(ListNode *dummy) {
  vector<int> res;
  for (ListNode *p = dummy->next; p != nullptr; p = p->next) {
    res.push_back(p->value);
  }
  return res;
}

int main() {
  ListNode dummy(0), a(1), b(2), c(3), d(4), e(5);
  dummy.next = &a;
  a.next = &b;
  b.next = &c;
  c.next = &d;
  d.next = &e;

  splice_after(&dummy, &c);  // Move 4 to the front: 4, 1, 2, 3, 5.
  vector<int> v = values(&dummy);
  vector<int> expected1{4, 1, 2, 3, 5};
  for (int i = 0; i < static_cast<int>(expected1.size()); i++) {
    assert(v[i] == expected1[i]);
  }

  splice_range_after(&d, &a, &e);  // Move 2, 3 after 4: 4, 2, 3, 1, 5.
  v = values(&dummy);
  vector<int> expected2{4, 2, 3, 1, 5};
  for (int i = 0; i < static_cast<int>(expected2.size()); i++) {
    assert(v[i] == expected2[i]);
  }
  return 0;
}
\end{lstlisting}
\subsection{LRU Cache}
Maintains a least-recently-used cache with fixed capacity. The cache supports lookups and updates by key, evicting the least recently used entry whenever an insertion would exceed capacity.

This implementation combines a doubly linked list (\inlinecode{std::list}) with a map from keys to list iterators. The front of the list stores the most recently used item. The operation \inlinecode{items.splice(items.begin(), items, it)} is the standard-library way to move an existing list node to the front in $O(1)$ time.

\begin{compactitem}
  \item \inlinecode{LRUCache(capacity)} constructs an empty cache holding at most \inlinecode{capacity} keys.
  \item \inlinecode{get(key, \&value)} returns whether \inlinecode{key} is present. If present, it stores the value in \inlinecode{value} and marks the key as most recently used.
  \item \inlinecode{put(key, value)} inserts or updates \inlinecode{key}, marking it as most recently used and evicting the least recently used key if needed.
  \item \inlinecode{size()} returns the number of keys currently stored.
\end{compactitem}

\textbf{Time Complexity}: $O(1)$ amortized per call to \inlinecode{get(key, \&value)} and \inlinecode{put(key, value)}.

\textbf{Space Complexity}: $O(n)$, where $n$ is the cache capacity.

\begin{lstlisting}[language=C++]
#include <list>
#include <unordered_map>
#include <utility>

template<class Key, class Value>
class LRUCache {
  using ListIter = typename std::list<std::pair<Key, Value>>::iterator;

  int cap;
  std::list<std::pair<Key, Value>> items;
  std::unordered_map<Key, ListIter> where;

 public:
  explicit LRUCache(int capacity) : cap(capacity) {}

  int size() const { return items.size(); }

  bool get(const Key &key, Value *value) {
    auto it = where.find(key);
    if (it == where.end()) {
      return false;
    }
    items.splice(items.begin(), items, it->second);
    it->second = items.begin();
    *value = it->second->second;
    return true;
  }

  void put(const Key &key, const Value &value) {
    auto it = where.find(key);
    if (it != where.end()) {
      it->second->second = value;
      items.splice(items.begin(), items, it->second);
      it->second = items.begin();
      return;
    }
    if (cap == 0) {
      return;
    }
    if (static_cast<int>(items.size()) == cap) {
      where.erase(items.back().first);
      items.pop_back();
    }
    items.push_front({key, value});
    where[key] = items.begin();
  }
};
\end{lstlisting}
\Needspace*{4\baselineskip}
\textbf{Example Usage}\par
\begin{lstlisting}[language=C++]
#include <cassert>
using namespace std;

int main() {
  LRUCache<int, int> cache(2);
  int value = 0;
  cache.put(1, 10);
  cache.put(2, 20);
  assert(cache.get(1, &value) && value == 10);
  cache.put(3, 30);  // Evicts key 2.
  assert(!cache.get(2, &value));
  assert(cache.get(3, &value) && value == 30);
  cache.put(1, 11);
  assert(cache.get(1, &value) && value == 11);
  return 0;
}
\end{lstlisting}

\section{Heaps and Priority Queues}
\setcounter{section}{2}
\setcounter{subsection}{0}
\subsection{Binary Heap}
Maintain a min-priority queue, that is, a collection of elements with support for querying and extraction of the minimum. This implementation requires an ordering on the set of possible elements defined by \inlinecode{operator\textless{}}. A binary min-heap implements a priority queue by inserting and deleting nodes into a binary tree such that the parent of any node is always less than its children.

\begin{compactitem}
  \item \inlinecode{BinaryHeap()} constructs an empty priority queue.
  \item \inlinecode{BinaryHeap(lo, hi)} constructs a priority queue from two ForwardIterators, consisting of elements in the range \inlinecode{[lo, hi)}.
  \item \inlinecode{size()} returns the size of the priority queue.
  \item \inlinecode{empty()} returns whether the priority queue is empty.
  \item \inlinecode{push(v)} inserts the value \inlinecode{v} into the priority queue.
  \item \inlinecode{pop()} removes the minimum element from the priority queue.
  \item \inlinecode{top()} returns the minimum element in the priority queue.
\end{compactitem}

\Needspace*{3\baselineskip}
\textbf{Time Complexity}:
\begin{compactitem}
  \item $O(1)$ per call to the first constructor, \inlinecode{size()}, \inlinecode{empty()}, and \inlinecode{top()}.
  \item $O(\log n)$ per call to \inlinecode{push()} and \inlinecode{pop()}, where $n$ is the number of elements in the priority queue.
  \item $O(n)$ per call to the second constructor on the distance between \inlinecode{lo} and \inlinecode{hi} (bottom-up heapify).
\end{compactitem}

\Needspace*{3\baselineskip}
\textbf{Space Complexity}:
\begin{compactitem}
  \item $O(n)$ for storage of the priority queue elements.
  \item $O(1)$ auxiliary for all operations.
\end{compactitem}

\begin{lstlisting}[language=C++]
#include <algorithm>
#include <stdexcept>
#include <vector>

template<class T>
class BinaryHeap {
  std::vector<T> heap;

  void sift_down(int i) {
    for (;;) {
      int child = 2 * i + 1;
      if (child >= static_cast<int>(heap.size())) {
        break;
      }
      if (child + 1 < static_cast<int>(heap.size()) && heap[child + 1] < heap[child]) {
        child++;
      }
      if (heap[child] < heap[i]) {
        std::swap(heap[i], heap[child]);
        i = child;
      } else {
        break;
      }
    }
  }

 public:
  BinaryHeap() {}

  template<class It>
  BinaryHeap(It lo, It hi) : heap(lo, hi) {
    for (int i = static_cast<int>(heap.size()) / 2 - 1; i >= 0; i--) {
      sift_down(i);
    }
  }

  int size() const { return heap.size(); }
  bool empty() const { return heap.empty(); }

  void push(const T &v) {
    heap.push_back(v);
    int i = heap.size() - 1;
    while (i > 0) {
      int parent = (i - 1) / 2;
      if (!(heap[i] < heap[parent])) {
        break;
      }
      std::swap(heap[i], heap[parent]);
      i = parent;
    }
  }

  void pop() {
    if (heap.empty()) {
      throw std::runtime_error("Cannot pop from empty heap.");
    }
    heap[0] = heap.back();
    heap.pop_back();
    if (!heap.empty()) {
      sift_down(0);
    }
  }

  T top() const {
    if (heap.empty()) {
      throw std::runtime_error("Cannot get top of empty heap.");
    }
    return heap[0];
  }
};
\end{lstlisting}
\Needspace*{4\baselineskip}
\textbf{Example Usage}\par
\begin{lstlisting}[language=C++]
#include <iostream>
using namespace std;

int main() {
  vector<int> a{0, 5, -1, 12};
  BinaryHeap<int> h(a.begin(), a.end());
  h.push(10);
  while (!h.empty()) {
    cout << h.top() << endl;
    h.pop();
  }
  return 0;
}
\end{lstlisting}
\begin{exampleoutput}
-1
0
5
10
12
\end{exampleoutput}
\subsection{Skew Heap}
Maintain a mergeable min-priority queue, that is, a collection of elements with support for querying and extraction of the minimum as well as efficient merging with other instances. This implementation requires an ordering on the set of possible elements defined by \inlinecode{operator\textless{}}. A skew heap attempts to maintain balance by unconditionally swapping all nodes in the merge path when merging.

\begin{compactitem}
  \item \inlinecode{SkewHeap()} constructs an empty priority queue.
  \item \inlinecode{SkewHeap(lo, hi)} constructs a priority queue from two ForwardIterators, consisting of elements in the range \inlinecode{[lo, hi)}.
  \item \inlinecode{size()} returns the size of the priority queue.
  \item \inlinecode{empty()} returns whether the priority queue is empty.
  \item \inlinecode{push(v)} inserts the value \inlinecode{v} into the priority queue.
  \item \inlinecode{pop()} removes the minimum element from the priority queue.
  \item \inlinecode{top()} returns the minimum element in the priority queue.
  \item \inlinecode{absorb(h)} inserts every value from \inlinecode{h} and sets \inlinecode{h} to the empty priority queue.
\end{compactitem}

\Needspace*{3\baselineskip}
\textbf{Time Complexity}:
\begin{compactitem}
  \item $O(1)$ per call to the first constructor, \inlinecode{size()}, \inlinecode{empty()}, and \inlinecode{top()}.
  \item $O(\log n)$ amortized auxiliary per call to \inlinecode{push()}, \inlinecode{pop()}, and \inlinecode{absorb()}, where $n$ is the number of elements in the priority queue.
  \item $O(n)$ per call to the second constructor, where $n$ is the distance between \inlinecode{lo} and \inlinecode{hi}.
\end{compactitem}

\Needspace*{3\baselineskip}
\textbf{Space Complexity}:
\begin{compactitem}
  \item $O(n)$ for storage of the priority queue elements.
  \item $O(\log n)$ amortized auxiliary stack space for \inlinecode{push()}, \inlinecode{pop()}, and \inlinecode{absorb()}.
  \item $O(1)$ auxiliary for all other operations.
\end{compactitem}

\begin{lstlisting}[language=C++]
#include <algorithm>
#include <cstddef>
#include <stdexcept>

template<class T>
class SkewHeap {
  struct Node {
    T value;
    Node *left, *right;

    explicit Node(const T &v) : value(v), left(nullptr), right(nullptr) {}
  } *root;

  int num_nodes;

  static Node *merge(Node *a, Node *b) {
    if (a == nullptr) {
      return b;
    }
    if (b == nullptr) {
      return a;
    }
    if (b->value < a->value) {
      std::swap(a, b);
    }
    std::swap(a->left, a->right);
    a->left = merge(b, a->left);
    return a;
  }

  static void clean_up(Node *n) {
    if (n != nullptr) {
      clean_up(n->left);
      clean_up(n->right);
      delete n;
    }
  }

 public:
  SkewHeap() : root(nullptr), num_nodes(0) {}

  template<class It>
  SkewHeap(It lo, It hi) : root(nullptr), num_nodes(0) {
    while (lo != hi) {
      push(*(lo++));
    }
  }

  ~SkewHeap() { clean_up(root); }
  int size() const { return num_nodes; }
  bool empty() const { return root == nullptr; }

  void push(const T &v) {
    root = merge(root, new Node(v));
    num_nodes++;
  }

  void pop() {
    if (empty()) {
      throw std::runtime_error("Cannot pop from empty heap.");
    }
    Node *tmp = root;
    root = merge(root->left, root->right);
    delete tmp;
    num_nodes--;
  }

  T top() const {
    if (empty()) {
      throw std::runtime_error("Cannot get top of empty heap.");
    }
    return root->value;
  }

  void absorb(SkewHeap &h) {
    root = merge(root, h.root);
    num_nodes += h.num_nodes;
    h.root = nullptr;
    h.num_nodes = 0;
  }
};
\end{lstlisting}
\Needspace*{4\baselineskip}
\textbf{Example Usage}\par
\begin{lstlisting}[language=C++]
#include <iostream>
using namespace std;

int main() {
  SkewHeap<int> h, h2;
  h.push(12);
  h.push(10);
  h2.push(5);
  h2.push(-1);
  h2.push(0);
  h.absorb(h2);
  while (!h.empty()) {
    cout << h.top() << endl;
    h.pop();
  }
  return 0;
}
\end{lstlisting}
\begin{exampleoutput}
-1
0
5
10
12
\end{exampleoutput}
\subsection{Pairing Heap}
Maintain a mergeable min-priority queue, that is, a collection of elements with support for querying and extraction of the minimum as well as efficient merging with other instances. This implementation requires an ordering on the set of possible elements defined by \inlinecode{operator\textless{}}. A pairing heap is a heap-ordered multi-way tree, using a two-pass merge to self-adjust during each deletion.

\begin{compactitem}
  \item \inlinecode{PairingHeap()} constructs an empty priority queue.
  \item \inlinecode{PairingHeap(lo, hi)} constructs a priority queue from two ForwardIterators, consisting of elements in the range \inlinecode{[lo, hi)}.
  \item \inlinecode{size()} returns the size of the priority queue.
  \item \inlinecode{empty()} returns whether the priority queue is empty.
  \item \inlinecode{push(v)} inserts the value \inlinecode{v} into the priority queue.
  \item \inlinecode{pop()} removes the minimum element from the priority queue.
  \item \inlinecode{top()} returns the minimum element in the priority queue.
  \item \inlinecode{absorb(h)} inserts every value from \inlinecode{h} and sets \inlinecode{h} to the empty priority queue.
\end{compactitem}

\Needspace*{3\baselineskip}
\textbf{Time Complexity}:
\begin{compactitem}
  \item $O(1)$ per call to the first constructor, \inlinecode{size()}, \inlinecode{empty()}, \inlinecode{top()}, \inlinecode{push()}, and \inlinecode{absorb()}.
  \item $O(\log n)$ amortized per call to \inlinecode{pop()}.
  \item $O(n)$ per call to the second constructor on the distance between \inlinecode{lo} and \inlinecode{hi}.
\end{compactitem}

\Needspace*{3\baselineskip}
\textbf{Space Complexity}:
\begin{compactitem}
  \item $O(n)$ for storage of the priority queue elements.
  \item $O(\log n)$ amortized auxiliary stack space for \inlinecode{pop()}.
  \item $O(1)$ auxiliary for all other operations.
\end{compactitem}

\begin{lstlisting}[language=C++]
#include <cstddef>
#include <stdexcept>

template<class T>
class PairingHeap {
  struct Node {
    T value;
    Node *left, *next;

    explicit Node(const T &v) : value(v), left(nullptr), next(nullptr) {}

    void add_child(Node *n) {
      if (left == nullptr) {
        left = n;
      } else {
        n->next = left;
        left = n;
      }
    }
  } *root;

  int num_nodes;

  static Node *merge(Node *a, Node *b) {
    if (a == nullptr) {
      return b;
    }
    if (b == nullptr) {
      return a;
    }
    if (a->value < b->value) {
      a->add_child(b);
      return a;
    }
    b->add_child(a);
    return b;
  }

  static Node *merge_pairs(Node *n) {
    if (n == nullptr || n->next == nullptr) {
      return n;
    }
    Node *a = n, *b = n->next, *c = n->next->next;
    a->next = b->next = nullptr;
    return merge(merge(a, b), merge_pairs(c));
  }

  static void clean_up(Node *n) {
    if (n != nullptr) {
      clean_up(n->left);
      clean_up(n->next);
      delete n;
    }
  }

 public:
  PairingHeap() : root(nullptr), num_nodes(0) {}

  template<class It>
  PairingHeap(It lo, It hi) : root(nullptr), num_nodes(0) {
    while (lo != hi) {
      push(*(lo++));
    }
  }

  ~PairingHeap() { clean_up(root); }
  int size() const { return num_nodes; }
  bool empty() const { return root == nullptr; }

  void push(const T &v) {
    root = merge(root, new Node(v));
    num_nodes++;
  }

  void pop() {
    if (empty()) {
      throw std::runtime_error("Cannot pop from empty heap.");
    }
    Node *tmp = root;
    root = merge_pairs(root->left);
    delete tmp;
    num_nodes--;
  }

  T top() const {
    if (empty()) {
      throw std::runtime_error("Cannot get top of empty heap.");
    }
    return root->value;
  }

  void absorb(PairingHeap &h) {
    root = merge(root, h.root);
    num_nodes += h.num_nodes;
    h.root = nullptr;
    h.num_nodes = 0;
  }
};
\end{lstlisting}
\Needspace*{4\baselineskip}
\textbf{Example Usage}\par
\begin{lstlisting}[language=C++]
#include <iostream>
using namespace std;

int main() {
  PairingHeap<int> h, h2;
  h.push(12);
  h.push(10);
  h2.push(5);
  h2.push(-1);
  h2.push(0);
  h.absorb(h2);
  while (!h.empty()) {
    cout << h.top() << endl;
    h.pop();
  }
  return 0;
}
\end{lstlisting}
\begin{exampleoutput}
-1
0
5
10
12
\end{exampleoutput}

\section{Dictionaries and Ordered Sets}
\setcounter{section}{3}
\setcounter{subsection}{0}
\subsection{Binary Search Tree}
Maintain a map, that is, a collection of key-value pairs such that each possible key appears at most once in the collection. This implementations requires an ordering on the set of possible keys defined by \inlinecode{operator\textless{}} on the key type. A binary search tree implements this map by inserting and deleting keys into a binary tree such that every node's left child has a lesser key and every node's right child has a greater key.

\begin{compactitem}
  \item \inlinecode{BST()} constructs an empty map.
  \item \inlinecode{size()} returns the size of the map.
  \item \inlinecode{empty()} returns whether the map is empty.
  \item \inlinecode{insert(k, v)} adds an entry with key \inlinecode{k} and value \inlinecode{v} to the map, returning \inlinecode{true} if an new entry was added or \inlinecode{false} if the key already exists (in which case the map is unchanged and the old value associated with the key is preserved).
  \item \inlinecode{erase(k)} removes the entry with key \inlinecode{k} from the map, returning \inlinecode{true} if the removal was successful or \inlinecode{false} if the key to be removed was not found.
  \item \inlinecode{find(k)} returns a pointer to a const value associated with key \inlinecode{k}, or \inlinecode{nullptr} if the key was not found.
  \item \inlinecode{walk(f)} calls the function \inlinecode{f(k, v)} on each entry of the map, in ascending order of keys.
\end{compactitem}

\Needspace*{3\baselineskip}
\textbf{Time Complexity}:
\begin{compactitem}
  \item $O(1)$ per call to the constructor, \inlinecode{size()}, and \inlinecode{empty()}.
  \item $O(n)$ per call to \inlinecode{insert()}, \inlinecode{erase()}, \inlinecode{find()}, and \inlinecode{walk()}, where $n$ is the number of nodes currently in the map.
\end{compactitem}

\Needspace*{3\baselineskip}
\textbf{Space Complexity}:
\begin{compactitem}
  \item $O(n)$ for storage of the map elements.
  \item $O(n)$ auxiliary stack space for \inlinecode{insert()}, \inlinecode{erase()}, and \inlinecode{walk()}.
  \item $O(1)$ auxiliary for all other operations.
\end{compactitem}

\begin{lstlisting}[language=C++]
#include <cstddef>

template<class K, class V>
class BST {
  struct Node {
    K key;
    V value;
    Node *left, *right;

    Node(const K &k, const V &v) : key(k), value(v), left(nullptr), right(nullptr) {}
  } *root;

  int num_nodes;

  static bool insert(Node *&n, const K &k, const V &v) {
    if (n == nullptr) {
      n = new Node(k, v);
      return true;
    }
    if (k < n->key) {
      return insert(n->left, k, v);
    } else if (n->key < k) {
      return insert(n->right, k, v);
    }
    return false;
  }

  static bool erase(Node *&n, const K &k) {
    if (n == nullptr) {
      return false;
    }
    if (k < n->key) {
      return erase(n->left, k);
    } else if (n->key < k) {
      return erase(n->right, k);
    }
    if (n->left != nullptr && n->right != nullptr) {
      Node *tmp = n->right, *parent = nullptr;
      while (tmp->left != nullptr) {
        parent = tmp;
        tmp = tmp->left;
      }
      n->key = tmp->key;
      n->value = tmp->value;
      if (parent != nullptr) {
        return erase(parent->left, parent->left->key);
      }
      return erase(n->right, n->right->key);
    }
    Node *tmp = (n->left != nullptr) ? n->left : n->right;
    delete n;
    n = tmp;
    return true;
  }

  template<class Fn>
  static void walk(Node *n, Fn f) {
    if (n != nullptr) {
      walk(n->left, f);
      f(n->key, n->value);
      walk(n->right, f);
    }
  }

  static void clean_up(Node *n) {
    if (n != nullptr) {
      clean_up(n->left);
      clean_up(n->right);
      delete n;
    }
  }

 public:
  BST() : root(nullptr), num_nodes(0) {}

  ~BST() { clean_up(root); }
  int size() const { return num_nodes; }
  bool empty() const { return root == nullptr; }

  bool insert(const K &k, const V &v) {
    if (insert(root, k, v)) {
      num_nodes++;
      return true;
    }
    return false;
  }

  bool erase(const K &k) {
    if (erase(root, k)) {
      num_nodes--;
      return true;
    }
    return false;
  }

  const V *find(const K &k) const {
    Node *n = root;
    while (n != nullptr) {
      if (k < n->key) {
        n = n->left;
      } else if (n->key < k) {
        n = n->right;
      } else {
        return &(n->value);
      }
    }
    return nullptr;
  }

  template<class Fn>
  void walk(Fn f) const {
    walk(root, f);
  }
};
\end{lstlisting}
\Needspace*{4\baselineskip}
\textbf{Example Usage}\par
\begin{lstlisting}[language=C++]
#include <cassert>
#include <iostream>
using namespace std;

int main() {
  BST<int, char> t;
  t.insert(2, 'b');
  t.insert(1, 'a');
  t.insert(3, 'c');
  t.insert(5, 'e');
  assert(t.insert(4, 'd'));
  assert(*t.find(4) == 'd');
  assert(!t.insert(4, 'd'));
  auto printch = [](int k, char v) { cout << v; };
  t.walk(printch);
  cout << endl;
  assert(t.erase(1));
  assert(!t.erase(1));
  assert(t.find(1) == nullptr);
  t.walk(printch);
  cout << endl;
  return 0;
}
\end{lstlisting}
\begin{exampleoutput}
abcde
bcde
\end{exampleoutput}
\subsection{Treap}
Maintain a map, that is, a collection of key-value pairs such that each possible key appears at most once in the collection. This implementations requires an ordering on the set of possible keys defined by \inlinecode{operator\textless{}} on the key type. A Treap is a binary search tree where each node holds a randomly assigned priority. The tree satisfies the BST property on keys and the min-heap property on priorities (lower priority value is closer to root). Since priorities are random, this keeps the tree balanced with high probability, making insertions and deletions run in $O(\log n)$.

\begin{compactitem}
  \item \inlinecode{Treap()} constructs an empty map.
  \item \inlinecode{size()} returns the size of the map.
  \item \inlinecode{empty()} returns whether the map is empty.
  \item \inlinecode{insert(k, v)} adds an entry with key \inlinecode{k} and value \inlinecode{v} to the map, returning \inlinecode{true} if an new entry was added or \inlinecode{false} if the key already exists (in which case the map is unchanged and the old value associated with the key is preserved).
  \item \inlinecode{erase(k)} removes the entry with key \inlinecode{k} from the map, returning \inlinecode{true} if the removal was successful or \inlinecode{false} if the key to be removed was not found.
  \item \inlinecode{find(k)} returns a pointer to a const value associated with key \inlinecode{k}, or \inlinecode{nullptr} if the key was not found.
  \item \inlinecode{walk(f)} calls the function \inlinecode{f(k, v)} on each entry of the map, in ascending order of keys.
\end{compactitem}

\Needspace*{3\baselineskip}
\textbf{Time Complexity}:
\begin{compactitem}
  \item $O(1)$ per call to the constructor, \inlinecode{size()}, and \inlinecode{empty()}.
  \item $O(\log n)$ on average per call to \inlinecode{insert()}, \inlinecode{erase()}, and \inlinecode{find()}, where $n$ is the number of entries currently in the map.
  \item $O(n)$ per call to \inlinecode{walk()}.
\end{compactitem}

\Needspace*{3\baselineskip}
\textbf{Space Complexity}:
\begin{compactitem}
  \item $O(n)$ for storage of the map elements.
  \item $O(\log n)$ auxiliary stack space on average for \inlinecode{insert()}, \inlinecode{erase()}, and \inlinecode{walk()}.
  \item $O(1)$ auxiliary for all other operations.
\end{compactitem}

\begin{lstlisting}[language=C++]
#include <random>

template<class K, class V>
class Treap {
  struct Node {
    static inline int rand32() {
      static std::mt19937 rng(std::random_device{}());
      return static_cast<int>(rng());
    }

    K key;
    V value;
    int priority;
    Node *left, *right;

    Node(const K &k, const V &v)
        : key(k), value(v), priority(rand32()), left(nullptr), right(nullptr) {}
  } *root;

  int num_nodes;

  static void rotate_left(Node *&n) {
    Node *tmp = n;
    n = n->right;
    tmp->right = n->left;
    n->left = tmp;
  }

  static void rotate_right(Node *&n) {
    Node *tmp = n;
    n = n->left;
    tmp->left = n->right;
    n->right = tmp;
  }

  static bool insert(Node *&n, const K &k, const V &v) {
    if (n == nullptr) {
      n = new Node(k, v);
      return true;
    }
    if (k < n->key && insert(n->left, k, v)) {
      if (n->left->priority < n->priority) {
        rotate_right(n);
      }
      return true;
    }
    if (n->key < k && insert(n->right, k, v)) {
      if (n->right->priority < n->priority) {
        rotate_left(n);
      }
      return true;
    }
    return false;
  }

  static bool erase(Node *&n, const K &k) {
    if (n == nullptr) {
      return false;
    }
    if (k < n->key) {
      return erase(n->left, k);
    } else if (n->key < k) {
      return erase(n->right, k);
    }
    if (n->left != nullptr && n->right != nullptr) {
      if (n->left->priority < n->right->priority) {
        rotate_right(n);
        return erase(n->right, k);
      }
      rotate_left(n);
      return erase(n->left, k);
    }
    Node *tmp = (n->left != nullptr) ? n->left : n->right;
    delete n;
    n = tmp;
    return true;
  }

  template<class Fn>
  static void walk(Node *n, Fn f) {
    if (n != nullptr) {
      walk(n->left, f);
      f(n->key, n->value);
      walk(n->right, f);
    }
  }

  static void clean_up(Node *n) {
    if (n != nullptr) {
      clean_up(n->left);
      clean_up(n->right);
      delete n;
    }
  }

 public:
  Treap() : root(nullptr), num_nodes(0) {}

  ~Treap() { clean_up(root); }
  int size() const { return num_nodes; }
  bool empty() const { return root == nullptr; }

  bool insert(const K &k, const V &v) {
    if (insert(root, k, v)) {
      num_nodes++;
      return true;
    }
    return false;
  }

  bool erase(const K &k) {
    if (erase(root, k)) {
      num_nodes--;
      return true;
    }
    return false;
  }

  const V *find(const K &k) const {
    Node *n = root;
    while (n != nullptr) {
      if (k < n->key) {
        n = n->left;
      } else if (n->key < k) {
        n = n->right;
      } else {
        return &(n->value);
      }
    }
    return nullptr;
  }

  template<class Fn>
  void walk(Fn f) const {
    walk(root, f);
  }
};
\end{lstlisting}
\Needspace*{4\baselineskip}
\textbf{Example Usage}\par
\begin{lstlisting}[language=C++]
#include <cassert>
#include <iostream>
using namespace std;

int main() {
  Treap<int, char> t;
  t.insert(2, 'b');
  t.insert(1, 'a');
  t.insert(3, 'c');
  t.insert(5, 'e');
  assert(t.insert(4, 'd'));
  assert(*t.find(4) == 'd');
  assert(!t.insert(4, 'd'));
  auto printch = [](int k, char v) { cout << v; };
  t.walk(printch);
  cout << endl;
  assert(t.erase(1));
  assert(!t.erase(1));
  assert(t.find(1) == nullptr);
  t.walk(printch);
  cout << endl;
  return 0;
}
\end{lstlisting}
\begin{exampleoutput}
abcde
bcde
\end{exampleoutput}
\subsection{AVL Tree}
Maintain a map, that is, a collection of key-value pairs such that each possible key appears at most once in the collection. This implementation requires an ordering on the set of possible keys defined by \inlinecode{operator\textless{}} on the key type. An AVL tree is a binary search tree balanced by height, guaranteeing $O(\log n)$ worst-case running time in insertions and deletions by making sure that the heights of the left and right subtrees at every node differ by at most 1.

\begin{compactitem}
  \item \inlinecode{AVLTree()} constructs an empty map.
  \item \inlinecode{size()} returns the size of the map.
  \item \inlinecode{empty()} returns whether the map is empty.
  \item \inlinecode{insert(k, v)} adds an entry with key \inlinecode{k} and value \inlinecode{v} to the map, returning \inlinecode{true} if an new entry was added or \inlinecode{false} if the key already exists (in which case the map is unchanged and the old value associated with the key is preserved).
  \item \inlinecode{erase(k)} removes the entry with key \inlinecode{k} from the map, returning \inlinecode{true} if the removal was successful or \inlinecode{false} if the key to be removed was not found.
  \item \inlinecode{find(k)} returns a pointer to a const value associated with key \inlinecode{k}, or \inlinecode{nullptr} if the key was not found.
  \item \inlinecode{walk(f)} calls the function \inlinecode{f(k, v)} on each entry of the map, in ascending order of keys.
\end{compactitem}

\Needspace*{3\baselineskip}
\textbf{Time Complexity}:
\begin{compactitem}
  \item $O(1)$ per call to the constructor, \inlinecode{size()}, and \inlinecode{empty()}.
  \item $O(\log n)$ per call to \inlinecode{insert()}, \inlinecode{erase()}, and \inlinecode{find()}, where $n$ is the number of entries currently in the map.
  \item $O(n)$ per call to \inlinecode{walk()}.
\end{compactitem}

\Needspace*{3\baselineskip}
\textbf{Space Complexity}:
\begin{compactitem}
  \item $O(n)$ for storage of the map elements.
  \item $O(\log n)$ auxiliary stack space for \inlinecode{insert()}, \inlinecode{erase()}, and \inlinecode{walk()}.
  \item $O(1)$ auxiliary for all other operations.
\end{compactitem}

\begin{lstlisting}[language=C++]
#include <algorithm>
#include <cstddef>

template<class K, class V>
class AVLTree {
  struct Node {
    K key;
    V value;
    int height;
    Node *left, *right;

    Node(const K &k, const V &v) : key(k), value(v), height(1), left(nullptr), right(nullptr) {}
  } *root;

  int num_nodes;

  static int height(Node *n) { return (n != nullptr) ? n->height : 0; }

  static void update_height(Node *n) {
    if (n != nullptr) {
      n->height = 1 + std::max(height(n->left), height(n->right));
    }
  }

  static void rotate_left(Node *&n) {
    Node *tmp = n;
    n = n->right;
    tmp->right = n->left;
    n->left = tmp;
    update_height(tmp);
    update_height(n);
  }

  static void rotate_right(Node *&n) {
    Node *tmp = n;
    n = n->left;
    tmp->left = n->right;
    n->right = tmp;
    update_height(tmp);
    update_height(n);
  }

  static int balance_factor(Node *n) {
    return (n != nullptr) ? (height(n->left) - height(n->right)) : 0;
  }

  static void rebalance(Node *&n) {
    if (n == nullptr) {
      return;
    }
    update_height(n);
    int bf = balance_factor(n);
    if (bf > 1 && balance_factor(n->left) >= 0) {
      rotate_right(n);
    } else if (bf > 1 && balance_factor(n->left) < 0) {
      rotate_left(n->left);
      rotate_right(n);
    } else if (bf < -1 && balance_factor(n->right) <= 0) {
      rotate_left(n);
    } else if (bf < -1 && balance_factor(n->right) > 0) {
      rotate_right(n->right);
      rotate_left(n);
    }
  }

  static bool insert(Node *&n, const K &k, const V &v) {
    if (n == nullptr) {
      n = new Node(k, v);
      return true;
    }
    if ((k < n->key && insert(n->left, k, v)) || (n->key < k && insert(n->right, k, v))) {
      rebalance(n);
      return true;
    }
    return false;
  }

  static bool erase(Node *&n, const K &k) {
    if (n == nullptr) {
      return false;
    }
    if (!(k < n->key || n->key < k)) {
      if (n->left != nullptr && n->right != nullptr) {
        Node *tmp = n->right, *parent = nullptr;
        while (tmp->left != nullptr) {
          parent = tmp;
          tmp = tmp->left;
        }
        n->key = tmp->key;
        n->value = tmp->value;
        if (parent != nullptr) {
          if (!erase(parent->left, parent->left->key)) {
            return false;
          }
        } else if (!erase(n->right, n->right->key)) {
          return false;
        }
      } else {
        Node *tmp = (n->left != nullptr) ? n->left : n->right;
        delete n;
        n = tmp;
      }
      rebalance(n);
      return true;
    }
    if ((k < n->key && erase(n->left, k)) || (n->key < k && erase(n->right, k))) {
      rebalance(n);
      return true;
    }
    return false;
  }

  template<class Fn>
  static void walk(Node *n, Fn f) {
    if (n != nullptr) {
      walk(n->left, f);
      f(n->key, n->value);
      walk(n->right, f);
    }
  }

  static void clean_up(Node *n) {
    if (n != nullptr) {
      clean_up(n->left);
      clean_up(n->right);
      delete n;
    }
  }

 public:
  AVLTree() : root(nullptr), num_nodes(0) {}

  ~AVLTree() { clean_up(root); }
  int size() const { return num_nodes; }
  bool empty() const { return root == nullptr; }

  bool insert(const K &k, const V &v) {
    if (insert(root, k, v)) {
      num_nodes++;
      return true;
    }
    return false;
  }

  bool erase(const K &k) {
    if (erase(root, k)) {
      num_nodes--;
      return true;
    }
    return false;
  }

  const V *find(const K &k) const {
    Node *n = root;
    while (n != nullptr) {
      if (k < n->key) {
        n = n->left;
      } else if (n->key < k) {
        n = n->right;
      } else {
        return &(n->value);
      }
    }
    return nullptr;
  }

  template<class Fn>
  void walk(Fn f) const {
    walk(root, f);
  }
};
\end{lstlisting}
\Needspace*{4\baselineskip}
\textbf{Example Usage}\par
\begin{lstlisting}[language=C++]
#include <cassert>
#include <iostream>
using namespace std;

int main() {
  AVLTree<int, char> t;
  t.insert(2, 'b');
  t.insert(1, 'a');
  t.insert(3, 'c');
  t.insert(5, 'e');
  assert(t.insert(4, 'd'));
  assert(*t.find(4) == 'd');
  assert(!t.insert(4, 'd'));
  auto printch = [](int k, char v) { cout << v; };
  t.walk(printch);
  cout << endl;
  assert(t.erase(1));
  assert(!t.erase(1));
  assert(t.find(1) == nullptr);
  t.walk(printch);
  cout << endl;
  return 0;
}
\end{lstlisting}
\begin{exampleoutput}
abcde
bcde
\end{exampleoutput}
\subsection{Red-Black Tree}
Maintain a map, that is, a collection of key-value pairs such that each possible key appears at most once in the collection. This implementation requires an ordering on the set of possible keys defined by \inlinecode{operator\textless{}} on the key type. A red black tree is a binary search tree balanced by coloring its nodes red or black, then constraining node colors on any simple path from the root to a leaf.

\begin{compactitem}
  \item \inlinecode{RedBlackTree()} constructs an empty map.
  \item \inlinecode{size()} returns the size of the map.
  \item \inlinecode{empty()} returns whether the map is empty.
  \item \inlinecode{insert(k, v)} adds an entry with key \inlinecode{k} and value \inlinecode{v} to the map, returning \inlinecode{true} if an new entry was added or \inlinecode{false} if the key already exists (in which case the map is unchanged and the old value associated with the key is preserved).
  \item \inlinecode{erase(k)} removes the entry with key \inlinecode{k} from the map, returning \inlinecode{true} if the removal was successful or \inlinecode{false} if the key to be removed was not found.
  \item \inlinecode{find(k)} returns a pointer to a const value associated with key \inlinecode{k}, or \inlinecode{nullptr} if the key was not found.
  \item \inlinecode{walk(f)} calls the function \inlinecode{f(k, v)} on each entry of the map, in ascending order of keys.
\end{compactitem}

\Needspace*{3\baselineskip}
\textbf{Time Complexity}:
\begin{compactitem}
  \item $O(1)$ per call to the constructor, \inlinecode{size()}, and \inlinecode{empty()}.
  \item $O(\log n)$ per call to \inlinecode{insert()}, \inlinecode{erase()}, and \inlinecode{find()}, where $n$ is the number of entries currently in the map.
  \item $O(n)$ per call to \inlinecode{walk()}.
\end{compactitem}

\Needspace*{3\baselineskip}
\textbf{Space Complexity}:
\begin{compactitem}
  \item $O(n)$ for storage of the map elements.
  \item $O(\log n)$ auxiliary stack space for \inlinecode{walk()}.
  \item $O(1)$ auxiliary for all other operations.
\end{compactitem}

\begin{lstlisting}[language=C++]
#include <algorithm>
#include <cstddef>

template<class K, class V>
class RedBlackTree {
  enum Color { RED, BLACK };
  struct Node {
    K key;
    V value;
    Color color;
    Node *left, *right, *parent;

    Node(const K &k, const V &v, Color c)
        : key(k), value(v), color(c), left(nullptr), right(nullptr), parent(nullptr) {}
  } *root, *LEAF_NIL;

  int num_nodes;

  void rotate_left(Node *n) {
    Node *tmp = n->right;
    if ((n->right = tmp->left) != LEAF_NIL) {
      n->right->parent = n;
    }
    if ((tmp->parent = n->parent) == LEAF_NIL) {
      root = tmp;
    } else if (n->parent->left == n) {
      n->parent->left = tmp;
    } else {
      n->parent->right = tmp;
    }
    tmp->left = n;
    n->parent = tmp;
  }

  void rotate_right(Node *n) {
    Node *tmp = n->left;
    if ((n->left = tmp->right) != LEAF_NIL) {
      n->left->parent = n;
    }
    if ((tmp->parent = n->parent) == LEAF_NIL) {
      root = tmp;
    } else if (n->parent->right == n) {
      n->parent->right = tmp;
    } else {
      n->parent->left = tmp;
    }
    tmp->right = n;
    n->parent = tmp;
  }

  void insert_fix(Node *n) {
    while (n->parent->color == RED) {
      Node *parent = n->parent;
      Node *grandparent = n->parent->parent;
      if (parent == grandparent->left) {
        Node *uncle = grandparent->right;
        if (uncle->color == RED) {
          grandparent->color = RED;
          parent->color = BLACK;
          uncle->color = BLACK;
          n = grandparent;
        } else {
          if (n == parent->right) {
            rotate_left(parent);
            n = parent;
            parent = n->parent;
          }
          rotate_right(grandparent);
          std::swap(parent->color, grandparent->color);
          n = parent;
        }
      } else if (parent == grandparent->right) {
        Node *uncle = grandparent->left;
        if (uncle->color == RED) {
          grandparent->color = RED;
          parent->color = BLACK;
          uncle->color = BLACK;
          n = grandparent;
        } else {
          if (n == parent->left) {
            rotate_right(parent);
            n = parent;
            parent = n->parent;
          }
          rotate_left(grandparent);
          std::swap(parent->color, grandparent->color);
          n = parent;
        }
      }
    }
    root->color = BLACK;
  }

  void replace(Node *n, Node *replacement) {
    if (n->parent == LEAF_NIL) {
      root = replacement;
    } else if (n == n->parent->left) {
      n->parent->left = replacement;
    } else {
      n->parent->right = replacement;
    }
    replacement->parent = n->parent;
  }

  void erase_fix(Node *n) {
    while (n != root && n->color == BLACK) {
      Node *parent = n->parent;
      if (n == parent->left) {
        Node *sibling = parent->right;
        if (sibling->color == RED) {
          sibling->color = BLACK;
          parent->color = RED;
          rotate_left(parent);
          sibling = parent->right;
        }
        if (sibling->left->color == BLACK && sibling->right->color == BLACK) {
          sibling->color = RED;
          n = parent;
        } else {
          if (sibling->right->color == BLACK) {
            sibling->left->color = BLACK;
            sibling->color = RED;
            rotate_right(sibling);
            sibling = parent->right;
          }
          sibling->color = parent->color;
          parent->color = BLACK;
          sibling->right->color = BLACK;
          rotate_left(parent);
          n = root;
        }
      } else {
        Node *sibling = parent->left;
        if (sibling->color == RED) {
          sibling->color = BLACK;
          parent->color = RED;
          rotate_right(parent);
          sibling = parent->left;
        }
        if (sibling->left->color == BLACK && sibling->right->color == BLACK) {
          sibling->color = RED;
          n = parent;
        } else {
          if (sibling->left->color == BLACK) {
            sibling->right->color = BLACK;
            sibling->color = RED;
            rotate_left(sibling);
            sibling = parent->left;
          }
          sibling->color = parent->color;
          parent->color = BLACK;
          sibling->left->color = BLACK;
          rotate_right(parent);
          n = root;
        }
      }
    }
    n->color = BLACK;
  }

  template<class Fn>
  void walk(Node *n, Fn f) const {
    if (n != LEAF_NIL) {
      walk(n->left, f);
      f(n->key, n->value);
      walk(n->right, f);
    }
  }

  void clean_up(Node *n) {
    if (n != LEAF_NIL) {
      clean_up(n->left);
      clean_up(n->right);
      delete n;
    }
  }

 public:
  RedBlackTree() : num_nodes(0) { root = LEAF_NIL = new Node(K(), V(), BLACK); }

  ~RedBlackTree() {
    clean_up(root);
    delete LEAF_NIL;
  }

  int size() const { return num_nodes; }
  bool empty() const { return num_nodes == 0; }

  bool insert(const K &k, const V &v) {
    Node *curr = root, *prev = LEAF_NIL;
    while (curr != LEAF_NIL) {
      prev = curr;
      if (k < curr->key) {
        curr = curr->left;
      } else if (curr->key < k) {
        curr = curr->right;
      } else {
        return false;
      }
    }
    Node *n = new Node(k, v, RED);
    n->parent = prev;
    if (prev == LEAF_NIL) {
      root = n;
    } else if (k < prev->key) {
      prev->left = n;
    } else {
      prev->right = n;
    }
    n->left = n->right = LEAF_NIL;
    insert_fix(n);
    num_nodes++;
    return true;
  }

  bool erase(const K &k) {
    Node *n = root;
    while (n != LEAF_NIL) {
      if (k < n->key) {
        n = n->left;
      } else if (n->key < k) {
        n = n->right;
      } else {
        break;
      }
    }
    if (n == LEAF_NIL) {
      return false;
    }
    Color color = n->color;
    Node *replacement;
    if (n->left == LEAF_NIL) {
      replacement = n->right;
      replace(n, n->right);
    } else if (n->right == LEAF_NIL) {
      replacement = n->left;
      replace(n, n->left);
    } else {
      Node *tmp = n->right;
      while (tmp->left != LEAF_NIL) {
        tmp = tmp->left;
      }
      color = tmp->color;
      replacement = tmp->right;
      if (tmp->parent == n) {
        replacement->parent = tmp;
      } else {
        replace(tmp, tmp->right);
        tmp->right = n->right;
        tmp->right->parent = tmp;
      }
      replace(n, tmp);
      tmp->left = n->left;
      tmp->left->parent = tmp;
      tmp->color = n->color;
    }
    delete n;
    if (color == BLACK) {
      erase_fix(replacement);
    }
    return true;
  }

  const V *find(const K &k) const {
    Node *n = root;
    while (n != LEAF_NIL) {
      if (k < n->key) {
        n = n->left;
      } else if (n->key < k) {
        n = n->right;
      } else {
        return &(n->value);
      }
    }
    return nullptr;
  }

  template<class Fn>
  void walk(Fn f) const {
    walk(root, f);
  }
};
\end{lstlisting}
\Needspace*{4\baselineskip}
\textbf{Example Usage}\par
\begin{lstlisting}[language=C++]
#include <cassert>
#include <iostream>
using namespace std;

int main() {
  RedBlackTree<int, char> t;
  t.insert(2, 'b');
  t.insert(1, 'a');
  t.insert(3, 'c');
  t.insert(5, 'e');
  assert(t.insert(4, 'd'));
  assert(*t.find(4) == 'd');
  assert(!t.insert(4, 'd'));
  auto printch = [](int k, char v) { cout << v; };
  t.walk(printch);
  cout << endl;
  assert(t.erase(1));
  assert(!t.erase(1));
  assert(t.find(1) == nullptr);
  t.walk(printch);
  cout << endl;
  return 0;
}
\end{lstlisting}
\begin{exampleoutput}
abcde
bcde
\end{exampleoutput}
\subsection{Splay Tree}
Maintain a map, that is, a collection of key-value pairs such that each possible key appears at most once in the collection. This implementation requires an ordering on the set of possible keys defined by \inlinecode{operator\textless{}} on the key type. A splay tree is a balanced binary search tree with the additional property that recently accessed elements are quick to access again.

\begin{compactitem}
  \item \inlinecode{SplayTree()} constructs an empty map.
  \item \inlinecode{size()} returns the size of the map.
  \item \inlinecode{empty()} returns whether the map is empty.
  \item \inlinecode{insert(k, v)} adds an entry with key \inlinecode{k} and value \inlinecode{v} to the map, returning \inlinecode{true} if an new entry was added or \inlinecode{false} if the key already exists (in which case the map is unchanged and the old value associated with the key is preserved).
  \item \inlinecode{erase(k)} removes the entry with key \inlinecode{k} from the map, returning \inlinecode{true} if the removal was successful or \inlinecode{false} if the key to be removed was not found.
  \item \inlinecode{find(k)} returns a pointer to a const value associated with key \inlinecode{k}, or \inlinecode{nullptr} if the key was not found.
  \item \inlinecode{walk(f)} calls the function \inlinecode{f(k, v)} on each entry of the map, in ascending order of keys.
\end{compactitem}

\Needspace*{3\baselineskip}
\textbf{Time Complexity}:
\begin{compactitem}
  \item $O(1)$ per call to the constructor, \inlinecode{size()}, and \inlinecode{empty()}.
  \item $O(\log n)$ per call to \inlinecode{insert()}, \inlinecode{erase()}, and \inlinecode{find()}, where $n$ is the number of entries currently in the map.
  \item $O(n)$ per call to \inlinecode{walk()}.
\end{compactitem}

\Needspace*{3\baselineskip}
\textbf{Space Complexity}:
\begin{compactitem}
  \item $O(n)$ for storage of the map elements.
  \item $O(\log n)$ auxiliary stack space for \inlinecode{insert()}, \inlinecode{erase()}, and \inlinecode{walk()}.
  \item $O(1)$ auxiliary for all other operations.
\end{compactitem}

\begin{lstlisting}[language=C++]
#include <cstddef>

template<class K, class V>
class SplayTree {
  struct Node {
    K key;
    V value;
    Node *left, *right;

    Node(const K &k, const V &v) : key(k), value(v), left(nullptr), right(nullptr) {}
  } *root;

  int num_nodes;

  static void rotate_left(Node *&n) {
    Node *tmp = n;
    n = n->right;
    tmp->right = n->left;
    n->left = tmp;
  }

  static void rotate_right(Node *&n) {
    Node *tmp = n;
    n = n->left;
    tmp->left = n->right;
    n->right = tmp;
  }

  static void splay(Node *&n, const K &k) {
    if (n == nullptr) {
      return;
    }
    if (k < n->key && n->left != nullptr) {
      if (k < n->left->key) {
        splay(n->left->left, k);
        rotate_right(n);
      } else if (n->left->key < k) {
        splay(n->left->right, k);
        if (n->left->right != nullptr) {
          rotate_left(n->left);
        }
      }
      if (n->left != nullptr) {
        rotate_right(n);
      }
    } else if (n->key < k && n->right != nullptr) {
      if (k < n->right->key) {
        splay(n->right->left, k);
        if (n->right->left != nullptr) {
          rotate_right(n->right);
        }
      } else if (n->right->key < k) {
        splay(n->right->right, k);
        rotate_left(n);
      }
      if (n->right != nullptr) {
        rotate_left(n);
      }
    }
  }

  static bool insert(Node *&n, const K &k, const V &v) {
    if (n == nullptr) {
      n = new Node(k, v);
      return true;
    }
    splay(n, k);
    if (k < n->key) {
      Node *tmp = new Node(k, v);
      tmp->left = n->left;
      tmp->right = n;
      n->left = nullptr;
      n = tmp;
    } else if (n->key < k) {
      Node *tmp = new Node(k, v);
      tmp->left = n;
      tmp->right = n->right;
      n->right = nullptr;
      n = tmp;
    } else {
      return false;
    }
    return true;
  }

  static bool erase(Node *&n, const K &k) {
    if (n == nullptr) {
      return false;
    }
    splay(n, k);
    if (k < n->key || n->key < k) {
      return false;
    }
    Node *tmp = n;
    if (n->left == nullptr) {
      n = n->right;
    } else {
      splay(n->left, k);
      n = n->left;
      n->right = tmp->right;
    }
    delete tmp;
    return true;
  }

  template<class Fn>
  static void walk(Node *n, Fn f) {
    if (n != nullptr) {
      walk(n->left, f);
      f(n->key, n->value);
      walk(n->right, f);
    }
  }

  static void clean_up(Node *n) {
    if (n != nullptr) {
      clean_up(n->left);
      clean_up(n->right);
      delete n;
    }
  }

 public:
  SplayTree() : root(nullptr), num_nodes(0) {}

  ~SplayTree() { clean_up(root); }
  int size() const { return num_nodes; }
  bool empty() const { return root == nullptr; }

  bool insert(const K &k, const V &v) {
    if (insert(root, k, v)) {
      num_nodes++;
      return true;
    }
    return false;
  }

  bool erase(const K &k) {
    if (erase(root, k)) {
      num_nodes--;
      return true;
    }
    return false;
  }

  const V *find(const K &k) {
    splay(root, k);
    return (k < root->key || root->key < k) ? nullptr : &(root->value);
  }

  template<class Fn>
  void walk(Fn f) const {
    walk(root, f);
  }
};
\end{lstlisting}
\Needspace*{4\baselineskip}
\textbf{Example Usage}\par
\begin{lstlisting}[language=C++]
#include <cassert>
#include <iostream>
using namespace std;

int main() {
  SplayTree<int, char> t;
  t.insert(2, 'b');
  t.insert(1, 'a');
  t.insert(3, 'c');
  t.insert(5, 'e');
  assert(t.insert(4, 'd'));
  assert(*t.find(4) == 'd');
  assert(!t.insert(4, 'd'));
  auto printch = [](int k, char v) { cout << v; };
  t.walk(printch);
  cout << endl;
  assert(t.erase(1));
  assert(!t.erase(1));
  assert(t.find(1) == nullptr);
  t.walk(printch);
  cout << endl;
  return 0;
}
\end{lstlisting}
\begin{exampleoutput}
abcde
bcde
\end{exampleoutput}
\subsection{Size Balanced Tree}
Maintain a map, that is, a collection of key-value pairs such that each possible key appears at most once in the collection. In addition, support queries for keys given their ranks as well as queries for the ranks of given keys. This implementation requires an ordering on the set of possible keys defined by \inlinecode{operator\textless{}} on the key type. A size balanced tree augments each nodes with the size of its subtree, using it to maintain balance and compute order statistics.

\begin{compactitem}
  \item \inlinecode{SBTree()} constructs an empty map.
  \item \inlinecode{size()} returns the size of the map.
  \item \inlinecode{empty()} returns whether the map is empty.
  \item \inlinecode{insert(k, v)} adds an entry with key \inlinecode{k} and value \inlinecode{v} to the map, returning \inlinecode{true} if an new entry was added or \inlinecode{false} if the key already exists (in which case the map is unchanged and the old value associated with the key is preserved).
  \item \inlinecode{erase(k)} removes the entry with key \inlinecode{k} from the map, returning \inlinecode{true} if the removal was successful or \inlinecode{false} if the key to be removed was not found.
  \item \inlinecode{find(k)} returns a pointer to a const value associated with key \inlinecode{k}, or \inlinecode{nullptr} if the key was not found.
  \item \inlinecode{select(r)} returns a key-value pair of the node with a key of zero-based rank \inlinecode{r} in the map, throwing an exception if the rank is not between 0 and \inlinecode{size() - 1}.
  \item \inlinecode{rank(k)} returns the zero-based rank of key \inlinecode{k} in the map, throwing an exception if the key was not found in the map.
  \item \inlinecode{walk(f)} calls the function \inlinecode{f(k, v)} on each entry of the map, in ascending order of keys.
\end{compactitem}

\Needspace*{3\baselineskip}
\textbf{Time Complexity}:
\begin{compactitem}
  \item $O(1)$ per call to the constructor, \inlinecode{size()}, and \inlinecode{empty()}.
  \item $O(\log n)$ per call to \inlinecode{insert()}, \inlinecode{erase()}, \inlinecode{find()}, \inlinecode{select()}, and \inlinecode{rank()}, where $n$ is the number of entries currently in the map.
  \item $O(n)$ per call to \inlinecode{walk()}.
\end{compactitem}

\Needspace*{3\baselineskip}
\textbf{Space Complexity}:
\begin{compactitem}
  \item $O(n)$ for storage of the map elements.
  \item $O(\log n)$ auxiliary stack space for \inlinecode{insert()}, \inlinecode{erase()}, and \inlinecode{walk()}.
  \item $O(1)$ auxiliary for all other operations.
\end{compactitem}

\begin{lstlisting}[language=C++]
#include <cstddef>
#include <stdexcept>
#include <utility>

template<class K, class V>
class SBTree {
  struct Node {
    K key;
    V value;
    int size;
    Node *left, *right;

    Node(const K &k, const V &v) : key(k), value(v), size(1), left(nullptr), right(nullptr) {}

    inline Node *&child(int c) { return (c == 0) ? left : right; }

    void update() {
      size = 1;
      if (left != nullptr) {
        size += left->size;
      }
      if (right != nullptr) {
        size += right->size;
      }
    }
  } *root;

  static inline int size(Node *n) { return (n == nullptr) ? 0 : n->size; }

  static void rotate(Node *&n, int c) {
    Node *tmp = n->child(c);
    n->child(c) = tmp->child(!c);
    tmp->child(!c) = n;
    n->update();
    tmp->update();
    n = tmp;
  }

  static void maintain(Node *&n, int c) {
    if (n == nullptr || n->child(c) == nullptr) {
      return;
    }
    Node *&tmp = n->child(c);
    if (size(tmp->child(c)) > size(n->child(!c))) {
      rotate(n, c);
    } else if (size(tmp->child(!c)) > size(n->child(!c))) {
      rotate(tmp, !c);
      rotate(n, c);
    } else {
      return;
    }
    maintain(n->left, 0);
    maintain(n->right, 1);
    maintain(n, 0);
    maintain(n, 1);
  }

  static bool insert(Node *&n, const K &k, const V &v) {
    if (n == nullptr) {
      n = new Node(k, v);
      return true;
    }
    bool found;
    if (k < n->key) {
      found = insert(n->left, k, v);
      maintain(n, 0);
    } else if (n->key < k) {
      found = insert(n->right, k, v);
      maintain(n, 1);
    } else {
      return false;
    }
    n->update();
    return found;
  }

  static bool erase(Node *&n, const K &k) {
    if (n == nullptr) {
      return false;
    }
    bool found;
    int c = (k < n->key);
    if (k < n->key) {
      found = erase(n->left, k);
    } else if (n->key < k) {
      found = erase(n->right, k);
    } else {
      if (n->right == nullptr || n->left == nullptr) {
        Node *tmp = n;
        n = (n->right == nullptr) ? n->left : n->right;
        delete tmp;
        return true;
      }
      Node *p = n->right;
      while (p->left != nullptr) {
        p = p->left;
      }
      n->key = p->key;
      found = erase(n->right, p->key);
    }
    maintain(n, c);
    n->update();
    return found;
  }

  static std::pair<K, V> select(Node *n, int r) {
    int rank = size(n->left);
    if (r < rank) {
      return select(n->left, r);
    } else if (r > rank) {
      return select(n->right, r - rank - 1);
    }
    return {n->key, n->value};
  }

  static int rank(Node *n, const K &k) {
    if (n == nullptr) {
      throw std::runtime_error("Cannot rank key that's not in tree.");
    }
    int r = size(n->left);
    if (k < n->key) {
      return rank(n->left, k);
    } else if (n->key < k) {
      return rank(n->right, k) + r + 1;
    }
    return r;
  }

  template<class Fn>
  static void walk(Node *n, Fn f) {
    if (n != nullptr) {
      walk(n->left, f);
      f(n->key, n->value);
      walk(n->right, f);
    }
  }

  static void clean_up(Node *n) {
    if (n != nullptr) {
      clean_up(n->left);
      clean_up(n->right);
      delete n;
    }
  }

 public:
  SBTree() : root(nullptr) {}

  ~SBTree() { clean_up(root); }
  int size() const { return size(root); }
  bool empty() const { return root == nullptr; }
  bool insert(const K &k, const V &v) { return insert(root, k, v); }
  bool erase(const K &k) { return erase(root, k); }
  int rank(const K &k) const { return rank(root, k); }

  const V *find(const K &k) const {
    Node *n = root;
    while (n != nullptr) {
      if (k < n->key) {
        n = n->left;
      } else if (n->key < k) {
        n = n->right;
      } else {
        return &(n->value);
      }
    }
    return nullptr;
  }

  std::pair<K, V> select(int r) const {
    if (r < 0 || r >= size(root)) {
      throw std::runtime_error("Select rank must be between 0 and size() - 1.");
    }
    return select(root, r);
  }

  template<class Fn>
  void walk(Fn f) const {
    walk(root, f);
  }
};
\end{lstlisting}
\Needspace*{4\baselineskip}
\textbf{Example Usage}\par
\begin{lstlisting}[language=C++]
#include <cassert>
#include <iostream>
using namespace std;

int main() {
  SBTree<int, char> t;
  t.insert(2, 'b');
  t.insert(1, 'a');
  t.insert(3, 'c');
  t.insert(5, 'e');
  assert(t.insert(4, 'd'));
  assert(*t.find(4) == 'd');
  assert(!t.insert(4, 'd'));
  auto printch = [](int k, char v) { cout << v; };
  t.walk(printch);
  cout << endl;
  assert(t.erase(1));
  assert(!t.erase(1));
  assert(t.find(1) == nullptr);
  t.walk(printch);
  cout << endl;
  assert(t.rank(2) == 0);
  assert(t.rank(3) == 1);
  assert(t.rank(5) == 3);
  assert(t.select(0).first == 2);
  assert(t.select(1).first == 3);
  assert(t.select(2).first == 4);
  return 0;
}
\end{lstlisting}
\begin{exampleoutput}
abcde
bcde
\end{exampleoutput}
\subsection{Interval Treap}
Maintain a map from closed, one-dimensional intervals to values while supporting efficient reporting of any or all entries that intersect with a given query interval. This implementation uses \inlinecode{std::pair} to represent intervals, requiring operators \inlinecode{\textless{}} and \inlinecode{==} to be defined on the numeric key type. A treap is used to process the entries, where keys are compared lexicographically as pairs.

\begin{compactitem}
  \item \inlinecode{IntervalTreap()} constructs an empty map.
  \item \inlinecode{size()} returns the size of the map.
  \item \inlinecode{empty()} returns whether the map is empty.
  \item \inlinecode{insert(lo, hi, v)} adds an entry with key \inlinecode{[lo, hi]} and value \inlinecode{v} to the map, returning \inlinecode{true} if a new interval was added or \inlinecode{false} if the interval already exists (in which case the map is unchanged and the old value associated with the key is preserved).
  \item \inlinecode{erase(lo, hi)} removes the entry with key \inlinecode{[lo, hi]} from the map, returning \inlinecode{true} if the removal was successful or \inlinecode{false} if the interval was not found.
  \item \inlinecode{find\_key(lo, hi)} returns a pointer to a const \inlinecode{std::pair} representing the key of some interval in the map which intersects with \inlinecode{[lo, hi]}, or \inlinecode{nullptr} if no such entry was found.
  \item \inlinecode{find\_value(lo, hi)} returns a pointer to a const value of some entry in the map with a key that intersects with \inlinecode{[lo, hi]}, or \inlinecode{nullptr} if no such entry was found.
  \item \inlinecode{find\_all(lo, hi, f)} calls the function \inlinecode{f(lo, hi, v)} on each entry in the map that overlaps with \inlinecode{[lo, hi]}, in lexicographically ascending order of intervals.
  \item \inlinecode{walk(f)} calls the function \inlinecode{f(lo, hi, v)} on each interval in the map, in lexicographically ascending order of intervals.
\end{compactitem}

\Needspace*{3\baselineskip}
\textbf{Time Complexity}:
\begin{compactitem}
  \item $O(1)$ per call to the constructor, \inlinecode{size()}, and \inlinecode{empty()}.
  \item $O(\log n)$ on average per call to \inlinecode{insert()}, \inlinecode{erase()}, and \inlinecode{find\_any()}, where $n$ is the number of intervals currently in the set.
  \item $O(\log n + m)$ on average per call to \inlinecode{find\_all()}, where $m$ is the number of intersecting intervals that are reported.
  \item $O(n)$ per call to \inlinecode{walk()}.
\end{compactitem}

\Needspace*{3\baselineskip}
\textbf{Space Complexity}:
\begin{compactitem}
  \item $O(n)$ for storage of the map elements.
  \item $O(1)$ auxiliary for \inlinecode{size()} and \inlinecode{empty()}.
  \item $O(\log n)$ auxiliary stack space on average for all other operations.
\end{compactitem}

\begin{lstlisting}[language=C++]
#include <cstdlib>
#include <utility>

template<class K, class V>
class IntervalTreap {
  using Interval = std::pair<K, K>;

  struct Node {
    static inline int rand32() { return (rand() & 0x7fff) | ((rand() & 0x7fff) << 15); }

    Interval interval;
    V value;
    K max;
    int priority;
    Node *left, *right;

    Node(const Interval &i, const V &v)
        : interval(i), value(v), max(i.second), priority(rand32()), left(nullptr), right(nullptr) {}

    void update() {
      max = interval.second;
      if (left != nullptr && left->max > max) {
        max = left->max;
      }
      if (right != nullptr && right->max > max) {
        max = right->max;
      }
    }
  } *root;

  int num_nodes;

  static void rotate_left(Node *&n) {
    Node *tmp = n;
    n = n->right;
    tmp->right = n->left;
    n->left = tmp;
    tmp->update();
  }

  static void rotate_right(Node *&n) {
    Node *tmp = n;
    n = n->left;
    tmp->left = n->right;
    n->right = tmp;
    tmp->update();
  }

  static bool insert(Node *&n, const Interval &i, const V &v) {
    if (n == nullptr) {
      n = new Node(i, v);
      return true;
    }
    if (i < n->interval && insert(n->left, i, v)) {
      if (n->left->priority < n->priority) {
        rotate_right(n);
      }
      n->update();
      return true;
    }
    if (i > n->interval && insert(n->right, i, v)) {
      if (n->right->priority < n->priority) {
        rotate_left(n);
      }
      n->update();
      return true;
    }
    return false;
  }

  static bool erase(Node *&n, const Interval &i) {
    if (n == nullptr) {
      return false;
    }
    if (i < n->interval) {
      return erase(n->left, i);
    }
    if (i > n->interval) {
      return erase(n->right, i);
    }
    if (n->left != nullptr && n->right != nullptr) {
      bool res;
      if (n->left->priority < n->right->priority) {
        rotate_right(n);
        res = erase(n->right, i);
      } else {
        rotate_left(n);
        res = erase(n->left, i);
      }
      n->update();
      return res;
    }
    Node *tmp = (n->left != nullptr) ? n->left : n->right;
    delete n;
    n = tmp;
    return true;
  }

  static Node *find_any(Node *n, const Interval &i) {
    if (n == nullptr) {
      return nullptr;
    }
    if (n->interval.first <= i.second && i.first <= n->interval.second) {
      return n;
    }
    if (n->left != nullptr && i.first <= n->left->max) {
      return find_any(n->left, i);
    }
    return find_any(n->right, i);
  }

  template<class Fn>
  static void find_all(Node *n, const Interval &i, Fn f) {
    if (n == nullptr || n->max < i.first) {
      return;
    }
    if (n->interval.first <= i.second && i.first <= n->interval.second) {
      f(n->interval.first, n->interval.second, n->value);
    }
    find_all(n->left, i, f);
    find_all(n->right, i, f);
  }

  template<class Fn>
  static void walk(Node *n, Fn f) {
    if (n != nullptr) {
      walk(n->left, f);
      f(n->interval.first, n->interval.second, n->value);
      walk(n->right, f);
    }
  }

  static void clean_up(Node *n) {
    if (n != nullptr) {
      clean_up(n->left);
      clean_up(n->right);
      delete n;
    }
  }

 public:
  IntervalTreap() : root(nullptr), num_nodes(0) {}

  ~IntervalTreap() { clean_up(root); }
  int size() const { return num_nodes; }
  bool empty() const { return root == nullptr; }

  bool insert(const K &lo, const K &hi, const V &v) {
    if (insert(root, Interval{lo, hi}, v)) {
      num_nodes++;
      return true;
    }
    return false;
  }

  bool erase(const K &lo, const K &hi) {
    if (erase(root, Interval{lo, hi})) {
      num_nodes--;
      return true;
    }
    return false;
  }

  const Interval *find_key(const K &lo, const K &hi) const {
    Node *n = find_any(root, Interval{lo, hi});
    return (n == nullptr) ? nullptr : &(n->interval);
  }

  const V *find_value(const K &lo, const K &hi) const {
    Node *n = find_any(root, Interval{lo, hi});
    return (n == nullptr) ? nullptr : &(n->value);
  }

  template<class Fn>
  void find_all(const K &lo, const K &hi, Fn f) const {
    find_all(root, Interval{lo, hi}, f);
  }

  template<class Fn>
  void walk(Fn f) const {
    walk(root, f);
  }
};
\end{lstlisting}
\Needspace*{4\baselineskip}
\textbf{Example Usage}\par
\begin{lstlisting}[language=C++]
#include <cassert>
#include <iostream>
using namespace std;

int main() {
  IntervalTreap<int, char> t;
  t.insert(15, 20, 'a');
  t.insert(10, 30, 'b');
  t.insert(17, 19, 'c');
  t.insert(5, 20, 'd');
  t.insert(12, 15, 'e');
  t.insert(10, 40, 'f');
  assert(t.size() == 6);
  assert(!t.insert(5, 20, 'x'));
  t.erase(17, 19);
  assert(t.size() == 5);
  assert(*t.find_key(3, 9) == (pair<int, int>{5, 20}));
  assert(*t.find_value(3, 9) == 'd');
  cout << "Intervals intersecting [16, 20]:";
  auto print = [](int lo, int hi, char v) { cout << " [" << lo << ", " << hi << "]"; };
  t.find_all(16, 20, print);
  cout << "\nAll intervals:";
  t.walk(print);
  cout << endl;
  return 0;
}
\end{lstlisting}
\begin{exampleoutput}
Intervals intersecting [16, 20]: [15, 20] [10, 30] [5, 20] [10, 40]
All intervals: [5, 20] [10, 30] [10, 40] [12, 15] [15, 20]
\end{exampleoutput}
\subsection{Hash Map}
Maintain a map, that is, a collection of key-value pairs such that each possible key appears at most once in the collection. This implementation requires \inlinecode{operator==} to be defined on the key type. A hash map implements a map by hashing keys into buckets using a hash function. This implementation resolves collisions by chaining entries hashed to the same bucket into a linked list.

\begin{compactitem}
  \item \inlinecode{HashMap()} constructs an empty map.
  \item \inlinecode{size()} returns the size of the map.
  \item \inlinecode{empty()} returns whether the map is empty.
  \item \inlinecode{insert(k, v)} adds an entry with key \inlinecode{k} and value \inlinecode{v} to the map, returning \inlinecode{true} if an new entry was added or \inlinecode{false} if the key already exists (in which case the map is unchanged and the old value associated with the key is preserved).
  \item \inlinecode{erase(k)} removes the entry with key \inlinecode{k} from the map, returning \inlinecode{true} if the removal was successful or \inlinecode{false} if the key to be removed was not found.
  \item \inlinecode{find(k)} returns a pointer to the value associated with key \inlinecode{k}, or \inlinecode{nullptr} if the key was not found.
  \item \inlinecode{operator[]} returns a reference to key \inlinecode{k}'s associated value (which may be modified), or if necessary, inserts and returns a new entry with the default constructed value if key \inlinecode{k} was not originally found.
  \item \inlinecode{walk(f)} calls the function \inlinecode{f(k, v)} on each entry of the map, in no guaranteed order.
\end{compactitem}

\Needspace*{3\baselineskip}
\textbf{Time Complexity}:
\begin{compactitem}
  \item $O(1)$ per call to the constructor, \inlinecode{size()}, and \inlinecode{empty()}.
  \item $O(1)$ amortized per call to \inlinecode{insert()}, \inlinecode{erase()}, \inlinecode{find()}, and \inlinecode{operator[]}.
  \item $O(n)$ per call to \inlinecode{walk()}, where $n$ is the number of entries in the map.
\end{compactitem}

\Needspace*{3\baselineskip}
\textbf{Space Complexity}:
\begin{compactitem}
  \item $O(n)$ for storage of the map elements.
  \item $O(n)$ auxiliary heap space for \inlinecode{insert()}.
  \item $O(1)$ auxiliary for all other operations.
\end{compactitem}

\begin{lstlisting}[language=C++]
#include <cstddef>
#include <list>
#include <vector>

template<class K, class V, class Hash>
class HashMap {
  struct HashMapEntry {
    K key;
    V value;

    HashMapEntry(const K &k, const V &v) : key(k), value(v) {}
  };

  std::vector<std::list<HashMapEntry>> table;
  int table_size, num_entries;

  void double_capacity_and_rehash() {
    std::vector<std::list<HashMapEntry>> old = std::move(table);
    table_size = 2 * table_size;
    table.assign(table_size, std::list<HashMapEntry>());
    num_entries = 0;
    for (auto &bucket : old) {
      for (auto &entry : bucket) {
        insert(entry.key, entry.value);
      }
    }
  }

 public:
  explicit HashMap(int size = 128) : table(size), table_size(size), num_entries(0) {}

  int size() const { return num_entries; }
  bool empty() const { return num_entries == 0; }

  bool insert(const K &k, const V &v) {
    if (find(k) != nullptr) {
      return false;
    }
    if (num_entries >= table_size) {
      double_capacity_and_rehash();
    }
    unsigned int i = Hash()(k) % table_size;
    table[i].emplace_back(k, v);
    num_entries++;
    return true;
  }

  bool erase(const K &k) {
    unsigned int i = Hash()(k) % table_size;
    auto it = table[i].begin();
    while (it != table[i].end() && !(it->key == k)) {
      ++it;
    }
    if (it == table[i].end()) {
      return false;
    }
    table[i].erase(it);
    num_entries--;
    return true;
  }

  V *find(const K &k) {
    unsigned int i = Hash()(k) % table_size;
    auto it = table[i].begin();
    while (it != table[i].end() && !(it->key == k)) {
      ++it;
    }
    if (it == table[i].end()) {
      return nullptr;
    }
    return &(it->value);
  }

  V &operator[](const K &k) {
    if (V *ptr = find(k); ptr != nullptr) {
      return *ptr;
    }
    if (num_entries >= table_size) {
      double_capacity_and_rehash();
    }
    unsigned int i = Hash()(k) % table_size;
    table[i].emplace_back(k, V());
    num_entries++;
    return table[i].back().value;
  }

  template<class Fn>
  void walk(Fn f) const {
    for (int i = 0; i < table_size; i++) {
      for (const auto &entry : table[i]) {
        f(entry.key, entry.value);
      }
    }
  }
};
\end{lstlisting}
\Needspace*{4\baselineskip}
\textbf{Example Usage}\par
\begin{lstlisting}[language=C++]
#include <cassert>
#include <iostream>
using namespace std;

struct ClassHash {
  unsigned int operator()(int k) { return ClassHash()(static_cast<unsigned int>(k)); }
  unsigned int operator()(long long k) { return ClassHash()(static_cast<unsigned long long>(k)); }

  // Knuth's one-to-one multiplicative method.
  unsigned int operator()(unsigned int k) {
    return k * 2654435761u;  // Or just return k.
  }

  // Jenkins's 64-bit hash.
  unsigned int operator()(unsigned long long k) {
    k += ~(k << 32);
    k ^= (k >> 22);
    k += ~(k << 13);
    k ^= (k >> 8);
    k += (k << 3);
    k ^= (k >> 15);
    k += ~(k << 27);
    k ^= (k >> 31);
    return k;
  }

  // Jenkins's one-at-a-time hash.
  unsigned int operator()(const std::string &k) {
    unsigned int hash = 0;
    for (char c : k) {
      hash += ((hash + static_cast<unsigned char>(c)) << 10);
      hash ^= (hash >> 6);
    }
    hash += (hash << 3);
    hash ^= (hash >> 11);
    return hash + (hash << 15);
  }
};

int main() {
  HashMap<string, char, ClassHash> m;
  m["foo"] = 'a';
  m.insert("bar", 'b');
  assert(m["foo"] == 'a');
  assert(m["bar"] == 'b');
  assert(m["baz"] == '\0');
  m["baz"] = 'c';
  m.walk([](const string &k, char v) { cout << v; });
  cout << endl;
  assert(m.erase("foo"));
  assert(m.size() == 2);
  assert(m["foo"] == '\0');
  assert(m.size() == 3);
  return 0;
}
\end{lstlisting}
\begin{exampleoutput}
cab
\end{exampleoutput}
\subsection{Skip List}
Maintain a map, that is, a collection of key-value pairs such that each possible key appears at most once in the collection. This implementation requires operators \inlinecode{\textless{}} and \inlinecode{==} to be defined on the key type. A skip list maintains a linked hierarchy of sorted subsequences with each successive subsequence skipping over fewer elements than the previous one.

\begin{compactitem}
  \item \inlinecode{SkipList()} constructs an empty map.
  \item \inlinecode{size()} returns the size of the map.
  \item \inlinecode{empty()} returns whether the map is empty.
  \item \inlinecode{insert(k, v)} adds an entry with key \inlinecode{k} and value \inlinecode{v} to the map, returning \inlinecode{true} if an new entry was added or \inlinecode{false} if the key already exists (in which case the map is unchanged and the old value associated with the key is preserved).
  \item \inlinecode{erase(k)} removes the entry with key \inlinecode{k} from the map, returning \inlinecode{true} if the removal was successful or \inlinecode{false} if the key to be removed was not found.
  \item \inlinecode{find(k)} returns a pointer to the value associated with key \inlinecode{k}, or \inlinecode{nullptr} if the key was not found.
  \item \inlinecode{operator[]} returns a reference to key \inlinecode{k}'s associated value (which may be modified), or if necessary, inserts and returns a new entry with the default constructed value if key \inlinecode{k} was not originally found.
  \item \inlinecode{walk(f)} calls the function \inlinecode{f(k, v)} on each entry of the map, in ascending order of keys.
\end{compactitem}

\Needspace*{3\baselineskip}
\textbf{Time Complexity}:
\begin{compactitem}
  \item $O(1)$ per call to the constructor, \inlinecode{size()}, and \inlinecode{empty()}.
  \item $O(\log n)$ on average per call to \inlinecode{insert()}, \inlinecode{erase()}, \inlinecode{find()}, and \inlinecode{operator[]}, where $n$ is the number of entries currently in the map.
  \item $O(n)$ per call to \inlinecode{walk()}.
\end{compactitem}

\Needspace*{3\baselineskip}
\textbf{Space Complexity}:
\begin{compactitem}
  \item $O(n)$ on average for storage of the map elements.
  \item $O(n)$ auxiliary heap space for \inlinecode{insert()} and \inlinecode{erase()}.
  \item $O(1)$ auxiliary for all other operations.
\end{compactitem}

\begin{lstlisting}[language=C++]
#include <random>
#include <vector>

template<class K, class V>
class SkipList {
  static const int MAX_LEVELS = 32;  // log2(max possible keys)

  struct Node {
    K key;
    V value;
    std::vector<Node *> next;

    Node(const K &k, const V &v, int levels) : key(k), value(v), next(levels, (Node *)nullptr) {}
  } *head;

  int num_nodes;

  static int random_level() {
    static std::mt19937 rng(std::random_device{}());
    static std::uniform_int_distribution<int> coin(0, 1);
    int level = 1;
    while (coin(rng) && level < MAX_LEVELS) {
      level++;
    }
    return level;
  }

  static int node_level(const std::vector<Node *> &v) {
    int i = 0;
    while (i < static_cast<int>(v.size()) && v[i] != nullptr) {
      i++;
    }
    return i + 1;
  }

 public:
  SkipList() : head(new Node(K(), V(), MAX_LEVELS)), num_nodes(0) {
    for (auto &ptr : head->next) {
      ptr = nullptr;
    }
  }

  ~SkipList() { delete head; }
  int size() const { return num_nodes; }
  bool empty() const { return num_nodes == 0; }

  bool insert(const K &k, const V &v) {
    std::vector<Node *> update(head->next);
    int curr_level = node_level(update);
    Node *n = head;
    for (int i = curr_level; i-- > 0;) {
      while (n->next[i] != nullptr && n->next[i]->key < k) {
        n = n->next[i];
      }
      update[i] = n;
    }
    n = n->next[0];
    if (n != nullptr && n->key == k) {
      return false;
    }
    int new_level = random_level();
    if (new_level > curr_level) {
      for (int i = curr_level; i < new_level; i++) {
        update[i] = head;
      }
    }
    n = new Node(k, v, new_level);
    for (int i = 0; i < new_level; i++) {
      n->next[i] = update[i]->next[i];
      update[i]->next[i] = n;
    }
    num_nodes++;
    return true;
  }

  bool erase(const K &k) {
    std::vector<Node *> update(head->next);
    Node *n = head;
    for (int i = node_level(update); i-- > 0;) {
      while (n->next[i] != nullptr && n->next[i]->key < k) {
        n = n->next[i];
      }
      update[i] = n;
    }
    n = n->next[0];
    if (n != nullptr && n->key == k) {
      for (int i = 0, levels = static_cast<int>(update.size()); i < levels; i++) {
        if (update[i]->next[i] != n) {
          break;
        }
        update[i]->next[i] = n->next[i];
      }
      delete n;
      num_nodes--;
      return true;
    }
    return false;
  }

  V *find(const K &k) const {
    Node *n = head;
    for (int i = node_level(n->next); i-- > 0;) {
      while (n->next[i] != nullptr && n->next[i]->key < k) {
        n = n->next[i];
      }
    }
    n = n->next[0];
    return (n != nullptr && n->key == k) ? &(n->value) : nullptr;
  }

  V &operator[](const K &k) {
    V *ptr = find(k);
    if (ptr != nullptr) {
      return *ptr;
    }
    insert(k, V());
    return *find(k);
  }

  template<class Fn>
  void walk(Fn f) const {
    Node *n = head->next[0];
    while (n != nullptr) {
      f(n->key, n->value);
      n = n->next[0];
    }
  }
};
\end{lstlisting}
\Needspace*{4\baselineskip}
\textbf{Example Usage}\par
\begin{lstlisting}[language=C++]
#include <cassert>
#include <iostream>
using namespace std;

int main() {
  SkipList<int, char> l;
  l.insert(2, 'b');
  l.insert(1, 'a');
  l.insert(3, 'c');
  l.insert(5, 'e');
  assert(l.insert(4, 'd'));
  assert(*l.find(4) == 'd');
  assert(!l.insert(4, 'd'));
  auto printch = [](int k, char v) { cout << v; };
  l.walk(printch);
  cout << endl;
  assert(l.erase(1));
  assert(!l.erase(1));
  assert(l.find(1) == nullptr);
  l.walk(printch);
  cout << endl;
  return 0;
}
\end{lstlisting}
\begin{exampleoutput}
abcde
bcde
\end{exampleoutput}

\section{Range Queries in One Dimension}
\setcounter{section}{4}
\setcounter{subsection}{0}
\subsection{Sparse Table (Range Minimum Query)}
Given a static array with indices from $0$ to $n - 1$, precompute a table that may later be used to perform range minimum queries on the array in constant time. This version is simplified to only work on integer arrays.

The dynamic programming state $dp[i][j]$ holds the index of the minimum value in the sub-array starting at $i$ and having length $2^j$. Each $dp[i][j]$ will always be equal to either $dp[i][j - 1]$ or $dp[i + 2^{j - 1}][j - 1]$, whichever of the indices corresponds to the smaller value in the array.

\Needspace*{3\baselineskip}
\textbf{Time Complexity}:
\begin{compactitem}
  \item $O(n \log n)$ per call to \inlinecode{build()}, where $n$ is the size of the array.
  \item $O(1)$ per call to \inlinecode{query()}.
\end{compactitem}

\Needspace*{3\baselineskip}
\textbf{Space Complexity}:
\begin{compactitem}
  \item $O(n \log n)$ for storage of the sparse table, where $n$ is the size of the array.
  \item $O(1)$ auxiliary for \inlinecode{query()}.
\end{compactitem}

\begin{lstlisting}[language=C++]
#include <vector>

std::vector<int> table;
std::vector<std::vector<int>> dp;

void build(const std::vector<int> &a) {
  int n = a.size();
  table.resize(n + 1);
  dp.assign(n, std::vector<int>());
  for (int i = 2; i <= n; i++) {
    table[i] = table[i >> 1] + 1;
  }
  for (int i = 0; i < n; i++) {
    dp[i].resize(table[n] + 1);
    dp[i][0] = i;
  }
  for (int j = 1; (1 << j) < n; j++) {
    for (int i = 0; i + (1 << j) <= n; i++) {
      int x = dp[i][j - 1];
      int y = dp[i + (1 << (j - 1))][j - 1];
      dp[i][j] = (a[x] < a[y]) ? x : y;
    }
  }
}

int query(const std::vector<int> &a, int lo, int hi) {
  int j = table[hi - lo];
  int x = dp[lo][j];
  int y = dp[hi - (1 << j) + 1][j];
  return (a[x] < a[y]) ? a[x] : a[y];
}
\end{lstlisting}
\Needspace*{4\baselineskip}
\textbf{Example Usage}\par
\begin{lstlisting}[language=C++]
#include <cassert>

int main() {
  std::vector<int> a{6, -2, 1, 8, 10};
  build(a);
  assert(query(a, 0, 3) == -2);
  return 0;
}
\end{lstlisting}
\subsection{Square Root Decomposition}
Maintain a fixed-size array (from 0 to \inlinecode{size() - 1}) while supporting dynamic queries of contiguous sub-arrays and dynamic updates of individual indices.

The query operation is defined by an associative aggregate function \inlinecode{combine(a, b)}. The default definition below assumes a numerical array type, supporting queries for the ``min'' of the target range. Another possible query operation is ``sum'', in which case \inlinecode{combine(a, b)} should return $a + b$.

The point update operation is defined by \inlinecode{apply\_delta(v, d)}, which returns the new value at a single updated index. The default definition below supports updates that ``set'' the chosen array index to a new value. Another possible update operation is ``increment'', in which case \inlinecode{apply\_delta(v, d)} should return $v + d$.

The operations supported by this data structure are identical to those of the point update segment tree found in this section.

\Needspace*{3\baselineskip}
\textbf{Time Complexity}:
\begin{compactitem}
  \item $O(n)$ per call to both constructors, where $n$ is the size of the array.
  \item $O(1)$ per call to \inlinecode{size()}.
  \item $O(\sqrt{n})$ per call to \inlinecode{at()}, \inlinecode{update()}, and \inlinecode{query()}.
\end{compactitem}

\Needspace*{3\baselineskip}
\textbf{Space Complexity}:
\begin{compactitem}
  \item $O(n)$ for storage of the array elements.
  \item $O(1)$ auxiliary for all operations.
\end{compactitem}

\begin{lstlisting}[language=C++]
#include <algorithm>
#include <cmath>
#include <vector>

template<class T>
class SqrtDecomposition {
  static T combine(const T &a, const T &b) { return std::min(a, b); }
  static T apply_delta(const T &v, const T &d) { return d; }

  int len, blocklen;
  std::vector<T> value, block;

  void init() {
    blocklen = std::max(1, static_cast<int>(sqrt(len)));
    int nblocks = (len + blocklen - 1) / blocklen;
    for (int i = 0; i < nblocks; i++) {
      T blockval = value[i * blocklen];
      int blockhi = std::min(len, (i + 1) * blocklen);
      for (int j = i * blocklen + 1; j < blockhi; j++) {
        blockval = combine(blockval, value[j]);
      }
      block.push_back(blockval);
    }
  }

 public:
  explicit SqrtDecomposition(int n, const T &v = T()) : len(n), value(n, v) { init(); }

  template<class It>
  SqrtDecomposition(It lo, It hi) : len(hi - lo), value(lo, hi) {
    init();
  }

  int size() const { return len; }
  T at(int i) const { return query(i, i); }

  T query(int lo, int hi) const {
    T res;
    int blocklo = ceil(static_cast<double>(lo) / blocklen), blockhi = (hi + 1) / blocklen - 1;
    if (blocklo > blockhi) {
      res = value[lo];
      for (int i = lo + 1; i <= hi; i++) {
        res = combine(res, value[i]);
      }
    } else {
      res = block[blocklo];
      for (int i = blocklo + 1; i <= blockhi; i++) {
        res = combine(res, block[i]);
      }
      for (int i = lo; i < blocklo * blocklen; i++) {
        res = combine(res, value[i]);
      }
      for (int i = (blockhi + 1) * blocklen; i <= hi; i++) {
        res = combine(res, value[i]);
      }
    }
    return res;
  }

  void update(int i, const T &d) {
    value[i] = apply_delta(value[i], d);
    int b = i / blocklen;
    int blockhi = std::min(len, (b + 1) * blocklen);
    block[b] = value[b * blocklen];
    for (int i = b * blocklen + 1; i < blockhi; i++) {
      block[b] = combine(block[b], value[i]);
    }
  }
};
\end{lstlisting}
\Needspace*{4\baselineskip}
\textbf{Example Usage}\par
\begin{lstlisting}[language=C++]
#include <cassert>
#include <iostream>
using namespace std;

int main() {
  vector<int> arr{6, -2, 1, 8, 10};
  SqrtDecomposition<int> sd(arr.begin(), arr.end());
  sd.update(2, 4);
  cout << "Values:";
  for (int i = 0; i < sd.size(); i++) {
    cout << " " << sd.at(i);
  }
  cout << endl;
  assert(sd.query(0, 3) == -2);
  return 0;
}
\end{lstlisting}
\begin{exampleoutput}
Values: 6 -2 4 8 10
\end{exampleoutput}
\subsection{Segment Tree (Point Update)}
Maintain a fixed-size array while supporting dynamic queries of contiguous sub-arrays and dynamic updates of individual indices.

The query operation is defined by an associative aggregate function \inlinecode{combine(a, b)}. The default code below assumes a numerical array type, defining queries for the ``min'' of the target range. Another possible query operation is ``sum'', in which case \inlinecode{combine(a, b)} should return $a + b$.

The point update operation is defined by \inlinecode{apply\_delta(v, d)}, which returns the new value at a single updated index. The default definition below supports updates that ``set'' the chosen array index to a new value. Another possible update operation is ``increment'', in which case \inlinecode{apply\_delta(v, d)} should return $v + d$.

\begin{compactitem}
  \item \inlinecode{SegTree(n, v)} constructs an array of size \inlinecode{n} with indices from 0 to \inlinecode{n - 1}, inclusive, and all values initialized to \inlinecode{v}.
  \item \inlinecode{SegTree(lo, hi)} constructs an array from two random-access iterators as a range \inlinecode{[lo, hi)}, initialized to the elements of the range in the same order.
  \item \inlinecode{size()} returns the size of the array.
  \item \inlinecode{at(i)} returns the value at index \inlinecode{i}.
  \item \inlinecode{query(lo, hi)} returns the result of \inlinecode{combine()} applied to all indices from \inlinecode{lo} to \inlinecode{hi}, inclusive. If \inlinecode{lo == hi}, then the single specified value is returned.
  \item \inlinecode{update(i, d)} assigns the value \inlinecode{v} at index \inlinecode{i} to \inlinecode{apply\_delta(v, d)}.
\end{compactitem}

\Needspace*{3\baselineskip}
\textbf{Time Complexity}:
\begin{compactitem}
  \item $O(n)$ per call to both constructors, where $n$ is the size of the array.
  \item $O(1)$ per call to \inlinecode{size()}.
  \item $O(\log n)$ per call to \inlinecode{at()}, \inlinecode{update()}, and \inlinecode{query()}.
\end{compactitem}

\Needspace*{3\baselineskip}
\textbf{Space Complexity}:
\begin{compactitem}
  \item $O(n)$ for storage of the array elements.
  \item $O(\log n)$ auxiliary stack space for \inlinecode{update()} and \inlinecode{query()}.
  \item $O(1)$ auxiliary for \inlinecode{size()}.
\end{compactitem}

\begin{lstlisting}[language=C++]
#include <algorithm>
#include <vector>

template<class T>
class SegTree {
  static T combine(const T &a, const T &b) { return std::min(a, b); }
  static T apply_delta(const T &v, const T &d) { return d; }

  int len;
  std::vector<T> value;

  void build(int i, int lo, int hi, const T &v) {
    if (lo == hi) {
      value[i] = v;
      return;
    }
    int mid = lo + (hi - lo) / 2;
    build(i * 2 + 1, lo, mid, v);
    build(i * 2 + 2, mid + 1, hi, v);
    value[i] = combine(value[i * 2 + 1], value[i * 2 + 2]);
  }

  template<class It>
  void build(int i, int lo, int hi, It arr) {
    if (lo == hi) {
      value[i] = *(arr + lo);
      return;
    }
    int mid = lo + (hi - lo) / 2;
    build(i * 2 + 1, lo, mid, arr);
    build(i * 2 + 2, mid + 1, hi, arr);
    value[i] = combine(value[i * 2 + 1], value[i * 2 + 2]);
  }

  T query(int i, int lo, int hi, int tgt_lo, int tgt_hi) const {
    if (lo == tgt_lo && hi == tgt_hi) {
      return value[i];
    }
    int mid = lo + (hi - lo) / 2;
    if (tgt_lo <= mid && mid < tgt_hi) {
      return combine(
          query(i * 2 + 1, lo, mid, tgt_lo, std::min(tgt_hi, mid)),
          query(i * 2 + 2, mid + 1, hi, std::max(tgt_lo, mid + 1), tgt_hi)
      );
    }
    if (tgt_lo <= mid) {
      return query(i * 2 + 1, lo, mid, tgt_lo, std::min(tgt_hi, mid));
    }
    return query(i * 2 + 2, mid + 1, hi, std::max(tgt_lo, mid + 1), tgt_hi);
  }

  void update(int i, int lo, int hi, int target, const T &d) {
    if (target < lo || target > hi) {
      return;
    }
    if (lo == hi) {
      value[i] = apply_delta(value[i], d);
      return;
    }
    int mid = lo + (hi - lo) / 2;
    update(i * 2 + 1, lo, mid, target, d);
    update(i * 2 + 2, mid + 1, hi, target, d);
    value[i] = combine(value[i * 2 + 1], value[i * 2 + 2]);
  }

 public:
  explicit SegTree(int n, const T &v = T()) : len(n), value(4 * len) {
    if (len > 0) {
      build(0, 0, len - 1, v);
    }
  }

  template<class It>
  SegTree(It lo, It hi) : len(hi - lo), value(4 * len) {
    if (len > 0) {
      build(0, 0, len - 1, lo);
    }
  }

  int size() const { return len; }
  T at(int i) const { return query(i, i); }
  T query(int lo, int hi) const { return query(0, 0, len - 1, lo, hi); }
  void update(int i, const T &d) { update(0, 0, len - 1, i, d); }
};
\end{lstlisting}
\Needspace*{4\baselineskip}
\textbf{Example Usage}\par
\begin{lstlisting}[language=C++]
#include <cassert>
#include <iostream>
using namespace std;

int main() {
  vector<int> arr{6, -2, 1, 8, 10};
  SegTree<int> t(arr.begin(), arr.end());
  t.update(2, 4);
  cout << "Values:";
  for (int i = 0; i < t.size(); i++) {
    cout << " " << t.at(i);
  }
  cout << endl;
  assert(t.query(0, 3) == -2);
  return 0;
}
\end{lstlisting}
\begin{exampleoutput}
Values: 6 -2 4 8 10
The minimum value in the range [0, 3] is -2.
\end{exampleoutput}
\subsection{Segment Tree (Range Update)}
Maintain a fixed-size array while supporting both dynamic queries and updates of contiguous subarrays via the lazy propagation technique.

The query operation is defined by an associative aggregate function \inlinecode{combine(a, b)}. The default code below assumes a numerical array type, defining queries for the ``min'' of the target range. Another possible query operation is ``sum'', in which case \inlinecode{combine(a, b)} should return $a + b$.

Range updates are defined by \inlinecode{apply\_delta(v, d, len)}, which applies an update delta \inlinecode{d} to an aggregate summary \inlinecode{v} representing \inlinecode{len} array values, and by \inlinecode{compose\_deltas(old, d)}, which combines a pending older delta with a newer delta in that order. These functions do not support arbitrary combinations: applying a delta to a combined segment must be equivalent to applying it to each child segment and then combining the results, and composed deltas must be equivalent to performing their updates sequentially. The default code below defines range assignment. For range increment, \inlinecode{compose\_deltas(old, d)} should return \inlinecode{old + d}; \inlinecode{apply\_delta(v, d, len)} should return \inlinecode{v + d} for range-min/range-max queries, and \inlinecode{v + d * len} for range-sum queries.

\begin{compactitem}
  \item \inlinecode{SegTree(n, v)} constructs an array of size \inlinecode{n} with indices from 0 to \inlinecode{n - 1}, inclusive, and all values initialized to \inlinecode{v}.
  \item \inlinecode{SegTree(lo, hi)} constructs an array from two random-access iterators as a range \inlinecode{[lo, hi)}, initialized to the elements of the range in the same order.
  \item \inlinecode{size()} returns the size of the array.
  \item \inlinecode{at(i)} returns the value at index \inlinecode{i}, where \inlinecode{i} is between 0 and \inlinecode{size() - 1}.
  \item \inlinecode{query(lo, hi)} returns the result of \inlinecode{combine()} applied to all indices from \inlinecode{lo} to \inlinecode{hi}, inclusive. If \inlinecode{lo == hi}, then the single specified value is returned.
  \item \inlinecode{update(i, d)} assigns the value \inlinecode{v} at index \inlinecode{i} to \inlinecode{apply\_delta(v, d)}.
  \item \inlinecode{update(lo, hi, d)} modifies the value at each array index from \inlinecode{lo} to \inlinecode{hi}, inclusive, by applying the delta \inlinecode{d} to each value.
\end{compactitem}

\Needspace*{3\baselineskip}
\textbf{Time Complexity}:
\begin{compactitem}
  \item $O(n)$ per call to both constructors, where $n$ is the size of the array.
  \item $O(1)$ per call to \inlinecode{size()}.
  \item $O(\log n)$ per call to \inlinecode{at()}, \inlinecode{update()}, and \inlinecode{query()}.
\end{compactitem}

\Needspace*{3\baselineskip}
\textbf{Space Complexity}:
\begin{compactitem}
  \item $O(n)$ for storage of the array elements.
  \item $O(\log n)$ auxiliary stack space for \inlinecode{update()} and \inlinecode{query()}.
  \item $O(1)$ auxiliary for \inlinecode{size()}.
\end{compactitem}

\begin{lstlisting}[language=C++]
#include <algorithm>
#include <vector>

template<class T>
class SegTree {
  static T combine(const T &a, const T &b) { return std::min(a, b); }
  static T apply_delta(const T &v, const T &d, long long len) { return d; }
  static T compose_deltas(const T &d1, const T &d2) { return d2; }

  int len;
  std::vector<T> value, delta;
  std::vector<bool> pending;

  void build(int i, int lo, int hi, const T &v) {
    if (lo == hi) {
      value[i] = v;
      return;
    }
    int mid = lo + (hi - lo) / 2;
    build(i * 2 + 1, lo, mid, v);
    build(i * 2 + 2, mid + 1, hi, v);
    value[i] = combine(value[i * 2 + 1], value[i * 2 + 2]);
  }

  template<class It>
  void build(int i, int lo, int hi, It arr) {
    if (lo == hi) {
      value[i] = *(arr + lo);
      return;
    }
    int mid = lo + (hi - lo) / 2;
    build(i * 2 + 1, lo, mid, arr);
    build(i * 2 + 2, mid + 1, hi, arr);
    value[i] = combine(value[i * 2 + 1], value[i * 2 + 2]);
  }

  void push_delta(int i, int lo, int hi) {
    if (pending[i]) {
      value[i] = apply_delta(value[i], delta[i], hi - lo + 1);
      if (lo != hi) {
        int l = 2 * i + 1, r = 2 * i + 2;
        delta[l] = pending[l] ? compose_deltas(delta[l], delta[i]) : delta[i];
        delta[r] = pending[r] ? compose_deltas(delta[r], delta[i]) : delta[i];
        pending[l] = pending[r] = true;
      }
      pending[i] = false;
    }
  }

  T query(int i, int lo, int hi, int tgt_lo, int tgt_hi) {
    push_delta(i, lo, hi);
    if (lo == tgt_lo && hi == tgt_hi) {
      return value[i];
    }
    int mid = lo + (hi - lo) / 2;
    if (tgt_lo <= mid && mid < tgt_hi) {
      return combine(
          query(i * 2 + 1, lo, mid, tgt_lo, std::min(tgt_hi, mid)),
          query(i * 2 + 2, mid + 1, hi, std::max(tgt_lo, mid + 1), tgt_hi)
      );
    }
    if (tgt_lo <= mid) {
      return query(i * 2 + 1, lo, mid, tgt_lo, std::min(tgt_hi, mid));
    }
    return query(i * 2 + 2, mid + 1, hi, std::max(tgt_lo, mid + 1), tgt_hi);
  }

  void update(int i, int lo, int hi, int tgt_lo, int tgt_hi, const T &d) {
    push_delta(i, lo, hi);
    if (hi < tgt_lo || lo > tgt_hi) {
      return;
    }
    if (tgt_lo <= lo && hi <= tgt_hi) {
      delta[i] = d;
      pending[i] = true;
      push_delta(i, lo, hi);
      return;
    }
    update(2 * i + 1, lo, lo + (hi - lo) / 2, tgt_lo, tgt_hi, d);
    update(2 * i + 2, lo + (hi - lo) / 2 + 1, hi, tgt_lo, tgt_hi, d);
    value[i] = combine(value[2 * i + 1], value[2 * i + 2]);
  }

 public:
  explicit SegTree(int n, const T &v = T())
      : len(n), value(4 * len), delta(4 * len), pending(4 * len, false) {
    if (len > 0) {
      build(0, 0, len - 1, v);
    }
  }

  template<class It>
  SegTree(It lo, It hi) : len(hi - lo), value(4 * len), delta(4 * len), pending(4 * len, false) {
    if (len > 0) {
      build(0, 0, len - 1, lo);
    }
  }

  int size() const { return len; }
  T at(int i) { return query(i, i); }
  T query(int lo, int hi) { return query(0, 0, len - 1, lo, hi); }
  void update(int i, const T &d) { update(0, 0, len - 1, i, i, d); }
  void update(int lo, int hi, const T &d) { update(0, 0, len - 1, lo, hi, d); }
};
\end{lstlisting}
\Needspace*{4\baselineskip}
\textbf{Example Usage}\par
\begin{lstlisting}[language=C++]
#include <cassert>
#include <iostream>
using namespace std;

int main() {
  vector<int> arr{6, -2, 1, 8, 10};
  SegTree<int> t(arr.begin(), arr.end());
  t.update(2, 4);
  cout << "Values:";
  for (int i = 0; i < t.size(); i++) {
    cout << " " << t.at(i);
  }
  cout << endl;
  assert(t.query(0, 3) == -2);
  t.update(0, 4, 5);
  t.update(3, 2);
  t.update(3, 1);
  cout << "Values:";
  for (int i = 0; i < t.size(); i++) {
    cout << " " << t.at(i);
  }
  cout << endl;
  assert(t.query(0, 3) == 1);
  return 0;
}
\end{lstlisting}
\begin{exampleoutput}
Values: 6 -2 4 8 10
Values: 5 5 5 1 5
\end{exampleoutput}
\subsection{Segment Tree (Compressed)}
Maintain an array over a large index range while supporting both dynamic queries and updates of contiguous subarrays via the lazy propagation technique. This implementation uses lazy initialization of nodes to conserve memory: only the nodes covering touched indices are ever allocated, so indices up to \inlinecode{N} are supported without preallocating the whole tree.

The query operation is defined by an associative aggregate function \inlinecode{combine(a, b)}. Since untouched nodes are implicit, \inlinecode{repeat\_value(v, len)} must return the aggregate summary of \inlinecode{len} copies of the initial value \inlinecode{v}. The default code below assumes a numerical array type, defining queries for the ``min'' of the target range. For range-sum queries, \inlinecode{combine(a, b)} should return $a + b$ and \inlinecode{repeat\_value(v, len)} should return $v \cdot len$.

Range updates are defined by \inlinecode{apply\_delta(v, d, len)}, which applies an update delta \inlinecode{d} to an aggregate summary \inlinecode{v} representing \inlinecode{len} array values, and by \inlinecode{compose\_deltas(old, d)}, which combines a pending older delta with a newer delta in that order. These functions do not support arbitrary combinations: applying a delta to a combined segment must be equivalent to applying it to each child segment and then combining the results, and composed deltas must be equivalent to performing their updates sequentially. The default code below defines range assignment. For range increment, \inlinecode{compose\_deltas(old, d)} should return \inlinecode{old + d}; \inlinecode{apply\_delta(v, d, len)} should return \inlinecode{v + d} for range-min/range-max queries, and \inlinecode{v + d * len} for range-sum queries.

\begin{compactitem}
  \item \inlinecode{SegTree(v)} constructs an array over indices 0 to \inlinecode{N}, inclusive, with every value implicitly initialized to \inlinecode{v}. Nodes are allocated lazily as indices are touched.
  \item \inlinecode{at(i)} returns the value at index \inlinecode{i}, where \inlinecode{i} is between 0 and \inlinecode{N}.
  \item \inlinecode{query(lo, hi)} returns the result of \inlinecode{combine()} applied to all indices from \inlinecode{lo} to \inlinecode{hi}, inclusive. If \inlinecode{lo == hi}, then the single specified value is returned.
  \item \inlinecode{update(i, d)} assigns the value \inlinecode{v} at index \inlinecode{i} to \inlinecode{apply\_delta(v, d)}.
  \item \inlinecode{update(lo, hi, d)} modifies the value at each array index from \inlinecode{lo} to \inlinecode{hi}, inclusive, by applying the delta \inlinecode{d} to each value.
\end{compactitem}

\Needspace*{3\baselineskip}
\textbf{Time Complexity}:
\begin{compactitem}
  \item $O(1)$ per call to the constructor.
  \item $O(\log N)$ per call to \inlinecode{at()}, \inlinecode{update()}, and \inlinecode{query()}.
\end{compactitem}

\Needspace*{3\baselineskip}
\textbf{Space Complexity}:
\begin{compactitem}
  \item $O(k \log\left(N\right))$ for storage after $k$ index updates.
  \item $O(\log N)$ auxiliary stack space for \inlinecode{update()} and \inlinecode{query()}.
\end{compactitem}

\begin{lstlisting}[language=C++]
#include <algorithm>
#include <cstddef>

template<class T>
class SegTree {
  static const int N = 1000000000;

  static T combine(const T &a, const T &b) { return std::min(a, b); }
  static T repeat_value(const T &v, long long len) { return v; }
  static T apply_delta(const T &v, const T &d, long long len) { return d; }
  static T compose_deltas(const T &d1, const T &d2) { return d2; }

  struct Node {
    T value, delta;
    bool pending;
    Node *left, *right;

    explicit Node(const T &v) : value(v), pending(false), left(nullptr), right(nullptr) {}
  } *root;

  T init;

  void update_delta(Node *&n, const T &d, long long len) {
    if (n == nullptr) {
      n = new Node(repeat_value(init, len));
    }
    n->delta = n->pending ? compose_deltas(n->delta, d) : d;
    n->pending = true;
  }

  void push_delta(Node *n, int lo, int hi) {
    if (n == nullptr) {
      return;
    }
    if (n->pending) {
      n->value = apply_delta(n->value, n->delta, hi - lo + 1);
      if (lo != hi) {
        int mid = lo + (hi - lo) / 2;
        update_delta(n->left, n->delta, mid - lo + 1);
        update_delta(n->right, n->delta, hi - mid);
      }
    }
    n->pending = false;
  }

  T query(Node *n, int lo, int hi, int tgt_lo, int tgt_hi) {
    if (n == nullptr) {
      return repeat_value(init, hi - lo + 1);
    }
    push_delta(n, lo, hi);
    if (lo == tgt_lo && hi == tgt_hi) {
      return n->value;
    }
    int mid = lo + (hi - lo) / 2;
    if (tgt_lo <= mid && mid < tgt_hi) {
      return combine(
          query(n->left, lo, mid, tgt_lo, std::min(tgt_hi, mid)),
          query(n->right, mid + 1, hi, std::max(tgt_lo, mid + 1), tgt_hi)
      );
    }
    if (tgt_lo <= mid) {
      return query(n->left, lo, mid, tgt_lo, std::min(tgt_hi, mid));
    }
    return query(n->right, mid + 1, hi, std::max(tgt_lo, mid + 1), tgt_hi);
  }

  void update(Node *&n, int lo, int hi, int tgt_lo, int tgt_hi, const T &d) {
    if (n == nullptr) {
      if (hi < tgt_lo || lo > tgt_hi) {
        return;
      }
      n = new Node(repeat_value(init, hi - lo + 1));
    } else {
      push_delta(n, lo, hi);
    }
    if (hi < tgt_lo || lo > tgt_hi) {
      return;
    }
    if (tgt_lo <= lo && hi <= tgt_hi) {
      n->delta = d;
      n->pending = true;
      push_delta(n, lo, hi);
      return;
    }
    int mid = lo + (hi - lo) / 2;
    update(n->left, lo, mid, tgt_lo, tgt_hi, d);
    update(n->right, mid + 1, hi, tgt_lo, tgt_hi, d);
    T left_value = (n->left != nullptr) ? n->left->value : repeat_value(init, mid - lo + 1);
    T right_value = (n->right != nullptr) ? n->right->value : repeat_value(init, hi - mid);
    n->value = combine(left_value, right_value);
  }

  void clean_up(Node *n) {
    if (n != nullptr) {
      clean_up(n->left);
      clean_up(n->right);
      delete n;
    }
  }

 public:
  explicit SegTree(const T &v = T()) : root(nullptr), init(v) {}

  ~SegTree() { clean_up(root); }
  T at(int i) { return query(i, i); }
  T query(int lo, int hi) { return query(root, 0, N, lo, hi); }
  void update(int i, const T &d) { return update(i, i, d); }
  void update(int lo, int hi, const T &d) { return update(root, 0, N, lo, hi, d); }
};
\end{lstlisting}
\Needspace*{4\baselineskip}
\textbf{Example Usage}\par
\begin{lstlisting}[language=C++]
#include <cassert>
#include <iostream>
using namespace std;

int main() {
  SegTree<int> t(0);
  t.update(0, 6);
  t.update(1, -2);
  t.update(2, 4);
  t.update(3, 8);
  t.update(4, 10);
  cout << "Values:";
  for (int i = 0; i < 5; i++) {
    cout << " " << t.at(i);
  }
  cout << endl;
  assert(t.query(0, 3) == -2);
  t.update(0, 4, 5);
  t.update(3, 2);
  t.update(3, 1);
  cout << "Values:";
  for (int i = 0; i < 5; i++) {
    cout << " " << t.at(i);
  }
  cout << endl;
  assert(t.query(0, 3) == 1);
  return 0;
}
\end{lstlisting}
\begin{exampleoutput}
Values: 6 -2 4 8 10
Values: 5 5 5 1 5
\end{exampleoutput}
\subsection{Segment Tree (Persistent)}
Maintains immutable versions of a segment tree under point updates. Each update creates a new root by copying only the nodes on the path from the root to the updated leaf, while all unchanged subtrees are shared with previous versions.

The implementation below stores sums over a fixed index range \inlinecode{[0, n - 1]}. Persistent segment trees are useful for offline order-statistics queries, rollback-like histories, and functional dynamic programming states.

\begin{compactitem}
  \item \inlinecode{PersistentSegTree(a)} constructs version 0 from array \inlinecode{a}.
  \item \inlinecode{update(version, index, value)} returns a new version where \inlinecode{a[index]} is set to \inlinecode{value}, leaving the old version unchanged.
  \item \inlinecode{query(version, lo, hi)} returns the sum of the inclusive range \inlinecode{[lo, hi]} in the specified version.
  \item \inlinecode{versions()} returns the number of stored roots.
\end{compactitem}

\Needspace*{3\baselineskip}
\textbf{Time Complexity}:
\begin{compactitem}
  \item $O(n)$ for construction.
  \item $O(\log n)$ per call to \inlinecode{update(version, index, value)} and \inlinecode{query(version, lo, hi)}.
\end{compactitem}

\Needspace*{3\baselineskip}
\textbf{Space Complexity}:
\begin{compactitem}
  \item $O(n + q \log n)$ for $q$ updates.
  \item $O(\log n)$ auxiliary stack space per update and query.
\end{compactitem}

\begin{lstlisting}[language=C++]
#include <vector>

class PersistentSegTree {
  struct node {
    long long sum;
    int left, right;

    node(long long sum = 0, int left = -1, int right = -1) : sum(sum), left(left), right(right) {}
  };

  int n;
  std::vector<node> tree;
  std::vector<int> root;

  int build(const std::vector<int> &a, int lo, int hi) {
    if (lo == hi) {
      tree.emplace_back(a[lo]);
      return static_cast<int>(tree.size()) - 1;
    }
    int mid = lo + (hi - lo) / 2;
    int left = build(a, lo, mid);
    int right = build(a, mid + 1, hi);
    tree.emplace_back(tree[left].sum + tree[right].sum, left, right);
    return static_cast<int>(tree.size()) - 1;
  }

  int update(int cur, int lo, int hi, int index, int value) {
    if (lo == hi) {
      tree.emplace_back(value);
      return static_cast<int>(tree.size()) - 1;
    }
    int mid = lo + (hi - lo) / 2;
    int left = tree[cur].left, right = tree[cur].right;
    if (index <= mid) {
      left = update(left, lo, mid, index, value);
    } else {
      right = update(right, mid + 1, hi, index, value);
    }
    tree.emplace_back(tree[left].sum + tree[right].sum, left, right);
    return static_cast<int>(tree.size()) - 1;
  }

  long long query(int cur, int lo, int hi, int qlo, int qhi) const {
    if (qlo <= lo && hi <= qhi) {
      return tree[cur].sum;
    }
    int mid = lo + (hi - lo) / 2;
    long long res = 0;
    if (qlo <= mid) {
      res += query(tree[cur].left, lo, mid, qlo, qhi);
    }
    if (mid < qhi) {
      res += query(tree[cur].right, mid + 1, hi, qlo, qhi);
    }
    return res;
  }

 public:
  explicit PersistentSegTree(const std::vector<int> &a) : n(a.size()) {
    if (n > 0) {
      root.push_back(build(a, 0, n - 1));
    }
  }

  int versions() const { return root.size(); }

  int update(int version, int index, int value) {
    root.push_back(update(root[version], 0, n - 1, index, value));
    return static_cast<int>(root.size()) - 1;
  }

  long long query(int version, int lo, int hi) const {
    return query(root[version], 0, n - 1, lo, hi);
  }
};
\end{lstlisting}
\Needspace*{4\baselineskip}
\textbf{Example Usage}\par
\begin{lstlisting}[language=C++]
#include <cassert>
using namespace std;

int main() {
  vector<int> a{1, 2, 3, 4};
  PersistentSegTree t(a);
  int v1 = t.update(0, 1, 10);
  int v2 = t.update(v1, 3, -1);
  assert(t.query(0, 0, 3) == 10);
  assert(t.query(v1, 0, 3) == 18);
  assert(t.query(v2, 1, 3) == 12);
  assert(t.versions() == 3);
  return 0;
}
\end{lstlisting}
\subsection{Cartesian Tree}
Builds the min Cartesian tree of an array in linear time. A Cartesian tree is a binary tree whose inorder traversal is the original array order and whose heap property is determined by array values. For a min Cartesian tree, each parent has value less than or equal to its children.

Cartesian trees connect arrays, stacks, range minimum queries, treaps, and largest rectangle style problems. Equal values are broken by position, so earlier equal values become ancestors of later equal values.

\begin{compactitem}
  \item \inlinecode{CartesianTree(a)} constructs the tree for array \inlinecode{a}.
  \item \inlinecode{root} stores the root index.
  \item \inlinecode{parent[i]}, \inlinecode{left[i]}, and \inlinecode{right[i]} store neighboring node indices, or $-1$ if absent.
\end{compactitem}

\textbf{Time Complexity}: $O(n)$ construction time, where $n$ is the array size.

\textbf{Space Complexity}: $O(n)$ for tree arrays and the monotone stack.

\begin{lstlisting}[language=C++]
#include <vector>

struct CartesianTree {
  int root;
  std::vector<int> parent, left, right;

  explicit CartesianTree(const std::vector<int> &a)
      : root(-1), parent(a.size(), -1), left(a.size(), -1), right(a.size(), -1) {
    std::vector<int> st;
    for (int i = 0, n = static_cast<int>(a.size()); i < n; i++) {
      int last = -1;
      while (!st.empty() && a[i] < a[st.back()]) {
        last = st.back();
        st.pop_back();
      }
      if (!st.empty()) {
        right[st.back()] = i;
        parent[i] = st.back();
      }
      if (last != -1) {
        left[i] = last;
        parent[last] = i;
      }
      st.push_back(i);
    }
    root = st.empty() ? -1 : st[0];
  }
};
\end{lstlisting}
\Needspace*{4\baselineskip}
\textbf{Example Usage}\par
\begin{lstlisting}[language=C++]
#include <cassert>
using namespace std;

void inorder(const CartesianTree &t, int u, vector<int> *order) {
  if (u == -1) {
    return;
  }
  inorder(t, t.left[u], order);
  order->push_back(u);
  inorder(t, t.right[u], order);
}

int main() {
  vector<int> a{3, 2, 6, 1, 9};
  CartesianTree t(a);
  assert(t.root == 3);
  assert(t.parent[1] == 3);
  assert(t.parent[4] == 3);

  vector<int> order;
  inorder(t, t.root, &order);
  for (int i = 0; i < 5; i++) {
    assert(order[i] == i);
  }
  return 0;
}
\end{lstlisting}
\subsection{Implicit Treap}
Maintain a dynamically-sized array using a balanced binary search tree while supporting both dynamic queries and updates of contiguous subarrays via the lazy propagation technique. A treap maintains a balanced binary tree structure by preserving the heap property on the randomly generated priority values of nodes, thereby making insertions and deletions run in $O(\log n)$ with high probability.

The query operation is defined by an associative aggregate function \inlinecode{combine(a, b)}. The default code below assumes a numerical array type, defining queries for the ``min'' of the target range. Another possible query operation is ``sum'', in which case \inlinecode{combine(a, b)} should return $a + b$.

Range updates are defined by \inlinecode{apply\_delta(v, d, len)}, which applies an update delta \inlinecode{d} to an aggregate summary \inlinecode{v} representing \inlinecode{len} array values, and by \inlinecode{compose\_deltas(old, d)}, which combines a pending older delta with a newer delta in that order. These functions do not support arbitrary combinations: applying a delta to a combined segment must be equivalent to applying it to each child segment and then combining the results, and composed deltas must be equivalent to performing their updates sequentially. The default code below defines range assignment. For range increment, \inlinecode{compose\_deltas(old, d)} should return \inlinecode{old + d}; \inlinecode{apply\_delta(v, d, len)} should return \inlinecode{v + d} for range-min/range-max queries, and \inlinecode{v + d * len} for range-sum queries.

This data structure shares every operation of one-dimensional segment trees in this section, with the additional operations \inlinecode{empty()}, \inlinecode{insert()}, \inlinecode{erase()}, \inlinecode{push\_back()}, and \inlinecode{pop\_back()} analogous to those of \inlinecode{std::vector} (here, \inlinecode{insert()} and \inlinecode{erase()} both take an index instead of an iterator).

\Needspace*{3\baselineskip}
\textbf{Time Complexity}:
\begin{compactitem}
  \item $O(n)$ per call to both constructors, where $n$ is the size of the array.
  \item $O(1)$ per call to \inlinecode{size()} and \inlinecode{empty()}.
  \item $O(\log n)$ on average per call to all other operations.
\end{compactitem}

\Needspace*{3\baselineskip}
\textbf{Space Complexity}:
\begin{compactitem}
  \item $O(n)$ for storage of the array elements.
  \item $O(1)$ auxiliary for \inlinecode{size()} and \inlinecode{empty()}.
  \item $O(\log n)$ auxiliary stack space for all other operations.
\end{compactitem}

\begin{lstlisting}[language=C++]
#include <cstdlib>

template<class T>
class ImplicitTreap {
  static T combine(const T &a, const T &b) { return a < b ? a : b; }
  static T apply_delta(const T &v, const T &d, long long len) { return d; }
  static T compose_deltas(const T &d1, const T &d2) { return d2; }

  struct Node {
    static inline int rand32() { return (rand() & 0x7fff) | ((rand() & 0x7fff) << 15); }

    T value, subtree_value, delta;
    bool pending;
    int size, priority;
    Node *left, *right;

    explicit Node(const T &v)
        : value(v),
          subtree_value(v),
          pending(false),
          size(1),
          priority(rand32()),
          left(nullptr),
          right(nullptr) {}
  } *root;

  static int size(Node *n) { return (n == nullptr) ? 0 : n->size; }

  static void update_value(Node *n) {
    if (n == nullptr) {
      return;
    }
    n->subtree_value = n->value;
    if (n->left != nullptr) {
      n->subtree_value = combine(n->subtree_value, n->left->subtree_value);
    }
    if (n->right != nullptr) {
      n->subtree_value = combine(n->subtree_value, n->right->subtree_value);
    }
    n->size = 1 + size(n->left) + size(n->right);
  }

  static void update_delta(Node *n, const T &d) {
    if (n != nullptr) {
      n->value = apply_delta(n->value, d, 1);
      n->subtree_value = apply_delta(n->subtree_value, d, n->size);
      n->delta = n->pending ? compose_deltas(n->delta, d) : d;
      n->pending = true;
    }
  }

  static void push_delta(Node *n) {
    if (n == nullptr || !n->pending) {
      return;
    }
    if (n->size > 1) {
      update_delta(n->left, n->delta);
      update_delta(n->right, n->delta);
    }
    n->pending = false;
  }

  static void merge(Node *&n, Node *left, Node *right) {
    push_delta(left);
    push_delta(right);
    if (left == nullptr) {
      n = right;
    } else if (right == nullptr) {
      n = left;
    } else if (left->priority < right->priority) {
      merge(left->right, left->right, right);
      n = left;
    } else {
      merge(right->left, left, right->left);
      n = right;
    }
    update_value(n);
  }

  static void split(Node *n, Node *&left, Node *&right, int i) {
    push_delta(n);
    if (n == nullptr) {
      left = right = nullptr;
    } else if (i <= size(n->left)) {
      split(n->left, left, n->left, i);
      right = n;
    } else {
      split(n->right, n->right, right, i - size(n->left) - 1);
      left = n;
    }
    update_value(n);
  }

  static void insert(Node *&n, Node *new_node, int i) {
    push_delta(n);
    if (n == nullptr) {
      n = new_node;
    } else if (new_node->priority < n->priority) {
      split(n, new_node->left, new_node->right, i);
      n = new_node;
    } else if (i <= size(n->left)) {
      insert(n->left, new_node, i);
    } else {
      insert(n->right, new_node, i - size(n->left) - 1);
    }
    update_value(n);
  }

  static void erase(Node *&n, int i) {
    push_delta(n);
    if (i == size(n->left)) {
      Node *left = n->left, *right = n->right;
      delete n;
      merge(n, left, right);
    } else if (i < size(n->left)) {
      erase(n->left, i);
    } else {
      erase(n->right, i - size(n->left) - 1);
    }
    update_value(n);
  }

  static Node *select(Node *n, int i) {
    push_delta(n);
    if (i < size(n->left)) {
      return select(n->left, i);
    }
    if (i > size(n->left)) {
      return select(n->right, i - size(n->left) - 1);
    }
    return n;
  }

  void clean_up(Node *&n) {
    if (n != nullptr) {
      clean_up(n->left);
      clean_up(n->right);
      delete n;
    }
  }

 public:
  explicit ImplicitTreap(int n = 0, const T &v = T()) : root(nullptr) {
    for (int i = 0; i < n; i++) {
      push_back(v);
    }
  }

  template<class It>
  ImplicitTreap(It lo, It hi) : root(nullptr) {
    for (; lo != hi; ++lo) {
      push_back(*lo);
    }
  }

  ~ImplicitTreap() { clean_up(root); }
  int size() const { return size(root); }
  bool empty() const { return root == nullptr; }
  void insert(int i, const T &v) { insert(root, new Node(v), i); }
  void erase(int i) { erase(root, i); }
  void push_back(const T &v) { insert(size(), v); }
  void pop_back() { erase(size() - 1); }
  T at(int i) const { return select(root, i)->value; }

  T query(int lo, int hi) {
    Node *l1, *r1, *l2, *r2, *t;
    split(root, l1, r1, hi + 1);
    split(l1, l2, r2, lo);
    T res = r2->subtree_value;
    merge(t, l2, r2);
    merge(root, t, r1);
    return res;
  }

  void update(int lo, int hi, const T &d) {
    Node *l1, *r1, *l2, *r2, *t;
    split(root, l1, r1, hi + 1);
    split(l1, l2, r2, lo);
    update_delta(r2, d);
    merge(t, l2, r2);
    merge(root, t, r1);
  }

  void update(int i, const T &d) { update(i, i, d); }
};
\end{lstlisting}
\Needspace*{4\baselineskip}
\textbf{Example Usage}\par
\begin{lstlisting}[language=C++]
#include <iostream>
#include <vector>
using namespace std;

void print(ImplicitTreap<int> &t) {
  cout << "Values:";
  for (int i = 0; i < t.size(); i++) {
    cout << " " << t.at(i);
  }
  cout << " (min: " << t.query(0, t.size() - 1) << ")" << endl;
}

int main() {
  vector<int> arr{99, -2, 1, 8, 10};
  ImplicitTreap<int> t(arr.begin(), arr.end());
  t.push_back(11);
  t.push_back(12);
  t.pop_back();
  print(t);
  t.insert(0, 90);
  t.erase(1);
  print(t);
  t.update(0, 1, 2);
  print(t);
  return 0;
}
\end{lstlisting}
\begin{exampleoutput}
Values: 99 -2 1 8 10 11 (min: -2)
Values: 90 -2 1 8 10 11 (min: -2)
Values: 2 2 1 8 10 11 (min: 1)
\end{exampleoutput}

\section{Range Queries in Two Dimensions}
\setcounter{section}{5}
\setcounter{subsection}{0}
\subsection{Quadtree (Point Update)}
Maintain a two-dimensional array while supporting dynamic queries of rectangular sub-arrays and dynamic updates of individual indices. This implementation uses lazy initialization of nodes to conserve memory while supporting large indices.

The query operation is defined by a commutative associative aggregate function \inlinecode{combine(a, b)}. Because untouched regions are implicit, \inlinecode{repeat\_value(v, area)} must return the aggregate summary of \inlinecode{area} copies of the initial value \inlinecode{v}. The default code below assumes a numerical array type, defining queries for the ``min'' of the target range. For rectangle-sum queries, \inlinecode{combine(a, b)} should return $a + b$ and \inlinecode{repeat\_value(v, area)} should return $v \cdot area$.

The point update operation is defined by \inlinecode{apply\_delta(v, d)}, which returns the new value at one updated cell. The default code below defines updates that ``set'' the chosen cell to a new value. Another possible update operation is ``increment'', in which case \inlinecode{apply\_delta(v, d)} should return $v + d$.

\begin{compactitem}
  \item \inlinecode{Quadtree(v)} constructs a two-dimensional array with rows from 0 to \inlinecode{R} and columns from 0 to \inlinecode{C}, inclusive. All values are implicitly initialized to \inlinecode{v}.
  \item \inlinecode{at(r, c)} returns the value at row \inlinecode{r}, column \inlinecode{c}.
  \item \inlinecode{query(r1, c1, r2, c2)} returns the result of \inlinecode{combine()} applied to every value in the rectangular region consisting of rows from \inlinecode{r1} to \inlinecode{r2}, inclusive, and columns from \inlinecode{c1} to \inlinecode{c2}, inclusive.
  \item \inlinecode{update(r, c, d)} assigns the value \inlinecode{v} at (\inlinecode{r}, \inlinecode{c}) to \inlinecode{apply\_delta(v, d)}.
\end{compactitem}

\Needspace*{3\baselineskip}
\textbf{Time Complexity}:
\begin{compactitem}
  \item $O(1)$ per call to the constructor.
  \item $O(\max\left(R, C\right))$ per call to \inlinecode{at()}, \inlinecode{update()}, and \inlinecode{query()}.
\end{compactitem}

\Needspace*{3\baselineskip}
\textbf{Space Complexity}:
\begin{compactitem}
  \item $O(n \log\left(\max\left(R, C\right)\right))$ for storage of the quadtree nodes after $n$ point updates.
  \item $O(\log\left(\max\left(R, C\right)\right))$ auxiliary stack space for \inlinecode{update()}, \inlinecode{query()}, and \inlinecode{at()}.
\end{compactitem}

\begin{lstlisting}[language=C++]
#include <algorithm>
#include <cstddef>

template<class T>
class Quadtree {
  static const int R = 1000000000;
  static const int C = 1000000000;

  static T combine(const T &a, const T &b) { return std::min(a, b); }
  static T repeat_value(const T &v, long long area) { return v; }
  static T apply_delta(const T &v, const T &d) { return d; }

  struct Node {
    T value;
    Node *child[4];

    explicit Node(const T &v) : value(v) {
      for (auto &c : child) {
        c = nullptr;
      }
    }
  };

  Node *root;
  T init;

  // Helper variables for query().
  int tgt_r1, tgt_c1, tgt_r2, tgt_c2;
  T res;
  bool found;

  static long long area(int r1, int c1, int r2, int c2) {
    return static_cast<long long>(r2 - r1 + 1) * (c2 - c1 + 1);
  }

  void append_result(const T &v) {
    res = found ? combine(res, v) : v;
    found = true;
  }

  T child_value(Node *n, int r1, int c1, int r2, int c2) {
    return (n != nullptr) ? n->value : repeat_value(init, area(r1, c1, r2, c2));
  }

  void query(Node *n, int r1, int c1, int r2, int c2) {
    if (tgt_r2 < r1 || tgt_r1 > r2 || tgt_c2 < c1 || tgt_c1 > c2) {
      return;
    }
    if (n == nullptr) {
      int rlen = std::min(r2, tgt_r2) - std::max(r1, tgt_r1) + 1;
      int clen = std::min(c2, tgt_c2) - std::max(c1, tgt_c1) + 1;
      append_result(repeat_value(init, static_cast<long long>(rlen) * clen));
      return;
    }
    if (tgt_r1 <= r1 && r2 <= tgt_r2 && tgt_c1 <= c1 && c2 <= tgt_c2) {
      append_result(n->value);
      return;
    }
    int rmid = r1 + (r2 - r1) / 2, cmid = c1 + (c2 - c1) / 2;
    query(n->child[0], r1, c1, rmid, cmid);
    query(n->child[1], rmid + 1, c1, r2, cmid);
    query(n->child[2], r1, cmid + 1, rmid, c2);
    query(n->child[3], rmid + 1, cmid + 1, r2, c2);
  }

  // Helper variables for update().
  int tgt_r, tgt_c;
  T delta;

  void update(Node *&n, int r1, int c1, int r2, int c2) {
    if (tgt_r < r1 || tgt_r > r2 || tgt_c < c1 || tgt_c > c2) {
      return;
    }
    if (n == nullptr) {
      n = new Node(repeat_value(init, area(r1, c1, r2, c2)));
    }
    if (r1 == r2 && c1 == c2) {
      n->value = apply_delta(n->value, delta);
      return;
    }
    int rmid = r1 + (r2 - r1) / 2, cmid = c1 + (c2 - c1) / 2;
    update(n->child[0], r1, c1, rmid, cmid);
    update(n->child[1], rmid + 1, c1, r2, cmid);
    update(n->child[2], r1, cmid + 1, rmid, c2);
    update(n->child[3], rmid + 1, cmid + 1, r2, c2);
    n->value = combine(
        combine(
            child_value(n->child[0], r1, c1, rmid, cmid),
            child_value(n->child[1], rmid + 1, c1, r2, cmid)
        ),
        combine(
            child_value(n->child[2], r1, cmid + 1, rmid, c2),
            child_value(n->child[3], rmid + 1, cmid + 1, r2, c2)
        )
    );
  }

  static void clean_up(Node *n) {
    if (n != nullptr) {
      for (auto *c : n->child) {
        clean_up(c);
      }
      delete n;
    }
  }

 public:
  explicit Quadtree(const T &v = T()) : root(nullptr), init(v) {}

  ~Quadtree() { clean_up(root); }
  T at(int r, int c) { return query(r, c, r, c); }

  T query(int r1, int c1, int r2, int c2) {
    tgt_r1 = r1;
    tgt_c1 = c1;
    tgt_r2 = r2;
    tgt_c2 = c2;
    found = false;
    query(root, 0, 0, R, C);
    return found ? res : repeat_value(init, area(r1, c1, r2, c2));
  }

  void update(int r, int c, const T &d) {
    tgt_r = r;
    tgt_c = c;
    delta = d;
    update(root, 0, 0, R, C);
  }
};
\end{lstlisting}
\Needspace*{4\baselineskip}
\textbf{Example Usage}\par
\begin{lstlisting}[language=C++]
#include <cassert>
#include <iostream>
using namespace std;

int main() {
  Quadtree<int> t(0);
  t.update(0, 0, 7);
  t.update(0, 1, 6);
  t.update(1, 0, 5);
  t.update(1, 1, 4);
  t.update(2, 1, 1);
  t.update(2, 2, 9);
  cout << "Values:" << endl;
  for (int i = 0; i < 3; i++) {
    for (int j = 0; j < 3; j++) {
      cout << t.at(i, j) << " ";
    }
    cout << endl;
  }
  assert(t.query(0, 0, 0, 1) == 6);
  assert(t.query(0, 0, 1, 0) == 5);
  assert(t.query(1, 1, 2, 2) == 0);
  assert(t.query(0, 0, 1000000000, 1000000000) == 0);
  t.update(500000000, 500000000, -100);
  assert(t.query(0, 0, 1000000000, 1000000000) == -100);
  return 0;
}
\end{lstlisting}
\begin{exampleoutput}
Values:
7 6 0
5 4 0
0 1 9
\end{exampleoutput}
\subsection{Quadtree (Range Update)}
Maintain a two-dimensional array while supporting both dynamic queries and updates of rectangular sub-arrays via the lazy propagation technique. This implementation uses lazy initialization of nodes to conserve memory while supporting large indices.

The query operation is defined by a commutative associative aggregate function \inlinecode{combine(a, b)}. Because untouched regions are implicit, \inlinecode{repeat\_value(v, area)} must return the aggregate summary of \inlinecode{area} copies of the initial value \inlinecode{v}. The default code below assumes a numerical array type, defining queries for the ``min'' of the target range. For rectangle-sum queries, \inlinecode{combine(a, b)} should return $a + b$ and \inlinecode{repeat\_value(v, area)} should return $v \cdot area$.

Range updates are defined by \inlinecode{apply\_delta(v, d, area)}, which applies an update delta \inlinecode{d} to an aggregate summary \inlinecode{v} representing \inlinecode{area} cells, and by \inlinecode{compose\_deltas(old, d)}, which combines a pending older delta with a newer delta in that order. These functions do not support arbitrary combinations: applying a delta to a combined region must be equivalent to applying it to each child region and then combining the results, and composed deltas must be equivalent to performing their updates sequentially. The default code below defines rectangle assignment. For rectangle increment, \inlinecode{compose\_deltas(old, d)} should return \inlinecode{old + d}; \inlinecode{apply\_delta(v, d, area)} should return \inlinecode{v + d} for range-min/range-max queries, and \inlinecode{v + d * area} for range-sum queries.

\begin{compactitem}
  \item \inlinecode{Quadtree(v)} constructs a two-dimensional array with rows from 0 to \inlinecode{R} and columns from 0 to \inlinecode{C}, inclusive. All values are implicitly initialized to \inlinecode{v}.
  \item \inlinecode{at(r, c)} returns the value at row \inlinecode{r}, column \inlinecode{c}.
  \item \inlinecode{query(r1, c1, r2, c2)} returns the result of \inlinecode{combine()} applied to every value in the rectangular region consisting of rows from \inlinecode{r1} to \inlinecode{r2} and columns from \inlinecode{c1} to \inlinecode{c2}, inclusive.
  \item \inlinecode{update(r, c, d)} assigns the value \inlinecode{v} at (\inlinecode{r}, \inlinecode{c}) to \inlinecode{apply\_delta(v, d)}.
  \item \inlinecode{update(r1, c1, r2, c2)} modifies the value at each index of the rectangular region consisting of rows from \inlinecode{r1} to \inlinecode{r2} and columns from \inlinecode{c1} to \inlinecode{c2}, inclusive, by applying the delta \inlinecode{d} to each value.
\end{compactitem}

\Needspace*{3\baselineskip}
\textbf{Time Complexity}:
\begin{compactitem}
  \item $O(1)$ per call to the constructor.
  \item $O(\max\left(R, C\right))$ per call to \inlinecode{at()}, \inlinecode{update()}, and \inlinecode{query()}.
\end{compactitem}

\Needspace*{3\baselineskip}
\textbf{Space Complexity}:
\begin{compactitem}
  \item $O(n \log\left(\max\left(R, C\right)\right))$ for storage of the quadtree nodes after $n$ point or rectangle updates.
  \item $O(\log\left(\max\left(R, C\right)\right))$ auxiliary stack space for \inlinecode{update()}, \inlinecode{query()}, and \inlinecode{at()}.
\end{compactitem}

\begin{lstlisting}[language=C++]
#include <algorithm>
#include <cstddef>

template<class T>
class Quadtree {
  static const int R = 1000000000;
  static const int C = 1000000000;

  static T combine(const T &a, const T &b) { return std::min(a, b); }
  static T repeat_value(const T &v, long long area) { return v; }
  static T apply_delta(const T &v, const T &d, long long area) { return d; }
  static T compose_deltas(const T &d1, const T &d2) { return d2; }

  struct Node {
    T value, delta;
    bool pending;
    Node *child[4];

    explicit Node(const T &v) : value(v), pending(false) {
      for (auto &c : child) {
        c = nullptr;
      }
    }
  } *root;

  T init;

  static long long area(int r1, int c1, int r2, int c2) {
    return static_cast<long long>(r2 - r1 + 1) * (c2 - c1 + 1);
  }

  void update_delta(Node *&n, const T &d, long long area) {
    if (n == nullptr) {
      n = new Node(repeat_value(init, area));
    }
    n->delta = n->pending ? compose_deltas(n->delta, d) : d;
    n->pending = true;
  }

  void push_delta(Node *n, int r1, int c1, int r2, int c2) {
    if (n->pending) {
      int rmid = r1 + (r2 - r1) / 2, cmid = c1 + (c2 - c1) / 2;
      long long cur_area = area(r1, c1, r2, c2);
      n->value = apply_delta(n->value, n->delta, cur_area);
      if (cur_area > 1) {
        int rlen = r2 - r1 + 1, clen = c2 - c1 + 1;
        int rlen1 = rmid - r1 + 1, rlen2 = rlen - rlen1;
        int clen1 = cmid - c1 + 1, clen2 = clen - clen1;
        update_delta(n->child[0], n->delta, static_cast<long long>(rlen1) * clen1);
        update_delta(n->child[1], n->delta, static_cast<long long>(rlen2) * clen1);
        update_delta(n->child[2], n->delta, static_cast<long long>(rlen1) * clen2);
        update_delta(n->child[3], n->delta, static_cast<long long>(rlen2) * clen2);
      }
      n->pending = false;
    }
  }

  // Helper variables for query() and update().
  int tgt_r1, tgt_c1, tgt_r2, tgt_c2;
  T res, delta;
  bool found;

  void append_result(const T &v) {
    res = found ? combine(res, v) : v;
    found = true;
  }

  T child_value(Node *n, int r1, int c1, int r2, int c2) {
    return (n != nullptr) ? n->value : repeat_value(init, area(r1, c1, r2, c2));
  }

  void query(Node *n, int r1, int c1, int r2, int c2) {
    if (tgt_r2 < r1 || tgt_r1 > r2 || tgt_c2 < c1 || tgt_c1 > c2) {
      return;
    }
    if (n == nullptr) {
      int rlen = std::min(r2, tgt_r2) - std::max(r1, tgt_r1) + 1;
      int clen = std::min(c2, tgt_c2) - std::max(c1, tgt_c1) + 1;
      append_result(repeat_value(init, static_cast<long long>(rlen) * clen));
      return;
    }
    push_delta(n, r1, c1, r2, c2);
    if (tgt_r1 <= r1 && r2 <= tgt_r2 && tgt_c1 <= c1 && c2 <= tgt_c2) {
      append_result(n->value);
      return;
    }
    int rmid = r1 + (r2 - r1) / 2, cmid = c1 + (c2 - c1) / 2;
    query(n->child[0], r1, c1, rmid, cmid);
    query(n->child[1], rmid + 1, c1, r2, cmid);
    query(n->child[2], r1, cmid + 1, rmid, c2);
    query(n->child[3], rmid + 1, cmid + 1, r2, c2);
  }

  void update(Node *&n, int r1, int c1, int r2, int c2) {
    if (n == nullptr) {
      if (tgt_r2 < r1 || tgt_r1 > r2 || tgt_c2 < c1 || tgt_c1 > c2) {
        return;
      }
      n = new Node(repeat_value(init, area(r1, c1, r2, c2)));
    } else {
      push_delta(n, r1, c1, r2, c2);
    }
    if (tgt_r2 < r1 || tgt_r1 > r2 || tgt_c2 < c1 || tgt_c1 > c2) {
      return;
    }
    if (tgt_r1 <= r1 && r2 <= tgt_r2 && tgt_c1 <= c1 && c2 <= tgt_c2) {
      n->delta = delta;
      n->pending = true;
      push_delta(n, r1, c1, r2, c2);
      return;
    }
    int rmid = r1 + (r2 - r1) / 2, cmid = c1 + (c2 - c1) / 2;
    update(n->child[0], r1, c1, rmid, cmid);
    update(n->child[1], rmid + 1, c1, r2, cmid);
    update(n->child[2], r1, cmid + 1, rmid, c2);
    update(n->child[3], rmid + 1, cmid + 1, r2, c2);
    n->value = combine(
        combine(
            child_value(n->child[0], r1, c1, rmid, cmid),
            child_value(n->child[1], rmid + 1, c1, r2, cmid)
        ),
        combine(
            child_value(n->child[2], r1, cmid + 1, rmid, c2),
            child_value(n->child[3], rmid + 1, cmid + 1, r2, c2)
        )
    );
  }

  static void clean_up(Node *n) {
    if (n != nullptr) {
      for (auto *c : n->child) {
        clean_up(c);
      }
      delete n;
    }
  }

 public:
  explicit Quadtree(const T &v = T()) : root(nullptr), init(v) {}

  ~Quadtree() { clean_up(root); }
  T at(int r, int c) { return query(r, c, r, c); }

  T query(int r1, int c1, int r2, int c2) {
    tgt_r1 = r1;
    tgt_c1 = c1;
    tgt_r2 = r2;
    tgt_c2 = c2;
    found = false;
    query(root, 0, 0, R, C);
    return found ? res : repeat_value(init, area(r1, c1, r2, c2));
  }

  void update(int r, int c, const T &d) { update(r, c, r, c, d); }

  void update(int r1, int c1, int r2, int c2, const T &d) {
    tgt_r1 = r1;
    tgt_c1 = c1;
    tgt_r2 = r2;
    tgt_c2 = c2;
    delta = d;
    update(root, 0, 0, R, C);
  }
};
\end{lstlisting}
\Needspace*{4\baselineskip}
\textbf{Example Usage}\par
\begin{lstlisting}[language=C++]
#include <cassert>
#include <iostream>
using namespace std;

int main() {
  Quadtree<int> t(0);
  t.update(0, 0, 7);
  t.update(0, 1, 6);
  t.update(1, 0, 5);
  t.update(1, 1, 4);
  t.update(2, 1, 1);
  t.update(2, 2, 9);
  cout << "Values:" << endl;
  for (int i = 0; i < 3; i++) {
    for (int j = 0; j < 3; j++) {
      cout << t.at(i, j) << " ";
    }
    cout << endl;
  }
  assert(t.query(0, 0, 0, 1) == 6);
  assert(t.query(0, 0, 1, 0) == 5);
  assert(t.query(1, 1, 2, 2) == 0);
  assert(t.query(0, 0, 1000000000, 1000000000) == 0);
  t.update(500000000, 500000000, -100);
  assert(t.query(0, 0, 1000000000, 1000000000) == -100);
  return 0;
}
\end{lstlisting}
\begin{exampleoutput}
Values:
7 6 0
5 4 0
0 1 9
\end{exampleoutput}
\subsection{2D Segment Tree}
Maintain a two-dimensional array while supporting dynamic queries of rectangular sub-arrays and dynamic updates of individual indices. This implementation uses lazy initialization of nodes to conserve memory while supporting large indices.

The query operation is defined by a commutative associative aggregate function \inlinecode{combine(a, b)}. Because untouched regions are implicit, \inlinecode{repeat\_value(v, area)} must return the aggregate summary of \inlinecode{area} copies of the initial value \inlinecode{v}. The default code below assumes a numerical array type, defining queries for the ``min'' of the target range. For rectangle-sum queries, \inlinecode{combine(a, b)} should return $a + b$ and \inlinecode{repeat\_value(v, area)} should return $v \cdot area$.

The point update operation is defined by \inlinecode{apply\_delta(v, d)}, which returns the new value at one updated cell. The default code below defines updates that ``set'' the chosen cell to a new value. Another possible update operation is ``increment'', in which case \inlinecode{apply\_delta(v, d)} should return $v + d$.

\begin{compactitem}
  \item \inlinecode{SegTree2D(v)} constructs a two-dimensional array with rows from 0 to \inlinecode{R} and columns from 0 to \inlinecode{C}, inclusive. All values are implicitly initialized to \inlinecode{v}.
  \item \inlinecode{at(r, c)} returns the value at row \inlinecode{r}, column \inlinecode{c}.
  \item \inlinecode{query(r1, c1, r2, c2)} returns the result of \inlinecode{combine()} applied to every value in the rectangular region consisting of rows from \inlinecode{r1} to \inlinecode{r2}, inclusive, and columns from \inlinecode{c1} to \inlinecode{c2}, inclusive.
  \item \inlinecode{update(r, c, d)} assigns the value \inlinecode{v} at (\inlinecode{r}, \inlinecode{c}) to \inlinecode{apply\_delta(v, d)}.
\end{compactitem}

\Needspace*{3\baselineskip}
\textbf{Time Complexity}:
\begin{compactitem}
  \item $O(1)$ per call to the constructor.
  \item $O(\log\left(R\right) \cdot \log\left(C\right))$ per call to \inlinecode{at()}, \inlinecode{update()}, and \inlinecode{query()}.
\end{compactitem}

\Needspace*{3\baselineskip}
\textbf{Space Complexity}:
\begin{compactitem}
  \item $O(n \log\left(R\right) \log\left(C\right))$ for storage after $n$ point updates.
  \item $O(\log\left(R\right) + \log\left(C\right))$ auxiliary stack space for \inlinecode{update()}, \inlinecode{query()}, and \inlinecode{at()}.
\end{compactitem}

\begin{lstlisting}[language=C++]
#include <algorithm>
#include <cstddef>

template<class T>
class SegTree2D {
  static const int R = 1000000000;
  static const int C = 1000000000;

  static T combine(const T &a, const T &b) { return std::min(a, b); }
  static T repeat_value(const T &v, long long area) { return v; }
  static T apply_delta(const T &v, const T &d) { return d; }

  struct InnerNode {
    T value;
    int low, high;
    InnerNode *left, *right;

    InnerNode(int lo, int hi, const T &v)
        : value(v), low(lo), high(hi), left(nullptr), right(nullptr) {}
  };

  struct OuterNode {
    InnerNode root;
    int low, high;
    OuterNode *left, *right;

    OuterNode(int lo, int hi, const T &v)
        : root(0, C, v), low(lo), high(hi), left(nullptr), right(nullptr) {}
  } *root;

  T init;

  // Helper variables for query() and update().
  int tgt_r1, tgt_c1, tgt_r2, tgt_c2;
  long long width;

  static long long length(int lo, int hi) { return hi - lo + 1LL; }

  void append_result(T &res, bool &found, const T &v) {
    res = found ? combine(res, v) : v;
    found = true;
  }

  T query(InnerNode *n, int qlo, int qhi, long long rows) {
    if (n == nullptr) {
      return repeat_value(init, rows * length(qlo, qhi));
    }
    int lo = n->low, hi = n->high, mid = lo + (hi - lo) / 2;
    T res;
    bool found = false;
    if (qlo < lo) {
      int missing_hi = std::min(qhi, lo - 1);
      append_result(res, found, repeat_value(init, rows * length(qlo, missing_hi)));
    }
    int overlap_lo = std::max(qlo, lo), overlap_hi = std::min(qhi, hi);
    if (overlap_lo <= overlap_hi) {
      if (overlap_lo == lo && overlap_hi == hi) {
        append_result(res, found, n->value);
      } else {
        if (overlap_lo <= mid) {
          append_result(res, found, query(n->left, overlap_lo, std::min(overlap_hi, mid), rows));
        }
        if (mid < overlap_hi) {
          append_result(
              res, found, query(n->right, std::max(overlap_lo, mid + 1), overlap_hi, rows)
          );
        }
      }
    }
    if (hi < qhi) {
      int missing_lo = std::max(qlo, hi + 1);
      append_result(res, found, repeat_value(init, rows * length(missing_lo, qhi)));
    }
    return res;
  }

  T query(OuterNode *n, int qlo, int qhi) {
    if (n == nullptr) {
      return repeat_value(init, width * length(qlo, qhi));
    }
    int lo = n->low, hi = n->high, mid = lo + (hi - lo) / 2;
    T res;
    bool found = false;
    if (qlo < lo) {
      int missing_hi = std::min(qhi, lo - 1);
      append_result(res, found, repeat_value(init, width * length(qlo, missing_hi)));
    }
    int overlap_lo = std::max(qlo, lo), overlap_hi = std::min(qhi, hi);
    if (overlap_lo <= overlap_hi) {
      if (overlap_lo == lo && overlap_hi == hi) {
        append_result(res, found, query(&(n->root), tgt_c1, tgt_c2, length(lo, hi)));
      } else {
        if (overlap_lo <= mid) {
          append_result(res, found, query(n->left, overlap_lo, std::min(overlap_hi, mid)));
        }
        if (mid < overlap_hi) {
          append_result(res, found, query(n->right, std::max(overlap_lo, mid + 1), overlap_hi));
        }
      }
    }
    if (hi < qhi) {
      int missing_lo = std::max(qlo, hi + 1);
      append_result(res, found, repeat_value(init, width * length(missing_lo, qhi)));
    }
    return res;
  }

  void update(InnerNode *n, int c, const T &d, bool leaf_row, long long rows) {
    int lo = n->low, hi = n->high, mid = lo + (hi - lo) / 2;
    if (lo == hi) {
      if (leaf_row) {
        n->value = apply_delta(n->value, d);
      } else {
        n->value = d;
      }
      return;
    }
    InnerNode *&target = (c <= mid) ? n->left : n->right;
    if (target == nullptr) {
      target = new InnerNode(c, c, repeat_value(init, rows));
    }
    if (target->low <= c && c <= target->high) {
      update(target, c, d, leaf_row, rows);
    } else {
      int split_lo = lo, split_hi = hi, split_mid = mid;
      do {
        if (c <= split_mid) {
          split_hi = split_mid;
        } else {
          split_lo = split_mid + 1;
        }
        split_mid = split_lo + (split_hi - split_lo) / 2;
      } while ((c <= split_mid) == (target->low <= split_mid));
      InnerNode *tmp =
          new InnerNode(split_lo, split_hi, repeat_value(init, rows * length(split_lo, split_hi)));
      if (target->low <= split_mid) {
        tmp->left = target;
      } else {
        tmp->right = target;
      }
      target = tmp;
      update(tmp, c, d, leaf_row, rows);
    }
    T left_value =
        (n->left != nullptr) ? n->left->value : repeat_value(init, rows * length(lo, mid));
    T right_value =
        (n->right != nullptr) ? n->right->value : repeat_value(init, rows * length(mid + 1, hi));
    n->value = combine(left_value, right_value);
  }

  void update(OuterNode *n, int r, int c, const T &d) {
    int lo = n->low, hi = n->high, mid = lo + (hi - lo) / 2;
    long long rows = length(lo, hi);
    if (lo == hi) {
      update(&(n->root), c, d, true, 1);
      return;
    }
    if (r <= mid) {
      if (n->left == nullptr) {
        n->left = new OuterNode(lo, mid, repeat_value(init, length(lo, mid) * length(0, C)));
      }
      update(n->left, r, c, d);
    } else {
      if (n->right == nullptr) {
        n->right =
            new OuterNode(mid + 1, hi, repeat_value(init, length(mid + 1, hi) * length(0, C)));
      }
      update(n->right, r, c, d);
    }
    T value = repeat_value(init, rows);
    if (n->left != nullptr || n->right != nullptr) {
      tgt_c1 = tgt_c2 = c;
      T left_value = (n->left != nullptr) ? query(&(n->left->root), c, c, length(lo, mid))
                                          : repeat_value(init, length(lo, mid));
      T right_value = (n->right != nullptr) ? query(&(n->right->root), c, c, length(mid + 1, hi))
                                            : repeat_value(init, length(mid + 1, hi));
      value = combine(left_value, right_value);
    }
    update(&(n->root), c, value, false, rows);
  }

  static void clean_up(InnerNode *n) {
    if (n != nullptr) {
      clean_up(n->left);
      clean_up(n->right);
      delete n;
    }
  }

  static void clean_up(OuterNode *n) {
    if (n != nullptr) {
      clean_up(n->root.left);
      clean_up(n->root.right);
      clean_up(n->left);
      clean_up(n->right);
      delete n;
    }
  }

 public:
  explicit SegTree2D(const T &v = T())
      : root(new OuterNode(0, R, repeat_value(v, length(0, R) * length(0, C)))), init(v) {}

  ~SegTree2D() { clean_up(root); }
  T at(int r, int c) { return query(r, c, r, c); }

  T query(int r1, int c1, int r2, int c2) {
    tgt_r1 = r1;
    tgt_c1 = c1;
    tgt_r2 = r2;
    tgt_c2 = c2;
    width = length(c1, c2);
    return query(root, r1, r2);
  }

  void update(int r, int c, const T &d) { update(root, r, c, d); }
};
\end{lstlisting}
\Needspace*{4\baselineskip}
\textbf{Example Usage}\par
\begin{lstlisting}[language=C++]
#include <cassert>
#include <iostream>
using namespace std;

int main() {
  SegTree2D<int> t(0);
  t.update(0, 0, 7);
  t.update(0, 1, 6);
  t.update(1, 0, 5);
  t.update(1, 1, 4);
  t.update(2, 1, 1);
  t.update(2, 2, 9);
  cout << "Values:" << endl;
  for (int i = 0; i < 3; i++) {
    for (int j = 0; j < 3; j++) {
      cout << t.at(i, j) << " ";
    }
    cout << endl;
  }
  assert(t.query(0, 0, 0, 1) == 6);
  assert(t.query(0, 0, 1, 0) == 5);
  assert(t.query(1, 1, 2, 2) == 0);
  assert(t.query(0, 0, 1000000000, 1000000000) == 0);
  t.update(500000000, 500000000, -100);
  assert(t.query(0, 0, 1000000000, 1000000000) == -100);
  return 0;
}
\end{lstlisting}
\begin{exampleoutput}
Values:
7 6 0
5 4 0
0 1 9
\end{exampleoutput}
\subsection{2D Range Tree}
Maintain a set of two-dimensional points while supporting queries for all points that fall inside given rectangular regions. This implementation uses \inlinecode{std::pair} to represent points, requiring operators \inlinecode{\textless{}} and \inlinecode{==} to be defined on the numeric template type.

\begin{compactitem}
  \item \inlinecode{RangeTree(lo, hi)} constructs a set from two random-access iterators to \inlinecode{std::pair} as a range \inlinecode{[lo, hi)} of points.
  \item \inlinecode{query(x1, y1, x2, y2, f)} calls the function \inlinecode{f(i, p)} on each point in the set that falls into the rectangular region consisting of rows from \inlinecode{x1} to \inlinecode{x2}, inclusive, and columns from \inlinecode{y1} to \inlinecode{y2}, inclusive. The first argument to \inlinecode{f} is the zero-based index of the point in the original range given to the constructor. The second argument is the point itself as an \inlinecode{std::pair}.
\end{compactitem}

\Needspace*{3\baselineskip}
\textbf{Time Complexity}:
\begin{compactitem}
  \item $O(n \log n)$ per call to the constructor, where $n$ is the number of points.
  \item $O(\log^{2}(n) + m)$ per call to \inlinecode{query()}, where $m$ is the number of points that are reported by the query.
\end{compactitem}

\Needspace*{3\baselineskip}
\textbf{Space Complexity}:
\begin{compactitem}
  \item $O(n \log n)$ for storage of the points.
  \item $O(\log^{2}(n))$ auxiliary stack space for \inlinecode{query()}.
\end{compactitem}

\begin{lstlisting}[language=C++]
#include <algorithm>
#include <iterator>
#include <utility>
#include <vector>

template<class T>
class RangeTree {
  using point = std::pair<T, T>;
  using colindex = std::pair<int, T>;

  std::vector<point> points;
  std::vector<std::vector<colindex>> columns;

  static inline bool comp1(const colindex &a, const colindex &b) { return a.second < b.second; }
  static inline bool comp2(const colindex &a, const T &v) { return a.second < v; }

  void build(int n, int lo, int hi) {
    if (points[lo].first == points[hi].first) {
      for (int i = lo; i <= hi; i++) {
        columns[n].emplace_back(i, points[i].second);
      }
      return;
    }
    int l = n * 2 + 1, r = n * 2 + 2, mid = lo + (hi - lo) / 2;
    build(l, lo, mid);
    build(r, mid + 1, hi);
    columns[n].resize(columns[l].size() + columns[r].size());
    std::merge(
        columns[l].begin(), columns[l].end(), columns[r].begin(), columns[r].end(),
        columns[n].begin(), comp1
    );
  }

  // Helper variables for query().
  T x1, y1, x2, y2;

  template<class Fn>
  void query(int n, int lo, int hi, Fn f) {
    if (points[hi].first < x1 || x2 < points[lo].first) {
      return;
    }
    if (!(points[lo].first < x1 || x2 < points[hi].first)) {
      if (!columns[n].empty() && !(y2 < y1)) {
        auto it = std::lower_bound(columns[n].begin(), columns[n].end(), y1, comp2);
        for (; it != columns[n].end() && it->second <= y2; ++it) {
          f(it->first, points[it->first]);
        }
      }
    } else if (lo != hi) {
      int mid = lo + (hi - lo) / 2;
      query(n * 2 + 1, lo, mid, f);
      query(n * 2 + 2, mid + 1, hi, f);
    }
  }

 public:
  template<class It>
  RangeTree(It lo, It hi) : points(lo, hi) {
    int n = std::distance(lo, hi);
    columns.resize(4 * n + 1);
    std::sort(points.begin(), points.end());
    build(0, 0, n - 1);
  }

  template<class Fn>
  void query(const T &x1, const T &y1, const T &x2, const T &y2, Fn f) {
    this->x1 = x1;
    this->y1 = y1;
    this->x2 = x2;
    this->y2 = y2;
    query(0, 0, points.size() - 1, f);
  }
};
\end{lstlisting}
\Needspace*{4\baselineskip}
\textbf{Example Usage}\par
\begin{lstlisting}[language=C++]
#include <iostream>
using namespace std;

void print(int i, const pair<int, int> &p) {
  cout << "(" << p.first << ", " << p.second << ") ";
}

int main() {
  vector<pair<int, int>> v{{1, 4},  {5, 4},  {2, 2},   {3, 1},   {6, -5},
                           {5, -1}, {3, -3}, {-1, -2}, {-1, -1}, {2, -1}};
  RangeTree<int> t(v.begin(), v.end());
  t.query(-1, -1, 2, 5, print);
  cout << endl;
  t.query(1, 1, 4, 8, print);
  cout << endl;
  return 0;
}
\end{lstlisting}
\begin{exampleoutput}
(-1, -1) (2, -1) (2, 2) (1, 4)
(1, 4) (2, 2) (3, 1)
\end{exampleoutput}
\subsection{K-d Tree (2D Range Query)}
Maintain a set of two-dimensional points while supporting queries for all points that fall inside given rectangular regions. This implementation uses \inlinecode{std::pair} to represent points, requiring operators \inlinecode{\textless{}} and \inlinecode{==} to be defined on the numeric template type.

\begin{compactitem}
  \item \inlinecode{KDTree(lo, hi)} constructs a set from two random-access iterators to \inlinecode{std::pair} as a range \inlinecode{[lo, hi)} of points.
  \item \inlinecode{query(x1, y1, x2, y2, f)} calls the function \inlinecode{f(i, p)} on each point in the set that falls into the rectangular region consisting of rows from \inlinecode{x1} to \inlinecode{x2}, inclusive, and columns from \inlinecode{y1} to \inlinecode{y2}, inclusive. The first argument to \inlinecode{f} is the zero-based index of the point in the original range given to the constructor. The second argument is the point itself as an \inlinecode{std::pair}.
\end{compactitem}

\Needspace*{3\baselineskip}
\textbf{Time Complexity}:
\begin{compactitem}
  \item $O(n \log n)$ per call to the constructor, where $n$ is the number of points.
  \item $O(\log\left(n\right) + m)$ on average per call to \inlinecode{query()}, where $m$ is the number of points that are reported by the query.
\end{compactitem}

\Needspace*{3\baselineskip}
\textbf{Space Complexity}:
\begin{compactitem}
  \item $O(n)$ for storage of the points.
  \item $O(\log n)$ auxiliary stack space for \inlinecode{query()}.
\end{compactitem}

\begin{lstlisting}[language=C++]
#include <algorithm>
#include <utility>
#include <vector>

template<class T>
class KDTree {
  using point = std::pair<T, T>;

  static inline bool comp1(const point &a, const point &b) { return a.first < b.first; }
  static inline bool comp2(const point &a, const point &b) { return a.second < b.second; }

  std::vector<point> tree, minp, maxp;
  std::vector<int> l_index, h_index;

  void build(int lo, int hi, bool div_x) {
    if (lo >= hi) {
      return;
    }
    int mid = lo + (hi - lo) / 2;
    std::nth_element(
        tree.begin() + lo, tree.begin() + mid, tree.begin() + hi, div_x ? comp1 : comp2
    );
    l_index[mid] = lo;
    h_index[mid] = hi;
    minp[mid].first = maxp[mid].first = tree[lo].first;
    minp[mid].second = maxp[mid].second = tree[lo].second;
    for (int i = lo + 1; i < hi; i++) {
      minp[mid].first = std::min(minp[mid].first, tree[i].first);
      minp[mid].second = std::min(minp[mid].second, tree[i].second);
      maxp[mid].first = std::max(maxp[mid].first, tree[i].first);
      maxp[mid].second = std::max(maxp[mid].second, tree[i].second);
    }
    build(lo, mid, !div_x);
    build(mid + 1, hi, !div_x);
  }

  // Helper variables for query().
  T x1, y1, x2, y2;

  template<class Fn>
  void query(int lo, int hi, Fn f) {
    if (lo >= hi) {
      return;
    }
    int mid = lo + (hi - lo) / 2;
    T ax = minp[mid].first, ay = minp[mid].second;
    T bx = maxp[mid].first, by = maxp[mid].second;
    if (x2 < ax || bx < x1 || y2 < ay || by < y1) {
      return;
    }
    if (!(ax < x1 || x2 < bx || ay < y1 || y2 < by)) {
      for (int i = l_index[mid]; i < h_index[mid]; i++) {
        f(tree[i]);
      }
      return;
    }
    query(lo, mid, f);
    query(mid + 1, hi, f);
    if (tree[mid].first < x1 || x2 < tree[mid].first || tree[mid].second < y1 ||
        y2 < tree[mid].second) {
      return;
    }
    f(tree[mid]);
  }

 public:
  template<class It>
  KDTree(It lo, It hi) : tree(lo, hi) {
    int n = std::distance(lo, hi);
    l_index.resize(n);
    h_index.resize(n);
    minp.resize(n);
    maxp.resize(n);
    build(0, n, true);
  }

  template<class Fn>
  void query(const T &x1, const T &y1, const T &x2, const T &y2, Fn f) {
    this->x1 = x1;
    this->y1 = y1;
    this->x2 = x2;
    this->y2 = y2;
    query(0, tree.size(), f);
  }
};
\end{lstlisting}
\Needspace*{4\baselineskip}
\textbf{Example Usage}\par
\begin{lstlisting}[language=C++]
#include <iostream>
using namespace std;

void print(const pair<int, int> &p) {
  cout << "(" << p.first << ", " << p.second << ") ";
}

int main() {
  vector<pair<int, int>> v{{1, 4},  {5, 4},  {2, 2},   {3, 1},   {6, -5},
                           {5, -1}, {3, -3}, {-1, -2}, {-1, -1}, {2, -1}};
  KDTree<int> t(v.begin(), v.end());
  t.query(-1, -1, 2, 5, print);
  cout << endl;
  t.query(1, 1, 4, 8, print);
  cout << endl;
  return 0;
}
\end{lstlisting}
\begin{exampleoutput}
(2, -1) (1, 4) (2, 2) (-1, -1)
(1, 4) (2, 2) (3, 1)
\end{exampleoutput}
\subsection{K-d Tree (Nearest Neighbor)}
Maintain a set of two-dimensional points while supporting queries for the closest point in the set to a given query point. This implementation uses \inlinecode{std::pair} to represent points, requiring operators \inlinecode{\textless{}}, \inlinecode{==}, \inlinecode{-}, and \inlinecode{long double} casting to be defined on the numeric template type.

\begin{compactitem}
  \item \inlinecode{KDTree(lo, hi)} constructs a set from two random-access iterators to \inlinecode{std::pair} as a range \inlinecode{[lo, hi)} of points.
  \item \inlinecode{nearest(x, y, can\_equal)} returns a point in the set that is closest to (\inlinecode{x}, \inlinecode{y}) by Euclidean distance. This may be equal to (\inlinecode{x}, \inlinecode{y}) only if \inlinecode{can\_equal} is \inlinecode{true}.
\end{compactitem}

\Needspace*{3\baselineskip}
\textbf{Time Complexity}:
\begin{compactitem}
  \item $O(n \log n)$ per call to the constructor, where $n$ is the number of points.
  \item $O(\log n)$ on average per call to \inlinecode{nearest()}.
\end{compactitem}

\Needspace*{3\baselineskip}
\textbf{Space Complexity}:
\begin{compactitem}
  \item $O(n)$ for storage of the points.
  \item $O(\log n)$ auxiliary stack space for \inlinecode{nearest()}.
\end{compactitem}

\begin{lstlisting}[language=C++]
#include <algorithm>
#include <cfloat>
#include <stdexcept>
#include <utility>
#include <vector>

template<class T>
class KDTree {
  using point = std::pair<T, T>;

  static inline bool comp1(const point &a, const point &b) { return a.first < b.first; }
  static inline bool comp2(const point &a, const point &b) { return a.second < b.second; }

  std::vector<point> tree;
  std::vector<bool> div_x;

  void build(int lo, int hi) {
    if (lo >= hi) {
      return;
    }
    int mid = lo + (hi - lo) / 2;
    T minx, maxx, miny, maxy;
    minx = maxx = tree[lo].first;
    miny = maxy = tree[lo].second;
    for (int i = lo + 1; i < hi; i++) {
      minx = std::min(minx, tree[i].first);
      miny = std::min(miny, tree[i].second);
      maxx = std::max(maxx, tree[i].first);
      maxy = std::max(maxy, tree[i].second);
    }
    div_x[mid] = !((maxx - minx) < (maxy - miny));
    std::nth_element(
        tree.begin() + lo, tree.begin() + mid, tree.begin() + hi, div_x[mid] ? comp1 : comp2
    );
    if (lo + 1 == hi) {
      return;
    }
    build(lo, mid);
    build(mid + 1, hi);
  }

  // Helper variables for nearest().
  long double min_dist;
  int id;

  void nearest(int lo, int hi, const T &x, const T &y, bool can_equal) {
    if (lo >= hi) {
      return;
    }
    int mid = lo + (hi - lo) / 2;
    T dx = x - tree[mid].first, dy = y - tree[mid].second;
    long double d = dx * static_cast<long double>(dx) + dy * static_cast<long double>(dy);
    if (d < min_dist && (can_equal || d != 0)) {
      min_dist = d;
      id = mid;
    }
    if (lo + 1 == hi) {
      return;
    }
    d = static_cast<long double>((div_x[mid] ? dx : dy));
    int l1 = lo, r1 = mid, l2 = mid + 1, r2 = hi;
    if (d > 0) {
      std::swap(l1, l2);
      std::swap(r1, r2);
    }
    nearest(l1, r1, x, y, can_equal);
    if (d * static_cast<long double>(d) < min_dist) {
      nearest(l2, r2, x, y, can_equal);
    }
  }

 public:
  template<class It>
  KDTree(It lo, It hi) : tree(lo, hi) {
    int n = std::distance(lo, hi);
    if (n <= 1) {
      throw std::runtime_error("K-d tree must be have at least 2 points.");
    }
    div_x.resize(n);
    build(0, n);
  }

  point nearest(const T &x, const T &y, bool can_equal = true) {
    min_dist = LDBL_MAX;
    nearest(0, tree.size(), x, y, can_equal);
    return tree[id];
  }
};
\end{lstlisting}
\Needspace*{4\baselineskip}
\textbf{Example Usage}\par
\begin{lstlisting}[language=C++]
#include <cassert>
using namespace std;

int main() {
  vector<pair<int, int>> p{{0, 2}, {0, 3}, {-1, 0}};
  KDTree<int> t(p.begin(), p.end());
  assert(t.nearest(0, 2, true) == (pair<int, int>{0, 2}));
  assert(t.nearest(0, 2, false) == (pair<int, int>{0, 3}));
  assert(t.nearest(0, 0) == (pair<int, int>{-1, 0}));
  assert(t.nearest(-10000, 0) == (pair<int, int>{-1, 0}));
  return 0;
}
\end{lstlisting}

\section{Fenwick Trees}
\setcounter{section}{6}
\setcounter{subsection}{0}
\subsection{Fenwick Tree (Simple)}
Maintain an array of numerical type, allowing for updates of individual indices (point update) and queries for the sum of contiguous sub-arrays (range queries). This implementation assumes that the array is 1-based (i.e. has valid indices from 1 to \inlinecode{n}, inclusive).

\begin{compactitem}
  \item \inlinecode{initialize(n)} resets the dat structure.
  \item \inlinecode{vals[i]} stores the value at index \inlinecode{i}.
  \item \inlinecode{add(i, x)} adds \inlinecode{x} to the value at index \inlinecode{i}.
  \item \inlinecode{set(i, x)} assigns the value at index \inlinecode{i} to \inlinecode{x}.
  \item \inlinecode{sum(hi)} returns the sum of all values at indices from 1 to \inlinecode{hi}, inclusive.
  \item \inlinecode{sum(lo, hi)} returns the sum of all values at indices from \inlinecode{lo} to \inlinecode{hi}, inclusive.
\end{compactitem}

\Needspace*{3\baselineskip}
\textbf{Time Complexity}:
\begin{compactitem}
  \item $O(n)$ per call to \inlinecode{initialize(n)}, where $n$ is the size of the array.
  \item $O(\log n)$ per call to all other operations.
\end{compactitem}

\Needspace*{3\baselineskip}
\textbf{Space Complexity}:
\begin{compactitem}
  \item $O(n)$ for storage of the array elements.
  \item $O(1)$ auxiliary for all operations.
\end{compactitem}

\begin{lstlisting}[language=C++]
#include <vector>

std::vector<int> vals, tree;

void initialize(int n) {
  vals.assign(n + 2, 0);
  tree.assign(n + 2, 0);
}

void add(int i, int x) {
  vals[i] += x;
  for (; i < static_cast<int>(tree.size()); i += i & -i) {
    tree[i] += x;
  }
}

void set(int i, int x) {
  add(i, x - vals[i]);
}

int sum(int hi) {
  int res = 0;
  for (; hi > 0; hi -= hi & -hi) {
    res += tree[hi];
  }
  return res;
}

int sum(int lo, int hi) {
  return sum(hi) - sum(lo - 1);
}
\end{lstlisting}
\Needspace*{4\baselineskip}
\textbf{Example Usage}\par
\begin{lstlisting}[language=C++]
#include <cassert>
#include <iostream>
using namespace std;

int main() {
  vector<int> v{10, 1, 2, 3, 4};
  int n = 5;
  initialize(n);
  for (int i = 1; i < n; i++) {
    set(i, v[i - 1]);
  }
  add(1, -5);
  cout << "Values: ";
  for (int i = 1; i < n; i++) {
    cout << vals[i] << " ";
  }
  cout << endl;
  assert(sum(2, 4) == 6);
  return 0;
}
\end{lstlisting}
\begin{exampleoutput}
Values: 5 1 2 3 4
\end{exampleoutput}
\subsection{Fenwick Tree (Range Update, Point Query)}
Maintain an array of numerical type, allowing for contiguous sub-arrays to be simultaneously incremented by arbitrary values (range update) and values at individual indices to be queried (point query). This implementation assumes that the array is 0-based (i.e. has valid indices from 0 to \inlinecode{size() - 1}, inclusive).

\begin{compactitem}
  \item \inlinecode{size()} returns the size of the array.
  \item \inlinecode{at(i)} returns the value at index \inlinecode{i}.
  \item \inlinecode{add(i, x)} adds \inlinecode{x} to the value at index \inlinecode{i}.
  \item \inlinecode{add(lo, hi, x)} adds \inlinecode{x} to the values at all indices from \inlinecode{lo} to \inlinecode{hi}, inclusive.
\end{compactitem}

\Needspace*{3\baselineskip}
\textbf{Time Complexity}:
\begin{compactitem}
  \item $O(n)$ per call to the constructor, where $n$ is the size of the array.
  \item $O(1)$ per call to \inlinecode{size()}.
  \item $O(\log n)$ per call to \inlinecode{at()} and both \inlinecode{add()} functions.
\end{compactitem}

\Needspace*{3\baselineskip}
\textbf{Space Complexity}:
\begin{compactitem}
  \item $O(n)$ for storage of the array elements.
  \item $O(1)$ auxiliary for all operations.
\end{compactitem}

\begin{lstlisting}[language=C++]
#include <vector>

template<class T>
class FenwickTree {
  int len;
  std::vector<T> tree;

 public:
  explicit FenwickTree(int n) : len(n), tree(n + 2) {}

  int size() const { return len; }

  T at(int i) const {
    T res = 0;
    for (i++; i > 0; i -= i & -i) {
      res += tree[i];
    }
    return res;
  }

  void add(int i, const T &x) {
    for (i++; i <= len + 1; i += i & -i) {
      tree[i] += x;
    }
  }

  void add(int lo, int hi, const T &x) {
    add(lo, x);
    add(hi + 1, -x);
  }
};
\end{lstlisting}
\Needspace*{4\baselineskip}
\textbf{Example Usage}\par
\begin{lstlisting}[language=C++]
#include <iostream>
using namespace std;

int main() {
  FenwickTree<int> t(5);
  t.add(0, 1, 5);
  t.add(1, 2, 5);
  t.add(2, 4, 10);
  cout << "Values: ";
  for (int i = 0; i < t.size(); i++) {
    cout << t.at(i) << " ";
  }
  cout << endl;
  return 0;
}
\end{lstlisting}
\begin{exampleoutput}
Values: 5 10 15 10 10
\end{exampleoutput}
\subsection{Fenwick Tree (Point Update, Range Query)}
Maintain an array of numerical type, allowing for updates of individual indices (point update) and queries for the sum of contiguous sub-arrays (range queries). This implementation assumes that the array is 0-based (i.e. has valid indices from 0 to \inlinecode{size() - 1}, inclusive).

\begin{compactitem}
  \item \inlinecode{size()} returns the size of the array.
  \item \inlinecode{at(i)} returns the value at index \inlinecode{i}.
  \item \inlinecode{add(i, x)} adds \inlinecode{x} to the value at index \inlinecode{i}.
  \item \inlinecode{set(i, x)} assigns the value at index \inlinecode{i} to \inlinecode{x}.
  \item \inlinecode{sum(hi)} returns the sum of all values at indices from 0 to \inlinecode{hi}, inclusive.
  \item \inlinecode{sum(lo, hi)} returns the sum of all values at indices from \inlinecode{lo} to \inlinecode{hi}, inclusive.
\end{compactitem}

\Needspace*{3\baselineskip}
\textbf{Time Complexity}:
\begin{compactitem}
  \item $O(n)$ per call to the constructor, where $n$ is the size of the array.
  \item $O(1)$ per call to \inlinecode{size()} and \inlinecode{at()}.
  \item $O(\log n)$ per call to \inlinecode{add()}, \inlinecode{set()}, and both \inlinecode{sum()} functions.
\end{compactitem}

\Needspace*{3\baselineskip}
\textbf{Space Complexity}:
\begin{compactitem}
  \item $O(n)$ for storage of the array elements.
  \item $O(1)$ auxiliary for all operations.
\end{compactitem}

\begin{lstlisting}[language=C++]
#include <vector>

template<class T>
class FenwickTree {
  int len;
  std::vector<T> vals, tree;

 public:
  explicit FenwickTree(int n) : len(n), vals(n + 1), tree(n + 1) {}

  int size() const { return len; }
  T at(int i) const { return vals[i + 1]; }

  void add(int i, const T &x) {
    vals[++i] += x;
    for (; i <= len; i += i & -i) {
      tree[i] += x;
    }
  }

  void set(int i, const T &x) {
    T inc = x - vals[i + 1];
    add(i, inc);
  }

  T sum(int hi) {
    T res = 0;
    for (hi++; hi > 0; hi -= hi & -hi) {
      res += tree[hi];
    }
    return res;
  }

  T sum(int lo, int hi) { return sum(hi) - sum(lo - 1); }
};
\end{lstlisting}
\Needspace*{4\baselineskip}
\textbf{Example Usage}\par
\begin{lstlisting}[language=C++]
#include <cassert>
#include <iostream>
using namespace std;

int main() {
  vector<int> a{10, 1, 2, 3, 4};
  FenwickTree<int> t(5);
  for (int i = 0; i < static_cast<int>(a.size()); i++) {
    t.set(i, a[i]);
  }
  t.add(0, -5);
  cout << "Values: ";
  for (int i = 0; i < t.size(); i++) {
    cout << t.at(i) << " ";
  }
  cout << endl;
  assert(t.sum(1, 3) == 6);
  return 0;
}
\end{lstlisting}
\begin{exampleoutput}
Values: 5 1 2 3 4
\end{exampleoutput}
\subsection{Fenwick Tree (Range Update, Range Query)}
Maintain an array of numerical type, allowing for contiguous sub-arrays to be simultaneously incremented by arbitrary values (range update) and queries for the sum of contiguous sub-arrays (range query). This implementation assumes that the array is 0-based (i.e. has valid indices from 0 to \inlinecode{size() - 1}, inclusive).

\begin{compactitem}
  \item \inlinecode{size()} returns the size of the array.
  \item \inlinecode{at(i)} returns the value at index \inlinecode{i}.
  \item \inlinecode{add(i, x)} increments the value at index \inlinecode{i} by \inlinecode{x}.
  \item \inlinecode{add(lo, hi, x)} adds \inlinecode{x} to the values at all indices from \inlinecode{lo} to \inlinecode{hi}, inclusive.
  \item \inlinecode{set(i, x)} assigns the value at index \inlinecode{i} to \inlinecode{x}.
  \item \inlinecode{sum(hi)} returns the sum of all values at indices from 0 to \inlinecode{hi}, inclusive.
  \item \inlinecode{sum(lo, hi)} returns the sum of all values at indices from \inlinecode{lo} to \inlinecode{hi}, inclusive.
\end{compactitem}

\Needspace*{3\baselineskip}
\textbf{Time Complexity}:
\begin{compactitem}
  \item $O(n)$ per call to the constructor, where $n$ is the size of the array.
  \item $O(1)$ per call to \inlinecode{size()}.
  \item $O(\log n)$ per call to all other operations.
\end{compactitem}

\Needspace*{3\baselineskip}
\textbf{Space Complexity}:
\begin{compactitem}
  \item $O(n)$ for storage of the array elements.
  \item $O(1)$ auxiliary for all operations.
\end{compactitem}

\begin{lstlisting}[language=C++]
#include <vector>

template<class T>
class FenwickTree {
  int len;
  std::vector<T> t1, t2;

  T sum(const std::vector<T> &tree, int i) {
    T res = 0;
    for (; i != 0; i -= i & -i) {
      res += tree[i];
    }
    return res;
  }

  void add(std::vector<T> &tree, int i, const T &x) {
    for (; i <= len + 1; i += i & -i) {
      tree[i] += x;
    }
  }

 public:
  explicit FenwickTree(int n) : len(n), t1(n + 2), t2(n + 2) {}

  int size() const { return len; }

  void add(int lo, int hi, const T &x) {
    lo++;
    hi++;
    add(t1, lo, x);
    add(t1, hi + 1, -x);
    add(t2, lo, x * (lo - 1));
    add(t2, hi + 1, -x * hi);
  }

  void add(int i, const T &x) { return add(i, i, x); }
  void set(int i, const T &x) { add(i, x - at(i)); }

  T sum(int hi) {
    hi++;
    return hi * sum(t1, hi) - sum(t2, hi);
  }

  T sum(int lo, int hi) { return sum(hi) - sum(lo - 1); }
  T at(int i) { return sum(i, i); }
};
\end{lstlisting}
\Needspace*{4\baselineskip}
\textbf{Example Usage}\par
\begin{lstlisting}[language=C++]
#include <cassert>
#include <iostream>
using namespace std;

int main() {
  vector<int> a{10, 1, 2, 3, 4};
  FenwickTree<int> t(5);
  for (int i = 0; i < t.size(); i++) {
    t.set(i, a[i]);
  }
  t.add(0, 2, 5);
  t.set(3, -5);
  cout << "Values: ";
  for (int i = 0; i < t.size(); i++) {
    cout << t.at(i) << " ";
  }
  cout << endl;
  assert(t.sum(0, 4) == 27);
  return 0;
}
\end{lstlisting}
\begin{exampleoutput}
Values: 15 6 7 -5 4
\end{exampleoutput}
\subsection{Fenwick Tree (Compressed)}
Maintain an array of numerical type, allowing for contiguous sub-arrays to be simultaneously incremented by arbitrary values (range update) and queries for the sum of contiguous sub-arrays (range query). This implementation uses \inlinecode{std::unordered\_map} for coordinate compression, allowing for large indices to be accessed with efficient space complexity. That is, all array indices from 0 to \inlinecode{N}, inclusive, are accessible.

\begin{compactitem}
  \item \inlinecode{at(i)} returns the value at index \inlinecode{i}.
  \item \inlinecode{add(i, x)} adds \inlinecode{x} to the value at index \inlinecode{i}.
  \item \inlinecode{add(lo, hi, x)} adds \inlinecode{x} to the values at all indices from \inlinecode{lo} to \inlinecode{hi}, inclusive.
  \item \inlinecode{set(i, x)} assigns the value at index \inlinecode{i} to \inlinecode{x}.
  \item \inlinecode{sum(hi)} returns the sum of all values at indices from 0 to \inlinecode{hi}, inclusive.
  \item \inlinecode{sum(lo, hi)} returns the sum of all values at indices from \inlinecode{lo} to \inlinecode{hi}, inclusive.
\end{compactitem}

\textbf{Time Complexity}: $O(\log n)$ per call to all member functions, where $n$ is the number of distinct indices that have been accessed.

\Needspace*{3\baselineskip}
\textbf{Space Complexity}:
\begin{compactitem}
  \item $O(n \log n)$ for storage of the array elements, where $n$ is the number of distinct indices that have been accessed across all of the operations so far.
  \item $O(1)$ auxiliary for all operations.
\end{compactitem}

\begin{lstlisting}[language=C++]
#include <unordered_map>

template<class T>
class FenwickTree {
  static const int N = 1000000001;
  std::unordered_map<int, T> tmul, tadd;

  void add_helper(int at, int mul, T add) {
    for (int i = at; i <= N; i |= i + 1) {
      tmul[i] += mul;
      tadd[i] += add;
    }
  }

 public:
  void add(int lo, int hi, const T &x) {
    add_helper(lo, x, -x * (lo - 1));
    add_helper(hi, -x, x * hi);
  }

  void add(int i, const T &x) { add(i, i, x); }
  void set(int i, const T &x) { add(i, x - at(i)); }

  T sum(int hi) {
    T mul = 0, add = 0;
    for (int i = hi; i >= 0; i = (i & (i + 1)) - 1) {
      if (auto it = tmul.find(i); it != tmul.end()) {
        mul += it->second;
      }
      if (auto it = tadd.find(i); it != tadd.end()) {
        add += it->second;
      }
    }
    return mul * hi + add;
  }

  T sum(int lo, int hi) { return sum(hi) - sum(lo - 1); }
  T at(int i) { return sum(i, i); }
};
\end{lstlisting}
\Needspace*{4\baselineskip}
\textbf{Example Usage}\par
\begin{lstlisting}[language=C++]
#include <cassert>
#include <iostream>
#include <vector>
using namespace std;

int main() {
  vector<int> a{10, 1, 2, 3, 4};
  FenwickTree<int> t;
  for (int i = 0; i < static_cast<int>(a.size()); i++) {
    t.set(i, a[i]);
  }
  t.add(0, 2, 5);
  t.set(3, -5);
  cout << "Values: ";
  for (int i = 0; i < 5; i++) {
    cout << t.at(i) << " ";
  }
  cout << endl;
  assert(t.sum(0, 4) == 27);
  t.add(500000001, 500000010, 3);
  t.add(500000011, 500000015, 5);
  t.set(500000000, 10);
  assert(t.sum(0, 1000000000) == 92);
  return 0;
}
\end{lstlisting}
\begin{exampleoutput}
Values: 15 6 7 -5 4
\end{exampleoutput}
\subsection{Fenwick Tree (Order Statistic)}
Maintains frequencies over a 1-based integer domain and finds order statistics by binary lifting on a Fenwick tree. This is useful for dynamic multisets of bounded integer values, coordinate-compressed kth-element queries, and online rank queries.

\begin{compactitem}
  \item \inlinecode{initialize(n)} resets all frequencies.
  \item \inlinecode{add(i, delta)} adds \inlinecode{delta} to the frequency of value/index \inlinecode{i}.
  \item \inlinecode{sum(i)} returns the total frequency over indices \inlinecode{[1, i]}.
  \item \inlinecode{kth(k)} returns the smallest index \inlinecode{i} such that \inlinecode{sum(i) \textgreater{}= k}. The argument \inlinecode{k} is 1-based and must satisfy \inlinecode{1 \textless{}= k \textless{}= sum(n)}.
\end{compactitem}

\Needspace*{3\baselineskip}
\textbf{Time Complexity}:
\begin{compactitem}
  \item $O(n)$ per call to \inlinecode{initialize()}.
  \item $O(\log n)$ per call to \inlinecode{add(i, delta)}, \inlinecode{sum(i)}, and \inlinecode{kth(k)}.
\end{compactitem}

\textbf{Space Complexity}: $O(n)$ for the Fenwick tree.

\begin{lstlisting}[language=C++]
#include <algorithm>
#include <vector>

std::vector<int> tree;

void initialize(int n) {
  tree.assign(n + 1, 0);
}

void add(int i, int delta) {
  for (; i < tree.size(); i += i & -i) {
    tree[i] += delta;
  }
}

int sum(int i) {
  int res = 0;
  for (; i > 0; i -= i & -i) {
    res += tree[i];
  }
  return res;
}

int kth(int k) {
  int idx = 0, bits = 1;
  while (bits * 2 < tree.size()) {
    bits *= 2;
  }
  for (; bits > 0; bits >>= 1) {
    int next = idx + bits;
    if (next < tree.size() && tree[next] < k) {
      idx = next;
      k -= tree[next];
    }
  }
  return idx + 1;
}
\end{lstlisting}
\Needspace*{4\baselineskip}
\textbf{Example Usage}\par
\begin{lstlisting}[language=C++]
#include <cassert>

int main() {
  initialize(100);
  add(2, 1);
  add(4, 3);
  add(7, 1);
  assert(sum(4) == 4);
  assert(kth(1) == 2);
  assert(kth(2) == 4);
  assert(kth(4) == 4);
  assert(kth(5) == 7);
  return 0;
}
\end{lstlisting}
\subsection{2D Fenwick Tree (Simple)}
Maintain a 2D array of numerical type, allowing for updates of individual cells in the matrix (point update) and queries for the sum of rectangular sub-matrices (range query). This implementation assumes that array dimensions are 1-based (i.e. rows have valid indices from 1 to \inlinecode{R}, inclusive, and columns have valid indices from 1 to \inlinecode{C}, inclusive).

\begin{compactitem}
  \item \inlinecode{initialize(R, C)} resets the data structure.
  \item \inlinecode{vals[r][c]} stores the value at index (\inlinecode{r}, \inlinecode{c}).
  \item \inlinecode{add(r, c, x)} adds \inlinecode{x} to the value at index (\inlinecode{r}, \inlinecode{c}).
  \item \inlinecode{set(r, c, x)} assigns \inlinecode{x} to the value at index (\inlinecode{r}, \inlinecode{c}).
  \item \inlinecode{sum(r, c)} returns the sum of the rectangle with upper-left corner (1, 1) and lower-right corner (\inlinecode{r}, \inlinecode{c}).
  \item \inlinecode{sum(r1, c1, r2, c2)} returns the sum of the rectangle with upper-left corner (\inlinecode{r1}, \inlinecode{c1}) and lower-right corner (\inlinecode{r2}, \inlinecode{c2}).
\end{compactitem}

\Needspace*{3\baselineskip}
\textbf{Time Complexity}:
\begin{compactitem}
  \item $O(R \cdot C)$ per call to \inlinecode{initialize(R, C)}.
  \item $O(\log\left(R\right) \cdot \log\left(C\right))$ per call to all other operations.
\end{compactitem}

\Needspace*{3\baselineskip}
\textbf{Space Complexity}:
\begin{compactitem}
  \item $O(R \cdot C)$ for storage of the array elements.
  \item $O(1)$ auxiliary for all operations.
\end{compactitem}

\begin{lstlisting}[language=C++]
#include <algorithm>
#include <vector>

std::vector<std::vector<int>> vals, tree;

void initialize(int R, int C) {
  vals.assign(R + 2, std::vector<int>(C + 2, 0));
  tree.assign(R + 2, std::vector<int>(C + 2, 0));
}

void add(int r, int c, int x) {
  vals[r][c] += x;
  for (int i = r; i < static_cast<int>(tree.size()); i += i & -i) {
    for (int j = c; j < static_cast<int>(tree[0].size()); j += j & -j) {
      tree[i][j] += x;
    }
  }
}

void set(int r, int c, int x) {
  add(r, c, x - vals[r][c]);
}

int sum(int r, int c) {
  int res = 0;
  for (int i = r; i > 0; i -= i & -i) {
    for (int j = c; j > 0; j -= j & -j) {
      res += tree[i][j];
    }
  }
  return res;
}

int sum(int r1, int c1, int r2, int c2) {
  return sum(r2, c2) + sum(r1 - 1, c1 - 1) - sum(r1 - 1, c2) - sum(r2, c1 - 1);
}
\end{lstlisting}
\Needspace*{4\baselineskip}
\textbf{Example Usage}\par
\begin{lstlisting}[language=C++]
#include <cassert>
#include <iostream>
using namespace std;

int main() {
  initialize(3, 3);
  set(1, 1, 5);
  set(1, 2, 6);
  set(2, 1, 7);
  add(3, 3, 9);
  add(2, 1, -4);
  cout << "Values:" << endl;
  for (int i = 1; i <= 3; i++) {
    for (int j = 1; j <= 3; j++) {
      cout << vals[i][j] << " ";
    }
    cout << endl;
  }
  assert(sum(1, 1, 1, 2) == 11);
  assert(sum(1, 1, 2, 1) == 8);
  assert(sum(1, 1, 3, 3) == 23);
  return 0;
}
\end{lstlisting}
\begin{exampleoutput}
Values:
5 6 0
3 0 0
0 0 9
\end{exampleoutput}
\subsection{2D Fenwick Tree (Compressed)}
Maintain a 2D array of numerical type, allowing for rectangular sub-matrices to be simultaneously incremented by arbitrary values (range update) and queries for the sum of rectangular sub-matrices (range query). This implementation uses \inlinecode{std::unordered\_map} for coordinate compression, allowing for large indices to be accessed with efficient space complexity. That is, rows have valid indices from 0 to \inlinecode{R}, inclusive, and columns have valid indices from 0 to \inlinecode{C}, inclusive.

\begin{compactitem}
  \item \inlinecode{add(r, c, x)} adds \inlinecode{x} to the value at index (\inlinecode{r}, \inlinecode{c}).
  \item \inlinecode{add(r1, c1, r2, c2, x)} adds \inlinecode{x} to all indices in the rectangle with upper-left corner (\inlinecode{r1}, \inlinecode{c1}) and lower-right corner (\inlinecode{r2}, \inlinecode{c2}).
  \item \inlinecode{set(r, c, x)} assigns \inlinecode{x} to the value at index (\inlinecode{r}, \inlinecode{c}).
  \item \inlinecode{sum(r, c)} returns the sum of the rectangle with upper-left corner (0, 0) and lower-right corner (\inlinecode{r}, \inlinecode{c}).
  \item \inlinecode{sum(r1, c1, r2, c2)} returns the sum of the rectangle with upper-left corner (\inlinecode{r1}, \inlinecode{c1}) and lower-right corner (\inlinecode{r2}, \inlinecode{c2}).
  \item \inlinecode{at(r, c)} returns the value at index (\inlinecode{r}, \inlinecode{c}).
\end{compactitem}

\textbf{Time Complexity}: $O(\log\left(R\right) \cdot \log\left(C\right))$ per call to all member functions.

\Needspace*{3\baselineskip}
\textbf{Space Complexity}:
\begin{compactitem}
  \item $O(n \cdot \log\left(R\right) \cdot \log\left(C\right))$ for storage of the array elements, where $n$ is the number of distinct indices that have been accessed across all of the operations so far.
  \item $O(1)$ auxiliary for all operations.
\end{compactitem}

\begin{lstlisting}[language=C++]
#include <unordered_map>
#include <utility>

template<class T>
class FenwickTree2D {
  static const int R = 1000000001;
  static const int C = 1000000001;

  std::unordered_map<long long, T> t1, t2, t3, t4;

  template<class Map>
  static T get(const Map &tree, int r, int c) {
    auto it = tree.find(static_cast<long long>(r) * C + c);
    return it == tree.end() ? T() : it->second;
  }

  template<class Map>
  static void add(Map &tree, int r, int c, const T &x) {
    for (int i = r + 1; i <= R; i += i & -i) {
      for (int j = c + 1; j <= C; j += j & -j) {
        tree[static_cast<long long>(i) * C + j] += x;
      }
    }
  }

  void add_helper(int r, int c, const T &x) {
    add(t1, 0, 0, x);
    add(t1, 0, c, -x);
    add(t2, 0, c, x * c);
    add(t1, r, 0, -x);
    add(t3, r, 0, x * r);
    add(t1, r, c, x);
    add(t2, r, c, -x * c);
    add(t3, r, c, -x * r);
    add(t4, r, c, x * r * c);
  }

 public:
  void add(int r1, int c1, int r2, int c2, const T &x) {
    add_helper(r2 + 1, c2 + 1, x);
    add_helper(r1, c2 + 1, -x);
    add_helper(r2 + 1, c1, -x);
    add_helper(r1, c1, x);
  }

  void add(int r, int c, const T &x) { add(r, c, r, c, x); }
  void set(int r, int c, const T &x) { add(r, c, x - at(r, c)); }

  T sum(int r, int c) {
    r++;
    c++;
    T s1 = 0, s2 = 0, s3 = 0, s4 = 0;
    for (int i = r; i > 0; i -= i & -i) {
      for (int j = c; j > 0; j -= j & -j) {
        s1 += get(t1, i, j);
        s2 += get(t2, i, j);
        s3 += get(t3, i, j);
        s4 += get(t4, i, j);
      }
    }
    return s1 * r * c + s2 * r + s3 * c + s4;
  }

  T sum(int r1, int c1, int r2, int c2) {
    return sum(r2, c2) + sum(r1 - 1, c1 - 1) - sum(r1 - 1, c2) - sum(r2, c1 - 1);
  }

  T at(int r, int c) { return sum(r, c, r, c); }
};
\end{lstlisting}
\Needspace*{4\baselineskip}
\textbf{Example Usage}\par
\begin{lstlisting}[language=C++]
#include <cassert>
#include <iostream>
using namespace std;

int main() {
  FenwickTree2D<int> t;
  t.set(0, 0, 5);
  t.set(0, 1, 6);
  t.set(1, 0, 7);
  t.add(2, 2, 9);
  t.add(1, 0, -4);
  t.add(1, 1, 2, 2, 5);
  cout << "Values:" << endl;
  for (int i = 0; i < 3; i++) {
    for (int j = 0; j < 3; j++) {
      cout << t.at(i, j) << " ";
    }
    cout << endl;
  }
  assert(t.sum(0, 0, 0, 1) == 11);
  assert(t.sum(0, 0, 1, 0) == 8);
  assert(t.sum(1, 1, 2, 2) == 29);
  t.set(500000000, 500000000, 100);
  assert(t.sum(0, 0, 1000000000, 1000000000) == 143);
  return 0;
}
\end{lstlisting}
\begin{exampleoutput}
Values:
5 6 0
3 5 5
0 5 14
\end{exampleoutput}

\section{Disjoint Sets and Tree Structures}
\setcounter{section}{7}
\setcounter{subsection}{0}
\subsection{Disjoint Set Forest (Simple)}
Maintain a set of elements partitioned into non-overlapping subsets. Each partition is assigned a unique representative known as the parent, or root. The following implements two well-known optimizations known as union-by-rank and path compression. This version is simplified to only work on integer elements.

\begin{compactitem}
  \item \inlinecode{initialize()} resets the data structure.
  \item \inlinecode{make\_set(u)} creates a new partition consisting of the single element \inlinecode{u}, which must not have been previously added to the data structure.
  \item \inlinecode{find\_root(u)} returns the unique representative of the partition containing \inlinecode{u}.
  \item \inlinecode{is\_united(u, v)} returns whether elements \inlinecode{u} and \inlinecode{v} belong to the same partition.
  \item \inlinecode{unite(u, v)} replaces the partitions containing \inlinecode{u} and \inlinecode{v} with a single new partition consisting of the union of elements in the original partitions.
\end{compactitem}

A precondition to the last three operations is that \inlinecode{make\_set()} must have been previously called on their arguments.

\Needspace*{3\baselineskip}
\textbf{Time Complexity}:
\begin{compactitem}
  \item $O(1)$ per call to \inlinecode{initialize()} and \inlinecode{make\_set()}.
  \item $O(\alpha\left(n\right))$ per call to \inlinecode{find\_root()}, \inlinecode{is\_united()}, and \inlinecode{unite()}, where $n$ is the number of elements that has been added via \inlinecode{make\_set()} so far, and $\alpha(n)$ is the extremely slow growing inverse of the Ackermann function (effectively a very small constant for all practical values of $n$).
\end{compactitem}

\Needspace*{3\baselineskip}
\textbf{Space Complexity}:
\begin{compactitem}
  \item $O(n)$ for storage of the disjoint set forest elements.
  \item $O(1)$ auxiliary for all operations.
\end{compactitem}

\begin{lstlisting}[language=C++]
#include <vector>

int num_sets;
std::vector<int> root, rank;

void initialize() {
  num_sets = 0;
  root.clear();
  rank.clear();
}

void make_set(int u) {
  if (u >= static_cast<int>(root.size())) {
    root.resize(u + 1);
    rank.resize(u + 1);
  }
  root[u] = u;
  rank[u] = 0;
  num_sets++;
}

int find_root(int u) {
  if (root[u] != u) {
    root[u] = find_root(root[u]);
  }
  return root[u];
}

bool is_united(int u, int v) {
  return find_root(u) == find_root(v);
}

void unite(int u, int v) {
  int ru = find_root(u), rv = find_root(v);
  if (ru == rv) {
    return;
  }
  num_sets--;
  if (rank[ru] < rank[rv]) {
    root[ru] = rv;
  } else {
    root[rv] = ru;
    if (rank[ru] == rank[rv]) {
      rank[ru]++;
    }
  }
}
\end{lstlisting}
\Needspace*{4\baselineskip}
\textbf{Example Usage}\par
\begin{lstlisting}[language=C++]
#include <cassert>

int main() {
  initialize();
  for (char c = 'a'; c <= 'g'; c++) {
    make_set(c);
  }
  unite('a', 'b');
  unite('b', 'f');
  unite('d', 'e');
  unite('d', 'g');
  assert(num_sets == 3);
  assert(is_united('a', 'b'));
  assert(!is_united('a', 'c'));
  assert(!is_united('b', 'g'));
  assert(is_united('e', 'g'));
  return 0;
}
\end{lstlisting}
\subsection{Disjoint Set Forest (Compressed)}
Maintain a set of elements partitioned into non-overlapping subsets using a collection of trees. Each partition is assigned a unique representative known as the parent, or root. The following implements two well-known optimizations known as union-by-rank and path compression. This version uses an \inlinecode{std::unordered\_map} for storage and coordinate compression (thus, element types must meet the requirements of key types for \inlinecode{std::unordered\_map}).

\begin{compactitem}
  \item \inlinecode{make\_set(u)} creates a new partition consisting of the single element \inlinecode{u}, which must not have been previously added to the data structure.
  \item \inlinecode{is\_united(u, v)} returns whether elements \inlinecode{u} and \inlinecode{v} belong to the same partition.
  \item \inlinecode{unite(u, v)} replaces the partitions containing \inlinecode{u} and \inlinecode{v} with a single new partition consisting of the union of elements in the original partitions.
  \item \inlinecode{get\_all\_sets()} returns all current partitions as a vector of vectors.
\end{compactitem}

A precondition to the last three operations is that \inlinecode{make\_set()} must have been previously called on their arguments.

\Needspace*{3\baselineskip}
\textbf{Time Complexity}:
\begin{compactitem}
  \item $O(1)$ per call to the constructor.
  \item $O(\log n)$ per call to \inlinecode{make\_set()}, where $n$ is the number of elements that have been added via \inlinecode{make\_set()} so far.
  \item $O(\alpha\left(n\right) \log n)$ per call to \inlinecode{is\_united()} and \inlinecode{unite()}, where $n$ is the number of elements that have been added via \inlinecode{make\_set()} so far, and $\alpha(n)$ is the extremely slow growing inverse of the Ackermann function (effectively a very small constant for all practical values of $n$).
  \item $O(n)$ per call to \inlinecode{get\_all\_sets()}.
\end{compactitem}

\Needspace*{3\baselineskip}
\textbf{Space Complexity}:
\begin{compactitem}
  \item $O(n)$ for storage of the disjoint set forest elements.
  \item $O(n)$ auxiliary heap space for \inlinecode{get\_all\_sets()}.
  \item $O(1)$ auxiliary for all other operations.
\end{compactitem}

\begin{lstlisting}[language=C++]
#include <unordered_map>
#include <vector>

template<class T>
class DisjointSetForest {
  int num_elements, num_sets;
  std::unordered_map<T, int> id;
  std::vector<int> root, rank;

  int find_root(int u) {
    if (root[u] != u) {
      root[u] = find_root(root[u]);
    }
    return root[u];
  }

 public:
  DisjointSetForest() : num_elements(0), num_sets(0) {}

  int size() const { return num_elements; }
  int sets() const { return num_sets; }

  void make_set(const T &u) {
    if (id.find(u) != id.end()) {
      return;
    }
    id[u] = num_elements;
    root.push_back(num_elements++);
    rank.push_back(0);
    num_sets++;
  }

  bool is_united(const T &u, const T &v) { return find_root(id[u]) == find_root(id[v]); }

  void unite(const T &u, const T &v) {
    int ru = find_root(id[u]), rv = find_root(id[v]);
    if (ru == rv) {
      return;
    }
    num_sets--;
    if (rank[ru] < rank[rv]) {
      root[ru] = rv;
    } else {
      root[rv] = ru;
      if (rank[ru] == rank[rv]) {
        rank[ru]++;
      }
    }
  }

  std::vector<std::vector<T>> get_all_sets() {
    std::unordered_map<int, std::vector<T>> tmp;
    for (auto &[key, val] : id) {
      tmp[find_root(val)].push_back(key);
    }
    std::vector<std::vector<T>> res;
    for (auto &[key, val] : tmp) {
      res.push_back(val);
    }
    return res;
  }
};
\end{lstlisting}
\Needspace*{4\baselineskip}
\textbf{Example Usage}\par
\begin{lstlisting}[language=C++]
#include <cassert>
#include <iostream>
using namespace std;

int main() {
  DisjointSetForest<char> dsf;
  for (char c = 'a'; c <= 'g'; c++) {
    dsf.make_set(c);
  }
  dsf.unite('a', 'b');
  dsf.unite('b', 'f');
  dsf.unite('d', 'e');
  dsf.unite('d', 'g');
  assert(dsf.size() == 7);
  assert(dsf.sets() == 3);
  auto s = dsf.get_all_sets();
  bool first_set = true;
  for (const auto &set : s) {
    cout << (first_set ? "[" : ", [");
    first_set = false;
    bool first_value = true;
    for (char value : set) {
      cout << (first_value ? "" : ", ") << value;
      first_value = false;
    }
    cout << "]";
  }
  cout << endl;
  return 0;
}
\end{lstlisting}
\begin{exampleoutput}
[a, b, f], [c], [d, e, g]
\end{exampleoutput}
\subsection{Disjoint Set Forest (Rollback)}
Maintains disjoint sets while supporting rollback to previous snapshots. Rollback DSU is useful for offline dynamic connectivity, divide-and-conquer over time, and backtracking search where unions must be undone.

Path compression is intentionally not used, because it is hard to undo. Union by size keeps tree height logarithmic, and every successful union records enough information to restore the previous state.

\begin{compactitem}
  \item \inlinecode{initialize(n)} creates singleton sets over elements \inlinecode{0} through \inlinecode{n - 1}.
  \item \inlinecode{find\_root(u)} returns the representative of the set containing \inlinecode{u}.
  \item \inlinecode{is\_united(u, v)} returns whether \inlinecode{u} and \inlinecode{v} are in the same set.
  \item \inlinecode{unite(u, v)} merges two sets and returns whether a merge occurred.
  \item \inlinecode{snapshot()} returns a token representing the current history size.
  \item \inlinecode{rollback(s)} undoes all changes made after snapshot \inlinecode{s}.
\end{compactitem}

\Needspace*{3\baselineskip}
\textbf{Time Complexity}:
\begin{compactitem}
  \item $O(n)$ per call to \inlinecode{initialize(n)}.
  \item $O(\log n)$ worst-case per call to \inlinecode{find\_root(u)}, \inlinecode{is\_united(u, v)}, and \inlinecode{unite(u, v)}.
  \item $O(1)$ per undone union during \inlinecode{rollback(s)}.
\end{compactitem}

\textbf{Space Complexity}: $O(n + q)$ for $q$ successful unions stored in the rollback history.

\begin{lstlisting}[language=C++]
#include <numeric>
#include <vector>

class RollbackDSU {
  struct Change {
    int child, parent, parent_size;

    Change(int child = -1, int parent = -1, int parent_size = 0)
        : child(child), parent(parent), parent_size(parent_size) {}
  };

  std::vector<int> root, size;
  std::vector<Change> history;
  int sets;

 public:
  RollbackDSU() : sets(0) {}

  explicit RollbackDSU(int n) { initialize(n); }

  void initialize(int n) {
    root.resize(n);
    size.assign(n, 1);
    history.clear();
    sets = n;
    std::iota(root.begin(), root.end(), 0);
  }

  int count_sets() const { return sets; }

  int find_root(int u) const {
    while (root[u] != u) {
      u = root[u];
    }
    return u;
  }

  bool is_united(int u, int v) const { return find_root(u) == find_root(v); }

  bool unite(int u, int v) {
    u = find_root(u);
    v = find_root(v);
    if (u == v) {
      return false;
    }
    if (size[u] > size[v]) {
      int tmp = u;
      u = v;
      v = tmp;
    }
    history.emplace_back(u, v, size[v]);
    root[u] = v;
    size[v] += size[u];
    sets--;
    return true;
  }

  int snapshot() const { return history.size(); }

  void rollback(int snapshot) {
    while (static_cast<int>(history.size()) > snapshot) {
      Change c = history.back();
      history.pop_back();
      root[c.child] = c.child;
      size[c.parent] = c.parent_size;
      sets++;
    }
  }
};
\end{lstlisting}
\Needspace*{4\baselineskip}
\textbf{Example Usage}\par
\begin{lstlisting}[language=C++]
#include <cassert>

int main() {
  RollbackDSU dsu(5);
  dsu.unite(0, 1);
  int s = dsu.snapshot();
  dsu.unite(1, 2);
  dsu.unite(3, 4);
  assert(dsu.is_united(0, 2));
  assert(dsu.count_sets() == 2);
  dsu.rollback(s);
  assert(dsu.is_united(0, 1));
  assert(!dsu.is_united(0, 2));
  assert(!dsu.is_united(3, 4));
  assert(dsu.count_sets() == 4);
  return 0;
}
\end{lstlisting}
\subsection{Lowest Common Ancestor (Sparse Table)}
Given a tree, determine the lowest common ancestor of any two nodes in the tree. The lowest common ancestor of two nodes $u$ and $v$ is the node that has the longest distance from the root while having both $u$ and $v$ as its descendant. A nodes is considered to be a descendant of itself. \inlinecode{build()} applies to a global, pre-populated adjacency list \inlinecode{adj} which must only consist of nodes numbered with integers between 0 (inclusive) and the total number of nodes (exclusive), as passed in the function argument.

\Needspace*{3\baselineskip}
\textbf{Time Complexity}:
\begin{compactitem}
  \item $O(n \log n)$ per call to \inlinecode{build()}, where $n$ is the number of nodes.
  \item $O(\log n)$ per call to \inlinecode{lca()}.
\end{compactitem}

\Needspace*{3\baselineskip}
\textbf{Space Complexity}:
\begin{compactitem}
  \item $O(n \log n)$ to store the sparse table, where $n$ is the number of nodes.
  \item $O(n)$ auxiliary stack space for \inlinecode{build()}.
  \item $O(1)$ auxiliary for \inlinecode{lca()}.
\end{compactitem}

\begin{lstlisting}[language=C++]
#include <vector>

std::vector<std::vector<int>> adj, dp;
int len, timer;
std::vector<int> tin, tout;

void dfs(int u, int p) {
  tin[u] = timer++;
  dp[u][0] = p;
  for (int i = 1; i < len; i++) {
    dp[u][i] = dp[dp[u][i - 1]][i - 1];
  }
  for (int v : adj[u]) {
    if (v != p) {
      dfs(v, u);
    }
  }
  tout[u] = timer++;
}

void build(int nodes, int root = 0) {
  len = 1;
  while ((1 << len) <= nodes) {
    len++;
  }
  dp.assign(nodes, std::vector<int>(len));
  tin.assign(nodes, 0);
  tout.assign(nodes, 0);
  timer = 0;
  dfs(root, root);
}

bool is_ancestor(int parent, int child) {
  return (tin[parent] <= tin[child]) && (tout[child] <= tout[parent]);
}

int lca(int u, int v) {
  if (is_ancestor(u, v)) {
    return u;
  }
  if (is_ancestor(v, u)) {
    return v;
  }
  for (int i = len - 1; i >= 0; i--) {
    if (!is_ancestor(dp[u][i], v)) {
      u = dp[u][i];
    }
  }
  return dp[u][0];
}
\end{lstlisting}
\Needspace*{4\baselineskip}
\textbf{Example Usage}\par
\begin{lstlisting}[language=C++]
#include <cassert>
using namespace std;

int main() {
  int nodes = 5;
  adj.resize(nodes);
  adj[0].push_back(1);
  adj[1].push_back(0);
  adj[1].push_back(2);
  adj[2].push_back(1);
  adj[3].push_back(1);
  adj[1].push_back(3);
  adj[0].push_back(4);
  adj[4].push_back(0);
  build(nodes, 0);
  assert(lca(3, 2) == 1);
  assert(lca(2, 4) == 0);
  return 0;
}
\end{lstlisting}
\subsection{Lowest Common Ancestor (Segment Tree)}
Given a tree, determine the lowest common ancestor of any two nodes in the tree. The lowest common ancestor of two nodes $u$ and $v$ is the node that has the longest distance from the root while having both $u$ and $v$ as its descendant. A nodes is considered to be a descendant of itself. \inlinecode{build()} applies to a global, pre-populated adjacency list \inlinecode{adj} which must only consist of nodes numbered with integers between 0 (inclusive) and the total number of nodes (exclusive), as passed in the function argument.

\Needspace*{3\baselineskip}
\textbf{Time Complexity}:
\begin{compactitem}
  \item $O(n \log n)$ per call to \inlinecode{build()}, where $n$ is the number of nodes.
  \item $O(\log n)$ per call to \inlinecode{lca()}.
\end{compactitem}

\Needspace*{3\baselineskip}
\textbf{Space Complexity}:
\begin{compactitem}
  \item $O(n)$ for storage of the segment tree, where $n$ is the number of nodes.
  \item $O(n \log n)$ auxiliary stack space for \inlinecode{build()}.
  \item $O(\log n)$ auxiliary stack space for \inlinecode{lca()}.
\end{compactitem}

\begin{lstlisting}[language=C++]
#include <algorithm>
#include <vector>

std::vector<std::vector<int>> adj;
int len, counter;
std::vector<int> depth, dfs_order, first, minpos;

void dfs(int u, int d) {
  depth[u] = d;
  dfs_order[counter++] = u;
  for (int v : adj[u]) {
    if (depth[v] == -1) {
      dfs(v, d + 1);
      dfs_order[counter++] = u;
    }
  }
}

void build(int n, int lo, int hi) {
  if (lo == hi) {
    minpos[n] = dfs_order[lo];
    return;
  }
  int lchild = 2 * n + 1, rchild = 2 * n + 2, mid = lo + (hi - lo) / 2;
  build(lchild, lo, mid);
  build(rchild, mid + 1, hi);
  minpos[n] = depth[minpos[lchild]] < depth[minpos[rchild]] ? minpos[lchild] : minpos[rchild];
}

void build(int nodes, int root) {
  depth.assign(nodes, -1);
  dfs_order.assign(2 * nodes, 0);
  first.assign(nodes, -1);
  minpos.assign(8 * nodes, 0);
  len = 2 * nodes - 1;
  counter = 0;
  dfs(root, 0);
  build(0, 0, len - 1);
  for (int i = 0; i < len; i++) {
    if (first[dfs_order[i]] == -1) {
      first[dfs_order[i]] = i;
    }
  }
}

int get_minpos(int a, int b, int n, int lo, int hi) {
  if (a == lo && b == hi) {
    return minpos[n];
  }
  int mid = lo + (hi - lo) / 2;
  if (a <= mid && mid < b) {
    int p1 = get_minpos(a, std::min(b, mid), 2 * n + 1, lo, mid);
    int p2 = get_minpos(std::max(a, mid + 1), b, 2 * n + 2, mid + 1, hi);
    return depth[p1] < depth[p2] ? p1 : p2;
  }
  if (a <= mid) {
    return get_minpos(a, std::min(b, mid), 2 * n + 1, lo, mid);
  }
  return get_minpos(std::max(a, mid + 1), b, 2 * n + 2, mid + 1, hi);
}

int lca(int u, int v) {
  return get_minpos(std::min(first[u], first[v]), std::max(first[u], first[v]), 0, 0, len - 1);
}
\end{lstlisting}
\Needspace*{4\baselineskip}
\textbf{Example Usage}\par
\begin{lstlisting}[language=C++]
#include <cassert>
using namespace std;

int main() {
  int nodes = 5;
  adj.resize(nodes);
  adj[0].push_back(1);
  adj[1].push_back(0);
  adj[1].push_back(2);
  adj[2].push_back(1);
  adj[3].push_back(1);
  adj[1].push_back(3);
  adj[0].push_back(4);
  adj[4].push_back(0);
  build(nodes, 0);
  assert(lca(3, 2) == 1);
  assert(lca(2, 4) == 0);
  return 0;
}
\end{lstlisting}
\subsection{Heavy Light Decomposition}
Maintain a tree with values associated with either its edges or nodes, while supporting both dynamic queries and dynamic updates of all values on a given path between two nodes in the tree. Heavy-light decomposition partitions the nodes of the tree into disjoint paths where all nodes have degree two, except the endpoints of a path which has degree one.

The query operation is defined by an associative \inlinecode{combine()} function which satisfies \inlinecode{combine(x, combine(y, z)) = combine(combine(x, y), z)} for all values \inlinecode{x}, \inlinecode{y}, and \inlinecode{z} in the tree. The default code below assumes a numerical tree type, defining queries for the ``min'' of the target range. Another possible query operation is ``sum'', in which case \inlinecode{combine(a, b)} should return $a + b$. For direction-independent path queries, \inlinecode{combine()} should also be commutative; otherwise, store enough information in each aggregate to combine paths in the required order.

The update operation is defined by \inlinecode{apply\_delta()} and \inlinecode{compose\_deltas()}. A delta must act on an aggregate summary of a path of length \inlinecode{len}: \inlinecode{apply\_delta(v, d, len)} returns the aggregate after applying update \inlinecode{d} to every element represented by aggregate \inlinecode{v}. Pending deltas are combined in chronological order by \inlinecode{compose\_deltas(old, d)}, meaning ``apply \inlinecode{old}, then apply \inlinecode{d}''. These hooks do not support arbitrary query/update pairings; the delta operation must distribute over \inlinecode{combine()}, and composed deltas must have the same effect as applying their updates sequentially. The default code below defines updates that ``set'' a path's edges or nodes to a new value. For range increment updates, \inlinecode{apply\_delta(v, d, len)} would return \inlinecode{v + d} for min/max queries, or \inlinecode{v + d * len} for sum queries, and \inlinecode{compose\_deltas(old, d)} would return \inlinecode{old + d}.

\begin{compactitem}
  \item \inlinecode{HeavyLight(adj, v)} constructs a new heavy light decomposition on a tree defined by the adjacency list \inlinecode{adj}, with all values initialized to \inlinecode{v}. The adjacency list must consist of only the integers from 0 to \inlinecode{adj.size() - 1}, inclusive. No duplicate edges should exist, and the graph must be connected.
  \item \inlinecode{query(u, v)} returns the result of \inlinecode{combine()} applied to all values on the path from node \inlinecode{u} to node \inlinecode{v}.
  \item \inlinecode{update(u, v, d)} modifies all values on the path from node \inlinecode{u} to node \inlinecode{v} by respectively applying the delta \inlinecode{d}.
\end{compactitem}

\Needspace*{3\baselineskip}
\textbf{Time Complexity}:
\begin{compactitem}
  \item $O(n)$ per call to the constructor, where $n$ is the number of nodes.
  \item $O(\log^{2} n)$ per call to \inlinecode{query()} and \inlinecode{update()};
\end{compactitem}

\Needspace*{3\baselineskip}
\textbf{Space Complexity}:
\begin{compactitem}
  \item $O(n)$ for storage of the decomposition.
  \item $O(n)$ auxiliary stack space for the constructor.
  \item $O(1)$ auxiliary for \inlinecode{query()} and \inlinecode{update()}.
\end{compactitem}

\begin{lstlisting}[language=C++]
#include <algorithm>
#include <stdexcept>
#include <vector>

template<class T>
class HeavyLight {
  // Set this to true to store values on edges, false to store values on nodes.
  static const bool VALUES_ON_EDGES = true;

  static T combine(const T &a, const T &b) { return std::min(a, b); }
  static T apply_delta(const T &v, const T &d, int len) { return d; }
  static T compose_deltas(const T &old, const T &d) { return d; }

  int counter, paths;
  std::vector<std::vector<T>> value, delta;
  std::vector<std::vector<bool>> pending;
  std::vector<std::vector<int>> len;
  std::vector<int> size, parent, tin, tout, path, pathlen, pathpos, pathroot;
  const std::vector<std::vector<int>> &adj;

  void dfs(int u, int p) {
    tin[u] = counter++;
    parent[u] = p;
    size[u] = 1;
    for (int v : adj[u]) {
      if (v != p) {
        dfs(v, u);
        size[u] += size[v];
      }
    }
    tout[u] = counter++;
  }

  int new_path(int u) {
    pathroot[paths] = u;
    return paths++;
  }

  void build_paths(int u, int path) {
    this->path[u] = path;
    pathpos[u] = pathlen[path]++;
    for (int v : adj[u]) {
      if (v != parent[u]) {
        build_paths(v, (2 * size[v] >= size[u]) ? path : new_path(v));
      }
    }
  }

  inline T applied_value(int path, int i) {
    return pending[path][i] ? apply_delta(value[path][i], delta[path][i], len[path][i])
                            : value[path][i];
  }

  void push_delta(int path, int i) {
    int d = 0;
    while ((i >> d) > 0) {
      d++;
    }
    for (d -= 2; d >= 0; d--) {
      int l = (i >> d), r = (l ^ 1), n = l / 2;
      if (pending[path][n]) {
        value[path][n] = applied_value(path, n);
        delta[path][l] =
            pending[path][l] ? compose_deltas(delta[path][l], delta[path][n]) : delta[path][n];
        delta[path][r] =
            pending[path][r] ? compose_deltas(delta[path][r], delta[path][n]) : delta[path][n];
        pending[path][l] = pending[path][r] = true;
        pending[path][n] = false;
      }
    }
  }

  bool query(int path, int u, int v, T *res) {
    push_delta(path, u += value[path].size() / 2);
    push_delta(path, v += value[path].size() / 2);
    bool found = false;
    for (; u <= v; u = (u + 1) / 2, v = (v - 1) / 2) {
      if ((u & 1) != 0) {
        T value = applied_value(path, u);
        *res = found ? combine(*res, value) : value;
        found = true;
      }
      if ((v & 1) == 0) {
        T value = applied_value(path, v);
        *res = found ? combine(*res, value) : value;
        found = true;
      }
    }
    return found;
  }

  void update(int path, int u, int v, const T &d) {
    push_delta(path, u += value[path].size() / 2);
    push_delta(path, v += value[path].size() / 2);
    int tu = -1, tv = -1;
    for (; u <= v; u = (u + 1) / 2, v = (v - 1) / 2) {
      if ((u & 1) != 0) {
        delta[path][u] = pending[path][u] ? compose_deltas(delta[path][u], d) : d;
        pending[path][u] = true;
        if (tu == -1) {
          tu = u;
        }
      }
      if ((v & 1) == 0) {
        delta[path][v] = pending[path][v] ? compose_deltas(delta[path][v], d) : d;
        pending[path][v] = true;
        if (tv == -1) {
          tv = v;
        }
      }
    }
    for (int i = tu; i > 1; i /= 2) {
      value[path][i / 2] = combine(applied_value(path, i), applied_value(path, i ^ 1));
    }
    for (int i = tv; i > 1; i /= 2) {
      value[path][i / 2] = combine(applied_value(path, i), applied_value(path, i ^ 1));
    }
  }

  inline bool is_ancestor(int parent, int child) {
    return (tin[parent] <= tin[child]) && (tout[child] <= tout[parent]);
  }

 public:
  explicit HeavyLight(const std::vector<std::vector<int>> &adj, const T &v = T())
      : counter(0),
        paths(0),
        size(adj.size()),
        parent(adj.size()),
        tin(adj.size()),
        tout(adj.size()),
        path(adj.size()),
        pathlen(adj.size()),
        pathpos(adj.size()),
        pathroot(adj.size()),
        adj(adj) {
    dfs(0, -1);
    build_paths(0, new_path(0));
    value.resize(paths);
    delta.resize(paths);
    pending.resize(paths);
    len.resize(paths);
    for (int i = 0; i < paths; i++) {
      int m = pathlen[i];
      value[i].assign(2 * m, v);
      delta[i].resize(2 * m);
      pending[i].assign(2 * m, false);
      len[i].assign(2 * m, 1);
      for (int j = 2 * m - 1; j > 1; j -= 2) {
        value[i][j / 2] = combine(value[i][j], value[i][j ^ 1]);
        len[i][j / 2] = len[i][j] + len[i][j ^ 1];
      }
    }
  }

  T query(int u, int v) {
    if (VALUES_ON_EDGES && u == v) {
      throw std::runtime_error("No edge between u and v to be queried.");
    }
    bool found = false;
    T res = T(), value;
    int root;
    while (!is_ancestor(root = pathroot[path[u]], v)) {
      if (query(path[u], 0, pathpos[u], &value)) {
        res = found ? combine(res, value) : value;
        found = true;
      }
      u = parent[root];
    }
    while (!is_ancestor(root = pathroot[path[v]], u)) {
      if (query(path[v], 0, pathpos[v], &value)) {
        res = found ? combine(res, value) : value;
        found = true;
      }
      v = parent[root];
    }
    if (query(
            path[u], std::min(pathpos[u], pathpos[v]) + static_cast<int>(VALUES_ON_EDGES),
            std::max(pathpos[u], pathpos[v]), &value
        )) {
      res = found ? combine(res, value) : value;
      found = true;
    }
    if (!found) {
      throw std::runtime_error("Unexpected error: No values found.");
    }
    return res;
  }

  void update(int u, int v, const T &d) {
    if (VALUES_ON_EDGES && u == v) {
      return;
    }
    int root;
    while (!is_ancestor(root = pathroot[path[u]], v)) {
      update(path[u], 0, pathpos[u], d);
      u = parent[root];
    }
    while (!is_ancestor(root = pathroot[path[v]], u)) {
      update(path[v], 0, pathpos[v], d);
      v = parent[root];
    }
    update(
        path[u], std::min(pathpos[u], pathpos[v]) + static_cast<int>(VALUES_ON_EDGES),
        std::max(pathpos[u], pathpos[v]), d
    );
  }
};
\end{lstlisting}
\Needspace*{4\baselineskip}
\textbf{Example Usage}\par
\begin{lstlisting}[language=C++]
#include <cassert>
using namespace std;

int main() {
  //     w=40      w=20      w=10
  // 0---------1---------2---------3
  //                     |
  //                     ----------4
  //                         w=30
  vector<vector<int>> adj(5);
  adj[0].push_back(1);
  adj[1].push_back(0);
  adj[1].push_back(2);
  adj[2].push_back(1);
  adj[2].push_back(3);
  adj[3].push_back(2);
  adj[2].push_back(4);
  adj[4].push_back(2);
  HeavyLight<int> hld(adj, 0);
  hld.update(0, 1, 40);
  hld.update(1, 2, 20);
  hld.update(2, 3, 10);
  hld.update(2, 4, 30);
  assert(hld.query(0, 3) == 10);
  assert(hld.query(2, 4) == 30);
  hld.update(3, 4, 5);
  assert(hld.query(1, 4) == 5);
  return 0;
}
\end{lstlisting}
\subsection{Link-Cut Tree}
Maintain a forest of trees with values associated with its nodes, while supporting both dynamic queries and dynamic updates of all values on any path between two nodes in a given tree. In addition, support testing of whether two nodes are connected in the forest, as well as the merging and spliting of trees by adding or removing specific edges. Link/cut forests divide each of its trees into vertex-disjoint paths, each represented by a splay tree.

The query operation is defined by an associative \inlinecode{combine()} function which satisfies \inlinecode{combine(x, combine(y, z)) = combine(combine(x, y), z)} for all values \inlinecode{x}, \inlinecode{y}, and \inlinecode{z} in the forest. The default code below assumes a numerical forest type, defining queries for the ``min'' of the target range. Another possible query operation is ``sum'', in which case \inlinecode{combine(a, b)} should return $a + b$. For direction-independent path queries, \inlinecode{combine()} should also be commutative; otherwise, store enough information in each aggregate to combine paths in the required order.

The update operation is defined by \inlinecode{apply\_delta()} and \inlinecode{compose\_deltas()}. A delta must act on an aggregate summary of a path of length \inlinecode{len}: \inlinecode{apply\_delta(v, d, len)} returns the aggregate after applying update \inlinecode{d} to every element represented by aggregate \inlinecode{v}. Pending deltas are combined in chronological order by \inlinecode{compose\_deltas(old, d)}, meaning ``apply \inlinecode{old}, then apply \inlinecode{d}''. These hooks do not support arbitrary query/update pairings; the delta operation must distribute over \inlinecode{combine()}, and composed deltas must have the same effect as applying their updates sequentially. The default code below defines updates that ``set'' a path's nodes to a new value. For range increment updates, \inlinecode{apply\_delta(v, d, len)} would return \inlinecode{v + d} for min/max queries, or \inlinecode{v + d * len} for sum queries, and \inlinecode{compose\_deltas(old, d)} would return \inlinecode{old + d}.

\begin{compactitem}
  \item \inlinecode{LinkCutForest()} constructs an empty forest with no trees.
  \item \inlinecode{size()} returns the number of nodes in the forest.
  \item \inlinecode{trees()} returns the number of trees in the forest.
  \item \inlinecode{make\_root(i, v)} creates a new tree in the forest consisting of a single node labeled with the integer \inlinecode{i} and value initialized to \inlinecode{v}.
  \item \inlinecode{is\_connected(a, b)} returns whether nodes \inlinecode{a} and \inlinecode{b} are connected.
  \item \inlinecode{link(a, b)} adds an edge between the nodes \inlinecode{a} and \inlinecode{b}, both of which must exist and not be connected.
  \item \inlinecode{cut(a, b)} removes the edge between the nodes \inlinecode{a} and \inlinecode{b}, both of which must exist and be connected.
  \item \inlinecode{query(a, b)} returns the result of \inlinecode{combine()} applied to all values on the path from the node \inlinecode{a} to node \inlinecode{b}.
  \item \inlinecode{update(a, b, d)} modifies all the values on the path from node \inlinecode{a} to node \inlinecode{b} by respectively applying the delta \inlinecode{d}.
\end{compactitem}

\Needspace*{3\baselineskip}
\textbf{Time Complexity}:
\begin{compactitem}
  \item $O(1)$ per call to the constructor, \inlinecode{size()}, and \inlinecode{trees()}.
  \item $O(\log n)$ amortized per call to all other operations, where $n$ is the number of nodes.
\end{compactitem}

\Needspace*{3\baselineskip}
\textbf{Space Complexity}:
\begin{compactitem}
  \item $O(n)$ for storage of the forest, where $n$ is the number of nodes.
  \item $O(1)$ auxiliary for all operations.
\end{compactitem}

\begin{lstlisting}[language=C++]
#include <algorithm>
#include <cstddef>
#include <stdexcept>
#include <unordered_map>

template<class T>
class LinkCutForest {
  static T combine(const T &a, const T &b) { return std::min(a, b); }
  static T apply_delta(const T &v, const T &d, int len) { return d; }
  static T compose_deltas(const T &old, const T &d) { return d; }

  struct Node {
    T value, subtree_value, delta;
    int size;
    bool rev, pending;
    Node *left, *right, *parent;

    explicit Node(const T &v)
        : value(v),
          subtree_value(v),
          size(1),
          rev(false),
          pending(false),
          left(nullptr),
          right(nullptr),
          parent(nullptr) {}

    inline bool is_root() const {
      return parent == nullptr || (parent->left != this && parent->right != this);
    }

    inline T get_subtree_value() const {
      return pending ? apply_delta(subtree_value, delta, size) : subtree_value;
    }

    void push() {
      if (rev) {
        rev = false;
        std::swap(left, right);
        if (left != nullptr) {
          left->rev = !left->rev;
        }
        if (right != nullptr) {
          right->rev = !right->rev;
        }
      }
      if (pending) {
        value = apply_delta(value, delta, 1);
        subtree_value = apply_delta(subtree_value, delta, size);
        if (left != nullptr) {
          left->delta = left->pending ? compose_deltas(left->delta, delta) : delta;
          left->pending = true;
        }
        if (right != nullptr) {
          right->delta = right->pending ? compose_deltas(right->delta, delta) : delta;
          right->pending = true;
        }
        pending = false;
      }
    }

    void update() {
      size = 1;
      subtree_value = value;
      if (left != nullptr) {
        subtree_value = combine(subtree_value, left->get_subtree_value());
        size += left->size;
      }
      if (right != nullptr) {
        subtree_value = combine(subtree_value, right->get_subtree_value());
        size += right->size;
      }
    }
  };

  int num_trees;
  std::unordered_map<int, Node *> nodes;

  static void connect(Node *child, Node *parent, bool is_left) {
    if (child != nullptr) {
      child->parent = parent;
    }
    if (is_left) {
      parent->left = child;
    } else {
      parent->right = child;
    }
  }

  static void rotate(Node *n) {
    Node *parent = n->parent, *grandparent = parent->parent;
    bool parent_is_root = parent->is_root(), is_left = (n == parent->left);
    connect(is_left ? n->right : n->left, parent, is_left);
    connect(parent, n, !is_left);
    if (parent_is_root) {
      if (n != nullptr) {
        n->parent = grandparent;
      }
    } else {
      connect(n, grandparent, parent == grandparent->left);
    }
    parent->update();
  }

  static void splay(Node *n) {
    while (!n->is_root()) {
      Node *parent = n->parent, *grandparent = parent->parent;
      if (!parent->is_root()) {
        grandparent->push();
      }
      parent->push();
      n->push();
      if (!parent->is_root()) {
        if ((n == parent->left) == (parent == grandparent->left)) {
          rotate(parent);
        } else {
          rotate(n);
        }
      }
      rotate(n);
    }
    n->push();
    n->update();
  }

  static Node *expose(Node *n) {
    Node *prev = nullptr;
    for (Node *curr = n; curr != nullptr; curr = curr->parent) {
      splay(curr);
      curr->left = prev;
      prev = curr;
    }
    splay(n);
    return prev;
  }

  // Helper variables.
  Node *u, *v;

  void get_uv(int a, int b) {
    auto it1 = nodes.find(a);
    auto it2 = nodes.find(b);
    if (it1 == nodes.end() || it2 == nodes.end()) {
      throw std::runtime_error("Queried node ID does not exist in forest.");
    }
    u = it1->second;
    v = it2->second;
  }

 public:
  LinkCutForest() : num_trees(0) {}

  ~LinkCutForest() {
    for (auto &[key, node] : nodes) {
      delete node;
    }
  }

  int size() const { return nodes.size(); }
  int trees() const { return num_trees; }

  void make_root(int i, const T &v = T()) {
    if (auto it = nodes.find(i); it != nodes.end()) {
      throw std::runtime_error("Cannot make a root with an existing ID.");
    }
    Node *n = new Node(v);
    expose(n);
    n->rev = !n->rev;
    nodes[i] = n;
    num_trees++;
  }

  bool is_connected(int a, int b) {
    get_uv(a, b);
    if (a == b) {
      return true;
    }
    expose(u);
    expose(v);
    return u->parent != nullptr;
  }

  void link(int a, int b) {
    if (is_connected(a, b)) {
      throw std::runtime_error("Cannot link nodes that are already connected.");
    }
    get_uv(a, b);
    expose(u);
    u->rev = !u->rev;
    u->parent = v;
    num_trees--;
  }

  void cut(int a, int b) {
    get_uv(a, b);
    expose(u);
    u->rev = !u->rev;
    expose(v);
    if (v->right != u || u->left != nullptr) {
      throw std::runtime_error("Cannot cut edge that does not exist.");
    }
    v->right->parent = nullptr;
    v->right = nullptr;
    num_trees++;
  }

  T query(int a, int b) {
    if (!is_connected(a, b)) {
      throw std::runtime_error("Cannot query nodes that are not connected.");
    }
    get_uv(a, b);
    expose(u);
    u->rev = !u->rev;
    expose(v);
    return v->get_subtree_value();
  }

  void update(int a, int b, const T &d) {
    if (!is_connected(a, b)) {
      throw std::runtime_error("Cannot update nodes that are not connected.");
    }
    get_uv(a, b);
    expose(u);
    u->rev = !u->rev;
    expose(v);
    v->delta = v->pending ? compose_deltas(v->delta, d) : d;
    v->pending = true;
  }
};
\end{lstlisting}
\Needspace*{4\baselineskip}
\textbf{Example Usage}\par
\begin{lstlisting}[language=C++]
#include <cassert>
using namespace std;

int main() {
  LinkCutForest<int> lcf;
  lcf.make_root(0, 10);
  lcf.make_root(1, 40);
  lcf.make_root(2, 20);
  lcf.make_root(3, 10);
  lcf.make_root(4, 30);
  assert(lcf.size() == 5);
  assert(lcf.trees() == 5);
  lcf.link(0, 1);
  lcf.link(1, 2);
  lcf.link(2, 3);
  lcf.link(2, 4);
  assert(lcf.trees() == 1);

  // v=10      v=40      v=20      v=10
  //  0---------1---------2---------3
  //                      |
  //                      ----------4
  //                               v=30
  assert(lcf.query(1, 4) == 20);
  lcf.update(1, 1, 100);
  lcf.update(2, 4, 100);

  // v=10     v=100     v=100      v=10
  //  0---------1---------2---------3
  //                      |
  //                      ----------4
  //                              v=100
  assert(lcf.query(4, 4) == 100);
  assert(lcf.query(0, 4) == 10);
  assert(lcf.query(3, 4) == 10);
  lcf.cut(1, 2);

  // v=10     v=100     v=100      v=0
  //  0---------1         2---------3
  //                      |
  //                      ----------4
  //                              v=100
  assert(lcf.trees() == 2);
  assert(!lcf.is_connected(1, 2));
  assert(!lcf.is_connected(0, 4));
  assert(lcf.is_connected(2, 3));
  return 0;
}
\end{lstlisting}

\chapter{Optimization}

\section{Search on Answer}
\setcounter{section}{1}
\setcounter{subsection}{0}
\subsection{Binary Search}
Binary search can be generally used to find the input value corresponding to any output value of a monotonic (strictly non-increasing or strictly non-decreasing) function in $O(\log n)$ time with respect to the domain size. This is a special case of finding the exact point at which any given monotonic Boolean function changes from true to false or vice versa. Unlike searching through an array, discrete binary search is not restricted by available memory, making it useful for handling infinitely large search spaces such as real number intervals.

\begin{compactitem}
  \item \inlinecode{binary\_search\_first\_true()} takes two integers \inlinecode{lo} and \inlinecode{hi} as boundaries for the search space \inlinecode{[lo, hi)} (i.e. including \inlinecode{lo}, but excluding \inlinecode{hi}) and returns the smallest integer \inlinecode{k} in \inlinecode{[lo, hi)} for which the predicate \inlinecode{pred(k)} tests true. If \inlinecode{pred(k)} tests false for every \inlinecode{k} in \inlinecode{[lo, hi)}, then \inlinecode{hi} is returned. This function must be used on a range in which there exists a constant \inlinecode{k} such that \inlinecode{pred(x)} tests false for every \inlinecode{x} in \inlinecode{[lo, k)} and true for every \inlinecode{x} in \inlinecode{[k, hi)}.
  \item \inlinecode{binary\_search\_last\_true()} takes two integers \inlinecode{lo} and \inlinecode{hi} as boundaries for the search space \inlinecode{[lo, hi)} (i.e. including \inlinecode{lo}, but excluding \inlinecode{hi}) and returns the largest integer \inlinecode{k} in \inlinecode{[lo, hi)} for which the predicate \inlinecode{pred(k)} tests true. If \inlinecode{pred(k)} tests false for every \inlinecode{k} in \inlinecode{[lo, hi)}, then \inlinecode{hi} is returned. This function must be used on a range in which there exists a constant \inlinecode{k} such that \inlinecode{pred(x)} tests true for every \inlinecode{x} in \inlinecode{[lo, k]} and false for every \inlinecode{x} in \inlinecode{(k, hi)}.
  \item \inlinecode{fbinary\_search()} is the equivalent of \inlinecode{binary\_search\_first\_true()} on floating point predicates. Since any interval of real numbers is dense, the exact target cannot be found due to floating point error. Instead, a value that is very close to the border between false and true is returned. The precision of the answer depends on the number of repetitions the function performs. Since each repetition bisects the search space, the absolute error of the answer is $1/(2^n)$ times the distance between \inlinecode{lo} and \inlinecode{hi} after n repetitions. Although it is possible to control the error by looping while \inlinecode{hi - lo} is greater than an arbitrary epsilon, it is simpler to let the loop run for a desired number of iterations until floating point arithmetic break down. 100 iterations is usually sufficient, since the search space will be reduced to $2^{-100}$ (roughly $10^{-30}$) times its original size. This implementation can be modified to find the "last true" point in the range by simply interchanging the assignments of \inlinecode{lo} and \inlinecode{hi} in the if-else statements.
\end{compactitem}

\textbf{Time Complexity}:
\begin{compactitem}
  \item $O(\log n)$ calls will be made to \inlinecode{pred()} in \inlinecode{binary\_search\_first\_true()} and \inlinecode{binary\_search\_last\_true()}, where $n$ is the distance between \inlinecode{lo} and \inlinecode{hi}.
  \item $O(n)$ calls will be made to \inlinecode{pred()} in \inlinecode{fbinary\_search()}, where $n$ is the number of iterations performed.
\end{compactitem}

\textbf{Space Complexity}: $O(1)$ auxiliary for all operations.

\begin{lstlisting}[language=C++]
template<class Int, class IntPredicate>
Int binary_search_first_true(Int lo, Int hi, IntPredicate pred) {  // 000[1]11
  Int mid, _hi = hi;
  while (lo < hi) {
    mid = lo + (hi - lo)/2;
    if (pred(mid)) {
      hi = mid;
    } else {
      lo = mid + 1;
    }
  }
  if (!pred(lo)) {
    return _hi;  // All false.
  }
  return lo;
}

template<class Int, class IntPredicate>
Int binary_search_last_true(Int lo, Int hi, IntPredicate pred) {  // 11[1]000
  Int mid, _hi = hi--;
  while (lo < hi) {
    mid = lo + (hi - lo + 1)/2;
    if (pred(mid)) {
      lo = mid;
    } else {
      hi = mid - 1;
    }
  }
  if (!pred(lo)) {
    return _hi;  // All false.
  }
  return lo;
}

template<class DoublePredicate>
double fbinary_search(double lo, double hi, DoublePredicate pred) {  // 000[1]11
  double mid;
  for (int i = 0; i < 100; i++) {
    mid = (lo + hi)/2.0;
    if (pred(mid)) {
      hi = mid;
    } else {
      lo = mid;
    }
  }
  return lo;
}
\end{lstlisting}
\textbf{Example Usage}\par
\begin{lstlisting}[language=C++]
#include <cassert>
#include <cmath>

// Simple predicate examples.
bool pred1(int x) { return x >= 3; }
bool pred2(int x) { return false; }
bool pred3(int x) { return x <= 5; }
bool pred4(int x) { return true; }
bool pred5(double x) { return x >= 1.2345; }

int main() {
  assert(binary_search_first_true(0, 7, pred1) == 3);
  assert(binary_search_first_true(0, 7, pred2) == 7);
  assert(binary_search_last_true(0, 7, pred3)  == 5);
  assert(binary_search_last_true(0, 7, pred4)  == 6);
  assert(fabs(fbinary_search(-10.0, 10.0, pred5) - 1.2345) < 1e-15);
  return 0;
}
\end{lstlisting}

\section{Unimodal and Continuous Search}
\setcounter{section}{2}
\setcounter{subsection}{0}
\subsection{Ternary Search}
Given a unimodal function \inlinecode{f(x)} taking a single \inlinecode{double} argument, find its global maximum or minimum point to a specified absolute error.

\inlinecode{ternary\_search\_min()} takes the domain \inlinecode{[lo, hi]} of a continuous function \inlinecode{f(x)} and returns a number \inlinecode{x} such that \inlinecode{f} is strictly decreasing on the interval \inlinecode{[lo, x]} and strictly increasing on the interval \inlinecode{[x, hi]}. For the function to be correct and deterministic, such an \inlinecode{x} must exist and be unique.

\inlinecode{ternary\_search\_max()} takes the domain \inlinecode{[lo, hi]} of a continuous function \inlinecode{f(x)} and returns a number \inlinecode{x} such that \inlinecode{f} is strictly increasing on the interval \inlinecode{[lo, x]} and strictly decreasing on the interval \inlinecode{[x, hi]}. For the function to be correct and deterministic, such an \inlinecode{x} must exist and be unique.

\textbf{Time Complexity}: $O(\log n)$ calls will be made to \inlinecode{f()}, where $n$ is the distance between \inlinecode{lo} and \inlinecode{hi} divided by the specified absolute error (epsilon).

\textbf{Space Complexity}: $O(1)$ auxiliary for both operations.

\begin{lstlisting}[language=C++]
template<class UnimodalFunction>
double ternary_search_min(double lo, double hi, UnimodalFunction f,
                          const double EPS = 1e-12) {
  double lthird, hthird;
  while (hi - lo > EPS) {
    lthird = lo + (hi - lo)/3;
    hthird = hi - (hi - lo)/3;
    if (f(lthird) < f(hthird)) {
      hi = hthird;
    } else {
      lo = lthird;
    }
  }
  return lo;
}

template<class UnimodalFunction>
double ternary_search_max(double lo, double hi, UnimodalFunction f,
                          const double EPS = 1e-12) {
  double lthird, hthird;
  while (hi - lo > EPS) {
    lthird = lo + (hi - lo)/3;
    hthird = hi - (hi - lo)/3;
    if (f(lthird) < f(hthird)) {
      lo = lthird;
    } else {
      hi = hthird;
    }
  }
  return hi;
}
\end{lstlisting}
\textbf{Example Usage}\par
\begin{lstlisting}[language=C++]
#include <cassert>
#include <cmath>

bool equal(double a, double b) {
  return fabs(a - b) < 1e-7;
}

// Parabola opening up with vertex at (-2, -24).
double f1(double x) {
  return 3*x*x + 12*x - 12;
}

// Parabola opening down with vertex at (2/19, 8366/95) ~ (0.105, 88.063).
double f2(double x) {
  return -5.7*x*x + 1.2*x + 88;
}

// Absolute value function shifted to the right by 30 units.
double f3(double x) {
  return fabs(x - 30);
}

int main() {
  assert(equal(ternary_search_min(-1000, 1000, f1), -2));
  assert(equal(ternary_search_max(-1000, 1000, f2), 2.0/19));
  assert(equal(ternary_search_min(-1000, 1000, f3), 30));
  return 0;
}
\end{lstlisting}
\subsection{Hill Climbing}
Given a continuous function \inlinecode{f(x, y)} to \inlinecode{double} and a (possibly arbitrary) initial guess (\inlinecode{x0}, \inlinecode{y0}), return a potential global minimum found through hill-climbing. Optionally, two \inlinecode{double} pointers \inlinecode{critical\_x} and \inlinecode{critical\_y} may be passed to store the input points to \inlinecode{f} at which the returned minimum value is attained.

Hill-climbing is a heuristic which starts at the guess, then considers taking a single step in each of a fixed number of directions. The direction with the best (in this case, minimum) value is chosen, and further steps are taken in it until the answer no longer improves. When this happens, the step size is reduced and the same process repeats until a desired absolute error is reached. The technique's success heavily depends on the behavior of \inlinecode{f} and the initial guess. Therefore, the result is not guaranteed to be the global minimum.

\textbf{Time Complexity}: $O(d \log n)$ call will be made to \inlinecode{f}, where $d$ is the number of directions considered at each position and $n$ is the search space that is approximately proportional to the maximum possible step size divided by the minimum possible step size.

\textbf{Space Complexity}: $O(1)$ auxiliary.

\begin{lstlisting}[language=C++]
#include <cstddef>
#include <cmath>

template<class ContinuousFunction>
double find_min(ContinuousFunction f, double x0, double y0,
                double *critical_x = NULL, double *critical_y = NULL,
                const double STEP_MIN = 1e-9, const double STEP_MAX = 1e6,
                const int NUM_DIRECTIONS = 6) {
  static const double PI = acos(-1.0);
  double x = x0, y = y0, res = f(x0, y0);
  for (double step = STEP_MAX; step > STEP_MIN; ) {
    double best = res, best_x = x, best_y = y;
    bool found = false;
    for (int i = 0; i < NUM_DIRECTIONS; i++) {
      double a = 2.0*PI*i / NUM_DIRECTIONS;
      double x2 = x + step*cos(a), y2 = y + step*sin(a);
      double value = f(x2, y2);
      if (best > value) {
        best_x = x2;
        best_y = y2;
        best = value;
        found = true;
      }
    }
    if (!found) {
      step /= 2.0;
    } else {
      x = best_x;
      y = best_y;
      res = best;
    }
  }
  if (critical_x != NULL && critical_y != NULL) {
    *critical_x = x;
    *critical_y = y;
  }
  return res;
}
\end{lstlisting}
\textbf{Example Usage}\par
\begin{lstlisting}[language=C++]
#include <cassert>
#include <cmath>

bool eq(double a, double b) {
  return fabs(a - b) < 1e-8;
}

// Paraboloid with global minimum at f(2, 3) = 0.
double f(double x, double y) {
  return (x - 2)*(x - 2) + (y - 3)*(y - 3);
}

int main() {
  double x, y;
  assert(eq(find_min(f, 0, 0, &x, &y), 0));
  assert(eq(x, 2) && eq(y, 3));
  return 0;
}
\end{lstlisting}

\section{Dynamic Programming Optimization}
\setcounter{section}{3}
\setcounter{subsection}{0}
\subsection{Convex Hull Trick (Semi-Dynamic)}
Given a set of pairs $(m, b)$ specifying lines of the form $y = mx + b$, process a set of x-coordinate queries each asking to find the minimum y-value when any of the given lines are evaluated at the specified x. For each \inlinecode{add\_line(m, b)} call, \inlinecode{m} must be the minimum \inlinecode{m} of all lines added so far. For each \inlinecode{query(x)} call, \inlinecode{x} must be the maximum \inlinecode{x} of all queries made so far.

The following implementation is a concise, semi-dynamic version of the convex hull optimization technique. It supports an an interlaced sequence of \inlinecode{add\_line()} and \inlinecode{query()} calls, as long as the preconditions of descending \inlinecode{m} and ascending \inlinecode{x} are satisfied. As a result, it may be necessary to sort the lines and queries before calling the functions. In that case, the overall time complexity will be dominated by the sorting step.

\textbf{Time Complexity}: $O(n)$ for any interlaced sequence of \inlinecode{add\_line()} and \inlinecode{query()} calls, where $n$ is the number of lines added. This is because the overall number of steps taken by \inlinecode{add\_line()} and \inlinecode{query()} are respectively bounded by the number of lines. Thus a single call to either \inlinecode{add\_line()} or \inlinecode{query()} will have an amortized $O(1)$ running time.

\textbf{Space Complexity}:
\begin{compactitem}
  \item $O(n)$ for storage of the lines.
  \item $O(1)$ auxiliary for \inlinecode{add\_line()} and \inlinecode{query()}.
\end{compactitem}

\begin{lstlisting}[language=C++]
#include <vector>

std::vector<long long> M, B;
int ptr = 0;

void add_line(long long m, long long b) {
  int len = M.size();
  while (len > 1 && (B[len - 2] - B[len - 1])*(m - M[len - 1]) >=
                    (B[len - 1] - b)*(M[len - 1] - M[len - 2])) {
    len--;
  }
  M.resize(len);
  B.resize(len);
  M.push_back(m);
  B.push_back(b);
}

long long query(long long x) {
  if (ptr >= (int)M.size()) {
    ptr = (int)M.size() - 1;
  }
  while (ptr + 1 < (int)M.size() &&
         M[ptr + 1]*x + B[ptr + 1] <= M[ptr]*x + B[ptr]) {
    ptr++;
  }
  return M[ptr]*x + B[ptr];
}
\end{lstlisting}
\textbf{Example Usage}\par
\begin{lstlisting}[language=C++]
#include <cassert>

int main() {
  add_line(3, 0);
  add_line(2, 1);
  add_line(1, 2);
  add_line(0, 6);
  assert(query(0) == 0);
  assert(query(1) == 3);
  assert(query(2) == 4);
  assert(query(3) == 5);
  return 0;
}
\end{lstlisting}
\subsection{Convex Hull Trick (Fully-Dynamic)}
Given a set of pairs $(m, b)$ specifying lines of the form $y = mx + b$, process a set of x-coordinate queries each asking to find the minimum y-value when any of the given lines are evaluated at the specified x. To instead have the queries optimize for maximum y-value, call the constructor with \inlinecode{query\_max=true}.

The following implementation is a fully dynamic variant of the convex hull optimization technique, using a self-balancing binary search tree (\inlinecode{std::set}) to support the ability to call \inlinecode{add\_line()} and \inlinecode{query()} in any desired order.

\textbf{Time Complexity}: $O(\log n)$ amortized per call to \inlinecode{add\_line()} and $O(\log n)$ per call to \inlinecode{query()}, where $n$ is the number of lines added.

\textbf{Space Complexity}:
\begin{compactitem}
  \item $O(n)$ for storage of the lines.
  \item $O(1)$ auxiliary for \inlinecode{add\_line()} and \inlinecode{query()}.
\end{compactitem}

\begin{lstlisting}[language=C++]
#include <cassert>
#include <climits>
#include <set>

class hull_optimizer {
  struct line {
    long long m, b;
    mutable long long xhi;
    bool is_query;

    line(long long m, long long b, long long xhi = 0, bool is_query = false)
        : m(m), b(b), xhi(xhi), is_query(is_query) {}

    bool operator<(const line &l) const {
      return l.is_query ? xhi < l.xhi : m < l.m;
    }
  };

  std::multiset<line> hull;
  bool query_max;

  typedef std::multiset<line>::iterator hulliter;

  static long long div_floor(long long a, long long b) {
    return a/b - ((a ^ b) < 0 && a % b);
  }

  bool update_border(hulliter x, hulliter y) {
    if (y == hull.end()) {
      x->xhi = LLONG_MAX;
      return false;
    }
    if (x->m == y->m) {
      x->xhi = (x->b > y->b) ? LLONG_MAX : LLONG_MIN;
    } else {
      x->xhi = div_floor(y->b - x->b, x->m - y->m);
    }
    return x->xhi >= y->xhi;
  }

 public:
  hull_optimizer(bool query_max = false) : query_max(query_max) {}

  void add_line(long long m, long long b) {
    if (!query_max) {
      m = -m;
      b = -b;
    }
    hulliter z = hull.insert(line(m, b));
    hulliter y = z++;
    hulliter x = y;
    while (update_border(y, z)) {
      z = hull.erase(z);
    }
    if (x != hull.begin() && update_border(--x, y)) {
      update_border(x, y = hull.erase(y));
    }
    while ((y = x) != hull.begin() && (--x)->xhi >= y->xhi) {
      update_border(x, hull.erase(y));
    }
  }

  long long query(long long x) const {
    assert(!hull.empty());
    line q(0, 0, x, true);
    hulliter it = hull.lower_bound(q);
    long long res = it->m*x + it->b;
    return query_max ? res : -res;
  }
};
\end{lstlisting}
\textbf{Example Usage}\par
\begin{lstlisting}[language=C++]
#include <cassert>

int main() {
  hull_optimizer h;
  h.add_line(3, 0);
  h.add_line(0, 6);
  h.add_line(1, 2);
  h.add_line(2, 1);
  assert(h.query(0) == 0);
  assert(h.query(2) == 4);
  assert(h.query(1) == 3);
  assert(h.query(3) == 5);
  return 0;
}
\end{lstlisting}

\section{Linear and Convex Optimization}
\setcounter{section}{4}
\setcounter{subsection}{0}
\subsection{Linear Programming (Simplex)}
Solves a linear programming problem using Dantzig's simplex algorithm. The canonical form of a linear programming problem is to maximize (or minimize) the dot product $cx$, subject to $ax \leq b$ and $x \geq 0$, where $x$ is a vector of unknowns to be solved, $c$ is a vector of coefficients, $a$ is a matrix of linear equation coefficients, and $b$ is a vector of boundary coefficients.

\begin{compactitem}
  \item \inlinecode{simplex\_solve(a, b, c, \&x)} solves the linear programming problem for an $m$ by $n$ matrix \inlinecode{a} of real values, a length $m$ vector \inlinecode{b}, and a length $n$ vector \inlinecode{c}, returning 0 if a solution was found or $-1$ if there are no solutions. If a solution is found, then the vector pointed to by \inlinecode{x} is populated with the solution vector of length $n$.
\end{compactitem}

\textbf{Time Complexity}: Polynomial (average) on the number of equations and unknowns, but exponential in the worst case.

\textbf{Space Complexity}: $O(m \cdot n)$ auxiliary heap space.

\begin{lstlisting}[language=C++]
#include <cmath>
#include <limits>
#include <vector>

template<class Matrix>
int simplex_solve(const Matrix &a, const std::vector<double> &b,
                  const std::vector<double> &c, std::vector<double> *x,
                  const bool MAXIMIZE = true, const double EPS = 1e-10) {
  int m = a.size(), n = c.size();
  Matrix t(m + 2, std::vector<double>(n + 2));
  t[1][1] = 0;
  for (int j = 1; j <= n; j++) {
    t[1][j + 1] = MAXIMIZE ? c[j - 1] : -c[j - 1];
  }
  for (int i = 1; i <= m; i++) {
    for (int j = 1; j <= n; j++) {
      t[i + 1][j + 1] = -a[i - 1][j - 1];
    }
    t[i + 1][1] = b[i - 1];
  }
  for (int j = 1; j <= n; j++) {
    t[0][j + 1] = j;
  }
  for (int i = n + 1; i <= m + n; i++) {
    t[i - n + 1][0] = i;
  }
  double p1 = 0, p2 = 0;
  bool done = true;
  do {
    double mn = std::numeric_limits<double>::max(), xmax = 0, v;
    for (int j = 2; j <= n + 1; j++) {
      if (t[1][j] > 0 && t[1][j] > xmax) {
        p2 = j;
        xmax = t[1][j];
      }
    }
    for (int i = 2; i <= m + 1; i++) {
      v = fabs(t[i][1] / t[i][p2]);
      if (t[i][p2] < 0 && mn > v) {
        mn = v;
        p1 = i;
      }
    }
    std::swap(t[p1][0], t[0][p2]);
    for (int i = 1; i <= m + 1; i++) {
      if (i != p1) {
        for (int j = 1; j <= n + 1; j++) {
          if (j != p2) {
            t[i][j] -= t[p1][j]*t[i][p2] / t[p1][p2];
          }
        }
      }
    }
    t[p1][p2] = 1.0 / t[p1][p2];
    for (int j = 1; j <= n + 1; j++) {
      if (j != p2) {
        t[p1][j] *= fabs(t[p1][p2]);
      }
    }
    for (int i = 1; i <= m + 1; i++) {
      if (i != p1) {
        t[i][p2] *= t[p1][p2];
      }
    }
    for (int i = 2; i <= m + 1; i++) {
      if (t[i][1] < 0) {
        return -1;
      }
    }
    done = true;
    for (int j = 2; j <= n + 1; j++) {
      if (t[1][j] > 0) {
        done = false;
      }
    }
  } while (!done);
  x->clear();
  for (int j = 1; j <= n; j++) {
    for (int i = 2; i <= m + 1; i++) {
      if (fabs(t[i][0] - j) < EPS) {
        x->push_back(t[i][1]);
      }
    }
  }
  return 0;
}
\end{lstlisting}
\textbf{Example Usage}\par
\begin{lstlisting}[language=C++]
#include <cassert>
#include <iostream>
using namespace std;

int main() {
  // Solve [x, y] that maximizes 3x + 4y, subject to x, y >= 0 and:
  //  -2x +    1y <=  0
  //   1x + 0.85y <=  9
  //   1x +    2y <= 14
  const int equations = 3, unknowns = 2;
  double a[equations][unknowns] = {{-2, 1}, {1, 0.85}, {1, 2}};
  double b[equations] = {0, 9, 14};
  double c[unknowns] = {3, 4};
  vector<vector<double>> va(equations, vector<double>(unknowns));
  vector<double> vb(b, b + equations), vc(c, c + unknowns), x;
  for (int i = 0; i < equations; i++) {
    for (int j = 0; j < unknowns; j++) {
      va[i][j] = a[i][j];
    }
  }
  assert(simplex_solve(va, vb, vc, &x) == 0);
  double maxval = 0;
  for (int i = 0; i < (int)x.size(); i++) {
    maxval += c[i]*x[i];
  }
  cout << "Solution = " << maxval << " at (" << x[0];
  for (int i = 1; i < (int)x.size(); i++) {
    cout << ", " << x[i];
  }
  cout << ")." << endl;
  return 0;
}
\end{lstlisting}
\begin{exampleoutput}
Solution = 33.3043 at (5.30435, 4.34783).
\end{exampleoutput}

\input{chapter4}
\chapter{\texorpdfstring{Optimization}{5. Optimization}}

\section{\texorpdfstring{Monotone Predicate Search}{5.1 Monotone Predicate Search}}
\setcounter{section}{1}
\setcounter{subsection}{0}
\subsection{Binary Search}
Binary search can be generally used to find the input value corresponding to any output value of a monotonic (strictly non-increasing or strictly non-decreasing) function in $O(\log n)$ time with respect to the domain size. This is a special case of finding the exact point at which any given monotonic Boolean function changes from true to false or vice versa. Each step evaluates the midpoint of the current interval and keeps the half that still contains the transition point, halving the search space until it closes on the answer. Unlike searching through an array, discrete binary search is not restricted by available memory, making it useful for handling infinitely large search spaces such as real number intervals.

\begin{compactitem}
  \item \inlinecode{binary\_search\_first\_true(lo, hi, pred)} takes signed integer boundaries for the search space $[{\inlinecodemath{lo}}, {\inlinecodemath{hi}})$ (i.e. including \inlinecode{lo}, but excluding \inlinecode{hi}) and returns the smallest integer \inlinecode{k} in $[{\inlinecodemath{lo}}, {\inlinecodemath{hi}})$ for which the predicate \inlinecode{pred(k)} tests true. If \inlinecode{pred(k)} tests false for the entire input range, then \inlinecode{hi} is returned. The caller must ensure \inlinecode{pred} is monotonic on the input range, i.e. returning all false for some (possibly empty) prefix, followed by all true in some (possibly empty) suffix. E.g., patterns \inlinecode{01}, \inlinecode{00}, and \inlinecode{11} are allowed, but \inlinecode{10} is disallowed.
  \item \inlinecode{binary\_search\_last\_true(lo, hi, pred)} takes signed integer boundaries for the search space $[{\inlinecodemath{lo}}, {\inlinecodemath{hi}})$ (i.e. including \inlinecode{lo}, but excluding \inlinecode{hi}) and returns the largest integer \inlinecode{k} in $[{\inlinecodemath{lo}}, {\inlinecodemath{hi}})$ for which the predicate \inlinecode{pred(k)} tests true. If \inlinecode{pred(k)} tests false for the entire input range, then the original \inlinecode{lo - 1} is returned, so \inlinecode{lo} must exceed the minimum representable value. The caller must ensure \inlinecode{pred} is monotonic on the input range, i.e. returning all true for some (possibly empty) prefix, followed by all false in some (possibly empty) suffix. E.g., patterns \inlinecode{10}, \inlinecode{00}, and \inlinecode{11} are allowed, but \inlinecode{01} is disallowed.
  \item \inlinecode{fbinary\_search(lo, hi, pred)} is the equivalent of \inlinecode{binary\_search\_first\_true()} on floating point predicates. Since any interval of real numbers is dense, the exact target cannot be found due to floating point error. Instead, a value that is very close to the border between false and true is returned. The precision of the answer depends on the number of repetitions the function performs. Since each repetition bisects the search space, the absolute error of the answer is $1/(2^n)$ times the distance between \inlinecode{lo} and \inlinecode{hi} after $n$ repetitions. Although the error can be controlled by looping until the distance shrinks to an arbitrary epsilon, it is simpler to let the loop run for a desired number of iterations until floating point arithmetic breaks down. $100$ iterations is usually sufficient, since the search space will be reduced to $2^{-100}$ (roughly $10^{-30}$) times its original size. This implementation can be modified to find the ``last true'' point by simply interchanging the assignments of \inlinecode{lo} and \inlinecode{hi} in the if-else statements.
\end{compactitem}

\Needspace*{3\baselineskip}
\textbf{Time Complexity}:
\begin{compactitem}
  \item $O(\log n)$ calls will be made to \inlinecode{pred()} in \inlinecode{binary\_search\_first\_true()} and \inlinecode{binary\_search\_last\_true()}, where $n$ is the distance between \inlinecode{lo} and \inlinecode{hi}.
  \item $O(n)$ calls will be made to \inlinecode{pred()} in \inlinecode{fbinary\_search()}, where $n$ is the number of iterations performed.
\end{compactitem}

\textbf{Space Complexity}: $O(1)$ auxiliary for all operations.

\begin{lstlisting}[language=C++]
template<typename Int, typename Pred>
Int binary_search_first_true(Int lo, Int hi, Pred pred) {  // 000[1]11
  while (lo < hi) {
    Int mid = lo + (hi - lo) / 2;
    if (pred(mid)) {
      hi = mid;
    } else {
      lo = mid + 1;
    }
  }
  return lo;  // hi if all false
}

template<typename Int, typename Pred>
Int binary_search_last_true(Int lo, Int hi, Pred pred) {  // 11[1]000
  while (lo < hi) {
    Int mid = lo + (hi - lo) / 2;
    if (pred(mid)) {
      lo = mid + 1;
    } else {
      hi = mid;
    }
  }
  return lo - 1;
}

template<typename Pred>
double fbinary_search(double lo, double hi, Pred pred) {  // 000[1]11
  double mid;
  for (int i = 0; i < 100; i++) {
    mid = (lo + hi) / 2.0;
    if (pred(mid)) {
      hi = mid;
    } else {
      lo = mid;
    }
  }
  return lo;
}
\end{lstlisting}
\Needspace*{4\baselineskip}
\textbf{Example Usage}\par
\begin{lstlisting}[language=C++]
#include <cassert>
#include <cmath>

int main() {
  assert(binary_search_first_true(0, 7, [](int x) { return x >= 3; }) == 3);
  assert(binary_search_first_true(0, 7, [](int x) { return true; }) == 0);
  assert(binary_search_first_true(0, 7, [](int x) { return false; }) == 7);
  assert(binary_search_first_true(4, 4, [](int x) { return x >= 4; }) == 4);

  assert(binary_search_last_true(0, 7, [](int x) { return x <= 5; }) == 5);
  assert(binary_search_last_true(0, 7, [](int x) { return true; }) == 6);
  assert(binary_search_last_true(0, 7, [](int x) { return false; }) == -1);
  assert(binary_search_last_true(4, 4, [](int x) { return x <= 4; }) == 3);

  double res = fbinary_search(-10.0, 10.0, [](double x) { return x >= 1.2345; });
  assert(fabs(res - 1.2345) < 1e-15);
  return 0;
}
\end{lstlisting}
\subsection{Exponential Search}
Finds the first value where a monotone Boolean predicate becomes true when an upper bound is not known in advance. Exponential search first grows the search range by powers of two until it brackets the transition point, then finishes with ordinary binary search.

This is useful for answer-search problems over unbounded or very large integer domains. The predicate must eventually become true, or the search may overflow or run forever unless an explicit limit is added.

\begin{compactitem}
  \item \inlinecode{exponential\_search\_first\_true(lo, pred)} returns the smallest integer \inlinecode{x} greater than or equal to \inlinecode{lo} such that \inlinecode{pred(x)} is true.
\end{compactitem}

\textbf{Time Complexity}: $O(\log n)$ calls to \inlinecode{pred()}, where $n$ is the distance from \inlinecode{lo} to the first true value.

\textbf{Space Complexity}: $O(1)$ auxiliary.

\begin{lstlisting}[language=C++]
#include <cstdint>

template<typename Int, typename Pred>
Int exponential_search_first_true(Int lo, Pred pred) {  // 000[1]11
  if (pred(lo)) {
    return lo;
  }
  Int step = 1;
  while (!pred(lo + step)) {
    step *= 2;
  }
  Int hi = lo + step;
  lo += step / 2;
  while (lo < hi) {
    Int mid = lo + (hi - lo) / 2;
    if (pred(mid)) {
      hi = mid;
    } else {
      lo = mid + 1;
    }
  }
  return lo;
}
\end{lstlisting}
\Needspace*{4\baselineskip}
\textbf{Example Usage}\par
\begin{lstlisting}[language=C++]
#include <cassert>

int main() {
  auto at_least_1000 = [](int x) { return x >= 1000; };
  assert(exponential_search_first_true(1000, at_least_1000) == 1000);
  assert(exponential_search_first_true(0, at_least_1000) == 1000);
  assert(exponential_search_first_true(5, at_least_1000) == 1000);
  auto square_large_enough = [](int64_t x) { return x * x >= 123456789LL; };
  assert(exponential_search_first_true(0LL, square_large_enough) == 11112);
  return 0;
}
\end{lstlisting}
\subsection{Local Minimum (Binary Search)}
Finds a non-strict local minimum in an array, that is, an index $i$ whose value is no greater than the values at its existing neighbors. Unlike ternary search, this does not require the array to be unimodal, but the result is only a local minimum and need not be the global minimum.

Compare the middle element with its right neighbor. If the sequence slopes downward, the right half contains a local minimum; otherwise, the middle element or the left half contains one. Each comparison therefore discards half of the remaining indices.

\begin{compactitem}
  \item \inlinecode{local\_min\_index(a)} returns the index of a non-strict local minimum in the nonempty array \inlinecode{a}. Either endpoint may be returned when it is no greater than its sole neighbor.
\end{compactitem}

\textbf{Time Complexity}: $O(\log n)$ comparisons, where $n$ is the size of \inlinecode{a}.

\textbf{Space Complexity}: $O(1)$ auxiliary.

\begin{lstlisting}[language=C++]
#include <cassert>
#include <vector>

template<typename T>
int local_min_index(const std::vector<T> &a) {
  assert(!a.empty());
  int lo = 0, hi = static_cast<int>(a.size()) - 1;
  while (lo < hi) {
    int mid = lo + (hi - lo) / 2;
    if (a[mid] > a[mid + 1]) {
      lo = mid + 1;
    } else {
      hi = mid;
    }
  }
  return lo;
}
\end{lstlisting}
\Needspace*{4\baselineskip}
\textbf{Example Usage}\par
\begin{lstlisting}[language=C++]
int main() {
  assert(local_min_index(std::vector<int>{9, 7, 3, 4, 8}) == 2);
  assert(local_min_index(std::vector<int>{1, 2, 3, 4}) == 0);

  std::vector<int> duplicates{4, 2, 2, 3};
  int i = local_min_index(duplicates);
  assert((i == 1 || i == 2) && duplicates[i] == 2);
  return 0;
}
\end{lstlisting}

\section{\texorpdfstring{Unimodal and Continuous Search}{5.2 Unimodal and Continuous Search}}
\setcounter{section}{2}
\setcounter{subsection}{0}
\subsection{Ternary and Golden Section Search}
Given a unimodal function $f$, find the position of its global minimum or maximum: a value $x$ such that the function strictly decreases up to $x$ and strictly increases after it (for a minimum), or the reverse (for a maximum). All three searches below exploit the same idea: evaluate two interior points, and the comparison reveals one side of the interval that cannot contain the optimum, so it can be discarded. They differ only in how the probes are placed and whether the domain is continuous or discrete. Each maximum routine is the minimum routine applied to the negated function, since the location of a maximum of $f$ is the location of a minimum of $-f$.

Ternary search splits the interval into thirds. It is the simplest to reason about, but it recomputes both probe points each step, spending two function evaluations per iteration while the interval shrinks to two-thirds.

Golden-section search places the probes at the golden ratio so that one probe is reused as an interior point of the next, smaller interval. It therefore performs only one new evaluation per iteration while shrinking the interval by a factor of $1/\varphi$ (about 0.618), and reaches a given accuracy in roughly 2.4 times fewer evaluations than ternary search. Prefer it when the function is expensive to evaluate.

The discrete version operates on an integer domain, narrowing by thirds until at most three candidates remain and then scanning them. It returns the index of the optimum rather than a real coordinate. As with ternary search on reals, the function must be strictly unimodal.

\begin{compactitem}
  \item \inlinecode{ternary\_search\_min(lo, hi, f, eps = 1e-12)} and \inlinecode{ternary\_search\_max(lo, hi, f, eps = 1e-12)} return a \inlinecode{double} \inlinecode{x} in $[{\inlinecodemath{lo}}, {\inlinecodemath{hi}}]$ optimizing the continuous unimodal function \inlinecode{f}, to within the optional absolute error \inlinecode{eps}.
  \item \inlinecode{golden\_section\_min(lo, hi, f, eps = 1e-12)} and \inlinecode{golden\_section\_max(lo, hi, f, eps = 1e-12)} do the same with one function evaluation per iteration.
  \item \inlinecode{discrete\_ternary\_min(lo, hi, f)} and \inlinecode{discrete\_ternary\_max(lo, hi, f)} return the integer index in range $[{\inlinecodemath{lo}}, {\inlinecodemath{hi}}]$ optimizing the unimodal function \inlinecode{f}, breaking ties toward the smaller index.
\end{compactitem}

\Needspace*{3\baselineskip}
\textbf{Time Complexity}:
\begin{compactitem}
  \item $O(\log\left(n / {\inlinecodemath{eps}}\right))$ calls to \inlinecode{f()} for the continuous searches, where $n$ is the distance between \inlinecode{lo} and \inlinecode{hi}. Golden-section makes about half as many per iteration as ternary search.
  \item $O(\log n)$ calls to \inlinecode{f()} for the discrete searches.
\end{compactitem}

\textbf{Space Complexity}: $O(1)$ auxiliary for all operations.

\begin{lstlisting}[language=C++]
template<typename Fn>
double ternary_search_min(double lo, double hi, Fn f, double eps = 1e-12) {
  while (hi - lo > eps) {
    double m1 = lo + (hi - lo) / 3, m2 = hi - (hi - lo) / 3;
    if (f(m1) < f(m2)) {
      hi = m2;
    } else {
      lo = m1;
    }
  }
  return (lo + hi) / 2;
}

template<typename Fn>
double ternary_search_max(double lo, double hi, Fn f, double eps = 1e-12) {
  return ternary_search_min(lo, hi, [&](double x) { return -f(x); }, eps);
}

template<typename Fn>
double golden_section_min(double lo, double hi, Fn f, double eps = 1e-12) {
  const double r = 0.6180339887498949;  // 1 / phi = (sqrt(5) - 1) / 2.
  double m1 = hi - r * (hi - lo), m2 = lo + r * (hi - lo);
  double f1 = f(m1), f2 = f(m2);
  while (hi - lo > eps) {
    if (f1 < f2) {
      hi = m2;
      m2 = m1;
      f2 = f1;
      m1 = hi - r * (hi - lo);
      f1 = f(m1);
    } else {
      lo = m1;
      m1 = m2;
      f1 = f2;
      m2 = lo + r * (hi - lo);
      f2 = f(m2);
    }
  }
  return (lo + hi) / 2;
}

template<typename Fn>
double golden_section_max(double lo, double hi, Fn f, double eps = 1e-12) {
  return golden_section_min(lo, hi, [&](double x) { return -f(x); }, eps);
}

template<typename Int, typename Fn>
Int discrete_ternary_min(Int lo, Int hi, Fn f) {
  while (hi - lo > 2) {
    Int m1 = lo + (hi - lo) / 3, m2 = hi - (hi - lo) / 3;
    if (f(m1) < f(m2)) {
      hi = m2 - 1;
    } else {
      lo = m1 + 1;
    }
  }
  Int best = lo;
  auto best_value = f(best);
  for (Int k = lo + 1; k <= hi; k++) {
    auto value = f(k);
    if (value < best_value) {
      best = k;
      best_value = value;
    }
  }
  return best;
}

template<typename Int, typename Fn>
Int discrete_ternary_max(Int lo, Int hi, Fn f) {
  return discrete_ternary_min(lo, hi, [&](Int x) { return -f(x); });
}
\end{lstlisting}
\Needspace*{4\baselineskip}
\textbf{Example Usage}\par
\begin{lstlisting}[language=C++]
#include <cassert>
#include <cmath>
#include <vector>
using namespace std;

bool EQ(double a, double b) {
  return fabs(a - b) < 1e-6;
}

int main() {
  // Parabola opening up with vertex at x = -2.
  auto up = [](double x) { return 3 * x * x + 12 * x - 12; };
  // Parabola opening down with vertex at x = 2/19.
  auto down = [](double x) { return -5.7 * x * x + 1.2 * x + 88; };

  assert(EQ(ternary_search_min(-1000, 1000, up), -2));
  assert(EQ(golden_section_min(-1000, 1000, up), -2));
  assert(EQ(ternary_search_max(-1000, 1000, down), 2.0 / 19));
  assert(EQ(golden_section_max(-1000, 1000, down), 2.0 / 19));

  // Discrete unimodal arrays: a valley and a peak.
  vector<int> valley{5, 3, 1, 2, 4, 6};
  vector<int> peak{1, 3, 7, 6, 2};
  auto v = [&](int i) { return valley[i]; };
  auto p = [&](int i) { return peak[i]; };
  assert(discrete_ternary_min(0, static_cast<int>(valley.size()) - 1, v) == 2);  // Value 1.
  assert(discrete_ternary_max(0, static_cast<int>(peak.size()) - 1, p) == 2);    // Value 7.
  return 0;
}
\end{lstlisting}
\subsection{2D Hill Climbing}
Given a continuous geometric objective $f(x, y)$ and a (possibly arbitrary) starting guess $(x_0, y_0)$, search for a small value using two-dimensional hill climbing.

The heuristic starts at the guess and evaluates one step in each of eight evenly spaced directions. It moves to the best improving neighbor, checks all eight directions again, and reduces the step size when no neighbor improves the answer. This repeats until the step size falls below the chosen threshold. Success depends heavily on the behavior of $f$ and the initial guess, so the result is not guaranteed to be the global minimum.

\begin{compactitem}
  \item \inlinecode{hill\_climb\_min(f, x0, y0, \&best\_x, \&best\_y, step\_min = 1e-9, step\_max = 1e6)} returns an approximate minimum value of function \inlinecode{f} reached by hill climbing from the starting guess $({\inlinecodemath{x0}}, {\inlinecodemath{y0}})$. If either optional pointer \inlinecode{best\_x} or \inlinecode{best\_y} is supplied, the corresponding coordinate of the point attaining the returned value is stored through it. The search starts with step size \inlinecode{step\_max} and stops below \inlinecode{step\_min}.
\end{compactitem}

\textbf{Time Complexity}: $O(k + s)$ calls to \inlinecode{f()}, where $k$ is the number of times the step size is reduced and $s$ is the total number of improving moves accepted.

\textbf{Space Complexity}: $O(1)$ auxiliary.

\begin{lstlisting}[language=C++]
template<typename Fn>
double hill_climb_min(
    Fn f, double x0, double y0, double *best_x = nullptr, double *best_y = nullptr,
    const double step_min = 1e-9, const double step_max = 1e6
) {
  static const double inv_sqrt2 = 0.7071067811865476;
  static const double dx[] = {1, inv_sqrt2, 0, -inv_sqrt2, -1, -inv_sqrt2, 0, inv_sqrt2};
  static const double dy[] = {0, inv_sqrt2, 1, inv_sqrt2, 0, -inv_sqrt2, -1, -inv_sqrt2};
  double x = x0, y = y0, res = f(x0, y0);
  for (double step = step_max; step > step_min;) {
    double next_value = res, next_x = x, next_y = y;
    bool found = false;
    for (int i = 0; i < 8; i++) {
      double x2 = x + step * dx[i], y2 = y + step * dy[i];
      double value = f(x2, y2);
      if (value < next_value) {
        next_x = x2;
        next_y = y2;
        next_value = value;
        found = true;
      }
    }
    if (!found) {
      step /= 2.0;
    } else {
      x = next_x;
      y = next_y;
      res = next_value;
    }
  }
  if (best_x != nullptr) {
    *best_x = x;
  }
  if (best_y != nullptr) {
    *best_y = y;
  }
  return res;
}
\end{lstlisting}
\Needspace*{4\baselineskip}
\textbf{Example Usage}\par
\begin{lstlisting}[language=C++]
#include <cassert>
#include <cmath>

bool EQ(double a, double b) {
  return fabs(a - b) < 1e-9;
}

// Paraboloid with global minimum at f(2, 3) = 0.
double f(double x, double y) {
  return (x - 2) * (x - 2) + (y - 3) * (y - 3);
}

int main() {
  double x, y;
  assert(EQ(hill_climb_min(f, 0, 0, &x, &y), 0));
  assert(EQ(x, 2) && EQ(y, 3));
  return 0;
}
\end{lstlisting}
\subsection{Simulated Annealing}
Approximately minimizes an energy function $E(s)$ over an arbitrary state space using simulated annealing. This randomized heuristic is useful when the states may be discrete, the energy has many local minima, and gradients or stronger structural properties are unavailable.

Starting from an initial state $s$, each iteration generates a random neighboring state $s'$. An improving move is always accepted, while a worsening move is accepted with probability $\exp(-(E(s') - E(s))/T)$ at temperature $T$. High temperatures encourage exploration across energy barriers; as the temperature decreases, the search increasingly favors improvements. The best state is tracked separately because the current state may later move to a worse one.

\begin{compactitem}
  \item \inlinecode{anneal\_min(initial, energy, rand\_neighbor, rng, ...)} returns \inlinecode{(best\_energy, best\_state)}. The callable \inlinecode{energy(state)} returns the state's energy, while \inlinecode{rand\_neighbor(state, temp, rng)} returns a randomly chosen nearby state. The optional temperature arguments default to \inlinecode{temp\_start = 1000}, \inlinecode{temp\_end = 1e-6}, and \inlinecode{cooling\_rate = 0.995}. \inlinecode{temp\_end} must be positive and less than \inlinecode{temp\_start}, while \inlinecode{cooling\_rate} must be strictly between 0 and 1.
\end{compactitem}

The geometric schedule repeatedly multiplies the temperature by \inlinecode{cooling\_rate}, performing $k = \lceil \log(T_{end}/T_{start}) / \log(r) \rceil$ iterations. Temperature has the same scale as energy: multiplying every energy by a constant requires multiplying both temperature bounds by the same constant to preserve acceptance probabilities. For continuous states, temperature can also control the neighbor step size. Slower cooling, independent restarts, and a time-based stopping condition often improve results, but simulated annealing never proves optimality.

This generic interface copies each proposed state and recomputes its energy. For large permutation states, adapt the loop to mutate and revert moves in place and update the energy from the affected terms instead.

\textbf{Time Complexity}: $O(k(C_E + C_N))$, where $k$ is the number of temperature levels and $C_E$ and $C_N$ are the costs of evaluating the energy and generating a neighboring state.

\textbf{Space Complexity}: $O(S)$ for the current, next, and best states, where $S$ is the size of one state.

\begin{lstlisting}[language=C++]
#include <cassert>
#include <cmath>
#include <random>
#include <utility>

template<typename State, typename Energy, typename Neighbor>
std::pair<double, State> anneal_min(
    State initial, Energy energy, Neighbor rand_neighbor, std::mt19937 &rng,
    const double temp_start = 1000, const double temp_end = 1e-6, const double cooling_rate = 0.995
) {
  assert(0 < temp_end && temp_end < temp_start);
  assert(0 < cooling_rate && cooling_rate < 1);
  std::uniform_real_distribution<double> unit(0.0, 1.0);
  State current = initial, best = initial;
  double current_energy = energy(current), best_energy = current_energy;
  for (double temp = temp_start; temp > temp_end; temp *= cooling_rate) {
    State next = rand_neighbor(current, temp, rng);
    double next_energy = energy(next);
    if (next_energy <= current_energy ||
        unit(rng) < std::exp((current_energy - next_energy) / temp)) {
      current = next;
      current_energy = next_energy;
      if (current_energy < best_energy) {
        best = current;
        best_energy = current_energy;
      }
    }
  }
  return {best_energy, best};
}
\end{lstlisting}
\Needspace*{4\baselineskip}
\textbf{Example Usage}\par
\begin{lstlisting}[language=C++]
#include <algorithm>
#include <vector>
using namespace std;

int main() {
  mt19937 rng(1234567);  // Set your own seed, or use random_device{}().

  // Escape the local minimum near x = 1 and find the lower minimum near x = -1.
  auto energy_1d = [](double x) { return (x * x - 1) * (x * x - 1) + 0.2 * x; };
  auto neighbor_1d = [](double x, double temp, mt19937 &gen) {
    uniform_real_distribution<double> move(-1.0, 1.0);
    return x + move(gen) * sqrt(temp);
  };
  auto [best_energy, best_x] = anneal_min(1.0, energy_1d, neighbor_1d, rng, 2.0, 1e-8);
  assert(best_energy < -0.2 && best_x < -1);

  // Approximate TSP with a permutation state and random 2-opt reversals.
  vector<pair<double, double>> points{{0, 0}, {1, 0}, {2, 0}, {2, 1},
                                      {2, 2}, {1, 2}, {0, 2}, {0, 1}};
  auto tour_length = [&](const vector<int> &tour) {
    double length = 0;
    for (int i = 0; i < static_cast<int>(tour.size()); i++) {
      auto [x1, y1] = points[tour[i]];
      auto [x2, y2] = points[tour[(i + 1) % tour.size()]];
      length += hypot(x1 - x2, y1 - y2);
    }
    return length;
  };
  auto reverse_segment = [](const vector<int> &tour, double, mt19937 &gen) {
    vector<int> next = tour;
    uniform_int_distribution<int> pick(0, static_cast<int>(tour.size()) - 1);
    int lo = pick(gen), hi = pick(gen);
    if (lo > hi) {
      swap(lo, hi);
    }
    reverse(next.begin() + lo, next.begin() + hi + 1);
    return next;
  };
  vector<int> initial_tour{0, 4, 2, 6, 1, 5, 3, 7};
  auto [best_length, best_tour] =
      anneal_min(initial_tour, tour_length, reverse_segment, rng, 5.0, 1e-6, 0.995);
  assert(fabs(best_length - 8) < 1e-9);
  assert(fabs(tour_length(best_tour) - best_length) < 1e-9);
  return 0;
}
\end{lstlisting}

\section{\texorpdfstring{Dynamic Programming Tricks}{5.3 Dynamic Programming Tricks}}
\setcounter{section}{3}
\setcounter{subsection}{0}
\subsection{Divide-and-Conquer DP Optimization}
Computes one layer of a dynamic program of the form \inlinecode{dp\_cur[i] = min(dp\_prev[k] + cost(k, i))}, where $0 \leq {\inlinecodemath{k}} \leq {\inlinecodemath{i}}$. Let \inlinecode{opt[i]} be the smallest index \inlinecode{k} attaining this minimum. Divide-and-conquer optimization applies when \inlinecode{opt[i]} $\leq$ \inlinecode{opt[i + 1]}, so the best transition index moves only to the right as \inlinecode{i} increases.

To compute \inlinecode{dp\_cur[l..r]}, evaluate its midpoint and record its best transition as \inlinecode{best\_k}. Monotonicity restricts every optimum in the left half to indices at most \inlinecode{best\_k}, and every optimum in the right half to indices at least \inlinecode{best\_k}. Recursively applying these bounds avoids checking every transition for every state. The caller must verify the required monotonicity property.

\begin{compactitem}
  \item \inlinecode{compute\_dp\_layer(dp\_prev, dp\_cur, l, r, opt\_l, opt\_r, cost)} fills \inlinecode{dp\_cur[l..r]} using candidate transition indices in [\inlinecode{opt\_l}, \inlinecode{opt\_r}]. The template parameter \inlinecode{cost} must be callable such that \inlinecode{cost(k, i)} returns the transition cost from previous state \inlinecode{k} to current state \inlinecode{i}.
\end{compactitem}

\textbf{Time Complexity}: $O(n \log n)$ calls to \inlinecode{cost(k, i)} per layer when each state has $O(n)$ possible transitions.

\textbf{Space Complexity}: $O(\log n)$ auxiliary stack space.

\begin{lstlisting}[language=C++]
#include <algorithm>
#include <cstdint>
#include <vector>

const int64_t INF = (1LL << 62);

template<typename Cost>
void compute_dp_layer(
    const std::vector<int64_t> &dp_prev, std::vector<int64_t> &dp_cur, int l, int r, int opt_l,
    int opt_r, Cost cost
) {
  if (l > r) {
    return;
  }
  int mid = l + (r - l) / 2;
  int64_t best = INF;
  int best_k = opt_l;
  int upper = std::min(mid, opt_r);
  for (int k = opt_l; k <= upper; k++) {
    int64_t candidate = dp_prev[k] + cost(k, mid);  // Overflow warning!
    if (candidate < best) {
      best = candidate;
      best_k = k;
    }
  }
  dp_cur[mid] = best;
  compute_dp_layer(dp_prev, dp_cur, l, mid - 1, opt_l, best_k, cost);
  compute_dp_layer(dp_prev, dp_cur, mid + 1, r, best_k, opt_r, cost);
}
\end{lstlisting}
\Needspace*{4\baselineskip}
\textbf{Example Usage}\par
\begin{lstlisting}[language=C++]
#include <cassert>
using namespace std;

struct SquareSegmentCost {
  vector<int64_t> prefix;

  explicit SquareSegmentCost(const vector<int> &a) : prefix(a.size() + 1) {
    for (int i = 0; i < static_cast<int>(a.size()); i++) {
      prefix[i + 1] = prefix[i] + a[i];
    }
  }

  int64_t operator()(int k, int i) const {
    // Cost of placing the half-open segment a[k..i) in one group.
    int64_t sum = prefix[i] - prefix[k];
    return sum * sum;
  }
};

int main() {
  vector<int> a{1, 2, 3, 4};
  int n = static_cast<int>(a.size());
  SquareSegmentCost cost(a);

  // Base layer for zero groups: only the empty prefix is reachable.
  vector<int64_t> dp_prev(n + 1, INF), dp_cur(n + 1, INF);
  dp_prev[0] = 0;

  // The first layer places each prefix a[0..j) into one group.
  compute_dp_layer(dp_prev, dp_cur, 1, n, 0, n - 1, cost);
  assert(dp_cur[4] == 100);  // One group: (1 + 2 + 3 + 4)^2.

  // The second layer chooses the best split between two groups.
  vector<int64_t> dp_two(n + 1, INF);
  compute_dp_layer(dp_cur, dp_two, 1, n, 0, n - 1, cost);
  assert(dp_two[4] == 52);  // (1 + 2 + 3)^2 + 4^2.
  return 0;
}
\end{lstlisting}
\subsection{SMAWK Row Minima}
Finds the position of the minimum in every row of a totally monotone matrix in time linear in its dimensions, without ever materializing the matrix. A matrix is totally monotone (for row minima) when the column index of each row's leftmost minimum never decreases as the row index increases, and this holds not just for the whole matrix but for every submatrix. Monge matrices, where $M_{i,j} + M_{i+1,j+1} \leq M_{i,j+1} + M_{i+1,j}$, are the most common source of such matrices and arise throughout dynamic programming on intervals and partitions.

The SMAWK algorithm interleaves two steps. REDUCE discards columns that cannot hold the minimum of any remaining row, leaving at most as many candidate columns as active rows. INTERPOLATE recurses on the odd-indexed active rows, then recovers each even-indexed active row's minimum by scanning only the candidate-column range bounded by its already-solved neighbors; total monotonicity guarantees those scans advance monotonically and cost $O(R + C)$ overall, where $R$ is the number of rows and $C$ is the number of columns in the original matrix. This removes the logarithmic factor of divide-and-conquer DP optimization, but unlike that technique it is purely offline: every matrix entry must be available on demand from the oracle, so it cannot be used when an entry depends on a row minimum computed earlier in the same pass.

\begin{compactitem}
  \item \inlinecode{smawk\_row\_minima(rows, cols, get)} returns a vector \inlinecode{arg} of length \inlinecode{rows}, where \inlinecode{arg[r]} is the column index minimizing \inlinecode{get(r, c)} over $0 \leq {\inlinecodemath{c}} < {\inlinecodemath{cols}}$ (on ties, the smallest qualifying column is chosen). The input matrix is supplied as a totally monotone function \inlinecode{get(r, c)} with bounds \inlinecode{rows} $\geq 0$ and \inlinecode{cols} $\geq 1$.
\end{compactitem}

\textbf{Time Complexity}: $O(R + C)$ evaluations of \inlinecode{get} per call, where $R$ is \inlinecode{rows} and $C$ is \inlinecode{cols}.

\textbf{Space Complexity}: $O(R + C)$ auxiliary.

\begin{lstlisting}[language=C++]
#include <numeric>
#include <vector>

template<typename Get>
std::vector<int> smawk_rec(
    const Get &get, const std::vector<int> &active_rows, const std::vector<int> &candidate_cols
) {
  int num_rows = static_cast<int>(active_rows.size());
  std::vector<int> arg(num_rows, -1);
  if (num_rows == 0) {
    return arg;
  }
  // Base case: a single row or column is cheap to scan directly.
  if (num_rows <= 2 || candidate_cols.size() <= 2) {
    for (int row_pos = 0; row_pos < num_rows; row_pos++) {
      for (int c : candidate_cols) {
        if (arg[row_pos] == -1 ||
            get(active_rows[row_pos], c) < get(active_rows[row_pos], arg[row_pos])) {
          arg[row_pos] = c;
        }
      }
    }
    return arg;
  }
  // REDUCE: keep only columns that can be the minimum of some row.
  std::vector<int> kept;
  if (static_cast<int>(candidate_cols.size()) > num_rows) {
    for (int c : candidate_cols) {
      while (!kept.empty()) {
        int r = active_rows[static_cast<int>(kept.size()) - 1];
        if (get(r, kept.back()) <= get(r, c)) {
          break;
        }
        kept.pop_back();
      }
      if (static_cast<int>(kept.size()) < num_rows) {
        kept.push_back(c);
      }
    }
  } else {
    kept = candidate_cols;
  }
  // INTERPOLATE: solve the odd rows recursively, then fill the even rows between known bounds.
  std::vector<int> odd_rows;
  for (int row_pos = 1; row_pos < num_rows; row_pos += 2) {
    odd_rows.push_back(active_rows[row_pos]);
  }
  std::vector<int> odd_arg = smawk_rec(get, odd_rows, kept);
  for (int odd_pos = 0; odd_pos < static_cast<int>(odd_arg.size()); odd_pos++) {
    arg[2 * odd_pos + 1] = odd_arg[odd_pos];
  }
  int col_pos = 0;
  for (int row_pos = 0; row_pos < num_rows; row_pos += 2) {
    int hi = (row_pos + 1 < num_rows) ? arg[row_pos + 1] : kept.back();
    int r = active_rows[row_pos];
    while (true) {
      int c = kept[col_pos];
      if (arg[row_pos] == -1 || get(r, c) < get(r, arg[row_pos])) {
        arg[row_pos] = c;
      }
      if (c == hi) {
        break;  // Stop at the upper bound; the next even row starts scanning from here.
      }
      col_pos++;
    }
  }
  return arg;
}

template<typename Get>
std::vector<int> smawk_row_minima(int rows, int cols, const Get &get) {
  std::vector<int> row_ids(rows), col_ids(cols);
  std::iota(row_ids.begin(), row_ids.end(), 0);
  std::iota(col_ids.begin(), col_ids.end(), 0);
  return smawk_rec(get, row_ids, col_ids);
}
\end{lstlisting}
\Needspace*{4\baselineskip}
\textbf{Example Usage}\par
\begin{lstlisting}[language=C++]
#include <cassert>
#include <vector>
using namespace std;

int main() {
  // An explicit totally monotone matrix; row minima sit at non-decreasing columns.
  vector<vector<int>> m{
      {10, 17, 13, 28, 23},
      {17, 22, 16, 29, 23},
      {24, 28, 22, 34, 24},
      {11, 13,  6, 17,  7},
      {45, 44, 32, 37, 23},
  };
  auto get = [&](int r, int c) { return m[r][c]; };
  vector<int> arg = smawk_row_minima(5, 5, get);

  // Cross-check each reported column against a brute-force row scan.
  for (int r = 0; r < 5; r++) {
    int best = 0;
    for (int c = 1; c < 5; c++) {
      if (m[r][c] < m[r][best]) {
        best = c;
      }
    }
    assert(m[r][arg[r]] == m[r][best]);  // Same minimum value (smallest index on ties).
  }
  return 0;
}
\end{lstlisting}
\subsection{Knuth DP Optimization}
Computes minimum costs for every half-open interval $[{\inlinecodemath{l}}, {\inlinecodemath{r}})$ using the recurrence \inlinecode{dp[l][r] = min(dp[l][k] + dp[k][r]) + cost(l, r)}. Each state chooses a split \inlinecode{k} strictly inside the interval and pays \inlinecode{cost(l, r)} after solving the two resulting subintervals.

A direct implementation checks every split and takes $O(n^{3})$. Knuth optimization applies when the smallest optimal splits satisfy \inlinecode{opt[l][r - 1]} $\leq$ \inlinecode{opt[l][r]} $\leq$ \inlinecode{opt[l + 1][r]}. The two neighboring states therefore bound the splits that must be checked for \inlinecode{dp[l][r]}, reducing the total number of candidate checks to $O(n^{2})$. Common applications include optimal binary search trees and some range merging problems. The caller must verify the required monotonicity property for the chosen \inlinecode{cost(l, r)}; the quadrangle inequality and interval monotonicity of the cost are common sufficient conditions.

\begin{compactitem}
  \item \inlinecode{knuth\_interval\_dp(n, cost, \&opt)} computes minimum costs for all half-open intervals $[{\inlinecodemath{l}}, {\inlinecodemath{r}})$ over \inlinecode{n} items. The template parameter \inlinecode{cost} must be callable such that \inlinecode{cost(l, r)} returns the interval cost added after choosing the best split. If the optional pointer \inlinecode{opt} is supplied, it is filled with the chosen split points.
\end{compactitem}

\textbf{Time Complexity}: $O(n^{2})$ calls to \inlinecode{cost(l, r)} and $O(n^{2})$ candidate split checks.

\textbf{Space Complexity}: $O(n^{2})$ for the \inlinecode{dp} and \inlinecode{opt} tables.

\begin{lstlisting}[language=C++]
#include <algorithm>
#include <cstdint>
#include <vector>

const int64_t INF = (1LL << 62);

template<typename Cost>
std::vector<std::vector<int64_t>> knuth_interval_dp(
    int n, Cost cost, std::vector<std::vector<int>> *out = nullptr
) {
  std::vector<std::vector<int64_t>> dp(n + 1, std::vector<int64_t>(n + 1, 0));
  std::vector<std::vector<int>> opt(n + 1, std::vector<int>(n + 1, 0));
  // Empty and one-item intervals have cost 0.
  for (int i = 0; i <= n; i++) {
    opt[i][i] = i;
    if (i < n) {
      opt[i][i + 1] = i + 1;
    }
  }
  for (int len = 2; len <= n; len++) {
    for (int l = 0; l + len <= n; l++) {
      int r = l + len;
      dp[l][r] = INF;
      int start = std::max(opt[l][r - 1], l + 1);
      int finish = std::min(opt[l + 1][r], r - 1);
      int64_t interval_cost = cost(l, r);
      for (int k = start; k <= finish; k++) {
        int64_t candidate = dp[l][k] + dp[k][r] + interval_cost;  // Overflow warning!
        if (candidate < dp[l][r]) {
          dp[l][r] = candidate;
          opt[l][r] = k;
        }
      }
    }
  }
  if (out != nullptr) {
    *out = opt;
  }
  return dp;
}
\end{lstlisting}
\Needspace*{4\baselineskip}
\textbf{Example Usage}\par
\begin{lstlisting}[language=C++]
#include <cassert>
using namespace std;

struct MergeCost {
  vector<int64_t> prefix;

  explicit MergeCost(const vector<int> &a) : prefix(a.size() + 1) {
    for (int i = 0; i < static_cast<int>(a.size()); i++) {
      prefix[i + 1] = prefix[i] + a[i];
    }
  }

  int64_t operator()(int l, int r) const {
    // Cost of merging the half-open interval a[l..r) into one group.
    return prefix[r] - prefix[l];
  }
};

int main() {
  vector<int> a{1, 2, 3, 4};
  vector<vector<int>> opt;
  vector<vector<int64_t>> dp = knuth_interval_dp(a.size(), MergeCost(a), &opt);
  assert(dp[0][4] == 19);  // Merge 1 + 2, then 3 + 3, then 6 + 4.
  assert(opt[0][4] == 3);  // The final merge splits [1, 2, 3] from [4].
  return 0;
}
\end{lstlisting}
\subsection{Lagrangian Relaxation (Aliens Trick)}
Given an optimization problem with a constraint on the number of chosen objects, Lagrangian relaxation moves that count into the objective by charging a penalty per object. If the relaxed problem can be solved quickly and the number chosen is monotone as the penalty changes, binary search recovers the best value for the desired count.

For maximization, solve the relaxed problem as \inlinecode{score - penalty*count}, breaking ties toward larger \inlinecode{count}. As \inlinecode{penalty} increases, the optimal count never increases. The largest penalty whose relaxed optimum still chooses at least \inlinecode{target\_count} objects gives the exact constrained score after adding back \inlinecode{penalty*target\_count}.

\begin{compactitem}
  \item \inlinecode{lagrangian\_maximize(target\_count, lo, hi, solve)} returns the maximum original score using exactly \inlinecode{target\_count} objects. The callable \inlinecode{solve(penalty)} must return \inlinecode{relaxed\_score} and \inlinecode{count} as a pair and must tie-break toward larger \inlinecode{count}. Penalties are searched over the half-open range $[{\inlinecodemath{lo}}, {\inlinecodemath{hi}})$; \inlinecode{lo} must choose at least \inlinecode{target\_count} objects, while \inlinecode{hi} must be greater than the last penalty that does so.
\end{compactitem}

This technique is often called the Aliens trick. Widen the penalty range when in doubt.

\textbf{Time Complexity}: $O(\log n)$ calls to \inlinecode{solve()} per call, where $n$ is the distance between \inlinecode{lo} and \inlinecode{hi}.

\textbf{Space Complexity}: $O(1)$ auxiliary.

\begin{lstlisting}[language=C++]
#include <cassert>
#include <cstdint>
#include <utility>
#include <vector>

template<typename Solve>
int64_t lagrangian_maximize(int target_count, int64_t lo, int64_t hi, Solve solve) {
  // The chosen count is nonincreasing, so find the last penalty that still chooses enough objects.
  while (lo < hi) {
    int64_t mid = lo + (hi - lo) / 2;
    if (solve(mid).second >= target_count) {
      lo = mid + 1;
    } else {
      hi = mid;
    }
  }
  int64_t penalty = lo - 1;
  auto [relaxed_score, count] = solve(penalty);
  assert(count >= target_count);
  return relaxed_score + penalty * target_count;  // Overflow warning!
}
\end{lstlisting}
\Needspace*{4\baselineskip}
\textbf{Example Usage}\par
\begin{lstlisting}[language=C++]
int main() {
  std::vector<int64_t> value{10, 7, 3, 2};
  auto choose_profitable = [&](int64_t penalty) {
    int64_t score = 0;
    int count = 0;
    for (int64_t x : value) {
      int64_t gain = x - penalty;
      if (gain >= 0) {  // Include zero gains to break ties toward larger counts.
        score += gain;
        count++;
      }
    }
    return std::pair<int64_t, int>{score, count};
  };

  assert(lagrangian_maximize(2, 0, 20, choose_profitable) == 17);
  assert(lagrangian_maximize(3, 0, 20, choose_profitable) == 20);
  return 0;
}
\end{lstlisting}
\subsection{Convex Hull Trick (Semi-Dynamic)}
Given a set of pairs $(m, b)$ specifying lines of the form $y = mx + b$, answer queries at specified $x$-coordinates, each asking for the minimum $y$-value over all given lines. This is useful for dynamic programming recurrences of the form \inlinecode{dp[i] = min(m[j] * x[i] + b[j])}. Only the lower envelope of the lines can ever answer a query, so each added line pops previously stored lines that it renders useless, and ascending queries advance a pointer along the envelope so every line is visited at most twice.

The following implementation is a concise, semi-dynamic version of the convex hull optimization technique. It supports an interlaced sequence of \inlinecode{add\_line()} and \inlinecode{query()} calls, as long as the preconditions of descending \inlinecode{m} and ascending \inlinecode{x} are satisfied. As a result, it may be necessary to sort the lines and queries before calling the functions. In that case, the overall time complexity will be dominated by the sorting step.

\begin{compactitem}
  \item \inlinecode{SemiDynamicCHT()} constructs an empty hull.
  \item \inlinecode{add\_line(m, b)} inserts the line $y = mx + b$. The caller must ensure that slope \inlinecode{m} is less than or equal to the slope of every line added so far.
  \item \inlinecode{query(x)} returns the minimum $y$-value among all inserted lines at coordinate \inlinecode{x}. The caller must ensure that query coordinates are nondecreasing across calls. At least one line must have been inserted.
\end{compactitem}

Overflow warning: \inlinecode{add\_line()} compares intersections by cross-multiplying slope and intercept differences, a product on the order of the squared coefficient magnitude. For large \inlinecode{m}/\inlinecode{b} (roughly beyond $10^9$ with \inlinecode{int64\_t}), cast that comparison to \inlinecode{\_\_int128}. \inlinecode{query()} evaluates \inlinecode{m * x + b}, which also needs to fit in \inlinecode{int64\_t}.

\textbf{Time Complexity}: $O(n + q)$ for any interlaced sequence containing $n$ calls to \inlinecode{add\_line()} and $q$ calls to \inlinecode{query()}. Each line is removed or passed by the query pointer at most once, so both operations take amortized $O(1)$ time.

\Needspace*{3\baselineskip}
\textbf{Space Complexity}:
\begin{compactitem}
  \item $O(n)$ for storage of the lines.
  \item $O(1)$ auxiliary for \inlinecode{add\_line()} and \inlinecode{query()}.
\end{compactitem}

\begin{lstlisting}[language=C++]
#include <cassert>
#include <cstdint>
#include <vector>

class SemiDynamicCHT {
  std::vector<int64_t> slopes, intercepts;
  int ptr = 0;

 public:
  void add_line(int64_t m, int64_t b) {
    int len = static_cast<int>(slopes.size());
    if (len > 0 && slopes.back() == m) {
      if (intercepts.back() <= b) {
        return;
      }
      slopes.pop_back();
      intercepts.pop_back();
      len--;
      if (ptr > len) {
        ptr = len;
      }
    }
    // Overflow warning!
    while (len > 1 && (intercepts[len - 2] - intercepts[len - 1]) * (m - slopes[len - 1]) >=
                          (intercepts[len - 1] - b) * (slopes[len - 1] - slopes[len - 2])) {
      len--;
    }
    slopes.resize(len);
    intercepts.resize(len);
    slopes.push_back(m);
    intercepts.push_back(b);
  }

  int64_t query(int64_t x) {
    assert(!slopes.empty());
    if (ptr >= static_cast<int>(slopes.size())) {
      ptr = static_cast<int>(slopes.size()) - 1;
    }
    // Overflow warning!
    while (ptr + 1 < static_cast<int>(slopes.size()) &&
           slopes[ptr + 1] * x + intercepts[ptr + 1] <= slopes[ptr] * x + intercepts[ptr]) {
      ptr++;
    }
    return slopes[ptr] * x + intercepts[ptr];
  }
};
\end{lstlisting}
\Needspace*{4\baselineskip}
\textbf{Example Usage}\par
\begin{lstlisting}[language=C++]
#include <cassert>

int main() {
  SemiDynamicCHT cht;
  cht.add_line(3, 0);
  cht.add_line(2, 1);
  cht.add_line(1, 2);
  cht.add_line(0, 6);
  assert(cht.query(0) == 0);
  assert(cht.query(1) == 3);
  assert(cht.query(2) == 4);
  assert(cht.query(3) == 5);
  cht.add_line(0, 7);  // Same slope, worse intercept; ignored.
  assert(cht.query(4) == 6);
  cht.add_line(0, 5);  // Same slope, better intercept; replaces the old constant line.
  assert(cht.query(5) == 5);
  return 0;
}
\end{lstlisting}
\subsection{Convex Hull Trick (Fully-Dynamic)}
Maintains a dynamic set of lines $y = mx + b$ and answers minimum value queries at integer points. Lines and queries may arrive in arbitrary order, making this useful for dynamic programming recurrences of the form \inlinecode{dp[i] = min(m[j] * x[i] + b[j])} without monotone slopes or query coordinates. The same interface can answer maximum queries when constructed with \inlinecode{query\_max = true}.

This fully dynamic convex hull stores the envelope in a self-balancing binary search tree (\inlinecode{std::set}). Inserting a line removes any neighbors it dominates, and each remaining line records the last coordinate where it is optimal, so a query is one tree lookup. Unlike a Li Chao tree, this structure requires no fixed coordinate domain and its complexity depends on the number of lines rather than the domain width. Prefer it when the domain is unknown or very large, or when maximum queries are needed; a Li Chao tree is often simpler when a manageable domain is known.

\begin{compactitem}
  \item \inlinecode{HullOptimizer(query\_max = false)} constructs an empty hull. By default, \inlinecode{query(x)} minimizes; if \inlinecode{query\_max} is true, \inlinecode{query(x)} maximizes.
  \item \inlinecode{add\_line(m, b)} inserts line $y = mx + b$. Lines may be added in any order.
  \item \inlinecode{query(x)} returns the best $y$-value among all inserted lines at coordinate \inlinecode{x}. At least one line must have been inserted, and query coordinates may be supplied in any order.
\end{compactitem}

Overflow warning: minimization negates the coefficients, border updates subtract slopes/intercepts, and \inlinecode{query()} evaluates \inlinecode{m * x + b}, so those intermediate values must fit in \inlinecode{int64\_t}.

\textbf{Time Complexity}: $O(\log n)$ amortized per call to \inlinecode{add\_line()} and $O(\log n)$ per call to \inlinecode{query()}, where $n$ is the number of lines added.

\Needspace*{3\baselineskip}
\textbf{Space Complexity}:
\begin{compactitem}
  \item $O(n)$ for storage of the lines.
  \item $O(1)$ auxiliary for \inlinecode{add\_line()} and \inlinecode{query()}.
\end{compactitem}

\begin{lstlisting}[language=C++]
#include <cassert>
#include <cstdint>
#include <set>

class HullOptimizer {
  struct Line {
    int64_t m, b;
    mutable int64_t xhi;
    bool is_query;

    Line(int64_t m, int64_t b, int64_t xhi = 0, bool is_query = false)
        : m(m), b(b), xhi(xhi), is_query(is_query) {}

    bool operator<(const Line &l) const { return l.is_query ? xhi < l.xhi : m < l.m; }
  };

  std::multiset<Line> hull;
  bool query_max;

  using hulliter = std::multiset<Line>::iterator;

  bool update_border(hulliter x, hulliter y) {
    if (y == hull.end()) {
      x->xhi = INT64_MAX;
      return false;
    }
    if (x->m == y->m) {
      x->xhi = (x->b > y->b) ? INT64_MAX : INT64_MIN;
    } else {
      int64_t a = y->b - x->b, b = x->m - y->m;  // Overflow warning!
      x->xhi = a / b - ((a ^ b) < 0 && a % b);
    }
    return x->xhi >= y->xhi;
  }

 public:
  explicit HullOptimizer(bool query_max = false) : query_max(query_max) {}

  void add_line(int64_t m, int64_t b) {
    if (!query_max) {
      m = -m;
      b = -b;
    }
    auto z = hull.insert(Line(m, b));
    auto y = z++;
    auto x = y;
    while (update_border(y, z)) {
      z = hull.erase(z);
    }
    if (x != hull.begin() && update_border(--x, y)) {
      update_border(x, y = hull.erase(y));
    }
    while ((y = x) != hull.begin() && (--x)->xhi >= y->xhi) {
      update_border(x, hull.erase(y));
    }
  }

  int64_t query(int64_t x) const {
    assert(!hull.empty());
    Line q(0, 0, x, true);
    auto it = hull.lower_bound(q);
    int64_t res = it->m * x + it->b;  // Overflow warning!
    return query_max ? res : -res;
  }
};
\end{lstlisting}
\Needspace*{4\baselineskip}
\textbf{Example Usage}\par
\begin{lstlisting}[language=C++]
#include <cassert>

int main() {
  HullOptimizer h;
  h.add_line(3, 0);
  h.add_line(0, 6);
  h.add_line(1, 2);
  h.add_line(2, 1);
  // Minimize among y = 3x, 6, x + 2, and 2x + 1.
  assert(h.query(0) == 0);
  assert(h.query(2) == 4);
  assert(h.query(1) == 3);
  assert(h.query(3) == 5);

  HullOptimizer mx(true);
  mx.add_line(3, 0);
  mx.add_line(0, 6);
  mx.add_line(1, 2);
  // Same interface can maximize when constructed with query_max = true.
  assert(mx.query(0) == 6);
  assert(mx.query(3) == 9);
  return 0;
}
\end{lstlisting}
\subsection{Li Chao Tree}
Maintains a dynamic set of lines $y = mx + b$ and answers minimum value queries at integer points. Lines and queries may arrive in arbitrary order, making this useful for dynamic programming recurrences of the form \inlinecode{dp[i] = min(m[j] * x[i] + b[j])} without monotone slopes or query coordinates.

A Li Chao tree fixes an inclusive coordinate domain $[{\inlinecodemath{lo}}, {\inlinecodemath{hi}}]$ and stores one line at each node of a dynamic segment tree. Inserting a line keeps the line that wins at the midpoint and recurses only into the half where the other line may still win; a query takes the best value along one root-to-leaf path. Prefer this structure when the domain is known and manageable. The fully dynamic convex hull in the previous section needs no fixed domain and also supports maximum queries, while a Li Chao tree has a simpler invariant and is easier to extend to lines active only on subintervals.

\begin{compactitem}
  \item \inlinecode{LiChaoTree(lo, hi)} constructs an empty tree over integer domain $[{\inlinecodemath{lo}}, {\inlinecodemath{hi}}]$.
  \item \inlinecode{add\_line(m, b)} inserts line $y = mx + b$. Lines may be added in any order.
  \item \inlinecode{query(x)} returns the minimum $y$-value among all inserted lines at coordinate \inlinecode{x}. At least one line must have been inserted, and query coordinates may be supplied in any order.
\end{compactitem}

Overflow warning: each comparison and query evaluates \inlinecode{m * x + b}, which must fit in \inlinecode{int64\_t}.

\textbf{Time Complexity}: $O(\log d)$ per call to \inlinecode{add\_line(m, b)} and \inlinecode{query(x)}, where $d$ is the distance between \inlinecode{lo} and \inlinecode{hi}.

\Needspace*{3\baselineskip}
\textbf{Space Complexity}:
\begin{compactitem}
  \item $O(n)$ node storage for $n$ lines inserted, since each insertion creates at most one node.
  \item $O(\log d)$ auxiliary stack space per operation.
\end{compactitem}

\begin{lstlisting}[language=C++]
#include <algorithm>
#include <cassert>
#include <cstdint>
#include <utility>

const int64_t INF = INT64_MAX / 4;

class LiChaoTree {
  struct Line {
    int64_t m, b;
    Line(int64_t m = 0, int64_t b = INF) : m(m), b(b) {}
    int64_t eval(int64_t x) const { return m * x + b; }  // Overflow warning!
  };

  struct Node {
    Line f;
    Node *left, *right;
    explicit Node(const Line &f) : f(f), left(nullptr), right(nullptr) {}
  };

  Node *root;
  int64_t lo, hi;

  static void clean_up(Node *n) {
    if (n == nullptr) {
      return;
    }
    clean_up(n->left);
    clean_up(n->right);
    delete n;
  }

  static void add_line(Node *&n, int64_t l, int64_t r, Line f) {
    if (n == nullptr) {
      n = new Node(f);
      return;
    }
    int64_t mid = l + (r - l) / 2;
    bool left_better = f.eval(l) < n->f.eval(l);
    bool mid_better = f.eval(mid) < n->f.eval(mid);
    if (mid_better) {
      std::swap(f, n->f);
    }
    if (l == r) {
      return;
    }
    if (left_better != mid_better) {
      add_line(n->left, l, mid, f);
    } else {
      add_line(n->right, mid + 1, r, f);
    }
  }

  static int64_t query(Node *n, int64_t l, int64_t r, int64_t x) {
    if (n == nullptr) {
      return INF;
    }
    int64_t res = n->f.eval(x);
    if (l == r) {
      return res;
    }
    int64_t mid = l + (r - l) / 2;
    if (x <= mid) {
      return std::min(res, query(n->left, l, mid, x));
    }
    return std::min(res, query(n->right, mid + 1, r, x));
  }

 public:
  LiChaoTree(int64_t lo, int64_t hi) : root(nullptr), lo(lo), hi(hi) {}

  ~LiChaoTree() { clean_up(root); }
  LiChaoTree(const LiChaoTree &) = delete;
  LiChaoTree &operator=(const LiChaoTree &) = delete;
  void add_line(int64_t m, int64_t b) { add_line(root, lo, hi, Line(m, b)); }

  int64_t query(int64_t x) const {
    assert(root != nullptr);
    assert(lo <= x && x <= hi);
    return query(root, lo, hi, x);
  }
};
\end{lstlisting}
\Needspace*{4\baselineskip}
\textbf{Example Usage}\par
\begin{lstlisting}[language=C++]
#include <cassert>

int main() {
  LiChaoTree h(-10, 10);
  h.add_line(3, 0);
  h.add_line(0, 6);
  h.add_line(1, 2);
  h.add_line(2, 1);
  // Minimize among y = 3x, 6, x + 2, and 2x + 1.
  assert(h.query(0) == 0);
  assert(h.query(2) == 4);
  assert(h.query(1) == 3);
  assert(h.query(3) == 5);
  return 0;
}
\end{lstlisting}

\section{\texorpdfstring{Numerical Methods}{5.4 Numerical Methods}}
\setcounter{section}{4}
\setcounter{subsection}{0}
\subsection{Root Finding (Bracketing)}
Finds an $x$ in an interval $[a, b]$ for a continuous function $f$ such that $f(x) = 0$. By the intermediate value theorem, a root must exist in $[a, b]$ if either endpoint is a root or the signs of $f(a)$ and $f(b)$ differ. Each bisection step evaluates the midpoint of the interval and keeps the half whose endpoints still differ in sign, so a root always remains bracketed. After $n$ iterations, the interval width is $(b - a) / 2^n$. Although it is possible to control the error by looping while $b - a$ is greater than an arbitrary epsilon, it is simpler to let the loop run for a desired number of iterations until floating point arithmetic breaks down. 100 iterations is usually sufficient, since the search space will be reduced to $2^{-100}$ (roughly $10^{-30}$) times its original size.

\begin{compactitem}
  \item \inlinecode{bisection\_root(f, a, b, iterations = 100)} returns a root in an interval $[{\inlinecodemath{a}}, {\inlinecodemath{b}}]$ for a continuous function $f$ where either endpoint is a root or the endpoint values have opposite signs, using the bisection method.
  \item \inlinecode{falsi\_illinois\_root(f, a, b, iterations = 100)} returns a root in an interval $[{\inlinecodemath{a}}, {\inlinecodemath{b}}]$ for a continuous function $f$ under the same endpoint conditions, using the Illinois algorithm variant of the false position (a.k.a. regula falsi) method.
\end{compactitem}

\textbf{Time Complexity}: $O(n)$ calls will be made to \inlinecode{f()} in \inlinecode{bisection\_root()} and \inlinecode{falsi\_illinois\_root()}, where $n$ is the number of iterations performed.

\textbf{Space Complexity}: $O(1)$ auxiliary for both operations.

\begin{lstlisting}[language=C++]
#include <cassert>

template<typename Fn>
double bisection_root(Fn f, double a, double b, const int iterations = 100) {
  double fa = f(a), fb = f(b);
  assert(a <= b && ((fa <= 0 && fb >= 0) || (fa >= 0 && fb <= 0)));
  if (fa == 0) {
    return a;
  }
  if (fb == 0) {
    return b;
  }
  double m = a;
  for (int i = 0; i < iterations; i++) {
    m = a + (b - a) / 2;
    double fm = f(m);
    if (fm == 0) {
      return m;
    }
    if ((fa < 0) == (fm < 0)) {
      a = m;
      fa = fm;
    } else {
      b = m;
    }
  }
  return m;
}

template<typename Fn>
double falsi_illinois_root(Fn f, double a, double b, const int iterations = 100) {
  double fa = f(a), fb = f(b);
  assert(a <= b && ((fa <= 0 && fb >= 0) || (fa >= 0 && fb <= 0)));
  if (fa == 0) {
    return a;
  }
  if (fb == 0) {
    return b;
  }
  double m = a;
  int side = 0;
  for (int i = 0; i < iterations; i++) {
    m = (fa * b - fb * a) / (fa - fb);
    double fm = f(m);
    if (fm == 0) {
      break;
    }
    if ((fb < 0) == (fm < 0)) {
      b = m;
      fb = fm;
      if (side < 0) {
        fa /= 2;
      }
      side = -1;
    } else {
      a = m;
      fa = fm;
      if (side > 0) {
        fb /= 2;
      }
      side = 1;
    }
  }
  return m;
}
\end{lstlisting}
\Needspace*{4\baselineskip}
\textbf{Example Usage}\par
\begin{lstlisting}[language=C++]
#include <cmath>

double f(double x) {
  return x * x - 4 * sin(x);
}

int main() {
  assert(fabs(f(bisection_root(f, 1, 3))) < 1e-10);
  assert(fabs(f(falsi_illinois_root(f, 1, 3))) < 1e-10);
  assert(bisection_root(f, 0, 1) == 0);
  assert(falsi_illinois_root(f, 0, 1) == 0);
  return 0;
}
\end{lstlisting}
\subsection{Root Finding (Iteration)}
Finds an $x$ for a continuous function $f$ such that $f(x) = 0$ using iterative approximation by an initial guess that is close to the answer. Each step follows the tangent line at the current guess down to its $x$-intercept, that is, $x \leftarrow x - f(x)/f'(x)$. Newton's method requires an explicit definition of the function's derivative while the secant method starts with two initial guesses and approximates the derivative using the secant slope from the previous iteration. For $n$ iterations and a good initial guess, the methods below compute approximately $2^n$ digits of precision, with the secant method converging approximately $1.6$ times slower than Newton's.

\begin{compactitem}
  \item \inlinecode{newton\_root(f, fprime, x0, eps = 1e-15, iterations = 100)} returns a root $x$ for a function \inlinecode{f} with derivative \inlinecode{fprime} using an initial guess \inlinecode{x0} which should be relatively close to $x$.
  \item \inlinecode{secant\_root(f, x0, x1, eps = 1e-15, iterations = 100)} returns a root $x$ for a function \inlinecode{f} using two initial guesses \inlinecode{x0} and \inlinecode{x1} which should be relatively close to $x$.
\end{compactitem}

\textbf{Time Complexity}: $O(n)$ calls will be made to \inlinecode{f()} in \inlinecode{newton\_root()} and \inlinecode{secant\_root()}, where $n$ is the number of iterations performed.

\textbf{Space Complexity}: $O(1)$ auxiliary for both operations.

\begin{lstlisting}[language=C++]
#include <cmath>
#include <stdexcept>

template<typename Fn, typename Deriv>
double newton_root(
    Fn f, Deriv fprime, double x0, const double eps = 1e-15, const int iterations = 100
) {
  double x = x0, error = eps + 1;
  for (int i = 0; error > eps && i < iterations; i++) {
    double xnew = x - f(x) / fprime(x);
    error = std::fabs(xnew - x);
    x = xnew;
  }
  if (error > eps) {
    throw std::runtime_error("Newton's method failed to converge.");
  }
  return x;
}

template<typename Fn>
double secant_root(
    Fn f, double x0, double x1, const double eps = 1e-15, const int iterations = 100
) {
  double xold = x0, fxold = f(x0), x = x1, error = eps + 1;
  for (int i = 0; error > eps && i < iterations; i++) {
    double fx = f(x);
    double xnew = x - fx * ((x - xold) / (fx - fxold));
    xold = x;
    fxold = fx;
    error = std::fabs(xnew - x);
    x = xnew;
  }
  if (error > eps) {
    throw std::runtime_error("Secant method failed to converge.");
  }
  return x;
}
\end{lstlisting}
\Needspace*{4\baselineskip}
\textbf{Example Usage}\par
\begin{lstlisting}[language=C++]
#include <cassert>

double f(double x) {
  return x * x - 4 * sin(x);
}

double fprime(double x) {
  return 2 * x - 4 * cos(x);
}

int main() {
  assert(fabs(f(newton_root(f, fprime, 3))) < 1e-10);
  assert(fabs(f(secant_root(f, 3, 2))) < 1e-10);
  auto g = [](double x) { return x * x - 2; };
  auto gprime = [](double x) { return 2 * x; };
  assert(fabs(g(newton_root(g, gprime, 1))) < 1e-10);
  return 0;
}
\end{lstlisting}
\subsection{Polynomial Root Finding (Differentiation)}
Finds every real root $x$ of a polynomial $p$ satisfying $p(x) = 0$ by recursively finding the roots of its derivative. Each adjacent pair of local extrema is searched using the bisection method.

The key observation is that derivative roots split the real line into intervals where the polynomial is monotone. A monotone interval contains at most one root, and if an endpoint is a root or the endpoint values have opposite signs, bisection isolates that root. Recursively solving the derivative therefore provides all interval boundaries needed to find every real root in the requested range, including roots of even multiplicity.

\begin{compactitem}
  \item \inlinecode{horner\_eval(p, x)} evaluates the polynomial \inlinecode{p} of degree $d$ (represented as a vector of size $d + 1$ where \inlinecode{p[i]} stores the coefficient for the $x^i$ term) at \inlinecode{x}, using Horner's method.
  \item \inlinecode{find\_one\_root(p, a, b, eps = 1e-15)} returns a root in the interval $[{\inlinecodemath{a}}, {\inlinecodemath{b}}]$ for a polynomial \inlinecode{p} where either endpoint is a root or the endpoint values have opposite signs, using the bisection method. If this precondition is not satisfied, then \inlinecode{NaN} is returned. The root is found to a tolerance of \inlinecode{eps} in absolute or relative error (whichever is reached first).
  \item \inlinecode{find\_all\_roots(p, a = -1e20, b = 1e20, eps = 1e-15)} returns a vector of all roots in the interval $[{\inlinecodemath{a}}, {\inlinecodemath{b}}]$ for a polynomial \inlinecode{p} using the bisection method. The roots are found to a tolerance of \inlinecode{eps} in absolute or relative error (whichever is reached first).
\end{compactitem}

\Needspace*{3\baselineskip}
\textbf{Time Complexity}:
\begin{compactitem}
  \item $O(n)$ per call to \inlinecode{horner\_eval()}, where $n$ is the degree of the polynomial.
  \item $O(nt)$ per call to \inlinecode{find\_one\_root()}, where $n$ is the degree of the polynomial and $t$ is the number of bisection iterations required to reach the requested tolerance.
  \item $O(n^{3} t)$ per call to \inlinecode{find\_all\_roots()}, where $n$ is the degree of the polynomial and $t$ is the number of bisection iterations required to reach the requested tolerance.
\end{compactitem}

\Needspace*{3\baselineskip}
\textbf{Space Complexity}:
\begin{compactitem}
  \item $O(1)$ auxiliary for \inlinecode{horner\_eval()} and \inlinecode{find\_one\_root()}.
  \item $O(n)$ auxiliary stack space and $O(n^{2})$ auxiliary heap space for \inlinecode{find\_all\_roots()}, where $n$ is the degree of the polynomial.
\end{compactitem}

\begin{lstlisting}[language=C++]
#include <cmath>
#include <limits>
#include <vector>

double horner_eval(const std::vector<double> &p, double x) {
  double res = p.back();
  for (int i = static_cast<int>(p.size()) - 2; i >= 0; i--) {
    res = res * x + p[i];
  }
  return res;
}

double find_one_root(const std::vector<double> &p, double a, double b, const double eps = 1e-15) {
  double pa = horner_eval(p, a), pb = horner_eval(p, b);
  if (std::fabs(pa) <= eps) {
    return a;
  }
  if (std::fabs(pb) <= eps) {
    return b;
  }
  bool paneg = pa < 0, pbneg = pb < 0;
  if (paneg == pbneg) {
    return std::numeric_limits<double>::quiet_NaN();
  }
  while (b - a > eps && a * (1 + eps) < b && a < b * (1 + eps)) {
    double m = a + (b - a) / 2;
    if ((horner_eval(p, m) < 0) == paneg) {
      a = m;
    } else {
      b = m;
    }
  }
  return a;
}

std::vector<double> find_all_roots(
    const std::vector<double> &p, double a = -1e20, double b = 1e20, const double eps = 1e-15
) {
  std::vector<double> pprime;
  pprime.reserve(p.size() > 0 ? p.size() - 1 : 0);
  for (int i = 1; i < static_cast<int>(p.size()); i++) {
    pprime.push_back(p[i] * i);
  }
  if (pprime.empty()) {
    return {};
  }
  std::vector<double> res, r = find_all_roots(pprime, a, b, eps);
  r.push_back(b);
  for (int i = 0; i < static_cast<int>(r.size()); i++) {
    double root = find_one_root(p, i == 0 ? a : r[i - 1], r[i], eps);
    if (!std::isnan(root) && (res.empty() || root != res.back())) {
      res.push_back(root);
    }
  }
  return res;
}
\end{lstlisting}
\Needspace*{4\baselineskip}
\textbf{Example Usage}\par
\begin{lstlisting}[language=C++]
#include <cassert>
using namespace std;

int main() {
  {  // -1 + 2x - 6x^2 + 2x^3
    vector<double> p{-1.0, 2.0, -6.0, 2.0};
    vector<double> roots = find_all_roots(p);
    assert(roots.size() == 1 && fabs(horner_eval(p, roots[0])) < 1e-10);
  }
  {  // -20 + 4x + 3x^2
    vector<double> p{-20.0, 4.0, 3.0};
    vector<double> roots = find_all_roots(p);
    assert(roots.size() == 2);
    assert(fabs(horner_eval(p, roots[0])) < 1e-10);
    assert(fabs(horner_eval(p, roots[1])) < 1e-10);
  }
  {  // (x - 1)^2
    vector<double> p{1.0, -2.0, 1.0};
    vector<double> roots = find_all_roots(p);
    assert(roots.size() == 1 && fabs(roots[0] - 1) < 1e-10);
  }
  return 0;
}
\end{lstlisting}
\subsection{Polynomial Root Finding (Laguerre)}
Finds every complex root $x$ for a polynomial $p$ with complex coefficients such that $p(x) = 0$. Laguerre's method is a root-finding iteration with reliable, near-cubic convergence from almost any starting point. Once a root is found it is removed by polynomial deflation, and the iteration repeats on the lower-degree quotient until every root has been extracted.

\begin{compactitem}
  \item \inlinecode{horner\_eval(p, x)} evaluates a complex polynomial $p$ of degree $d$ (represented as a vector of size $d + 1$ where \inlinecode{p[i]} stores the complex coefficient for the $x^i$ term) at \inlinecode{x}, using Horner's method, returning a pair where the first value is the final result $p(x)$ and the second value is the quotient polynomial $q$ from the identity $p(t) = (t - x)q(t) + p(x)$.
  \item \inlinecode{find\_one\_root(p, x0, eps = 1e-15, iterations = 10000)} returns a complex root $x$ for polynomial \inlinecode{p} (represented as a vector of size $d + 1$ where \inlinecode{p[i]} stores the complex coefficient for the $x^i$ term) using an initial guess \inlinecode{x0} which should be relatively close to $x$. The root is found to a tolerance of \inlinecode{eps} in absolute or relative error (whichever is reached first).
  \item \inlinecode{find\_all\_roots(p, eps = 1e-15, iterations = 10000)} returns a vector of all complex roots for a complex polynomial \inlinecode{p}. The roots are found to a tolerance of \inlinecode{eps} in absolute or relative error (whichever is reached first).
\end{compactitem}

\Needspace*{3\baselineskip}
\textbf{Time Complexity}:
\begin{compactitem}
  \item $O(n)$ per call to \inlinecode{horner\_eval()}, where $n$ is the degree of the polynomial.
  \item $O(nt)$ per call to \inlinecode{find\_one\_root()}, where $n$ is the degree of the polynomial and $t$ is the number of iterations performed.
  \item $O(n^{2} t)$ per call to \inlinecode{find\_all\_roots()}, where $n$ is the degree of the polynomial and $t$ is the number of iterations performed per root.
\end{compactitem}

\Needspace*{3\baselineskip}
\textbf{Space Complexity}:
\begin{compactitem}
  \item $O(n)$ auxiliary for \inlinecode{horner\_eval()} and \inlinecode{find\_one\_root()}, where $n$ is the degree of the polynomial.
  \item $O(n)$ auxiliary for \inlinecode{find\_all\_roots()}, where $n$ is the degree of the polynomial.
\end{compactitem}

\begin{lstlisting}[language=C++]
#include <algorithm>
#include <cmath>
#include <complex>
#include <random>
#include <utility>
#include <vector>

using cdouble = std::complex<double>;
using cpoly = std::vector<cdouble>;

std::pair<cdouble, cpoly> horner_eval(const cpoly &p, const cdouble &x) {
  int n = static_cast<int>(p.size());
  cpoly b(std::max(1, n - 1));
  for (int i = n - 1; i > 0; i--) {
    b[i - 1] = p[i] + (i < n - 1 ? b[i] * x : 0);
  }
  return {p[0] + b[0] * x, b};
}

cpoly derivative(const cpoly &p) {
  int n = static_cast<int>(p.size());
  cpoly res(std::max(1, n - 1));
  for (int i = 1; i < n; i++) {
    res[i - 1] = p[i] * cdouble(i);
  }
  return res;
}

int comp(const cdouble &a, const cdouble &b, const double eps = 1e-15) {
  double diff = std::abs(a) - std::abs(b);
  return (diff < -eps) ? -1 : (diff > eps ? 1 : 0);
}

cdouble find_one_root(
    const cpoly &p, const cdouble &x0, const double eps = 1e-15, const int iterations = 10000
) {
  cdouble x = x0;
  int n = static_cast<int>(p.size()) - 1;
  cpoly p1 = derivative(p), p2 = derivative(p1);
  for (int i = 0; i < iterations; i++) {
    cdouble y0 = horner_eval(p, x).first;
    if (comp(y0, 0, eps) == 0) {
      break;
    }
    cdouble g = horner_eval(p1, x).first / y0;
    cdouble h = g * g - horner_eval(p2, x).first / y0;
    cdouble r = std::sqrt(cdouble(n - 1) * (h * cdouble(n) - g * g));
    cdouble d1 = g + r, d2 = g - r;
    cdouble a = cdouble(n) / (comp(d1, d2, eps) > 0 ? d1 : d2);
    x -= a;
    if (comp(a, 0, eps) == 0) {
      break;
    }
  }
  return x;
}

std::vector<cdouble> find_all_roots(
    const cpoly &p, const double eps = 1e-15, const int iterations = 10000
) {
  std::vector<cdouble> res;
  cpoly q = p;
  if (q.size() <= 1) {
    return res;
  }
  static std::mt19937 rng(std::random_device{}());
  std::uniform_real_distribution<double> unit(0.0, 1.0);
  while (q.size() > 2) {
    cdouble z(unit(rng), unit(rng));
    z = find_one_root(p, find_one_root(q, z, eps, iterations), eps, iterations);
    q = horner_eval(q, z).second;
    res.push_back(z);
  }
  res.push_back(-q[0] / q[1]);
  return res;
}
\end{lstlisting}
\Needspace*{4\baselineskip}
\textbf{Example Usage}\par
\begin{lstlisting}[language=C++]
#include <cstdio>
using namespace std;

void print_roots(vector<cdouble> x) {
  sort(x.begin(), x.end(), [](const cdouble &a, const cdouble &b) {
    return abs(a.real() - b.real()) > 0.5e-5 ? a.real() < b.real() : a.imag() < b.imag();
  });
  for (const auto &z : x) {
    double real = abs(z.real()) < 0.5e-5 ? 0 : z.real();
    double imag = abs(z.imag()) < 0.5e-5 ? 0 : z.imag();
    printf("(%.5lf, %.5lf)\n", real, imag);
  }
}

int main() {
  {  // 140 - 13x - 8x^2 + x^3 = (x + 4)(x - 5)(x - 7)
    printf("Roots of 140 - 13x - 8x^2 + x^3:\n");
    cpoly p;
    p.push_back(140);
    p.push_back(-13);
    p.push_back(-8);
    p.push_back(1);
    print_roots(find_all_roots(p));
  }
  {  // (-24+36i) + (-26+12i)x + (-30+40i)x^2 + (-26+12i)x^3 + (-6+4i)x^4
    // = ((2 + 3i)x + 6)(x + i)(2x + (6 + 4i))(xi + 1):
    printf("Roots of ((2 + 3i)x + 6)(x + i)(2x + (6 + 4i))(xi + 1):\n");
    cpoly p;
    p.emplace_back(-24, 36);
    p.emplace_back(-26, 12);
    p.emplace_back(-30, 40);
    p.emplace_back(-26, 12);
    p.emplace_back(-6, 4);
    print_roots(find_all_roots(p));
  }
  return 0;
}
\end{lstlisting}
\begin{exampleoutput}
Roots of 140 - 13x - 8x^2 + x^3:
(-4.00000, 0.00000)
(5.00000, 0.00000)
(7.00000, 0.00000)
Roots of ((2 + 3i)x + 6)(x + i)(2x + (6 + 4i))(xi + 1):
(-3.00000, -2.00000)
(-0.92308, 1.38462)
(0.00000, -1.00000)
(0.00000, 1.00000)
\end{exampleoutput}
\subsection{Polynomial Root Finding (Ehrlich-Aberth)}
Finds every complex root $x$ for a polynomial $p$ such that $p(x) = 0$. The Ehrlich-Aberth method is a simultaneous iteration: every root estimate is updated using both the Newton correction and a repulsion term from the other estimates.

This routine is intended for well-scaled polynomials when all complex roots are wanted. If only real roots are needed, a real-root isolator such as derivative recursion with bisection is usually more reliable. Multiple or tightly clustered roots may converge slowly, and very large or very small root scales may require rescaling the input or using multiprecision arithmetic.

\begin{compactitem}
  \item \inlinecode{eval\_with\_derivative(p, x)} returns a pair $(p(x), p'(x))$ for a polynomial $p$ given as a vector \inlinecode{p} where \inlinecode{p[i]} stores the coefficient for the $x^i$ term.
  \item \inlinecode{find\_all\_roots(p, eps = ROOT\_EPS, iterations = 2000)} returns a vector of all complex roots for a complex polynomial given by the vector of coefficients \inlinecode{p}. A \inlinecode{vector\textless{}dbl\textgreater{}} overload is provided for polynomials with real coefficients. The roots are found to a tolerance of \inlinecode{eps} in absolute or relative error (whichever is reached first), and zero roots are removed exactly before the simultaneous iteration starts.
\end{compactitem}

\Needspace*{3\baselineskip}
\textbf{Time Complexity}:
\begin{compactitem}
  \item $O(n)$ per call to \inlinecode{eval\_with\_derivative(p, x)}, where $n$ is the degree of the polynomial.
  \item $O(n^{2} t)$ per call to \inlinecode{find\_all\_roots(p, eps, iterations)}, where $n$ is the degree of the polynomial and $t$ is the number of iterations required to reach the desired precision.
\end{compactitem}

\Needspace*{3\baselineskip}
\textbf{Space Complexity}:
\begin{compactitem}
  \item $O(1)$ auxiliary for \inlinecode{eval\_with\_derivative(p, x)}.
  \item $O(n)$ auxiliary for \inlinecode{find\_all\_roots(p, eps, iterations)}, where $n$ is the degree of the polynomial.
\end{compactitem}

\begin{lstlisting}[language=C++]
#include <algorithm>
#include <cmath>
#include <complex>
#include <utility>
#include <vector>

using dbl = long double;
using cdbl = std::complex<dbl>;
using cpoly = std::vector<cdbl>;

const dbl PI = acosl(-1.0L);
const dbl ZERO_EPS = 1e-30L;   // Treat coefficients and denominators this small as zero.
const dbl ROOT_EPS = 1e-18L;   // Stop once root updates are this small relative to the root.
const dbl CHECK_EPS = 1e-12L;  // Residual tolerance used by the example assertions.

bool is_zero(const cdbl &z, const dbl eps = ZERO_EPS) {
  return std::abs(z) <= eps;
}

bool is_finite(const cdbl &z) {
  return std::isfinite(z.real()) && std::isfinite(z.imag());
}

std::pair<cdbl, cdbl> eval_with_derivative(const cpoly &p, const cdbl &x) {
  cdbl value = p.back(), derivative = 0;
  for (int i = static_cast<int>(p.size()) - 2; i >= 0; i--) {
    derivative = derivative * x + value;
    value = value * x + p[i];
  }
  return {value, derivative};
}

dbl root_bound(const cpoly &p) {
  dbl res = 0;
  int n = static_cast<int>(p.size()) - 1;
  for (int i = 0; i < n; i++) {
    dbl ratio = std::abs(p[i] / p.back());
    if (ratio > 0) {
      res = std::max(res, powl(ratio, 1.0L / (n - i)));
    }
  }
  return 2 * std::max(static_cast<dbl>(1), res);
}

bool root_less(const cdbl &a, const cdbl &b) {
  if (a.real() != b.real()) {
    return a.real() < b.real();
  }
  return a.imag() < b.imag();
}

cpoly find_all_roots(cpoly p, const dbl eps = ROOT_EPS, const int iterations = 2000) {
  while (!p.empty() && is_zero(p.back())) {
    p.pop_back();
  }
  cpoly roots;
  while (p.size() > 1 && is_zero(p[0])) {
    roots.push_back(0);
    p.erase(p.begin());
  }
  if (p.size() <= 1) {
    return roots;
  }
  dbl scale = 0;
  for (int i = 0; i < static_cast<int>(p.size()); i++) {
    scale = std::max(scale, std::abs(p[i]));
  }
  for (int i = 0; i < static_cast<int>(p.size()); i++) {
    p[i] /= scale;
  }
  int n = static_cast<int>(p.size()) - 1;
  if (n == 1) {
    roots.push_back(-p[0] / p[1]);
    return roots;
  }
  cpoly z(n);
  dbl radius = root_bound(p), offset = PI / (2 * n);
  dbl max_radius = 2 * radius;
  for (int i = 0; i < n; i++) {
    dbl angle = offset + 2 * PI * i / n;
    z[i] = radius * cdbl(cosl(angle), sinl(angle));
  }
  for (int it = 0; it < iterations; it++) {
    bool done = true;
    cpoly next = z;
    for (int i = 0; i < n; i++) {
      auto [fx, dfx] = eval_with_derivative(p, z[i]);
      if (std::abs(fx) <= eps) {
        continue;
      }
      cdbl repulsion = 0;
      for (int j = 0; j < n; j++) {
        if (i != j) {
          cdbl diff = z[i] - z[j];
          if (std::abs(diff) <= eps * eps) {
            dbl angle = 2 * PI * (i + 1) / (n + 1);
            diff = eps * (1 + std::abs(z[i])) * cdbl(cosl(angle), sinl(angle));
          }
          repulsion += cdbl(1) / diff;
        }
      }
      cdbl denom = dfx - fx * repulsion;
      if (is_zero(denom)) {
        done = false;
        continue;
      }
      cdbl step = fx / denom;
      if (!is_finite(step) && !is_zero(dfx)) {
        step = fx / dfx;
      }
      if (!is_finite(step)) {
        done = false;
        continue;
      }
      dbl limit = 2 * max_radius;
      if (std::abs(step) > limit) {
        step *= limit / std::abs(step);
      }
      next[i] = z[i] - step;
      if (!is_finite(next[i])) {
        next[i] = z[i];
        done = false;
        continue;
      }
      if (std::abs(next[i]) > max_radius) {
        next[i] *= max_radius / std::abs(next[i]);
      }
      if (std::abs(step) > eps * (1 + std::abs(next[i]))) {
        done = false;
      }
    }
    z = next;
    if (done) {
      break;
    }
  }
  for (int i = 0; i < n; i++) {
    roots.emplace_back(z[i]);
  }
  std::sort(roots.begin(), roots.end(), root_less);
  return roots;
}

std::vector<cdbl> find_all_roots(
    const std::vector<dbl> &p, const dbl eps = ROOT_EPS, const int iterations = 2000
) {
  cpoly q(p.begin(), p.end());
  return find_all_roots(q, eps, iterations);
}
\end{lstlisting}
\Needspace*{4\baselineskip}
\textbf{Example Usage}\par
\begin{lstlisting}[language=C++]
#include <cassert>
#include <cstdio>
using namespace std;

cdbl eval(const cpoly &p, const cdbl &x) {
  return eval_with_derivative(p, x).first;
}

void print_roots(vector<cdbl> x) {
  sort(x.begin(), x.end(), [](const cdbl &a, const cdbl &b) {
    return abs(a.real() - b.real()) > 0.5e-5 ? a.real() < b.real() : a.imag() < b.imag();
  });
  for (const auto &z : x) {
    dbl real = abs(z.real()) < 0.5e-5 ? 0 : z.real();
    dbl imag = abs(z.imag()) < 0.5e-5 ? 0 : z.imag();
    printf("(%.5Lf, %.5Lf)\n", real, imag);
  }
}

int main() {
  {  // -1 + 2x - 6x^2 + 2x^3
    printf("Roots of -1 + 2x - 6x^2 + 2x^3:\n");
    vector<dbl> p{-1.0, 2.0, -6.0, 2.0};
    vector<cdbl> roots = find_all_roots(p);
    assert(roots.size() == 3);
    for (const auto &root : roots) {
      assert(abs(eval(cpoly(p.begin(), p.end()), root)) < CHECK_EPS);
    }
    print_roots(roots);
  }
  {  // (-24+36i) + (-26+12i)x + (-30+40i)x^2 + (-26+12i)x^3 + (-6+4i)x^4
    // = ((2 + 3i)x + 6)(x + i)(2x + (6 + 4i))(xi + 1):
    printf("Roots of ((2 + 3i)x + 6)(x + i)(2x + (6 + 4i))(xi + 1):\n");
    cpoly p;
    p.emplace_back(-24, 36);
    p.emplace_back(-26, 12);
    p.emplace_back(-30, 40);
    p.emplace_back(-26, 12);
    p.emplace_back(-6, 4);
    vector<cdbl> roots = find_all_roots(p);
    assert(roots.size() == 4);
    for (const auto &root : roots) {
      assert(abs(eval(p, root)) < CHECK_EPS);
    }
    print_roots(roots);
  }
  return 0;
}
\end{lstlisting}
\begin{exampleoutput}
Roots of -1 + 2x - 6x^2 + 2x^3:
(0.15098, -0.40314)
(0.15098, 0.40314)
(2.69805, 0.00000)
Roots of ((2 + 3i)x + 6)(x + i)(2x + (6 + 4i))(xi + 1):
(-3.00000, -2.00000)
(-0.92308, 1.38462)
(0.00000, -1.00000)
(0.00000, 1.00000)
\end{exampleoutput}
\subsection{Integration (Simpson)}
Computes the definite integral from $a$ to $b$ for a continuous function $f$ using the composite Simpson's rule: the interval is split into $n$ equal subintervals and each consecutive pair is approximated by a parabola, giving the approximation $\int_{a}^{b} f(x) \,dx \approx$ $h/3 \cdot [f(x_0) + 4f(x_1) + 2f(x_2) + \ldots + 4f(x_{n-1}) + f(x_n)]$ with step $h = (b - a)/n$.

Unlike adaptive quadrature, the step size here is uniform, so a feature narrower than \inlinecode{h} may be under-resolved. The benefit is a simple, non-recursive routine with predictable cost that is not fooled by the premature-termination heuristic that adaptive Simpson's rule can suffer on functions whose coarse sample points happen to fit a parabola.

\begin{compactitem}
  \item \inlinecode{integrate(f, a, b, n = 1000000)} returns the definite integral of a function \inlinecode{f} from \inlinecode{a} to \inlinecode{b} using \inlinecode{n} subintervals, where \inlinecode{n} must be even. Larger \inlinecode{n} trades runtime for accuracy; the leading error scales with \inlinecode{(b - a) * h\textasciicircum{}4} and the maximum magnitude of the fourth derivative of $f$ over the interval.
\end{compactitem}

\textbf{Time Complexity}: $O(n)$ per call, requiring $n + 1$ evaluations of $f$.

\textbf{Space Complexity}: $O(1)$ auxiliary.

\begin{lstlisting}[language=C++]
#include <cassert>

template<typename Fn>
double integrate(Fn f, double a, double b, int n = 1000000) {
  assert(n > 0 && n % 2 == 0);
  double h = (b - a) / n, s = f(a) + f(b);
  for (int i = 1; i < n; i++) {
    s += f(a + h * i) * (i & 1 ? 4 : 2);
  }
  return s * h / 3;
}
\end{lstlisting}
\Needspace*{4\baselineskip}
\textbf{Example Usage}\par
\begin{lstlisting}[language=C++]
#include <cmath>
using namespace std;

double f(double x) {
  return sin(x);
}

int main() {
  const double PI = acos(-1.0);
  assert(fabs(integrate(f, 0.0, PI / 2) - 1) < 1e-10);
  assert(fabs(integrate(f, PI / 2, 0.0) + 1) < 1e-10);
  assert(fabs(integrate(f, 1.0, 1.0)) < 1e-12);
  return 0;
}
\end{lstlisting}
\subsection{Linear Programming (Simplex)}
Solves a linear programming problem using Dantzig's simplex algorithm. The canonical form of a linear programming problem is to maximize (or minimize) the dot product $cx$, subject to $ax \leq b$ and $x \geq 0$, where $x$ is a vector of unknowns to be solved, $c$ is a vector of coefficients, $a$ is a matrix of linear equation coefficients, and $b$ is a vector of boundary coefficients.

The simplex algorithm walks between vertices of the feasible polytope, at each step pivoting to an adjacent vertex that improves the objective, until no improving move remains (an optimum is reached) or an unbounded improving direction is found.

The implementation uses a two-phase simplex tableau. Phase 1 introduces an artificial variable to find a feasible starting basis, which is essential when some constraint bounds are negative. Phase 2 then optimizes the original objective. The \inlinecode{basis} and \inlinecode{nonbasis} arrays track which variable each tableau row or column currently represents, making solution recovery independent of pivot order.

\begin{compactitem}
  \item \inlinecode{simplex\_solve(a, b, c, \&x, maximize = true, eps = 1e-10)} solves the linear programming problem for an $m$ by $n$ matrix \inlinecode{a} of real values, a length $m$ vector \inlinecode{b}, and a length $n$ vector \inlinecode{c}, returning $0$ if an optimum was found, $-1$ if there are no feasible solutions, or $1$ if the objective is unbounded. Set \inlinecode{maximize = false} to minimize instead, and use \inlinecode{eps} as the pivot and feasibility tolerance. If an optimum is found, \inlinecode{x} is assigned a vector of length \inlinecode{n}, where \inlinecode{x[j]} is the value of the $j$-th variable.
\end{compactitem}

\textbf{Time Complexity}: $O(pmn)$, where $p$ is the total number of pivots across both phases. Each pivot takes $O(mn)$, and $p$ can be exponential in the worst case, although it is usually much smaller in practice.

\textbf{Space Complexity}: $O(m \cdot n)$ auxiliary.

\begin{lstlisting}[language=C++]
#include <cassert>
#include <cmath>
#include <utility>
#include <vector>

int simplex_solve(
    const std::vector<std::vector<double>> &a, const std::vector<double> &b,
    const std::vector<double> &c, std::vector<double> *x, const bool maximize = true,
    const double eps = 1e-10
) {
  int m = static_cast<int>(a.size()), n = static_cast<int>(c.size());
  assert(x != nullptr && b.size() == a.size());
  std::vector<int> basis(m), nonbasis(n + 1);
  std::vector<std::vector<double>> tab(m + 2, std::vector<double>(n + 2));
  for (int i = 0; i < m; i++) {
    for (int j = 0; j < n; j++) {
      tab[i][j] = a[i][j];
    }
    basis[i] = n + i;
    tab[i][n] = -1;
    tab[i][n + 1] = b[i];
  }
  for (int j = 0; j < n; j++) {
    nonbasis[j] = j;
    tab[m][j] = maximize ? -c[j] : c[j];
  }
  nonbasis[n] = -1;
  tab[m + 1][n] = 1;
  auto pivot = [&](int row, int col) {
    double inv = 1.0 / tab[row][col];
    for (int i = 0; i < m + 2; i++) {
      if (i == row || fabs(tab[i][col]) <= eps) {
        continue;
      }
      double factor = tab[i][col] * inv;
      for (int j = 0; j < n + 2; j++) {
        if (j != col) {
          tab[i][j] -= factor * tab[row][j];
        }
      }
      tab[i][col] = -factor;
    }
    for (int j = 0; j < n + 2; j++) {
      if (j != col) {
        tab[row][j] *= inv;
      }
    }
    tab[row][col] = inv;
    std::swap(basis[row], nonbasis[col]);
  };
  auto simplex = [&](int phase) {
    int objective = (phase == 1 ? m + 1 : m);
    while (true) {
      int col = -1;
      for (int j = 0; j <= n; j++) {
        if (phase == 2 && nonbasis[j] == -1) {
          continue;
        }
        if (col == -1 || tab[objective][j] < tab[objective][col] - eps ||
            (fabs(tab[objective][j] - tab[objective][col]) <= eps && nonbasis[j] < nonbasis[col])) {
          col = j;
        }
      }
      if (tab[objective][col] >= -eps) {
        return true;
      }
      int row = -1;
      for (int i = 0; i < m; i++) {
        if (tab[i][col] <= eps) {
          continue;
        }
        double cur = tab[i][n + 1] / tab[i][col];
        double best = row == -1 ? 0 : tab[row][n + 1] / tab[row][col];
        if (row == -1 || cur < best - eps || (fabs(cur - best) <= eps && basis[i] < basis[row])) {
          row = i;
        }
      }
      if (row == -1) {
        return false;
      }
      pivot(row, col);
    }
  };
  int row = 0;
  for (int i = 1; i < m; i++) {
    if (tab[i][n + 1] < tab[row][n + 1]) {
      row = i;
    }
  }
  if (m > 0 && tab[row][n + 1] < -eps) {
    pivot(row, n);
    if (!simplex(1) || fabs(tab[m + 1][n + 1]) > eps) {
      return -1;
    }
    for (int i = 0; i < m; i++) {
      if (basis[i] != -1) {
        continue;
      }
      int col = -1;
      for (int j = 0; j <= n; j++) {
        if (fabs(tab[i][j]) > eps && (col == -1 || nonbasis[j] < nonbasis[col])) {
          col = j;
        }
      }
      if (col != -1) {
        pivot(i, col);
      }
    }
  }
  if (!simplex(2)) {
    return 1;
  }
  x->assign(n, 0);
  for (int i = 0; i < m; i++) {
    if (basis[i] < n) {
      (*x)[basis[i]] = tab[i][n + 1];
    }
  }
  return 0;
}
\end{lstlisting}
\Needspace*{4\baselineskip}
\textbf{Example Usage}\par
\begin{lstlisting}[language=C++]
#include <iostream>
using namespace std;

int main() {
  // Solve [x, y] that maximizes 3x + 4y, subject to x, y >= 0 and:
  //  -2x +    1y <=  0
  //   1x + 0.85y <=  9
  //   1x +    2y <= 14
  vector<vector<double>> va{{-2, 1}, {1, 0.85}, {1, 2}};
  vector<double> vb{0, 9, 14};
  vector<double> vc{3, 4};
  vector<double> x;
  assert(simplex_solve(va, vb, vc, &x) == 0);
  double maxval = 0;
  for (int i = 0; i < static_cast<int>(x.size()); i++) {
    maxval += vc[i] * x[i];
  }
  cout << "Solution = " << maxval << " at (" << x[0];
  for (int i = 1; i < static_cast<int>(x.size()); i++) {
    cout << ", " << x[i];
  }
  cout << ")." << endl;

  // Feasible even though the first RHS is negative: x >= 1 and x <= 3.
  vector<vector<double>> negative_rhs{{-1}, {1}};
  vector<double> negative_bounds{-1, 3};
  vector<double> one_var{1};
  assert(simplex_solve(negative_rhs, negative_bounds, one_var, &x) == 0);
  assert(fabs(x[0] - 3) < 1e-9);

  // Infeasible: x <= -1 contradicts x >= 0.
  vector<vector<double>> infeasible{{1}};
  vector<double> bad_bounds{-1};
  assert(simplex_solve(infeasible, bad_bounds, one_var, &x) == -1);

  // Unbounded: maximize x subject only to x >= 0.
  vector<vector<double>> unbounded{{-1}};
  vector<double> lower_bound{0};
  assert(simplex_solve(unbounded, lower_bound, one_var, &x) == 1);
  return 0;
}
\end{lstlisting}
\begin{exampleoutput}
Solution = 33.3043 at (5.30435, 4.34783).
\end{exampleoutput}

\chapter{\texorpdfstring{Mathematics}{6. Mathematics}}

\section{\texorpdfstring{Math Utilities}{6.1 Math Utilities}}
\setcounter{section}{1}
\setcounter{subsection}{0}
\subsection{Floating-Point Comparison}
\label{sec:6.1.1}
Floating-point results carry rounding error, so two values that should be equal in exact arithmetic usually differ in their last bits. Comparing them with \inlinecode{==} is therefore a bug in almost every numeric setting, and the helpers below compare within a tolerance instead. This section also collects the mathematical constants and the sign inspections that the rest of the anthology uses.

\begin{compactitem}
  \item \inlinecode{M\_PI}, \inlinecode{M\_E}, \inlinecode{M\_PHI}, \inlinecode{M\_INF}, and \inlinecode{M\_NAN} are the usual constants, computed rather than assumed so that no platform-specific macro is required.
  \item \inlinecode{EQ()}, \inlinecode{NE()}, \inlinecode{LT()}, \inlinecode{GT()}, \inlinecode{LE()}, and \inlinecode{GE()} relationally compare two values $x$ and $y$. Arguments may be of different types; the common type governs behavior. If the common type is a floating-point type, exactly equal values (including same-signed infinities) compare equal; otherwise absolute-error epsilon comparison is used. Values within \inlinecode{EPS} of each other are considered equal, and \inlinecode{LT}/\inlinecode{GT}/\inlinecode{LE}/\inlinecode{GE} shift the boundary by \inlinecode{EPS} accordingly. Otherwise exact comparison is used (\inlinecode{==}, \inlinecode{\textless{}}, etc.). The branch is selected with \inlinecode{if constexpr}, so exact types without floating-point arithmetic (e.g. integers, \inlinecode{Modular}, \inlinecode{Rational}) compose as well.
  \item \inlinecode{rEQ(ref, val)} returns whether \inlinecode{val} equals reference \inlinecode{ref} within relative error \inlinecode{EPS}. The tolerance scales with $|{\inlinecodemath{ref}}|$, so \inlinecode{rEQ(ref, val)} is NOT the same as \inlinecode{rEQ(val, ref)}. Use this when one argument is a known exact value and the other is a computed approximation. Degenerates to exact comparison when \inlinecode{ref} is $0$, since the tolerance collapses to $0$; use \inlinecode{EQ} near zero.
  \item \inlinecode{rEQ\_sym(a, b)} is the symmetric (commutative) variant: tolerance scales with $\max(|{\inlinecodemath{a}}|, |{\inlinecodemath{b}}|)$, so the result is the same regardless of argument order. Still degenerates near zero when both arguments are close to $0$. For both relative comparisons, exactly equal infinities compare equal, while unequal infinities and comparisons between finite and non-finite values compare unequal. Non-floating common types use exact equality.
  \item \inlinecode{sgn(x)} returns $-1$ (if $x < 0$), $0$ (if $x = 0$), or $1$ (if $x > 0$). Unlike \inlinecode{std::signbit()} or \inlinecode{std::copysign()}, this does not handle the sign of \inlinecode{NaN}.
  \item \inlinecode{sign\_bit(x)} is analogous to \inlinecode{std::signbit()}, returning whether the sign bit of the floating point number is set to true. If so, then \inlinecode{x} is considered ``negative.'' Note that this works as expected on \inlinecode{+0.0}, \inlinecode{-0.0}, \inlinecode{Inf}, \inlinecode{-Inf}, \inlinecode{NaN}, as well as \inlinecode{-NaN}. Warning: This assumes that the sign bit is the leading (most significant) bit in the internal IEEE representation and that bytes are stored in little-endian order.
  \item \inlinecode{copy\_sign(x, y)} is analogous to \inlinecode{std::copysign()}, returning a number with the magnitude of \inlinecode{x} but the sign of \inlinecode{y}.
\end{compactitem}

\textbf{Time Complexity}: $O(1)$ per call to all operations.

\textbf{Space Complexity}: $O(1)$ auxiliary for all operations.

\begin{lstlisting}[style=cpp]
#include <algorithm>
#include <climits>
#include <cmath>
#include <limits>
#include <type_traits>

#ifndef M_PI
const double M_PI = std::acos(-1.0);  // Or std::numbers::pi in C++20 and later.
#endif
#ifndef M_E
const double M_E = std::exp(1.0);  // or std::numbers::e and std::numbers::e_v<> in C++20 and later
#endif
const double M_PHI = (1.0 + std::sqrt(5.0)) / 2.0;
const double M_INF = std::numeric_limits<double>::infinity();
const double M_NAN = std::numeric_limits<double>::quiet_NaN();

const double EPS = 1e-9;

template<typename T, typename U, typename C = std::common_type_t<T, U>>
bool EQ(T a, U b) {
  if constexpr (std::is_floating_point_v<C>) return C(a) == C(b) || std::fabs(C(a) - C(b)) <= EPS;
  return C(a) == C(b);
}

template<typename T, typename U, typename C = std::common_type_t<T, U>>
bool LT(T a, U b) {
  if constexpr (std::is_floating_point_v<C>) return C(a) < C(b) - EPS;
  return C(a) < C(b);
}

template<typename T, typename U> bool NE(T a, U b) { return !EQ(a, b); }
template<typename T, typename U> bool GT(T a, U b) { return LT(b, a); }
template<typename T, typename U> bool LE(T a, U b) { return !LT(b, a); }
template<typename T, typename U> bool GE(T a, U b) { return !LT(a, b); }

template<typename T, typename U, typename C = std::common_type_t<T, U>>
bool rEQ(T ref, U val) {
  C x = C(ref), y = C(val);
  if (x == y) return true;
  if constexpr (!std::is_floating_point_v<C>) {
    return false;
  }
  if (!std::isfinite(x) || !std::isfinite(y)) return false;
  return std::fabs(x - y) <= EPS * std::fabs(x);
}

template<typename T, typename U, typename C = std::common_type_t<T, U>>
bool rEQ_sym(T a, U b) {
  C x = C(a), y = C(b);
  if (x == y) return true;
  if constexpr (!std::is_floating_point_v<C>) {
    return false;
  }
  if (!std::isfinite(x) || !std::isfinite(y)) return false;
  return std::fabs(x - y) <= EPS * std::max(std::fabs(x), std::fabs(y));
}

template<typename T>
int sgn(const T &x) {
  return (T{0} < x) - (x < T{0});
}

template<typename Dbl>
bool sign_bit(Dbl x) {
  return (((unsigned char *)&x)[sizeof(x) - 1] >> (CHAR_BIT - 1)) & 1;
}

template<typename Dbl>
Dbl copy_sign(Dbl x, Dbl y) {
  return sign_bit(y) ? -std::fabs(x) : std::fabs(x);
}
\end{lstlisting}
\Needspace*{4\baselineskip}
\textbf{Example Usage}\par
\begin{lstlisting}[style=cpp]
#include <cassert>
using namespace std;

int main() {
  assert(EQ(M_PI, 3.14159265359));
  assert(EQ(M_E, 2.718281828459) && EQ(M_PHI, 1.61803398875));
  assert(EQ(M_INF, M_INF) && rEQ(M_INF, M_INF) && rEQ_sym(M_INF, M_INF));
  assert(!rEQ(M_INF, 0.0) && !rEQ_sym(M_INF, -M_INF));
  assert(!rEQ(1000000000000LL, 1000000000001LL));

  double x = -12345.6789;
  assert((-M_INF < x) && (x < M_INF));
  assert((M_INF + x == M_INF) && (M_INF - x == M_INF));
  assert((M_NAN != x) && (M_NAN != M_INF) && (M_NAN != M_NAN));
  assert(!(M_NAN < x) && !(M_NAN > x) && !(M_NAN <= x) && !(M_NAN >= x));
  assert(isnan(0.0 * M_INF) && isnan(M_INF - M_INF));

  assert(sgn(x) == -1 && sgn(0.0) == 0 && sgn(5678) == 1);
  assert(sign_bit(x) && !sign_bit(0.0) && sign_bit(-0.0));
  assert(!sign_bit(M_INF) && sign_bit(-M_INF));
  assert(!sign_bit(M_NAN) && sign_bit(-M_NAN));
  assert(copy_sign(1.0, +2.0) == +1.0 && copy_sign(M_INF, -2.0) == -M_INF);
  assert(copy_sign(1.0, -2.0) == -1.0 && sign_bit(copy_sign(M_NAN, -2.0)));
  return 0;
}
\end{lstlisting}
\subsection{Rounding}
\label{sec:6.1.2}
Rounding is only ambiguous at a tie, when a value sits exactly halfway between two integers, and the rule chosen for that case is what distinguishes the functions below. The choice matters more than it appears: always rounding ties away from zero biases a sum of many rounded values upward, which is why financial and statistical work usually rounds ties to even instead, spreading them evenly between the two neighbors.

\begin{compactitem}
  \item \inlinecode{floor0(x)} returns \inlinecode{x} rounded down, symmetrically towards zero. This function is analogous to \inlinecode{std::trunc()}.
  \item \inlinecode{ceil0(x)} returns \inlinecode{x} rounded up, symmetrically away from zero. It mirrors \inlinecode{floor0()} and has no \inlinecode{\textless{}cmath\textgreater{}} equivalent, since \inlinecode{std::ceil()} rounds towards positive infinity rather than outward.
  \item \inlinecode{round\_half\_up(x)} returns \inlinecode{x} rounded half up, towards positive infinity.
  \item \inlinecode{round\_half\_down(x)} returns \inlinecode{x} rounded half down, towards negative infinity.
  \item \inlinecode{round\_half\_to0(x)} returns \inlinecode{x} rounded half down, symmetrically towards zero.
  \item \inlinecode{round\_half\_from0(x)} returns \inlinecode{x} rounded half up, symmetrically away from zero. This function is analogous to \inlinecode{std::round()}.
  \item \inlinecode{round\_half\_even(x, eps = 1e-9)} returns \inlinecode{x} rounded half to even, using \inlinecode{eps} to detect ties for banker's rounding.
  \item \inlinecode{round\_half\_alternate(x)} returns \inlinecode{x} rounded, where ties are broken by alternating rounds towards positive and negative infinity.
  \item \inlinecode{round\_half\_alternate0(x)} returns \inlinecode{x} rounded, where ties are broken by alternating symmetric rounds towards and away from zero.
  \item \inlinecode{round\_half\_random(x)} returns \inlinecode{x} rounded, where ties are broken randomly.
  \item \inlinecode{round\_n\_places(x, n, round)} returns \inlinecode{x} rounded to \inlinecode{n} digits after the decimal, using the specified rounding function \inlinecode{round(x)}.
\end{compactitem}

Every function propagates \inlinecode{+0.0}, \inlinecode{-0.0}, \inlinecode{Inf}, \inlinecode{-Inf}, \inlinecode{NaN}, and \inlinecode{-NaN} exactly with the same sign, as \inlinecode{std::round()} does. The four built on adding or subtracting $0.5$ before flooring are otherwise exact except at $0.5 - 2^{-54}$, which rounds up to $1$ because the sum is already $1$ before flooring; prefer \inlinecode{std::round()} or \inlinecode{round\_half\_even()} where that last bit matters.

\textbf{Time Complexity}: $O(1)$ per call to all operations.

\textbf{Space Complexity}: $O(1)$ auxiliary for all operations.

\begin{lstlisting}[style=cpp]
#include <chrono>
#include <cmath>
#include <random>

template<typename Dbl>
Dbl floor0(const Dbl &x) {
  Dbl res = std::floor(std::fabs(x));
  return std::signbit(x) ? -res : res;
}

template<typename Dbl>
Dbl ceil0(const Dbl &x) {
  Dbl res = std::ceil(std::fabs(x));
  return std::signbit(x) ? -res : res;
}

template<typename Dbl>
Dbl round_half_up(const Dbl &x) {
  return std::copysign(std::floor(x + 0.5), x);  // copysign() only ever fixes a zero result.
}

template<typename Dbl>
Dbl round_half_down(const Dbl &x) {
  return std::copysign(std::ceil(x - 0.5), x);  // copysign() only ever fixes a zero result.
}

template<typename Dbl>
Dbl round_half_to0(const Dbl &x) {
  Dbl res = round_half_down(std::fabs(x));
  return std::signbit(x) ? -res : res;
}

template<typename Dbl>
Dbl round_half_from0(const Dbl &x) {
  Dbl res = round_half_up(std::fabs(x));
  return std::signbit(x) ? -res : res;
}

template<typename Dbl>
Dbl round_half_even(const Dbl &x, const Dbl &eps = 1e-9) {
  if (std::signbit(x)) {
    return -round_half_even(-x, eps);
  }
  Dbl ipart;
  std::modf(x, &ipart);
  if (std::fabs(x - (ipart + 0.5)) < eps) {  // exactly halfway: break the tie towards even
    return (std::fmod(ipart, 2.0) < eps) ? ipart : ceil0(ipart + 0.5);
  }
  return round_half_from0(x);
}

template<typename Dbl>
Dbl round_half_alternate(const Dbl &x) {
  Dbl up = round_half_up(x), down = round_half_down(x);
  if (std::isnan(x) || up == down) {  // A NaN must not consume a toggle: NaN == NaN is false.
    return up;
  }
  static bool round_up = false;
  return (round_up = !round_up) ? up : down;
}

template<typename Dbl>
Dbl round_half_alternate0(const Dbl &x) {
  Dbl away = round_half_from0(x), toward = round_half_to0(x);
  if (std::isnan(x) || away == toward) {  // A NaN must not consume a toggle.
    return away;
  }
  static bool round_away = false;
  return (round_away = !round_away) ? away : toward;
}

template<typename Dbl>
Dbl round_half_random(const Dbl &x) {
  Dbl away = round_half_from0(x), toward = round_half_to0(x);
  if (std::isnan(x) || away == toward) {
    return away;
  }
  static std::mt19937 rng(std::chrono::steady_clock::now().time_since_epoch().count());
  return (rng() % 2 == 0) ? away : toward;
}

template<typename Dbl, typename RoundFn>
Dbl round_n_places(const Dbl &x, unsigned int n, RoundFn round) {
  Dbl scale = std::pow(Dbl{10}, n);
  return round(x * scale) / scale;
}
\end{lstlisting}
\Needspace*{4\baselineskip}
\textbf{Example Usage}\par
\begin{lstlisting}[style=cpp]
#include <cassert>
#include <limits>
using namespace std;

bool EQ(double a, double b) {
  return fabs(a - b) < 1e-9;
}

int main() {
  assert(EQ(floor0(1.5), 1.0) && EQ(ceil0(1.5), 2.0));
  assert(EQ(floor0(-1.5), -1.0) && EQ(ceil0(-1.5), -2.0));
  // Both round outward at any fraction, unlike std::round() which rounds to the nearest.
  assert(EQ(floor0(1.8), 1.0) && EQ(ceil0(1.2), 2.0) && EQ(ceil0(-1.2), -2.0));

  // NaN, the infinities, and the sign of zero are all preserved, as std::round() does.
  const double M_INF = numeric_limits<double>::infinity();
  const double M_NAN = numeric_limits<double>::quiet_NaN();
  assert(isnan(round_half_even(M_NAN)) && isinf(ceil0(-M_INF)));
  assert(signbit(round_half_from0(-0.0)) && !signbit(round_half_down(0.0)));
  assert(EQ(round_half_up(+1.5), +2) && EQ(round_half_down(+1.5), +1));
  assert(EQ(round_half_up(-1.5), -1) && EQ(round_half_down(-1.5), -2));
  assert(EQ(round_half_to0(+1.5), +1) && EQ(round_half_from0(+1.5), +2));
  assert(EQ(round_half_to0(-1.5), -1) && EQ(round_half_from0(-1.5), -2));
  assert(EQ(round_half_even(+1.5), +2) && EQ(round_half_even(-1.5), -2));
  assert(EQ(round_half_even(3.1), 3) && EQ(round_half_even(3.4), 3));  // Non-ties round normally.

  double alt1 = round_half_alternate(+1.5);
  assert(EQ(round_half_alternate(+1.2), +1));  // Non-ties do not consume an alternating turn.
  double alt2 = round_half_alternate(+1.5);
  assert(!EQ(alt1, alt2));
  double alt01 = round_half_alternate0(-1.5);
  assert(EQ(round_half_alternate0(-1.2), -1));
  double alt02 = round_half_alternate0(-1.5);
  assert(!EQ(alt01, alt02));

  assert(EQ(round_n_places(-1.23456, 3, round_half_to0<double>), -1.235));
  return 0;
}
\end{lstlisting}
\subsection{Base Conversion}
\label{sec:6.1.3}
A positional numeral is a sequence of digits weighted by powers of its base, so converting between bases is repeated division: divide by the target base, record the remainder as the next digit, and continue with the quotient. Doing that division on a digit vector rather than on a machine integer is what lets \inlinecode{convert\_base()} handle values far larger than any built-in type. Roman numerals are the counterexample that shows what positional notation buys, being additive with a subtractive special case and having no digit for zero.

\begin{compactitem}
  \item \inlinecode{to\_base(x, b = 10)} returns the digits of the unsigned integer \inlinecode{x} in base \inlinecode{b}, where index $0$ of the result stores the least significant digit.
  \item \inlinecode{to\_roman(x)} returns the Roman numeral representation of the unsigned integer \inlinecode{x} as a string.
  \item \inlinecode{from\_roman(s)} returns the value of the Roman numeral \inlinecode{s}, which must consist of Roman digits. Subtractive pairs such as \inlinecode{"IX"} are handled, but the numeral is not otherwise validated, so a malformed one such as \inlinecode{"IC"} simply returns the value its symbols imply.
  \item \inlinecode{convert\_base(d, a, b)} converts an integer in base \inlinecode{a} as a vector \inlinecode{d} of digits (where \inlinecode{d[0]} is the least significant digit) to base \inlinecode{b} as a vector of digits (again with index $0$ holding the least significant digit). This uses repeated long division, so the value itself does not need to fit in a machine integer.
\end{compactitem}

Overflow warning: Each intermediate \inlinecode{rem*a + digit} must fit in \inlinecode{uint64\_t}.

\Needspace*{3\baselineskip}
\textbf{Time Complexity}:
\begin{compactitem}
  \item $O(\log_{b}\left(x + 1\right) + 1)$ per call to \inlinecode{to\_base()}.
  \item $O(x / 1000 + 1)$ per call to \inlinecode{to\_roman()} due to the repeated \inlinecode{M} prefix.
  \item $O(n)$ per call to \inlinecode{from\_roman()}, where $n$ is the length of the numeral.
  \item $O(d \cdot e)$ per call to \inlinecode{convert\_base(d, a, b)}, where $d$ is the number of input digits and $e$ is the number of output digits.
\end{compactitem}

\Needspace*{3\baselineskip}
\textbf{Space Complexity}:
\begin{compactitem}
  \item $O(e)$ auxiliary for \inlinecode{to\_base(x, b)}, where $e$ is the number of output digits.
  \item $O(x / 1000 + 1)$ auxiliary for \inlinecode{to\_roman(x)}.
  \item $O(1)$ auxiliary for \inlinecode{from\_roman()}.
  \item $O(d + e)$ auxiliary for \inlinecode{convert\_base(d, a, b)}.
\end{compactitem}

\begin{lstlisting}[style=cpp]
#include <algorithm>
#include <cassert>
#include <cstdint>
#include <map>
#include <string>
#include <utility>
#include <vector>

std::vector<int> to_base(unsigned int x, int b = 10) {
  assert(b >= 2);
  std::vector<int> res;
  do {  // do-while so that a value of 0 yields the single digit {0}
    res.push_back(x % b);
    x /= b;
  } while (x != 0);
  return res;
}

std::string to_roman(unsigned int x) {
  static const std::string h[] = {"", "C", "CC", "CCC", "CD", "D", "DC", "DCC", "DCCC", "CM"};
  static const std::string t[] = {"", "X", "XX", "XXX", "XL", "L", "LX", "LXX", "LXXX", "XC"};
  static const std::string o[] = {"", "I", "II", "III", "IV", "V", "VI", "VII", "VIII", "IX"};
  std::string prefix(x / 1000, 'M');
  x %= 1000;
  return prefix + h[x / 100] + t[x / 10 % 10] + o[x % 10];
}

unsigned int from_roman(const std::string &s) {
  static const std::map<char, unsigned int> VALUES = {{'I', 1},   {'V', 5},   {'X', 10},  {'L', 50},
                                                      {'C', 100}, {'D', 500}, {'M', 1000}};
  unsigned int res = 0;
  for (int i = 0; i < static_cast<int>(s.size()); i++) {
    auto it = VALUES.find(s[i]);
    assert(it != VALUES.end());
    auto next = i + 1 < static_cast<int>(s.size()) ? VALUES.find(s[i + 1]) : VALUES.end();
    // A smaller symbol before a larger one is subtracted, as the IV in XIV.
    res += next != VALUES.end() && it->second < next->second ? -it->second : it->second;
  }
  return res;
}

std::vector<int> convert_base(const std::vector<int> &d, int a, int b) {
  assert(a >= 2 && b >= 2);
  assert(std::all_of(d.begin(), d.end(), [a](int x) { return 0 <= x && x < a; }));
  std::vector<int> cur = d, res;
  auto trim = [](std::vector<int> &v) {
    while (v.size() > 1 && v.back() == 0) {
      v.pop_back();
    }
  };
  trim(cur);
  if (cur.empty() || (cur.size() == 1 && cur[0] == 0)) {
    return {0};
  }
  while (!(cur.size() == 1 && cur[0] == 0)) {
    std::vector<int> q(cur.size());
    uint64_t rem = 0;
    for (int i = static_cast<int>(cur.size()) - 1; i >= 0; i--) {
      uint64_t x = rem * static_cast<uint64_t>(a) + static_cast<uint64_t>(cur[i]);
      q[i] = static_cast<int>(x / static_cast<uint64_t>(b));
      rem = x % static_cast<uint64_t>(b);
    }
    res.push_back(static_cast<int>(rem));
    trim(q);
    cur = std::move(q);
  }
  return res;
}
\end{lstlisting}
\Needspace*{4\baselineskip}
\textbf{Example Usage}\par
\begin{lstlisting}[style=cpp]
#include <cassert>
using namespace std;

int main() {
  assert(to_roman(1234) == "MCCXXXIV");
  assert(to_roman(5678) == "MMMMMDCLXXVIII");
  assert(from_roman("MCCXXXIV") == 1234);
  assert(from_roman("IX") == 9 && from_roman("MMXXVI") == 2026);

  vector<int> digits{6, 5, 4, 3, 2, 1};
  assert(convert_base(to_base(123456, 20), 20, 10) == digits);
  assert(convert_base(vector<int>{0, 0, 0}, 10, 2) == vector<int>{0});

  vector<int> big_decimal(30, 9);  // 10^30 - 1, larger than uint64_t.
  assert(convert_base(convert_base(big_decimal, 10, 2), 2, 10) == big_decimal);
  return 0;
}
\end{lstlisting}
\subsection{Date and Time}
\label{sec:6.1.4}
Date questions become ordinary arithmetic once dates are converted to a day count. These functions use the proleptic Gregorian calendar, which applies the modern leap rule to every year, including those before the calendar was adopted, and they count days from the Unix epoch of 1970-01-01.

The conversion follows Howard Hinnant's closed-form algorithm, which shifts the year to start in March so that the leap day falls at the end of the year and never interrupts the month lengths. The remaining months then have lengths that repeat in a pattern whose running total is exactly $\lfloor (153m + 2)/5 \rfloor$ days, and whole $400$-year eras of $146097$ days handle the leap rule without a single division-by-$100$ special case. Both directions are branch-light formulas rather than loops over years, so a date thousands of years away costs the same as tomorrow.

\begin{compactitem}
  \item \inlinecode{is\_leap\_year(y)} returns whether year \inlinecode{y} has a February 29.
  \item \inlinecode{days\_in\_month(y, m)} returns the length of month \inlinecode{m} of year \inlinecode{y}, where \inlinecode{m} is in $[1, 12]$.
  \item \inlinecode{days\_from\_civil(y, m, d)} returns the number of days from 1970-01-01 to the date \inlinecode{y}-\inlinecode{m}-\inlinecode{d}, which is negative for earlier dates. The date must be valid.
  \item \inlinecode{civil\_from\_days(days)} returns the tuple (\inlinecode{year}, \inlinecode{month}, \inlinecode{day}) of the date that many days from 1970-01-01, inverting \inlinecode{days\_from\_civil()}.
  \item \inlinecode{day\_of\_week(y, m, d)} returns the weekday of the date, from $0$ for Sunday to $6$ for Saturday.
  \item \inlinecode{days\_between(y1, m1, d1, y2, m2, d2)} returns the number of days from the first date to the second, which is negative when the second date is earlier.
\end{compactitem}

Unix time counts seconds from the same epoch, and is exactly $86400$ times the day count plus the seconds elapsed within the day. That decomposition is exact only because Unix time ignores leap seconds: it is not a count of elapsed SI seconds but a label that repeats one value during a positive leap second, which is what keeps every day $86400$ seconds long.

\begin{compactitem}
  \item \inlinecode{seconds\_from\_hms(h, m, s)} returns the seconds elapsed since midnight, and \inlinecode{hms\_from\_seconds(seconds)} returns the tuple (\inlinecode{hour}, \inlinecode{minute}, \inlinecode{second}) inverting it.
  \item \inlinecode{timestamp\_from\_civil(y, m, d, hh = 0, mm = 0, ss = 0)} returns the Unix timestamp of that instant.
  \item \inlinecode{civil\_from\_timestamp(t)} returns the \inlinecode{DateTime} of a Unix timestamp, which may be negative for instants before the epoch. Its day count must fit in \inlinecode{int}.
\end{compactitem}

ISO 8601 numbers weeks from the one containing the first Thursday of the year, and assigns a date to the year of its own week's Thursday. Early January therefore often belongs to the previous ISO year and late December to the next, which is the mistake this replaces: 2021-01-01 is ISO week $53$ of $2020$.

\begin{compactitem}
  \item \inlinecode{iso\_week\_date(y, m, d)} returns the tuple (\inlinecode{iso\_year}, \inlinecode{week}, \inlinecode{weekday}) of the date, where \inlinecode{week} counts from $1$ and \inlinecode{weekday} runs from $1$ for Monday to $7$ for Sunday.
\end{compactitem}

Easter follows the anonymous Gregorian computus, which locates the paschal full moon from the year's place in the $19$-year Metonic cycle, corrects for the century's leap and lunar drift, then advances to the following Sunday.

\begin{compactitem}
  \item \inlinecode{easter(y)} returns the (\inlinecode{month}, \inlinecode{day}) of Easter Sunday in the Gregorian year \inlinecode{y}, which must be at least $1583$.
\end{compactitem}

Dates are ordinary \inlinecode{int} values with no range checking, and the era arithmetic stays within \inlinecode{int} for any year within a few million of the epoch. The proleptic calendar disagrees with history before a country's adoption of the Gregorian one in or after 1582, where problems usually specify Julian rules instead: the simpler leap rule of every fourth year.

\textbf{Time Complexity}: $O(1)$ per call to all operations.

\textbf{Space Complexity}: $O(1)$ auxiliary for all operations.

\begin{lstlisting}[style=cpp]
#include <cassert>
#include <climits>
#include <cstdint>
#include <tuple>
#include <utility>

struct DateTime {
  int year, month, day, hour, minute, second;
};

bool is_leap_year(int y) {
  return y % 4 == 0 && (y % 100 != 0 || y % 400 == 0);
}

int days_in_month(int y, int m) {
  assert(m >= 1 && m <= 12);
  static const int lengths[] = {31, 28, 31, 30, 31, 30, 31, 31, 30, 31, 30, 31};
  return m == 2 && is_leap_year(y) ? 29 : lengths[m - 1];
}

int days_from_civil(int y, int m, int d) {
  assert(1 <= m && m <= 12 && 1 <= d && d <= days_in_month(y, m));
  y -= m <= 2;  // Start the year in March, so that February 29 is the final day of a year.
  int era = (y >= 0 ? y : y - 399) / 400;
  int yoe = y - era * 400;  // Year of era, in [0, 399].
  int doy =
      (153 * (m + (m > 2 ? -3 : 9)) + 2) / 5 + d - 1;  // Day of March-based year, in [0, 365].
  int doe = yoe * 365 + yoe / 4 - yoe / 100 + doy;     // Day of era, in [0, 146096].
  return era * 146097 + doe - 719468;                  // 719468 is the day of era of 1970-01-01.
}

std::tuple<int, int, int> civil_from_days(int days) {
  days += 719468;
  int era = (days >= 0 ? days : days - 146096) / 146097;
  int doe = days - era * 146097;
  int yoe = (doe - doe / 1460 + doe / 36524 - doe / 146096) / 365;
  int y = yoe + era * 400;
  int doy = doe - (365 * yoe + yoe / 4 - yoe / 100);
  int mp = (5 * doy + 2) / 153;  // Month index of the March-based year, in [0, 11].
  int d = doy - (153 * mp + 2) / 5 + 1;
  int m = mp + (mp < 10 ? 3 : -9);
  return {y + (m <= 2), m, d};
}

int day_of_week(int y, int m, int d) {
  // 1970-01-01 was a Thursday, which is weekday 4.
  return ((days_from_civil(y, m, d) + 4) % 7 + 7) % 7;
}

int days_between(int y1, int m1, int d1, int y2, int m2, int d2) {
  return days_from_civil(y2, m2, d2) - days_from_civil(y1, m1, d1);
}

int seconds_from_hms(int h, int m, int s) {
  assert(h >= 0 && h < 24 && m >= 0 && m < 60 && s >= 0 && s < 60);
  return (h * 60 + m) * 60 + s;
}

std::tuple<int, int, int> hms_from_seconds(int seconds) {
  assert(seconds >= 0 && seconds < 86400);
  return {seconds / 3600, seconds / 60 % 60, seconds % 60};
}

int64_t timestamp_from_civil(int y, int m, int d, int hh = 0, int mm = 0, int ss = 0) {
  return 86400LL * days_from_civil(y, m, d) + seconds_from_hms(hh, mm, ss);
}

DateTime civil_from_timestamp(int64_t t) {
  int64_t days = t / 86400, rest = t % 86400;
  if (rest < 0) {  // Truncating division rounds toward zero, but the day count must floor.
    rest += 86400;
    days--;
  }
  assert(INT_MIN <= days && days <= INT_MAX);
  auto [y, m, d] = civil_from_days(static_cast<int>(days));
  auto [hh, mm, ss] = hms_from_seconds(static_cast<int>(rest));
  return {y, m, d, hh, mm, ss};
}

std::tuple<int, int, int> iso_week_date(int y, int m, int d) {
  int days = days_from_civil(y, m, d);
  int weekday = ((days + 3) % 7 + 7) % 7 + 1;  // The epoch was a Thursday, ISO weekday 4.
  int thursday = days - weekday + 4;           // This week's Thursday fixes the ISO year.
  int iso_year = std::get<0>(civil_from_days(thursday));
  return {iso_year, (thursday - days_from_civil(iso_year, 1, 1)) / 7 + 1, weekday};
}

std::pair<int, int> easter(int y) {
  assert(y >= 1583);
  int a = y % 19;                // Position in the Metonic cycle.
  int b = y / 100, c = y % 100;  // Century and year within the century.
  int d = b / 4, e = b % 4;      // Leap corrections for the century.
  int f = (b + 8) / 25, g = (b - f + 1) / 3;
  int h = (19 * a + b - d - g + 15) % 30;  // Days from March 21 to the paschal full moon.
  int i = c / 4, k = c % 4;
  int l = (32 + 2 * e + 2 * i - h - k) % 7;  // Days from the full moon to the next Sunday.
  int mm = (a + 11 * h + 22 * l) / 451;      // Correction for the two latest possible dates.
  return {(h + l - 7 * mm + 114) / 31, (h + l - 7 * mm + 114) % 31 + 1};
}
\end{lstlisting}
\Needspace*{4\baselineskip}
\textbf{Example Usage}\par
\begin{lstlisting}[style=cpp]
#include <cassert>
using namespace std;

int main() {
  assert(is_leap_year(2024) && is_leap_year(2000));
  assert(!is_leap_year(2023) && !is_leap_year(1900));
  assert(days_in_month(2024, 2) == 29 && days_in_month(2023, 2) == 28);
  assert(days_in_month(2023, 12) == 31 && days_in_month(2023, 4) == 30);

  assert(days_from_civil(1970, 1, 1) == 0);
  assert(days_from_civil(1969, 12, 31) == -1);
  assert(days_from_civil(2000, 1, 1) == 10957);
  assert(civil_from_days(0) == make_tuple(1970, 1, 1));
  assert(civil_from_days(-1) == make_tuple(1969, 12, 31));
  assert(civil_from_days(10957) == make_tuple(2000, 1, 1));

  assert(day_of_week(1970, 1, 1) == 4);   // A Thursday.
  assert(day_of_week(2000, 1, 1) == 6);   // A Saturday.
  assert(day_of_week(1969, 7, 20) == 0);  // A Sunday.

  // 2024 was a leap year, so February contributed 29 days.
  assert(days_between(2024, 1, 1, 2025, 1, 1) == 366);
  assert(days_between(2023, 1, 1, 2024, 1, 1) == 365);
  assert(days_between(2024, 3, 1, 2024, 2, 1) == -29);

  assert(seconds_from_hms(23, 31, 30) == 84690);
  assert(hms_from_seconds(84690) == make_tuple(23, 31, 30));
  assert(timestamp_from_civil(1970, 1, 1) == 0);
  assert(timestamp_from_civil(2009, 2, 13, 23, 31, 30) == 1234567890);

  auto [y, mo, d, hh, mm, ss] = civil_from_timestamp(1234567890);
  assert(y == 2009 && mo == 2 && d == 13 && hh == 23 && mm == 31 && ss == 30);
  auto before = civil_from_timestamp(-1);  // One second before the epoch.
  assert(before.year == 1969 && before.month == 12 && before.day == 31);
  assert(before.hour == 23 && before.minute == 59 && before.second == 59);

  assert(iso_week_date(1970, 1, 1) == make_tuple(1970, 1, 4));
  // An ISO year runs from the Monday of the week holding the first Thursday, so it straddles
  // the calendar year on both sides.
  assert(iso_week_date(2019, 12, 30) == make_tuple(2020, 1, 1));
  assert(iso_week_date(2021, 1, 1) == make_tuple(2020, 53, 5));
  assert(iso_week_date(2026, 8, 17) == make_tuple(2026, 34, 1));

  assert(easter(2000) == make_pair(4, 23));
  assert(easter(2024) == make_pair(3, 31));
  assert(easter(2025) == make_pair(4, 20));
  auto [month, day] = easter(2026);
  assert(day_of_week(2026, month, day) == 0);  // Easter always falls on a Sunday.

  // The century rule, where 2000 is a leap year but 1900 is not.
  assert(days_between(2000, 2, 28, 2000, 3, 1) == 2);
  assert(days_between(1900, 2, 28, 1900, 3, 1) == 1);
  assert(civil_from_days(days_from_civil(2000, 2, 29)) == make_tuple(2000, 2, 29));
  assert(civil_from_days(days_from_civil(1, 1, 1)) == make_tuple(1, 1, 1));
  assert(civil_from_days(days_from_civil(-44, 3, 15)) == make_tuple(-44, 3, 15));
  return 0;
}
\end{lstlisting}

\section{\texorpdfstring{Combinatorics}{6.2 Combinatorics}}
\setcounter{section}{2}
\setcounter{subsection}{0}
\subsection{Combinatorial Calculations}
\label{sec:6.2.1}
The following functions implement common operations in combinatorics. All size and count inputs must be nonnegative, and all moduli must be positive. All return values and table entries are computed as 64-bit integers modulo an input argument $m$ or $p$.

\begin{compactitem}
  \item \inlinecode{factorial(n, m = MOD)} returns $n! \bmod m$.
  \item \inlinecode{factorial\_without\_p(n, p = MOD)} returns $n!$ modulo the prime $p$ after removing every factor of $p$ from the factorial.
  \item \inlinecode{legendre(n, p)} returns the exponent of the prime $p$ in the factorization of $n!$, where $p$ must be at least $2$.
  \item \inlinecode{binomial\_table(n, m = MOD)} returns the first $n + 1$ rows of Pascal's triangle as a two-dimensional vector $t$ such that $t[i][j] = \binom{i}{j} \bmod m$.
  \item \inlinecode{permute(n, k, m = MOD)} returns $(n \mathbin{\text{permute}} k) \bmod m$.
  \item \inlinecode{choose(n, k, p = MOD)} returns $\binom{n}{k} \bmod p$, where $p$ is prime and $n < p$.
  \item \inlinecode{choose\_lucas(n, k, p = MOD)} returns $\binom{n}{k} \bmod p$ for any 64-bit \inlinecode{n} and \inlinecode{k}, where $p$ is prime. It is only practical for a small $p$, since each base-$p$ digit costs a call to \inlinecode{choose()}.
  \item \inlinecode{multichoose(n, k, p = MOD)} returns $(n \mathbin{\text{multichoose}} k) \bmod p$, where $p$ is prime and $n + k - 1 < p$ when $n > 0$.
  \item \inlinecode{catalan(n, p = MOD)} returns the $n$-th Catalan number mod $p$, where $p$ is prime and $2n < p$.
  \item \inlinecode{partitions(n, m = MOD)} returns the number of partitions of $n$, mod $m$.
  \item \inlinecode{partitions(n, k, m = MOD)} returns the number of partitions of $n$ into $k$ parts, mod $m$.
  \item \inlinecode{stirling1(n, k, m = MOD)} returns the $(n, k)$ unsigned Stirling number of the 1st kind mod $m$.
  \item \inlinecode{stirling2(n, k, m = MOD)} returns the $(n, k)$ Stirling number of the 2nd kind mod $m$.
  \item \inlinecode{eulerian1(n, k, m = MOD)} returns the $(n, k)$ Eulerian number of the 1st kind mod $m$, where $n > k$.
  \item \inlinecode{eulerian2(n, k, m = MOD)} returns the $(n, k)$ Eulerian number of the 2nd kind mod $m$, where $n > k$.
\end{compactitem}

Legendre's formula sums $\lfloor n/p^i \rfloor$ over every $i \geq 1$, counting how many of $1, \dots, n$ are divisible by each power of $p$, which totals the exponent of $p$ in $n!$. It answers questions such as how many trailing zeros $n!$ has in base $p$. Paired with Kummer's theorem, which equates the exponent of $p$ in $\binom{n}{k}$ with the number of carries when adding $k$ and $n - k$ in base $p$, it also decides when a binomial coefficient is divisible by $p$.

Lucas' theorem reduces a binomial coefficient modulo a prime to a product over base-$p$ digits: $\binom{n}{k} \equiv \prod_i \binom{n_i}{k_i} \pmod p$, where $n_i$ and $k_i$ are the $i$-th digits of $n$ and $k$ in base $p$. The product is $0$ as soon as some $k_i$ exceeds $n_i$. For a composite modulus, factor it into prime powers, apply the generalized version of the theorem to each factor, and recombine the residues with the Chinese remainder theorem of section \secref{6.3.2}.

Every routine here recomputes its answer from scratch, which is what suits a handful of values or a modulus that varies between calls. When a solution instead needs thousands of binomial coefficients against one fixed modulus, precompute factorials and inverse factorials: the combinatorics tables of section \secref{6.3.3} amortize $O(n)$ of table growth and then answer each query in $O(1)$.

Overflow warning: Modular products use ordinary \inlinecode{int64\_t} multiplication, so the square of the chosen modulus must fit in \inlinecode{int64\_t}.

\Needspace*{3\baselineskip}
\textbf{Time Complexity}:
\begin{compactitem}
  \item $O(n)$ per call to \inlinecode{factorial()}.
  \item $O(p \log_{p}\left(n\right))$ per call to \inlinecode{factorial\_without\_p()}.
  \item $O(\log_{p}\left(n\right))$ per call to \inlinecode{legendre()}.
  \item $O(n^{2})$ per call to \inlinecode{binomial\_table()}.
  \item $O(k)$ per call to \inlinecode{permute()}.
  \item $O(\min\left(k, n - k\right))$ per call to \inlinecode{choose()}.
  \item $O(p \log_{p}\left(n\right))$ per call to \inlinecode{choose\_lucas()}.
  \item $O(k)$ per call to \inlinecode{multichoose()}.
  \item $O(n)$ per call to \inlinecode{catalan()}.
  \item $O(n^{2})$ per call to \inlinecode{partitions(n, m)}.
  \item $O(n \cdot k)$ per call to \inlinecode{partitions(n, k, m)}, \inlinecode{stirling1()}, \inlinecode{stirling2()}, \inlinecode{eulerian1()}, and \inlinecode{eulerian2()}.
\end{compactitem}

\Needspace*{3\baselineskip}
\textbf{Space Complexity}:
\begin{compactitem}
  \item $O(n^{2})$ auxiliary for \inlinecode{binomial\_table(n, m)}.
  \item $O(n \cdot k)$ auxiliary for \inlinecode{partitions(n, k, m)}, \inlinecode{stirling1(n, k, m)}, \inlinecode{stirling2(n, k, m)}, \inlinecode{eulerian1(n, k, m)}, and \inlinecode{eulerian2(n, k, m)}.
  \item $O(1)$ auxiliary for all other operations.
\end{compactitem}

\begin{lstlisting}[style=cpp]
#include <cassert>
#include <cstdint>
#include <vector>

const int64_t MOD = 1000000007;

int64_t factorial(int n, int64_t m = MOD) {
  assert(n >= 0 && m > 0);
  int64_t res = 1;
  for (int i = 2; i <= n; i++) {
    res = (res * i) % m;
  }
  return res % m;
}

int64_t factorial_without_p(int64_t n, int64_t p = MOD) {
  assert(n >= 0 && p >= 2);
  int64_t res = 1;
  while (n > 1) {
    if (n / p % 2 == 1) {
      res = res * (p - 1) % p;
    }
    int64_t max = n % p;
    for (int64_t i = 2; i <= max; i++) {
      res = (res * i) % p;
    }
    n /= p;
  }
  return res % p;
}

int64_t legendre(int64_t n, int64_t p) {
  assert(n >= 0 && p >= 2);
  int64_t res = 0;
  for (int64_t q = n / p; q > 0; q /= p) {  // Dividing repeatedly avoids overflowing p^i.
    res += q;
  }
  return res;
}

std::vector<std::vector<int64_t>> binomial_table(int n, int64_t m = MOD) {
  assert(n >= 0 && m > 0);
  std::vector<std::vector<int64_t>> t(n + 1);
  for (int i = 0; i <= n; i++) {
    for (int j = 0; j <= i; j++) {
      if (i < 2 || j == 0 || i == j) {
        t[i].push_back(1);
      } else {
        t[i].push_back((t[i - 1][j - 1] + t[i - 1][j]) % m);
      }
    }
  }
  return t;
}

int64_t permute(int n, int k, int64_t m = MOD) {
  assert(n >= 0 && k >= 0 && m > 0);
  if (n < k) {
    return 0;
  }
  int64_t res = 1;
  for (int i = 0; i < k; i++) {
    res = res * (n - i) % m;
  }
  return res % m;
}

int64_t powmod(int64_t x, int64_t n, int64_t m) {
  int64_t res = 1 % m;
  for (x %= m; n > 0; n >>= 1) {
    if (n & 1) {
      res = res * x % m;
    }
    x = x * x % m;
  }
  return res;
}

int64_t choose(int64_t n, int64_t k, int64_t p = MOD) {
  assert(n >= 0 && k >= 0 && p >= 2 && n < p);
  if (n < k) {
    return 0;
  }
  if (k > n - k) {
    k = n - k;
  }
  int64_t num = 1, den = 1;
  for (int64_t i = 0; i < k; i++) {
    num = num * (n - i) % p;
  }
  for (int64_t i = 1; i <= k; i++) {
    den = den * i % p;
  }
  return num * powmod(den, p - 2, p) % p;
}

int64_t choose_lucas(int64_t n, int64_t k, int64_t p = MOD) {
  assert(n >= 0 && p >= 2);
  if (k < 0 || k > n) {
    return 0;
  }
  int64_t res = 1;
  while (n > 0) {  // Since k never exceeds n, k is also exhausted when n reaches 0.
    int64_t ni = n % p, ki = k % p;
    if (ki > ni) {
      return 0;
    }
    res = res * choose(ni, ki, p) % p;
    n /= p;
    k /= p;
  }
  return res;
}

int64_t multichoose(int n, int k, int64_t p = MOD) {
  assert(n >= 0 && k >= 0 && p >= 2);
  if (n == 0) {
    return k == 0;
  }
  return choose(static_cast<int64_t>(n) + k - 1, k, p);
}

int64_t catalan(int n, int64_t p = MOD) {
  assert(n >= 0 && p >= 2);
  return choose(2LL * n, n, p) * powmod(n + 1, p - 2, p) % p;
}

int64_t partitions(int n, int64_t m = MOD) {
  assert(n >= 0 && m > 0);
  std::vector<int64_t> t(n + 1);
  t[0] = 1;
  for (int i = 1; i <= n; i++) {
    for (int j = i; j <= n; j++) {
      t[j] = (t[j] + t[j - i]) % m;
    }
  }
  return t[n] % m;
}

int64_t partitions(int n, int k, int64_t m = MOD) {
  assert(n >= 0 && k >= 0 && m > 0);
  std::vector<std::vector<int64_t>> t(n + 1, std::vector<int64_t>(k + 1));
  t[0][0] = 1;
  for (int i = 1; i <= n; i++) {
    for (int j = 1, h = k < i ? k : i; j <= h; j++) {
      t[i][j] = (t[i - 1][j - 1] + t[i - j][j]) % m;
    }
  }
  return t[n][k] % m;
}

int64_t stirling1(int n, int k, int64_t m = MOD) {
  assert(n >= 0 && k >= 0 && m > 0);
  std::vector<std::vector<int64_t>> t(n + 1, std::vector<int64_t>(k + 1));
  t[0][0] = 1;
  for (int i = 1; i <= n; i++) {
    for (int j = 1; j <= k; j++) {
      t[i][j] = (i - 1) * t[i - 1][j] % m;
      t[i][j] = (t[i][j] + t[i - 1][j - 1]) % m;
    }
  }
  return t[n][k] % m;
}

int64_t stirling2(int n, int k, int64_t m = MOD) {
  assert(n >= 0 && k >= 0 && m > 0);
  std::vector<std::vector<int64_t>> t(n + 1, std::vector<int64_t>(k + 1));
  t[0][0] = 1;
  for (int i = 1; i <= n; i++) {
    for (int j = 1; j <= k; j++) {
      t[i][j] = j * t[i - 1][j] % m;
      t[i][j] = (t[i][j] + t[i - 1][j - 1]) % m;
    }
  }
  return t[n][k] % m;
}

int64_t eulerian1(int n, int k, int64_t m = MOD) {
  assert(n > k && k >= 0 && m > 0);
  if (k > n - 1 - k) {
    k = n - 1 - k;
  }
  std::vector<std::vector<int64_t>> t(n + 1, std::vector<int64_t>(k + 1, 1));
  for (int j = 1; j <= k; j++) {
    t[0][j] = 0;
  }
  for (int i = 1; i <= n; i++) {
    for (int j = 1; j <= k; j++) {
      t[i][j] = (i - j) * t[i - 1][j - 1] % m;
      t[i][j] = (t[i][j] + ((j + 1) * t[i - 1][j] % m)) % m;
    }
  }
  return t[n][k] % m;
}

int64_t eulerian2(int n, int k, int64_t m = MOD) {
  assert(n > k && k >= 0 && m > 0);
  std::vector<std::vector<int64_t>> t(n + 1, std::vector<int64_t>(k + 1, 1));
  for (int i = 1; i <= n; i++) {
    for (int j = 1; j <= k; j++) {
      if (i == j) {
        t[i][j] = 0;
      } else {
        t[i][j] = (j + 1) * t[i - 1][j] % m;
        t[i][j] = ((2 * i - 1 - j) * t[i - 1][j - 1] % m + t[i][j]) % m;
      }
    }
  }
  return t[n][k] % m;
}
\end{lstlisting}
\Needspace*{4\baselineskip}
\textbf{Example Usage}\par
\begin{lstlisting}[style=cpp]
#include <cassert>

int main() {
  auto t = binomial_table(10);
  for (int i = 0; i < static_cast<int>(t.size()); i++) {
    for (int j = 0; j < static_cast<int>(t[i].size()); j++) {
      assert(t[i][j] == choose(i, j));
    }
  }
  assert(factorial(10) == 3628800);
  assert(factorial_without_p(123456) == 639390503);
  assert(factorial_without_p(7, 5) == 3);  // 7! / 5 = 1008 = 3 (mod 5).
  assert(legendre(10, 2) == 8);            // 10! = 2^8 * 3^4 * 5^2 * 7.
  assert(legendre(100, 5) == 24);          // 100! ends in 24 zeros.
  assert(permute(10, 4) == 5040);
  assert(choose(20, 7) == 77520);
  assert(choose_lucas(20, 7) == 77520);
  auto mod7 = binomial_table(30, 7);
  for (int i = 0; i <= 30; i++) {
    for (int j = 0; j <= i; j++) {
      assert(choose_lucas(i, j, 7) == mod7[i][j]);
    }
  }
  // By Lucas' theorem mod 2, a binomial coefficient is odd exactly when k is a submask of n.
  assert(choose_lucas((1LL << 40) - 1, 12345, 2) == 1);
  assert(choose_lucas(1LL << 40, 1, 2) == 0);
  assert(multichoose(20, 7) == 657800);
  assert(catalan(10) == 16796);
  assert(partitions(4) == 5);
  assert(partitions(0, 0) == 1);
  assert(partitions(100, 5) == 38225);
  assert(stirling1(4, 2) == 11);
  assert(stirling2(4, 3) == 6);
  assert(eulerian1(9, 5) == 88234);
  assert(eulerian2(8, 3) == 195800);
  return 0;
}
\end{lstlisting}
\subsection{Enumerating Arrangements}
\label{sec:6.2.2}
For the purposes of this section, we define a ``size $k$ arrangement of $n$'' to be a permutation of a size $k$ subset of the integers in $[0, n)$, for $0 \leq k \leq n$. There are $n \mathbin{\text{permute}} k$ possible arrangements, but $n^k$ possible arrangements if repeated values are allowed.

For arrangements without repeats, each prefix fixes some used values and leaves a shrinking set of choices for the remaining positions. Ranking counts how many unused smaller choices could have been placed at the current position, multiplied by the number of possible suffix arrangements; unranking performs the inverse search by subtracting those block sizes. With repeats allowed, the sequence is just a base-$n$ counter.

\begin{compactitem}
  \item \inlinecode{next\_arrangement(n, a)} tries to rearrange \inlinecode{a} to the next lexicographically greater arrangement, returning true if such an arrangement exists or false if the array is already in descending order (in which case \inlinecode{a} is unchanged). The input \inlinecode{a} must consist of distinct integers in the range $[0, {\inlinecodemath{n}})$.
  \item \inlinecode{arrangement\_by\_rank(n, k, r)} returns the size $k$ arrangement of $n$ which is lexicographically ranked $r$ out of all size $k$ arrangements of $n$, where $r$ is a 0-based rank in the range $[0, n \mathbin{\text{permute}} k)$.
  \item \inlinecode{rank\_by\_arrangement(n, a)} returns an integer representing the 0-based rank of arrangement \inlinecode{a}, which must consist of distinct integers in the range $[0, {\inlinecodemath{n}})$.
  \item \inlinecode{next\_arrangement\_with\_repeats(n, a)} tries to rearrange \inlinecode{a} to the next lexicographically greater arrangement with repeats, returning true if such an arrangement exists or false if the array is already in descending order (in which case \inlinecode{a} is unchanged). The input \inlinecode{a} must consist of integers in the range $[0, {\inlinecodemath{n}})$. If \inlinecode{a} were interpreted as a $k$ digit integer in base $n$, this function could be thought of as incrementing the integer.
\end{compactitem}

Overflow warning: Exact counts and ranks used by the ranking operations must fit in \inlinecode{int64\_t}.

\Needspace*{3\baselineskip}
\textbf{Time Complexity}:
\begin{compactitem}
  \item $O(n \cdot k)$ per call to \inlinecode{next\_arrangement()}, \inlinecode{arrangement\_by\_rank()}, and \inlinecode{rank\_by\_arrangement(n, a)}, where $k$ is the size of \inlinecode{a}.
  \item $O(k)$ per call to \inlinecode{next\_arrangement\_with\_repeats()}.
\end{compactitem}

\Needspace*{3\baselineskip}
\textbf{Space Complexity}:
\begin{compactitem}
  \item $O(n)$ auxiliary for \inlinecode{next\_arrangement()}, \inlinecode{arrangement\_by\_rank()}, and \inlinecode{rank\_by\_arrangement()}.
  \item $O(1)$ auxiliary for \inlinecode{next\_arrangement\_with\_repeats()}.
\end{compactitem}

\begin{lstlisting}[style=cpp]
#include <algorithm>
#include <cassert>
#include <cstdint>
#include <numeric>
#include <vector>

bool next_arrangement(int n, std::vector<int> &a) {
  int k = static_cast<int>(a.size());
  assert(n >= 0 && k <= n);
  std::vector<char> used(n);
  for (int i = 0; i < k; i++) {
    assert(0 <= a[i] && a[i] < n && !used[a[i]]);
    used[a[i]] = true;
  }
  for (int i = k - 1; i >= 0; i--) {
    used[a[i]] = false;
    for (int j = a[i] + 1; j < n; j++) {
      if (!used[j]) {
        a[i++] = j;
        used[j] = true;
        for (int x = 0; i < k; x++) {
          if (!used[x]) {
            a[i++] = x;
          }
        }
        return true;
      }
    }
  }
  return false;
}

int64_t n_permute_k(int n, int k) {
  assert(n >= 0 && 0 <= k && k <= n);
  int64_t res = 1;
  for (int i = 0; i < k; i++) {
    res *= n - i;
  }
  return res;
}

std::vector<int> arrangement_by_rank(int n, int k, int64_t r) {
  assert(0 <= r && r < n_permute_k(n, k));
  std::vector<int> values(n), res(k);
  std::iota(values.begin(), values.end(), 0);
  for (int i = 0; i < k; i++) {
    int64_t count = n_permute_k(n - 1 - i, k - 1 - i);
    int pos = r / count;
    res[i] = values[pos];
    values.erase(values.begin() + pos);
    r %= count;
  }
  return res;
}

int64_t rank_by_arrangement(int n, const std::vector<int> &a) {
  int k = static_cast<int>(a.size());
  assert(n >= 0 && k <= n);
  assert(std::all_of(a.begin(), a.end(), [n](int x) { return 0 <= x && x < n; }));
  int64_t res = 0;
  std::vector<char> used(n);
  for (int i = 0; i < k; i++) {
    assert(!used[a[i]]);
    int count = 0;
    for (int j = 0; j < a[i]; j++) {
      if (!used[j]) {
        count++;
      }
    }
    res += count * n_permute_k(n - i - 1, k - i - 1);
    used[a[i]] = true;
  }
  return res;
}

bool next_arrangement_with_repeats(int n, std::vector<int> &a) {
  int k = static_cast<int>(a.size());
  assert(n >= 0 && std::all_of(a.begin(), a.end(), [n](int x) { return 0 <= x && x < n; }));
  for (int i = k - 1; i >= 0; i--) {
    if (a[i] < n - 1) {
      a[i]++;
      std::fill(a.begin() + i + 1, a.end(), 0);
      return true;
    }
  }
  return false;
}
\end{lstlisting}
\Needspace*{4\baselineskip}
\textbf{Example Usage}\par
\begin{lstlisting}[style=cpp]
#include <cassert>
#include <iostream>
using namespace std;

template<typename It>
void print_range(It lo, It hi) {
  cout << "{";
  for (; lo != hi; ++lo) {
    cout << *lo << (lo == hi - 1 ? "" : ",");
  }
  cout << "} ";
}

int main() {
  {
    int n = 4, k = 3;
    vector<int> a{0, 1, 2};
    cout << n << " permute " << k << " arrangements:" << endl;
    int count = 0;
    do {
      print_range(a.begin(), a.end());
      assert(a == arrangement_by_rank(n, k, count));
      assert(rank_by_arrangement(n, a) == count);
      count++;
      if (count == 10 || count == 20) {
        cout << endl;
      }
    } while (next_arrangement(n, a));
    assert(count == 24);
    cout << endl;
  }
  {
    int n = 4, k = 2;
    vector<int> a{0, 0};
    cout << endl << n << "^" << k << " arrangements with repeats:" << endl;
    int count = 0;
    do {
      print_range(a.begin(), a.end());
      count++;
      if (count == 13) {
        cout << endl;
      }
    } while (next_arrangement_with_repeats(n, a));
    assert(count == 16);
    cout << endl;
  }
  return 0;
}
\end{lstlisting}
\begin{exampleoutput}
4 permute 3 arrangements:
{0,1,2} {0,1,3} {0,2,1} {0,2,3} {0,3,1} {0,3,2} {1,0,2} {1,0,3} {1,2,0} {1,2,3}
{1,3,0} {1,3,2} {2,0,1} {2,0,3} {2,1,0} {2,1,3} {2,3,0} {2,3,1} {3,0,1} {3,0,2}
{3,1,0} {3,1,2} {3,2,0} {3,2,1}

4^2 arrangements with repeats:
{0,0} {0,1} {0,2} {0,3} {1,0} {1,1} {1,2} {1,3} {2,0} {2,1} {2,2} {2,3} {3,0}
{3,1} {3,2} {3,3}
\end{exampleoutput}
\subsection{Enumerating Permutations}
\label{sec:6.2.3}
These routines enumerate and analyze reorderings of a sequence of $n$ elements; the values need not be distinct except where noted.

The lexicographic successor is found by locating the rightmost ascent, increasing that position by the smallest possible amount, then reversing the suffix into its minimum order. Ranking and unranking permutations use the factorial number system: each position contributes the number of unused smaller values before the chosen one, scaled by the number of possible suffix permutations. The cycle decomposition views a permutation as a function on indices and follows each unvisited orbit until it returns to its start. Sorting by swaps is that same decomposition read differently: a cycle of length $k$ is resolved by $k - 1$ swaps and no fewer, so the whole array needs $n$ minus its number of cycles.

\begin{compactitem}
  \item \inlinecode{next\_permutation2(lo, hi, comp = std::less\textless{}\textgreater{}())} is analogous to \inlinecode{std::next\_permutation()}, taking two BidirectionalIterators as a range $[{\inlinecodemath{lo}}, {\inlinecodemath{hi}})$ for which the function tries to rearrange to the next lexicographically greater permutation according to \inlinecode{comp}. The function returns true if such a permutation exists, or false if the range is already in reverse comparator order, in which case it is rearranged into comparator order.
  \item \inlinecode{next\_permutation(a, comp = std::less\textless{}\textgreater{}())} is analogous to \inlinecode{next\_permutation2()}, except that it takes a vector instead of a range.
  \item \inlinecode{next\_permutation\_mask(x)} returns the next integer having the same number of 1-bits. Treating each 1-bit as whether to take its corresponding item generates combinations of a set of $n$ items. It returns $0$ if no successor fits in \inlinecode{uint64\_t}.
  \item \inlinecode{permutation\_by\_rank(n, r)} returns the permutation of the integers in the range $[0, {\inlinecodemath{n}})$ which is lexicographically ranked $r$, where $r$ is a 0-based rank in the range $[0, n!)$.
  \item \inlinecode{rank\_by\_permutation(a)} returns an integer representing the 0-based rank of permutation \inlinecode{a}, which must be a permutation of the integers $[0, n)$.
  \item \inlinecode{permutation\_cycles(a)} returns the orbits of the index mapping \inlinecode{i} $\mapsto$ \inlinecode{a[i]}. For example, $\{3, 1, 0, 2\}$ decomposes into cycles $\{0, 3, 2\}$ and $\{1\}$ because following the mapping from index $0$ gives $0 \to 3 \to 2 \to 0$, while index $1$ maps to itself.
  \item \inlinecode{min\_swaps\_to\_sort(a)} returns the fewest swaps of any two elements that sort \inlinecode{a}, whose values must be distinct. This is unrelated to the inversion count of section \secref{1.1.5}, which measures the number of swaps when only adjacent elements may be exchanged.
\end{compactitem}

The two ranking operations require $n \leq 20$, so every possible rank fits in \inlinecode{int64\_t}.

\Needspace*{3\baselineskip}
\textbf{Time Complexity}:
\begin{compactitem}
  \item $O(n)$ per call to \inlinecode{next\_permutation2()} and \inlinecode{next\_permutation()}, where $n$ is the input size.
  \item $O(n^{2})$ per call to \inlinecode{permutation\_by\_rank()} and \inlinecode{rank\_by\_permutation()}.
  \item $O(1)$ per call to \inlinecode{next\_permutation\_mask()}.
  \item $O(n)$ per call to \inlinecode{permutation\_cycles()}.
  \item $O(n \log n)$ per call to \inlinecode{min\_swaps\_to\_sort()}, from ranking the values.
\end{compactitem}

\Needspace*{3\baselineskip}
\textbf{Space Complexity}:
\begin{compactitem}
  \item $O(1)$ auxiliary for \inlinecode{next\_permutation2()} and \inlinecode{next\_permutation()}.
  \item $O(n)$ auxiliary for \inlinecode{permutation\_by\_rank()}, \inlinecode{rank\_by\_permutation()}, \inlinecode{permutation\_cycles()}, and \inlinecode{min\_swaps\_to\_sort()}.
\end{compactitem}

\begin{lstlisting}[style=cpp]
#include <algorithm>
#include <cassert>
#include <cstdint>
#include <functional>
#include <numeric>
#include <utility>
#include <vector>

template<typename It, typename Compare = std::less<>>
bool next_permutation2(It lo, It hi, Compare comp = Compare{}) {
  if (lo == hi) {
    return false;
  }
  It i = lo;
  if (++i == hi) {
    return false;
  }
  i = hi;
  --i;
  while (true) {
    It j = i;
    if (comp(*--i, *j)) {
      It k = hi;
      while (!comp(*i, *--k)) {
      }
      std::iter_swap(i, k);
      std::reverse(j, hi);
      return true;
    }
    if (i == lo) {
      std::reverse(lo, hi);
      return false;
    }
  }
}

template<typename T, typename Compare = std::less<>>
bool next_permutation(std::vector<T> &a, Compare comp = Compare{}) {
  int n = static_cast<int>(a.size());
  for (int i = n - 2; i >= 0; i--) {
    if (comp(a[i], a[i + 1])) {
      for (int j = n - 1;; j--) {
        if (comp(a[i], a[j])) {
          std::swap(a[i++], a[j]);
          for (j = n - 1; i < j; i++, j--) {
            std::swap(a[i], a[j]);
          }
          return true;
        }
      }
    }
  }
  std::reverse(a.begin(), a.end());
  return false;
}

uint64_t next_permutation_mask(uint64_t x) {
  if (x == 0) {
    return 0;
  }
  uint64_t s = x & -x, r = x + s;
  if (r == 0) {
    return 0;
  }
  return r | (((x ^ r) >> 2) / s);
}

std::vector<int> permutation_by_rank(int n, int64_t r) {
  assert(0 <= n && n <= 20 && r >= 0);
  std::vector<int64_t> factorial(n + 1, 1);
  std::vector<int> values(n), res(n);
  for (int i = 1; i <= n; i++) {
    factorial[i] = i * factorial[i - 1];
  }
  assert(r < factorial[n]);
  std::iota(values.begin(), values.end(), 0);
  for (int i = 0; i < n; i++) {
    int pos = r / factorial[n - 1 - i];
    res[i] = values[pos];
    values.erase(values.begin() + pos);
    r %= factorial[n - 1 - i];
  }
  return res;
}

int64_t rank_by_permutation(const std::vector<int> &a) {
  int n = static_cast<int>(a.size());
  assert(n <= 20);
  std::vector<char> seen(n);
  for (int x : a) {
    assert(0 <= x && x < n && !seen[x]);
    seen[x] = true;
  }
  if (n == 0) {
    return 0;
  }
  std::vector<int64_t> factorial(n);
  factorial[0] = 1;
  for (int i = 1; i < n; i++) {
    factorial[i] = i * factorial[i - 1];
  }
  int64_t res = 0;
  for (int i = 0; i < n; i++) {
    int v = a[i];
    for (int j = 0; j < i; j++) {
      if (a[j] < a[i]) {
        v--;
      }
    }
    res += v * factorial[n - 1 - i];
  }
  return res;
}

using Cycles = std::vector<std::vector<int>>;

Cycles permutation_cycles(const std::vector<int> &a) {
  int n = static_cast<int>(a.size());
  std::vector<char> visit(n);
  for (int x : a) {
    assert(0 <= x && x < n && !visit[x]);
    visit[x] = true;
  }
  std::fill(visit.begin(), visit.end(), false);
  Cycles res;
  for (int i = 0; i < n; i++) {
    if (!visit[i]) {
      int j = i;
      std::vector<int> curr;
      do {
        curr.push_back(j);
        visit[j] = true;
        j = a[j];
      } while (j != i);
      res.push_back(std::move(curr));
    }
  }
  return res;
}

int min_swaps_to_sort(const std::vector<int> &a) {
  int n = static_cast<int>(a.size());
  std::vector<int> order(n);
  std::iota(order.begin(), order.end(), 0);
  std::sort(order.begin(), order.end(), [&](int i, int j) { return a[i] < a[j]; });
  assert(std::adjacent_find(order.begin(), order.end(), [&](int i, int j) {
           return a[i] == a[j];
         }) == order.end());
  return n - static_cast<int>(permutation_cycles(order).size());
}
\end{lstlisting}
\Needspace*{4\baselineskip}
\textbf{Example Usage}\par
\begin{lstlisting}[style=cpp]
#include <bitset>
#include <cassert>
#include <iostream>
using namespace std;

template<typename It>
void print_range(It lo, It hi) {
  cout << "{";
  for (; lo != hi; ++lo) {
    cout << *lo << (lo == hi - 1 ? "" : ",");
  }
  cout << "} ";
}

int main() {
  {
    const int n = 4;
    vector<int> a{0, 1, 2, 3}, b = a, c = a;
    cout << "Permutations of [0, " << n << "):" << endl;
    int count = 0;
    do {
      print_range(a.begin(), a.end());
      assert(b == a);
      assert(c == a);
      assert(permutation_by_rank(n, count) == a);
      assert(rank_by_permutation(a) == count);
      count++;
      if (count == 8 || count == 16) {
        cout << endl;
      }
      std::next_permutation(b.begin(), b.end());
      next_permutation(c);
    } while (next_permutation(a));
    assert(count == 24);
    assert((a == vector<int>{0, 1, 2, 3}));
    cout << endl;
  }
  {  // Permutations of binary digits.
    const int n = 5;
    cout << "\nPermutations of 2 zeros and 3 ones:" << endl;
    uint64_t lo = bitset<5>(string("00111")).to_ullong();
    uint64_t hi = bitset<6>(string("100011")).to_ullong();
    int count = 0;
    do {
      cout << bitset<n>(lo).to_string() << " ";
      count++;
    } while ((lo = next_permutation_mask(lo)) != hi);
    assert(count == 10);
    assert(next_permutation_mask(1ULL << 63) == 0);
    cout << endl;
  }
  {
    vector<int> a{3, 2, 1}, b = a;
    assert(next_permutation(a, greater<int>()));
    assert(next_permutation2(b.begin(), b.end(), greater<int>()));
    assert((a == vector<int>{3, 1, 2} && b == a));
  }
  {  // Decomposition into cycles.
    vector<int> a{3, 1, 0, 2};
    cout << "\nDecomposition of {3,1,0,2} into cycles:" << endl;
    Cycles c = permutation_cycles(a);
    assert((c == Cycles{{0, 3, 2}, {1}}));
    for (const auto &cycle : c) {
      print_range(cycle.begin(), cycle.end());
    }
    cout << endl;
    // Two cycles over four elements, so two swaps sort it.
    assert(min_swaps_to_sort(a) == 2);
    assert(min_swaps_to_sort({50, 20, 40, 10}) == 1);  // Values need not be a permutation.
    assert(min_swaps_to_sort({1, 2, 3}) == 0);
    assert(min_swaps_to_sort({}) == 0);
  }
  return 0;
}
\end{lstlisting}
\begin{exampleoutput}
Permutations of [0, 4):
{0,1,2,3} {0,1,3,2} {0,2,1,3} {0,2,3,1} {0,3,1,2} {0,3,2,1} {1,0,2,3} {1,0,3,2}
{1,2,0,3} {1,2,3,0} {1,3,0,2} {1,3,2,0} {2,0,1,3} {2,0,3,1} {2,1,0,3} {2,1,3,0}
{2,3,0,1} {2,3,1,0} {3,0,1,2} {3,0,2,1} {3,1,0,2} {3,1,2,0} {3,2,0,1} {3,2,1,0}

Permutations of 2 zeros and 3 ones:
00111 01011 01101 01110 10011 10101 10110 11001 11010 11100

Decomposition of {3,1,0,2} into cycles:
{0,3,2} {1}
\end{exampleoutput}
\subsection{Enumerating Combinations}
\label{sec:6.2.4}
A combination is a subset of size $k$ chosen from a total of $n$ (not necessarily distinct) elements, where order does not matter.

The lexicographic successor advances the rightmost chosen position that can still move right, then packs every later position as far left as possible. Ranking and unranking use the combinatorial number system: at each position, count how many combinations would be skipped by choosing a smaller next value, then either add that count to the rank or subtract it while searching for the requested rank. Bitmask successors use the same order as increasing integers with a fixed popcount.

\begin{compactitem}
  \item \inlinecode{next\_combination(lo, mid, hi, comp = std::less\textless{}\textgreater{}())} takes random-access iterators \inlinecode{lo}, \inlinecode{mid}, and \inlinecode{hi} as a range $[{\inlinecodemath{lo}}, {\inlinecodemath{hi}})$ of $n$ elements for which the function will rearrange such that the $k$ elements in $[{\inlinecodemath{lo}}, {\inlinecodemath{mid}})$ become the next lexicographically greater combination. The function returns true if such a combination exists, or false if $[{\inlinecodemath{lo}}, {\inlinecodemath{mid}})$ already consists of the lexicographically greatest combination of the elements in $[{\inlinecodemath{lo}}, {\inlinecodemath{hi}})$, in which case the range is reset to its first combination. The comparator \inlinecode{comp} defines the element ordering.
  \item \inlinecode{next\_combination(n, a)} rearranges \inlinecode{a} to become the next lexicographically greater combination of distinct integers in the range $[0, {\inlinecodemath{n}})$. The vector \inlinecode{a} must be sorted and contain distinct integers in the range $[0, {\inlinecodemath{n}})$.
  \item \inlinecode{next\_combination\_mask(x)} interprets the bits of an integer \inlinecode{x} as a mask with 1-bits specifying the chosen items for a combination and returns the mask of the next lexicographically greater combination (that is, the lowest integer greater than \inlinecode{x} with the same number of 1 bits). Note that this does not generate combinations in the same order as \inlinecode{next\_combination()}, nor does it work if the corresponding $n$ items are not distinct (in that case, duplicate combinations will be generated). It returns $0$ if no successor fits in \inlinecode{uint64\_t}.
  \item \inlinecode{combination\_by\_rank(n, k, r)} returns the combination of $k$ distinct integers in the range $[0, {\inlinecodemath{n}})$ that is lexicographically ranked $r$, where $r$ is a 0-based rank in the range $[0, \binom{n}{k})$.
  \item \inlinecode{rank\_by\_combination(n, a)} returns an integer representing the 0-based rank of combination \inlinecode{a}, which must contain sorted distinct integers in $[0, {\inlinecodemath{n}})$.
  \item \inlinecode{next\_combination\_with\_repeats(n, a)} rearranges \inlinecode{a} to become the next lexicographically greater combination of not necessarily distinct integers in the range $[0, {\inlinecodemath{n}})$. The vector \inlinecode{a} must be sorted. Note that there is a total of $n \mathbin{\text{multichoose}} k$ combinations if repetition is allowed, where $n \mathbin{\text{multichoose}} k = \binom{n + k - 1}{k}$.
\end{compactitem}

Overflow warning: Exact counts and ranks used by the ranking operations must fit in \inlinecode{int64\_t}.

\Needspace*{3\baselineskip}
\textbf{Time Complexity}:
\begin{compactitem}
  \item $O(n)$ per call to \inlinecode{next\_combination(lo, mid, hi)}, where $n$ is the distance between \inlinecode{lo} and \inlinecode{hi}.
  \item $O(k)$ per call to \inlinecode{next\_combination(n, a)} and \inlinecode{next\_combination\_with\_repeats(n, a)}, where $k$ is the size of \inlinecode{a}.
  \item $O(1)$ per call to \inlinecode{next\_combination\_mask()}.
  \item $O(n \cdot k)$ per call to \inlinecode{combination\_by\_rank()} and \inlinecode{rank\_by\_combination()}.
\end{compactitem}

\Needspace*{3\baselineskip}
\textbf{Space Complexity}:
\begin{compactitem}
  \item $O(k)$ auxiliary for \inlinecode{combination\_by\_rank(n, k, r)}.
  \item $O(1)$ auxiliary for all other operations.
\end{compactitem}

\begin{lstlisting}[style=cpp]
#include <algorithm>
#include <cassert>
#include <cstdint>
#include <functional>
#include <iterator>
#include <numeric>
#include <utility>
#include <vector>

template<typename It, typename Compare = std::less<>>
bool next_combination(It lo, It mid, It hi, Compare comp = Compare{}) {
  using T = typename std::iterator_traits<It>::value_type;
  if (lo == mid || mid == hi) {
    return false;
  }
  It l = mid - 1, h = hi - 1;
  int len1 = 1, len2 = 1;
  while (l != lo && !comp(*l, *h)) {
    --l;
    ++len1;
  }
  if (l == lo && !comp(*l, *h)) {
    std::rotate(lo, mid, hi);
    return false;
  }
  for (; mid < h; ++len2) {
    if (!comp(*l, *--h)) {
      ++h;
      break;
    }
  }
  if (len1 == 1 || len2 == 1) {
    std::iter_swap(l, h);
  } else if (len1 == len2) {
    std::swap_ranges(l, mid, h);
  } else {
    std::iter_swap(l++, h++);
    int total = (--len1) + (--len2), gcd = total;
    for (int i = len1; i != 0;) {
      std::swap(gcd %= i, i);
    }
    int skip = total / gcd - 1;
    for (int i = 0; i < gcd; i++) {
      It curr = (i < len1) ? (l + i) : (h + (i - len1));
      int k = i;
      T prev = *curr;
      for (int j = 0; j < skip; j++) {
        k = (k + len1) % total;
        It next = (k < len1) ? (l + k) : (h + (k - len1));
        *curr = *next;
        curr = next;
      }
      *curr = prev;
    }
  }
  return true;
}

bool next_combination(int n, std::vector<int> &a) {
  int k = static_cast<int>(a.size());
  for (int i = k - 1; i >= 0; i--) {
    if (a[i] < n - k + i) {
      a[i]++;
      while (++i < k) {
        a[i] = a[i - 1] + 1;
      }
      return true;
    }
  }
  return false;
}

uint64_t next_combination_mask(uint64_t x) {
  if (x == 0) {
    return 0;
  }
  uint64_t s = x & -x, r = x + s;
  if (r == 0) {
    return 0;
  }
  return r | (((x ^ r) >> 2) / s);
}

int64_t n_choose_k(int64_t n, int64_t k) {
  assert(n >= 0 && 0 <= k && k <= n);
  if (k > n - k) {
    k = n - k;
  }
  int64_t res = 1;
  for (int i = 0; i < k; i++) {
    int64_t num = n - i, den = i + 1;
    int64_t g = std::gcd(num, den);
    num /= g;
    den /= g;
    res /= den;
    res *= num;  // Overflow warning.
  }
  return res;
}

std::vector<int> combination_by_rank(int n, int k, int64_t r) {
  assert(0 <= r && r < n_choose_k(n, k));
  std::vector<int> res(k);
  int count = n;
  for (int i = 0; i < k; i++) {
    int j = 1;
    for (;; j++) {
      int64_t am = n_choose_k(count - j, k - 1 - i);
      if (r < am) {
        break;
      }
      r -= am;
    }
    res[i] = (i > 0) ? (res[i - 1] + j) : (j - 1);
    count -= j;
  }
  return res;
}

int64_t rank_by_combination(int n, const std::vector<int> &a) {
  int k = static_cast<int>(a.size());
  assert(n >= 0 && k <= n);
  assert(std::all_of(a.begin(), a.end(), [n](int x) { return 0 <= x && x < n; }));
  assert(std::adjacent_find(a.begin(), a.end(), std::greater_equal<int>()) == a.end());
  int64_t res = 0;
  int prev = -1;
  for (int i = 0; i < k; i++) {
    for (int j = prev + 1; j < a[i]; j++) {
      res += n_choose_k(n - 1 - j, k - 1 - i);
    }
    prev = a[i];
  }
  return res;
}

bool next_combination_with_repeats(int n, std::vector<int> &a) {
  int k = static_cast<int>(a.size());
  for (int i = k - 1; i >= 0; i--) {
    if (a[i] < n - 1) {
      for (++a[i]; ++i < k;) {
        a[i] = a[i - 1];
      }
      return true;
    }
  }
  return false;
}
\end{lstlisting}
\Needspace*{4\baselineskip}
\textbf{Example Usage}\par
\begin{lstlisting}[style=cpp]
#include <cassert>
#include <iostream>
using namespace std;

template<typename It>
void print_range(It lo, It hi) {
  cout << "{";
  for (; lo != hi; ++lo) {
    cout << *lo << (lo == hi - 1 ? "" : ",");
  }
  cout << "} ";
}

int main() {
  {
    int k = 3;
    string s = "11234";
    cout << "\"" << s << "\" choose " << k << ":" << endl;
    int count = 0;
    do {
      cout << s.substr(0, k) << " ";
      count++;
    } while (next_combination(s.begin(), s.begin() + k, s.end()));
    assert(count == 7);  // Duplicate '1' values collapse equal output strings.
    cout << endl;
  }
  {  // Unordered combinations using masks.
    int n = 5, k = 3;
    string char_set = "abcde";  // Must be distinct.
    cout << "\n\"" << char_set << "\" choose " << k << " with masks:" << endl;
    uint64_t mask = (1ULL << k) - 1, limit = 1ULL << n;
    int count = 0;
    do {
      for (int i = 0; i < n; i++) {
        if ((mask >> i) & 1) {
          cout << char_set[i];
        }
      }
      cout << " ";
      count++;
      mask = next_combination_mask(mask);
    } while (mask < limit);
    assert(count == 10);
    assert(next_combination_mask(1ULL << 63) == 0);
    cout << endl;
  }
  {
    string s = "4321";
    assert(next_combination(s.begin(), s.begin() + 2, s.end(), greater<char>()));
    assert(s.substr(0, 2) == "42");
  }
  {  // Combinations of distinct integers in [0, n).
    int n = 5, k = 3;
    vector<int> a{0, 1, 2};
    cout << endl << n << " choose " << k << ":" << endl;
    int count = 0;
    do {
      print_range(a.begin(), a.end());
      assert(a == combination_by_rank(n, k, count));
      assert(rank_by_combination(n, a) == count);
      count++;
    } while (next_combination(n, a));
    assert(count == 10);
    assert(n_choose_k(66, 33) == 7219428434016265740LL);
    cout << endl;
  }
  {  // Combinations with repeats.
    int n = 3, k = 2;
    vector<int> a{0, 0};
    cout << endl << n << " multichoose " << k << ":" << endl;
    int count = 0;
    do {
      print_range(a.begin(), a.end());
      count++;
    } while (next_combination_with_repeats(n, a));
    assert(count == 6);
    cout << endl;
  }
  return 0;
}
\end{lstlisting}
\begin{exampleoutput}
"11234" choose 3:
112 113 114 123 124 134 234

"abcde" choose 3 with masks:
abc abd acd bcd abe ace bce ade bde cde

5 choose 3:
{0,1,2} {0,1,3} {0,1,4} {0,2,3} {0,2,4} {0,3,4} {1,2,3} {1,2,4} {1,3,4} {2,3,4}

3 multichoose 2:
{0,0} {0,1} {0,2} {1,1} {1,2} {2,2}
\end{exampleoutput}
\subsection{Enumerating Partitions}
\label{sec:6.2.5}
A partition of a natural number $n$ is a way to write $n$ as a sum of positive integers where the order of the addends does not matter.

The lexicographic successor edits the suffix of a non-increasing partition: increase the rightmost part that can grow without breaking the order, then refill the remaining sum with ones. Ranking and unranking use a dynamic-programming table that counts partitions of a remaining sum by their next part. This lets the algorithms skip whole blocks of partitions sharing a common prefix.

\begin{compactitem}
  \item \inlinecode{next\_partition(p)} takes a reference to a vector \inlinecode{p} of positive integers as a partition of $n$ for which the function will re-assign to become the next lexicographically greater partition. The function returns true if such a partition exists, or false if \inlinecode{p} already consists of the lexicographically greatest partition (i.e. the single integer $n$).
  \item \inlinecode{partition\_by\_rank(n, r)} returns the partition of $n$ that is lexicographically ranked $r$ if addends in each partition were sorted in non-increasing order, where $r$ is a 0-based rank in the range $[0, p(n))$ where $p(n)$ is the number of partitions of $n$.
  \item \inlinecode{rank\_by\_partition(p)} returns an integer representing the 0-based rank of the partition specified by vector \inlinecode{p}, which must consist of positive integers sorted in non-increasing order.
  \item \inlinecode{generate\_increasing\_partitions(n, f)} calls the function \inlinecode{f(lo, hi)} on strictly increasing partitions of $n$ in lexicographically increasing order of partition, where \inlinecode{lo} and \inlinecode{hi} are random-access iterators to a range $[{\inlinecodemath{lo}}, {\inlinecodemath{hi}})$ of integers. Note that non-strictly increasing partitions like $\{1, 1, 1, 1\}$ are skipped.
\end{compactitem}

Overflow warning: Exact partition counts and ranks must fit in \inlinecode{int64\_t}.

\Needspace*{3\baselineskip}
\textbf{Time Complexity}:
\begin{compactitem}
  \item $O(n)$ per call to \inlinecode{next\_partition()}.
  \item $O(n^{2})$ per call to \inlinecode{partition\_by\_rank()} and \inlinecode{rank\_by\_partition()}.
  \item $O(p(n))$ per call to \inlinecode{generate\_increasing\_partitions()}, where $p(n)$ is the number of partitions of $n$.
\end{compactitem}

\Needspace*{3\baselineskip}
\textbf{Space Complexity}:
\begin{compactitem}
  \item $O(1)$ auxiliary for \inlinecode{next\_partition()}.
  \item $O(n^{2})$ auxiliary heap space for \inlinecode{partition\_function()}, \inlinecode{partition\_by\_rank()}, and \inlinecode{rank\_by\_partition()}.
  \item $O(n)$ auxiliary stack space for \inlinecode{generate\_increasing\_partitions()}.
\end{compactitem}

\begin{lstlisting}[style=cpp]
#include <algorithm>
#include <cassert>
#include <cstdint>
#include <functional>
#include <numeric>
#include <vector>

bool next_partition(std::vector<int> &p) {
  int n = static_cast<int>(p.size());
  if (n <= 1) {
    return false;
  }
  int s = p[n - 1] - 1, i = n - 2;
  p.pop_back();
  for (; i > 0 && p[i] == p[i - 1]; i--) {
    s += p[i];
    p.pop_back();
  }
  for (p[i]++; s > 0; s--) {
    p.push_back(1);
  }
  return true;
}

int64_t partition_function(int a, int b) {
  assert(a >= 0 && 0 <= b && b <= a);
  static std::vector<std::vector<int64_t>> p(1, std::vector<int64_t>(1, 1));
  if (a >= static_cast<int>(p.size())) {
    int old = static_cast<int>(p.size());
    p.resize(a + 1);
    for (int i = 0; i < old; i++) {
      p[i].resize(a + 1);
    }
    for (int i = old; i <= a; i++) {
      p[i].resize(a + 1);
      for (int j = 1; j <= i; j++) {
        p[i][j] = p[i - 1][j - 1] + p[i - j][j];
      }
    }
  }
  return p[a][b];
}

std::vector<int> partition_by_rank(int n, int64_t r) {
  assert(n >= 0 && r >= 0);
  std::vector<int> res;
  for (int i = n, j; i > 0; i -= j) {
    for (j = 1;; j++) {
      int64_t count = partition_function(i, j);
      if (r < count) {
        break;
      }
      r -= count;
    }
    res.push_back(j);
  }
  return res;
}

int64_t rank_by_partition(const std::vector<int> &p) {
  assert(std::all_of(p.begin(), p.end(), [](int x) { return x > 0; }));
  assert(std::is_sorted(p.begin(), p.end(), std::greater<int>()));
  int64_t res = 0;
  int sum = std::accumulate(p.begin(), p.end(), 0);
  for (int x : p) {
    for (int j = 0; j < x; j++) {
      res += partition_function(sum, j);
    }
    sum -= x;
  }
  return res;
}

template<typename Fn>
void generate_increasing_partitions(int left, int prev, int i, std::vector<int> &p, Fn f) {
  if (left == 0) {
    f(p.begin(), p.begin() + i);
    return;
  }
  for (p[i] = prev + 1; p[i] <= left; p[i]++) {
    generate_increasing_partitions(left - p[i], p[i], i + 1, p, f);
  }
}

template<typename Fn>
void generate_increasing_partitions(int n, Fn f) {
  std::vector<int> p(n);
  generate_increasing_partitions(n, 0, 0, p, f);
}
\end{lstlisting}
\Needspace*{4\baselineskip}
\textbf{Example Usage}\par
\begin{lstlisting}[style=cpp]
#include <cassert>
#include <iostream>
using namespace std;

template<typename It>
void print_range(It lo, It hi) {
  cout << "{";
  for (; lo != hi; ++lo) {
    cout << *lo << (lo == hi - 1 ? "" : ",");
  }
  cout << "} ";
}

int main() {
  // Exercise incremental growth of the cached dynamic-programming table.
  assert(partition_function(3, 1) == 1);
  assert(partition_function(8, 3) == 5);
  {
    int n = 4;
    vector<int> a(n, 1);
    cout << "Partitions of " << n << ":" << endl;
    int count = 0;
    do {
      print_range(a.begin(), a.end());
      assert(a == partition_by_rank(n, count));
      assert(rank_by_partition(a) == count);
      count++;
    } while (next_partition(a));
    assert(count == 5);
    cout << endl;
  }
  {
    int n = 8;
    cout << "\nIncreasing partitions of " << n << ":" << endl;
    int count = 0;
    generate_increasing_partitions(n, [&](auto lo, auto hi) {
      print_range(lo, hi);
      count++;
    });
    assert(count == 6);
    cout << endl;
  }
  return 0;
}
\end{lstlisting}
\begin{exampleoutput}
Partitions of 4:
{1,1,1,1} {2,1,1} {2,2} {3,1} {4}

Increasing partitions of 8:
{1,2,5} {1,3,4} {1,7} {2,6} {3,5} {8}
\end{exampleoutput}
\subsection{Generic Combinatorial Enumeration}
\label{sec:6.2.6}
Ranks, unranks, and enumerates combinatorial sequences in lexicographic order using only a prefix counting oracle. This is useful when the objects are too numerous to precompute, but there is a simple dynamic-programming formula for the number of valid completions after a fixed prefix.

\inlinecode{Enumerator(range, length, count)} considers values in $[0, {\inlinecodemath{range}})$ at each position of a length-\inlinecode{length} sequence. The callback \inlinecode{count(prefix)} returns the number of valid full sequences beginning with \inlinecode{prefix}; it may capture a precomputed dynamic-programming table or other state. The factory functions below supply this oracle for the principal combinatorial families, while the class remains available directly for custom families. The range and length must be nonnegative; the factory arguments must satisfy $0 \leq k \leq n$.

\begin{compactitem}
  \item \inlinecode{arrangements(n, k)} returns an enumerator for the length-$k$ arrangements of $[0, n)$.
  \item \inlinecode{permutations(n)} returns an enumerator for the permutations of $[0, n)$.
  \item \inlinecode{combinations(n, k)} returns an enumerator for the size-$k$ subsets of $[0, n)$, represented as strictly increasing sequences.
  \item \inlinecode{partitions(n)} returns an enumerator for the additive partitions of $n$, represented as non-increasing length-$n$ sequences padded with zeros.
  \item \inlinecode{to\_rank(a)} returns the 0-based rank of the valid combinatorial sequence \inlinecode{a}.
  \item \inlinecode{from\_rank(r)} returns a combinatorial sequence of integers that is lexicographically ranked \inlinecode{r}, where \inlinecode{r} is a 0-based rank in the range $[0, {\inlinecodemath{total\_count()}})$.
  \item \inlinecode{total\_count()} returns the number of valid sequences.
  \item \inlinecode{enumerate(f)} calls the function \inlinecode{f(a)} on every specified combinatorial sequence \inlinecode{a} in lexicographically increasing order.
\end{compactitem}

Overflow warning: All exact counts and ranks must fit in \inlinecode{int64\_t}.

\Needspace*{3\baselineskip}
\textbf{Time Complexity}:
\begin{compactitem}
  \item $O(A \cdot L \cdot C)$ per call to \inlinecode{to\_rank()} and \inlinecode{from\_rank()}, where $A$ is \inlinecode{range}, $L$ is \inlinecode{length}, and $C$ is the cost of \inlinecode{count()}.
  \item $O(T \cdot A \cdot L \cdot C)$ per call to \inlinecode{enumerate()}, excluding the cost of the callback \inlinecode{f}, where $T$ is \inlinecode{total\_count()}.
  \item $O(A \cdot (\min\left(L, A - L\right) + 1))$ per call to \inlinecode{combinations()} and $O(L^{2})$ per call to \inlinecode{partitions()}; the other factories take $O(1)$ time per call.
\end{compactitem}

\Needspace*{3\baselineskip}
\textbf{Space Complexity}:
\begin{compactitem}
  \item $O(L)$ auxiliary for \inlinecode{to\_rank()}, \inlinecode{from\_rank()}, and \inlinecode{enumerate()}, excluding output storage.
  \item $O(S)$ object storage, where $S$ is the state captured by \inlinecode{count()}.
  \item $O(A \cdot (\min\left(L, A - L\right) + 1))$ object storage for \inlinecode{combinations()} and $O(L^{2})$ for \inlinecode{partitions()}.
\end{compactitem}

\begin{lstlisting}[style=cpp]
#include <algorithm>
#include <cassert>
#include <cstdint>
#include <numeric>
#include <utility>
#include <vector>

template<typename Count>
class Enumerator {
  int range, length;
  Count count;

 public:
  Enumerator(int range, int length, Count count)
      : range(range), length(length), count(std::move(count)) {
    assert(range >= 0 && length >= 0);
  }

  int64_t total_count() { return count({}); }

  int64_t to_rank(const std::vector<int> &a) {
    assert(static_cast<int>(a.size()) == length);
    assert(std::all_of(a.begin(), a.end(), [&](int x) { return 0 <= x && x < range; }));
    int64_t res = 0;
    std::vector<int> prefix;
    for (int i = 0; i < static_cast<int>(a.size()); i++) {
      prefix.push_back(0);
      for (; prefix.back() < a[i]; prefix.back()++) {
        res += count(prefix);
      }
    }
    return res;
  }

  std::vector<int> from_rank(int64_t r) {
    assert(0 <= r && r < total_count());
    std::vector<int> a;
    for (int i = 0; i < length; i++) {
      a.push_back(0);
      for (; a.back() < range; a.back()++) {
        int64_t curr = count(a);
        if (r < curr) {
          break;
        }
        r -= curr;
      }
    }
    return a;
  }

  // Accepts any callable f(a), including capturing lambdas and functors.
  template<typename Fn>
  void enumerate(Fn f) {
    int64_t total = total_count();
    for (int64_t i = 0; i < total; i++) {
      std::vector<int> curr = from_rank(i);
      f(curr);
    }
  }
};

auto arrangements(int n, int k) {
  assert(n >= 0 && 0 <= k && k <= n);
  auto count = [=](const std::vector<int> &prefix) -> int64_t {
    int len = static_cast<int>(prefix.size());
    for (int i = 0; i < len - 1; i++) {
      if (prefix[i] == prefix[len - 1]) {
        return 0;
      }
    }
    int64_t res = 1;
    for (int i = 0; i < k - len; i++) {
      res *= n - len - i;  // Overflow warning.
    }
    return res;
  };
  return Enumerator(n, k, count);
}

auto permutations(int n) {
  return arrangements(n, n);
}

auto combinations(int n, int k) {
  assert(n >= 0 && 0 <= k && k <= n);
  std::vector<std::vector<int64_t>> table(n + 1, std::vector<int64_t>(std::min(k, n - k) + 1));
  int max_col = static_cast<int>(table[0].size()) - 1;
  for (int i = 0; i <= n; i++) {
    for (int j = 0; j <= std::min(i, max_col); j++) {
      table[i][j] = (j == 0) ? 1 : table[i - 1][j - 1] + table[i - 1][j];  // Overflow warning.
    }
  }
  auto count = [n, k, table = std::move(table)](const std::vector<int> &pref) -> int64_t {
    int len = static_cast<int>(pref.size());
    if (len >= 2 && pref[len - 1] <= pref[len - 2]) {
      return 0;
    }
    int rem = len == 0 ? n : n - pref.back() - 1;
    int choose = k - len;
    return (choose < 0 || choose > rem) ? 0 : table[rem][std::min(choose, rem - choose)];
  };
  return Enumerator(n, k, std::move(count));
}

auto partitions(int n) {
  assert(n >= 0);
  std::vector<std::vector<int64_t>> table(n + 1, std::vector<int64_t>(n + 1));
  std::fill(table[0].begin(), table[0].end(), 1);
  for (int sum = 1; sum <= n; sum++) {
    for (int max_part = 1; max_part <= n; max_part++) {
      table[sum][max_part] = table[sum][max_part - 1];
      if (max_part <= sum) {
        table[sum][max_part] += table[sum - max_part][max_part];  // Overflow warning.
      }
    }
  }
  auto count = [n, table = std::move(table)](const std::vector<int> &pref) -> int64_t {
    int len = static_cast<int>(pref.size());
    int sum = std::accumulate(pref.begin(), pref.end(), 0);
    if (sum == n) {
      return 1;
    }
    if (sum > n || (len > 0 && pref.back() == 0) || (len >= 2 && pref[len - 1] > pref[len - 2])) {
      return 0;
    }
    return table[n - sum][len == 0 ? n : pref.back()];
  };
  return Enumerator(n + 1, n, std::move(count));
}
\end{lstlisting}
\Needspace*{4\baselineskip}
\textbf{Example Usage}\par
\begin{lstlisting}[style=cpp]
#include <iostream>
using namespace std;

void print_sequence(const vector<int> &a) {
  cout << "{";
  for (int i = 0; i < static_cast<int>(a.size()); i++) {
    cout << a[i] << (i + 1 == static_cast<int>(a.size()) ? "" : ",");
  }
  cout << "} ";
}

int main() {
  cout << "3 permute 2 arrangements:" << endl;
  auto arr = arrangements(3, 2);
  int count = 0;
  arr.enumerate([&](const auto &a) {
    print_sequence(a);
    count++;
  });
  assert(count == 6);
  assert((arr.from_rank(5) == vector<int>{2, 1}));
  assert(arr.to_rank(vector<int>{2, 1}) == 5);

  cout << "\n\nPermutations of [0, 3):" << endl;
  auto perm = permutations(3);
  count = 0;
  perm.enumerate([&](const auto &a) {
    print_sequence(a);
    count++;
  });
  assert(count == 6);

  cout << "\n\n4 choose 3 combinations:" << endl;
  auto comb = combinations(4, 3);
  count = 0;
  comb.enumerate([&](const auto &a) {
    print_sequence(a);
    count++;
  });
  assert(count == 4);
  assert((comb.from_rank(3) == vector<int>{1, 2, 3}));
  assert(comb.to_rank(vector<int>{1, 2, 3}) == 3);
  assert(combinations(67, 1).total_count() == 67);
  assert(combinations(67, 67).total_count() == 1);
  assert(combinations(66, 33).total_count() == 7219428434016265740LL);

  cout << "\n\nPartitions of 4:" << endl;
  auto part = partitions(4);
  count = 0;
  part.enumerate([&](const auto &a) {
    print_sequence(a);
    count++;
  });
  assert(count == 5);

  int length = 4;
  auto count_completions = [length](const vector<int> &pref) -> int64_t {
    if (static_cast<int>(pref.size()) > length) {
      return 0;
    }
    for (int i = 0; i < static_cast<int>(pref.size()); i++) {
      if (pref[i] < 0 || pref[i] > 1 || (i > 0 && pref[i - 1] == 1 && pref[i] == 1)) {
        return 0;
      }
    }
    int64_t after_zero = 1, after_one = 1;
    for (int i = static_cast<int>(pref.size()); i < length; i++) {
      int64_t next_zero = after_zero + after_one;  // Overflow warning.
      after_one = after_zero;
      after_zero = next_zero;
    }
    return pref.empty() || pref.back() == 0 ? after_zero : after_one;
  };
  Enumerator e(2, length, count_completions);
  cout << "\n\nLength-4 binary strings without consecutive ones:" << endl;
  count = 0;
  e.enumerate([&](const auto &a) {
    print_sequence(a);
    count++;
  });
  assert(count == 8 && e.total_count() == 8);
  assert((e.from_rank(7) == vector<int>{1, 0, 1, 0}));
  assert(e.to_rank(vector<int>{1, 0, 1, 0}) == 7);
  cout << endl;
  return 0;
}
\end{lstlisting}
\begin{exampleoutput}
3 permute 2 arrangements:
{0,1} {0,2} {1,0} {1,2} {2,0} {2,1}

Permutations of [0, 3):
{0,1,2} {0,2,1} {1,0,2} {1,2,0} {2,0,1} {2,1,0}

4 choose 3 combinations:
{0,1,2} {0,1,3} {0,2,3} {1,2,3}

Partitions of 4:
{1,1,1,1} {2,1,1,0} {2,2,0,0} {3,1,0,0} {4,0,0,0}

Length-4 binary strings without consecutive ones:
{0,0,0,0} {0,0,0,1} {0,0,1,0} {0,1,0,0} {0,1,0,1} {1,0,0,0} {1,0,0,1} {1,0,1,0}
\end{exampleoutput}

\section{\texorpdfstring{Number Theory}{6.3 Number Theory}}
\setcounter{section}{3}
\setcounter{subsection}{0}
\subsection{Mod Helpers and Binary Exponentiation}
\label{sec:6.3.1}
Modular arithmetic replaces every value by its remainder modulo \inlinecode{m}, which keeps intermediate results bounded and is what makes counting problems with astronomically large answers tractable. The operations are straightforward except for two hazards, both handled below: a sum or product can overflow the underlying integer type before the remainder is taken, and a negative operand leaves a negative remainder under C++ truncating division. Binary exponentiation squares the base while walking the bits of the exponent, so a power costs $O(\log n)$ multiplications rather than $n$; the same loop computes powers in any monoid, which is how the matrix exponentiation of section \secref{6.5.1} and the linear recurrence of section \secref{6.6.6} are evaluated.

\begin{compactitem}
  \item \inlinecode{addmod(a, b, m)} and \inlinecode{submod(a, b, m)} respectively return addition and subtraction modulo \inlinecode{m}, each result in $[0, {\inlinecodemath{m}})$. Both operands may be negative or unreduced; they are normalized before arithmetic to avoid signed overflow. The modulus \inlinecode{m} must be positive.
  \item \inlinecode{mulmod(a, b, m)} returns \inlinecode{a} multiplied by \inlinecode{b}, modulo \inlinecode{m}. This is done in a way to avoid overflow: on compilers with \inlinecode{\_\_uint128\_t} it uses one wide product, while the portable fallback uses double-and-add multiplication. The fallback is slower by a logarithmic factor, but avoids relying on nonstandard 128-bit integers. Unlike \inlinecode{addmod}/\inlinecode{submod}, this takes unsigned operands and supports a full 64-bit modulus.
  \item \inlinecode{powmod(x, n, m)} returns \inlinecode{x} raised to the power \inlinecode{n}, modulo \inlinecode{m}.
\end{compactitem}

Paste these when a solution needs a handful of modular operations; when modular arithmetic pervades one instead, the \inlinecode{Modular} value type of section \secref{6.3.3} overloads the operators so expressions read normally. Only \inlinecode{mulmod()} has no counterpart there, since a value type multiplying in \inlinecode{int64\_t} overflows once the modulus passes about three billion, which is why it underlies the Miller-Rabin test and Pollard's rho of section \secref{6.3.9}.

\Needspace*{3\baselineskip}
\textbf{Time Complexity}:
\begin{compactitem}
  \item $O(1)$ per call to \inlinecode{addmod()} and \inlinecode{submod()}.
  \item $O(1)$ per call to \inlinecode{mulmod()} with \inlinecode{\_\_uint128\_t}, or $O(\log b)$ with the portable fallback, where $b$ is the second argument.
  \item $O(\log n)$ calls to \inlinecode{mulmod()} per call to \inlinecode{powmod(x, n, m)}.
\end{compactitem}

\textbf{Space Complexity}: $O(1)$ auxiliary for all operations.

\begin{lstlisting}[style=cpp]
#include <cassert>
#include <cstdint>

int64_t addmod(int64_t a, int64_t b, int64_t m) {
  assert(m > 0);
  if ((a %= m) < 0) a += m;
  if ((b %= m) < 0) b += m;
  return a >= m - b ? a - (m - b) : a + b;
}

int64_t submod(int64_t a, int64_t b, int64_t m) {
  assert(m > 0);
  if ((a %= m) < 0) a += m;
  if ((b %= m) < 0) b += m;
  return a >= b ? a - b : m - (b - a);
}

uint64_t mulmod(uint64_t a, uint64_t b, uint64_t m) {
  assert(m > 0);
#if defined(__SIZEOF_INT128__)
  return static_cast<uint64_t>(static_cast<__uint128_t>(a) * b % m);
#else
  uint64_t res = 0, cur = a % m;
  for (; b > 0; b >>= 1) {
    if (b & 1) {
      res = res >= m - cur ? res - (m - cur) : res + cur;
    }
    cur = cur >= m - cur ? cur - (m - cur) : cur + cur;
  }
  return res;
#endif
}

uint64_t powmod(uint64_t x, uint64_t n, uint64_t m) {
  assert(m > 0);
  uint64_t res = 1 % m;
  for (; n > 0; n >>= 1) {
    if (n & 1) {
      res = mulmod(res, x, m);
    }
    x = mulmod(x, x, m);
  }
  return res;
}
\end{lstlisting}
\Needspace*{4\baselineskip}
\textbf{Example Usage}\par
\begin{lstlisting}[style=cpp]
#include <cassert>

int main() {
  assert(addmod(7, 8, 10) == 5 && submod(2, 5, 10) == 7);
  // Negative and unreduced operands are normalized into [0, m).
  assert(addmod(-3, -4, 10) == 3 && submod(-3, 4, 10) == 3);
  assert(addmod(25, -7, 10) == 8);
  assert(addmod(INT64_MAX - 1, INT64_MAX - 1, INT64_MAX) == INT64_MAX - 2);

  assert(mulmod(INT64_MAX - 1, INT64_MAX - 1, INT64_MAX) == 1);
  assert(powmod(2, 10, 1000000007) == 1024);
  assert(powmod(2, 62, 1000000) == 387904);
  assert(powmod(10001, 10001, 100000) == 10001);
  return 0;
}
\end{lstlisting}
\subsection{GCD, LCM, Mod Inverse, Chinese Remainder}
\label{sec:6.3.2}
Core integer and modular-arithmetic operations built around the Euclidean algorithm. Extended GCD produces Bezout coefficients, which solve linear Diophantine equations and modular inverses, while the Chinese remainder routines combine compatible congruences into one residue class.

\begin{compactitem}
  \item \inlinecode{gcd(a, b)} returns the greatest common divisor of \inlinecode{a} and \inlinecode{b} using the Euclidean algorithm. This is mainly for educational purposes, as \inlinecode{std::gcd(a, b)} from \inlinecode{\textless{}numeric\textgreater{}} is available as of C++17 (\inlinecode{\_\_gcd(a, b)} from \inlinecode{\textless{}algorithm\textgreater{}} in C++14 and earlier).
  \item \inlinecode{lcm(a, b)} returns the least common multiple of \inlinecode{a} and \inlinecode{b}. This implementation is mainly for educational purposes, as \inlinecode{std::lcm(a, b)} from \inlinecode{\textless{}numeric\textgreater{}} is available as of C++17.
  \item \inlinecode{extended\_gcd(a, b)} returns a pair $(x, y)$ of integers such that $\gcd(a, b) = ax + by$.
  \item \inlinecode{diophantine(a, b, c, \&g, \&x, \&y)} solves the linear Diophantine equation $ax + by = c$, returning whether a solution exists (one does if and only if $\gcd(a, b)$ divides $c$). \inlinecode{g} is always set to $\gcd(a, b)$, while (\inlinecode{x}, \inlinecode{y}) is set only on success to a particular solution bounded by $\max(|a|, |b|, |c|)$ in magnitude. Every other solution is $({\inlinecodemath{x}} + t({\inlinecodemath{b}}/{\inlinecodemath{g}}), {\inlinecodemath{y}} - t({\inlinecodemath{a}}/{\inlinecodemath{g}}))$ for an integer $t$. A 128-bit intermediate is used where available to keep the scaling step overflow-free.
  \item \inlinecode{mod(a, m)} returns the least nonnegative residue of \inlinecode{a} modulo \inlinecode{m}, that is, the unique value in $[0, {\inlinecodemath{m}})$ congruent to \inlinecode{a}, where \inlinecode{m} must be positive. Unlike the C++ remainder operator \inlinecode{\%}, whose result follows the sign of \inlinecode{a}, the result is never negative.
  \item \inlinecode{mod\_inverse(a, m)} returns an integer $x$ such that $ax \equiv 1 \pmod m$, where the arguments must satisfy $m > 0$ and $\gcd(a, m) = 1$.
  \item \inlinecode{mod\_inverse\_table(p)} returns a vector \inlinecode{v} where \inlinecode{v[i]} is the modular inverse of \inlinecode{i} modulo the prime \inlinecode{p} for every valid index \inlinecode{i}.
  \item \inlinecode{crt(r1, m1, r2, m2, \&r, \&m)} merges the two congruences $x \equiv r_1 \pmod{m_1}$ and $x \equiv r_2 \pmod{m_2}$ for arbitrary moduli (not necessarily coprime). It returns whether the system is consistent, and on success sets \inlinecode{m} to \inlinecode{lcm(m\_1, m\_2)} and \inlinecode{r} to the unique solution in $[0, {\inlinecodemath{m}})$. Both moduli must be positive, and their least common multiple must fit in \inlinecode{int64\_t}. Fold it pairwise to merge more than two congruences.
  \item \inlinecode{garner\_restore(a, p)} returns the smallest nonnegative solution whose residue modulo \inlinecode{p[i]} is \inlinecode{a[i]} at every valid index. The entries of \inlinecode{p} must be pairwise coprime (unlike \inlinecode{crt}, which allows shared factors) and positive, and \inlinecode{a} and \inlinecode{p} must have equal sizes. The exact solution is unique modulo the product of all moduli, so that product and the final answer must fit in \inlinecode{int64\_t}.
  \item \inlinecode{garner\_restore\_mod(a, p, m)} returns that same CRT solution modulo \inlinecode{m}. This is the right variant when the product of the moduli is too large to fit in \inlinecode{int64\_t}; \inlinecode{m} must be positive.
\end{compactitem}

Overflow warning: Every intermediate and result of the templated signed-integer helpers must be representable by \inlinecode{Int}; in particular, the minimum value of \inlinecode{Int} cannot be negated or divided by $-1$.

\Needspace*{3\baselineskip}
\textbf{Time Complexity}:
\begin{compactitem}
  \item $O(\log M)$ per call to \inlinecode{gcd()}, \inlinecode{lcm()}, \inlinecode{extended\_gcd()}, \inlinecode{diophantine()}, \inlinecode{mod\_inverse()}, and \inlinecode{crt()}, where $M$ is the largest relevant input magnitude.
  \item $O(1)$ per call to \inlinecode{mod()}.
  \item $O(p)$ per call to \inlinecode{mod\_inverse\_table()}.
  \item $O(n^{2})$ per call to \inlinecode{garner\_restore()} and \inlinecode{garner\_restore\_mod()}.
\end{compactitem}

\Needspace*{3\baselineskip}
\textbf{Space Complexity}:
\begin{compactitem}
  \item $O(p)$ auxiliary for \inlinecode{mod\_inverse\_table()}.
  \item $O(n)$ auxiliary for \inlinecode{garner\_restore()} and \inlinecode{garner\_restore\_mod()}.
  \item $O(1)$ auxiliary for all other operations.
\end{compactitem}

\begin{lstlisting}[style=cpp]
#include <algorithm>
#include <cassert>
#include <cstdint>
#include <type_traits>
#include <utility>
#include <vector>

template<typename Int>
Int gcd(Int a, Int b) {
  while (b != 0) {
    Int t = b;
    b = a % b;
    a = t;
  }
  return (a < 0 ? -a : a);
}

template<typename Int>
Int lcm(Int a, Int b) {
  if (a == 0 || b == 0) {
    return 0;
  }
  Int res = a / gcd(a, b) * b;
  return (res < 0 ? -res : res);
}

template<typename Int>
std::pair<Int, Int> extended_gcd(Int a, Int b) {
  Int x = 1, y = 0, x1 = 0, y1 = 1;
  while (b != 0) {
    Int q = a / b, prev_x1 = x1, prev_y1 = y1, prev_b = b;
    x1 = x - q * x1;
    y1 = y - q * y1;
    b = a - q * b;
    x = prev_x1;
    y = prev_y1;
    a = prev_b;
  }
  return (a > 0) ? std::pair<Int, Int>{x, y} : std::pair<Int, Int>{-x, -y};
}

template<typename Int>
bool diophantine(Int a, Int b, Int c, Int *g, Int *x, Int *y) {
  if (a == 0 && b == 0) {
    *g = *x = *y = 0;
    return c == 0;
  }
  if (a == 0) {
    *g = (b < 0) ? -b : b;
    if (c % b != 0) {
      return false;
    }
    *x = 0;
    *y = c / b;
    return true;
  }
  if (b == 0) {
    *g = (a < 0) ? -a : a;
    if (c % a != 0) {
      return false;
    }
    *x = c / a;
    *y = 0;
    return true;
  }
  *g = gcd(a, b);
  if (c % *g != 0) {
    return false;
  }
  // Absorb the bulk of c into a*dx + b*dy, then scale the base solution by the small remainder so
  // the reported (x, y) stay bounded by max(|a|, |b|, |c|). The scaled product is taken modulo b
  // (resp. a) through a 128-bit intermediate where available, to avoid overflow.
  std::pair<Int, Int> base = extended_gcd(a, b);
  Int dx = c / a;
  c -= dx * a;
  Int dy = c / b;
  c -= dy * b;
  Int f = c / *g;
#if defined(__SIZEOF_INT128__)
  __extension__ typedef std::conditional_t<sizeof(Int) <= 4, int64_t, __int128> wide_t;
#else
  using wide_t = int64_t;
#endif
  *x = dx + static_cast<Int>(static_cast<wide_t>(base.first) * f % b);
  *y = dy + static_cast<Int>(static_cast<wide_t>(base.second) * f % a);
  return true;
}

template<typename Int>
Int mod(Int a, Int m) {
  assert(m > 0);
  Int r = a % m;
  return (r < 0) ? (r + m) : r;
}

template<typename Int>
Int mod_inverse(Int a, Int m) {
  assert(m > 0 && gcd(a, m) == 1);
  return mod(extended_gcd(a, m).first, m);
}

std::vector<int> mod_inverse_table(int p) {
  assert(p >= 2);
  std::vector<int> res(p);
  res[1] = 1;
  for (int i = 2; i < p; i++) {
    res[i] = (p - (p / i) * res[p % i] % p) % p;
  }
  return res;
}

bool crt(int64_t r1, int64_t m1, int64_t r2, int64_t m2, int64_t *r, int64_t *m) {
  r1 = mod(r1, m1);
  r2 = mod(r2, m2);
  int64_t g, x, y;
  if (!diophantine(m1, -m2, r2 - r1, &g, &x, &y)) {
    return false;
  }
  *m = m1 / g * m2;
#if defined(__SIZEOF_INT128__)
  __extension__ typedef __int128 int128_t;
  *r = static_cast<int64_t>(
      (static_cast<int128_t>(r1) + static_cast<int128_t>(m1) * mod(x, m2 / g)) % *m
  );
#else
  *r = mod(r1 + m1 * mod(x, m2 / g), *m);  // Overflow warning.
#endif
  return true;
}

std::vector<int64_t> garner_digits(const std::vector<int> &a, const std::vector<int> &p) {
  assert(a.size() == p.size());
  assert(std::all_of(p.begin(), p.end(), [](int m) { return m > 0; }));
  int n = static_cast<int>(a.size());
  std::vector<int64_t> x(n);
  for (int i = 0; i < n; i++) {
    x[i] = mod(static_cast<int64_t>(a[i]), static_cast<int64_t>(p[i]));
  }
  for (int i = 0; i < n; i++) {
    for (int j = 0; j < i; j++) {
      // Reduce mod p[i] each step; otherwise x[i] compounds to ~p^i and overflows int64_t.
      x[i] = mod_inverse(static_cast<int64_t>(p[j]), static_cast<int64_t>(p[i])) * (x[i] - x[j]);
      x[i] = (x[i] % p[i] + p[i]) % p[i];
    }
  }
  return x;
}

int64_t garner_restore(const std::vector<int> &a, const std::vector<int> &p) {
  assert(a.size() == p.size());
  int n = static_cast<int>(a.size());
  if (n == 0) {
    return 0;
  }
  std::vector<int64_t> x = garner_digits(a, p);
  int64_t res = x[0], m = 1;
  for (int i = 1; i < n; i++) {
    m *= p[i - 1];
    res += x[i] * m;
  }
  return res;
}

int64_t garner_restore_mod(const std::vector<int> &a, const std::vector<int> &p, int64_t m) {
  assert(a.size() == p.size() && m > 0);
  int n = static_cast<int>(a.size());
  if (n == 0) {
    return 0;
  }
  std::vector<int64_t> x = garner_digits(a, p);
  int64_t res = 0;
  for (int i = n - 1; i >= 0; i--) {
#if defined(__SIZEOF_INT128__)
    __extension__ typedef __int128 int128_t;
    res = static_cast<int64_t>((static_cast<int128_t>(res) * p[i] + x[i]) % m);
#else
    res = (res * p[i] + x[i]) % m;  // Requires the intermediate product to fit without int128.
#endif
  }
  return res;
}
\end{lstlisting}
\Needspace*{4\baselineskip}
\textbf{Example Usage}\par
\begin{lstlisting}[style=cpp]
#include <cassert>
using namespace std;

int main() {
  assert(mod(-5, 3) == 1);  // Negative dividends wrap up into [0, m).
  assert(mod(5, 3) == 2);
  assert(mod(-6, 3) == 0);
  {
    for (int a = -20; a <= 20; a++) {
      for (int b = -20; b <= 20; b++) {
        int g = gcd(a, b);
        auto [x, y] = extended_gcd(a, b);
        assert(g == a * x + b * y);
        if (g == 1 && b > 1) {
          assert(mod(a * mod_inverse(a, b), b) == 1);
        }
      }
    }
  }
  {
    int p = 17;
    auto res = mod_inverse_table(p);
    for (int i = 0; i < p; i++) {
      if (i > 0) {
        assert(mod(i * res[i], p) == 1);
      }
    }
  }
  {
    vector<int> a{2, 3, 1}, m{3, 4, 5};
    int x = garner_restore(a, m);
    assert(x == garner_restore_mod(a, m, 1000000007));
    int n = static_cast<int>(a.size());
    for (int i = 0; i < n; i++) {
      assert(mod(x, m[i]) == a[i]);
    }
    assert(x == 11);
  }
  assert(garner_restore(vector<int>{-1}, vector<int>{5}) == 4);
#if defined(__SIZEOF_INT128__)
  {  // Garner modulo another number works even when the product of CRT moduli is too large.
    vector<int> p{1000000007, 1000000009, 1000000033};
    __extension__ typedef __int128 int128_t;
    int128_t x = (int128_t{1} << 80) + 123456789;
    vector<int> a;
    for (int m : p) {
      a.push_back(static_cast<int>(x % m));
    }
    int64_t m = 998244353;
    assert(garner_restore_mod(a, p, m) == static_cast<int64_t>(x % m));
  }
#endif
  {  // Diophantine: exhaustively verify solvability and that solutions satisfy a*x + b*y == c.
    for (int a = -8; a <= 8; a++) {
      for (int b = -8; b <= 8; b++) {
        for (int c = -8; c <= 8; c++) {
          int g, x, y, gg = gcd(a, b);
          bool ok = diophantine(a, b, c, &g, &x, &y);
          assert(ok == (gg == 0 ? c == 0 : c % gg == 0));
          assert(g == gg);  // g is set even when no solution exists.
          if (ok) {
            assert(a * x + b * y == c);
          }
        }
      }
    }
  }
  {  // CRT: merge two congruences, including non-coprime and inconsistent moduli.
    for (int m1 = 1; m1 <= 12; m1++) {
      for (int m2 = 1; m2 <= 12; m2++) {
        for (int r1 = 0; r1 < m1; r1++) {
          for (int r2 = 0; r2 < m2; r2++) {
            int64_t r, m;
            bool ok = crt(r1, m1, r2, m2, &r, &m);
            int limit = lcm(m1, m2), want = -1;
            for (int v = 0; v < limit && want < 0; v++) {
              if (v % m1 == r1 && v % m2 == r2) {
                want = v;
              }
            }
            assert(ok == (want >= 0));
            if (ok) {
              assert(m == limit && r == want);
            }
          }
        }
      }
    }
  }
  return 0;
}
\end{lstlisting}
\subsection{Modular Integer and Combinatorics}
\label{sec:6.3.3}
Provides a modular integer value type together with lazy factorial-based combinatorics tables. Both are common contest helpers for dynamic programming, counting, polynomial operations, and any calculation whose answers are taken modulo some number such as $10^9 + 7$. Where the helpers of section \secref{6.3.1} suit a few scattered operations, a value type is what makes an entire solution read in ordinary arithmetic notation while every operation reduces silently.

The modulus is a template parameter rather than a member, which is a deliberate speed choice: a compile-time constant lets the compiler replace each division by a multiplication and a shift, while a modulus only known at run time forces a hardware divide costing several times as much. Storing it in a \inlinecode{static} member instead makes the same class work with a runtime modulus at that price, and is worth doing when the modulus is read from the input. A modulus above roughly $3 \times 10^9$ needs the \inlinecode{mulmod()} of section \secref{6.3.1} in place of the plain products used here.

The \inlinecode{Modular\textless{}MOD\textgreater{}} value type wraps arithmetic modulo the positive compile-time constant \inlinecode{MOD}. Its signed \inlinecode{auto} argument determines the storage type: \inlinecode{Modular\textless{}1000000007\textgreater{}} uses 32-bit storage, while \inlinecode{Modular\textless{}(1LL \textless{}\textless{} 61) - 1\textgreater{}} uses 64-bit storage. The implementation follows the conventional \inlinecode{Mint} wrapper: construction normalizes values, while hidden friend operators support mixed expressions such as \inlinecode{2 + Mint(3)} through implicit conversion.

\begin{compactitem}
  \item \inlinecode{Modular\textless{}MOD\textgreater{}(x = 0)} constructs the residue class of integer \inlinecode{x} modulo \inlinecode{MOD}.
  \item \inlinecode{Modular\textless{}MOD\textgreater{}::mod()} returns the modulus \inlinecode{MOD}.
  \item \inlinecode{value()} and \inlinecode{operator()()} return the stored representative in $[0, {\inlinecodemath{MOD}})$.
  \item Explicit casts to \inlinecode{int}, \inlinecode{long long}, \inlinecode{double}, and \inlinecode{long double} convert that stored representative to the corresponding primitive.
  \item \inlinecode{pow(n)} returns this value raised to nonnegative integer exponent \inlinecode{n}.
  \item \inlinecode{inv()} returns the multiplicative inverse, asserting the value is coprime to \inlinecode{MOD}.
  \item Operators \inlinecode{+}, \inlinecode{-}, \inlinecode{*}, \inlinecode{/}, comparison, increment, decrement, and stream I/O are overloaded. Division likewise requires the divisor to be coprime to \inlinecode{MOD}.
\end{compactitem}

Overflow warning: Construction, multiplication, and inverses widen through \inlinecode{int64\_t} for 32-bit storage or \inlinecode{\_\_int128} for 64-bit storage. Thus $2p$ must fit the storage type and $(p - 1)^2$ must fit the intermediate type. On compilers without \inlinecode{\_\_int128}, only 32-bit moduli are supported. Estimating a 64-bit product with \inlinecode{long double} is not used because it silently loses the needed precision on platforms where \inlinecode{long double} is only 64 bits.

\Needspace*{3\baselineskip}
\textbf{Time Complexity}:
\begin{compactitem}
  \item $O(1)$ per addition, subtraction, multiplication, comparison, and stream output.
  \item $O(\log n)$ per call to \inlinecode{pow()}.
  \item $O(\log m)$ per call to \inlinecode{inv()} and division, where $m = {\inlinecodemath{MOD}}$.
\end{compactitem}

\textbf{Space Complexity}: $O(1)$ auxiliary for Modular arithmetic.

\begin{lstlisting}[style=cpp]
#include <cassert>
#include <cstdint>
#include <iostream>
#include <limits>
#include <type_traits>
#include <utility>
#include <vector>

template<typename T>
T inverse(T a, T m) {
  T original_m = m;
  T u = 0, v = 1;
  while (a != 0) {
    T q = m / a;
    m -= q * a;
    std::swap(a, m);
    u -= q * v;
    std::swap(u, v);
  }
  assert(m == 1);
  u %= original_m;
  return u < 0 ? u + original_m : u;
}

template<auto MOD>
class Modular {
  using T = decltype(MOD);
  static_assert(std::is_integral_v<T> && std::is_signed_v<T>);
  static_assert(MOD > 0 && MOD <= std::numeric_limits<T>::max() / 2);

  // Widen through int64_t for 32-bit storage, or __int128 for 64-bit storage where it exists
  // (__extension__ silences -pedantic). Without __int128, only 32-bit moduli are supported.
#if defined(__SIZEOF_INT128__)
  __extension__ typedef std::conditional_t<sizeof(T) <= 4, int64_t, __int128> wide_t;
#else
  using wide_t = int64_t;
  static_assert(sizeof(T) <= 4, "64-bit moduli require __int128, which is unavailable");
#endif

  T v;

  static T normalize(wide_t x) {
    if (-static_cast<wide_t>(mod()) <= x && x < static_cast<wide_t>(mod())) {
      T y = static_cast<T>(x);
      return y < 0 ? y + mod() : y;
    }
    x %= mod();
    if (x < 0) {
      x += mod();
    }
    return static_cast<T>(x);
  }

 public:
  Modular(wide_t x = 0) : v(normalize(x)) {}

  static T mod() { return MOD; }
  T value() const { return v; }
  T operator()() const { return v; }
  explicit operator int() const { return static_cast<int>(v); }
  explicit operator long long() const { return static_cast<long long>(v); }
  explicit operator double() const { return static_cast<double>(v); }
  explicit operator long double() const { return static_cast<long double>(v); }

  Modular pow(int64_t n) const {
    assert(n >= 0);
    Modular base = *this, res = 1;
    while (n > 0) {
      if (n & 1) {
        res *= base;
      }
      base *= base;
      n >>= 1;
    }
    return res;
  }

  Modular inv() const {
    assert(v != 0);
    return Modular(inverse(static_cast<wide_t>(v), static_cast<wide_t>(mod())));
  }

  Modular &operator+=(const Modular &other) {
    v += other.v;
    if (v >= mod()) {
      v -= mod();
    }
    return *this;
  }

  Modular &operator-=(const Modular &other) {
    v -= other.v;
    if (v < 0) {
      v += mod();
    }
    return *this;
  }

  Modular &operator*=(const Modular &other) {
    v = normalize(static_cast<wide_t>(v) * other.v);
    return *this;
  }

  Modular operator-() const { return Modular(-v); }
  Modular &operator++() { return *this += 1; }
  Modular &operator--() { return *this -= 1; }
  Modular &operator/=(const Modular &other) { return *this *= other.inv(); }
  Modular operator++(int) { Modular res = *this; ++*this; return res; }
  Modular operator--(int) { Modular res = *this; --*this; return res; }
  friend Modular operator+(Modular a, const Modular &b) { return a += b; }
  friend Modular operator-(Modular a, const Modular &b) { return a -= b; }
  friend Modular operator*(Modular a, const Modular &b) { return a *= b; }
  friend Modular operator/(Modular a, const Modular &b) { return a /= b; }
  friend bool operator==(const Modular &a, const Modular &b) { return a.v == b.v; }
  friend bool operator!=(const Modular &a, const Modular &b) { return !(a == b); }
  friend bool operator<(const Modular &a, const Modular &b) { return a.v < b.v; }
  friend std::ostream &operator<<(std::ostream &os, const Modular &x) { return os << x.v; }

  friend std::istream &operator>>(std::istream &is, Modular &x) {
    int64_t value;
    is >> value;
    x = Modular(value);
    return is;
  }
};
\end{lstlisting}
\inlinecode{ModCombinatorics\textless{}Mint\textgreater{}} maintains lazy factorial and inverse-factorial tables for a modular type \inlinecode{Mint}. It assumes the modulus is prime and all requested factorials are invertible. Growing the tables costs one modular inverse rather than one per entry: after extending the factorials by repeated multiplication, a single inversion of the largest factorial seeds a backward pass in which each inverse factorial follows from the next by one multiplication.

\begin{compactitem}
  \item \inlinecode{factorial(n)} returns $n!$ using a lazy factorial table.
  \item \inlinecode{choose(n, k)} returns $\binom n k$ using lazy factorial and inverse-factorial tables.
  \item \inlinecode{permute(n, k)} returns the number of ordered selections of \inlinecode{k} distinct elements from \inlinecode{n}.
  \item \inlinecode{multichoose(n, k)} returns the number of size-\inlinecode{k} multisets drawn from \inlinecode{n} types.
\end{compactitem}

Tables pay off once a solution asks for many values against one fixed modulus, which is the usual counting or dynamic programming case. For a handful of values, or for a modulus that changes between calls, the table-free routines of section \secref{6.2.1} compute the same quantities in $O(\min\left(k, n - k\right))$ per call and allocate nothing.

\textbf{Time Complexity}: $O(n)$ total table growth to answer calls with arguments up to $n$, then $O(1)$ per call.

\textbf{Space Complexity}: $O(n)$ for the factorial and inverse-factorial tables.

\begin{lstlisting}[style=cpp]
template<typename Mint>
class ModCombinatorics {
  std::vector<Mint> fact, inv_fact;

  void ensure(int n) {
    auto old = static_cast<int>(fact.size());
    if (old > n) {
      return;
    }
    fact.resize(n + 1);
    inv_fact.resize(n + 1);
    for (int i = old; i <= n; i++) {
      fact[i] = fact[i - 1] * i;
    }
    inv_fact[n] = fact[n].inv();
    for (int i = n; i > old; i--) {
      inv_fact[i - 1] = inv_fact[i] * i;
    }
  }

 public:
  ModCombinatorics() : fact(1, Mint{1}), inv_fact(1, Mint{1}) {}

  Mint factorial(int n) {
    ensure(n);
    return fact[n];
  }

  Mint choose(int n, int k) {
    if (k < 0 || k > n) {
      return 0;
    }
    ensure(n);
    return fact[n] * inv_fact[k] * inv_fact[n - k];
  }

  Mint permute(int n, int k) {
    if (k < 0 || k > n) {
      return 0;
    }
    ensure(n);
    return fact[n] * inv_fact[n - k];
  }

  Mint multichoose(int n, int k) { return n == 0 ? Mint(k == 0) : choose(n + k - 1, k); }
};

using Mint = Modular<1000000007>;
\end{lstlisting}
\Needspace*{4\baselineskip}
\textbf{Example Usage}\par
\begin{lstlisting}[style=cpp]
#include <cassert>
#include <sstream>
using namespace std;

int main() {
  Mint a = 1000000008LL;
  Mint b = -2;
  assert(a.value() == 1);
  assert(b() == Mint::mod() - 2);
  assert(static_cast<int>(a) == 1);
  assert(static_cast<long long>(a) == 1LL);
  assert(static_cast<double>(a) == 1.0);
  assert(a + b == -1);
  assert(2 + Mint(3) == 5);
  assert(Mint(2) * 3 == 6);
  assert(Mint(2).pow(10) == 1024);
  assert(Mint(5) / 5 == 1);

  Mint x = 0;
  x++;
  ++x;
  assert(x == 2);
  stringstream ss("1000000008");
  ss >> x;
  assert(x == 1);

  ModCombinatorics<Mint> comb;
  assert(comb.factorial(5) == 120);
  assert(comb.choose(5, 2) == 10);
  assert(comb.permute(5, 2) == 20);
  assert(comb.multichoose(3, 2) == 6);
  assert(comb.multichoose(0, 0) == 1);

  // Each distinct modulus is just another instantiation.
  using Mint2 = Modular<998244353>;
  assert(Mint2(2).pow(10) == 1024);
  assert(Mint2(-1) == Mint2::mod() - 1);

#if defined(__SIZEOF_INT128__)
  // A 64-bit modulus (deduced as long long) automatically widens through __int128.
  using Mint3 = Modular<(1LL << 61) - 1>;
  assert(Mint3(2).pow(62) == 2);  // 2^62 mod (2^61 - 1)
  assert(Mint3(3).inv() * 3 == 1);
#endif
  return 0;
}
\end{lstlisting}
\subsection{Primitive Roots}
\label{sec:6.3.4}
A primitive root modulo $m$ is a number $g$ whose powers generate every invertible residue modulo $m$. Equivalently, the multiplicative order of $g$ modulo $m$ is exactly $\phi(m)$. Primitive roots are useful for constructing generators of finite multiplicative groups, choosing roots for number theoretic transforms, discrete logarithm setups, and randomized modular hashing or testing schemes that need a full-period multiplicative generator.

Primitive roots do not exist for every $m$. They exist exactly for $m \in \{1, 2, 4, p^k, 2p^k\}$, where $p$ is an odd prime and $k \geq 1$. This implementation first checks that classification by factoring $m$, then factors $\phi(m)$. A candidate $g$ is primitive exactly when $g^{\phi(m)} \equiv 1 \pmod m$ and $g^{\phi(m)/q} \not\equiv 1 \pmod m$ for every prime factor $q$ of $\phi(m)$.

For a prime modulus $p$, this reduces to the common contest test: factor $p - 1$, then scan for the first $g$ such that $g^{(p - 1)/q} \not\equiv 1 \pmod p$ for every prime factor $q$ of $p - 1$.

\begin{compactitem}
  \item \inlinecode{primitive\_root(m)} returns the smallest primitive root modulo \inlinecode{m}, or $-1$ if no primitive root exists.
\end{compactitem}

Modular multiplication uses \inlinecode{\_\_uint128\_t} when available for speed, with a slower portable double-and-add fallback to avoid overflow on compilers without 128-bit integers.

\textbf{Time Complexity}: $O(\sqrt{m} + g \cdot k \cdot \log m)$ per call, where $g$ is the answer candidate tested and $k$ is the number of distinct prime factors of $\phi(m)$.

\textbf{Space Complexity}: $O(k)$ auxiliary.

\begin{lstlisting}[style=cpp]
#include <cassert>
#include <cstdint>
#include <numeric>
#include <set>
#include <utility>
#include <vector>

int64_t mulmod(int64_t a, int64_t b, int64_t m) {
  assert(a >= 0 && b >= 0 && m > 0);
#if defined(__SIZEOF_INT128__)
  return static_cast<int64_t>(static_cast<__uint128_t>(a) * b % m);
#else
  int64_t res = 0;
  for (a %= m; b > 0; b >>= 1) {
    if (b & 1) {
      res = res >= m - a ? res - (m - a) : res + a;
    }
    a = a >= m - a ? a - (m - a) : a + a;
  }
  return res;
#endif
}

int64_t powmod(int64_t x, int64_t n, int64_t m) {
  int64_t res = 1 % m;
  for (x %= m; n > 0; n >>= 1) {
    if (n & 1) {
      res = mulmod(res, x, m);
    }
    x = mulmod(x, x, m);
  }
  return res;
}

std::vector<std::pair<int64_t, int>> factorize(int64_t n) {
  std::vector<std::pair<int64_t, int>> res;
  for (int64_t p = 2; p <= n / p; p += (p == 2 ? 1 : 2)) {
    if (n % p != 0) {
      continue;
    }
    int count = 0;
    while (n % p == 0) {
      n /= p;
      count++;
    }
    res.emplace_back(p, count);
  }
  if (n > 1) {
    res.emplace_back(n, 1);
  }
  return res;
}

int64_t primitive_root(int64_t m) {
  if (m <= 0) {
    return -1;
  }
  if (m == 1 || m == 2 || m == 4) {
    return m - 1;
  }
  std::vector<std::pair<int64_t, int>> factors = factorize(m);
  bool has_root = (factors.size() == 1 && factors[0].first != 2) ||
                  (factors.size() == 2 && factors[0].first == 2 && factors[0].second == 1);
  if (!has_root) {
    return -1;
  }
  int64_t phi = m;
  std::set<int64_t> phi_factors;
  for (auto [p, e] : factors) {
    phi = phi / p * (p - 1);
    if (e > 1) {
      phi_factors.insert(p);
    }
    for (auto [q, count] : factorize(p - 1)) {
      phi_factors.insert(q);
    }
  }
  for (int64_t g = 1; g < m; g++) {
    if (std::gcd(g, m) != 1 || powmod(g, phi, m) != 1) {
      continue;
    }
    bool ok = true;
    for (int64_t q : phi_factors) {
      if (powmod(g, phi / q, m) == 1) {
        ok = false;
        break;
      }
    }
    if (ok) {
      return g;
    }
  }
  return -1;
}
\end{lstlisting}
\Needspace*{4\baselineskip}
\textbf{Example Usage}\par
\begin{lstlisting}[style=cpp]
#include <cassert>

int main() {
  assert(primitive_root(1) == 0);
  assert(primitive_root(2) == 1);
  assert(primitive_root(4) == 3);

  assert(primitive_root(17) == 3);         // 3 has order phi(17) = 16.
  assert(primitive_root(998244353) == 3);  // Common NTT prime.
  assert(primitive_root(9) == 2);          // Odd prime power.
  assert(primitive_root(18) == 5);         // Twice an odd prime power.

  assert(primitive_root(8) == -1);
  assert(primitive_root(12) == -1);
  return 0;
}
\end{lstlisting}
\subsection{Discrete Logarithm (Baby-Step Giant-Step)}
\label{sec:6.3.5}
Given integers $a$, $b$, and a modulus $m$, the discrete logarithm problem asks for the smallest nonnegative integer $x$ with $a^x$ congruent to $b$ modulo $m$. This is the modular analogue of a logarithm and underlies Diffie-Hellman key exchange, order computations in the multiplicative group, and many number-theory reductions.

The baby-step giant-step method writes $x = np - q$ for a block size $n$ near the square root of the modulus. It tabulates the ``baby steps'' $b \cdot a^q$ for $q$ in $[0, n]$, then walks the ``giant steps'' $a^{np}$ for increasing $p$ until one lands in the table, yielding $x = np - q$. This trades the linear scan of all exponents for $O(\sqrt{m})$ work and memory. A short reduction loop first strips common factors between $a$ and $m$, so the routine works for any modulus, not only when $a$ and $m$ are coprime.

\begin{compactitem}
  \item \inlinecode{discrete\_log(a, b, m)} returns the smallest nonnegative \inlinecode{x} with \inlinecode{a\textasciicircum{}x} congruent to \inlinecode{b} modulo \inlinecode{m}, or $-1$ if no such \inlinecode{x} exists. Requires \inlinecode{m} $\geq 1$; \inlinecode{a} and \inlinecode{b} are reduced modulo \inlinecode{m} internally.
\end{compactitem}

\textbf{Time Complexity}: $O(\sqrt{m})$ expected per call with constant-time modular multiplication, and $O(\sqrt{m} \cdot \log m)$ with the portable double-and-add fallback.

\textbf{Space Complexity}: $O(\sqrt{m})$ auxiliary for the table.

\begin{lstlisting}[style=cpp]
#include <cassert>
#include <cmath>
#include <cstdint>
#include <numeric>
#include <unordered_map>

int64_t mulmod(int64_t a, int64_t b, int64_t m) {
  assert(a >= 0 && b >= 0 && m > 0);
#if defined(__SIZEOF_INT128__)
  return static_cast<int64_t>(static_cast<__uint128_t>(a) * b % m);
#else
  int64_t res = 0;
  for (a %= m; b > 0; b >>= 1) {
    if (b & 1) {
      res = res >= m - a ? res - (m - a) : res + a;
    }
    a = a >= m - a ? a - (m - a) : a + a;
  }
  return res;
#endif
}

int64_t discrete_log(int64_t a, int64_t b, int64_t m) {
  assert(m >= 1);
  a %= m;
  b %= m;
  if (a < 0) {
    a += m;
  }
  if (b < 0) {
    b += m;
  }
  // Strip common factors of a and m so that a becomes invertible; add counts the exponents removed.
  int64_t k = 1 % m, add = 0, g;
  while ((g = std::gcd(a, m)) > 1) {
    if (b == k) {
      return add;
    }
    if (b % g != 0) {
      return -1;
    }
    b /= g;
    m /= g;
    add++;
    k = mulmod(k, a / g, m);
  }
  int64_t n = static_cast<int64_t>(std::sqrt(static_cast<double>(m))) + 1;
  int64_t an = 1;
  for (int64_t i = 0; i < n; i++) {
    an = mulmod(an, a, m);  // an = a^n, the giant step.
  }
  std::unordered_map<int64_t, int64_t> baby;
  int64_t cur = b;
  for (int64_t q = 0; q <= n; q++) {
    baby[cur] = q;  // Keep the largest q for each value to make the answer minimal.
    cur = mulmod(cur, a, m);
  }
  cur = k;
  for (int64_t p = 1; p <= n; p++) {
    cur = mulmod(cur, an, m);
    auto it = baby.find(cur);
    if (it != baby.end()) {
      return n * p - it->second + add;
    }
  }
  return -1;
}
\end{lstlisting}
\Needspace*{4\baselineskip}
\textbf{Example Usage}\par
\begin{lstlisting}[style=cpp]
#include <cassert>

int main() {
  // 2^x = 9 (mod 11): 2^6 = 64 = 9, and 6 is the smallest such exponent.
  assert(discrete_log(2, 9, 11) == 6);

  assert(discrete_log(3, 1, 7) == 0);    // a^0 = 1 always.
  assert(discrete_log(2, 3, 5) == 3);    // 2^3 = 8 = 3 (mod 5).
  assert(discrete_log(2, 1, 6) == 0);    // Works for composite, non-coprime moduli.
  assert(discrete_log(4, 6, 8) == -1);   // No power of 4 is 6 (mod 8).
  assert(discrete_log(5, 33, 58) == 9);  // 5^9 = 33 (mod 58).
  return 0;
}
\end{lstlisting}
\subsection{Modular Square Root (Tonelli-Shanks)}
\label{sec:6.3.6}
Given an integer $a$ and an odd prime $p$, find a modular square root: an $x$ with $x^2$ congruent to $a$ modulo $p$. Exactly half of the nonzero residues modulo a prime are squares (quadratic residues), so a solution exists only for those; when it does, the two roots are $x$ and $p - x$.

The Tonelli-Shanks algorithm handles the general case. Euler's criterion, $a^{(p-1)/2}$ modulo $p$, first tests whether $a$ is a residue at all. When $p$ is congruent to $3 \pmod{4}$, the root is simply $a^{(p+1)/4}$. Otherwise the algorithm factors $p - 1$ as $s \cdot 2^e$, picks any quadratic non-residue to generate the $2$-power part of the group, and repeatedly squares a running value to find the order of the current error term, correcting the candidate root until the error becomes $1$.

\begin{compactitem}
  \item \inlinecode{mod\_sqrt(a, p)} returns the smaller of the two square roots of \inlinecode{a} modulo the odd prime \inlinecode{p}, or $-1$ if \inlinecode{a} is not a quadratic residue. The value $0$ maps to $0$. The argument \inlinecode{a} is reduced modulo \inlinecode{p} internally. The modulus \inlinecode{p} must be an odd prime.
\end{compactitem}

\textbf{Time Complexity}: $O(\log^{2} p + z \log p)$ per call, where $z$ is the smallest integer at least $2$ that is a quadratic non-residue modulo $p$; thus the unconditional worst-case bound is $O(p \log p)$.

\textbf{Space Complexity}: $O(1)$ auxiliary.

\begin{lstlisting}[style=cpp]
#include <cassert>
#include <cstdint>

int64_t mulmod(int64_t a, int64_t b, int64_t m) {
  assert(a >= 0 && b >= 0 && m > 0);
#if defined(__SIZEOF_INT128__)
  return static_cast<int64_t>(static_cast<__uint128_t>(a) * b % m);
#else
  int64_t res = 0;
  for (a %= m; b > 0; b >>= 1) {
    if (b & 1) {
      res = res >= m - a ? res - (m - a) : res + a;
    }
    a = a >= m - a ? a - (m - a) : a + a;
  }
  return res;
#endif
}

int64_t powmod(int64_t x, int64_t n, int64_t m) {
  int64_t res = 1 % m;
  for (x %= m; n > 0; n >>= 1) {
    if (n & 1) {
      res = mulmod(res, x, m);
    }
    x = mulmod(x, x, m);
  }
  return res;
}

int64_t mod_sqrt(int64_t a, int64_t p) {
  a %= p;
  if (a < 0) {
    a += p;
  }
  if (a == 0) {
    return 0;
  }
  if (powmod(a, (p - 1) / 2, p) != 1) {
    return -1;  // a is a quadratic non-residue.
  }
  if (p % 4 == 3) {
    int64_t x = powmod(a, p / 4 + 1, p);
    return x < p - x ? x : p - x;
  }
  // Write p - 1 = s * 2^e with s odd.
  int64_t s = p - 1, e = 0;
  while (s % 2 == 0) {
    s /= 2;
    e++;
  }
  // Any quadratic non-residue generates the 2-power part of the group.
  int64_t z = 2;
  while (powmod(z, (p - 1) / 2, p) != p - 1) {
    z++;
  }
  int64_t x = powmod(a, (s + 1) / 2, p);  // Candidate root, off by a 2-power factor.
  int64_t b = powmod(a, s, p);            // Error term, a power of two in order.
  int64_t g = powmod(z, s, p);            // Generator of the relevant 2-group.
  int64_t r = e;
  while (true) {
    int64_t t = b, order = 0;
    while (order < r && t != 1) {
      t = mulmod(t, t, p);
      order++;
    }
    if (order == 0) {
      return x < p - x ? x : p - x;
    }
    int64_t gs = powmod(g, 1LL << (r - order - 1), p);
    g = mulmod(gs, gs, p);
    x = mulmod(x, gs, p);
    b = mulmod(b, g, p);
    r = order;
  }
}
\end{lstlisting}
\Needspace*{4\baselineskip}
\textbf{Example Usage}\par
\begin{lstlisting}[style=cpp]
#include <cassert>

int main() {
  int64_t r = mod_sqrt(2, 113);
  assert(r != -1 && mulmod(r, r, 113) == 2);
  assert(r == 51);  // 51^2 = 2601 = 2 (mod 113); the other root is 62.

  assert(mod_sqrt(0, 7) == 0);
  assert(mod_sqrt(5, 41) != -1);  // p = 1 (mod 4) exercises the full algorithm.
  assert(mod_sqrt(4, 7) == 2);    // p = 3 (mod 4) fast path.
  assert(mod_sqrt(3, 7) == -1);   // 3 is a non-residue modulo 7.

  // Every residue's reported root squares back to it.
  for (int64_t a = 0; a < 41; a++) {
    int64_t x = mod_sqrt(a, 41);
    if (x != -1) {
      assert(mulmod(x, x, 41) == a);
    }
  }
  return 0;
}
\end{lstlisting}
\subsection{Jacobi Symbol}
\label{sec:6.3.7}
The Jacobi symbol $(a / n)$, defined for an odd positive integer $n$, generalizes the Legendre symbol. For an odd prime $p$ the Legendre symbol $(a / p)$ is $0$ when $p$ divides $a$, $1$ when $a$ is a nonzero quadratic residue modulo $p$, and $-1$ when $a$ is a non-residue. The Jacobi symbol extends this to composite $n$ by multiplying the Legendre symbols of the prime factors of $n$, counted with multiplicity. It is computed without factoring $n$, using the law of quadratic reciprocity together with the supplementary rules for $-1$ and $2$, in the style of the binary GCD.

The symbol is fully multiplicative in $a$ and in $n$. For prime $n$ it exactly decides quadratic residuosity, so it is the fast way to test whether a modular square root exists. For composite $n$, beware that $(a / n) = 1$ does not guarantee that $a$ is a quadratic residue; this asymmetry is what the Solovay-Strassen primality test exploits.

\begin{compactitem}
  \item \inlinecode{jacobi(a, n)} returns the Jacobi symbol (\inlinecode{a} / \inlinecode{n}) as one of $-1$, $0$, or $1$. The modulus \inlinecode{n} must be a positive odd integer; \inlinecode{a} is reduced modulo \inlinecode{n} internally and may be negative. The result is $0$ exactly when \inlinecode{a} and \inlinecode{n} share a common factor.
\end{compactitem}

\textbf{Time Complexity}: $O(\log n)$ arithmetic operations per call.

\textbf{Space Complexity}: $O(1)$ auxiliary.

\begin{lstlisting}[style=cpp]
#include <cassert>
#include <cstdint>
#include <utility>

int jacobi(int64_t a, int64_t n) {
  assert(n > 0 && n % 2 == 1);
  a %= n;
  if (a < 0) {
    a += n;
  }
  int result = 1;
  while (a != 0) {
    while (a % 2 == 0) {
      a /= 2;
      // (2 / n) is -1 when n is 3 or 5 modulo 8.
      if (n % 8 == 3 || n % 8 == 5) {
        result = -result;
      }
    }
    std::swap(a, n);
    // Quadratic reciprocity flips the sign when both are 3 modulo 4.
    if (a % 4 == 3 && n % 4 == 3) {
      result = -result;
    }
    a %= n;
  }
  return n == 1 ? result : 0;
}
\end{lstlisting}
\Needspace*{4\baselineskip}
\textbf{Example Usage}\par
\begin{lstlisting}[style=cpp]
int64_t legendre_via_euler(int64_t a, int64_t p) {
  int64_t r = 1 % p;
  for (int64_t b = ((a % p) + p) % p, e = (p - 1) / 2; e > 0; e >>= 1) {
    if (e & 1) {
      r = r * b % p;
    }
    b = b * b % p;
  }
  return r == p - 1 ? -1 : r;  // r is 0, 1, or p - 1.
}

int main() {
  assert(jacobi(2, 15) == 1);  // (2/15) = (2/3)(2/5) = (-1)(-1).
  assert(jacobi(7, 15) == -1);
  assert(jacobi(3, 15) == 0);  // gcd(3, 15) = 3.
  assert(jacobi(1001, 9907) == -1);

  // For odd primes the Jacobi symbol equals the Legendre symbol from Euler's criterion.
  for (int64_t p : {3, 5, 7, 11, 13, 101, 9973}) {
    for (int64_t a = 0; a < 50; a++) {
      assert(jacobi(a, p) == legendre_via_euler(a, p));
    }
  }
  return 0;
}
\end{lstlisting}
\subsection{Prime Generation}
\label{sec:6.3.8}
Generate prime numbers using the sieve of Eratosthenes, which repeatedly marks the multiples of each prime as composite. The linear sieve technique optimizes the classic sieve by finding the least prime factor of each integer. When processing \inlinecode{i}, it marks \inlinecode{i * p} for each known prime \inlinecode{p} no greater than \inlinecode{least[i]}. Every composite \inlinecode{x} is therefore generated exactly once: writing \inlinecode{x = i*p}, where \inlinecode{p} is its least prime factor, gives \inlinecode{p} $\leq$ \inlinecode{least[i]}, while any representation using a larger prime is skipped. This single visit per composite gives the $O(n)$ running time.

The segmented sieve technique computes primes in a high but narrow window without sieving everything below it. Since every composite at most \inlinecode{hi} has a prime factor at most $\sqrt{{\inlinecodemath{hi}}}$, it first sieves the primes up to $\sqrt{{\inlinecodemath{hi}}}$, then uses only those primes to mark composites within $[{\inlinecodemath{lo}}, {\inlinecodemath{hi}}]$ (each prime $p$ starting at the first of its multiples that is at least both $p^2$ and \inlinecode{lo}). This needs only $O(w + \sqrt{{\inlinecodemath{hi}}})$ space, where $w = {\inlinecodemath{hi}} - {\inlinecodemath{lo}} + 1$, instead of the $O({\inlinecodemath{hi}})$ space that sieving all of $[2, {\inlinecodemath{hi}}]$ would require.

\begin{compactitem}
  \item \inlinecode{sieve(n)} returns a vector of all the primes less than or equal to \inlinecode{n}.
  \item \inlinecode{linear\_sieve(n, \&least\_out)} returns a vector of all primes less than or equal to \inlinecode{n} in linear time. If \inlinecode{least\_out} is not null, it is filled so that \inlinecode{least\_out[x]} is the least prime factor of \inlinecode{x} for every \inlinecode{x} $\geq 2$.
  \item \inlinecode{segmented\_sieve(lo, hi)} returns a vector of all the primes in the range $[{\inlinecodemath{lo}}, {\inlinecodemath{hi}}]$.
\end{compactitem}

\Needspace*{3\baselineskip}
\textbf{Time Complexity}:
\begin{compactitem}
  \item $O(n \log \log n)$ per call to \inlinecode{sieve(n)}.
  \item $O(n)$ per call to \inlinecode{linear\_sieve()}.
  \item $O((w + \sqrt{{\inlinecodemath{hi}}}) \cdot \log\left(\log\left({\inlinecodemath{hi}}\right)\right))$ per call to \inlinecode{segmented\_sieve(lo, hi)}, where $w = {\inlinecodemath{hi}} - {\inlinecodemath{lo}} + 1$.
\end{compactitem}

\Needspace*{3\baselineskip}
\textbf{Space Complexity}:
\begin{compactitem}
  \item $O(n)$ auxiliary for \inlinecode{sieve(n)}.
  \item $O(n)$ auxiliary for \inlinecode{linear\_sieve()}.
  \item $O(w + \sqrt{{\inlinecodemath{hi}}})$ auxiliary for \inlinecode{segmented\_sieve(lo, hi)}.
\end{compactitem}

\begin{lstlisting}[style=cpp]
#include <algorithm>
#include <cmath>
#include <cstdint>
#include <vector>

std::vector<int> sieve(int n) {
  if (n < 2) {
    return {};
  }
  std::vector<char> prime(n + 1, true);
  for (int i = 2; i <= n / i; i++) {
    if (prime[i]) {
      for (int64_t j = 1LL * i * i; j <= n; j += i) {
        prime[j] = false;
      }
    }
  }
  std::vector<int> res;
  for (int i = 2; i <= n; i++) {
    if (prime[i]) {
      res.push_back(i);
    }
  }
  return res;
}

std::vector<int> linear_sieve(int n, std::vector<int> *least_out = nullptr) {
  if (n < 2) {
    if (least_out != nullptr) {
      least_out->assign(std::max(n + 1, 0), 0);
    }
    return {};
  }
  std::vector<int> least(n + 1), primes;
  for (int i = 2; i <= n; i++) {
    if (least[i] == 0) {
      least[i] = i;
      primes.push_back(i);
    }
    for (int p : primes) {
      if (p > least[i] || i > n / p) {
        break;
      }
      least[i * p] = p;
    }
  }
  if (least_out != nullptr) {
    *least_out = least;
  }
  return primes;
}

std::vector<int> segmented_sieve(int lo, int hi) {
  if (hi < 2 || lo > hi) {
    return {};
  }
  int sqrt_hi = std::sqrt(hi);
  while (1LL * (sqrt_hi + 1) * (sqrt_hi + 1) <= hi) {
    sqrt_hi++;
  }
  int fourth_root_hi = std::sqrt(sqrt_hi);
  while (1LL * (fourth_root_hi + 1) * (fourth_root_hi + 1) <= sqrt_hi) {
    fourth_root_hi++;
  }
  std::vector<char> prime1(sqrt_hi + 1, true), prime2(hi - lo + 1, true);
  for (int i = 2; i <= fourth_root_hi; i++) {
    if (prime1[i]) {
      for (int64_t j = 1LL * i * i; j <= sqrt_hi; j += i) {
        prime1[j] = false;
      }
    }
  }
  for (int i = 2; i <= sqrt_hi; i++) {
    if (prime1[i]) {
      int64_t first_multiple = ((static_cast<int64_t>(lo) + i - 1) / i) * i;
      int64_t start = std::max(static_cast<int64_t>(i) * i, first_multiple);
      for (int64_t j = start; j <= hi; j += i) {
        prime2[j - lo] = false;
      }
    }
  }
  std::vector<int> res;
  for (int i = (lo > 1) ? lo : 2; i <= hi; i++) {
    if (prime2[i - lo]) {
      res.push_back(i);
    }
  }
  return res;
}
\end{lstlisting}
\Needspace*{4\baselineskip}
\textbf{Example Usage}\par
\begin{lstlisting}[style=cpp]
#include <cassert>
#include <ctime>
#include <iostream>
using namespace std;

int main() {
  int pmax = 10000000;
  time_t start;
  double delta;

  start = clock();
  auto p = sieve(pmax);
  delta = static_cast<double>((clock() - start)) / CLOCKS_PER_SEC;
  cout << "sieve(n=" << pmax << "): " << delta << "s" << endl;
  assert((sieve(1) == vector<int>{}));
  assert((sieve(10) == vector<int>{2, 3, 5, 7}));

  vector<int> least;
  assert(linear_sieve(pmax, &least) == p);
  assert(least[999983] == 999983);
  assert(least[999984] == 2);

  int lo = 1000000000, hi = 1005000000;
  start = clock();
  p = segmented_sieve(lo, hi);
  delta = static_cast<double>((clock() - start)) / CLOCKS_PER_SEC;
  cout << "segmented_sieve([" << lo << ", " << hi << "]): " << delta << "s" << endl;
  assert((segmented_sieve(14, 16) == vector<int>{}));
  assert((segmented_sieve(17, 19) == vector<int>{17, 19}));
  assert(!p.empty() && p.front() >= lo && p.back() <= hi);
  return 0;
}
\end{lstlisting}
\begin{exampleoutput}
sieve(n=10000000): 0.042641s
segmented_sieve([1000000000, 1005000000]): 0.01969s
\end{exampleoutput}
\subsection{Prime Tests and Factorization}
\label{sec:6.3.9}
Provides primality tests, prime factorization, and divisor generation. Alongside generic trial division, it includes deterministic Miller-Rabin and Pollard's rho for signed 64-bit integers. Prime factorizations are represented as sorted vectors of (\inlinecode{prime}, \inlinecode{exponent}) pairs. For $0$ and $1$, the prime factorization is empty.

Trial division tests primality and computes prime factorizations by trying possible divisors through the square root.

\begin{compactitem}
  \item \inlinecode{is\_prime\_slow(n)} returns whether the integer \inlinecode{n} is prime using trial division.
  \item \inlinecode{factorize\_slow(n)} returns the prime factorization of \inlinecode{n} using trial division.
\end{compactitem}

\textbf{Time Complexity}: $O(\sqrt{n})$ per call to \inlinecode{is\_prime\_slow()} and \inlinecode{factorize\_slow()}.

\textbf{Space Complexity}: $O(1)$ auxiliary for \inlinecode{is\_prime\_slow()} and $O(f)$ for the returned factorization, where $f$ is the number of compressed prime factors.

\begin{lstlisting}[style=cpp]
#include <algorithm>
#include <cassert>
#include <chrono>
#include <cstdint>
#include <numeric>
#include <random>
#include <utility>
#include <vector>

template<typename Int>
bool is_prime_slow(Int n) {
  if (n == 2 || n == 3) {
    return true;
  }
  if (n < 2 || n % 2 == 0 || n % 3 == 0) {
    return false;
  }
  for (Int i = 5, w = 4; i <= n / i; i += w) {
    if (n % i == 0) {
      return false;
    }
    w = 6 - w;
  }
  return true;
}

template<typename Int>
std::vector<std::pair<Int, int>> factorize_slow(Int n) {
  if (n <= 1) {
    return {};
  }
  std::vector<std::pair<Int, int>> res;
  for (Int i = 2; i <= n / i; i++) {
    if (n % i == 0) {
      int exponent = 0;
      do {
        n /= i;
        exponent++;
      } while (n % i == 0);
      res.emplace_back(i, exponent);
    }
  }
  if (n > 1) {
    res.emplace_back(n, 1);
  }
  return res;
}
\end{lstlisting}
The Miller-Rabin test writes $n - 1 = d \cdot 2^r$ and repeatedly squares $a^d$; a base that breaks the pattern required of a prime proves that $n$ is composite.

\begin{compactitem}
  \item \inlinecode{is\_probable\_prime(n, k = 20)} returns whether \inlinecode{n} is probably prime using \inlinecode{k} random Miller-Rabin bases. The result is guaranteed correct when \inlinecode{n} is prime. For composite \inlinecode{n}, a true result has one-sided error probability at most $(1 / 4)^k$.
  \item \inlinecode{is\_prime(n)} returns whether the signed 64-bit integer \inlinecode{n} is prime using deterministic Miller-Rabin bases sufficient up to and including $2^{63} - 1$.
\end{compactitem}

Modular multiplication uses \inlinecode{\_\_uint128\_t} when available, with a slower portable double-and-add fallback that avoids overflow on compilers without 128-bit integers.

\textbf{Time Complexity}: $O(k \cdot \log n)$ modular multiplications per call to \inlinecode{is\_probable\_prime()} and $O(\log n)$ per call to \inlinecode{is\_prime()}. The portable multiplication fallback adds another $O(\log n)$ factor.

\textbf{Space Complexity}: $O(1)$ auxiliary.

\begin{lstlisting}[style=cpp]
uint64_t mulmod(uint64_t a, uint64_t b, uint64_t m) {
  assert(m > 0);
#if defined(__SIZEOF_INT128__)
  return static_cast<uint64_t>(static_cast<__uint128_t>(a) * b % m);
#else
  uint64_t res = 0, cur = a % m;
  for (; b > 0; b >>= 1) {
    if (b & 1) {
      res = res >= m - cur ? res - (m - cur) : res + cur;
    }
    cur = cur >= m - cur ? cur - (m - cur) : cur + cur;
  }
  return res;
#endif
}

uint64_t powmod(uint64_t x, uint64_t n, uint64_t m) {
  assert(m > 0);
  uint64_t res = 1 % m;
  for (; n > 0; n >>= 1) {
    if (n & 1) {
      res = mulmod(res, x, m);
    }
    x = mulmod(x, x, m);
  }
  return res;
}

uint64_t rand64u() {
  static std::mt19937_64 rng(std::chrono::steady_clock::now().time_since_epoch().count());
  return rng();
}

bool is_probable_prime(int64_t n, int k = 20) {
  if (n == 2 || n == 3) {
    return true;
  }
  if (n < 2 || n % 2 == 0 || n % 3 == 0) {
    return false;
  }
  uint64_t s = n - 1, p = n - 1;
  while (!(s & 1)) {
    s >>= 1;
  }
  for (int i = 0; i < k; i++) {
    uint64_t x, r = powmod(rand64u() % (n - 3) + 2, s, n);  // random witness base in [2, n - 2]
    for (x = s; x != p && r != 1 && r != p; x <<= 1) {
      r = mulmod(r, r, n);
    }
    if (r != p && !(x & 1)) {
      return false;
    }
  }
  return true;
}

bool is_prime(int64_t n) {
  if (n < 2) {
    return false;
  }
  static const int small_primes[] = {2, 3, 5, 7, 11, 13, 17, 19, 23, 29};
  for (int p : small_primes) {
    if (n % p == 0) {
      return n == p;
    }
  }
  if (n < 31 * 31) {
    return true;
  }
  uint64_t t;
  int s = 0;
  for (t = n - 1; !(t & 1); t >>= 1) {
    s++;
  }
  static const uint64_t bases[] = {2, 325, 9375, 28178, 450775, 9780504, 1795265022};
  for (uint64_t a : bases) {
    if (a % n == 0) {
      continue;
    }
    uint64_t r = powmod(a, t, n);
    if (r == 1) {
      continue;
    }
    bool ok = false;
    for (int j = 0; j < s && !ok; j++) {
      ok |= (r == static_cast<uint64_t>(n) - 1);
      r = mulmod(r, r, n);
    }
    if (!ok) {
      return false;
    }
  }
  return true;
}
\end{lstlisting}
Pollard's rho iterates a pseudorandom map modulo $n$; two iterates that collide modulo a hidden prime factor reveal that factor through a GCD with their difference. The complete factorization routine first removes small prime factors with a cached linear sieve, then uses Miller-Rabin and Pollard's rho on the remaining cofactor. Factorizations remain sorted and compressed as (\inlinecode{prime}, \inlinecode{exponent}) pairs.

\begin{compactitem}
  \item \inlinecode{rho\_factor(n)} returns a factor of \inlinecode{n} that is not necessarily prime using Pollard's rho with Brent's optimization. If \inlinecode{n} is prime, then \inlinecode{n} itself is returned. The algorithm may need to be retried to find a nontrivial factor, as done in \inlinecode{factorize\_rho()}.
  \item \inlinecode{factorize\_rho(n)} returns the prime factorization of a 64-bit integer using Miller-Rabin and Pollard's rho without initial trial division.
  \item \inlinecode{factorize(n, small\_prime\_limit = 1000000)} returns the prime factorization of a 64-bit integer \inlinecode{n}, first testing primes through \inlinecode{small\_prime\_limit} and then falling back to Pollard's rho. It supports integers up to and including $2^{63} - 1$.
  \item \inlinecode{divisors\_from\_factors(factors)} returns all divisors from a compressed factorization.
  \item \inlinecode{divisors(n)} returns all divisors of a 64-bit integer using \inlinecode{factorize()} followed by \inlinecode{divisors\_from\_factors()}.
\end{compactitem}

\Needspace*{3\baselineskip}
\textbf{Time Complexity}:
\begin{compactitem}
  \item $O(\sqrt{p})$ expected modular multiplications per call to \inlinecode{rho\_factor(n)}, where $p$ is the smallest prime factor of \inlinecode{n}. For composite \inlinecode{n}, this is at most $O(n^{1/4})$.
  \item $O(n^{1/4})$ expected per call to \inlinecode{factorize\_rho()}. \inlinecode{factorize()} additionally performs $O(L/\log L)$ trial divisions, where $L$ is \inlinecode{small\_prime\_limit}, and takes $O(L)$ time if the cache is rebuilt.
  \item $O(d \log d)$ per call to \inlinecode{divisors\_from\_factors()}, where $d$ is the number of divisors generated.
  \item $O(L/\log L + n^{1/4} + d \log d)$ expected per call to \inlinecode{divisors()}, plus $O(L)$ if the cache is rebuilt.
\end{compactitem}

\Needspace*{3\baselineskip}
\textbf{Space Complexity}:
\begin{compactitem}
  \item $O(c)$ cached space after sieving through $c$.
  \item $O(f + \log n)$ auxiliary for recursive factorization, where $f$ is the number of compressed prime factors returned.
  \item $O(d)$ output space and $O(\log d)$ auxiliary stack space for \inlinecode{divisors\_from\_factors()} and \inlinecode{divisors()}.
\end{compactitem}

\begin{lstlisting}[style=cpp]
using Factors = std::vector<std::pair<int64_t, int>>;

int64_t rho_factor(int64_t n) {
  if (n % 2 == 0) {
    return 2;
  }
  uint64_t y = rand64u() % (n - 1) + 1;
  uint64_t c = rand64u() % (n - 1) + 1;
  uint64_t m = rand64u() % (n - 1) + 1;
  uint64_t g = 1, r = 1, q = 1, ys = 0, x = 0;
  for (r = 1; g == 1; r <<= 1) {
    x = y;
    for (uint64_t i = 0; i < r; i++) {
      y = (mulmod(y, y, n) + c) % n;
    }
    for (uint64_t k = 0; k < r && g == 1; k += m) {
      ys = y;
      int64_t lim = std::min(m, r - k);
      for (int j = 0; j < lim; j++) {
        y = (mulmod(y, y, n) + c) % n;
        q = mulmod(q, (x > y) ? (x - y) : (y - x), n);
      }
      g = std::gcd(q, n);
    }
  }
  if (g == static_cast<uint64_t>(n)) {
    do {
      ys = (mulmod(ys, ys, n) + c) % n;
      g = std::gcd((x > ys) ? (x - ys) : (ys - x), n);
    } while (g <= 1);
  }
  return g;
}

Factors factorize_rho(int64_t n) {
  std::vector<int64_t> factors;
  auto collect = [&](auto &&collect, int64_t value) {
    if (value <= 1) {
      return;
    }
    if (is_prime(value)) {
      factors.push_back(value);
      return;
    }
    int64_t p;
    do {
      p = rho_factor(value);
    } while (p == value);
    collect(collect, p);
    collect(collect, value / p);
  };
  collect(collect, n);
  std::sort(factors.begin(), factors.end());
  Factors res;
  for (int64_t p : factors) {
    if (res.empty() || res.back().first != p) {
      res.emplace_back(p, 1);
    } else {
      res.back().second++;
    }
  }
  return res;
}

static const auto &cached_sieve(int n) {
  struct Cache {
    int limit = 1;
    std::vector<int> least{0, 1}, primes;
  };
  static Cache cache;
  if (n <= cache.limit) {
    return cache;
  }
  cache.limit = std::max(1, n);
  cache.least.assign(cache.limit + 1, 0);
  cache.primes.clear();
  for (int i = 2; i <= cache.limit; i++) {
    if (cache.least[i] == 0) {
      cache.least[i] = i;
      cache.primes.push_back(i);
    }
    for (int p : cache.primes) {
      if (p > cache.least[i] || i > cache.limit / p) {
        break;
      }
      cache.least[i * p] = p;
    }
  }
  return cache;
}

Factors factorize(int64_t n, int small_prime_limit = 1000000) {
  assert(small_prime_limit >= 1);
  if (n <= 1) {
    return {};
  }
  const auto &sieve = cached_sieve(small_prime_limit);
  if (n <= small_prime_limit) {
    Factors res;
    while (n > 1) {
      int p = sieve.least[n], cnt = 0;
      do {
        n /= p;
        cnt++;
      } while (n > 1 && sieve.least[n] == p);
      res.emplace_back(p, cnt);
    }
    return res;
  }
  Factors res;
  const std::vector<int> &primes = sieve.primes;
  for (int p : primes) {
    if (p > small_prime_limit || static_cast<int64_t>(p) > n / p) {
      break;
    }
    if (n % p == 0) {
      int cnt = 0;
      do {
        n /= p;
        cnt++;
      } while (n % p == 0);
      res.emplace_back(p, cnt);
    }
  }
  Factors tail = factorize_rho(n);
  res.insert(res.end(), tail.begin(), tail.end());
  return res;
}

std::vector<int64_t> divisors_from_factors(const Factors &factors) {
  std::vector<int64_t> res{1};
  for (const auto &factor : factors) {
    int old_size = static_cast<int>(res.size());
    int64_t power = 1;
    for (int e = 1; e <= factor.second; e++) {
      power *= factor.first;
      for (int i = 0; i < old_size; i++) {
        res.push_back(res[i] * power);
      }
    }
  }
  std::sort(res.begin(), res.end());
  return res;
}

std::vector<int64_t> divisors(int64_t n) {
  if (n <= 1) {
    return (n < 1) ? std::vector<int64_t>() : std::vector<int64_t>(1, 1);
  }
  return divisors_from_factors(factorize(n));
}
\end{lstlisting}
\Needspace*{4\baselineskip}
\textbf{Example Usage}\par
\begin{lstlisting}[style=cpp]
#include <ctime>
#include <iomanip>
#include <iostream>
#include <set>
using namespace std;

void validate(int64_t n, const Factors &factors) {
  if (n <= 1) {
    assert(factors.empty());
    return;
  }
  if (is_prime(n)) {
    assert((factors == Factors{{n, 1}}));
    return;
  }
  int64_t prod = 1;
  for (const auto &factor : factors) {
    assert(is_prime(factor.first));
    for (int i = 0; i < factor.second; i++) {
      prod *= factor.first;
    }
  }
  assert(prod == n);
}

int main() {
  {  // Small primality tests.
    vector<pair<int, bool>> cases{
        {-10, false}, {-1, false}, {0, false}, {1, false}, // non-primes
        {2, true}, {3, true}, {4, false}, {5, true},
        {9, false},    // divisible by 3
        {25, false},   // square of small prime
        {29, true},    // last small-prime table entry
        {31, true},    // just past small-prime table
        {49, false},   // square, not divisible by 2 or 3
        {121, false},  // 11^2
        {169, false},  // 13^2
        {961, false},  // 31^2, boundary for n < 31*31 shortcut
        {997, true},   // prime near 1000
        {1001, false}, // 7*11*13
        {2047, false}, // 23*89, pseudoprime to base 2
        {341, false},  // 11*31, Fermat pseudoprime to base 2
        {561, false},  // Carmichael number
    };
    for (const auto &[n, expected] : cases) {
      assert(is_prime_slow(n) == expected);
      assert(is_probable_prime(n) == expected);
      assert(is_prime(n) == expected);
    }
  }
  {  // Primality test benchmark.
    vector<int64_t> nums{
        772023803LL, 2147483647LL, 5705234089LL, 6339503641LL, 999966000289LL,
    };
    cout << "Primality test:" << endl;
    cout << setw(15) << left << "num:";
    cout << setw(18) << left << "is_prime_slow():";
    cout << setw(20) << left << "is_probable_prime():";
    cout << setw(13) << right << "is_prime():" << endl;
    cout << fixed << setprecision(3);
    for (int64_t n : nums) {
      clock_t start = clock();
      bool p1 = is_prime_slow(n);
      double t1 = static_cast<double>(clock() - start) / CLOCKS_PER_SEC;
      start = clock();
      bool p2 = is_probable_prime(n);
      double t2 = static_cast<double>(clock() - start) / CLOCKS_PER_SEC;
      start = clock();
      bool p3 = is_prime(n);
      double t3 = static_cast<double>(clock() - start) / CLOCKS_PER_SEC;
      assert(p1 == p2 && p1 == p3);
      cout << setw(12) << left << n << right;
      cout << setw(14) << 1000 * t1 << "ms";
      cout << setw(16) << 1000 * t2 << "ms";
      cout << setw(16) << 1000 * t3 << "ms" << endl;
    }
    cout << endl;
  }
  {  // Small factorization tests.
    for (int64_t i = 1; i <= 10000; i++) {
      auto v1 = factorize_slow(i);
      auto v2 = factorize(i);
      validate(i, v1);
      assert(v1 == v2);
      auto d = divisors(i);
      set<int> s(d.begin(), d.end());
      assert(d.size() == s.size());
      for (int j = 1; j <= i; j++) {
        if (i % j == 0) {
          assert(s.count(j));
        }
      }
    }
  }
  {  // Compressed factors are convenient for building divisors.
    Factors factors{{2, 3}, {3, 2}, {5, 1}};
    vector<int64_t> divs = divisors_from_factors(factors);
    assert(
        (divs == vector<int64_t>{
                     1,  2,  3,  4,  5,  6,  8,  9,  10, 12,  15,  18,
                     20, 24, 30, 36, 40, 45, 60, 72, 90, 120, 180, 360,
                 })
    );
    assert(factorize(360) == factors);
    assert(divisors(360) == divs);
  }
  {  // Large factorization tests.
    const vector<int64_t> nums{
        (1LL << 62),                    // high power of two
        9223372036854775807LL,          // 2^63 - 1, many small factors
        999983LL * 999983,              // square of a large prime
        1900009LL * 1910009 * 2541547,  // three medium-size prime factors
        50000017LL * 50001037,          // 16-digit semiprime
    };
    cout << "Factorization test:" << endl;
    cout << setw(24) << left << "num:";
    cout << setw(20) << left << "factorize_slow():";
    cout << "factorize():" << endl;
    cout << fixed << setprecision(3);
    for (int64_t n : nums) {
      clock_t start = clock();
      Factors factors1 = factorize_slow(n);
      double t1 = static_cast<double>(clock() - start) / CLOCKS_PER_SEC;
      start = clock();
      Factors factors2 = factorize(n);
      double t2 = static_cast<double>(clock() - start) / CLOCKS_PER_SEC;
      validate(n, factors1);
      validate(n, factors2);
      assert(factors1 == factors2);
      cout << setw(26) << left << n << right << setw(8) << 1000 * t1 << "ms";
      cout << setw(16) << 1000 * t2 << "ms" << endl;
      cout << left;
    }
  }
  return 0;
}
\end{lstlisting}
\begin{exampleoutput}
Primality test:
num:           is_prime_slow():  is_probable_prime():  is_prime():
772023803            0.006ms           0.005ms           0.001ms
2147483647           0.009ms           0.006ms           0.002ms
5705234089           0.014ms           0.000ms           0.001ms
6339503641           0.013ms           0.001ms           0.001ms
999966000289         0.171ms           0.001ms           0.001ms

Factorization test:
num:                    factorize_slow():   factorize():
4611686018427387904          0.002ms           0.002ms
9223372036854775807          0.049ms           0.007ms
999966000289                 0.503ms           0.043ms
9223361212852495307          0.952ms           0.096ms
2500052700017629            24.587ms           0.194ms
\end{exampleoutput}
\subsection{Euler's Totient Function}
\label{sec:6.3.10}
Euler's totient function $\phi(n)$ returns the number of positive integers less than or equal to $n$ that are relatively prime to $n$; that is, the number of integers $k$ in the range $[1, n]$ for which $\gcd(n, k) = 1$.

The function is multiplicative, with $\phi(p^k) = p^k - p^{k-1}$ for a prime power, which yields the product formula $\phi(n) = n \prod_{p \mid n} (1 - 1/p)$ taken over the distinct primes $p$ dividing $n$. \inlinecode{phi(n)} applies this formula directly, factoring \inlinecode{n} by trial division. \inlinecode{phi\_table(n)} computes the whole range at once with a sieve of Eratosthenes: starting from \inlinecode{v[i] = i}, it sweeps each prime $p$ across its multiples \inlinecode{j} and folds in the factor $(1 - 1/p)$ by replacing \inlinecode{v[j]} with ${\inlinecodemath{v[j]}} - {\inlinecodemath{v[j]}} / p$. This is faster than the alternative that uses the divisor-sum identity $\sum_{d \mid n} \phi(d) = n$ to subtract each \inlinecode{v[i]} from its multiples, which costs $O(n \log n)$.

The overload \inlinecode{phi\_table(lo, hi)} computes the totients of a high but narrow range with a segmented sieve, analogous to the segmented prime sieve. It sieves the primes up to $\sqrt{{\inlinecodemath{hi}}}$, applies each one's $(1 - 1/p)$ factor to its multiples within $[{\inlinecodemath{lo}}, {\inlinecodemath{hi}}]$ while dividing that prime out of a parallel copy of every value, and finally folds in any single leftover prime factor greater than $\sqrt{{\inlinecodemath{hi}}}$ (a number at most \inlinecode{hi} can have at most one such factor).

\begin{compactitem}
  \item \inlinecode{phi(n)} returns Euler's totient of \inlinecode{n}.
  \item \inlinecode{phi\_table(n)} returns a vector \inlinecode{v} of length \inlinecode{n + 1} such that \inlinecode{v[i]} stores \inlinecode{phi(i)} for every \inlinecode{i} in the range $[0, {\inlinecodemath{n}}]$.
  \item \inlinecode{phi\_table(lo, hi)} returns a vector \inlinecode{v} of length ${\inlinecodemath{hi}} - {\inlinecodemath{lo}} + 1$ such that \inlinecode{v[i - lo]} stores \inlinecode{phi(i)} for every \inlinecode{i} in the range $[{\inlinecodemath{lo}}, {\inlinecodemath{hi}}]$. This overload assumes \inlinecode{lo} $\geq 0$; if \inlinecode{hi} \textless{} \inlinecode{lo}, it returns an empty vector.
\end{compactitem}

\Needspace*{3\baselineskip}
\textbf{Time Complexity}:
\begin{compactitem}
  \item $O(\sqrt{n})$ per call to \inlinecode{phi(n)}.
  \item $O(n \log \log n)$ per call to \inlinecode{phi\_table(n)}.
  \item $O((w + \sqrt{{\inlinecodemath{hi}}}) \cdot \log\left(\log\left({\inlinecodemath{hi}}\right)\right))$ per call to \inlinecode{phi\_table(lo, hi)}, where $w = {\inlinecodemath{hi}} - {\inlinecodemath{lo}} + 1$.
\end{compactitem}

\Needspace*{3\baselineskip}
\textbf{Space Complexity}:
\begin{compactitem}
  \item $O(1)$ auxiliary for \inlinecode{phi(n)}.
  \item $O(n)$ auxiliary for \inlinecode{phi\_table(n)}.
  \item $O(w + \sqrt{{\inlinecodemath{hi}}})$ auxiliary for \inlinecode{phi\_table(lo, hi)}.
\end{compactitem}

\begin{lstlisting}[style=cpp]
#include <algorithm>
#include <cmath>
#include <cstdint>
#include <numeric>
#include <vector>

int phi(int n) {
  int res = n;
  for (int i = 2; i <= n / i; i++) {
    if (n % i == 0) {
      while (n % i == 0) {
        n /= i;
      }
      res -= res / i;
    }
  }
  if (n > 1) {
    res -= res / n;
  }
  return res;
}

std::vector<int> phi_table(int n) {
  std::vector<int> res(n + 1);
  std::iota(res.begin(), res.end(), 0);
  for (int i = 2; i <= n; i++) {
    if (res[i] == i) {  // i is prime, so apply its (1 - 1/i) factor to every multiple.
      for (int j = i; j <= n; j += i) {
        res[j] -= res[j] / i;
      }
    }
  }
  return res;
}

std::vector<int> phi_table(int lo, int hi) {
  if (hi < lo) {
    return {};
  }
  int root = static_cast<int>(std::sqrt(hi));
  while (1LL * (root + 1) * (root + 1) <= hi) {
    root++;
  }
  // Sieve the primes up to sqrt(hi); any larger prime divisor is the one leftover factor.
  std::vector<char> composite(root + 1);
  std::vector<int> primes;
  for (int i = 2; i <= root; i++) {
    if (!composite[i]) {
      primes.push_back(i);
      for (int64_t j = 1LL * i * i; j <= root; j += i) {
        composite[j] = true;
      }
    }
  }
  std::vector<int> res(hi - lo + 1);
  std::vector<int64_t> remaining(hi - lo + 1);
  for (int i = lo; i <= hi; i++) {
    res[i - lo] = i;
    remaining[i - lo] = i;
  }
  for (int p : primes) {
    // Start at the first multiple of p at least lo, but never below p itself (this skips 0).
    int64_t start = std::max(1LL * p, ((1LL * lo + p - 1) / p) * p);
    for (int64_t j = start; j <= hi; j += p) {
      int idx = static_cast<int>(j - lo);
      res[idx] -= res[idx] / p;
      while (remaining[idx] % p == 0) {
        remaining[idx] /= p;
      }
    }
  }
  for (int i = lo; i <= hi; i++) {
    int idx = i - lo;
    if (remaining[idx] > 1) {  // a single prime factor greater than sqrt(hi) is left over
      res[idx] -= static_cast<int>(res[idx] / remaining[idx]);
    }
  }
  return res;
}
\end{lstlisting}
\Needspace*{4\baselineskip}
\textbf{Example Usage}\par
\begin{lstlisting}[style=cpp]
#include <cassert>
using namespace std;

int main() {
  assert(phi(1) == 1);
  assert(phi(9) == 6);
  assert(phi(1234567) == 1224720);
  const int n = 1000;
  vector<int> pt = phi_table(n);
  for (int i = 0; i <= n; i++) {
    assert(pt[i] == phi(i));
  }
  // The segmented overload agrees with phi(), including over a low range (with 0 and 1) and a
  // high, narrow window where many values have a prime factor exceeding sqrt(hi).
  vector<int> lphi = phi_table(0, 50);
  for (int i = 0; i <= 50; i++) {
    assert(lphi[i] == phi(i));
  }
  int lo = 1000000000, hi = 1000000100;
  vector<int> hphi = phi_table(lo, hi);
  for (int i = lo; i <= hi; i++) {
    assert(hphi[i - lo] == phi(i));
  }
  return 0;
}
\end{lstlisting}
\subsection{Divisor Functions}
\label{sec:6.3.11}
The divisor-counting function $d(n)$ counts the positive divisors of $n$, and the divisor-sum function $\sigma(n)$ adds them up. Both are multiplicative, so each is determined by the prime factorization $n = \prod p_i^{e_i}$: choosing an exponent in $[0, e_i]$ for every prime gives $d(n) = \prod (e_i + 1)$, and summing the resulting geometric series for each prime gives $\sigma(n) = \prod (1 + p_i + \dots + p_i^{e_i})$. A single value is therefore computed by trial division, factoring $n$ as the divisors are accumulated.

\begin{compactitem}
  \item \inlinecode{divisor\_count(n)} returns $d(n)$, the number of positive divisors of \inlinecode{n}, which must be positive.
  \item \inlinecode{divisor\_sum(n)} returns $\sigma(n)$, the sum of the positive divisors of \inlinecode{n}, which must be positive.
\end{compactitem}

For every value in a range, factoring each number separately is wasteful. Instead, sweep over each candidate divisor $d$ and add its contribution to every multiple of $d$. The total work is the harmonic sum $n/1 + n/2 + \dots + n/n$, which is $O(n \log n)$. The same sweep computes any divisor-indexed accumulation, such as the divisor sums used by the Mobius inversion of section \secref{6.3.12}.

\begin{compactitem}
  \item \inlinecode{divisor\_count\_table(n)} returns a vector \inlinecode{v} of length \inlinecode{n + 1} such that \inlinecode{v[i]} is the number of divisors of \inlinecode{i} for every \inlinecode{i} in $[1, {\inlinecodemath{n}}]$, with \inlinecode{v[0]} left as $0$.
  \item \inlinecode{divisor\_sum\_table(n)} returns a vector \inlinecode{v} of length \inlinecode{n + 1} such that \inlinecode{v[i]} is the sum of the divisors of \inlinecode{i} for every \inlinecode{i} in $[1, {\inlinecodemath{n}}]$, with \inlinecode{v[0]} left as $0$.
\end{compactitem}

The aliquot sum $\sigma(n) - n$ adds the proper divisors, marking \inlinecode{n} perfect when it equals \inlinecode{n} and pairing amicable numbers when each is the other's. For an \inlinecode{n} too large for trial division, factor it with Pollard's rho in section \secref{6.3.9} and apply the product formulas; to tabulate a high narrow window, sieve it with the primes up to $\sqrt{{\inlinecodemath{hi}}}$ as in section \secref{6.3.8}.

Overflow warning: $\sigma(n)$ can be several times larger than \inlinecode{n}, so \inlinecode{divisor\_sum()} overflows \inlinecode{int64\_t} somewhat before \inlinecode{n} itself does. Trial division binds first in practice, since $\sqrt{n}$ divisions are already too slow well below that point; factor with section \secref{6.3.9} instead.

\Needspace*{3\baselineskip}
\textbf{Time Complexity}:
\begin{compactitem}
  \item $O(\sqrt{n})$ per call to \inlinecode{divisor\_count(n)} and \inlinecode{divisor\_sum(n)}.
  \item $O(n \log n)$ per call to \inlinecode{divisor\_count\_table(n)} and \inlinecode{divisor\_sum\_table(n)}.
\end{compactitem}

\Needspace*{3\baselineskip}
\textbf{Space Complexity}:
\begin{compactitem}
  \item $O(1)$ auxiliary for \inlinecode{divisor\_count()} and \inlinecode{divisor\_sum()}.
  \item $O(n)$ for the vectors returned by \inlinecode{divisor\_count\_table(n)} and \inlinecode{divisor\_sum\_table(n)}.
\end{compactitem}

\begin{lstlisting}[style=cpp]
#include <cassert>
#include <cstdint>
#include <vector>

int64_t divisor_count(int64_t n) {
  assert(n > 0);
  int64_t res = 1;
  for (int64_t p = 2; p <= n / p; p++) {
    if (n % p == 0) {
      int e = 0;
      for (; n % p == 0; n /= p) {
        e++;
      }
      res *= e + 1;
    }
  }
  if (n > 1) {  // A prime cofactor larger than the square root is left over.
    res *= 2;
  }
  return res;
}

int64_t divisor_sum(int64_t n) {
  assert(n > 0);
  int64_t res = 1;
  for (int64_t p = 2; p <= n / p; p++) {
    if (n % p == 0) {
      int64_t power = 1, series = 1;
      for (; n % p == 0; n /= p) {
        power *= p;
        series += power;
      }
      res *= series;  // Overflow warning.
    }
  }
  if (n > 1) {
    res *= n + 1;
  }
  return res;
}

std::vector<int> divisor_count_table(int n) {
  assert(n >= 0);
  std::vector<int> res(n + 1);
  for (int d = 1; d <= n; d++) {
    for (int i = d; i <= n; i += d) {
      res[i]++;
    }
  }
  return res;
}

std::vector<int64_t> divisor_sum_table(int n) {
  assert(n >= 0);
  std::vector<int64_t> res(n + 1);
  for (int d = 1; d <= n; d++) {
    for (int i = d; i <= n; i += d) {
      res[i] += d;
    }
  }
  return res;
}
\end{lstlisting}
\Needspace*{4\baselineskip}
\textbf{Example Usage}\par
\begin{lstlisting}[style=cpp]
#include <cassert>

int main() {
  assert(divisor_count(1) == 1);
  assert(divisor_count(12) == 6);  // 1, 2, 3, 4, 6, and 12.
  assert(divisor_count(999999937) == 2);
  assert(divisor_sum(1) == 1);
  assert(divisor_sum(12) == 28);
  assert(divisor_sum(6) == 12);  // 6 is perfect, since its aliquot sum is 6.

  // 220 and 284 are the smallest amicable pair.
  assert(divisor_sum(220) - 220 == 284);
  assert(divisor_sum(284) - 284 == 220);

  auto counts = divisor_count_table(100);
  auto sums = divisor_sum_table(100);
  assert(counts[0] == 0 && sums[0] == 0);
  for (int i = 1; i <= 100; i++) {
    assert(counts[i] == divisor_count(i));
    assert(sums[i] == divisor_sum(i));
  }
  return 0;
}
\end{lstlisting}
\subsection{Mobius Function and Inclusion-Exclusion}
\label{sec:6.3.12}
The Mobius function $\mu(n)$ is the multiplicative function defined as $1$ if $n = 1$, $0$ if $n$ is divisible by the square of a prime, and $(-1)^k$ if $n$ is a product of $k$ distinct primes. It is the coefficient that appears in Mobius inversion: if $g(n) = \sum_{d \mid n} f(d)$ for all $n$, then $f(n) = \sum_{d \mid n} \mu(d) \, g(n / d)$. This formalizes the principle of inclusion-exclusion in the divisor lattice, letting one recover a summand function from its divisor sums.

Inclusion-exclusion itself counts the size of a union by alternately adding and subtracting the sizes of intersections. A canonical number-theoretic instance is counting integers in $[1, n]$ that are coprime to a modulus $m$: subtract the multiples of each prime factor of $m$, add back those counted twice, and so on. The signed terms are precisely the values of $\mu$ over the squarefree divisors of $m$, so \inlinecode{count\_coprime(n, m)} equals $\sum_{d \mid m} \mu(d) \lfloor n / d \rfloor$.

\begin{compactitem}
  \item \inlinecode{mobius(n)} returns $\mu(n)$ for a single positive \inlinecode{n} by trial division.
  \item \inlinecode{mobius\_sieve(n)} returns a vector of length $n + 1$ where index $i$ holds $\mu(i)$, computed together with a linear sieve, where \inlinecode{n} must be nonnegative.
  \item \inlinecode{count\_coprime(n, m)} returns the number of integers in $[1, {\inlinecodemath{n}}]$ that are coprime to \inlinecode{m}, via inclusion-exclusion over the distinct prime factors of \inlinecode{m}, where \inlinecode{n} must be nonnegative and \inlinecode{m} must be positive.
\end{compactitem}

\Needspace*{3\baselineskip}
\textbf{Time Complexity}:
\begin{compactitem}
  \item $O(\sqrt{n})$ per call to \inlinecode{mobius()}.
  \item $O(n)$ per call to \inlinecode{mobius\_sieve()}.
  \item $O(\sqrt{m} + 2^{k} \cdot k)$ per call to \inlinecode{count\_coprime()}, where $k$ is the number of distinct prime factors of $m$ (at most 15 for $m$ within 64-bit range).
\end{compactitem}

\Needspace*{3\baselineskip}
\textbf{Space Complexity}:
\begin{compactitem}
  \item $O(n)$ auxiliary for \inlinecode{mobius\_sieve()}.
  \item $O(k)$ auxiliary for \inlinecode{count\_coprime()}.
  \item $O(1)$ auxiliary for \inlinecode{mobius()}.
\end{compactitem}

\begin{lstlisting}[style=cpp]
#include <cassert>
#include <cstdint>
#include <vector>

int mobius(int n) {
  assert(n > 0);
  if (n == 1) {
    return 1;
  }
  int res = 1;
  for (int p = 2; p <= n / p; p++) {
    if (n % p == 0) {
      n /= p;
      if (n % p == 0) {
        return 0;  // p^2 divides the original n.
      }
      res = -res;
    }
  }
  if (n > 1) {
    res = -res;  // One remaining prime factor larger than its square root.
  }
  return res;
}

std::vector<int> mobius_sieve(int n) {
  assert(n >= 0);
  std::vector<int> mu(n + 1), primes;
  std::vector<char> composite(n + 1);
  if (n >= 1) {
    mu[1] = 1;
  }
  for (int i = 2; i <= n; i++) {
    if (!composite[i]) {
      primes.push_back(i);
      mu[i] = -1;
    }
    for (int p : primes) {
      if (p > n / i) {
        break;
      }
      composite[i * p] = true;
      if (i % p == 0) {
        mu[i * p] = 0;  // p^2 divides i * p.
        break;
      }
      mu[i * p] = -mu[i];
    }
  }
  return mu;
}

int64_t count_coprime(int64_t n, int64_t m) {
  assert(n >= 0 && m > 0);
  std::vector<int64_t> primes;
  for (int64_t p = 2; p <= m / p; p++) {
    if (m % p == 0) {
      primes.push_back(p);
      while (m % p == 0) {
        m /= p;
      }
    }
  }
  if (m > 1) {
    primes.push_back(m);
  }
  int k = static_cast<int>(primes.size());
  int64_t total = 0;
  for (int mask = 0; mask < (1 << k); mask++) {
    int64_t product = 1;
    int bits = 0;
    for (int i = 0; i < k; i++) {
      if ((mask >> i) & 1) {
        product *= primes[i];
        bits++;
      }
    }
    total += (bits & 1 ? -1 : 1) * (n / product);
  }
  return total;
}
\end{lstlisting}
\Needspace*{4\baselineskip}
\textbf{Example Usage}\par
\begin{lstlisting}[style=cpp]
#include <cassert>
using namespace std;

int main() {
  assert(mobius(1) == 1);
  assert(mobius(2) == -1);
  assert(mobius(6) == 1);    // 2 * 3, two distinct primes.
  assert(mobius(12) == 0);   // Divisible by 2^2.
  assert(mobius(30) == -1);  // 2 * 3 * 5, three distinct primes.

  assert((mobius_sieve(12) == vector<int>{0, 1, -1, -1, 0, -1, 1, -1, 0, 0, 1, -1, 0}));

  assert(count_coprime(10, 6) == 3);  // 1, 5, 7
  assert(count_coprime(100, 1) == 100);
  return 0;
}
\end{lstlisting}

\section{\texorpdfstring{Arbitrary Precision Arithmetic}{6.4 Arbitrary Precision Arithmetic}}
\setcounter{section}{4}
\setcounter{subsection}{0}
\subsection{Big Integer (Simple)}
\label{sec:6.4.1}
Perform simple arithmetic operations on arbitrary precision big integers whose digits are internally represented as an \inlinecode{std::string} in little-endian order.

This version is intentionally small and decimal-oriented. Addition and subtraction are the usual schoolbook digit scans with carry or borrow. Multiplication keeps a shifted copy of the left operand for each digit of the right operand, adds that row once per digit value, then shifts by one decimal place. Division is long division: it scans the dividend from most significant digit to least, maintains a running remainder, and repeatedly subtracts the divisor to discover each quotient digit.

\begin{compactitem}
  \item \inlinecode{BigInt(n = 0)} constructs a big integer from an integer \inlinecode{n}.
  \item \inlinecode{BigInt(s)} constructs a big integer from a string \inlinecode{s}, which must strictly consist of a sequence of numeric digits, optionally preceded by a minus sign.
  \item \inlinecode{to\_string()} returns the string representation of the big integer.
  \item \inlinecode{comp(a, b)} returns $-1$, $0$, or $1$ depending on whether the big integers \inlinecode{a} and \inlinecode{b} compare less, equal, or greater, respectively.
  \item \inlinecode{add(a, b)} returns the sum of big integers \inlinecode{a} and \inlinecode{b}.
  \item \inlinecode{sub(a, b)} returns the difference of big integers \inlinecode{a} and \inlinecode{b}.
  \item \inlinecode{mul(a, b)} returns the product of big integers \inlinecode{a} and \inlinecode{b}.
  \item \inlinecode{div(a, b)} returns the quotient of big integers \inlinecode{a} and \inlinecode{b}.
\end{compactitem}

\Needspace*{3\baselineskip}
\textbf{Time Complexity}:
\begin{compactitem}
  \item $O(n)$ per call to the constructor, \inlinecode{to\_string()}, \inlinecode{comp()}, \inlinecode{add()}, and \inlinecode{sub()}, where $n$ is the total number of digits in the arguments and result for each operation.
  \item $O(n \cdot m)$ per call to \inlinecode{mul(a, b)} and \inlinecode{div(a, b)} where $n$ is the number of digits in \inlinecode{a} and $m$ is the number of digits in \inlinecode{b}.
\end{compactitem}

\Needspace*{3\baselineskip}
\textbf{Space Complexity}:
\begin{compactitem}
  \item $O(n)$ for storage of the big integer, where $n$ is the number of digits.
  \item $O(n)$ auxiliary for \inlinecode{to\_string()}, \inlinecode{add()}, \inlinecode{sub()}, \inlinecode{mul()}, and \inlinecode{div()}, where $n$ is the total number of digits in the arguments and result for each operation.
\end{compactitem}

\begin{lstlisting}[style=cpp]
#include <algorithm>
#include <cassert>
#include <cctype>
#include <cstddef>
#include <cstdint>
#include <stdexcept>
#include <string>

class BigInt {
  std::string digits;
  int sign;

  void normalize() {
    std::size_t pos = digits.find_last_not_of('0');
    if (pos != std::string::npos) {
      digits.erase(pos + 1);
    } else {
      digits = "0";  // An all-zero (or empty) magnitude collapses to a single "0".
    }
    if (digits.size() == 1 && digits[0] == '0') {
      sign = 1;
    }
  }

  static int comp(const std::string &a, const std::string &b, int asign, int bsign) {
    if (asign != bsign) {
      return asign < bsign ? -1 : 1;
    }
    if (a.size() != b.size()) {
      return a.size() < b.size() ? -asign : asign;
    }
    for (int i = static_cast<int>(a.size()) - 1; i >= 0; i--) {
      if (a[i] != b[i]) {
        return a[i] < b[i] ? -asign : asign;
      }
    }
    return 0;
  }

  static BigInt add(const std::string &a, const std::string &b, int asign, int bsign) {
    if (asign != bsign) {
      return (asign == 1) ? sub(a, b, asign, 1) : sub(b, a, bsign, 1);
    }
    BigInt res;
    res.sign = asign;
    res.digits.resize(std::max(a.size(), b.size()) + 1, '0');
    for (int i = 0, carry = 0; i < static_cast<int>(res.digits.size()); i++) {
      int d = carry;
      if (i < static_cast<int>(a.size())) {
        d += a[i] - '0';
      }
      if (i < static_cast<int>(b.size())) {
        d += b[i] - '0';
      }
      res.digits[i] = '0' + (d % 10);
      carry = d / 10;
    }
    res.normalize();
    return res;
  }

  static BigInt sub(const std::string &a, const std::string &b, int asign, int bsign) {
    if (asign == -1 || bsign == -1) {
      return add(a, b, asign, -bsign);
    }
    BigInt res;
    if (comp(a, b, asign, bsign) < 0) {
      res = sub(b, a, bsign, asign);
      res.sign = -1;
      return res;
    }
    res.digits.assign(a.size(), '0');
    for (int i = 0, borrow = 0; i < static_cast<int>(res.digits.size()); i++) {
      int d = (i < static_cast<int>(b.size()) ? a[i] - b[i] : a[i] - '0') - borrow;
      if (a[i] > '0') {
        borrow = 0;
      }
      if (d < 0) {
        d += 10;
        borrow = 1;
      }
      res.digits[i] = '0' + (d % 10);
    }
    res.normalize();
    return res;
  }

 public:
  BigInt(int64_t n = 0) {
    sign = (n < 0) ? -1 : 1;
    if (n == 0) {
      digits = "0";
      return;
    }
    for (; n != 0; n /= 10) {
      int digit = static_cast<int>(n % 10);
      digits += '0' + (digit < 0 ? -digit : digit);
    }
    normalize();
  }

  BigInt(const std::string &s) {
    if (s.empty() || (s[0] == '-' && s.size() == 1)) {
      throw std::runtime_error("Invalid string format to construct BigInt.");
    }
    digits.assign(s.rbegin(), s.rend());
    if (s[0] == '-') {
      sign = -1;
      digits.erase(digits.size() - 1);
    } else {
      sign = 1;
    }
    if (digits.find_first_not_of("0123456789") != std::string::npos) {
      throw std::runtime_error("Invalid string format to construct BigInt.");
    }
    normalize();
  }

  std::string to_string() const {
    return (sign < 0 ? "-" : "") + std::string(digits.rbegin(), digits.rend());
  }

  friend int comp(const BigInt &a, const BigInt &b) {
    return comp(a.digits, b.digits, a.sign, b.sign);
  }

  friend BigInt add(const BigInt &a, const BigInt &b) {
    return add(a.digits, b.digits, a.sign, b.sign);
  }

  friend BigInt sub(const BigInt &a, const BigInt &b) {
    return sub(a.digits, b.digits, a.sign, b.sign);
  }

  friend BigInt mul(const BigInt &a, const BigInt &b) {
    BigInt res, row(a);
    for (char dc : b.digits) {
      for (int j = 0; j < (dc - '0'); j++) {
        res = add(res.digits, row.digits, res.sign, row.sign);
      }
      if (row.digits.size() > 1 || row.digits[0] != '0') {
        row.digits.insert(0, "0");
      }
    }
    res.sign = a.sign * b.sign;
    res.normalize();
    return res;
  }

  friend BigInt div(const BigInt &a, const BigInt &b) {
    assert(b.digits.size() != 1 || b.digits[0] != '0');
    BigInt res, row;
    res.digits.assign(a.digits.size(), '0');
    row.digits.clear();  // Running remainder; kept normalized and nonnegative below.
    for (int i = static_cast<int>(a.digits.size()) - 1; i >= 0; i--) {
      row.digits.insert(row.digits.begin(), a.digits[i]);
      row.normalize();  // Strip the spurious high-order zero left by the previous remainder.
      // Subtract while the remainder is >= b. A strict ">" undercounts the quotient digit when the
      // remainder equals b, and leaves the remainder too large, overflowing the next digit past 9.
      while (comp(row.digits, b.digits, 1, 1) >= 0) {
        res.digits[i]++;
        row = sub(row.digits, b.digits, 1, 1);
      }
    }
    res.sign = a.sign * b.sign;
    res.normalize();
    return res;
  }
};
\end{lstlisting}
\Needspace*{4\baselineskip}
\textbf{Example Usage}\par
\begin{lstlisting}[style=cpp]
int main() {
  BigInt a("-9899819294989142124"), b("12398124981294214");
  assert(add(a, b).to_string() == "-9887421170007847910");
  assert(sub(a, b).to_string() == "-9912217419970436338");
  assert(mul(a, b).to_string() == "-122739196911503356525379735104870536");
  assert(div(a, b).to_string() == "-798");
  assert(comp(a, b) == -1 && comp(a, a) == 0 && comp(b, a) == 1);
  assert(BigInt(INT64_MIN).to_string() == "-9223372036854775808");
  return 0;
}
\end{lstlisting}
\subsection{Big Integer}
\label{sec:6.4.2}
Perform operations on arbitrary precision big integers internally represented as a vector of base-\inlinecode{BASE} digits in little-endian order. \inlinecode{BASE} defaults to $10^9$ and must be a power of 10 no larger than $10^9$. Multiplication selects schoolbook, Karatsuba, or FFT from the operand size; \inlinecode{MUL\_CUTOFF} is the padded transform length where FFT takes over and normally needs no adjustment. Division likewise splits into a quadratic loop for small divisors and a recursive one past \inlinecode{DIV\_CUTOFF} limbs, so it inherits whichever multiplication the size selects. Typical arithmetic operations involving mixed numeric primitives and strings are supported through implicit construction and hidden friend operators, as long as at least one operand is a \inlinecode{BigInt} at any given level of evaluation.

Addition, subtraction, and comparisons are performed using the standard linear algorithms. Multiplication first converts limbs from base \inlinecode{BASE} to the narrower base \inlinecode{MUL\_BASE} so coefficient products fit safely, then multiplies with schoolbook, Karatsuba, or (if beyond \inlinecode{MUL\_CUTOFF}) complex FFT convolution according to the operand size, converting the result back. Division and modulo are computed together by normalized long division: scale the operands so the divisor's leading limb is large, estimate each quotient limb from the top one or two limbs of the running remainder, correct the estimate by adding the divisor back if necessary, then unscale the remainder. Beyond \inlinecode{DIV\_CUTOFF} limbs, that loop becomes the base case of Burnikel-Ziegler recursion, which halves both operands and recovers each half of the quotient from a division of the same shape, paying one multiplication per level rather than one estimate per limb.

\begin{compactitem}
  \item \inlinecode{BigInt()} constructs zero, and \inlinecode{BigInt(n)} constructs a big integer from an integer \inlinecode{n}.
  \item \inlinecode{BigInt(s)} constructs a big integer from a C string or an \inlinecode{std::string} \inlinecode{s}.
  \item \inlinecode{operator=} is defined to copy from another big integer or to assign from a 64-bit integer primitive.
  \item \inlinecode{size()} returns the number of digits in the base-10 representation.
  \item Operators \inlinecode{\textgreater{}\textgreater{}} and \inlinecode{\textless{}\textless{}} are defined to support stream-based input and output.
  \item \inlinecode{to\_string()}, \inlinecode{to\_llong()}, \inlinecode{to\_double()}, and \inlinecode{to\_ldouble()} return the big integer converted to an \inlinecode{std::string}, \inlinecode{int64\_t}, \inlinecode{double}, and \inlinecode{long double} respectively. For the latter three data types, overflow behavior is based on that of inputting from \inlinecode{std::istream}. Equivalent conversions are also available through explicit casts to \inlinecode{int}, \inlinecode{long long}, \inlinecode{double}, and \inlinecode{long double}.
  \item \inlinecode{to\_arithmetic\textless{}T\textgreater{}()} returns the value converted to arithmetic type \inlinecode{T} by streaming the base-10 representation into it, and is what the three conversions above are defined in terms of.
  \item \inlinecode{abs()} returns the absolute value.
  \item \inlinecode{comp(v)} returns $-1$, $0$, or $1$ depending on whether the big integer compares less than, equal to, or greater than \inlinecode{v}, respectively.
  \item Operators \inlinecode{\textless{}}, \inlinecode{\textgreater{}}, \inlinecode{\textless{}=}, \inlinecode{\textgreater{}=}, \inlinecode{==}, \inlinecode{!=}, \inlinecode{+}, \inlinecode{-}, \inlinecode{*}, \inlinecode{/}, \inlinecode{\%}, \inlinecode{++}, \inlinecode{--}, \inlinecode{+=}, \inlinecode{-=}, \inlinecode{*=}, \inlinecode{/=}, and \inlinecode{\%=} are defined analogous to those on integer primitives.
  \item \inlinecode{div(v)} returns a pair consisting of the quotient and remainder after dividing by \inlinecode{v}.
  \item \inlinecode{pow(n)} returns the big integer raised to the nonnegative power \inlinecode{n} using binary exponentiation.
  \item \inlinecode{mul\_pow10(n)} returns the big integer multiplied by $10^{{\inlinecodemath{n}}}$ for nonnegative \inlinecode{n}, which costs a limb shift rather than a multiplication.
  \item \inlinecode{sqrt()} returns the integral part of the square root using a digit-by-digit algorithm in the internal base.
  \item \inlinecode{nth\_root(n)} returns the integral part of the \inlinecode{n}-th root using Newton's method. Negative inputs require odd \inlinecode{n}.
  \item \inlinecode{rand(d)} returns a random, positive big integer with \inlinecode{d} digits.
\end{compactitem}

\Needspace*{3\baselineskip}
\textbf{Time Complexity}:
\begin{compactitem}
  \item $O(1)$ per call to \inlinecode{size()}.
  \item $O(d)$ per call to the constructors, \inlinecode{to\_string()}, \inlinecode{to\_llong()}, \inlinecode{to\_double()}, \inlinecode{to\_ldouble()}, \inlinecode{to\_arithmetic()}, \inlinecode{abs()}, \inlinecode{comp()}, \inlinecode{rand()}, and all comparison and arithmetic operators except multiplication, division, and modulo, where $d$ is the total number of digits in the arguments and result for each operation.
  \item $O(d \cdot \log\left(d\right) \cdot \log\left(\log\left(d\right)\right))$ per call to multiplication operations once operands reach \inlinecode{MUL\_CUTOFF}, and $O(d^{1.585})$ below it.
  \item $O((d/m) \cdot M(m) \cdot \log m)$ per call to division and modulo operations once the divisor reaches \inlinecode{DIV\_CUTOFF}, and $O(d \cdot m)$ below it, where $d$ and $m$ are the number of digits in the dividend and divisor respectively, and $M(m)$ is the cost of one multiplication of two $m$-digit big integers. Both forms consume the dividend in $d/m$ blocks of the divisor's size; only the per-block cost differs.
  \item $O(M(d) \log n)$ per call to \inlinecode{pow()}, where $n$ is the exponent and $d$ is the maximum digit length reached during exponentiation.
  \item $O(d + n)$ per call to \inlinecode{mul\_pow10()}, where $n$ is the power of ten.
  \item $O(d^{2})$ per call to \inlinecode{sqrt()}.
  \item $O(M(d) \cdot \log\left(d\right) \cdot \log\left(d \cdot n\right))$ per call to \inlinecode{nth\_root()}, where $n$ is the root. Newton's iteration starts within a factor of ten of the answer and doubles the correct digits each step, so it converges in $O(\log d)$ steps, each costing one exponentiation and one division.
\end{compactitem}

\Needspace*{3\baselineskip}
\textbf{Space Complexity}:
\begin{compactitem}
  \item $O(d)$ for storage of the big integer.
  \item $O(d)$ auxiliary for negation, addition, subtraction, multiplication, division, \inlinecode{abs()}, \inlinecode{sqrt()}, \inlinecode{pow()}, and \inlinecode{nth\_root()}.
  \item $O(1)$ auxiliary for all other operations.
\end{compactitem}

\begin{lstlisting}[style=cpp]
#include <algorithm>
#include <cassert>
#include <chrono>
#include <cmath>
#include <complex>
#include <cstdint>
#include <cstring>
#include <iomanip>
#include <istream>
#include <iterator>
#include <ostream>
#include <random>
#include <sstream>
#include <string>
#include <tuple>
#include <utility>
#include <vector>

class BigInt {
  static const int BASE = 1000000000, BASE_DIGITS = 9;  // BASE must equal 10^BASE_DIGITS.
  // Multiplication rebases to half as many decimal digits per limb, so that coefficient products
  // stay exact both in int64_t and within the FFT's double precision.
  static const int MUL_BASE = 10000, MUL_BASE_DIGITS = 4;  // Keep consistent with BASE_DIGITS.
  static const int MUL_CUTOFF = 256;  // Padded MUL_BASE length past which FFT beats Karatsuba.
  static const int DIV_CUTOFF = 64;   // Divisor limb count past which recursion wins.
  static_assert(
      BASE >= 100 && BASE <= 1000000000,
      "BASE must be between 10^2 and 10^9; below that MUL_BASE_DIGITS would round down to zero"
  );
  static_assert(
      MUL_BASE_DIGITS == BASE_DIGITS / 2 && MUL_BASE <= 10000,
      "MUL_BASE must hold half of BASE's digits and stay exact under the FFT"
  );

  using vint = std::vector<int>;
  using vint64 = std::vector<int64_t>;
  using vcd = std::vector<std::complex<double>>;

  vint digits;
  int sign;

  void normalize() {
    while (!digits.empty() && digits.back() == 0) {
      digits.pop_back();
    }
    if (digits.empty()) {
      sign = 1;
    }
  }

  void read(int n, const char *s) {
    sign = 1;
    digits.clear();
    int pos = 0;
    while (pos < n && (s[pos] == '-' || s[pos] == '+')) {
      if (s[pos] == '-') {
        sign = -sign;
      }
      pos++;
    }
    for (int i = n - 1; i >= pos; i -= BASE_DIGITS) {
      int x = 0;
      for (int j = std::max(pos, i - BASE_DIGITS + 1); j <= i; j++) {
        x = x * 10 + s[j] - '0';
      }
      digits.push_back(x);
    }
    normalize();
  }

  static int comp(const vint &a, const vint &b, int asign, int bsign) {
    if (asign != bsign) {
      return asign < bsign ? -1 : 1;
    }
    if (a.size() != b.size()) {
      return a.size() < b.size() ? -asign : asign;
    }
    for (int i = static_cast<int>(a.size()) - 1; i >= 0; i--) {
      if (a[i] != b[i]) {
        return a[i] < b[i] ? -asign : asign;
      }
    }
    return 0;
  }

  static BigInt add(const vint &a, const vint &b, int asign, int bsign) {
    if (asign != bsign) {
      return (asign == 1) ? sub(a, b, asign, 1) : sub(b, a, bsign, 1);
    }
    BigInt res;
    res.digits = a;
    res.sign = asign;
    int carry = 0, size = static_cast<int>(std::max(a.size(), b.size()));
    for (int i = 0; i < size || carry; i++) {
      if (i == static_cast<int>(res.digits.size())) {
        res.digits.push_back(0);
      }
      res.digits[i] += carry + (i < static_cast<int>(b.size()) ? b[i] : 0);
      carry = (res.digits[i] >= BASE) ? 1 : 0;
      if (carry) {
        res.digits[i] -= BASE;
      }
    }
    return res;
  }

  static BigInt sub(const vint &a, const vint &b, int asign, int bsign) {
    if (asign == -1 || bsign == -1) {
      return add(a, b, asign, -bsign);
    }
    BigInt res;
    if (comp(a, b, asign, bsign) < 0) {
      res = sub(b, a, bsign, asign);
      res.sign = -1;
      return res;
    }
    res.digits = a;
    res.sign = asign;
    for (int i = 0, borrow = 0; i < static_cast<int>(a.size()) || borrow; i++) {
      res.digits[i] -= borrow + (i < static_cast<int>(b.size()) ? b[i] : 0);
      borrow = res.digits[i] < 0;
      if (borrow) {
        res.digits[i] += BASE;
      }
    }
    res.normalize();
    return res;
  }

  static vint convert_base(const vint &digits, int l1, int l2) {
    vint64 p(std::max(l1, l2) + 1);
    p[0] = 1;
    for (int i = 1; i < static_cast<int>(p.size()); i++) {
      p[i] = p[i - 1] * 10;
    }
    vint res;
    int64_t curr = 0;
    for (int i = 0, curr_digits = 0; i < static_cast<int>(digits.size()); i++) {
      curr += digits[i] * p[curr_digits];
      curr_digits += l1;
      while (curr_digits >= l2) {
        res.push_back(static_cast<int>((curr % p[l2])));
        curr /= p[l2];
        curr_digits -= l2;
      }
    }
    res.push_back(static_cast<int>(curr));
    while (!res.empty() && res.back() == 0) {
      res.pop_back();
    }
    return res;
  }

  template<typename It>
  static vint64 karatsuba(It alo, It ahi, It blo, It bhi) {
    int n = std::distance(alo, ahi), k = n / 2;
    vint64 res(n * 2);
    if (n <= 32) {
      for (int i = 0; i < n; i++) {
        for (int j = 0; j < n; j++) {
          res[i + j] += alo[i] * blo[j];
        }
      }
      return res;
    }
    auto a1b1 = karatsuba(alo, alo + k, blo, blo + k);
    auto a2b2 = karatsuba(alo + k, ahi, blo + k, bhi);
    vint64 a2(alo + k, ahi), b2(blo + k, bhi);
    for (int i = 0; i < k; i++) {
      a2[i] += alo[i];
      b2[i] += blo[i];
    }
    auto r = karatsuba(a2.begin(), a2.end(), b2.begin(), b2.end());
    for (int i = 0; i < static_cast<int>(a1b1.size()); i++) {
      r[i] -= a1b1[i];
      res[i] += a1b1[i];
    }
    for (int i = 0; i < static_cast<int>(a2b2.size()); i++) {
      r[i] -= a2b2[i];
      res[i + n] += a2b2[i];
    }
    for (int i = 0; i < static_cast<int>(r.size()); i++) {
      res[i + k] += r[i];
    }
    return res;
  }

  template<typename It>
  static vcd fft(It lo, It hi, bool invert = false) {
    int n = std::distance(lo, hi), k = 0, high1 = -1;
    while ((1 << k) < n) {
      k++;
    }
    std::vector<int> rev(n);
    for (int i = 1; i < n; i++) {
      if (!(i & (i - 1))) {
        high1++;
      }
      rev[i] = rev[i ^ (1 << high1)];
      rev[i] |= (1 << (k - high1 - 1));
    }
    vcd roots(n), res(n);
    for (int i = 0; i < n; i++) {
      double alpha = 2 * 3.14159265358979323846 * i / n;
      roots[i] = std::complex<double>(std::cos(alpha), std::sin(alpha));
      res[i] = *(lo + rev[i]);
    }
    for (int len = 1; len < n; len <<= 1) {
      vcd tmp(n);
      int rstep = static_cast<int>(roots.size()) / (len << 1);
      for (int pdest = 0; pdest < n; pdest += len) {
        int p = pdest;
        for (int i = 0; i < len; i++) {
          std::complex<double> c = roots[i * rstep] * res[p + len];
          tmp[pdest] = res[p] + c;
          tmp[pdest + len] = res[p] - c;
          pdest++;
          p++;
        }
      }
      res.swap(tmp);
    }
    if (invert) {
      for (int i = 0; i < static_cast<int>(res.size()); i++) {
        res[i] /= n;
      }
      std::reverse(res.begin() + 1, res.end());
    }
    return res;
  }

  // Schoolbook long division of nonnegative values, quadratic in the limb counts.
  std::pair<BigInt, BigInt> div_schoolbook(const BigInt &v) const {
    int norm = BASE / (v.digits.back() + 1);
    BigInt an = *this * norm, bn = v * norm, q, r;
    q.digits.resize(an.digits.size());
    for (int i = static_cast<int>(an.digits.size()) - 1; i >= 0; i--) {
      r *= BASE;
      r += an.digits[i];
      int s1 = (r.digits.size() <= bn.digits.size()) ? 0 : r.digits[bn.digits.size()];
      int s2 = (r.digits.size() <= bn.digits.size() - 1) ? 0 : r.digits[bn.digits.size() - 1];
      int d = (static_cast<int64_t>(s1) * BASE + s2) / bn.digits.back();
      for (r -= bn * d; r < 0; r += bn) {
        d--;
      }
      q.digits[i] = d;
    }
    q.normalize();
    r.normalize();
    return {q, r / norm};
  }

  // Burnikel-Ziegler recursive division of nonnegative values. Balanced halves replace the
  // quadratic inner loop, paying one multiplication per level instead of one estimate per limb, so
  // division inherits whichever multiplication the operand size selects.
  std::pair<BigInt, BigInt> div_recursive(const BigInt &v) const {
    auto shifted = [](const BigInt &x, int k) {  // Multiplies by BASE^k.
      BigInt res(x);
      if (!res.digits.empty()) {
        res.digits.insert(res.digits.begin(), k, 0);
      }
      return res;
    };
    auto slice = [](const BigInt &x, int lo, int hi) {  // Limbs [lo, hi), or 0 if empty range.
      BigInt res;
      lo = std::max(lo, 0);
      hi = std::min(hi, static_cast<int>(x.digits.size()));
      if (lo < hi) {
        res.digits.assign(x.digits.begin() + lo, x.digits.begin() + hi);
      }
      res.normalize();
      return res;
    };
    // Divides at most 2n limbs by exactly n limbs whose leading limb is at least BASE / 2, where
    // the quotient is known to be below BASE^n. Each half of the quotient comes from one
    // three-by-two block division, which itself divides two half-size blocks by one.
    auto rec = [&](auto &&rec, const BigInt &a, const BigInt &b,
                   int n) -> std::pair<BigInt, BigInt> {
      if (n % 2 != 0 || n < DIV_CUTOFF) {
        return a.div_schoolbook(b);
      }
      int k = n / 2;
      BigInt hi = slice(b, k, 2 * k), lo = slice(b, 0, k);
      auto block = [&](const BigInt &x) {  // At most 3k limbs by b, with quotient below BASE^k.
        BigInt q, r;
        if (slice(x, 2 * k, 3 * k).comp(hi) < 0) {
          std::tie(q, r) = rec(rec, slice(x, k, 3 * k), hi, k);
        } else {
          // The estimate saturates: the leading limbs already reach the divisor, so take the
          // largest quotient the block can hold and recover the matching remainder.
          q = shifted(BigInt(1), k) - 1;
          r = slice(x, k, 3 * k) - hi * q;
        }
        r = shifted(r, k) + slice(x, 0, k) - q * lo;
        while (r < 0) {  // At most twice, since the divisor's leading limb is at least BASE / 2.
          r += b;
          q -= 1;
        }
        return std::make_pair(q, r);
      };
      auto top = block(slice(a, k, 4 * k));
      auto bottom = block(shifted(top.second, k) + slice(a, 0, k));
      return {shifted(top.first, k) + bottom.first, bottom.second};
    };
    // Scale so the divisor's leading limb exceeds BASE / 2, then pad it to a multiple of a power
    // of two so that every recursive split stays even until the blocks fall under DIV_CUTOFF, and
    // consume the dividend one block at a time.
    int scale = BASE / (v.digits.back() + 1);
    BigInt a = *this * scale, b = v * scale;
    int limbs = static_cast<int>(b.digits.size()), block = 1;
    while (block * DIV_CUTOFF < limbs) {
      block *= 2;
    }
    int n = (limbs + block - 1) / block * block, sigma = n - limbs;
    a = shifted(a, sigma);
    b = shifted(b, sigma);
    BigInt q, r;
    for (int i = (static_cast<int>(a.digits.size()) + n - 1) / n - 1; i >= 0; i--) {
      auto step = rec(rec, shifted(r, n) + slice(a, i * n, (i + 1) * n), b, n);
      q = shifted(q, n) + step.first;
      r = step.second;
    }
    // Both scalings scale the remainder and leave the quotient alone.
    return {q, slice(r, sigma, static_cast<int>(r.digits.size())) / scale};
  }

 public:
  BigInt() : sign(1) {}
  BigInt(int v) { *this = static_cast<int64_t>(v); }
  BigInt(int64_t v) { *this = v; }
  BigInt(const char *s) { read(strlen(s), s); }
  BigInt(const std::string &s) { read(s.size(), s.c_str()); }

  BigInt &operator=(int64_t v) {
    sign = v < 0 ? -1 : 1;
    digits.clear();
    for (; v != 0; v /= BASE) {
      int limb = static_cast<int>(v % BASE);
      digits.push_back(limb < 0 ? -limb : limb);
    }
    return *this;
  }

  int size() const {
    return digits.empty() ? 1
                          : static_cast<int>(std::to_string(digits.back()).size()) +
                                BASE_DIGITS * (static_cast<int>(digits.size()) - 1);
  }

  friend std::istream &operator>>(std::istream &in, BigInt &v) {
    std::string s;
    in >> s;
    v.read(s.size(), s.c_str());
    return in;
  }

  friend std::ostream &operator<<(std::ostream &out, const BigInt &v) {
    return out << v.to_string();
  }

  std::string to_string() const {
    std::ostringstream oss;
    if (sign == -1) {
      oss << '-';
    }
    oss << (digits.empty() ? 0 : digits.back());
    for (int i = static_cast<int>(digits.size()) - 2; i >= 0; i--) {
      oss << std::setw(BASE_DIGITS) << std::setfill('0') << digits[i];
    }
    return oss.str();
  }

  template<typename T>
  T to_arithmetic() const {
    std::stringstream ss(to_string());
    T res;
    ss >> res;
    return res;
  }

  int64_t to_llong() const { return to_arithmetic<int64_t>(); }
  double to_double() const { return to_arithmetic<double>(); }
  long double to_ldouble() const { return to_arithmetic<long double>(); }

  explicit operator int() const { return static_cast<int>(to_llong()); }
  explicit operator long long() const { return static_cast<long long>(to_llong()); }
  explicit operator double() const { return to_double(); }
  explicit operator long double() const { return to_ldouble(); }

  int comp(const BigInt &v) const { return comp(digits, v.digits, sign, v.sign); }

  // The comparison and binary arithmetic operators are hidden friends, so a raw integer operand on
  // either side converts through the implicit constructor.
  friend bool operator<(const BigInt &a, const BigInt &b) { return a.comp(b) < 0; }
  friend bool operator>(const BigInt &a, const BigInt &b) { return a.comp(b) > 0; }
  friend bool operator<=(const BigInt &a, const BigInt &b) { return a.comp(b) <= 0; }
  friend bool operator>=(const BigInt &a, const BigInt &b) { return a.comp(b) >= 0; }
  friend bool operator==(const BigInt &a, const BigInt &b) { return a.comp(b) == 0; }
  friend bool operator!=(const BigInt &a, const BigInt &b) { return a.comp(b) != 0; }

  BigInt abs() const {
    BigInt res(*this);
    res.sign = 1;
    return res;
  }

  BigInt operator-() const {
    BigInt res(*this);
    if (!digits.empty()) res.sign = -sign;
    return res;
  }

  friend BigInt operator+(const BigInt &a, const BigInt &b) {
    return add(a.digits, b.digits, a.sign, b.sign);
  }

  friend BigInt operator-(const BigInt &a, const BigInt &b) {
    return sub(a.digits, b.digits, a.sign, b.sign);
  }

  BigInt &operator*=(int v) {
    int64_t factor = v;
    if (factor < 0) {
      sign = -sign;
      factor = -factor;
    }
    int64_t carry = 0;
    for (int i = 0; i < static_cast<int>(digits.size()) || carry; i++) {
      if (i == static_cast<int>(digits.size())) {
        digits.push_back(0);
      }
      int64_t curr = digits[i] * factor + carry;
      carry = curr / BASE;
      digits[i] = static_cast<int>((curr % BASE));
    }
    normalize();
    return *this;
  }

  BigInt operator*(int v) const {
    BigInt res(*this);
    res *= v;
    return res;
  }

  friend BigInt operator*(const BigInt &u, const BigInt &v) {
    vint a = convert_base(u.digits, BASE_DIGITS, MUL_BASE_DIGITS);
    vint b = convert_base(v.digits, BASE_DIGITS, MUL_BASE_DIGITS);
    int n = 1;
    while (n < 2 * static_cast<int>(std::max(a.size(), b.size()))) {
      n <<= 1;
    }
    a.resize(n, 0);
    b.resize(n, 0);
    vint64 c;
    if (n < MUL_CUTOFF) {
      c = karatsuba(a.begin(), a.end(), b.begin(), b.end());
    } else {
      auto at = fft(a.begin(), a.end()), bt = fft(b.begin(), b.end());
      for (int i = 0; i < n; i++) {
        at[i] *= bt[i];
      }
      at = fft(at.begin(), at.end(), true);
      c.resize(n);
      for (int i = 0; i < n; i++) {
        c[i] = at[i].real() + 0.5;
      }
    }
    BigInt res;
    res.sign = u.sign * v.sign;
    for (int i = 0, carry = 0; i < static_cast<int>(c.size()); i++) {
      int64_t d = c[i] + carry;
      res.digits.push_back(d % MUL_BASE);
      carry = d / MUL_BASE;
    }
    res.digits = convert_base(res.digits, MUL_BASE_DIGITS, BASE_DIGITS);
    res.normalize();
    return res;
  }

  BigInt &operator/=(int v) {
    assert(v != 0);
    int64_t divisor = v;
    if (divisor < 0) {
      sign = -sign;
      divisor = -divisor;
    }
    int64_t rem = 0;
    for (int i = static_cast<int>(digits.size()) - 1; i >= 0; i--) {
      int64_t curr = digits[i] + rem * static_cast<int64_t>(BASE);
      digits[i] = static_cast<int>((curr / divisor));
      rem = curr % divisor;
    }
    normalize();
    return *this;
  }

  BigInt operator/(int v) const {
    BigInt res(*this);
    res /= v;
    return res;
  }

  int operator%(int v) const {
    assert(v != 0);
    int64_t divisor = v;
    if (divisor < 0) {
      divisor = -divisor;
    }
    int64_t m = 0;
    for (int i = static_cast<int>(digits.size()) - 1; i >= 0; i--) {
      m = (digits[i] + m * static_cast<int64_t>(BASE)) % divisor;
    }
    return static_cast<int>(m * sign);
  }

  std::pair<BigInt, BigInt> div(const BigInt &v) const {
    assert(v != 0);
    if (comp(digits, v.digits, 1, 1) < 0) {
      return {BigInt(0), *this};
    }
    auto res = static_cast<int>(v.digits.size()) < DIV_CUTOFF ? abs().div_schoolbook(v.abs())
                                                              : abs().div_recursive(v.abs());
    res.first.sign = sign * v.sign;
    res.second.sign = sign;
    res.first.normalize();
    res.second.normalize();
    return res;
  }

  friend BigInt operator/(const BigInt &a, const BigInt &b) { return a.div(b).first; }
  friend BigInt operator%(const BigInt &a, const BigInt &b) { return a.div(b).second; }
  BigInt operator++(int){ BigInt t(*this); operator++(); return t; }
  BigInt operator--(int){ BigInt t(*this); operator--(); return t; }
  BigInt &operator++() { *this = *this + BigInt(1); return *this; }
  BigInt &operator--() { *this = *this - BigInt(1); return *this; }
  BigInt &operator+=(const BigInt &v) { *this = *this + v; return *this; }
  BigInt &operator-=(const BigInt &v) { *this = *this - v; return *this; }
  BigInt &operator*=(const BigInt &v) { *this = *this * v; return *this; }
  BigInt &operator/=(const BigInt &v) { *this = *this / v; return *this; }
  BigInt &operator%=(const BigInt &v) { *this = *this % v; return *this; }

  BigInt pow(int n) const {
    assert(n >= 0);
    if (n == 0) {
      return BigInt(1);
    }
    if (*this == 0) {
      return BigInt(0);
    }
    BigInt x(*this), res(1);
    for (; n != 0; n >>= 1) {
      if (n & 1) {
        res *= x;
      }
      x *= x;
    }
    return res;
  }

  // Multiplying by a power of ten is a limb shift, since BASE is a power of ten by construction.
  BigInt mul_pow10(int n) const {
    assert(n >= 0);
    static const int POW10[] = {1, 10, 100, 1000, 10000, 100000, 1000000, 10000000, 100000000};
    BigInt res(*this);
    res.digits.insert(res.digits.begin(), n / BASE_DIGITS, 0);
    res *= POW10[n % BASE_DIGITS];  // Also normalizes, which is what clears a shifted zero.
    return res;
  }

  BigInt sqrt() const {
    assert(sign != -1);
    BigInt v(*this);
    while (v.digits.empty() || v.digits.size() % 2 == 1) {
      v.digits.push_back(0);
    }
    int n = static_cast<int>(v.digits.size());
    int ldig =
        static_cast<int>(::sqrt(static_cast<double>(v.digits[n - 1]) * BASE + v.digits[n - 2]));
    int norm = BASE / (ldig + 1);
    v *= norm;
    v *= norm;
    while (v.digits.empty() || v.digits.size() % 2 == 1) {
      v.digits.push_back(0);
    }
    BigInt r(static_cast<int64_t>(v.digits[n - 1]) * BASE + v.digits[n - 2]);
    int q = ldig =
        static_cast<int>(::sqrt(static_cast<double>(v.digits[n - 1]) * BASE + v.digits[n - 2]));
    BigInt res;
    for (int j = n / 2 - 1; j >= 0; j--) {
      for (;; q--) {
        BigInt r1 =
            (r - (res * 2 * BASE + q) * q) * BASE * BASE +
            (j > 0 ? static_cast<int64_t>(v.digits[2 * j - 1]) * BASE + v.digits[2 * j - 2] : 0);
        if (r1 >= 0) {
          r = r1;
          break;
        }
      }
      res = res * BASE + q;
      if (j > 0) {
        int sz1 = static_cast<int>(res.digits.size());
        int sz2 = static_cast<int>(r.digits.size());
        int d1 = (sz1 + 2 < sz2) ? r.digits[sz1 + 2] : 0;
        int d2 = (sz1 + 1 < sz2) ? r.digits[sz1 + 1] : 0;
        int d3 = (sz1 < sz2) ? r.digits[sz1] : 0;
        q = (static_cast<int64_t>(d1) * BASE * BASE + static_cast<int64_t>(d2) * BASE + d3) /
            (ldig * 2);
      }
    }
    res.normalize();
    return res / norm;
  }

  BigInt nth_root(int n) const {
    assert(n > 0 && (sign != -1 || n % 2 == 1));
    if (*this == 0) {
      return BigInt(0);
    }
    // Newton's iteration, started just above the answer and decreasing to it.
    BigInt magnitude = abs();
    BigInt x = BigInt(10).pow(static_cast<int>(std::ceil(static_cast<double>(size()) / n))), y;
    while (x.comp(y = (x * (n - 1) + magnitude / x.pow(n - 1)) / n) > 0) {
      x = y;
    }
    return sign == -1 ? -x : x;
  }

  static BigInt rand(int d) {
    if (d == 0) {
      return BigInt(0);
    }
    static std::mt19937 rng(std::chrono::steady_clock::now().time_since_epoch().count());
    std::uniform_int_distribution<int> first_digit(1, 9), digit(0, 9);
    std::string s(1, static_cast<char>('0' + first_digit(rng)));
    for (int i = 1; i < d; i++) {
      s += static_cast<char>('0' + digit(rng));
    }
    return BigInt(s);
  }

  friend int comp(const BigInt &a, const BigInt &b) { return a.comp(b); }
  friend BigInt abs(const BigInt &v) { return v.abs(); }
  friend BigInt pow(const BigInt &v, int n) { return v.pow(n); }
  friend BigInt sqrt(const BigInt &v) { return v.sqrt(); }
  friend BigInt nth_root(const BigInt &v, int n) { return v.nth_root(n); }
};
\end{lstlisting}
\Needspace*{4\baselineskip}
\textbf{Example Usage}\par
\begin{lstlisting}[style=cpp]
#include <climits>

int main() {
  BigInt a("-9899819294989142124"), b("12398124981294214");
  assert(a + b == "-9887421170007847910");
  assert(a - b == "-9912217419970436338");
  assert(a * b == "-122739196911503356525379735104870536");
  assert(a / b == "-798");

  // Raw integers and strings work on either side of comparisons and arithmetic.
  assert("12398124981294214" == b);
  assert(5 + BigInt(7) == 12);
  assert(-100 < b && 5 % BigInt(3) == 2);

  assert(BigInt(20).pow(12345).size() == 16062);
  assert(BigInt("9812985918924981892491829").nth_root(4) == 1769906);
  assert(BigInt(-8).nth_root(3) == -2);
  for (int i = -100; i <= 100; i++) {
    if (i >= 0) {
      assert(BigInt(i).sqrt() == static_cast<int>(sqrt(i)));
    }
    for (int j = -100; j <= 100; j++) {
      assert(BigInt(i) + BigInt(j) == i + j);
      assert(BigInt(i) - BigInt(j) == i - j);
      assert(BigInt(i) * BigInt(j) == i * j);
      if (j != 0) {
        assert(BigInt(i) / BigInt(j) == i / j);
      }
      if (0 < i && i <= 10 && 0 < j && j <= 10) {
        assert(BigInt(i).nth_root(j) == static_cast<int64_t>((pow(i, 1.0 / j) + 1E-5)));
        int64_t p = 1;
        for (int k = 0; k < j; k++) {
          p *= i;
        }
        assert(BigInt(i).pow(j) == p);
      }
    }
  }
  std::mt19937 rng(1234567);  // Fixed seed for reproducibility.
  std::uniform_int_distribution<int> len_dist(1, 100);
  for (int i = 0; i < 20; i++) {
    int d = len_dist(rng);
    BigInt val(BigInt::rand(d)), root(val.sqrt()), lower(root * root), upper(root + 1);
    upper *= upper;
    assert(lower <= val && val < upper);
    std::uniform_int_distribution<int> divisor_len_dist(1, d);
    BigInt divisor(BigInt::rand(divisor_len_dist(rng)) + 1), quotient(val / divisor);
    lower = quotient * divisor;
    upper = divisor * (quotient + 1);
    assert(val >= lower && val < upper);
  }
  BigInt x(-6);
  assert(x.to_string() == "-6");
  assert(x.to_llong() == -6LL);
  assert(x.to_double() == -6.0);
  assert(x.to_ldouble() == -6.0);
  assert(static_cast<int>(x) == -6);
  assert(static_cast<long long>(x) == -6LL);
  assert(static_cast<double>(x) == -6.0);
  assert(BigInt(INT64_MIN).to_string() == "-9223372036854775808");
  assert(BigInt(INT64_MIN).to_llong() == INT64_MIN);
  assert(-BigInt(0) == 0);
  assert(BigInt(3) * INT_MIN == "-6442450944");
  assert(BigInt("6442450944") / INT_MIN == -3);
  assert(BigInt("6442450945") % INT_MIN == 1);
  return 0;
}
\end{lstlisting}
\subsection{Big Decimal}
\label{sec:6.4.3}
Performs exact decimal arithmetic using the full \inlinecode{BigInt} implementation from \secref{6.4.2} as its coefficient type, together with a nonnegative scale. Unlike binary floating point, every parsed decimal is represented exactly: for example, \inlinecode{BigDecimal("1.20")} stores coefficient $120$ with scale $2$. Addition and subtraction preserve the larger operand scale, multiplication adds the scales, and formatting preserves trailing zeros unless \inlinecode{normalized()} is called. Only finite, well-formed decimal strings are supported; \inlinecode{NaN} and infinities are not represented.

Division and square root generally have nonterminating decimal expansions, so both operations require an explicit result scale. The rounding mode defaults to \inlinecode{HalfAway}, which rounds ties away from zero. The other modes are \inlinecode{TowardZero}, \inlinecode{AwayFromZero}, \inlinecode{Floor}, \inlinecode{Ceil}, and \inlinecode{HalfEven}.

\begin{compactitem}
  \item \inlinecode{BigDecimal(n = 0)} constructs a decimal from integer \inlinecode{n}.
  \item \inlinecode{BigDecimal(s)} parses decimal string \inlinecode{s}, optionally written in scientific notation, and preserves trailing zeros in its mantissa.
  \item \inlinecode{decimal\_places()}, \inlinecode{precision()}, and \inlinecode{sign()} return the scale, the number of coefficient digits, and the sign as $-1$, $0$, or $1$, respectively.
  \item \inlinecode{unscaled\_value()} returns the underlying \inlinecode{BigInt} coefficient.
  \item \inlinecode{abs()} and \inlinecode{normalized()} return the absolute value and a value with trailing fractional zeros removed.
  \item \inlinecode{move\_point(places)} moves the decimal point right by \inlinecode{places} positions, or left if \inlinecode{places} is negative.
  \item \inlinecode{rescale(new\_scale, mode = Round::HalfAway)} returns the value with exactly \inlinecode{new\_scale} decimal places, rounding if the scale is reduced.
  \item \inlinecode{trunc()}, \inlinecode{floor()}, and \inlinecode{ceil()} return the corresponding integer as a \inlinecode{BigInt}.
  \item \inlinecode{divide(v, result\_scale, mode = Round::HalfAway)} divides by nonzero \inlinecode{v} and rounds to exactly \inlinecode{result\_scale} decimal places.
  \item \inlinecode{sqrt(result\_scale, mode = Round::HalfAway)} returns the nonnegative square root rounded to exactly \inlinecode{result\_scale} decimal places.
  \item \inlinecode{pow(n)} returns the value raised to nonnegative integer power \inlinecode{n}.
  \item \inlinecode{to\_string()} returns ordinary decimal notation with the stored number of decimal places.
  \item Stream operators and comparisons are defined, as are \inlinecode{+}, \inlinecode{-}, \inlinecode{*}, \inlinecode{\%}, and their compound assignments.
\end{compactitem}

Overflow warning: The stored scale and every decimal-shift magnitude must fit in \inlinecode{int}.

\Needspace*{3\baselineskip}
\textbf{Time Complexity}:
\begin{compactitem}
  \item $O(1)$ per call to \inlinecode{decimal\_places()}, \inlinecode{precision()}, \inlinecode{sign()}, and \inlinecode{unscaled\_value()}.
  \item $O(d)$ per call to parsing, formatting, normalization, addition, subtraction, comparisons, and decimal-point movement, where $d$ is the number of relevant decimal digits.
  \item $O(M(d) \log n)$ per call to \inlinecode{pow(n)}, where $M(d)$ is the multiplication cost for the largest intermediate coefficient.
  \item $O(d^{2})$ per call to square root, matching the \inlinecode{BigInt} operation used here.
  \item Division, remainder, and rescaling to fewer places inherit the \inlinecode{BigInt} division cost of \secref{6.4.2}: $O((d/m) \cdot M(m) \cdot \log m)$ once the divisor reaches its \inlinecode{DIV\_CUTOFF} and $O(d \cdot m)$ below it, where $d$ and $m$ are the number of digits in the dividend and divisor.
\end{compactitem}

\textbf{Space Complexity}: $O(d)$ for stored digits, returned values, and arithmetic results.

\begin{lstlisting}[style=cpp]
#define main bigint_example
#include "6.4.2_Big_Integer.cpp"
#undef main

#include <climits>

class BigDecimal {
 public:
  enum class Round { TowardZero, AwayFromZero, Floor, Ceil, HalfAway, HalfEven };

 private:
  BigInt coeff;
  int scale;

  static int checked_scale(int64_t scale) {
    assert(0 <= scale && scale <= INT_MAX);
    return static_cast<int>(scale);
  }

  BigDecimal(BigInt coeff, int scale) : coeff(std::move(coeff)), scale(scale) {}

  BigInt scaled(int new_scale) const { return coeff.mul_pow10(new_scale - scale); }
  static BigInt pow10(int n) { return BigInt(10).pow(n); }

  static bool round_up(bool negative, int half_cmp, const BigInt &q, Round mode) {
    return mode == Round::AwayFromZero || (mode == Round::Floor && negative) ||
           (mode == Round::Ceil && !negative) || (mode == Round::HalfAway && half_cmp >= 0) ||
           (mode == Round::HalfEven && (half_cmp > 0 || (half_cmp == 0 && q % 2 != 0)));
  }

  static BigInt rounded_quotient(const BigInt &num, const BigInt &den, Round mode) {
    assert(den > 0);
    bool negative = num < 0;
    auto [q, r] = num.abs().div(den);
    bool increment = r != 0 && round_up(negative, (2 * r).comp(den), q, mode);
    return negative ? -(q + increment) : q + increment;
  }

 public:
  BigDecimal(int64_t n = 0) : coeff(n), scale(0) {}

  explicit BigDecimal(const std::string &s) : coeff(0), scale(0) {
    assert(s.size() <= INT_MAX);
    size_t e = s.find_first_of("eE");
    std::string mantissa = s.substr(0, e);
    int exponent = e == std::string::npos ? 0 : std::stoi(s.substr(e + 1));
    int pos = 0;
    bool negative = false, point = false;
    if (pos < static_cast<int>(mantissa.size()) && (mantissa[pos] == '+' || mantissa[pos] == '-')) {
      negative = mantissa[pos++] == '-';
    }
    std::string digits;
    for (; pos < static_cast<int>(mantissa.size()); pos++) {
      if (mantissa[pos] == '.') {
        assert(!point);
        point = true;
      } else {
        assert('0' <= mantissa[pos] && mantissa[pos] <= '9');
        digits += mantissa[pos];
        scale += point;
      }
    }
    assert(!digits.empty());
    coeff = negative ? -BigInt(digits) : BigInt(digits);
    int64_t adjusted_scale = static_cast<int64_t>(scale) - exponent;
    if (adjusted_scale < 0) {
      coeff = coeff.mul_pow10(checked_scale(-adjusted_scale));
      scale = 0;
    } else {
      scale = checked_scale(adjusted_scale);
    }
  }

  int decimal_places() const { return scale; }
  int precision() const { return coeff.size(); }
  int sign() const { return (coeff > 0) - (coeff < 0); }
  const BigInt &unscaled_value() const { return coeff; }
  BigDecimal abs() const { return BigDecimal(coeff.abs(), scale); }

  BigDecimal normalized() const {
    BigInt c = coeff;
    int s = scale;
    while (s > 0 && c % 10 == 0) {
      c /= 10;
      s--;
    }
    return BigDecimal(c, s);
  }

  BigDecimal move_point(int places) const {
    int64_t new_scale = static_cast<int64_t>(scale) - places;
    return new_scale >= 0 ? BigDecimal(coeff, checked_scale(new_scale))
                          : BigDecimal(coeff.mul_pow10(checked_scale(-new_scale)), 0);
  }

  BigDecimal rescale(int new_scale, Round mode = Round::HalfAway) const {
    assert(new_scale >= 0);
    return new_scale >= scale
               ? BigDecimal(scaled(new_scale), new_scale)
               : BigDecimal(rounded_quotient(coeff, pow10(scale - new_scale), mode), new_scale);
  }

  BigInt trunc() const { return rescale(0, Round::TowardZero).coeff; }
  BigInt floor() const { return rescale(0, Round::Floor).coeff; }
  BigInt ceil() const { return rescale(0, Round::Ceil).coeff; }

  BigDecimal divide(const BigDecimal &v, int result_scale, Round mode = Round::HalfAway) const {
    assert(v.coeff != 0 && result_scale >= 0);
    BigInt num = coeff.abs(), den = v.coeff.abs();
    int64_t shift = static_cast<int64_t>(v.scale) + result_scale - scale;
    (shift >= 0 ? num : den) =
        (shift >= 0 ? num : den).mul_pow10(checked_scale(shift < 0 ? -shift : shift));
    if ((coeff < 0) != (v.coeff < 0)) {
      num = -num;
    }
    return BigDecimal(rounded_quotient(num, den, mode), result_scale);
  }

  BigDecimal sqrt(int result_scale, Round mode = Round::HalfAway) const {
    assert(coeff >= 0 && result_scale >= 0);
    BigInt num = coeff, den = 1;
    int64_t shift = 2LL * result_scale - scale;
    (shift >= 0 ? num : den) =
        (shift >= 0 ? num : den).mul_pow10(checked_scale(shift < 0 ? -shift : shift));
    BigInt q = (num / den).sqrt();
    bool exact = q * q * den == num;
    int half_cmp = (4 * num).comp(den * (4 * q * q + 4 * q + 1));
    return BigDecimal(q + (!exact && round_up(false, half_cmp, q, mode)), result_scale);
  }

  std::string to_string() const {
    std::string s = coeff.abs().to_string();
    if (static_cast<int>(s.size()) <= scale) {
      s.insert(0, scale + 1 - s.size(), '0');
    }
    if (scale > 0) {
      s.insert(s.size() - scale, 1, '.');
    }
    return coeff < 0 ? "-" + s : s;
  }

  friend std::istream &operator>>(std::istream &in, BigDecimal &v) {
    std::string s;
    if (in >> s) {
      v = BigDecimal(s);
    }
    return in;
  }

  friend std::ostream &operator<<(std::ostream &out, const BigDecimal &v) {
    return out << v.to_string();
  }

  friend bool operator<(const BigDecimal &a, const BigDecimal &b) {
    int s = std::max(a.scale, b.scale);
    return a.scaled(s) < b.scaled(s);
  }

  friend bool operator==(const BigDecimal &a, const BigDecimal &b) {
    int s = std::max(a.scale, b.scale);
    return a.scaled(s) == b.scaled(s);
  }

  friend bool operator>(const BigDecimal &a, const BigDecimal &b) { return b < a; }
  friend bool operator<=(const BigDecimal &a, const BigDecimal &b) { return !(b < a); }
  friend bool operator>=(const BigDecimal &a, const BigDecimal &b) { return !(a < b); }
  friend bool operator!=(const BigDecimal &a, const BigDecimal &b) { return !(a == b); }

  friend BigDecimal operator+(const BigDecimal &a, const BigDecimal &b) {
    int s = std::max(a.scale, b.scale);
    return BigDecimal(a.scaled(s) + b.scaled(s), s);
  }

  friend BigDecimal operator-(const BigDecimal &a, const BigDecimal &b) { return a + -b; }

  friend BigDecimal operator*(const BigDecimal &a, const BigDecimal &b) {
    return BigDecimal(a.coeff * b.coeff, checked_scale(static_cast<int64_t>(a.scale) + b.scale));
  }

  friend BigDecimal operator%(const BigDecimal &a, const BigDecimal &b) {
    assert(b.coeff != 0);
    int s = std::max(a.scale, b.scale);
    return BigDecimal(a.scaled(s) % b.scaled(s), s);
  }

  BigDecimal operator-() const { return BigDecimal(-coeff, scale); }

  BigDecimal pow(int n) const {
    assert(n >= 0);
    return BigDecimal(coeff.pow(n), checked_scale(static_cast<int64_t>(scale) * n));
  }

  BigDecimal &operator+=(const BigDecimal &v) { return *this = *this + v; }
  BigDecimal &operator-=(const BigDecimal &v) { return *this = *this - v; }
  BigDecimal &operator*=(const BigDecimal &v) { return *this = *this * v; }
  BigDecimal &operator%=(const BigDecimal &v) { return *this = *this % v; }
};
\end{lstlisting}
\Needspace*{4\baselineskip}
\textbf{Example Usage}\par
\begin{lstlisting}[style=cpp]
using namespace std;

int main() {
  using Round = BigDecimal::Round;

  assert(BigDecimal().to_string() == "0");
  assert(BigDecimal(-12).to_string() == "-12");
  assert(BigDecimal("+.50").to_string() == "0.50");
  assert(BigDecimal("1.20e3").to_string() == "1200");
  assert(BigDecimal("1.20e-2").to_string() == "0.0120");

  BigDecimal a("1.20"), b("3.4");
  assert(a.decimal_places() == 2 && a.precision() == 3);
  assert(a.unscaled_value() == 120);
  assert(a.sign() == 1 && BigDecimal(0).sign() == 0 && BigDecimal(-1).sign() == -1);
  assert(BigDecimal("-1.20").abs().to_string() == "1.20");
  assert(a.normalized().to_string() == "1.2");
  assert(a == BigDecimal("1.2") && a != b);
  assert(a < b && b > a && a <= BigDecimal("1.200") && b >= a);

  assert((a + b).to_string() == "4.60");
  assert((b - a).to_string() == "2.20");
  assert((a * b).to_string() == "4.080");
  assert((-a).to_string() == "-1.20");
  assert((BigDecimal("5.50") % 2).to_string() == "1.50");

  BigDecimal x = a;
  x += b;
  assert(x.to_string() == "4.60");
  x -= b;
  assert(x.to_string() == "1.20");
  x *= b;
  assert(x.to_string() == "4.080");
  x %= 1.0;
  assert(x.to_string() == "0.080");

  assert(a.move_point(2).to_string() == "120");
  assert(a.move_point(-1).to_string() == "0.120");
  assert(a.rescale(4).to_string() == "1.2000");
  assert(BigDecimal("1.29").rescale(1, Round::TowardZero).to_string() == "1.2");
  assert(BigDecimal("1.21").rescale(1, Round::AwayFromZero).to_string() == "1.3");
  assert(BigDecimal("-1.21").rescale(1, Round::Floor).to_string() == "-1.3");
  assert(BigDecimal("-1.21").rescale(1, Round::Ceil).to_string() == "-1.2");
  assert(BigDecimal("1.25").rescale(1, Round::HalfAway).to_string() == "1.3");
  assert(BigDecimal("2.685").rescale(2, Round::HalfEven).to_string() == "2.68");

  assert(BigDecimal(1).divide(8, 4).to_string() == "0.1250");
  assert(BigDecimal(2).divide(3, 4).to_string() == "0.6667");
  assert(BigDecimal(-1).divide(8, 2).to_string() == "-0.13");
  assert(BigDecimal(1).divide(8, 2, Round::HalfEven).to_string() == "0.12");

  assert(BigDecimal("-1.2").trunc() == -1);
  assert(BigDecimal("-1.2").floor() == -2);
  assert(BigDecimal("-1.2").ceil() == -1);
  assert(a.pow(3).to_string() == "1.728000");

  assert(BigDecimal(2).sqrt(5).to_string() == "1.41421");
  assert(BigDecimal("2.25").sqrt(2).to_string() == "1.50");
  assert(BigDecimal("2.25").sqrt(0, Round::HalfAway).to_string() == "2");
  assert(BigDecimal("6.25").sqrt(0, Round::HalfEven).to_string() == "2");
  assert(BigDecimal(2).sqrt(0, Round::Ceil).to_string() == "2");
  for (int i = 0; i <= 100; i++) {
    assert(BigDecimal(i * i).sqrt(0) == BigDecimal(i));
  }

  stringstream in("-12.340"), out;
  BigDecimal streamed;
  in >> streamed;
  out << streamed;
  assert(out.str() == "-12.340");

  assert((BigDecimal("12345678901234567890.12") * 10).to_string() == "123456789012345678901.20");
  return 0;
}
\end{lstlisting}
\subsection{Rational Numbers}
\label{sec:6.4.4}
Perform operations on rational numbers internally represented as two integers: a numerator and a denominator. The template integer type must support streamed input/output, comparisons, and arithmetic operations.

\begin{compactitem}
  \item \inlinecode{Rational\textless{}Int\textgreater{}(n)} constructs a rational number with numerator \inlinecode{n} and denominator $1$.
  \item \inlinecode{Rational\textless{}Int\textgreater{}(n, d)} constructs a rational number with numerator \inlinecode{n} and denominator \inlinecode{d}.
  \item \inlinecode{operator\textgreater{}\textgreater{}} inputs a rational number using the next integer from the stream as the numerator and 1 as the denominator.
  \item \inlinecode{operator\textless{}\textless{}} outputs the rational number as a string consisting of possibly a minus sign followed by the numerator, followed by a slash, followed by the denominator.
  \item \inlinecode{to\_string()}, \inlinecode{to\_llong()}, \inlinecode{to\_double()}, and \inlinecode{to\_ldouble()} return the rational converted to an \inlinecode{std::string}, \inlinecode{int64\_t}, \inlinecode{double}, and \inlinecode{long double} respectively with potential truncation or approximation for the primitive types. Equivalent conversions are also available through explicit casts to \inlinecode{int}, \inlinecode{long long}, \inlinecode{double}, and \inlinecode{long double}.
  \item \inlinecode{to\_arithmetic\textless{}T\textgreater{}()} is the general form behind the three numeric conversions above, returning the rational as any arithmetic type \inlinecode{T} readable from a stream. Integral \inlinecode{T} truncates toward zero.
  \item \inlinecode{abs()}, \inlinecode{floor()}, and \inlinecode{ceil()} return the absolute value, floor, and ceiling.
  \item Operators \inlinecode{\textless{}}, \inlinecode{\textgreater{}}, \inlinecode{\textless{}=}, \inlinecode{\textgreater{}=}, \inlinecode{==}, \inlinecode{!=}, \inlinecode{+}, \inlinecode{-}, \inlinecode{*}, \inlinecode{/}, \inlinecode{\%}, \inlinecode{++}, \inlinecode{--}, \inlinecode{+=}, \inlinecode{-=}, \inlinecode{*=}, \inlinecode{/=}, and \inlinecode{\%=} are defined analogous to those on numerical primitives. The comparisons and binary arithmetic are hidden friends, so a raw integer operand works on either side.
\end{compactitem}

Overflow warning: Internal operations do not check for overflow. Comparisons and arithmetic cross-multiply numerators and denominators, so instantiate with a wider integer type (such as \inlinecode{\_\_int128}, or even \inlinecode{BigInt}) if the values may grow large.

\Needspace*{3\baselineskip}
\textbf{Time Complexity}:
\begin{compactitem}
  \item $O(\log s)$ operations on \inlinecode{Int} per call to the constructor \inlinecode{Rational(n, d)}, where $s = |n| + |d|$, spent reducing to lowest terms.
  \item $O(1)$ operations on \inlinecode{Int} per call to everything else.
\end{compactitem}

\Needspace*{3\baselineskip}
\textbf{Space Complexity}:
\begin{compactitem}
  \item Two \inlinecode{Int} values for storage of the rational number.
  \item $O(1)$ auxiliary \inlinecode{Int} values for all operations.
\end{compactitem}

\begin{lstlisting}[style=cpp]
#include <cassert>
#include <cstdint>
#include <istream>
#include <ostream>
#include <sstream>
#include <string>

template<typename Int>
class Rational {
  Int num, den;

 public:
  Rational() : num(0), den(1) {}
  Rational(const Int &n) : num(n), den(1) {}

  Rational(const Int &n, const Int &d) : num(n), den(d) {
    assert(den != 0);
    if (den < 0) {
      num = -num;
      den = -den;
    }
    Int a(num < 0 ? -num : num), b(den), tmp;
    while (a != 0 && b != 0) {
      tmp = a % b;
      a = b;
      b = tmp;
    }
    Int gcd = (b == 0) ? a : b;
    num /= gcd;
    den /= gcd;
  }

  friend std::istream &operator>>(std::istream &in, Rational &r) {
    in >> r.num;
    r.den = 1;
    return in;
  }

  friend std::ostream &operator<<(std::ostream &out, const Rational &r) {
    out << r.num << "/" << r.den;
    return out;
  }

  std::string to_string() const {
    std::stringstream ss;
    ss << num << "/" << den;
    return ss.str();
  }

  // Round-trips through a stringstream so that any Int supporting streamed I/O (e.g. a big-integer
  // type) converts, even without a direct cast to the target type. Integral T truncates.
  template<typename T>
  T to_arithmetic() const {
    std::stringstream ss;
    ss << num << " " << den;
    T n, d;
    ss >> n >> d;
    return n / d;
  }

  int64_t to_llong() const { return to_arithmetic<int64_t>(); }
  double to_double() const { return to_arithmetic<double>(); }
  long double to_ldouble() const { return to_arithmetic<long double>(); }

  explicit operator int() const { return static_cast<int>(to_llong()); }
  explicit operator long long() const { return static_cast<long long>(to_llong()); }
  explicit operator double() const { return to_double(); }
  explicit operator long double() const { return to_ldouble(); }

  // The comparison and binary arithmetic operators are hidden friends, so a raw integer operand on
  // either side converts through the implicit constructor.
  friend bool operator<(const Rational &a, const Rational &b) {
    return a.num * b.den < b.num * a.den;
  }

  friend bool operator==(const Rational &a, const Rational &b) {
    return a.num == b.num && a.den == b.den;
  }

  friend bool operator>(const Rational &a, const Rational &b) { return b < a; }
  friend bool operator<=(const Rational &a, const Rational &b) { return !(b < a); }
  friend bool operator>=(const Rational &a, const Rational &b) { return !(a < b); }
  friend bool operator!=(const Rational &a, const Rational &b) { return !(a == b); }

  Rational abs() const { return Rational(num < 0 ? -num : num, den); }
  friend Rational abs(const Rational &r) { return r.abs(); }
  Int floor() const { return num < 0 ? -((-num + den - 1) / den) : num / den; }
  Int ceil() const { return num < 0 ? -(-num / den) : (num + den - 1) / den; }

  friend Rational operator+(const Rational &a, const Rational &b) {
    return Rational(a.num * b.den + b.num * a.den, a.den * b.den);
  }

  friend Rational operator-(const Rational &a, const Rational &b) {
    return Rational(a.num * b.den - b.num * a.den, a.den * b.den);
  }

  friend Rational operator*(const Rational &a, const Rational &b) {
    return Rational(a.num * b.num, a.den * b.den);
  }

  friend Rational operator/(const Rational &a, const Rational &b) {
    return Rational(a.num * b.den, a.den * b.num);
  }

  friend Rational operator%(const Rational &a, const Rational &b) {
    return a - b * Rational(a.num * b.den / (b.num * a.den), 1);
  }

  Rational operator-() const { return Rational(-num, den); }
  Rational operator++(int) { Rational t(*this); operator++(); return t; }
  Rational operator--(int) { Rational t(*this); operator--(); return t; }
  Rational &operator++() { *this = *this + 1; return *this; }
  Rational &operator--() { *this = *this - 1; return *this; }
  Rational &operator+=(const Rational &r) { *this = *this + r; return *this; }
  Rational &operator-=(const Rational &r) { *this = *this - r; return *this; }
  Rational &operator*=(const Rational &r) { *this = *this * r; return *this; }
  Rational &operator/=(const Rational &r) { *this = *this / r; return *this; }
  Rational &operator%=(const Rational &r) { *this = *this % r; return *this; }
};
\end{lstlisting}
\Needspace*{4\baselineskip}
\textbf{Example Usage}\par
\begin{lstlisting}[style=cpp]
#include <cmath>
using namespace std;

bool EQ(double a, double b) {
  return fabs(a - b) < 1e-9;
}

int main() {
  using Rational = Rational<int64_t>;
  assert(Rational(-21, 1) % 2 == -1);
  Rational r(Rational(-53, 10) % Rational(-17, 10));
  assert(EQ(r.to_ldouble(), fmod(-5.3, -1.7)));
  assert(r.to_string() == "-1/5");
  assert(static_cast<int>(Rational(7, 2)) == 3);
  assert(static_cast<long long>(Rational(7, 2)) == 3LL);
  assert(EQ(static_cast<double>(Rational(1, 2)), 0.5));

  // Raw integers work on either side of comparisons and arithmetic.
  assert(2 + Rational(1, 2) == Rational(5, 2));
  assert(1 < Rational(3, 2) && Rational(3, 2) < 2);
  assert(2 % Rational(3, 4) == Rational(1, 2));

  // Construction reduces the fraction and moves any sign onto the numerator.
  assert(Rational(2, 4).to_string() == "1/2");
  assert(Rational(1, -2).to_string() == "-1/2");
  assert(Rational(5).to_string() == "5/1");

  // Arithmetic between rationals reduces the result too.
  assert(Rational(1, 2) + Rational(1, 3) == Rational(5, 6));
  assert(Rational(1, 2) - Rational(1, 3) == Rational(1, 6));
  assert(Rational(2, 3) * Rational(3, 4) == Rational(1, 2));
  assert(Rational(1, 2) / Rational(3, 4) == Rational(2, 3));

  // Rounding takes a different branch on each sign, and is exact on whole numbers.
  assert(Rational(7, 2).floor() == 3 && Rational(7, 2).ceil() == 4);
  assert(Rational(-7, 2).floor() == -4 && Rational(-7, 2).ceil() == -3);
  assert(Rational(4, 2).floor() == 2 && Rational(4, 2).ceil() == 2);
  assert(Rational(-7, 2).abs() == Rational(7, 2) && abs(Rational(-7, 2)) == Rational(7, 2));

  // The named conversions and the general form behind them.
  assert(Rational(7, 2).to_llong() == 3 && EQ(Rational(7, 2).to_double(), 3.5));
  assert(Rational(7, 2).to_arithmetic<int>() == 3);

  // Stream input reads one integer, and output prints the reduced fraction.
  Rational streamed;
  istringstream in("7");
  in >> streamed;
  assert(streamed == 7);
  ostringstream out;
  out << Rational(-3, 6);
  assert(out.str() == "-1/2");
  return 0;
}
\end{lstlisting}

\section{\texorpdfstring{Linear Algebra}{6.5 Linear Algebra}}
\setcounter{section}{5}
\setcounter{subsection}{0}
\subsection{Matrix Operations}
\label{sec:6.5.1}
Basic matrix operations defined on a two-dimensional vector of numeric values. Matrix arithmetic also supports modular value types such as \inlinecode{Modular\textless{}MOD\textgreater{}} from \secref{6.3.3}. The element type must be constructible from $0$ and $1$ and support the arithmetic operators used by the requested operation.

\begin{compactitem}
  \item \inlinecode{make\_matrix\textless{}T\textgreater{}(m, n, v)} constructs and returns an $m$ by $n$ matrix with 0-based indices (row indices $[0, {\inlinecodemath{m}})$ and column indices $[0, {\inlinecodemath{n}})$), where every value is initialized to \inlinecode{v}.
  \item \inlinecode{identity\_matrix\textless{}T\textgreater{}(n)} returns the \inlinecode{n} by \inlinecode{n} identity matrix, that is, a matrix where each \inlinecode{a[i][j]} equals $1$ if \inlinecode{i} = \inlinecode{j}, or $0$ otherwise.
  \item \inlinecode{rows(a)} returns the number of rows in matrix \inlinecode{a}.
  \item \inlinecode{cols(a)} returns the number of columns in matrix \inlinecode{a}.
  \item \inlinecode{a[i][j]} may be used to access or modify the specified entry of \inlinecode{a}.
  \item Operators \inlinecode{\textless{}}, \inlinecode{\textgreater{}}, \inlinecode{\textless{}=}, \inlinecode{\textgreater{}=}, \inlinecode{==}, and \inlinecode{!=} define lexicographical comparison based on that of \inlinecode{std::vector}.
  \item Operators \inlinecode{+}, \inlinecode{-}, \inlinecode{*}, \inlinecode{/}, \inlinecode{+=}, \inlinecode{-=}, \inlinecode{*=}, and \inlinecode{/=} define scalar addition, subtraction, multiplication, and division involving a matrix and a numeric scalar value.
  \item \inlinecode{a * v} returns the matrix-vector product as a vector.
  \item Operators \inlinecode{*} and \inlinecode{*=} define matrix multiplication.
\end{compactitem}

Multiplication is the schoolbook triple loop, which is $O(m \cdot n \cdot p)$ and is what contest constraints assume. Strassen's algorithm reaches $O(n^{2.807})$ by splitting each matrix into quarters and combining them with seven recursive products instead of eight, and Toom-Cook style refinements go lower still, but the additions that replace the saved multiplication carry large constants and worsen numerical stability. The crossover sits in the hundreds of rows even in tuned libraries, so the plain loop below is the right default; reach for a blocked or cache-oblivious ordering before reaching for a subcubic algorithm.

Exponentiation uses iterative binary exponentiation: keep an accumulated result, square the base each round, and multiply it when the current exponent bit is set. The power sum uses the block identity that $\begin{bmatrix} A & A \\ 0 & I \end{bmatrix}^p$ has $A + A^2 + \dots + A^p$ in its upper-right block.

\begin{compactitem}
  \item Operators \inlinecode{\textasciicircum{}} and \inlinecode{\textasciicircum{}=} define iterative binary matrix exponentiation of a square matrix \inlinecode{a} by a \inlinecode{uint64\_t} power \inlinecode{p}.
  \item \inlinecode{power\_sum(a, p)} returns the power sum of a square matrix \inlinecode{a} up to a \inlinecode{uint64\_t} power \inlinecode{p}, that is, $a + a^2 + \ldots + a^p$.
\end{compactitem}

\Needspace*{3\baselineskip}
\textbf{Time Complexity}:
\begin{compactitem}
  \item $O(m \cdot n)$ per construction, output, comparison, or scalar-arithmetic operation on $m$ by $n$ matrices.
  \item $O(1)$ per call to \inlinecode{rows()} and \inlinecode{cols()}.
  \item $O(m \cdot n)$ per matrix-matrix addition or subtraction of $m$ by $n$ matrices.
  \item $O(n^{3} \log p)$ per exponentiation of an $n$ by $n$ matrix to power $p$.
  \item $O(n^{3} \log p)$ per call to \inlinecode{power\_sum()} for an $n$ by $n$ matrix and power $p$.
  \item $O(m \cdot n)$ per multiplication of an $m$ by $n$ matrix by a vector of length $n$.
  \item $O(m \cdot n \cdot k)$ per multiplication of an $m$ by $n$ matrix by an $n$ by $k$ matrix.
\end{compactitem}

\Needspace*{3\baselineskip}
\textbf{Space Complexity}:
\begin{compactitem}
  \item $O(1)$ auxiliary for \inlinecode{rows()}, \inlinecode{cols()}, \inlinecode{a[i][j]} access, comparison and scalar compound operators.
  \item $O(n^{2})$ auxiliary for exponentiation of an $n$ by $n$ matrix to power $p$.
  \item $O(n^{2})$ auxiliary for \inlinecode{power\_sum()} of an $n$ by $n$ matrix up to power $p$.
  \item $O(m)$ for the vector returned by matrix-vector multiplication.
  \item $O(m \cdot n)$ auxiliary for all non-in-place operations returning an $m$ by $n$ matrix.
\end{compactitem}

\begin{lstlisting}[style=cpp]
#include <cassert>
#include <cstdint>
#include <iomanip>
#include <ostream>
#include <vector>

template<typename T>
using Matrix = std::vector<std::vector<T>>;

template<typename T = int>
Matrix<T> make_matrix(int m, int n, const T &v = T{}) {
  assert(m >= 0 && n >= 0);
  return Matrix<T>(m, std::vector<T>(n, v));
}

template<typename T = int>
Matrix<T> identity_matrix(int n) {
  Matrix<T> res = make_matrix<T>(n, n);
  for (int i = 0; i < n; i++) {
    res[i][i] = 1;
  }
  return res;
}

template<typename T>
int rows(const Matrix<T> &a) {
  return static_cast<int>(a.size());
}

template<typename T>
int cols(const Matrix<T> &a) {
  return a.empty() ? 0 : static_cast<int>(a[0].size());
}

template<typename T>
std::ostream &operator<<(std::ostream &out, const Matrix<T> &a) {
  auto flags = out.flags();
  auto precision = out.precision();
  out << std::fixed << std::setprecision(5);
  for (int i = 0; i < rows(a); i++) {
    for (int j = 0; j < cols(a); j++) {
      out << std::setw(10) << a[i][j];
    }
    out << std::endl;
  }
  out.flags(flags);
  out.precision(precision);
  return out;
}

template<typename T, typename U>
Matrix<T> &operator+=(Matrix<T> &a, const U &v) {
  for (int i = 0; i < rows(a); i++) {
    for (int j = 0; j < cols(a); j++) {
      a[i][j] += v;
    }
  }
  return a;
}

template<typename T, typename U>
Matrix<T> &operator-=(Matrix<T> &a, const U &v) {
  for (int i = 0; i < rows(a); i++) {
    for (int j = 0; j < cols(a); j++) {
      a[i][j] -= v;
    }
  }
  return a;
}

template<typename T, typename U>
Matrix<T> &operator*=(Matrix<T> &a, const U &v) {
  for (int i = 0; i < rows(a); i++) {
    for (int j = 0; j < cols(a); j++) {
      a[i][j] *= v;
    }
  }
  return a;
}

template<typename T, typename U>
Matrix<T> &operator/=(Matrix<T> &a, const U &v) {
  for (int i = 0; i < rows(a); i++) {
    for (int j = 0; j < cols(a); j++) {
      a[i][j] /= v;
    }
  }
  return a;
}

template<typename T>
Matrix<T> &operator+=(Matrix<T> &a, const Matrix<T> &b) {
  assert(rows(a) == rows(b) && cols(a) == cols(b));
  for (int i = 0; i < rows(a); i++) {
    for (int j = 0; j < cols(a); j++) {
      a[i][j] += b[i][j];
    }
  }
  return a;
}

template<typename T>
Matrix<T> &operator-=(Matrix<T> &a, const Matrix<T> &b) {
  assert(rows(a) == rows(b) && cols(a) == cols(b));
  for (int i = 0; i < rows(a); i++) {
    for (int j = 0; j < cols(a); j++) {
      a[i][j] -= b[i][j];
    }
  }
  return a;
}

template<typename T>
Matrix<T> operator+(const Matrix<T> &a, const Matrix<T> &b) {
  Matrix<T> c(a);
  return c += b;
}

template<typename T>
Matrix<T> operator-(const Matrix<T> &a, const Matrix<T> &b) {
  Matrix<T> c(a);
  return c -= b;
}

template<typename T>
std::vector<T> operator*(const Matrix<T> &a, const std::vector<T> &v) {
  assert(cols(a) == static_cast<int>(v.size()));
  std::vector<T> res(rows(a));
  for (int i = 0; i < rows(a); i++) {
    for (int j = 0; j < cols(a); j++) {
      res[i] += a[i][j] * v[j];
    }
  }
  return res;
}

template<typename T>
Matrix<T> operator*(const Matrix<T> &a, const Matrix<T> &b) {
  assert(cols(a) == rows(b));
  Matrix<T> res = make_matrix<T>(rows(a), cols(b));
  // The k loop sits outside j so the inner loop walks res[i] and b[k] contiguously instead of
  // striding down a column of b. Each res[i][j] still accumulates in increasing k, so results are
  // unchanged, including for floating-point types.
  for (int i = 0; i < rows(a); i++) {
    for (int k = 0; k < rows(b); k++) {
      for (int j = 0; j < cols(b); j++) {
        res[i][j] += a[i][k] * b[k][j];
      }
    }
  }
  return res;
}

template<typename T>
Matrix<T> &operator*=(Matrix<T> &a, const Matrix<T> &b) { return a = a * b; }

template<typename T, typename U>
Matrix<T> operator+(const Matrix<T> &a, const U &v) { Matrix<T> m(a); return m += v; }

template<typename T, typename U>
Matrix<T> operator-(const Matrix<T> &a, const U &v) { Matrix<T> m(a); return m -= v; }

template<typename T, typename U>
Matrix<T> operator*(const Matrix<T> &a, const U &v) { Matrix<T> m(a); return m *= v; }

template<typename T, typename U>
Matrix<T> operator/(const Matrix<T> &a, const U &v) { Matrix<T> m(a); return m /= v; }

template<typename T, typename U>
Matrix<T> operator+(const U &v, const Matrix<T> &a) { return a + v; }

template<typename T, typename U>
Matrix<T> operator*(const U &v, const Matrix<T> &a) { return a * v; }

template<typename T, typename U>
Matrix<T> operator-(const U &v, const Matrix<T> &a) {
  Matrix<T> m = make_matrix<T>(rows(a), cols(a), v);
  return m -= a;
}

template<typename T>
Matrix<T> operator^(Matrix<T> a, uint64_t p) {
  assert(rows(a) == cols(a));
  Matrix<T> res = identity_matrix<T>(rows(a));
  while (p > 0) {
    if (p & 1) {
      res *= a;
    }
    p >>= 1;
    if (p > 0) {
      a *= a;
    }
  }
  return res;
}

template<typename T>
Matrix<T> &operator^=(Matrix<T> &a, uint64_t p) {
  return a = a ^ p;
}

template<typename T>
Matrix<T> power_sum(const Matrix<T> &a, uint64_t p) {
  assert(rows(a) == cols(a));
  int n = rows(a);
  if (p == 0) {
    return make_matrix<T>(n, n);
  }
  // The upper-right block of [[A, A], [0, I]]^p is A + A^2 + ... + A^p.
  Matrix<T> block = make_matrix<T>(2 * n, 2 * n);
  for (int i = 0; i < n; i++) {
    block[i + n][i + n] = 1;
    for (int j = 0; j < n; j++) {
      block[i][j] = a[i][j];
      block[i][j + n] = a[i][j];
    }
  }
  block = block ^ p;
  Matrix<T> res = make_matrix<T>(n, n);
  for (int i = 0; i < n; i++) {
    for (int j = 0; j < n; j++) {
      res[i][j] = block[i][j + n];
    }
  }
  return res;
}
\end{lstlisting}
\Needspace*{4\baselineskip}
\textbf{Example Usage}\par
\begin{lstlisting}[style=cpp]
using namespace std;

int main() {
  using matrix = Matrix<int>;
  matrix m = make_matrix(5, 5, 10) + 10;
  vector<int> v{1, 2, 3, 4, 5};
  assert((m * v == vector<int>{300, 300, 300, 300, 300}));

  m[0][0] += 5;
  assert(m[0][0] == 25 && m[1][1] == 20);
  assert((m ^ 0) == identity_matrix(5));
  assert((identity_matrix<int>(3) ^ 5) == identity_matrix<int>(3));
  assert(power_sum(m, 3) == m + m * m + (m ^ 3));

  Matrix<double> d = make_matrix<double>(2, 2, 0.5);
  assert(rows(d) == 2 && cols(d) == 2);
  assert((d + 0.25)[0][0] == 0.75);
  return 0;
}
\end{lstlisting}
\subsection{Row Reduction}
\label{sec:6.5.2}
Converts a matrix to reduced row echelon form using Gaussian elimination to solve a system of linear equations and compute rank. Each round finds a row with a nonzero entry in the current leading column, normalizes that row, and subtracts multiples of it from every other row to clear the column. Floating-point matrices use the largest-magnitude candidate for partial pivoting, though elimination can still be prone to rounding error; use LU decomposition for greater numerical stability. Exact field types such as \inlinecode{Modular\textless{}MOD\textgreater{}} instead use any nonzero pivot and perform exact arithmetic. Such types must support construction from $0$ and $1$, equality, addition, subtraction, multiplication, and division.

\begin{compactitem}
  \item \inlinecode{row\_reduce(a)} assigns the matrix \inlinecode{a} to its reduced row echelon form, returning a reference to the modified argument itself.
  \item \inlinecode{matrix\_rank(a)} returns the rank of matrix \inlinecode{a}, i.e. the number of nonzero rows after row reduction.
  \item \inlinecode{solve\_system(a, b, \&x)} solves the system of linear equations $ax = b$ given an $m$ by $n$ matrix \inlinecode{a} over the real numbers or an exact field, and a length $m$ vector \inlinecode{b}, returning $0$ if there is one solution, $-1$ if there are zero solutions, or $-2$ if there are infinite solutions. If there is exactly one solution, then the output vector \inlinecode{x} is populated with the solution of length $n$.
\end{compactitem}

\Needspace*{3\baselineskip}
\textbf{Time Complexity}:
\begin{compactitem}
  \item $O(m \cdot n \cdot \min\left(m, n\right))$ per call to \inlinecode{row\_reduce()} and \inlinecode{matrix\_rank()}, where $m$ and $n$ are the numbers of rows and columns of \inlinecode{a}, respectively.
  \item $O(m \cdot n \cdot \min\left(m, n\right))$ per call to \inlinecode{solve\_system()}.
\end{compactitem}

\Needspace*{3\baselineskip}
\textbf{Space Complexity}:
\begin{compactitem}
  \item $O(1)$ auxiliary for \inlinecode{row\_reduce()}.
  \item $O(m \cdot n)$ auxiliary for \inlinecode{matrix\_rank()} and \inlinecode{solve\_system()}.
\end{compactitem}

\begin{lstlisting}[style=cpp]
#include <algorithm>
#include <cmath>
#include <type_traits>
#include <utility>
#include <vector>

const double EPS = 1e-9;

template<typename T>
bool rref_is_zero(const T &v) {
  if constexpr (std::is_floating_point_v<T>) {
    return std::fabs(v) < EPS;
  }
  return v == T{0};
}

template<typename Matrix>
Matrix &row_reduce(Matrix &a) {
  if (a.empty()) {
    return a;
  }
  using T = std::decay_t<decltype(a[0][0])>;
  int rows = static_cast<int>(a.size()), cols = static_cast<int>(a[0].size());
  for (int r = 0, lead = 0; r < rows && lead < cols; lead++) {
    int pivot = r;
    if constexpr (std::is_floating_point_v<T>) {
      for (int i = r + 1; i < rows; i++) {
        if (std::fabs(a[i][lead]) > std::fabs(a[pivot][lead])) {
          pivot = i;
        }
      }
    } else {
      while (pivot < rows && rref_is_zero(a[pivot][lead])) {
        pivot++;
      }
    }
    if (pivot == rows || rref_is_zero(a[pivot][lead])) {
      continue;
    }
    std::swap(a[pivot], a[r]);
    auto lv = a[r][lead];
    for (int j = 0; j < cols; j++) {
      a[r][j] /= lv;
    }
    for (int i = 0; i < rows; i++) {
      if (i != r) {
        lv = a[i][lead];
        for (int j = 0; j < cols; j++) {
          a[i][j] -= lv * a[r][j];
        }
      }
    }
    std::fill(a[r].begin(), a[r].begin() + lead, 0);
    a[r][lead] = 1;
    r++;
  }
  return a;
}

template<typename Matrix>
int matrix_rank(Matrix a) {
  row_reduce(a);
  int rank = 0;
  for (int i = 0; i < static_cast<int>(a.size()); i++) {
    for (int j = 0; j < static_cast<int>(a[i].size()); j++) {
      if (!rref_is_zero(a[i][j])) {
        rank++;
        break;
      }
    }
  }
  return rank;
}

template<typename Matrix, typename T>
int solve_system(const Matrix &a, const std::vector<T> &b, std::vector<T> *x) {
  if (x == nullptr || a.empty() || a.size() != b.size()) {
    return -1;
  }
  int rows = static_cast<int>(a.size()), cols = static_cast<int>(a[0].size());
  Matrix m(a);
  for (int i = 0; i < rows; i++) {
    m[i].push_back(b[i]);
  }
  row_reduce(m);
  int rank = 0;
  for (int i = 0; i < rows; i++) {
    int lead = -1;
    for (int j = 0; j < cols && lead < 0; j++) {
      if (!rref_is_zero(m[i][j])) {
        lead = j;
      }
    }
    if (lead < 0) {
      // A zero coefficient row with a nonzero constant means 0 = nonzero (no solution).
      if (!rref_is_zero(m[i][cols])) {
        return -1;
      }
    } else {
      rank++;
    }
  }
  // A unique solution requires a pivot in every column; fewer means free variables remain. This
  // also catches trailing free columns, which a per-row "lead > i" test alone would miss.
  if (rank < cols) {
    return -2;
  }
  x->resize(cols);
  for (int i = 0; i < cols; i++) {
    (*x)[i] = m[i][cols];
  }
  return 0;
}
\end{lstlisting}
\Needspace*{4\baselineskip}
\textbf{Example Usage}\par
\begin{lstlisting}[style=cpp]
#include <cassert>
using namespace std;

int main() {
  vector<vector<double>> a{{-1, 2, 5}, {1, 0, -6}, {-4, 2, 2}};
  vector<double> b{3, 1, -2};
  vector<double> x;
  assert(solve_system(a, b, &x) == 0);
  assert(matrix_rank(a) == 3);
  assert(matrix_rank(vector<vector<double>>{{1, 2, 3}, {2, 4, 6}, {0, 1, 1}}) == 2);
  assert(solve_system(vector<vector<double>>{{0, 0}}, vector<double>{1}, &x) == -1);
  assert(solve_system(vector<vector<double>>{{1, 0}}, vector<double>{1}, &x) == -2);
  for (int i = 0; i < static_cast<int>(a.size()); i++) {
    double sum = 0;
    for (int j = 0; j < static_cast<int>(a[i].size()); j++) {
      sum += a[i][j] * x[j];
    }
    assert(fabs(sum - b[i]) < EPS);
  }
  return 0;
}
\end{lstlisting}
\subsection{Determinant and Inverse}
\label{sec:6.5.3}
Computes the determinant and inverse of a square matrix using Gaussian elimination. The inverse of a matrix $a$ is another matrix $b$ such that $ab$ equals the identity matrix. The inverse of $a$ exists if and only if the determinant of $a$ is nonzero. In this case, $a$ is called invertible or non-singular. The determinant falls out of elimination as the product of the pivots, negated once per row swap. The inverse is found by appending the identity matrix and row-reducing the combined matrix: the operations that turn $a$ into the identity turn the identity into the inverse.

Floating-point matrices use the largest-magnitude pivot in each column, although Gaussian elimination can still suffer from rounding error; use LU decomposition for greater numerical stability. Exact field types such as \inlinecode{Modular\textless{}MOD\textgreater{}} use any nonzero pivot and perform exact arithmetic. These types must support construction from $0$ and $1$, equality, addition, subtraction, multiplication, and division. For an integer matrix whose determinant is wanted exactly, the Bareiss algorithm runs a fraction-free elimination: every intermediate value is itself a determinant of a submatrix, so the divisions are always exact and the arithmetic stays in integers.

\begin{compactitem}
  \item \inlinecode{det\_naive(a)} returns the determinant of an $n$ by $n$ matrix \inlinecode{a}, using the classic divide-and-conquer algorithm by Laplace expansions. It is division-free and computes in the matrix's element type, so an integer matrix yields an exact integer determinant; but at $O(n!)$ it is only practical for tiny $n$, so prefer \inlinecode{det\_bareiss} for exact integer determinants.
  \item \inlinecode{det(a, eps = 1e-10)} returns the determinant of an $n$ by $n$ floating-point or exact-field matrix \inlinecode{a} using Gaussian elimination. Floating-point pivots within \inlinecode{eps} of zero are treated as singular; exact fields use equality with zero.
  \item \inlinecode{det\_bareiss(a)} returns the exact determinant of an integer matrix \inlinecode{a} using fraction-free elimination, with no rounding error.
  \item \inlinecode{inverse(a, eps = 1e-10)} returns the inverse of square matrix \inlinecode{a} with floating-point or exact-field elements, or \inlinecode{std::nullopt} if the matrix is singular. Floating-point pivots within \inlinecode{eps} of zero are treated as singular.
\end{compactitem}

Overflow warning: In \inlinecode{det\_bareiss()}, stored entries are minors of \inlinecode{a}, but the products in each update can be larger. They use 128-bit intermediates when available; otherwise those products must fit in \inlinecode{int64\_t}.

\Needspace*{3\baselineskip}
\textbf{Time Complexity}:
\begin{compactitem}
  \item $O(n!)$ per call to \inlinecode{det\_naive()}, where $n$ is the dimension of the matrix.
  \item $O(n^{3})$ per call to \inlinecode{det()}, \inlinecode{det\_bareiss()}, and \inlinecode{inverse()} where $n$ is the dimension of the matrix.
\end{compactitem}

\Needspace*{3\baselineskip}
\textbf{Space Complexity}:
\begin{compactitem}
  \item $O(n)$ auxiliary stack space and $O(n^{3})$ auxiliary heap space for \inlinecode{det\_naive()}, where $n$ is the dimension of the matrix.
  \item $O(n^{2})$ auxiliary heap space for \inlinecode{det()}, \inlinecode{det\_bareiss()}, and \inlinecode{inverse()}.
\end{compactitem}

\begin{lstlisting}[style=cpp]
#include <cassert>
#include <cmath>
#include <cstdint>
#include <optional>
#include <type_traits>
#include <utility>
#include <vector>

template<typename SquareMatrix>
auto det_naive(const SquareMatrix &a) {
  int n = static_cast<int>(a.size());
  assert(n == 0 || static_cast<int>(a[0].size()) == n);
  using T = std::decay_t<decltype(a[0][0])>;
  if (n == 0) {
    return T{1};
  }
  if (n == 1) {
    return a[0][0];
  }
  if (n == 2) {
    return a[0][0] * a[1][1] - a[0][1] * a[1][0];
  }
  T res = 0;
  SquareMatrix temp(n - 1, typename SquareMatrix::value_type(n - 1));
  for (int p = 0; p < n; p++) {
    int h = 0, k = 0;
    for (int i = 1; i < n; i++) {
      for (int j = 0; j < n; j++) {
        if (j == p) {
          continue;
        }
        temp[h][k++] = a[i][j];
        if (k == n - 1) {
          h++;
          k = 0;
        }
      }
    }
    res += (p % 2 == 0 ? 1 : -1) * a[0][p] * det_naive(temp);
  }
  return res;
}

template<typename SquareMatrix>
int elimination_pivot(const SquareMatrix &a, int row, int col, double eps) {
  using T = std::decay_t<decltype(a[0][0])>;
  int n = static_cast<int>(a.size());
  if constexpr (std::is_floating_point_v<T>) {
    int pivot = row;
    for (int r = row + 1; r < n; r++) {
      if (std::fabs(a[r][col]) > std::fabs(a[pivot][col])) {
        pivot = r;
      }
    }
    return std::fabs(a[pivot][col]) <= eps ? -1 : pivot;
  }
  for (int r = row; r < n; r++) {
    if (a[r][col] != T{0}) {
      return r;
    }
  }
  return -1;
}

template<typename SquareMatrix>
auto det(const SquareMatrix &a, double eps = 1e-10) {
  int n = static_cast<int>(a.size());
  assert(n == 0 || static_cast<int>(a[0].size()) == n);
  using T = std::decay_t<decltype(a[0][0])>;
  SquareMatrix b(a);
  T res = 1;
  for (int i = 0; i < n; i++) {
    int p = elimination_pivot(b, i, i, eps);
    if (p == -1) {
      return T{0};
    }
    if (p != i) {
      std::swap(b[p], b[i]);
      res = T{0} - res;
    }
    res *= b[i][i];
    for (int j = i + 1; j < n; j++) {
      T z = b[j][i] / b[i][i];
      for (int k = i; k < n; k++) {
        b[j][k] -= z * b[i][k];
      }
    }
  }
  return res;
}

int64_t det_bareiss(std::vector<std::vector<int64_t>> a) {
  int n = static_cast<int>(a.size());
  assert(n == 0 || static_cast<int>(a[0].size()) == n);
  int64_t prev = 1, sign = 1;
  for (int k = 0; k < n; k++) {
    if (a[k][k] == 0) {  // Swap in a nonzero pivot; each swap flips the sign.
      int p = -1;
      for (int r = k + 1; r < n; r++) {
        if (a[r][k] != 0) {
          p = r;
          break;
        }
      }
      if (p == -1) {
        return 0;
      }
      std::swap(a[k], a[p]);
      sign = -sign;
    }
    for (int i = k + 1; i < n; i++) {
      for (int j = k + 1; j < n; j++) {
#if defined(__SIZEOF_INT128__)
        __extension__ typedef __int128 int128_t;
        int128_t numerator =
            static_cast<int128_t>(a[i][j]) * a[k][k] - static_cast<int128_t>(a[i][k]) * a[k][j];
        a[i][j] = static_cast<int64_t>(numerator / prev);
#else
        a[i][j] = (a[i][j] * a[k][k] - a[i][k] * a[k][j]) / prev;  // Overflow warning.
#endif
      }
    }
    prev = a[k][k];
  }
  return n == 0 ? 1 : sign * a[n - 1][n - 1];
}

template<typename SquareMatrix>
std::optional<SquareMatrix> inverse(SquareMatrix a, double eps = 1e-10) {
  int n = static_cast<int>(a.size());
  assert(n == 0 || static_cast<int>(a[0].size()) == n);
  using T = std::decay_t<decltype(a[0][0])>;
  for (int i = 0; i < n; i++) {
    a[i].resize(2 * n);
    for (int j = n; j < n * 2; j++) {
      a[i][j] = (i == j - n ? 1 : 0);
    }
  }
  for (int i = 0; i < n; i++) {
    int p = elimination_pivot(a, i, i, eps);
    if (p == -1) {
      return std::nullopt;
    }
    std::swap(a[p], a[i]);
    T pivot = a[i][i];
    for (int j = i; j < n * 2; j++) {
      a[i][j] /= pivot;
    }
    for (int j = 0; j < n; j++) {
      if (i != j) {
        T factor = a[j][i];
        for (int k = 0; k < n * 2; k++) {
          a[j][k] -= factor * a[i][k];
        }
      }
    }
  }
  for (int i = 0; i < n; i++) {
    a[i].erase(a[i].begin(), a[i].begin() + n);
  }
  return a;
}
\end{lstlisting}
\Needspace*{4\baselineskip}
\textbf{Example Usage}\par
\begin{lstlisting}[style=cpp]
#include <cassert>
using namespace std;

int main() {
  vector<vector<double>> a{{6, 1, 1}, {4, -2, 5}, {2, 8, 7}};
  int n = static_cast<int>(a.size());
  vector<vector<double>> res(n, vector<double>(n));
  assert(fabs(det(a) - det_naive(a)) < 1e-10);

  // Bareiss gives the determinant of an integer matrix exactly, with no rounding.
  vector<vector<int64_t>> ai{{6, 1, 1}, {4, -2, 5}, {2, 8, 7}};
  assert(det_naive(ai) == -306);  // Division-free, so also exact in the int64_t element type.
  assert(det_naive(vector<vector<int64_t>>{}) == 1);
  assert(det_bareiss(ai) == -306);
  assert(det_bareiss({{2, 0, 0}, {0, 3, 0}, {0, 0, 5}}) == 30);
  assert(det_bareiss({{1, 2}, {2, 4}}) == 0);  // Singular.
  auto inv = inverse(a);
  assert(inv);
  assert(!inverse(vector<vector<double>>{{1, 2}, {2, 4}}));
  for (int i = 0; i < n; i++) {
    for (int j = 0; j < n; j++) {
      for (int k = 0; k < n; k++) {
        res[i][j] += a[i][k] * (*inv)[k][j];
      }
    }
  }
  for (int i = 0; i < n; i++) {
    for (int j = 0; j < n; j++) {
      assert(fabs(res[i][j] - (i == j ? 1 : 0)) < 1e-10);
    }
  }
  return 0;
}
\end{lstlisting}
\subsection{LU Decomposition}
\label{sec:6.5.4}
The LU decomposition of a matrix $A$ with row-partial pivoting is a factorization of $A$ (after some rows are possibly permuted by a permutation matrix $p$) as a product of a lower triangular matrix $L$ and an upper triangular matrix $U$. This factorization can be used to tackle many common problems in linear algebra such as solving systems of linear equations and computing determinants. An improvement on basic row reduction, LU decomposition by row-partial pivoting keeps the relative magnitude of matrix values small, thus reducing the relative error due to rounding in computed solutions. These routines require floating-point matrix elements; use \secref{6.5.3} for exact-field or integer determinant calculations.

\begin{compactitem}
  \item \inlinecode{lu\_decompose(a, \&perm, eps = 1e-10)} assigns floating-point matrix \inlinecode{a} to merged LU decomposition matrix \inlinecode{lu}; it returns $0$ for even row-swap parity, $1$ for odd parity, or $-1$ for a degenerate matrix (i.e. singular for square matrices). The merged matrix stores \inlinecode{l[i][j]} in \inlinecode{lu[i][j]} when \inlinecode{i \textgreater{} j}, and \inlinecode{u[i][j]} otherwise. The diagonal entries \inlinecode{l[i][i]} are always $1$, so they are not explicitly stored. Access general entries of the lower and upper triangular matrices via \inlinecode{getl(lu, i, j)} and \inlinecode{getu(lu, i, j)}. Optional output vector \inlinecode{perm} represents $p$: \inlinecode{perm[i]} is the only column containing $1$ in row \inlinecode{i}. Left-multiplying \inlinecode{a} by this permutation matrix gives the product of the separate lower and upper triangular matrices, not the merged storage matrix \inlinecode{lu} itself.
  \item \inlinecode{solve\_system(a, b, \&x, eps = 1e-10)} solves the system of linear equations $Ax = b$ given an $m$ by $n$ floating-point matrix \inlinecode{a} and a length $m$ vector \inlinecode{b}, returning $0$ if there is one solution or $-1$ if there are zero or infinite solutions. If there is exactly one solution, then the output vector \inlinecode{x} is populated with the solution of length $n$; otherwise, \inlinecode{x} is unchanged.
  \item \inlinecode{det(a)} returns the determinant of an $n$ by $n$ floating-point matrix \inlinecode{a} using LU decomposition.
  \item \inlinecode{inverse(a)} returns the inverse of the $n$ by $n$ floating-point matrix \inlinecode{a}, or \inlinecode{std::nullopt} if \inlinecode{a} is singular.
\end{compactitem}

\Needspace*{3\baselineskip}
\textbf{Time Complexity}:
\begin{compactitem}
  \item $O(m \cdot n \cdot \min\left(m, n\right))$ per call to \inlinecode{lu\_decompose(a)}, where $m$ and $n$ are the numbers of rows and columns of \inlinecode{a}, respectively.
  \item $O(m \cdot n^{2})$ per call to \inlinecode{solve\_system()}, which requires $m \geq n$.
  \item $O(n^{3})$ per call to \inlinecode{det(a)} and \inlinecode{inverse(a)}, where $n$ is the length of square matrix \inlinecode{a}.
\end{compactitem}

\Needspace*{3\baselineskip}
\textbf{Space Complexity}:
\begin{compactitem}
  \item $O(1)$ auxiliary for \inlinecode{lu\_decompose()}.
  \item $O(n^{2})$ auxiliary for \inlinecode{det()} and \inlinecode{inverse()}.
  \item $O(m \cdot n)$ auxiliary for \inlinecode{solve\_system()}.
\end{compactitem}

\begin{lstlisting}[style=cpp]
#include <algorithm>
#include <cassert>
#include <cmath>
#include <numeric>
#include <optional>
#include <type_traits>
#include <vector>

template<typename Matrix>
int lu_decompose(Matrix &a, std::vector<int> *perm = nullptr, double eps = 1e-10) {
  int rows = static_cast<int>(a.size());
  int cols = a.empty() ? 0 : static_cast<int>(a[0].size());
  int parity = 0;
  if (perm != nullptr) {
    perm->resize(rows);
    std::iota(perm->begin(), perm->end(), 0);
  }
  for (int i = 0; i < rows && i < cols; i++) {
    int pi = i;
    for (int k = i + 1; k < rows; k++) {
      if (std::fabs(a[k][i]) > std::fabs(a[pi][i])) {
        pi = k;
      }
    }
    if (std::fabs(a[pi][i]) <= eps) {
      return -1;
    }
    if (pi != i) {
      if (perm != nullptr) {
        std::iter_swap(perm->begin() + i, perm->begin() + pi);
      }
      std::iter_swap(a.begin() + i, a.begin() + pi);
      parity = 1 - parity;
    }
    for (int j = i + 1; j < rows; j++) {
      a[j][i] /= a[i][i];
      for (int k = i + 1; k < cols; k++) {
        a[j][k] -= a[j][i] * a[i][k];
      }
    }
  }
  return parity;
}

template<typename Matrix>
auto getl(const Matrix &lu, int i, int j) {
  using T = std::decay_t<decltype(lu[0][0])>;
  return i > j ? lu[i][j] : T(i == j);
}

template<typename Matrix>
auto getu(const Matrix &lu, int i, int j) {
  using T = std::decay_t<decltype(lu[0][0])>;
  return i <= j ? lu[i][j] : T{0};
}

template<typename Matrix, typename T>
int solve_system(const Matrix &a, const std::vector<T> &b, std::vector<T> *x, double eps = 1e-10) {
  int rows = static_cast<int>(a.size());
  int cols = a.empty() ? 0 : static_cast<int>(a[0].size());
  if (x == nullptr || a.empty() || a.size() != b.size() || rows < cols) {
    return -1;
  }
  Matrix lu(a);
  std::vector<int> perm;
  int status = lu_decompose(lu, &perm, eps);
  if (status < 0) {
    return status;
  }
  std::vector<T> solution(cols);
  for (int i = 0; i < cols; i++) {
    solution[i] = b[perm[i]];
    for (int k = 0; k < i; k++) {
      solution[i] -= getl(lu, i, k) * solution[k];
    }
  }
  for (int i = cols - 1; i >= 0; i--) {
    for (int k = i + 1; k < cols; k++) {
      solution[i] -= getu(lu, i, k) * solution[k];
    }
    solution[i] /= getu(lu, i, i);
  }
  for (int i = 0; i < rows; i++) {
    T val = 0;
    for (int j = 0; j < cols; j++) {
      val += a[i][j] * solution[j];
    }
    // Mixed absolute/relative tolerance: dividing by b[i] alone would skip the check for negative
    // b[i] (ratio goes negative) and divide by zero when b[i] == 0.
    if (std::fabs(val - b[i]) > eps * (1.0 + std::fabs(b[i]))) {
      return -1;
    }
  }
  x->swap(solution);
  return 0;
}

template<typename Matrix>
auto det(const Matrix &a) {
  int n = static_cast<int>(a.size());
  assert(n == 0 || static_cast<int>(a[0].size()) == n);
  using T = std::decay_t<decltype(a[0][0])>;
  Matrix lu(a);
  int status = lu_decompose(lu);
  if (status < 0) {
    return T{0};
  }
  T res = 1;
  for (int i = 0; i < n; i++) {
    res *= lu[i][i];
  }
  return status == 0 ? res : -res;
}

template<typename SquareMatrix>
std::optional<SquareMatrix> inverse(SquareMatrix a) {
  int n = static_cast<int>(a.size());
  assert(n == 0 || static_cast<int>(a[0].size()) == n);
  std::vector<int> perm;
  int status = lu_decompose(a, &perm);
  if (status < 0) {
    return std::nullopt;
  }
  SquareMatrix ia(n, typename SquareMatrix::value_type(n, 0));
  for (int j = 0; j < n; j++) {
    for (int i = 0; i < n; i++) {
      if (perm[i] == j) {
        ia[i][j] = 1.0;
      } else {
        ia[i][j] = 0.0;
      }
      for (int k = 0; k < i; k++) {
        ia[i][j] -= getl(a, i, k) * ia[k][j];
      }
    }
    for (int i = n - 1; i >= 0; i--) {
      for (int k = i + 1; k < n; k++) {
        ia[i][j] -= getu(a, i, k) * ia[k][j];
      }
      ia[i][j] /= getu(a, i, i);
    }
  }
  return ia;
}
\end{lstlisting}
\Needspace*{4\baselineskip}
\textbf{Example Usage}\par
\begin{lstlisting}[style=cpp]
#include <cassert>
using namespace std;

bool EQ(double a, double b) {
  return fabs(a - b) < 1e-9;
}

int main() {
  {  // Solve a system.
    vector<vector<double>> a{{-1, 2, 5}, {1, 0, -6}, {-4, 2, 2}};
    vector<double> b{3, 1, -2};
    vector<double> x;
    assert(solve_system(a, b, &x) == 0);
    for (int i = 0; i < static_cast<int>(a.size()); i++) {
      double sum = 0;
      for (int j = 0; j < static_cast<int>(a[0].size()); j++) {
        sum += a[i][j] * x[j];
      }
      assert(EQ(sum, b[i]));
    }
  }
  {  // Find the determinant.
    vector<vector<double>> a{{1, 3, 5}, {2, 4, 7}, {1, 1, 0}};
    assert(EQ(det(a), 4));
  }
  {  // Find the inverse.
    vector<vector<double>> a{{6, 1, 1}, {4, -2, 5}, {2, 8, 7}};
    auto inv = inverse(a);
    int n = static_cast<int>(a.size());
    vector<vector<double>> res(n, vector<double>(n));
    assert(inv);
    for (int i = 0; i < n; i++) {
      for (int j = 0; j < n; j++) {
        for (int k = 0; k < n; k++) {
          res[i][j] += a[i][k] * (*inv)[k][j];
        }
      }
    }
    for (int i = 0; i < n; i++) {
      for (int j = 0; j < n; j++) {
        assert(EQ(res[i][j], i == j ? 1 : 0));
      }
    }
  }
  {  // Singular matrices have no inverse.
    vector<vector<double>> singular{{1, 2}, {2, 4}};
    assert(!inverse(singular));
    // Failed operations leave their output arguments unchanged.
    vector<vector<double>> inconsistent{{1}, {1}};
    vector<double> b{1, 2}, x{99, 100};
    assert(solve_system(inconsistent, b, &x) == -1);
    assert((x == vector<double>{99, 100}));
  }
  return 0;
}
\end{lstlisting}
\subsection{QR Decomposition and Least Squares}
\label{sec:6.5.5}
The QR decomposition factors an $m$ by $n$ matrix $A$ with $m \geq n$ into $A = QR$, where $Q$ has orthonormal columns and $R$ is upper triangular. Orthonormal columns are what make the factorization useful: $Q$ preserves lengths, so multiplying by $Q^T$ rewrites a problem in a new basis without distorting distances, and a triangular system is solved by back substitution.

This implementation uses Householder reflections. A Householder reflection is the mirror across the hyperplane orthogonal to a chosen vector $v$, written $I - 2vv^T/(v^Tv)$, and $v$ is chosen so that the reflection maps the part of a column below the diagonal onto the axis, zeroing it in one step. Applying one reflection per column leaves $R$, and accumulating the reflections gives $Q$. Reflecting is stable because it never scales anything: the alternative of orthogonalizing column by column with Gram-Schmidt loses orthogonality as the columns approach linear dependence, whereas reflections keep it to within rounding error.

The overdetermined system $Ax = b$ usually has no solution, so least squares asks for the $x$ whose residual is shortest. Because $Q$ preserves lengths, $\|Ax - b\|$ equals $\|Rx - Q^Tb\|$, and the entries of $Q^Tb$ below row $n$ cannot be changed by any $x$; minimizing therefore means solving the square triangular system formed by the first $n$ rows, which back substitution does directly.

\begin{compactitem}
  \item \inlinecode{qr\_decompose(a, \&q, \&r)} factors the $m$ by $n$ matrix \inlinecode{a} with $m \geq n$ into an $m$ by $m$ orthogonal matrix \inlinecode{q} and an $m$ by $n$ upper triangular matrix \inlinecode{r}, whose product is \inlinecode{a}.
  \item \inlinecode{least\_squares(a, b, eps = 1e-10)} returns the vector $x$ minimizing the Euclidean norm of ${\inlinecodemath{a}}x - {\inlinecodemath{b}}$, or \inlinecode{std::nullopt} if the columns of \inlinecode{a} are linearly dependent to within \inlinecode{eps}. The matrix \inlinecode{a} must have at least as many rows as columns, and \inlinecode{b} must have one entry per row.
\end{compactitem}

Fitting any linear model reduces to this: one row per data point, one column per basis function. Solving the normal equations $A^TAx = A^Tb$ by the row reduction of section \secref{6.5.2} gives the same answer but squares the condition number, losing roughly twice as many digits.

\Needspace*{3\baselineskip}
\textbf{Time Complexity}:
\begin{compactitem}
  \item $O(m^{2} \cdot n)$ per call to \inlinecode{qr\_decompose()}, where $m$ and $n$ are the numbers of rows and columns.
  \item $O(m \cdot n^{2})$ per call to \inlinecode{least\_squares()}, since it never forms \inlinecode{q} explicitly.
\end{compactitem}

\Needspace*{3\baselineskip}
\textbf{Space Complexity}:
\begin{compactitem}
  \item $O(m^{2} + m \cdot n)$ for the matrices returned by \inlinecode{qr\_decompose()}.
  \item $O(m \cdot n)$ auxiliary and $O(n)$ for the vector returned by \inlinecode{least\_squares()}.
\end{compactitem}

\begin{lstlisting}[style=cpp]
#include <cassert>
#include <cmath>
#include <optional>
#include <vector>

using Matrix = std::vector<std::vector<double>>;

// Reflects the trailing rows of the given column onto the axis, returning the reflection vector.
std::vector<double> householder_vector(const Matrix &a, int col) {
  int m = static_cast<int>(a.size());
  std::vector<double> v(m);
  double norm = 0;
  for (int i = col; i < m; i++) {
    v[i] = a[i][col];
    norm += v[i] * v[i];
  }
  norm = std::sqrt(norm);
  // Reflect away from a[col][col] so that the subtraction below never cancels.
  v[col] += a[col][col] >= 0 ? norm : -norm;
  return v;
}

void qr_decompose(const Matrix &a, Matrix *q, Matrix *r) {
  int m = static_cast<int>(a.size()), n = a.empty() ? 0 : static_cast<int>(a[0].size());
  assert(m >= n && q != nullptr && r != nullptr);
  *r = a;
  q->assign(m, std::vector<double>(m));
  for (int i = 0; i < m; i++) {
    (*q)[i][i] = 1;
  }
  for (int col = 0; col < n; col++) {
    std::vector<double> v = householder_vector(*r, col);
    double vv = 0;
    for (int i = col; i < m; i++) {
      vv += v[i] * v[i];
    }
    if (vv == 0) {
      continue;  // The column is already zero below the diagonal.
    }
    for (int j = 0; j < n; j++) {  // Apply the reflection to R from the left.
      double dot = 0;
      for (int i = col; i < m; i++) {
        dot += v[i] * (*r)[i][j];
      }
      for (int i = col; i < m; i++) {
        (*r)[i][j] -= 2 * dot / vv * v[i];
      }
    }
    for (int j = 0; j < m; j++) {  // Accumulate the reflection into Q from the right.
      double dot = 0;
      for (int i = col; i < m; i++) {
        dot += v[i] * (*q)[j][i];
      }
      for (int i = col; i < m; i++) {
        (*q)[j][i] -= 2 * dot / vv * v[i];
      }
    }
  }
}

std::optional<std::vector<double>> least_squares(
    const Matrix &a, const std::vector<double> &b, double eps = 1e-10
) {
  int m = static_cast<int>(a.size()), n = a.empty() ? 0 : static_cast<int>(a[0].size());
  assert(m >= n && static_cast<int>(b.size()) == m);
  Matrix r = a;
  std::vector<double> rhs = b;
  for (int col = 0; col < n; col++) {
    std::vector<double> v = householder_vector(r, col);
    double vv = 0;
    for (int i = col; i < m; i++) {
      vv += v[i] * v[i];
    }
    if (vv == 0) {
      continue;
    }
    for (int j = col; j < n; j++) {
      double dot = 0;
      for (int i = col; i < m; i++) {
        dot += v[i] * r[i][j];
      }
      for (int i = col; i < m; i++) {
        r[i][j] -= 2 * dot / vv * v[i];
      }
    }
    double dot = 0;  // The same reflection applied to the right-hand side.
    for (int i = col; i < m; i++) {
      dot += v[i] * rhs[i];
    }
    for (int i = col; i < m; i++) {
      rhs[i] -= 2 * dot / vv * v[i];
    }
  }
  std::vector<double> x(n);
  for (int i = n - 1; i >= 0; i--) {  // Back substitution over the triangular top block.
    if (std::fabs(r[i][i]) <= eps) {
      return std::nullopt;
    }
    double sum = rhs[i];
    for (int j = i + 1; j < n; j++) {
      sum -= r[i][j] * x[j];
    }
    x[i] = sum / r[i][i];
  }
  return x;
}
\end{lstlisting}
\Needspace*{4\baselineskip}
\textbf{Example Usage}\par
\begin{lstlisting}[style=cpp]
#include <cmath>
using namespace std;

bool EQ(double a, double b) {
  return fabs(a - b) < 1e-9;
}

int main() {
  Matrix a{{12, -51, 4}, {6, 167, -68}, {-4, 24, -41}};
  Matrix q, r;
  qr_decompose(a, &q, &r);
  for (int i = 0; i < 3; i++) {
    for (int j = 0; j < 3; j++) {
      double product = 0, orthogonal = 0;
      for (int k = 0; k < 3; k++) {
        product += q[i][k] * r[k][j];
        orthogonal += q[k][i] * q[k][j];
      }
      assert(EQ(product, a[i][j]));                // The factors multiply back to a.
      assert(EQ(orthogonal, i == j ? 1.0 : 0.0));  // The columns of q are orthonormal.
      assert(i <= j || EQ(r[i][j], 0.0));          // r is upper triangular.
    }
  }

  // An exactly solvable square system.
  auto exact = least_squares(Matrix{{2, 1}, {1, 3}}, vector<double>{5, 10});
  assert(exact.has_value() && EQ((*exact)[0], 1.0) && EQ((*exact)[1], 3.0));

  // Fit y = c0 + c1*x to four collinear points, then to points with one outlier.
  Matrix design{{1, 0}, {1, 1}, {1, 2}, {1, 3}};
  auto line = least_squares(design, vector<double>{1, 3, 5, 7});
  assert(line.has_value() && EQ((*line)[0], 1.0) && EQ((*line)[1], 2.0));
  auto noisy = least_squares(design, vector<double>{1, 3, 5, 8});
  assert(noisy.has_value() && EQ((*noisy)[0], 0.8) && EQ((*noisy)[1], 2.3));

  // Dependent columns have no unique least-squares solution.
  assert(!least_squares(Matrix{{1, 2}, {2, 4}, {3, 6}}, vector<double>{1, 2, 3}).has_value());
  return 0;
}
\end{lstlisting}
\subsection{Characteristic Polynomial}
\label{sec:6.5.6}
Computes the characteristic polynomial of a square matrix over a prime field: the coefficients of $\det(xI - A)$, a monic degree-$n$ polynomial whose roots are the eigenvalues of $A$. Its constant term recovers the determinant (up to sign), its second-highest coefficient is the negated trace, and evaluating it at a scalar gives $\det(xI - A)$ for that $x$. It is the standard tool behind eigenvalue counting, matrix-power recurrences via Cayley-Hamilton, and problems that reduce to a determinant in an unknown.

A direct expansion would cost a determinant per coefficient. Instead the matrix is first reduced to upper Hessenberg form, where every entry more than one place below the diagonal is zero. This reduction uses similarity transforms (a row operation paired with the inverse column operation), which leave the characteristic polynomial unchanged. The characteristic polynomial of a Hessenberg matrix then follows from a short recurrence over its leading principal submatrices, adding one row and column at a time, for an overall cost of $O(n^{3})$. The arithmetic is exact modulo a prime, since the Hessenberg reduction divides by pivots and is not numerically stable over the reals.

\begin{compactitem}
  \item \inlinecode{characteristic\_polynomial(a)} returns a vector $c$ of length $n + 1$ with $\det(xI - {\inlinecodemath{a}}) = c_0 + c_1 x + ... + c_n x^n$ modulo \inlinecode{MOD}, where $c_n = 1$. The matrix \inlinecode{a} must be square; its entries are reduced modulo \inlinecode{MOD} internally.
\end{compactitem}

\textbf{Time Complexity}: $O(n^{3})$ per call, where $n$ is the dimension of the matrix.

\textbf{Space Complexity}: $O(n^{2})$ auxiliary.

\begin{lstlisting}[style=cpp]
#include <cassert>
#include <cstdint>
#include <utility>
#include <vector>

const int64_t MOD = 998244353;

int64_t powmod(int64_t x, int64_t n, int64_t m) {
  int64_t res = 1 % m;
  for (x %= m; n > 0; n >>= 1) {
    if (n & 1) {
      res = res * x % m;
    }
    x = x * x % m;
  }
  return res;
}

void to_hessenberg(std::vector<std::vector<int64_t>> &m) {
  int n = static_cast<int>(m.size());
  for (int j = 0; j + 2 < n; j++) {
    int pivot = -1;
    for (int i = j + 1; i < n; i++) {
      if (m[i][j] != 0) {
        pivot = i;
        break;
      }
    }
    if (pivot == -1) {
      continue;
    }
    if (pivot != j + 1) {  // Move the pivot to the subdiagonal via a similarity swap.
      std::swap(m[pivot], m[j + 1]);
      for (int r = 0; r < n; r++) {
        std::swap(m[r][pivot], m[r][j + 1]);
      }
    }
    int64_t pinv = powmod(m[j + 1][j], MOD - 2, MOD);
    for (int k = j + 2; k < n; k++) {
      int64_t u = m[k][j] * pinv % MOD;
      if (u == 0) {
        continue;
      }
      for (int c = 0; c < n; c++) {  // Row k -= u * row (j + 1).
        m[k][c] = (m[k][c] - u * m[j + 1][c] % MOD + MOD) % MOD;
      }
      for (int r = 0; r < n; r++) {  // Inverse column operation keeps it a similarity.
        m[r][j + 1] = (m[r][j + 1] + u * m[r][k]) % MOD;
      }
    }
  }
}

std::vector<int64_t> characteristic_polynomial(std::vector<std::vector<int64_t>> a) {
  int n = static_cast<int>(a.size());
  assert(n == 0 || static_cast<int>(a[0].size()) == n);
  for (auto &row : a) {
    for (int64_t &x : row) {
      x = (x % MOD + MOD) % MOD;
    }
  }
  to_hessenberg(a);
  std::vector<std::vector<int64_t>> p(n + 1);
  p[0] = {1};
  for (int k = 0; k < n; k++) {
    p[k + 1].assign(k + 2, 0);
    for (int i = 0; i < static_cast<int>(p[k].size()); i++) {
      p[k + 1][i + 1] = (p[k + 1][i + 1] + p[k][i]) % MOD;                // x * p[k].
      p[k + 1][i] = (p[k + 1][i] - a[k][k] * p[k][i] % MOD + MOD) % MOD;  // -a[k][k] * p[k].
    }
    int64_t t = 1;
    for (int i = 0; i < k; i++) {
      t = t * a[k - i][k - i - 1] % MOD;  // Running product of subdiagonal entries.
      int64_t coeff = t * a[k - i - 1][k] % MOD;
      for (int j = 0; j < static_cast<int>(p[k - i - 1].size()); j++) {
        p[k + 1][j] = (p[k + 1][j] - coeff * p[k - i - 1][j] % MOD + MOD) % MOD;
      }
    }
  }
  return p[n];
}
\end{lstlisting}
\Needspace*{4\baselineskip}
\textbf{Example Usage}\par
\begin{lstlisting}[style=cpp]
#include <cassert>
using namespace std;

int main() {
  // [[1, 2], [3, 4]]: det(xI - A) = x^2 - 5x - 2.
  vector<vector<int64_t>> a{{1, 2}, {3, 4}};
  assert((characteristic_polynomial(a) == vector<int64_t>{MOD - 2, MOD - 5, 1}));  // -2, -5, 1.

  // det(xI - B) = (x - 3)(x^2 - 7x + 6) = x^3 - 10x^2 + 27x - 18.
  vector<vector<int64_t>> b{{2, 0, 1}, {0, 3, 0}, {4, 1, 5}};
  assert((characteristic_polynomial(b) == vector<int64_t>{MOD - 18, 27, MOD - 10, 1}));
  return 0;
}
\end{lstlisting}
\subsection{Eigenvalues (Jacobi)}
\label{sec:6.5.7}
An eigenvector of a square matrix is a direction the matrix does not rotate, only rescales, and the scale factor is its eigenvalue. For a real symmetric matrix, the spectral theorem guarantees that a full set of real eigenvalues exists and that eigenvectors for distinct eigenvalues are orthogonal, so the matrix is diagonalized by an orthogonal change of basis.

The Jacobi eigenvalue algorithm builds that basis one rotation at a time. It repeatedly picks the largest off-diagonal entry and applies a rotation in the plane of that entry's row and column, choosing the angle that sets the entry to zero. A rotation is a similarity transform, so it preserves the eigenvalues, and it also preserves the sum of squares of all entries; since zeroing an entry removes its square from the off-diagonal total and moves it onto the diagonal, every rotation strictly reduces the off-diagonal weight. Earlier zeros are partially refilled by later rotations, but the total decays geometrically, and the matrix converges to a diagonal one whose entries are the eigenvalues. Accumulating the rotations gives the eigenvectors.

\begin{compactitem}
  \item \inlinecode{jacobi\_eigen(a, eps = 1e-12, iterations = 100)} returns the pair (\inlinecode{values}, \inlinecode{vectors}) for a real symmetric $n$ by $n$ matrix \inlinecode{a}. The eigenvalues in \inlinecode{values} are sorted in decreasing order, and \inlinecode{vectors[i]} is a unit eigenvector for \inlinecode{values[i]}, so the eigenvectors are rows rather than columns. Each iteration rotates away the currently largest off-diagonal entry, stopping early once every such entry is smaller than \inlinecode{eps}; if the \inlinecode{iterations} limit is reached first, the current approximation is returned. \inlinecode{iterations} must be nonnegative. The matrix must be symmetric, which is not checked; an unsymmetric one silently produces nonsense.
\end{compactitem}

A general matrix instead needs the roots of the characteristic polynomial of section \secref{6.5.6}, found including complex ones by the Ehrlich-Aberth iteration of section \secref{5.4.5}, though accuracy degrades with degree since polynomial roots are far more sensitive to coefficients than eigenvalues are to the matrix. When only the dominant eigenvalue is wanted, power iteration finds it with no decomposition at all, by repeatedly multiplying a random vector and renormalizing.

\textbf{Time Complexity}: $O(I \cdot n^{2})$ per call for at most $I$ iterations, where $n$ is the matrix dimension. Each iteration searches $O(n^{2})$ entries for the largest off-diagonal one and applies one $O(n)$ rotation.

\textbf{Space Complexity}: $O(n^{2})$ auxiliary, and $O(n^{2})$ for the returned eigenvectors.

\begin{lstlisting}[style=cpp]
#include <algorithm>
#include <cassert>
#include <cmath>
#include <numeric>
#include <utility>
#include <vector>

using Matrix = std::vector<std::vector<double>>;

std::pair<std::vector<double>, Matrix> jacobi_eigen(
    Matrix a, double eps = 1e-12, int iterations = 100
) {
  int n = static_cast<int>(a.size());
  assert((n == 0 || static_cast<int>(a[0].size()) == n) && iterations >= 0);
  Matrix vectors(n, std::vector<double>(n));
  for (int i = 0; i < n; i++) {
    vectors[i][i] = 1;
  }
  while (iterations--) {
    int p = 0, q = 1;
    double largest = 0;
    for (int i = 0; i < n; i++) {
      for (int j = i + 1; j < n; j++) {
        if (std::fabs(a[i][j]) > largest) {
          largest = std::fabs(a[i][j]);
          p = i;
          q = j;
        }
      }
    }
    if (n < 2 || largest <= eps) {
      break;
    }
    // Choose the rotation angle that zeroes a[p][q], using the numerically safe form of
    // tan(theta) rather than computing the angle itself.
    double theta = (a[q][q] - a[p][p]) / (2 * a[p][q]);
    double t = (theta >= 0 ? 1 : -1) / (std::fabs(theta) + std::sqrt(theta * theta + 1));
    double c = 1 / std::sqrt(t * t + 1), s = t * c;
    for (int k = 0; k < n; k++) {
      double akp = a[k][p], akq = a[k][q];
      a[k][p] = c * akp - s * akq;
      a[k][q] = s * akp + c * akq;
    }
    for (int k = 0; k < n; k++) {
      double apk = a[p][k], aqk = a[q][k];
      a[p][k] = c * apk - s * aqk;
      a[q][k] = s * apk + c * aqk;
      double vpk = vectors[p][k], vqk = vectors[q][k];
      vectors[p][k] = c * vpk - s * vqk;
      vectors[q][k] = s * vpk + c * vqk;
    }
  }
  std::vector<int> order(n);
  std::iota(order.begin(), order.end(), 0);
  std::sort(order.begin(), order.end(), [&](int x, int y) { return a[x][x] > a[y][y]; });
  std::vector<double> values(n);
  Matrix sorted(n);
  for (int i = 0; i < n; i++) {
    values[i] = a[order[i]][order[i]];
    sorted[i] = vectors[order[i]];
  }
  return {values, sorted};
}
\end{lstlisting}
\Needspace*{4\baselineskip}
\textbf{Example Usage}\par
\begin{lstlisting}[style=cpp]
#include <cassert>
#include <cmath>
using namespace std;

bool EQ(double a, double b) {
  return fabs(a - b) < 1e-9;
}

int main() {
  // A symmetric matrix with eigenvalues 3, 2, and 1.
  Matrix a{{2, 0, -1}, {0, 2, 0}, {-1, 0, 2}};
  auto [values, vectors] = jacobi_eigen(a);
  assert(EQ(values[0], 3) && EQ(values[1], 2) && EQ(values[2], 1));

  for (int i = 0; i < 3; i++) {
    double norm = 0;
    for (int k = 0; k < 3; k++) {
      norm += vectors[i][k] * vectors[i][k];
    }
    assert(EQ(norm, 1));  // Eigenvectors come back normalized.
    for (int r = 0; r < 3; r++) {
      double image = 0;
      for (int k = 0; k < 3; k++) {
        image += a[r][k] * vectors[i][k];
      }
      assert(EQ(image, values[i] * vectors[i][r]));  // The defining equation a*v = lambda*v.
    }
  }

  // A diagonal matrix is already solved, and a 1-by-1 matrix is its own eigenvalue.
  auto diagonal = jacobi_eigen(Matrix{{5, 0}, {0, -4}});
  assert(EQ(diagonal.first[0], 5) && EQ(diagonal.first[1], -4));
  assert(EQ(jacobi_eigen(Matrix{{7}}).first[0], 7));
  return 0;
}
\end{lstlisting}
\subsection{Sparse Matrix}
\label{sec:6.5.8}
Provides a sparse matrix type together with arithmetic and elimination routines. A sparse matrix stores only nonzero entries, which is useful when the matrix is large but most values are zero.

The matrix keeps both row maps and column maps in sync: \inlinecode{row(i)} can iterate all nonzeros in row \inlinecode{i}, while \inlinecode{col(j)} can iterate all nonzeros in column \inlinecode{j}. That bidirectional storage is especially useful for sparse elimination or graph-like matrix operations, where both row updates and column pivot lookups are common.

The class treats \inlinecode{T\{\}} as additive zero and requires \inlinecode{value == T\{\}} to be a valid zero test.

\begin{compactitem}
  \item \inlinecode{SparseMatrix\textless{}T\textgreater{}(rows, cols)} constructs a \inlinecode{rows} by \inlinecode{cols} sparse matrix.
  \item \inlinecode{num\_rows()} returns the number of rows.
  \item \inlinecode{num\_cols()} returns the number of columns.
  \item \inlinecode{nonzeros()} returns the number of stored nonzero entries.
  \item \inlinecode{get(i, j)} returns the value at row \inlinecode{i}, column \inlinecode{j}, or $0$ if the entry is not stored.
  \item \inlinecode{set(i, j, value)} assigns an entry. Assigning zero erases it from both maps.
  \item \inlinecode{add(i, j, delta)} adds \inlinecode{delta} to an entry. If the result becomes zero, the entry is erased.
  \item \inlinecode{row(i)} returns a map from column index to value for the nonzero entries in row \inlinecode{i}.
  \item \inlinecode{col(j)} returns a map from row index to value for the nonzero entries in column \inlinecode{j}.
  \item \inlinecode{swap\_rows(i, j)} swaps two rows while keeping the column maps synchronized.
  \item \inlinecode{transpose()} transposes the matrix in place.
  \item \inlinecode{a * x} returns the matrix-vector product with vector \inlinecode{x}.
  \item \inlinecode{a + b} and \inlinecode{a - b} return the sum and difference of two matrices with identical dimensions.
  \item \inlinecode{a * b} returns the matrix product, requiring \inlinecode{a.num\_cols()} to equal \inlinecode{b.num\_rows()}. Only pairs of stored entries are ever multiplied.
  \item \inlinecode{a * k} returns \inlinecode{a} scaled by the value \inlinecode{k}. Scaling by zero yields an empty matrix.
  \item Operators \inlinecode{+=}, \inlinecode{-=}, and \inlinecode{*=} assign the corresponding result back into the left operand. In all of these operators, entries that cancel to zero are not stored.
\end{compactitem}

\Needspace*{3\baselineskip}
\textbf{Time Complexity}:
\begin{compactitem}
  \item $O(\log d)$ per call to \inlinecode{get()}, \inlinecode{set()}, and \inlinecode{add()}, where $d$ is the number of nonzeros in the touched row or column.
  \item $O(1)$ per call to \inlinecode{transpose()} and $O(z)$ per matrix-vector product, where $z$ is the number of stored nonzero entries.
  \item $O((r_i + r_j) \cdot \log d)$ per call to \inlinecode{swap\_rows(i, j)}, where $r_i$ and $r_j$ are the sizes of the two rows.
  \item $O((z_a + z_b) \cdot \log d)$ per call to \inlinecode{operator+} and \inlinecode{operator-}, where $z_a$ and $z_b$ are the operand nonzero counts.
  \item $O(z \cdot \log d)$ per call to scalar \inlinecode{operator*}, where $z$ is the number of stored nonzero entries.
  \item $O(f \cdot \log d)$ per call to \inlinecode{operator*}, where $f$ is the number of scalar products accumulated, that is, the sum of \inlinecode{b.row(k).size()} over every stored entry \inlinecode{a[i][k]}.
  \item $O(1)$ per call to \inlinecode{num\_rows()}, \inlinecode{num\_cols()}, \inlinecode{nonzeros()}, \inlinecode{row()}, and \inlinecode{col()}.
\end{compactitem}

\Needspace*{3\baselineskip}
\textbf{Space Complexity}:
\begin{compactitem}
  \item $O(z)$ storage.
  \item $O(1)$ auxiliary for \inlinecode{get()}, \inlinecode{set()}, \inlinecode{add()}, \inlinecode{row()}, \inlinecode{col()}, and \inlinecode{transpose()}.
  \item $O(r_i + r_j)$ auxiliary for \inlinecode{swap\_rows(i, j)}.
  \item $O(m)$ for the vector returned by matrix-vector multiplication, where $m$ is the number of rows.
  \item $O(c)$ auxiliary for \inlinecode{operator*}, where $c$ is the largest number of nonzeros in one result row.
\end{compactitem}

\begin{lstlisting}[style=cpp]
#include <cassert>
#include <map>
#include <unordered_set>
#include <utility>
#include <vector>

template<typename T>
class SparseMatrix {
  int rows = 0, cols = 0, stored = 0;
  std::vector<std::map<int, T>> rvals, cvals;

  bool is_zero(const T &value) const { return value == T{}; }

 public:
  SparseMatrix(int rows, int cols) : rows(rows), cols(cols) {
    assert(rows >= 0 && cols >= 0);
    rvals.resize(rows);
    cvals.resize(cols);
  }

  int num_rows() const { return rows; }
  int num_cols() const { return cols; }
  int nonzeros() const { return stored; }

  T get(int i, int j) const {
    auto it = rvals[i].find(j);
    return it == rvals[i].end() ? T{} : it->second;
  }

  void set(int i, int j, const T &value) {
    auto it = rvals[i].find(j);
    if (is_zero(value)) {
      if (it != rvals[i].end()) {
        rvals[i].erase(it);
        cvals[j].erase(i);
        stored--;
      }
      return;
    }
    if (it == rvals[i].end()) {
      stored++;
    }
    rvals[i][j] = value;
    cvals[j][i] = value;
  }

  void add(int i, int j, const T &delta) {
    if (is_zero(delta)) {
      return;
    }
    set(i, j, get(i, j) + delta);
  }

  const std::map<int, T> &row(int i) const { return rvals[i]; }
  const std::map<int, T> &col(int j) const { return cvals[j]; }

  void swap_rows(int i, int j) {
    if (i == j) {
      return;
    }
    std::unordered_set<int> touched;
    for (const auto &[col, value] : rvals[i]) {
      touched.insert(col);
    }
    for (const auto &[col, value] : rvals[j]) {
      touched.insert(col);
    }
    for (int col : touched) {
      T x = get(i, col), y = get(j, col);
      set(i, col, y);
      set(j, col, x);
    }
  }

  void transpose() {
    std::swap(rows, cols);
    std::swap(rvals, cvals);
  }

  std::vector<T> operator*(const std::vector<T> &x) const {
    assert(static_cast<int>(x.size()) == cols);
    std::vector<T> res(rows);
    for (int i = 0; i < rows; i++) {
      for (const auto &[j, value] : rvals[i]) {
        res[i] += value * x[j];
      }
    }
    return res;
  }

  SparseMatrix operator+(const SparseMatrix &b) const {
    assert(rows == b.rows && cols == b.cols);
    SparseMatrix res(rows, cols);
    for (int i = 0; i < rows; i++) {
      for (const auto &[j, value] : rvals[i]) {
        res.set(i, j, value);
      }
      for (const auto &[j, value] : b.rvals[i]) {
        res.add(i, j, value);  // Cancellation to zero erases the entry.
      }
    }
    return res;
  }

  SparseMatrix operator-(const SparseMatrix &b) const {
    assert(rows == b.rows && cols == b.cols);
    SparseMatrix res(rows, cols);
    for (int i = 0; i < rows; i++) {
      for (const auto &[j, value] : rvals[i]) {
        res.set(i, j, value);
      }
      for (const auto &[j, value] : b.rvals[i]) {
        res.add(i, j, T{} - value);
      }
    }
    return res;
  }

  // Scaling by zero cancels every entry, so the result stores nothing at all.
  SparseMatrix operator*(const T &k) const {
    SparseMatrix res(rows, cols);
    if (is_zero(k)) {
      return res;
    }
    for (int i = 0; i < rows; i++) {
      for (const auto &[j, value] : rvals[i]) {
        res.set(i, j, value * k);
      }
    }
    return res;
  }

  // Row-by-row product: each stored a[i][k] scatters row k of b into an accumulator for row i, so
  // only nonzero-by-nonzero pairs are ever multiplied.
  SparseMatrix operator*(const SparseMatrix &b) const {
    assert(cols == b.rows);
    SparseMatrix res(rows, b.cols);
    for (int i = 0; i < rows; i++) {
      std::map<int, T> acc;
      for (const auto &[k, aik] : rvals[i]) {
        for (const auto &[j, bkj] : b.rvals[k]) {
          acc[j] = acc[j] + aik * bkj;
        }
      }
      for (const auto &[j, value] : acc) {
        res.set(i, j, value);
      }
    }
    return res;
  }

  SparseMatrix &operator+=(const SparseMatrix &b) { return *this = *this + b; }
  SparseMatrix &operator-=(const SparseMatrix &b) { return *this = *this - b; }
  SparseMatrix &operator*=(const SparseMatrix &b) { return *this = *this * b; }
  SparseMatrix &operator*=(const T &k) { return *this = *this * k; }
};
\end{lstlisting}
The elimination routines require field-like arithmetic because they divide by pivots, so use types such as \inlinecode{double}, rational numbers, or modular integers. For floating-point matrices, replace \inlinecode{is\_zero()} with an epsilon comparison if small roundoff values should be erased. At each pivot column, sparse elimination chooses the available row with the fewest stored entries to reduce fill-in, although an unfriendly matrix can still become dense.

\begin{compactitem}
  \item \inlinecode{row\_reduce(a, limit)} converts columns $[0, {\inlinecodemath{limit}})$ of \inlinecode{a} to sparse row echelon form, returning the rank found in those columns.
  \item \inlinecode{det(a)} returns the determinant of a square sparse matrix.
  \item \inlinecode{sparse\_rank(a)} returns the rank of a sparse matrix.
  \item \inlinecode{solve\_system(a, b, \&x)} solves the system $ax = b$, returning $0$ for one solution, $-1$ for no solution, or $-2$ for infinitely many solutions. When one solution exists, \inlinecode{x} is populated.
\end{compactitem}

\Needspace*{3\baselineskip}
\textbf{Time Complexity}:
\begin{compactitem}
  \item $O(f \cdot \log d)$ for sparse elimination, where $f$ is the number of entry updates performed after fill-in and $d$ is a touched row or column size. In the worst case this is still cubic.
  \item $O(f \cdot \log d)$ per call to \inlinecode{det()}, \inlinecode{sparse\_rank()}, and \inlinecode{solve\_system()}.
\end{compactitem}

\textbf{Space Complexity}: $O(z + m)$ auxiliary for elimination, in addition to the copied matrix used by \inlinecode{det()} and \inlinecode{sparse\_rank()} or the augmented matrix used by \inlinecode{solve\_system()}.

\begin{lstlisting}[style=cpp]
template<typename T>
int row_reduce(SparseMatrix<T> &a, int limit) {
  assert(0 <= limit && limit <= a.num_cols());
  int r = 0;
  for (int c = 0; c < limit && r < a.num_rows(); c++) {
    int pivot = -1;
    for (const auto &[i, value] : a.col(c)) {
      if (i >= r && (pivot == -1 || a.row(i).size() < a.row(pivot).size())) {
        pivot = i;
      }
    }
    if (pivot == -1) {
      continue;
    }
    a.swap_rows(r, pivot);
    T inv_pivot = T{1} / a.get(r, c);
    std::vector<std::pair<int, T>> pivot_row(a.row(r).begin(), a.row(r).end());
    std::vector<int> touched_rows;
    for (const auto &[i, value] : a.col(c)) {
      if (i > r) {
        touched_rows.push_back(i);
      }
    }
    for (int i : touched_rows) {
      T coeff = -a.get(i, c) * inv_pivot;
      for (const auto &[j, value] : pivot_row) {
        if (j != c) {
          a.add(i, j, coeff * value);
        }
      }
      a.set(i, c, T{});
    }
    r++;
  }
  return r;
}

template<typename T>
T det(SparseMatrix<T> a) {
  assert(a.num_rows() == a.num_cols());
  int swaps = 0, r = 0, n = a.num_rows();
  T det = 1;
  for (int c = 0; c < n && r < n; c++) {
    int pivot = -1;
    for (const auto &[i, value] : a.col(c)) {
      if (i >= r && (pivot == -1 || a.row(i).size() < a.row(pivot).size())) {
        pivot = i;
      }
    }
    if (pivot == -1) {
      return T{};
    }
    if (pivot != r) {
      a.swap_rows(r, pivot);
      swaps++;
    }
    T pivot_value = a.get(r, c);
    det *= pivot_value;
    T inv_pivot = T{1} / pivot_value;
    std::vector<std::pair<int, T>> pivot_row(a.row(r).begin(), a.row(r).end());
    std::vector<int> touched_rows;
    for (const auto &[i, value] : a.col(c)) {
      if (i > r) {
        touched_rows.push_back(i);
      }
    }
    for (int i : touched_rows) {
      T coeff = -a.get(i, c) * inv_pivot;
      for (const auto &[j, value] : pivot_row) {
        if (j != c) {
          a.add(i, j, coeff * value);
        }
      }
      a.set(i, c, T{});
    }
    r++;
  }
  return (r == n ? (swaps % 2 == 0 ? det : -det) : T{});
}

template<typename T>
int sparse_rank(SparseMatrix<T> a) {
  return row_reduce(a, a.num_cols());
}

template<typename T>
int solve_system(const SparseMatrix<T> &a, const std::vector<T> &b, std::vector<T> *x) {
  if (x == nullptr || a.num_rows() != static_cast<int>(b.size())) {
    return -1;
  }
  int rows = a.num_rows(), cols = a.num_cols();
  SparseMatrix<T> aug(rows, cols + 1);
  for (int i = 0; i < rows; i++) {
    for (const auto &[j, value] : a.row(i)) {
      aug.set(i, j, value);
    }
    aug.set(i, cols, b[i]);
  }
  row_reduce(aug, cols);
  std::vector<int> pivot_col;
  for (int i = 0; i < rows; i++) {
    int lead = -1;
    for (const auto &[j, value] : aug.row(i)) {
      if (j < cols) {
        lead = j;
        break;
      }
    }
    if (lead == -1) {
      if (aug.get(i, cols) != T{}) {
        return -1;
      }
    } else {
      pivot_col.push_back(lead);
    }
  }
  if (static_cast<int>(pivot_col.size()) < cols) {
    return -2;
  }
  x->assign(cols, T{});
  for (int row = static_cast<int>(pivot_col.size()) - 1; row >= 0; row--) {
    int c = pivot_col[row];
    T rhs = aug.get(row, cols);
    for (const auto &[j, value] : aug.row(row)) {
      if (j > c && j < cols) {
        rhs -= value * (*x)[j];
      }
    }
    (*x)[c] = rhs / aug.get(row, c);
  }
  return 0;
}
\end{lstlisting}
\Needspace*{4\baselineskip}
\textbf{Example Usage}\par
\begin{lstlisting}[style=cpp]
#include <cassert>
#include <cstdint>
using namespace std;

int main() {
  SparseMatrix<int64_t> a(3, 4);
  a.set(0, 1, 5);
  a.set(2, 3, 7);
  a.add(0, 1, -2);
  a.add(1, 2, 4);

  assert(a.num_rows() == 3 && a.num_cols() == 4);
  assert(a.nonzeros() == 3);
  assert(a.get(0, 1) == 3);
  assert(a.get(0, 0) == 0);
  assert(a.row(0).begin()->first == 1);
  assert(a.col(2).begin()->first == 1);

  vector<int64_t> x{10, 20, 30, 40};
  assert((a * x == vector<int64_t>{60, 120, 280}));

  a.add(0, 1, -3);  // Entries that become zero are removed from both row and column maps.
  assert(a.get(0, 1) == 0);
  assert(a.nonzeros() == 2);
  assert(a.row(0).empty());
  assert(a.col(1).empty());

  a.transpose();
  assert(a.num_rows() == 4 && a.num_cols() == 3);
  assert(a.get(3, 2) == 7);
  assert(a.get(2, 1) == 4);

  //  p = [1 2]   q = [0 3]   p + q = [1 5]   p * q = [4 3]
  //      [0 1]       [2 0]           [2 1]           [2 0]
  SparseMatrix<int64_t> p(2, 2), q(2, 2);
  p.set(0, 0, 1);
  p.set(0, 1, 2);
  p.set(1, 1, 1);
  q.set(0, 1, 3);
  q.set(1, 0, 2);
  SparseMatrix<int64_t> sum = p + q;
  assert(sum.get(0, 0) == 1 && sum.get(0, 1) == 5);
  assert(sum.get(1, 0) == 2 && sum.get(1, 1) == 1);
  SparseMatrix<int64_t> product = p * q;
  assert(product.get(0, 0) == 4 && product.get(0, 1) == 3);
  assert(product.get(1, 0) == 2 && product.get(1, 1) == 0);
  assert(product.nonzeros() == 3);  // The zero entry is never stored.

  // Cancelling entries are dropped rather than stored as explicit zeros.
  SparseMatrix<int64_t> negated(2, 2);
  negated.set(0, 0, -1);
  negated.set(0, 1, -2);
  negated.set(1, 1, -1);
  assert((p + negated).nonzeros() == 0);
  assert((p - p).nonzeros() == 0);

  SparseMatrix<int64_t> difference = p - q;
  assert(difference.get(0, 0) == 1 && difference.get(0, 1) == -1);
  assert(difference.get(1, 0) == -2 && difference.get(1, 1) == 1);

  SparseMatrix<int64_t> scaled = p * 3;
  assert(scaled.get(0, 0) == 3 && scaled.get(0, 1) == 6 && scaled.get(1, 1) == 3);
  assert((p * 0LL).nonzeros() == 0);

  SparseMatrix<int64_t> acc = p;
  acc += q;
  assert(acc.get(0, 1) == 5 && acc.get(1, 0) == 2);
  acc -= q;
  acc *= 2;
  assert(acc.get(0, 0) == 2 && acc.get(0, 1) == 4 && acc.get(1, 1) == 2);

  SparseMatrix<double> d(3, 3);
  d.set(0, 0, 2);
  d.set(0, 1, 1);
  d.set(1, 1, 3);
  d.set(2, 0, 1);
  d.set(2, 2, 4);
  double determinant = det(d);
  assert(determinant > 24 - 1e-9 && determinant < 24 + 1e-9);
  assert(sparse_rank(d) == 3);

  SparseMatrix<double> sys(3, 3);
  sys.set(0, 0, -1);
  sys.set(0, 1, 2);
  sys.set(0, 2, 5);
  sys.set(1, 0, 1);
  sys.set(1, 2, -6);
  sys.set(2, 0, -4);
  sys.set(2, 1, 2);
  sys.set(2, 2, 2);
  vector<double> rhs{3, 1, -2}, sol;
  assert(solve_system(sys, rhs, &sol) == 0);
  for (int i = 0; i < sys.num_rows(); i++) {
    double sum = 0;
    for (const auto &[j, value] : sys.row(i)) {
      sum += value * sol[j];
    }
    assert(sum > rhs[i] - 1e-9 && sum < rhs[i] + 1e-9);
  }
  return 0;
}
\end{lstlisting}
\subsection{Matrix-Tree Theorem}
\label{sec:6.5.9}
Kirchhoff's matrix-tree theorem counts the spanning trees of an undirected graph. Form the Laplacian matrix $L = D - A$, where $D$ is the diagonal matrix of degrees and $A$ is the adjacency matrix (with multiplicities for parallel edges). Every cofactor of $L$ is equal, and that common value is the number of spanning trees. Equivalently, delete any one row and the matching column from $L$ and take the determinant of what remains.

The determinant is evaluated modulo a prime by Gaussian elimination, so the count is returned modulo \inlinecode{MOD}. Spanning-tree counts grow astronomically, so a modular count is the usual contest form; run the routine under a few different primes and combine with the Chinese remainder theorem to recover a large exact value. The same construction handles weighted graphs (sum edge weights into the Laplacian to get the weighted tree count) and, with the directed Laplacian and a fixed deleted row, counts rooted spanning arborescences.

\begin{compactitem}
  \item \inlinecode{count\_spanning\_trees(n, edges)} returns the number of spanning trees of the undirected graph on nodes $[0, {\inlinecodemath{n}})$ with the given edge list, modulo \inlinecode{MOD}. Parallel edges are honored; self-loops do not affect the count. The result is $0$ when the graph is disconnected, and $1$ when \inlinecode{n} is $1$.
\end{compactitem}

\textbf{Time Complexity}: $O(n^{3} + m)$ per call, where $n$ is the number of nodes and $m$ is the number of edges.

\textbf{Space Complexity}: $O(n^{2})$ auxiliary.

\begin{lstlisting}[style=cpp]
#include <algorithm>
#include <cassert>
#include <cstdint>
#include <utility>
#include <vector>

const int64_t MOD = 998244353;

int64_t powmod(int64_t x, int64_t n, int64_t m) {
  int64_t res = 1 % m;
  for (x %= m; n > 0; n >>= 1) {
    if (n & 1) {
      res = res * x % m;
    }
    x = x * x % m;
  }
  return res;
}

int64_t count_spanning_trees(int n, const std::vector<std::pair<int, int>> &edges) {
  assert(n >= 0);
  assert(std::all_of(edges.begin(), edges.end(), [n](auto edge) {
    auto [u, v] = edge;
    return 0 <= u && u < n && 0 <= v && v < n;
  }));
  if (n <= 1) {
    return n == 1 ? 1 : 0;
  }
  // Build the Laplacian, then drop the last row and column to form an (n-1) by (n-1) cofactor.
  int m = n - 1;
  std::vector<std::vector<int64_t>> lap(m, std::vector<int64_t>(m));
  for (const auto &e : edges) {
    int u = e.first, v = e.second;
    if (u == v) {
      continue;  // Self-loops contribute nothing to the Laplacian.
    }
    if (u < m) {
      lap[u][u] = (lap[u][u] + 1) % MOD;
    }
    if (v < m) {
      lap[v][v] = (lap[v][v] + 1) % MOD;
    }
    if (u < m && v < m) {
      lap[u][v] = (lap[u][v] - 1 + MOD) % MOD;
      lap[v][u] = (lap[v][u] - 1 + MOD) % MOD;
    }
  }
  // Determinant modulo MOD by Gaussian elimination.
  int64_t det = 1;
  for (int col = 0; col < m; col++) {
    int pivot = -1;
    for (int r = col; r < m; r++) {
      if (lap[r][col] != 0) {
        pivot = r;
        break;
      }
    }
    if (pivot == -1) {
      return 0;  // Singular cofactor: the graph is disconnected.
    }
    if (pivot != col) {
      std::swap(lap[pivot], lap[col]);
      det = (MOD - det) % MOD;
    }
    det = det * lap[col][col] % MOD;
    int64_t pinv = powmod(lap[col][col], MOD - 2, MOD);
    for (int r = col + 1; r < m; r++) {
      int64_t factor = lap[r][col] * pinv % MOD;
      if (factor == 0) {
        continue;
      }
      for (int k = col; k < m; k++) {
        lap[r][k] = (lap[r][k] - factor * lap[col][k]) % MOD;
        if (lap[r][k] < 0) {
          lap[r][k] += MOD;
        }
      }
    }
  }
  return det;
}
\end{lstlisting}
\Needspace*{4\baselineskip}
\textbf{Example Usage}\par
\begin{lstlisting}[style=cpp]
#include <cassert>
using namespace std;

int main() {
  // A path 0-1-2 and a single triangle each have a known number of spanning trees.
  assert(count_spanning_trees(3, {{0, 1}, {1, 2}}) == 1);          // A tree has exactly one.
  assert(count_spanning_trees(3, {{0, 1}, {1, 2}, {2, 0}}) == 3);  // Triangle: drop any one edge.

  // Complete graphs: Cayley's formula gives n^(n-2) spanning trees.
  auto complete = [](int n) {
    vector<pair<int, int>> e;
    for (int i = 0; i < n; i++) {
      for (int j = i + 1; j < n; j++) {
        e.emplace_back(i, j);
      }
    }
    return count_spanning_trees(n, e);
  };
  assert(complete(4) == 16);   // 4^2.
  assert(complete(5) == 125);  // 5^3.

  // Disconnected graphs have no spanning tree.
  assert(count_spanning_trees(4, {{0, 1}, {2, 3}}) == 0);
  return 0;
}
\end{lstlisting}

\section{\texorpdfstring{Polynomials and Linear Recurrences}{6.6 Polynomials and Linear Recurrences}}
\setcounter{section}{6}
\setcounter{subsection}{0}
\subsection{Number Theoretic Transform}
\label{sec:6.6.1}
The number theoretic transform (NTT) is the discrete Fourier transform carried out in modular arithmetic instead of over the complex numbers. It multiplies two polynomials (equivalently, computes the convolution of two sequences) modulo a prime in $O(n \log n)$ time, with no floating point error. The modulus must be of the form $c \cdot 2^k + 1$ so that the ring of integers modulo it contains a $2^k$-th root of unity, which limits transform sizes to powers of two up to $2^k$. The code below uses the common prime $998244353 = 119 \cdot 2^{23} + 1$ with primitive root $3$, supporting transforms of any power-of-two length up to $2^{23}$.

The transform replaces a coefficient vector with its evaluations at the powers of a root of unity; pointwise multiplying two such evaluation vectors and transforming back gives their convolution. This is the same divide-and-conquer butterfly as the fast Fourier transform, with the root of unity and all arithmetic taken modulo the prime.

\begin{compactitem}
  \item \inlinecode{ntt(a, invert = false)} transforms the vector \inlinecode{a} in place, whose length must be a power of two. The forward transform uses \inlinecode{invert = false}; the inverse uses \inlinecode{invert = true}.
  \item \inlinecode{convolve(a, b)} returns the convolution of \inlinecode{a} and \inlinecode{b} modulo the prime, that is, the coefficients of the product of the two polynomials whose coefficients are the inputs. The result has length ${\inlinecodemath{a.size()}} + {\inlinecodemath{b.size()}} - 1$, or is empty if either input is empty. Input values are assumed to lie in $[0, {\inlinecodemath{MOD}})$. Small inputs use direct multiplication to avoid NTT overhead.
\end{compactitem}

\Needspace*{3\baselineskip}
\textbf{Time Complexity}:
\begin{compactitem}
  \item $O(n \log n)$ per call to \inlinecode{ntt()}, where $n$ is the length of the vector.
  \item $O(|a||b|)$ per call to \inlinecode{convolve()} on small inputs and $O(n \log n)$ otherwise, where $n$ is the padded transform length.
\end{compactitem}

\Needspace*{3\baselineskip}
\textbf{Space Complexity}:
\begin{compactitem}
  \item $O(1)$ auxiliary for \inlinecode{ntt()}, transforming in place.
  \item $O(n)$ auxiliary for \inlinecode{convolve()}.
\end{compactitem}

\begin{lstlisting}[style=cpp]
#include <algorithm>
#include <cassert>
#include <cstdint>
#include <utility>
#include <vector>

const int64_t MOD = 998244353;
const int64_t ROOT = 3;
const int MAX_POWER_OF_TWO = 23;
const int NAIVE_CUTOFF = 150;

int64_t powmod(int64_t x, int64_t n, int64_t m) {
  int64_t res = 1 % m;
  for (x %= m; n > 0; n >>= 1) {
    if (n & 1) {
      res = res * x % m;
    }
    x = x * x % m;
  }
  return res;
}

void ntt(std::vector<int64_t> &a, bool invert = false) {
  int n = static_cast<int>(a.size());
  assert(n > 0 && (n & (n - 1)) == 0 && n <= (1 << MAX_POWER_OF_TWO));
  for (int i = 1, j = 0; i < n; i++) {
    int bit = n >> 1;
    for (; j & bit; bit >>= 1) {
      j ^= bit;
    }
    j ^= bit;
    if (i < j) {
      std::swap(a[i], a[j]);
    }
  }
  for (int len = 2; len <= n; len <<= 1) {
    int64_t root = powmod(ROOT, (MOD - 1) / len, MOD);
    if (invert) {
      root = powmod(root, MOD - 2, MOD);
    }
    for (int i = 0; i < n; i += len) {
      int64_t w = 1;
      for (int k = 0; k < len / 2; k++) {
        int64_t u = a[i + k], v = a[i + k + len / 2] * w % MOD;
        a[i + k] = (u + v) % MOD;
        a[i + k + len / 2] = (u - v + MOD) % MOD;
        w = w * root % MOD;
      }
    }
  }
  if (invert) {
    int64_t n_inv = powmod(n, MOD - 2, MOD);
    for (int64_t &x : a) {
      x = x * n_inv % MOD;
    }
  }
}

std::vector<int64_t> convolve(std::vector<int64_t> a, std::vector<int64_t> b) {
  if (a.empty() || b.empty()) {
    return {};
  }
  int result_size = static_cast<int>(a.size() + b.size()) - 1;
  if (std::min(a.size(), b.size()) < NAIVE_CUTOFF) {
    std::vector<int64_t> res(result_size);
    for (int i = 0; i < static_cast<int>(a.size()); i++) {
      for (int j = 0; j < static_cast<int>(b.size()); j++) {
        res[i + j] = (res[i + j] + a[i] * b[j]) % MOD;
      }
    }
    return res;
  }
  int n = 1;
  while (n < result_size) {
    n <<= 1;
  }
  assert(n <= (1 << MAX_POWER_OF_TWO));
  a.resize(n);
  b.resize(n);
  ntt(a);
  ntt(b);
  for (int i = 0; i < n; i++) {
    a[i] = a[i] * b[i] % MOD;
  }
  ntt(a, true);
  a.resize(result_size);
  return a;
}
\end{lstlisting}
\Needspace*{4\baselineskip}
\textbf{Example Usage}\par
\begin{lstlisting}[style=cpp]
#include <cassert>
using namespace std;

int main() {
  // (1 + 2x + 3x^2)(4 + 5x + 6x^2) = 4 + 13x + 28x^2 + 27x^3 + 18x^4.
  vector<int64_t> a{1, 2, 3}, b{4, 5, 6};
  assert((convolve(a, b) == vector<int64_t>{4, 13, 28, 27, 18}));
  assert(convolve({}, b).empty());
  assert((convolve(vector<int64_t>{5}, vector<int64_t>{7}) == vector<int64_t>{35}));
  vector<int64_t> ones(160, 1);
  vector<int64_t> d = convolve(ones, ones);
  assert(d[0] == 1 && d[159] == 160 && d[318] == 1);
  return 0;
}
\end{lstlisting}
\subsection{Fast Fourier Transform}
\label{sec:6.6.2}
The fast Fourier transform (FFT) multiplies two polynomials, or equivalently convolves two sequences, in $O(n \log n)$ time by evaluating them at the complex roots of unity. Like the number theoretic transform of the previous section, it replaces a coefficient vector with its evaluations at the powers of a root of unity, multiplies two such evaluation vectors pointwise, and transforms back. The difference is the arithmetic: the FFT works over the complex numbers with $e^{2\pi i / n}$ as the root of unity, so it convolves arbitrary real coefficients rather than residues modulo a prime.

The trade-off is precision. Complex arithmetic in double precision accumulates rounding error, so the result is recovered by rounding each output to the nearest integer; this is reliable as long as the true coefficients stay well within the roughly 15 significant decimal digits a \inlinecode{double} can hold. When exact results modulo a prime are required, or coefficients are large, prefer the number theoretic transform; use the FFT for real-valued convolution or big-integer multiplication.

\begin{compactitem}
  \item \inlinecode{fft(a, invert = false)} transforms the complex vector \inlinecode{a} in place, whose length must be a power of two. The forward transform uses \inlinecode{invert = false}; the inverse uses \inlinecode{invert = true}.
  \item \inlinecode{convolve(a, b)} returns the convolution of two integer sequences \inlinecode{a} and \inlinecode{b}, that is, the coefficients of the product of the two polynomials whose coefficients are the inputs. The result has length ${\inlinecodemath{a.size()}} + {\inlinecodemath{b.size()}} - 1$, or is empty if either input is empty, with each entry rounded to the nearest integer. Small inputs use direct multiplication to avoid FFT overhead.
\end{compactitem}

Overflow warning: The exact coefficients and the direct branch's intermediate products and sums must fit in \inlinecode{int64\_t}.

\Needspace*{3\baselineskip}
\textbf{Time Complexity}:
\begin{compactitem}
  \item $O(n \log n)$ per call to \inlinecode{fft()}, where $n$ is the length of the vector.
  \item $O(|a||b|)$ per call to \inlinecode{convolve()} on small inputs and $O(n \log n)$ otherwise, where $n$ is the padded transform length.
\end{compactitem}

\Needspace*{3\baselineskip}
\textbf{Space Complexity}:
\begin{compactitem}
  \item $O(1)$ auxiliary for \inlinecode{fft()}, transforming in place.
  \item $O(n)$ auxiliary for \inlinecode{convolve()}.
\end{compactitem}

\begin{lstlisting}[style=cpp]
#include <algorithm>
#include <cassert>
#include <cmath>
#include <complex>
#include <cstdint>
#include <utility>
#include <vector>

using cdbl = std::complex<double>;
const double PI = std::acos(-1.0);
const int NAIVE_CUTOFF = 150;

void fft(std::vector<cdbl> &a, bool invert = false) {
  int n = static_cast<int>(a.size());
  assert(n > 0 && (n & (n - 1)) == 0);
  for (int i = 1, j = 0; i < n; i++) {
    int bit = n >> 1;
    for (; j & bit; bit >>= 1) {
      j ^= bit;
    }
    j ^= bit;
    if (i < j) {
      std::swap(a[i], a[j]);
    }
  }
  for (int len = 2; len <= n; len <<= 1) {
    double angle = 2 * PI / len * (invert ? -1 : 1);
    cdbl root(std::cos(angle), std::sin(angle));
    for (int i = 0; i < n; i += len) {
      cdbl w(1);
      for (int k = 0; k < len / 2; k++) {
        cdbl u = a[i + k], v = a[i + k + len / 2] * w;
        a[i + k] = u + v;
        a[i + k + len / 2] = u - v;
        w *= root;
      }
    }
  }
  if (invert) {
    for (cdbl &x : a) {
      x /= n;
    }
  }
}

std::vector<int64_t> convolve(const std::vector<int64_t> &a, const std::vector<int64_t> &b) {
  if (a.empty() || b.empty()) {
    return {};
  }
  int result_size = static_cast<int>(a.size() + b.size() - 1);
  if (std::min(a.size(), b.size()) < NAIVE_CUTOFF) {
    std::vector<int64_t> res(result_size);
    for (int i = 0; i < static_cast<int>(a.size()); i++) {
      for (int j = 0; j < static_cast<int>(b.size()); j++) {
        res[i + j] += a[i] * b[j];  // Overflow warning.
      }
    }
    return res;
  }
  int n = 1;
  while (n < result_size) {
    n <<= 1;
  }
  std::vector<cdbl> fa(a.begin(), a.end()), fb(b.begin(), b.end());
  fa.resize(n);
  fb.resize(n);
  fft(fa);
  fft(fb);
  for (int i = 0; i < n; i++) {
    fa[i] *= fb[i];
  }
  fft(fa, true);
  std::vector<int64_t> res(result_size);
  for (int i = 0; i < result_size; i++) {
    res[i] = std::llround(fa[i].real());
  }
  return res;
}
\end{lstlisting}
\Needspace*{4\baselineskip}
\textbf{Example Usage}\par
\begin{lstlisting}[style=cpp]
using namespace std;

int main() {
  // (1 + 2x + 3x^2)(4 + 5x + 6x^2) = 4 + 13x + 28x^2 + 27x^3 + 18x^4.
  vector<int64_t> a{1, 2, 3}, b{4, 5, 6};
  assert((convolve(a, b) == vector<int64_t>{4, 13, 28, 27, 18}));
  assert(convolve({}, b).empty());
  assert((convolve(vector<int64_t>{5}, vector<int64_t>{7}) == vector<int64_t>{35}));
  vector<int64_t> ones(160, 1);
  vector<int64_t> d = convolve(ones, ones);
  assert(d[0] == 1 && d[159] == 160 && d[318] == 1);
  return 0;
}
\end{lstlisting}
\subsection{Polynomial Operations}
\label{sec:6.6.3}
Polynomial operations treat a vector \inlinecode{a} as coefficients of $a_0 + a_1 x + \dots + a_{n-1} x^{n-1}$ over a field. For formal power series operations, vectors represent coefficients modulo $x^n$; no notion of analytic convergence is involved. This section implements the core algebra needed for formal power series and polynomial division modulo the prime $998244353$. Multiplication uses the number theoretic transform, so inverse and division are fast enough for large polynomials while still keeping the API close to the underlying coefficient vectors.

The code aliases one modular field element as \inlinecode{Coeff} and one coefficient vector as \inlinecode{Poly}. Input coefficients must lie in $[0, {\inlinecodemath{MOD}})$.

\begin{compactitem}
  \item \inlinecode{eval(a, x)} returns the value of \inlinecode{a} at \inlinecode{x} modulo \inlinecode{MOD}; \inlinecode{x} is reduced modulo \inlinecode{MOD}.
  \item \inlinecode{add(a, b)} returns \inlinecode{a + b}.
  \item \inlinecode{subtract(a, b)} returns \inlinecode{a - b}.
  \item \inlinecode{multiply(a, b)} returns \inlinecode{a * b}.
  \item \inlinecode{derivative(a)} returns the formal derivative of \inlinecode{a}.
  \item \inlinecode{antiderivative(a)} returns the formal antiderivative of \inlinecode{a} with constant term zero.
\end{compactitem}

For formal power series, the inverse modulo $x^n$ is a series $b$ satisfying $a \cdot b \equiv 1 \pmod{x^n}$. Series division uses this inverse, and logarithm follows from $(\log a)' = a'/a$. Exponential and square root use Newton iteration, doubling the number of correct coefficients each round. Power factors $a = x^t c b$, where $b(0) = 1$, and applies $a^k = x^{tk} c^k \exp(k \log b)$. To extend the normalized square-root implementation, factor out the lowest power of $x$, normalize its constant term, then multiply back its modular square root (see \secref{6.3.5}); an odd lowest nonzero power has no square root.

\begin{compactitem}
  \item \inlinecode{inverse(a, n)} returns the first \inlinecode{n} coefficients of \inlinecode{1 / a}, requiring \inlinecode{a[0]} to be nonzero.
  \item \inlinecode{series\_divide(a, b, n)} returns the first \inlinecode{n} coefficients of the formal power series \inlinecode{a / b}, requiring \inlinecode{b[0]} to be nonzero.
  \item \inlinecode{log(a, n)} returns the first \inlinecode{n} coefficients of $\log(a)$, requiring \inlinecode{a[0]} to be $1$.
  \item \inlinecode{exp(a, n)} returns the first \inlinecode{n} coefficients of $\exp(a)$, requiring \inlinecode{a[0]} to be $0$.
  \item \inlinecode{sqrt(a, n)} returns the first \inlinecode{n} coefficients of the square root of \inlinecode{a} whose constant term is $1$, requiring \inlinecode{a[0]} to be $1$.
  \item \inlinecode{power(a, k, n)} returns the first \inlinecode{n} coefficients of $a^k$ for a nonnegative \inlinecode{uint64\_t} exponent \inlinecode{k}; \inlinecode{k = 0} returns the constant series $1$.
\end{compactitem}

Euclidean polynomial division uses the reversal trick: reverse both polynomials, compute a truncated series inverse of the reversed divisor, multiply, truncate, and reverse back.

\begin{compactitem}
  \item \inlinecode{div(a, b)} returns the polynomial quotient of \inlinecode{a / b}, requiring \inlinecode{b} to be nonzero.
  \item \inlinecode{mod(a, b)} returns the polynomial remainder of \inlinecode{a / b}, requiring \inlinecode{b} to be nonzero.
\end{compactitem}

\Needspace*{3\baselineskip}
\textbf{Time Complexity}:
\begin{compactitem}
  \item $O(n)$ per call to \inlinecode{eval()}, \inlinecode{add()}, \inlinecode{subtract()}, \inlinecode{derivative()}, and \inlinecode{antiderivative()}.
  \item $O(|a||b|)$ per call to \inlinecode{multiply()} on small inputs and $O(n \log n)$ otherwise, where $n$ is the padded transform length.
  \item $O(n \log n)$ per call to \inlinecode{inverse()} and \inlinecode{series\_divide()}, where $n$ is the requested length.
  \item $O(n \log n)$ per call to \inlinecode{log()}, \inlinecode{exp()}, and \inlinecode{sqrt()}, where $n$ is the requested length. The exponential carries the largest constant, since every Newton round computes a logarithm.
  \item $O(n \log n + \log k)$ per call to \inlinecode{power()}, where $n$ is the requested length.
  \item $O(n \log n)$ per call to \inlinecode{div()} and \inlinecode{mod()}, where $n$ is the padded multiplication length.
\end{compactitem}

\textbf{Space Complexity}: $O(1)$ auxiliary per call to \inlinecode{eval()} and $O(n)$ auxiliary for all other operations.

\begin{lstlisting}[style=cpp]
#include <algorithm>
#include <cassert>
#include <cstdint>
#include <utility>
#include <vector>

using Coeff = int64_t;
using Poly = std::vector<Coeff>;

const Coeff MOD = 998244353;
const Coeff ROOT = 3;
const int MAX_POWER_OF_TWO = 23;
const int NAIVE_CUTOFF = 150;

Coeff powmod(Coeff x, uint64_t n, Coeff m = MOD) {
  Coeff res = 1 % m;
  for (x %= m; n > 0; n >>= 1) {
    if (n & 1) {
      res = res * x % m;
    }
    x = x * x % m;
  }
  return res;
}

void ntt(Poly &a, bool invert) {
  int n = static_cast<int>(a.size());
  assert(n > 0 && (n & (n - 1)) == 0 && n <= (1 << MAX_POWER_OF_TWO));
  for (int i = 1, j = 0; i < n; i++) {
    int bit = n >> 1;
    for (; j & bit; bit >>= 1) {
      j ^= bit;
    }
    j ^= bit;
    if (i < j) {
      std::swap(a[i], a[j]);
    }
  }
  for (int len = 2; len <= n; len <<= 1) {
    Coeff root = powmod(ROOT, (MOD - 1) / len);
    if (invert) {
      root = powmod(root, MOD - 2);
    }
    for (int i = 0; i < n; i += len) {
      Coeff w = 1;
      for (int k = 0; k < len / 2; k++) {
        Coeff u = a[i + k], v = a[i + k + len / 2] * w % MOD;
        a[i + k] = (u + v) % MOD;
        a[i + k + len / 2] = (u - v + MOD) % MOD;
        w = w * root % MOD;
      }
    }
  }
  if (invert) {
    Coeff n_inv = powmod(n, MOD - 2);
    for (Coeff &x : a) {
      x = x * n_inv % MOD;
    }
  }
}

void trim(Poly &a) {
  while (!a.empty() && a.back() == 0) {
    a.pop_back();
  }
}

Coeff eval(const Poly &a, Coeff x) {
  x = (x % MOD + MOD) % MOD;
  Coeff res = 0;
  for (auto it = a.rbegin(); it != a.rend(); ++it) {
    res = (res * x + *it) % MOD;
  }
  return res;
}

Poly add(const Poly &a, const Poly &b) {
  Poly res(std::max(a.size(), b.size()));
  for (int i = 0; i < static_cast<int>(res.size()); i++) {
    Coeff x = i < static_cast<int>(a.size()) ? a[i] : 0;
    Coeff y = i < static_cast<int>(b.size()) ? b[i] : 0;
    res[i] = (x + y) % MOD;
  }
  trim(res);
  return res;
}

Poly subtract(const Poly &a, const Poly &b) {
  Poly res(std::max(a.size(), b.size()));
  for (int i = 0; i < static_cast<int>(res.size()); i++) {
    Coeff x = i < static_cast<int>(a.size()) ? a[i] : 0;
    Coeff y = i < static_cast<int>(b.size()) ? b[i] : 0;
    res[i] = (x - y + MOD) % MOD;
  }
  trim(res);
  return res;
}

Poly multiply(Poly a, Poly b) {
  if (a.empty() || b.empty()) {
    return {};
  }
  int result_size = static_cast<int>(a.size() + b.size() - 1);
  if (std::min(a.size(), b.size()) < NAIVE_CUTOFF) {
    Poly res(result_size);
    for (int i = 0; i < static_cast<int>(a.size()); i++) {
      for (int j = 0; j < static_cast<int>(b.size()); j++) {
        res[i + j] = (res[i + j] + a[i] * b[j]) % MOD;
      }
    }
    trim(res);
    return res;
  }
  int n = 1;
  while (n < result_size) {
    n <<= 1;
  }
  assert(n <= (1 << MAX_POWER_OF_TWO));
  a.resize(n);
  b.resize(n);
  ntt(a, false);
  ntt(b, false);
  for (int i = 0; i < n; i++) {
    a[i] = a[i] * b[i] % MOD;
  }
  ntt(a, true);
  a.resize(result_size);
  trim(a);
  return a;
}

Poly derivative(const Poly &a) {
  if (a.size() <= 1) {
    return {};
  }
  Poly res(a.size() - 1);
  for (int i = 1; i < static_cast<int>(a.size()); i++) {
    res[i - 1] = a[i] * i % MOD;
  }
  trim(res);
  return res;
}

Poly antiderivative(const Poly &a) {
  Poly res(a.size() + 1);
  for (int i = 0; i < static_cast<int>(a.size()); i++) {
    res[i + 1] = a[i] * powmod(i + 1, MOD - 2) % MOD;
  }
  return res;
}

Poly inverse(const Poly &a, int n) {
  assert(n >= 0 && !a.empty() && a[0] != 0);
  if (n == 0) {
    return {};
  }
  Poly res{powmod(a[0], MOD - 2)};
  while (static_cast<int>(res.size()) < n) {
    int len = static_cast<int>(res.size()) << 1;
    Poly cut(a.begin(), a.begin() + std::min(static_cast<int>(a.size()), len));
    Poly prod = multiply(multiply(res, res), cut);
    res.resize(len);
    for (int i = len / 2; i < std::min(len, static_cast<int>(prod.size())); i++) {
      res[i] = (MOD - prod[i]) % MOD;
    }
  }
  res.resize(n);
  return res;
}

Poly series_divide(const Poly &a, const Poly &b, int n) {
  assert(n >= 0 && !b.empty() && b[0] != 0);
  if (n == 0) {
    return {};
  }
  Poly cut(a.begin(), a.begin() + std::min(static_cast<int>(a.size()), n));
  Poly res = multiply(cut, inverse(b, n));
  res.resize(n);
  return res;
}

Poly log(const Poly &a, int n) {
  assert(n >= 0 && !a.empty() && a[0] == 1);
  if (n == 0) {
    return {};
  }
  // Since (log a)' = a'/a, integrating the truncated quotient recovers log a.
  Poly cut(a.begin(), a.begin() + std::min(static_cast<int>(a.size()), n));
  Poly res = multiply(derivative(cut), inverse(cut, n));
  res.resize(n - 1);
  res = antiderivative(res);
  res.resize(n);
  return res;
}

Poly exp(const Poly &a, int n) {
  assert(n >= 0 && (a.empty() || a[0] == 0));
  if (n == 0) {
    return {};
  }
  // Newton iteration on log(f) - a = 0 gives f <- f*(1 - log(f) + a), doubling the number of
  // correct coefficients each round.
  Poly res{1};
  while (static_cast<int>(res.size()) < n) {
    int len = static_cast<int>(res.size()) << 1;
    Poly cut(a.begin(), a.begin() + std::min(static_cast<int>(a.size()), len));
    Poly correction = subtract(cut, log(res, len));
    correction.resize(std::max(1, static_cast<int>(correction.size())));
    correction[0] = (correction[0] + 1) % MOD;
    res = multiply(res, correction);
    res.resize(len);
  }
  res.resize(n);
  return res;
}

Poly sqrt(const Poly &a, int n) {
  assert(n >= 0 && !a.empty() && a[0] == 1);
  if (n == 0) {
    return {};
  }
  // Newton iteration on f^2 - a = 0 gives f <- (f + a/f) / 2.
  const Coeff INV2 = powmod(2, MOD - 2);
  Poly res{1};
  while (static_cast<int>(res.size()) < n) {
    int len = static_cast<int>(res.size()) << 1;
    Poly cut(a.begin(), a.begin() + std::min(static_cast<int>(a.size()), len));
    res = add(res, multiply(cut, inverse(res, len)));
    res.resize(len);
    for (Coeff &c : res) {
      c = c * INV2 % MOD;
    }
  }
  res.resize(n);
  return res;
}

Poly power(const Poly &a, uint64_t k, int n) {
  assert(n >= 0);
  if (n == 0) {
    return {};
  }
  Poly res(n);
  if (k == 0) {
    res[0] = 1;
    return res;
  }
  int limit = std::min(static_cast<int>(a.size()), n);
  int first = 0;
  while (first < limit && a[first] == 0) {
    first++;
  }
  if (first == limit || (first > 0 && k > static_cast<uint64_t>((n - 1) / first))) {
    return res;
  }
  int shift = static_cast<int>(static_cast<uint64_t>(first) * k);
  int len = n - shift;
  int terms = std::min(static_cast<int>(a.size()) - first, len);
  Poly normalized(a.begin() + first, a.begin() + first + terms);
  Coeff leading = normalized[0];
  Coeff leading_inv = powmod(leading, MOD - 2);
  for (Coeff &c : normalized) {
    c = c * leading_inv % MOD;
  }
  Poly exponent = log(normalized, len);
  Coeff k_mod = static_cast<Coeff>(k % static_cast<uint64_t>(MOD));
  for (Coeff &c : exponent) {
    c = c * k_mod % MOD;
  }
  Poly body = exp(exponent, len);
  Coeff leading_power = powmod(leading, k);
  for (int i = 0; i < len; i++) {
    res[i + shift] = body[i] * leading_power % MOD;
  }
  return res;
}

Poly div(Poly a, Poly b) {
  trim(a);
  trim(b);
  assert(!b.empty());
  if (a.size() < b.size()) {
    return {};
  }
  int quotient_size = static_cast<int>(a.size() - b.size() + 1);
  std::reverse(a.begin(), a.end());
  std::reverse(b.begin(), b.end());
  b.resize(quotient_size);
  Poly q = multiply(a, inverse(b, quotient_size));
  q.resize(quotient_size);
  std::reverse(q.begin(), q.end());
  trim(q);
  return q;
}

Poly mod(const Poly &a, const Poly &b) {
  Poly q = div(a, b);
  Poly res = subtract(a, multiply(q, b));
  if (res.size() >= b.size()) {
    res.resize(b.size() - 1);
  }
  trim(res);
  return res;
}
\end{lstlisting}
\Needspace*{4\baselineskip}
\textbf{Example Usage}\par
\begin{lstlisting}[style=cpp]
#include <cassert>

int main() {
  // 5 + 7x + 11x^2 evaluates to 63 at x = 2 and 9 at x = -1.
  Poly p{5, 7, 11};
  assert(eval(p, 2) == 63 && eval(p, -1) == 9);

  // (1 + 2x + 3x^2)(4 + 5x + 6x^2) = 4 + 13x + 28x^2 + 27x^3 + 18x^4.
  Poly a{1, 2, 3}, b{4, 5, 6};
  assert((add(a, b) == Poly{5, 7, 9}));
  assert((subtract(b, a) == Poly{3, 3, 3}));
  assert((multiply(a, b) == Poly{4, 13, 28, 27, 18}));

  // d/dx (5 + 7x + 11x^2) = 7 + 22x; integrating back chooses constant term 0.
  assert((derivative(p) == Poly{7, 22}));
  assert((antiderivative(derivative(p)) == Poly{0, 7, 11}));

  // (1 - x)^-1 = 1 + x + x^2 + x^3 + ... modulo x^6.
  assert((inverse({1, MOD - 1}, 6) == Poly{1, 1, 1, 1, 1, 1}));

  // (1 + x) / (1 - x) = 1 + 2x + 2x^2 + ... modulo x^5.
  assert((series_divide({1, 1}, {1, MOD - 1}, 5) == Poly{1, 2, 2, 2, 2}));

  // (x^3 - 1) / (x - 1) = x^2 + x + 1, remainder 0.
  Poly dividend{MOD - 1, 0, 0, 1};
  Poly divisor{MOD - 1, 1};
  assert((div(dividend, divisor) == Poly{1, 1, 1}));
  assert(mod(dividend, divisor).empty());

  // exp(x) = 1 + x + x^2/2 + x^3/6 + ..., checked by clearing each denominator.
  Poly exp_x = exp({0, 1}, 4);
  assert(exp_x[0] == 1 && exp_x[1] == 1);
  assert(exp_x[2] * 2 % MOD == 1 && exp_x[3] * 6 % MOD == 1);

  // log(1 + x) = x - x^2/2 + x^3/3 - ...
  Poly log_1x = log({1, 1}, 4);
  assert(log_1x[0] == 0 && log_1x[1] == 1);
  assert((log_1x[2] * 2 + 1) % MOD == 0 && log_1x[3] * 3 % MOD == 1);

  // Truncated to the same length, exp and log invert each other.
  Poly series{1, 5, 9, 2, 7};
  Poly logged = log(series, 5);
  assert(logged[0] == 0);
  assert(exp(logged, 5) == series);

  // A series square root squares back to the original, modulo x^5.
  Poly root = sqrt(series, 5);
  Poly root_squared = multiply(root, root);
  root_squared.resize(5);
  assert(root_squared == series);

  // (2x + x^2)^3 = 8x^3 + 12x^4 + 6x^5 + x^6.
  assert((power({0, 2, 1}, 3, 7) == Poly{0, 0, 0, 8, 12, 6, 1}));
  assert((power({}, 0, 4) == Poly{1, 0, 0, 0}));

  Poly ones(160, 1);
  Poly square = multiply(ones, ones);
  assert(square[0] == 1 && square[159] == 160 && square[318] == 1);
  return 0;
}
\end{lstlisting}
\subsection{Polynomial Interpolation}
\label{sec:6.6.4}
Recovers the unique polynomial of degree less than $n$ that passes through $n$ given points, working modulo a prime. The $x$ coordinates must be distinct; the $y$ coordinates may be anything. Two operations are provided: reconstructing the full coefficient vector, and evaluating the interpolating polynomial at a single point without ever forming the coefficients.

Coefficient recovery uses Newton's divided differences. It first replaces the sample values by their successive divided differences in place, then expands the Newton basis $1$, $(x - x_0)$, $(x - x_0)(x - x_1)$, $\ldots$ into ordinary coefficients by multiplying through one factor at a time. Point evaluation instead applies the Lagrange formula directly, summing each sample weighted by the product of $(t - x_j) / (x_k - x_j)$ over the other nodes. Coefficient recovery performs $O(n^{2})$ modular inverses, while point evaluation performs $n$.

\begin{compactitem}
  \item \inlinecode{interpolate(x, y)} returns the coefficients $a_0, ..., a_{n-1}$ (constant term first) of the polynomial $P$ with $P({\inlinecodemath{x[i]}}) = {\inlinecodemath{y[i]}}$, modulo \inlinecode{MOD}. The entries of \inlinecode{x} must be distinct modulo \inlinecode{MOD}.
  \item \inlinecode{interpolate\_at(x, y, t)} returns $P({\inlinecodemath{t}})$ modulo \inlinecode{MOD} for that same polynomial, without building its coefficients.
\end{compactitem}

\textbf{Time Complexity}: $O(n^{2} \log p)$ per call to \inlinecode{interpolate()} and $O(n^{2} + n \log p)$ per call to \inlinecode{interpolate\_at()}, where $n$ is the number of points and $p$ is \inlinecode{MOD}.

\textbf{Space Complexity}: $O(n)$ auxiliary.

\begin{lstlisting}[style=cpp]
#include <cassert>
#include <cstdint>
#include <vector>

const int64_t MOD = 998244353;

int64_t powmod(int64_t x, int64_t n, int64_t m) {
  int64_t res = 1 % m;
  for (x %= m; n > 0; n >>= 1) {
    if (n & 1) {
      res = res * x % m;
    }
    x = x * x % m;
  }
  return res;
}

int64_t modnorm(int64_t a) { return (a %= MOD) < 0 ? a + MOD : a; }
int64_t addmod(int64_t a, int64_t b) { return modnorm(modnorm(a) + modnorm(b)); }
int64_t submod(int64_t a, int64_t b) { return modnorm(modnorm(a) - modnorm(b)); }
int64_t inv(int64_t a) { return powmod(modnorm(a), MOD - 2, MOD); }

void validate_points(const std::vector<int64_t> &x) {
  for (int i = 0; i < static_cast<int>(x.size()); i++) {
    for (int j = 0; j < i; j++) {
      assert(modnorm(x[i]) != modnorm(x[j]));
    }
  }
}

std::vector<int64_t> interpolate(std::vector<int64_t> x, std::vector<int64_t> y) {
  assert(x.size() == y.size());
  validate_points(x);
  int n = static_cast<int>(x.size());
  for (int i = 0; i < n; i++) {
    x[i] = modnorm(x[i]);
    y[i] = modnorm(y[i]);
  }
  // Replace y by its divided differences in place.
  for (int k = 0; k < n; k++) {
    for (int i = k + 1; i < n; i++) {
      y[i] = submod(y[i], y[k]) * inv(submod(x[i], x[k])) % MOD;
    }
  }
  // Expand the Newton basis into ordinary coefficients.
  std::vector<int64_t> res(n), temp(n);
  temp[0] = 1;
  for (int k = 0; k < n; k++) {
    int64_t last = 0;
    for (int i = 0; i < n; i++) {
      res[i] = addmod(res[i], y[k] * temp[i] % MOD);
      int64_t old = temp[i];
      temp[i] = last;
      last = old;
      temp[i] = submod(temp[i], last * x[k] % MOD);
    }
  }
  return res;
}

int64_t interpolate_at(const std::vector<int64_t> &x, const std::vector<int64_t> &y, int64_t t) {
  assert(x.size() == y.size());
  validate_points(x);
  int n = static_cast<int>(x.size());
  t = modnorm(t);
  int64_t res = 0;
  for (int k = 0; k < n; k++) {
    int64_t num = 1, den = 1;
    for (int j = 0; j < n; j++) {
      if (j != k) {
        num = num * submod(t, x[j]) % MOD;
        den = den * submod(x[k], x[j]) % MOD;
      }
    }
    res = addmod(res, modnorm(y[k]) * num % MOD * inv(den) % MOD);
  }
  return res;
}
\end{lstlisting}
\Needspace*{4\baselineskip}
\textbf{Example Usage}\par
\begin{lstlisting}[style=cpp]
#include <cassert>
using namespace std;

int main() {
  // P(x) = x^2 + 2x + 3 sampled at x = 0, 1, 2 gives y = 3, 6, 11.
  vector<int64_t> x{0, 1, 2}, y{3, 6, 11};
  assert((interpolate(x, y) == vector<int64_t>{3, 2, 1}));
  assert(interpolate_at(x, y, 5) == 38);  // 25 + 10 + 3.

  // Non-consecutive nodes for P(x) = 2x^3 - x + 4.
  auto poly = [](int64_t t) { return modnorm(2 * t * t * t - t + 4); };
  vector<int64_t> xs{2, 5, 7, 10}, ys;
  for (int64_t v : xs) {
    ys.push_back(poly(v));
  }
  assert((interpolate(xs, ys) == vector<int64_t>{4, MOD - 1, 0, 2}));  // 4 - x + 2x^3.
  assert(interpolate_at(xs, ys, 13) == poly(13));
  return 0;
}
\end{lstlisting}
\subsection{Linear Recurrence (Berlekamp-Massey)}
\label{sec:6.6.5}
The Berlekamp-Massey algorithm finds the shortest linear recurrence that generates a given sequence over a field. Given the first terms of a sequence taken modulo a prime, it returns coefficients $\Lambda_0, \dots, \Lambda_{L-1}$ of the minimal recurrence $s_i = \Lambda_0 \, s_{i-1} + \Lambda_1 \, s_{i-2} + \dots + \Lambda_{L-1} \, s_{i-L} \pmod{p}$, holding for every $i \geq L$. It scans the sequence once, maintaining a current candidate recurrence and correcting it whenever a term fails to match the prediction, using the last position where a correction was needed. The order $L$ never decreases, and $2L$ terms suffice to determine a recurrence of order $L$. This is the standard tool for guessing a linear recurrence from sampled values, after which the sequence can be extended or its $n$-th term found by fast methods.

An order-$L$ recurrence generally needs at least $2L$ terms to be determined, so the result is most trustworthy when \inlinecode{s} is comfortably longer than twice the returned length. If the returned $L$ is close to $n / 2$, the sequence is under-determined and the recurrence may be a spurious overfit.

The code operates modulo the prime $998244353$; change \inlinecode{MOD} to use a different prime field.

\begin{compactitem}
  \item \inlinecode{berlekamp\_massey(s)} returns the coefficients of the shortest recurrence that the sequence \inlinecode{s} satisfies, as described above. The returned length is the order $L$ of the recurrence. An empty vector is returned when \inlinecode{s} is entirely zero (no recurrence is needed). Entries are reduced modulo \inlinecode{MOD} internally.
\end{compactitem}

\textbf{Time Complexity}: $O(n^{2} + L \log p)$ per call, where $n$ is the length of \inlinecode{s}, $L$ is the order of the recurrence found, and $p$ is the modulus.

\textbf{Space Complexity}: $O(n)$ auxiliary.

\begin{lstlisting}[style=cpp]
#include <cstdint>
#include <vector>

const int64_t MOD = 998244353;

int64_t powmod(int64_t x, int64_t n, int64_t m) {
  int64_t res = 1 % m;
  for (x %= m; n > 0; n >>= 1) {
    if (n & 1) {
      res = res * x % m;
    }
    x = x * x % m;
  }
  return res;
}

std::vector<int64_t> berlekamp_massey(const std::vector<int64_t> &s) {
  int n = static_cast<int>(s.size()), len = 0, m = 0;
  if (n == 0) {
    return {};
  }
  std::vector<int64_t> cur(n), last(n), prev;
  cur[0] = last[0] = 1;
  int64_t inv_last_delta = 1;
  for (int i = 0; i < n; i++) {
    m++;
    int64_t delta = (s[i] % MOD + MOD) % MOD;
    for (int j = 1; j <= len; j++) {
      int64_t value = (s[i - j] % MOD + MOD) % MOD;
      delta = (delta + cur[j] * value) % MOD;
    }
    if (delta == 0) {
      continue;
    }
    prev = cur;
    int64_t coeff = delta * inv_last_delta % MOD;
    for (int j = m; j < n; j++) {
      cur[j] = (cur[j] - coeff * last[j - m] % MOD + MOD) % MOD;
    }
    if (2 * len > i) {
      continue;
    }
    len = i + 1 - len;
    last = prev;
    inv_last_delta = powmod(delta, MOD - 2, MOD);  // Refresh only when the length increases.
    m = 0;
  }
  cur.resize(len + 1);
  cur.erase(cur.begin());
  for (int64_t &x : cur) {
    x = (MOD - x) % MOD;
  }
  return cur;
}
\end{lstlisting}
\Needspace*{4\baselineskip}
\textbf{Example Usage}\par
\begin{lstlisting}[style=cpp]
#include <cassert>
using namespace std;

int main() {
  assert(berlekamp_massey({}).empty());
  // Fibonacci: s[i] = s[i-1] + s[i-2].
  vector<int64_t> fib{1, 1, 2, 3, 5, 8, 13, 21};
  assert((berlekamp_massey(fib) == vector<int64_t>{1, 1}));

  // Geometric: s[i] = 2 * s[i-1].
  vector<int64_t> geo{1, 2, 4, 8, 16, 32};
  assert((berlekamp_massey(geo) == vector<int64_t>{2}));
  assert((berlekamp_massey(vector<int64_t>{1, -1, 1, -1, 1, -1}) == vector<int64_t>{MOD - 1}));

  // Tribonacci: s[i] = s[i-1] + s[i-2] + s[i-3].
  vector<int64_t> trib{0, 1, 1, 2, 4, 7, 13, 24, 44};
  assert((berlekamp_massey(trib) == vector<int64_t>{1, 1, 1}));
  return 0;
}
\end{lstlisting}
\subsection{Linear Recurrence (Kitamasa)}
\label{sec:6.6.6}
Computes the $n$-th term of a linear recurrence modulo a prime in logarithmic time. Given a recurrence of order $L$, $s_i = c_0 s_{i-1} + c_1 s_{i-2} + \ldots + c_{L-1} s_{i-L}$, together with the first $L$ terms, the naive approach unrolls $n$ steps. Kitamasa's method instead jumps directly to index $n$ by working with the characteristic polynomial. The coefficient layout matches the output of Berlekamp-Massey, so the two compose directly: guess the recurrence from sampled values, then jump to any index.

The key identity is that the $n$-th term is a fixed linear combination of the first $L$ terms, with weights given by $x^n$ reduced modulo the characteristic polynomial $f(x) = x^L - c_0 x^{L-1} - ... - c_{L-1}$. The reduction $x^L = c_0 x^{L-1} + ... + c_{L-1}$ lets any product of two degree-$<L$ polynomials be folded back to degree $<L$, so $x^n \mod f$ is obtained by binary exponentiation, and the $n$-th term is the dot product of its coefficients with the initial terms.

Each doubling step is one polynomial multiplication followed by one remainder, both taken from section \secref{6.6.3}, where the multiplication is a number theoretic transform and the remainder is a Newton-iterated division. Squaring and reducing by hand instead costs $O(L^{2})$ per step, which is simpler to write but asymptotically worse; \inlinecode{MOD} is inherited from that section and must stay transform-friendly.

\begin{compactitem}
  \item \inlinecode{kth\_term(rec, init, n)} returns $s_n$ modulo \inlinecode{MOD}, where \inlinecode{rec} holds the coefficients $c_0, ..., c_{L-1}$ and \inlinecode{init} holds at least $L$ initial terms $s_0, ..., s_{L-1}$. Indexing is 0-based and \inlinecode{n} $\geq 0$. Values are reduced modulo \inlinecode{MOD} internally.
\end{compactitem}

\textbf{Time Complexity}: $O(L \log L \log n)$ per call.

\textbf{Space Complexity}: $O(L)$ auxiliary.

\begin{lstlisting}[style=cpp]
#include <cassert>
#include <cstdint>
#include <vector>

#define main polynomial_example
#include "6.6.3_Polynomial_Operations.cpp"
#undef main

int64_t kth_term(const std::vector<int64_t> &rec, const std::vector<int64_t> &init, int64_t n) {
  assert(n >= 0 && init.size() >= rec.size());
  int L = static_cast<int>(rec.size());
  if (n < static_cast<int64_t>(init.size())) {
    return ((init[n] % MOD) + MOD) % MOD;
  }
  if (L == 0) {
    return 0;
  }
  // The characteristic polynomial f(x) = x^L - c_0 x^(L-1) - ... - c_(L-1), indexed by degree.
  Poly f(L + 1);
  f[L] = 1;
  for (int j = 0; j < L; j++) {
    f[L - 1 - j] = ((-rec[j]) % MOD + MOD) % MOD;
  }
  Poly result{1}, base = mod(Poly{0, 1}, f);  // The polynomials x^0 and x, each modulo f.
  for (int64_t e = n; e > 0; e >>= 1) {
    if (e & 1) {
      result = mod(multiply(result, base), f);
    }
    base = mod(multiply(base, base), f);
  }
  int64_t ans = 0;
  for (int j = 0; j < L && j < static_cast<int>(result.size()); j++) {
    ans = (ans + result[j] * (((init[j] % MOD) + MOD) % MOD)) % MOD;
  }
  return ans;
}
\end{lstlisting}
\Needspace*{4\baselineskip}
\textbf{Example Usage}\par
\begin{lstlisting}[style=cpp]
using namespace std;

int main() {
  // Fibonacci s_i = s_{i-1} + s_{i-2} with s_0 = s_1 = 1: 1, 1, 2, 3, 5, 8, ...
  vector<int64_t> fib_rec{1, 1}, fib_init{1, 1};
  assert(kth_term(fib_rec, fib_init, 0) == 1);
  assert(kth_term(fib_rec, fib_init, 9) == 55);
  assert(kth_term(fib_rec, fib_init, 10) == 89);

  // Tribonacci s_i = s_{i-1} + s_{i-2} + s_{i-3} with s_0,s_1,s_2 = 0,1,1.
  vector<int64_t> trib_rec{1, 1, 1}, trib_init{0, 1, 1};
  assert(kth_term(trib_rec, trib_init, 8) == 44);

  // Geometric s_i = 3 s_{i-1}, s_0 = 1: matches 3^n modulo MOD.
  vector<int64_t> geo_rec{3}, geo_init{1};
  int64_t pw = 1;
  for (int i = 0; i < 50; i++) {
    assert(kth_term(geo_rec, geo_init, i) == pw);
    pw = pw * 3 % MOD;
  }
  assert(kth_term({-1}, {1}, 5) == MOD - 1);
  return 0;
}
\end{lstlisting}

\section{\texorpdfstring{Probability and Statistics}{6.7 Probability and Statistics}}
\setcounter{section}{7}
\setcounter{subsection}{0}
\subsection{Probability Distributions}
\label{sec:6.7.1}
The following functions evaluate the standard probability distributions. For a discrete distribution, the probability mass function gives the probability of an exact outcome and the cumulative distribution function gives the probability of that outcome or any smaller one. For a continuous distribution, the density integrates to those probabilities rather than being one, so only the cumulative function is a probability. Counts must be nonnegative, and every probability argument must lie in $[0, 1]$.

Every mass function here is evaluated in logarithms and exponentiated once at the end. Computing $\binom{n}{k} p^k (1 - p)^{n-k}$ directly overflows the binomial coefficient and underflows the powers long before the product itself becomes small, whereas the logarithms stay near the magnitude of the answer. \inlinecode{std::lgamma} supplies $\ln n!$ as $\ln \Gamma(n + 1)$ without forming $n!$, so the factorials never exist, and outcomes are ordered so that no cancellation occurs.

Each cumulative function sums its mass function term by term, which keeps the code short and the error small but costs one term per outcome. For a single tail probability at a very large count the normal approximation with a continuity correction is the usual substitute, while \inlinecode{quantile()} inverts any of these by the bisection of section \secref{5.1.1}.

\begin{compactitem}
  \item \inlinecode{log\_factorial(n)} returns $\ln(n!)$, and \inlinecode{log\_choose(n, k)} returns $\ln \binom{n}{k}$, which is $-\infty$ when \inlinecode{k} lies outside $[0, {\inlinecodemath{n}}]$.
  \item \inlinecode{binomial\_pmf(n, k, p)} and \inlinecode{binomial\_cdf(n, k, p)} evaluate the number of successes in \inlinecode{n} independent trials that each succeed with probability \inlinecode{p}.
  \item \inlinecode{geometric\_pmf(k, p)} and \inlinecode{geometric\_cdf(k, p)} evaluate the number of trials up to and including the first success, so \inlinecode{k} is at least $1$ and \inlinecode{p} must be positive.
  \item \inlinecode{negative\_binomial\_pmf(k, r, p)} and \inlinecode{negative\_binomial\_cdf(k, r, p)} evaluate the number of trials up to and including the \inlinecode{r}-th success, so \inlinecode{k} is at least \inlinecode{r}, which is at least $1$.
  \item \inlinecode{poisson\_pmf(k, lambda)} and \inlinecode{poisson\_cdf(k, lambda)} evaluate the number of events in one unit of time when events arrive independently at nonnegative mean rate \inlinecode{lambda}.
  \item \inlinecode{hypergeometric\_pmf(n, k, draws, hits)} and \inlinecode{hypergeometric\_cdf(n, k, draws, hits)} evaluate the number of marked items among \inlinecode{draws} items taken without replacement from a population of \inlinecode{n} items of which \inlinecode{k} are marked.
  \item \inlinecode{normal\_pdf(x, mu = 0, sigma = 1)} and \inlinecode{normal\_cdf(x, mu = 0, sigma = 1)} evaluate the normal distribution of mean \inlinecode{mu} and positive standard deviation \inlinecode{sigma}, defaulting to the standard normal distribution.
  \item \inlinecode{exponential\_pdf(x, lambda)} and \inlinecode{exponential\_cdf(x, lambda)} evaluate the waiting time until the next event when events arrive at positive mean rate \inlinecode{lambda}.
  \item \inlinecode{quantile(cdf, p, lo, hi, iterations = 100)} returns the smallest $x$ in [\inlinecode{lo}, \inlinecode{hi}] whose cumulative probability \inlinecode{cdf(x)} reaches \inlinecode{p}, for any nondecreasing \inlinecode{cdf}. The bounds must satisfy \inlinecode{lo} $\leq$ \inlinecode{hi}, \inlinecode{cdf(hi)} must reach \inlinecode{p}, and \inlinecode{iterations} must be positive.
\end{compactitem}

Sampling is left to \inlinecode{\textless{}random\textgreater{}}, whose \inlinecode{std::binomial\_distribution}, \inlinecode{std::poisson\_distribution}, \inlinecode{std::normal\_distribution}, and their siblings draw from exactly these families. What the standard library does not provide, and this section supplies, is the ability to evaluate and invert them.

Overflow warning: \inlinecode{log\_choose()} is exact only while \inlinecode{std::lgamma} resolves neighboring integers, which holds well past the range where a \inlinecode{double} probability retains any precision.

\Needspace*{3\baselineskip}
\textbf{Time Complexity}:
\begin{compactitem}
  \item $O(1)$ per call to every mass and density function, and to \inlinecode{log\_factorial()} and \inlinecode{log\_choose()}.
  \item $O(k)$ per call to \inlinecode{binomial\_cdf()}, \inlinecode{poisson\_cdf()}, and \inlinecode{negative\_binomial\_cdf()}, where $k$ is the outcome whose tail is summed, and $O(h)$ per call to \inlinecode{hypergeometric\_cdf()}, where $h$ is its \inlinecode{hits} argument.
  \item $O(1)$ per call to \inlinecode{geometric\_cdf()}, \inlinecode{normal\_cdf()}, and \inlinecode{exponential\_cdf()}, which have closed forms.
  \item $O(n)$ calls to \inlinecode{cdf()} per call to \inlinecode{quantile()}, where $n$ is the number of iterations.
\end{compactitem}

\textbf{Space Complexity}: $O(1)$ auxiliary for all operations.

\begin{lstlisting}[style=cpp]
#include <algorithm>
#include <cassert>
#include <cmath>
#include <limits>

double log_factorial(int n) {
  assert(n >= 0);
  return std::lgamma(n + 1.0);
}

double log_choose(int n, int k) {
  assert(n >= 0);
  if (k < 0 || k > n) {
    return -std::numeric_limits<double>::infinity();
  }
  return log_factorial(n) - log_factorial(k) - log_factorial(n - k);
}

double binomial_pmf(int n, int k, double p) {
  assert(n >= 0 && 0 <= p && p <= 1);
  if (k < 0 || k > n) {
    return 0;
  }
  if (p <= 0) {
    return k == 0 ? 1 : 0;
  }
  if (p >= 1) {
    return k == n ? 1 : 0;
  }
  return std::exp(log_choose(n, k) + k * std::log(p) + (n - k) * std::log1p(-p));
}

double binomial_cdf(int n, int k, double p) {
  assert(n >= 0 && 0 <= p && p <= 1);
  double res = 0;
  int last = std::min(k, n);
  for (int i = 0; i <= last; i++) {
    res += binomial_pmf(n, i, p);
    if (i == last) {
      break;
    }
  }
  return std::min(res, 1.0);
}

double geometric_pmf(int k, double p) {
  assert(0 < p && p <= 1);
  return k < 1 ? 0 : p * std::pow(1 - p, k - 1);
}

double geometric_cdf(int k, double p) {
  assert(0 < p && p <= 1);
  return k < 1 ? 0 : -std::expm1(k * std::log1p(-p));
}

double negative_binomial_pmf(int k, int r, double p) {
  assert(r >= 1 && 0 < p && p <= 1);
  if (k < r) {
    return 0;
  }
  if (p >= 1) {
    return k == r ? 1 : 0;
  }
  return std::exp(log_choose(k - 1, r - 1) + r * std::log(p) + (k - r) * std::log1p(-p));
}

double negative_binomial_cdf(int k, int r, double p) {
  assert(r >= 1 && 0 < p && p <= 1);
  if (k < r) {
    return 0;
  }
  double res = 0;
  for (int i = r;; i++) {
    res += negative_binomial_pmf(i, r, p);
    if (i == k) {
      break;
    }
  }
  return std::min(res, 1.0);
}

double poisson_pmf(int k, double lambda) {
  assert(lambda >= 0);
  if (k < 0) {
    return 0;
  }
  if (lambda <= 0) {
    return k == 0 ? 1 : 0;
  }
  return std::exp(-lambda + k * std::log(lambda) - log_factorial(k));
}

double poisson_cdf(int k, double lambda) {
  assert(lambda >= 0);
  if (k < 0) {
    return 0;
  }
  double res = 0;
  for (int i = 0;; i++) {
    res += poisson_pmf(i, lambda);
    if (i == k) {
      break;
    }
  }
  return std::min(res, 1.0);
}

double hypergeometric_pmf(int n, int k, int draws, int hits) {
  assert(n >= 0 && 0 <= k && k <= n && 0 <= draws && draws <= n);
  if (hits < 0 || hits > draws || hits > k || draws - hits > n - k) {
    return 0;
  }
  return std::exp(log_choose(k, hits) + log_choose(n - k, draws - hits) - log_choose(n, draws));
}

double hypergeometric_cdf(int n, int k, int draws, int hits) {
  assert(n >= 0 && 0 <= k && k <= n && 0 <= draws && draws <= n);
  if (hits < 0) {
    return 0;
  }
  double res = 0;
  int last = std::min({hits, k, draws});
  for (int i = 0; i <= last; i++) {
    res += hypergeometric_pmf(n, k, draws, i);
    if (i == last) {
      break;
    }
  }
  return std::min(res, 1.0);
}

double normal_pdf(double x, double mu = 0, double sigma = 1) {
  assert(sigma > 0);
  static const double INV_SQRT_2PI = 0.3989422804014327;
  double z = (x - mu) / sigma;
  return INV_SQRT_2PI / sigma * std::exp(-0.5 * z * z);
}

double normal_cdf(double x, double mu = 0, double sigma = 1) {
  assert(sigma > 0);
  static const double INV_SQRT_2 = 0.7071067811865476;
  return 0.5 * std::erfc(-(x - mu) * INV_SQRT_2 / sigma);
}

double exponential_pdf(double x, double lambda) {
  assert(lambda > 0);
  return x < 0 ? 0 : lambda * std::exp(-lambda * x);
}

double exponential_cdf(double x, double lambda) {
  assert(lambda > 0);
  return x < 0 ? 0 : -std::expm1(-lambda * x);
}

template<typename Cdf>
double quantile(Cdf cdf, double p, double lo, double hi, int iterations = 100) {
  assert(0 <= p && p <= 1 && lo <= hi && iterations > 0);
  if (cdf(lo) >= p) {
    return lo;
  }
  assert(cdf(hi) >= p);
  for (int i = 0; i < iterations; i++) {
    double mid = lo + (hi - lo) / 2;
    if (cdf(mid) >= p) {
      hi = mid;
    } else {
      lo = mid;
    }
  }
  return hi;
}
\end{lstlisting}
\Needspace*{4\baselineskip}
\textbf{Example Usage}\par
\begin{lstlisting}[style=cpp]
#include <cassert>
using namespace std;

bool EQ(double a, double b) {
  return fabs(a - b) < 1e-9;
}

int main() {
  // Three fair coin flips give two heads with probability 3/8.
  assert(EQ(binomial_pmf(3, 2, 0.5), 0.375));
  assert(EQ(binomial_cdf(3, 3, 0.5), 1));
  // Large counts stay finite where the binomial coefficient alone would not.
  assert(EQ(binomial_cdf(1000, 1000, 0.5), 1));
  assert(binomial_pmf(1000, 500, 0.5) > 0.02);

  // A fair die shows its first six on the third roll with probability 25/216.
  assert(EQ(geometric_pmf(3, 1.0 / 6), 25.0 / 216));
  assert(EQ(geometric_cdf(1, 0.25), 0.25));
  assert(EQ(negative_binomial_pmf(3, 2, 0.5), 0.25));

  assert(EQ(poisson_pmf(0, 2), exp(-2.0)));
  assert(EQ(poisson_cdf(1, 1), 2 * exp(-1.0)));

  // Two marked items among five drawn from ten, of which four are marked.
  assert(EQ(hypergeometric_pmf(10, 4, 5, 2), 10.0 / 21));
  assert(EQ(hypergeometric_cdf(10, 4, 5, 4), 1));

  assert(EQ(normal_cdf(0), 0.5));
  assert(fabs(normal_cdf(1.959963984540054) - 0.975) < 1e-9);
  assert(EQ(exponential_cdf(0, 3), 0));

  // Inverting a cumulative function recovers the value that produced it.
  assert(
      fabs(quantile([](double x) { return normal_cdf(x); }, 0.975, -10, 10) - 1.959963984540054) <
      1e-9
  );
  return 0;
}
\end{lstlisting}
\subsection{Expected Value DP}
\label{sec:6.7.2}
A random process that moves between states accumulates an expected cost that satisfies one linear equation per state. If leaving state $u$ costs $c_u$ and sends the process to state $v$ with probability $p_{uv}$, then linearity of expectation gives $E[u] = c_u + \sum_v p_{uv} E[v]$, where the expected cost of a state is the immediate cost plus the average of the expected costs of its successors. Counting steps is the special case $c_u = 1$ for every non-absorbing state, and an absorbing state is one with no outgoing transitions, so that $E = c$ there.

When the transition graph is acyclic, the equations can be evaluated instead of solved: process the states in reverse topological order and every successor on the right-hand side is already known. This is ordinary backward induction, and it is the reason expected-value problems are usually solved from the goal backwards rather than from the start forwards.

A self-loop is the one cycle this method still handles. If the process stays in $u$ with probability $q$, then $E[u] = c_u + q E[u] + \sum_{v \neq u} p_{uv} E[v]$, and isolating $E[u]$ gives $E[u] = (c_u + \sum_{v \neq u} p_{uv} E[v]) / (1 - q)$. Equivalently, the process retries until it leaves, so the expected number of attempts is the geometric mean $1/(1 - q)$. Self-loops appear whenever a move can be rejected and repeated, such as a die roll that would overshoot the goal or a draw that repeats a coupon already held.

\begin{compactitem}
  \item \inlinecode{expected\_cost(trans, cost)} returns a vector \inlinecode{e} such that \inlinecode{e[u]} is the expected total cost accumulated from state \inlinecode{u} until the process reaches a state with no outgoing transitions. States are the indices of \inlinecode{trans}, each entry \inlinecode{trans[u]} lists the outgoing transitions of \inlinecode{u} as pairs of (\inlinecode{next\_state}, \inlinecode{probability}), and \inlinecode{cost[u]} is the cost charged on leaving \inlinecode{u}. Transition destinations must be valid state indices, and their nonnegative probabilities must sum to at most $1$, with any missing probability treated as ending the process.
\end{compactitem}

Aside from self-loops the transition graph must be acyclic, since a genuine cycle leaves the equations mutually dependent and they must then be solved as a linear system; two states that reach each other call for the absorbing Markov chains of section \secref{6.7.3} instead.

\textbf{Time Complexity}: $O(n + m)$ per call, where $n$ is the number of states and $m$ is the total number of transitions.

\Needspace*{3\baselineskip}
\textbf{Space Complexity}:
\begin{compactitem}
  \item $O(n + m)$ auxiliary.
  \item $O(n)$ for the returned vector.
\end{compactitem}

\begin{lstlisting}[style=cpp]
#include <cassert>
#include <utility>
#include <vector>

std::vector<double> expected_cost(
    const std::vector<std::vector<std::pair<int, double>>> &trans, const std::vector<double> &cost
) {
  int n = static_cast<int>(trans.size());
  assert(static_cast<int>(cost.size()) == n);
  std::vector<int> indeg(n);
  for (int u = 0; u < n; u++) {
    double prob = 0;
    for (auto [v, p] : trans[u]) {
      assert(v >= 0 && v < n && p >= 0);
      prob += p;
      if (v != u) {  // Self-loops are resolved algebraically, so they do not order the states.
        indeg[v]++;
      }
    }
    assert(prob <= 1 + 1e-12);
    (void)prob;
  }
  std::vector<int> order, ready;
  order.reserve(n);
  for (int u = 0; u < n; u++) {
    if (indeg[u] == 0) {
      ready.push_back(u);
    }
  }
  while (!ready.empty()) {
    int u = ready.back();
    ready.pop_back();
    order.push_back(u);
    for (auto [v, p] : trans[u]) {
      if (v != u && --indeg[v] == 0) {
        ready.push_back(v);
      }
    }
  }
  assert(static_cast<int>(order.size()) == n);  // Cycles other than self-loops are unsupported.
  std::vector<double> res(n);
  for (int i = n - 1; i >= 0; i--) {
    int u = order[i];
    double stay = 0, total = cost[u];
    for (auto [v, p] : trans[u]) {
      if (v == u) {
        stay += p;
      } else {
        total += p * res[v];
      }
    }
    assert(stay < 1);  // A state that never leaves itself has infinite expected cost.
    res[u] = total / (1 - stay);
  }
  return res;
}
\end{lstlisting}
\Needspace*{4\baselineskip}
\textbf{Example Usage}\par
\begin{lstlisting}[style=cpp]
#include <cmath>
using namespace std;

bool EQ(double a, double b) {
  return fabs(a - b) < 1e-9;
}

int main() {
  // Roll a fair die to advance along squares [0, n], where a roll that would pass the last square
  // is rerolled. State n is absorbing, and every roll costs one step.
  const int n = 2;
  vector<vector<pair<int, double>>> board(n + 1);
  for (int i = 0; i < n; i++) {
    for (int face = 1; face <= 6; face++) {
      board[i].emplace_back(i + face <= n ? i + face : i, 1.0 / 6);
    }
  }
  vector<double> steps(n + 1, 1);
  steps[n] = 0;
  auto rolls = expected_cost(board, steps);
  assert(EQ(rolls[n], 0));
  assert(EQ(rolls[1], 6.0));  // Only one of six faces advances from square 1.
  assert(EQ(rolls[0], 6.0));
  // The same equation before the self-loop is eliminated: four of six faces reroll square 0.
  assert(EQ(rolls[0], 1 + rolls[1] / 6 + rolls[2] / 6 + 4.0 / 6 * rolls[0]));

  // The coupon collector, where state k is the number of distinct coupons already held. A draw
  // repeats a held coupon with probability k/m, which is a self-loop.
  const int m = 8;
  vector<vector<pair<int, double>>> coupons(m + 1);
  for (int k = 0; k < m; k++) {
    coupons[k].emplace_back(k, static_cast<double>(k) / m);
    coupons[k].emplace_back(k + 1, static_cast<double>(m - k) / m);
  }
  vector<double> draws(m + 1, 1);
  draws[m] = 0;
  auto collected = expected_cost(coupons, draws);
  double harmonic = 0;
  for (int i = 1; i <= m; i++) {
    harmonic += 1.0 / i;
  }
  assert(EQ(collected[0], m * harmonic));  // The classic m*H_m coupon collector bound.
  return 0;
}
\end{lstlisting}
\subsection{Absorbing Markov Chains}
\label{sec:6.7.3}
An absorbing Markov chain splits its states into absorbing states, which are never left once entered, and transient states, which are left forever with probability $1$. Ordering the transient states first splits the transition matrix into four blocks: $Q$ holds the probabilities between transient states, $R$ holds the probabilities from transient states into absorbing ones, and the two remaining blocks are a zero matrix and an identity matrix, since an absorbing state only ever moves to itself.

The expected number of visits to transient state $j$ starting from transient state $i$ is the entry $(i, j)$ of $I + Q + Q^2 + \dots$, since the entry $(i, j)$ of $Q^k$ is the probability of being at $j$ after exactly $k$ steps. Because absorption is certain, the powers of $Q$ tend to zero and the series converges to the fundamental matrix $N = (I - Q)^{-1}$. Every other quantity follows from $N$: summing a row counts the expected visits to all transient states, which is the expected number of steps before absorption, and multiplying by $R$ takes one final step into an absorbing state.

\begin{compactitem}
  \item \inlinecode{fundamental\_matrix(q)} returns $N = (I - {\inlinecodemath{q}})^{-1}$ for the transient-to-transient matrix \inlinecode{q}, as a matrix whose entry in row $i$ and column $j$ is the expected number of visits to transient state $j$ when starting at transient state $i$, counting the starting visit itself.
  \item \inlinecode{absorption\_steps(q)} returns a vector $t$ where $t_i$ is the expected number of steps taken from transient state $i$ until absorption.
  \item \inlinecode{absorption\_probabilities(q, r)} returns a matrix $B$ where $B_{i,k}$ is the probability that a chain starting at transient state $i$ is eventually absorbed at absorbing state $k$, given the transient-to-absorbing matrix \inlinecode{r}.
\end{compactitem}

Every transient state must reach an absorbing one, which is exactly what makes $I - Q$ invertible; a state that cannot is not transient but part of a closed group to be modeled as one absorbing state, so the zero pivot is asserted rather than reported. Unlike the backward induction of section \secref{6.7.2}, cycles are no obstacle here, since the system is solved rather than evaluated.

\Needspace*{3\baselineskip}
\textbf{Time Complexity}:
\begin{compactitem}
  \item $O(n^{3})$ per call to \inlinecode{fundamental\_matrix()} and \inlinecode{absorption\_steps()}, where $n$ is the number of transient states.
  \item $O(n^{3} + n^{2} \cdot k)$ per call to \inlinecode{absorption\_probabilities()}, where $k$ is the number of absorbing states.
\end{compactitem}

\Needspace*{3\baselineskip}
\textbf{Space Complexity}:
\begin{compactitem}
  \item $O(n^{2})$ auxiliary for \inlinecode{fundamental\_matrix()} and \inlinecode{absorption\_steps()}, and $O(n^{2} + n \cdot k)$ for \inlinecode{absorption\_probabilities()}.
  \item $O(n^{2})$ for the returned matrix from \inlinecode{fundamental\_matrix()}, $O(n)$ for the returned vector from \inlinecode{absorption\_steps()}, and $O(n \cdot k)$ for the returned matrix from \inlinecode{absorption\_probabilities()}.
\end{compactitem}

\begin{lstlisting}[style=cpp]
#include <cassert>
#include <cmath>
#include <utility>
#include <vector>

const double EPS = 1e-9;

std::vector<std::vector<double>> fundamental_matrix(const std::vector<std::vector<double>> &q) {
  int n = static_cast<int>(q.size());
  assert(n == 0 || static_cast<int>(q[0].size()) == n);
  // Reduce the augmented matrix [I - Q | I] to [I | N] by Gauss-Jordan elimination.
  std::vector<std::vector<double>> a(n, std::vector<double>(2 * n));
  for (int i = 0; i < n; i++) {
    for (int j = 0; j < n; j++) {
      a[i][j] = (i == j ? 1.0 : 0.0) - q[i][j];
    }
    a[i][n + i] = 1;
  }
  for (int col = 0; col < n; col++) {
    int pivot = col;
    for (int i = col + 1; i < n; i++) {
      if (fabs(a[i][col]) > fabs(a[pivot][col])) {
        pivot = i;
      }
    }
    assert(fabs(a[pivot][col]) > EPS);  // Some transient state never reaches an absorbing state.
    std::swap(a[col], a[pivot]);
    double scale = a[col][col];
    for (int j = col; j < 2 * n; j++) {
      a[col][j] /= scale;
    }
    for (int i = 0; i < n; i++) {
      if (i != col && a[i][col] != 0) {
        double factor = a[i][col];
        for (int j = col; j < 2 * n; j++) {
          a[i][j] -= factor * a[col][j];
        }
      }
    }
  }
  std::vector<std::vector<double>> res(n, std::vector<double>(n));
  for (int i = 0; i < n; i++) {
    for (int j = 0; j < n; j++) {
      res[i][j] = a[i][n + j];
    }
  }
  return res;
}

std::vector<double> absorption_steps(const std::vector<std::vector<double>> &q) {
  auto visits = fundamental_matrix(q);
  int n = static_cast<int>(visits.size());
  std::vector<double> res(n);
  for (int i = 0; i < n; i++) {
    for (int j = 0; j < n; j++) {
      res[i] += visits[i][j];
    }
  }
  return res;
}

std::vector<std::vector<double>> absorption_probabilities(
    const std::vector<std::vector<double>> &q, const std::vector<std::vector<double>> &r
) {
  assert(q.size() == r.size());
  auto visits = fundamental_matrix(q);
  int n = static_cast<int>(visits.size()), k = n == 0 ? 0 : static_cast<int>(r[0].size());
  std::vector<std::vector<double>> res(n, std::vector<double>(k));
  for (int i = 0; i < n; i++) {
    for (int j = 0; j < n; j++) {
      for (int c = 0; c < k; c++) {
        res[i][c] += visits[i][j] * r[j][c];
      }
    }
  }
  return res;
}
\end{lstlisting}
\Needspace*{4\baselineskip}
\textbf{Example Usage}\par
\begin{lstlisting}[style=cpp]
#include <algorithm>
#include <cmath>
using namespace std;

bool EQ(double a, double b) {
  return fabs(a - b) < 1e-9;
}

int main() {
  // A gambler with i of 4 dollars bets a dollar on a fair coin until going broke or reaching 4.
  // Transient states 0, 1, and 2 stand for holding 1, 2, and 3 dollars, and the two absorbing
  // states are ruin and victory. States 0 and 1 can each reach the other, so the transition graph
  // is cyclic and backward induction alone does not apply.
  vector<vector<double>> q{{0, 0.5, 0}, {0.5, 0, 0.5}, {0, 0.5, 0}};
  vector<vector<double>> r{{0.5, 0}, {0, 0}, {0, 0.5}};

  auto visits = fundamental_matrix(q);
  assert(EQ(visits[0][0], 1.5) && EQ(visits[0][1], 1.0) && EQ(visits[0][2], 0.5));
  assert(EQ(visits[1][0], 1.0) && EQ(visits[1][1], 2.0) && EQ(visits[1][2], 1.0));
  assert(EQ(visits[2][0], 0.5) && EQ(visits[2][1], 1.0) && EQ(visits[2][2], 1.5));

  // Starting with i dollars, the game lasts i*(4 - i) rounds on average.
  auto steps = absorption_steps(q);
  assert(EQ(steps[0], 3.0) && EQ(steps[1], 4.0) && EQ(steps[2], 3.0));

  // A fair game is won with probability i/4.
  auto odds = absorption_probabilities(q, r);
  assert(EQ(odds[0][1], 0.25) && EQ(odds[1][1], 0.5) && EQ(odds[2][1], 0.75));
  assert(all_of(odds.begin(), odds.end(), [](const auto &row) {
    return EQ(row[0] + row[1], 1.0);  // Absorption is certain.
  }));
  return 0;
}
\end{lstlisting}
\subsection{Stationary Distribution}
\label{sec:6.7.4}
A Markov chain with no absorbing states never settles on a single state, but the fraction of time it spends in each state does settle. That limit is the stationary distribution: the row vector $\pi$ with $\pi P = \pi$ and $\sum_i \pi_i = 1$, which is the one distribution left unchanged by a step of the chain. It exists and is unique whenever the chain is irreducible, meaning every state can reach every other, and it is the limit of the state distribution from any starting point once the chain is also aperiodic.

The defining equations are one short of determining $\pi$. Each column of $\pi P = \pi$ says that the probability flowing into a state equals the probability flowing out, and those $n$ equations sum to the trivial identity $1 = 1$, leaving the system rank $n - 1$ and any scalar multiple of a solution also a solution. Substituting the normalization $\sum_i \pi_i = 1$ for one of them pins down the scale and makes the system uniquely solvable by ordinary elimination.

\begin{compactitem}
  \item \inlinecode{stationary\_distribution(p, eps = 1e-10)} returns the stationary distribution of the $n$ by $n$ row-stochastic transition matrix \inlinecode{p}, where \inlinecode{p[i][j]} is the probability of stepping from state \inlinecode{i} to state \inlinecode{j}. It returns \inlinecode{std::nullopt} when the system is singular to within \inlinecode{eps}, which happens when the chain has more than one closed recurrent class and therefore no unique stationary distribution. A reducible chain may still have a unique distribution when all states eventually enter the same closed class.
\end{compactitem}

Two quantities follow at once: the expected number of steps to return to state $i$ is $1/\pi_i$, and on an undirected graph walked by a uniformly random incident edge, $\pi_v$ is proportional to the degree of $v$, which is why such a walk is uniform in the limit on a regular graph. Where the absorbing chains of section \secref{6.7.3} ask what happens before the process stops, this asks what happens when it never does.

\textbf{Time Complexity}: $O(n^{3})$ per call, where $n$ is the number of states.

\textbf{Space Complexity}: $O(n^{2})$ auxiliary and $O(n)$ for the returned distribution.

\begin{lstlisting}[style=cpp]
#include <cassert>
#include <cmath>
#include <optional>
#include <utility>
#include <vector>

std::optional<std::vector<double>> stationary_distribution(
    const std::vector<std::vector<double>> &p, double eps = 1e-10
) {
  int n = static_cast<int>(p.size());
  assert(n == 0 || static_cast<int>(p[0].size()) == n);
  if (n == 0) {
    return std::vector<double>{};
  }
  // Build the transpose of P - I, then replace its last row with the normalization equation.
  std::vector<std::vector<double>> a(n, std::vector<double>(n + 1));
  for (int i = 0; i < n; i++) {
    for (int j = 0; j < n; j++) {
      a[i][j] = p[j][i] - (i == j ? 1.0 : 0.0);
    }
  }
  for (int j = 0; j < n; j++) {
    a[n - 1][j] = 1;
  }
  a[n - 1][n] = 1;
  for (int col = 0; col < n; col++) {
    int pivot = col;
    for (int i = col + 1; i < n; i++) {
      if (fabs(a[i][col]) > fabs(a[pivot][col])) {
        pivot = i;
      }
    }
    if (fabs(a[pivot][col]) <= eps) {
      return std::nullopt;  // Multiple closed classes make the stationary distribution nonunique.
    }
    std::swap(a[col], a[pivot]);
    double scale = a[col][col];  // Saved, since the loop below overwrites it on its first step.
    for (int j = col; j <= n; j++) {
      a[col][j] /= scale;
    }
    for (int i = 0; i < n; i++) {
      if (i != col && a[i][col] != 0) {
        double factor = a[i][col];
        for (int j = col; j <= n; j++) {
          a[i][j] -= factor * a[col][j];
        }
      }
    }
  }
  std::vector<double> res(n);
  for (int i = 0; i < n; i++) {
    res[i] = a[i][n];
  }
  return res;
}
\end{lstlisting}
\Needspace*{4\baselineskip}
\textbf{Example Usage}\par
\begin{lstlisting}[style=cpp]
#include <algorithm>
#include <cassert>
#include <cmath>
using namespace std;

bool EQ(double a, double b) {
  return fabs(a - b) < 1e-9;
}

int main() {
  // A two-state chain that leaves state 0 with probability 1/4 and state 1 with probability 1/2,
  // so it rests in state 0 twice as often.
  auto two = stationary_distribution({{0.75, 0.25}, {0.5, 0.5}});
  assert(two.has_value());
  assert(EQ((*two)[0], 2.0 / 3) && EQ((*two)[1], 1.0 / 3));

  // A random walk on a path of three vertices, where the middle vertex has twice the degree and so
  // twice the stationary probability. Its expected return time is the reciprocal, 2.
  auto path = stationary_distribution({{0, 1, 0}, {0.5, 0, 0.5}, {0, 1, 0}});
  assert(path.has_value());
  assert(EQ((*path)[0], 0.25) && EQ((*path)[1], 0.5) && EQ((*path)[2], 0.25));
  assert(EQ(1 / (*path)[1], 2.0));

  // A doubly stochastic chain is uniform in the limit no matter the transition probabilities.
  auto cycle = stationary_distribution({{0, 0.3, 0.7}, {0.7, 0, 0.3}, {0.3, 0.7, 0}});
  assert(cycle.has_value());
  assert(all_of(cycle->begin(), cycle->end(), [](double p) { return EQ(p, 1.0 / 3); }));

  // Two states that never reach each other admit no unique distribution.
  assert(!stationary_distribution({{1, 0}, {0, 1}}).has_value());

  // Reducibility alone is not an obstruction: state 0 is transient and state 1 is the sole class.
  auto reducible = stationary_distribution({{0, 1}, {0, 1}});
  assert(reducible.has_value() && EQ((*reducible)[0], 0) && EQ((*reducible)[1], 1));
  return 0;
}
\end{lstlisting}
\subsection{Statistical Learning}
\label{sec:6.7.5}
The following routines fit a model to observed data and use it to describe or predict unobserved data. Each takes its input as a \inlinecode{Points} matrix whose rows are samples and whose columns are features, so a row is one observation of the same measurements taken in the same order. Supervised routines additionally take a \inlinecode{labels} vector, parallel to the rows, holding each sample's class as an integer in $[0, {\inlinecodemath{classes}})$.

Features measured on different scales distort every distance here, so call \inlinecode{standardize()} before \inlinecode{kmeans()}, \inlinecode{knn\_classify()}, or \inlinecode{logistic\_regression()} and apply the returned pair to later queries. Naive Bayes is exempt, since it fits a separate variance per feature. None of these choose \inlinecode{k} or \inlinecode{l2} for you: hold out part of the data and keep the value predicting it best.

\begin{lstlisting}[style=cpp]
#include <algorithm>
#include <cassert>
#include <cmath>
#include <cstddef>
#include <limits>
#include <numeric>
#include <random>
#include <utility>
#include <vector>

using Points = std::vector<std::vector<double>>;

double sqdist(const std::vector<double> &a, const std::vector<double> &b) {
  assert(a.size() == b.size());
  double res = 0;
  for (size_t i = 0; i < a.size(); i++) {
    res += (a[i] - b[i]) * (a[i] - b[i]);
  }
  return res;
}
\end{lstlisting}
Standardizing subtracts each column's mean and divides by its standard deviation, which is what makes one Euclidean distance meaningful across features recorded in different units. The population standard deviation is used, dividing by $n$ rather than $n - 1$, since these columns are the data being fit and not a sample standing in for a larger population. A column that never varies would give a zero divisor, so it reports $1$ instead and standardizes to all zeros, contributing nothing to any distance, which is exactly what a constant feature deserves.

\begin{compactitem}
  \item \inlinecode{column\_mean(points)} and \inlinecode{column\_stddev(points, mean)} return the per-column mean and population standard deviation, the latter reporting $1$ for a column that never varies.
  \item \inlinecode{standardize(points)} rescales every column in place to zero mean and unit variance and returns the (\inlinecode{mean}, \inlinecode{stddev}) it used, while \inlinecode{standardize(x, mean, stddev)} applies a returned pair to one later sample.
\end{compactitem}

\textbf{Time Complexity}: $O(n \cdot d)$ per call over $n$ rows of $d$ columns, and $O(d)$ for one sample.

\textbf{Space Complexity}: $O(d)$ auxiliary.

\begin{lstlisting}[style=cpp]
std::vector<double> column_mean(const Points &points) {
  assert(!points.empty());
  int n = static_cast<int>(points.size()), d = static_cast<int>(points[0].size());
  std::vector<double> res(d);
  for (int i = 0; i < n; i++) {
    for (int j = 0; j < d; j++) {
      res[j] += points[i][j];
    }
  }
  for (int j = 0; j < d; j++) {
    res[j] /= n;
  }
  return res;
}

std::vector<double> column_stddev(const Points &points, const std::vector<double> &mean) {
  assert(!points.empty() && mean.size() == points[0].size());
  int n = static_cast<int>(points.size()), d = static_cast<int>(points[0].size());
  std::vector<double> res(d);
  for (int i = 0; i < n; i++) {
    for (int j = 0; j < d; j++) {
      res[j] += (points[i][j] - mean[j]) * (points[i][j] - mean[j]);
    }
  }
  for (int j = 0; j < d; j++) {
    res[j] = std::sqrt(res[j] / n);
    if (res[j] < 1e-12) {
      res[j] = 1;
    }
  }
  return res;
}

void standardize(
    std::vector<double> &x, const std::vector<double> &mean, const std::vector<double> &stddev
) {
  assert(x.size() == mean.size() && x.size() == stddev.size());
  assert(std::all_of(stddev.begin(), stddev.end(), [](double value) { return value > 0; }));
  for (size_t j = 0; j < x.size(); j++) {
    x[j] = (x[j] - mean[j]) / stddev[j];
  }
}

std::pair<std::vector<double>, std::vector<double>> standardize(Points &points) {
  std::vector<double> mean = column_mean(points), stddev = column_stddev(points, mean);
  for (std::vector<double> &row : points) {
    standardize(row, mean, stddev);
  }
  return std::make_pair(mean, stddev);
}
\end{lstlisting}
$k$-means partitions the rows into $k$ clusters so that each row lies nearer to its own cluster's mean than to any other. Lloyd's algorithm alternates the two halves of that condition: assign every row to the nearest centroid, then move every centroid to the mean of the rows assigned to it. Each half can only lower the total squared distance, so the process converges, but only to a local optimum that the starting centroids decide.

Choosing those starting centroids is called seeding, and $k$-means++ is the standard rule: take the first uniformly at random, then draw each remaining one with probability proportional to its squared distance from the nearest already chosen. That spreads them out instead of letting several land in one dense region, bounding the expected cost at $O(\log k)$ times optimal, which uniform choice cannot promise. The loop stops once an assignment pass changes nothing, and a cluster that loses all its rows keeps its previous centroid.

\begin{compactitem}
  \item \inlinecode{kmeans(points, k, rng, iterations = 100)} returns (\inlinecode{centroids}, \inlinecode{assignment}), partitioning the rows into \inlinecode{k} clusters, where \inlinecode{assignment[i]} is the cluster of row \inlinecode{i} and \inlinecode{centroids[c]} is the mean of cluster \inlinecode{c}. The number of rows must be at least \inlinecode{k}; both \inlinecode{k} and \inlinecode{iterations} must be positive.
\end{compactitem}

\textbf{Time Complexity}: $O(i \cdot n \cdot k \cdot d)$ per call for \inlinecode{iterations} $i$, $n$ rows, and $d$ columns.

\textbf{Space Complexity}: $O(n + k \cdot d)$ auxiliary.

\begin{lstlisting}[style=cpp]
std::pair<Points, std::vector<int>> kmeans(
    const Points &points, int k, std::mt19937 &rng, int iterations = 100
) {
  int n = static_cast<int>(points.size());
  assert(k > 0 && k <= n && iterations > 0);
  Points centroids;
  centroids.push_back(points[std::uniform_int_distribution<int>(0, n - 1)(rng)]);
  std::vector<double> best(n);
  for (int i = 0; i < n; i++) {
    best[i] = sqdist(points[i], centroids[0]);
  }
  while (static_cast<int>(centroids.size()) < k) {
    std::discrete_distribution<int> pick(best.begin(), best.end());
    centroids.push_back(points[pick(rng)]);
    for (int i = 0; i < n; i++) {
      best[i] = std::min(best[i], sqdist(points[i], centroids.back()));
    }
  }
  std::vector<int> assignment(n, -1);
  for (int it = 0; it < iterations; it++) {
    bool changed = false;
    for (int i = 0; i < n; i++) {
      int nearest = 0;
      double nearest_dist = sqdist(points[i], centroids[0]);
      for (int c = 1; c < k; c++) {
        double dist = sqdist(points[i], centroids[c]);
        if (dist < nearest_dist) {
          nearest = c;
          nearest_dist = dist;
        }
      }
      if (assignment[i] != nearest) {
        assignment[i] = nearest;
        changed = true;
      }
    }
    if (!changed) {
      break;
    }
    Points sums(k, std::vector<double>(points[0].size()));
    std::vector<int> counts(k);
    for (int i = 0; i < n; i++) {
      for (size_t j = 0; j < points[i].size(); j++) {
        sums[assignment[i]][j] += points[i][j];
      }
      counts[assignment[i]]++;
    }
    for (int c = 0; c < k; c++) {
      if (counts[c] > 0) {
        for (size_t j = 0; j < sums[c].size(); j++) {
          centroids[c][j] = sums[c][j] / counts[c];
        }
      }
    }
  }
  return std::make_pair(centroids, assignment);
}
\end{lstlisting}
$k$-nearest neighbors classifies a query by letting the $k$ nearest training rows vote, which fits no model at all: the training data is the model, and all the work happens at query time. Small $k$ follows the data closely and is swayed by noise, while large $k$ smooths across genuine boundaries, so $k$ trades variance against bias. Averaging the neighbors' values rather than counting their votes turns the same routine into regression.

Selecting the $k$ nearest with \inlinecode{std::nth\_element} costs linear time rather than the $O(n \log n)$ of a full sort, since the neighbors' relative order is irrelevant to the vote. Scanning every row makes each query linear in the data, which the k-d tree of section \secref{2.6.5} improves to logarithmic on low-dimensional data; above roughly ten features the two converge and the scan below is preferable.

\begin{compactitem}
  \item \inlinecode{knn\_classify(points, labels, query, k)} returns the majority label among the \inlinecode{k} rows nearest to \inlinecode{query}, breaking a tie toward the smaller label. \inlinecode{k} must lie in $[1, n]$ for \inlinecode{n} rows.
\end{compactitem}

\textbf{Time Complexity}: $O(n \cdot d)$ per call, for $n$ rows of $d$ columns.

\textbf{Space Complexity}: $O(n)$ auxiliary.

\begin{lstlisting}[style=cpp]
int knn_classify(
    const Points &points, const std::vector<int> &labels, const std::vector<double> &query, int k
) {
  int n = static_cast<int>(points.size());
  assert(n > 0 && static_cast<int>(labels.size()) == n && query.size() == points[0].size());
  assert(k >= 1 && k <= n);
  assert(std::all_of(labels.begin(), labels.end(), [](int label) { return label >= 0; }));
  std::vector<int> order(n);
  std::iota(order.begin(), order.end(), 0);
  std::vector<double> dist(n);
  for (int i = 0; i < n; i++) {
    dist[i] = sqdist(points[i], query);
  }
  std::nth_element(order.begin(), order.begin() + (k - 1), order.end(), [&](int a, int b) {
    return dist[a] < dist[b];
  });
  std::vector<int> votes(*std::max_element(labels.begin(), labels.end()) + 1);
  for (int i = 0; i < k; i++) {
    votes[labels[order[i]]]++;
  }
  return static_cast<int>(std::max_element(votes.begin(), votes.end()) - votes.begin());
}
\end{lstlisting}
Naive Bayes picks the class maximizing $P(c) \prod_j P(x_j \mid c)$, which is Bayes' rule with the denominator dropped because it does not depend on the class. The independence it assumes among features given the class is almost never true, and yet the classifier is accurate anyway: the products it computes are poorly calibrated as probabilities, but the class that maximizes them is usually still the right one, since the errors tend to move every class the same way.

The Gaussian variant below models each feature within each class by a normal distribution, so fitting is one pass accumulating a mean and variance per (class, feature) pair. Scores accumulate as logarithms, both because a product of hundreds of densities underflows and because the logarithm of a normal density is a plain quadratic. A feature that is constant within a class would give a zero variance and an infinite score, so variances are floored at a small fraction of the overall spread.

\begin{compactitem}
  \item \inlinecode{NaiveBayes(points, labels, classes)} fits a Gaussian naive Bayes classifier, \inlinecode{predict(x)} returns its most probable class for a sample \inlinecode{x}, and \inlinecode{log\_likelihood(x, c)} returns the unnormalized log probability that \inlinecode{x} belongs to class \inlinecode{c}. A class carried by no row is given zero prior, so it scores $-\infty$ and is never predicted.
\end{compactitem}

\Needspace*{3\baselineskip}
\textbf{Time Complexity}:
\begin{compactitem}
  \item $O(n \cdot d)$ per construction.
  \item $O(c \cdot d)$ per call to \inlinecode{predict()} and $O(d)$ per call to \inlinecode{log\_likelihood()}, for $n$ rows of $d$ columns and $c$ classes.
\end{compactitem}

\textbf{Space Complexity}: $O(c \cdot d)$ for storage.

\begin{lstlisting}[style=cpp]
class NaiveBayes {
  std::vector<double> log_prior;
  Points mean, var;

 public:
  NaiveBayes(const Points &points, const std::vector<int> &labels, int classes) {
    assert(!points.empty() && labels.size() == points.size() && classes > 0);
    assert(std::all_of(labels.begin(), labels.end(), [&](int label) {
      return 0 <= label && label < classes;
    }));
    int n = static_cast<int>(points.size()), d = static_cast<int>(points[0].size());
    std::vector<int> counts(classes);
    mean.assign(classes, std::vector<double>(d));
    var.assign(classes, std::vector<double>(d));
    log_prior.assign(classes, 0);
    for (int i = 0; i < n; i++) {
      counts[labels[i]]++;
      for (int j = 0; j < d; j++) {
        mean[labels[i]][j] += points[i][j];
      }
    }
    for (int c = 0; c < classes; c++) {
      if (counts[c] == 0) {  // A class with no rows has zero prior, so it is never predicted.
        log_prior[c] = -std::numeric_limits<double>::infinity();
        continue;
      }
      log_prior[c] = std::log(static_cast<double>(counts[c]) / n);
      for (int j = 0; j < d; j++) {
        mean[c][j] /= counts[c];
      }
    }
    double spread = 0;
    for (int i = 0; i < n; i++) {
      for (int j = 0; j < d; j++) {
        double diff = points[i][j] - mean[labels[i]][j];
        var[labels[i]][j] += diff * diff;
        spread += diff * diff;
      }
    }
    double floor_var = 1e-9 * spread / n + 1e-12;
    for (int c = 0; c < classes; c++) {
      for (int j = 0; j < d; j++) {
        var[c][j] = counts[c] == 0 ? floor_var : std::max(var[c][j] / counts[c], floor_var);
      }
    }
  }

  double log_likelihood(const std::vector<double> &x, int c) const {
    assert(c >= 0 && c < static_cast<int>(log_prior.size()) && x.size() == mean[c].size());
    static const double LOG_2PI = 1.8378770664093453;
    double res = log_prior[c];
    for (size_t j = 0; j < x.size(); j++) {
      double diff = x[j] - mean[c][j];
      res -= 0.5 * (LOG_2PI + std::log(var[c][j]) + diff * diff / var[c][j]);
    }
    return res;
  }

  int predict(const std::vector<double> &x) const {
    int best = 0;
    double best_score = log_likelihood(x, 0);
    for (size_t c = 1; c < log_prior.size(); c++) {
      double score = log_likelihood(x, static_cast<int>(c));
      if (score > best_score) {
        best = static_cast<int>(c);
        best_score = score;
      }
    }
    return best;
  }
};
\end{lstlisting}
Logistic regression models the probability of class $1$ as $\sigma(w \cdot x)$ for the logistic function $\sigma(z) = 1/(1 + e^{-z})$, which maps any real score into $(0, 1)$. Fitting minimizes the average negative log-likelihood, a convex function whose gradient is $X^{\top}(\sigma(Xw) - y)/n$: each row pushes the weights by its own features, scaled by how wrong the current prediction is on it. Convexity makes every local minimum global, so plain gradient descent suffices.

Unlike the least squares of section \secref{6.5.5} there is no closed form, and on perfectly separable data the weights diverge as every predicted probability is driven to $0$ or $1$. The \inlinecode{l2} penalty adds $\lambda\|w\|^2$ to keep them finite, excluding the intercept so that shifting all labels does not distort it. Section \secref{5.2.4} supplies the adaptive steps worth substituting when convergence is slow.

\begin{compactitem}
  \item \inlinecode{logistic\_predict(w, x)} returns the probability that sample \inlinecode{x} belongs to class $1$ under a weight vector \inlinecode{w} whose index $0$ is the intercept.
  \item \inlinecode{logistic\_regression(points, labels, ...)} returns such a weight vector, fit by gradient descent over labels that are $0$ or $1$. Optional parameters default to \inlinecode{rate = 0.1}, \inlinecode{iterations = 1000}, and \inlinecode{l2 = 0}.
\end{compactitem}

\textbf{Time Complexity}: $O(i \cdot n \cdot d)$ per call for \inlinecode{iterations} $i$, $n$ rows, and $d$ columns.

\textbf{Space Complexity}: $O(d)$ auxiliary.

\begin{lstlisting}[style=cpp]
double sigmoid(double z) {
  return z >= 0 ? 1 / (1 + std::exp(-z)) : std::exp(z) / (1 + std::exp(z));
}

double logistic_predict(const std::vector<double> &w, const std::vector<double> &x) {
  assert(w.size() == x.size() + 1);
  double z = w[0];
  for (size_t j = 0; j < x.size(); j++) {
    z += w[j + 1] * x[j];
  }
  return sigmoid(z);
}

std::vector<double> logistic_regression(
    const Points &points, const std::vector<int> &labels, double rate = 0.1, int iterations = 1000,
    double l2 = 0
) {
  assert(!points.empty() && labels.size() == points.size());
  assert(rate > 0 && iterations > 0 && l2 >= 0);
  assert(std::all_of(labels.begin(), labels.end(), [](int label) {
    return label == 0 || label == 1;
  }));
  int n = static_cast<int>(points.size()), d = static_cast<int>(points[0].size());
  std::vector<double> w(d + 1), grad(d + 1);
  for (int it = 0; it < iterations; it++) {
    std::fill(grad.begin(), grad.end(), 0.0);
    for (int i = 0; i < n; i++) {
      double error = logistic_predict(w, points[i]) - labels[i];
      grad[0] += error;
      for (int j = 0; j < d; j++) {
        grad[j + 1] += error * points[i][j];
      }
    }
    w[0] -= rate * grad[0] / n;
    for (int j = 1; j <= d; j++) {
      w[j] -= rate * (grad[j] / n + l2 * w[j]);
    }
  }
  return w;
}
\end{lstlisting}
\Needspace*{4\baselineskip}
\textbf{Example Usage}\par
\begin{lstlisting}[style=cpp]
#include <cassert>
using namespace std;

int main() {
  // Two well-separated blobs: the first three rows lie near (0, 0) and the last three near
  // (10, 10).
  Points points{{0, 0}, {1, 0}, {0, 1}, {10, 10}, {11, 10}, {10, 11}};
  vector<int> labels{0, 0, 0, 1, 1, 1};

  // Standardizing puts both columns on one scale; the same pair transforms later queries.
  Points scaled(points);
  pair<vector<double>, vector<double>> shift = standardize(scaled);
  vector<double> query{9, 9};
  standardize(query, shift.first, shift.second);
  assert(knn_classify(scaled, labels, query, 3) == 1);

  mt19937 rng(1234567);
  pair<Points, vector<int>> clustered = kmeans(points, 2, rng);
  assert(clustered.first.size() == 2);
  assert(clustered.second[0] == clustered.second[1] && clustered.second[1] == clustered.second[2]);
  assert(clustered.second[3] == clustered.second[4] && clustered.second[4] == clustered.second[5]);
  assert(clustered.second[0] != clustered.second[3]);

  assert(knn_classify(points, labels, {0.5, 0.5}, 3) == 0);
  assert(knn_classify(points, labels, {9, 9}, 3) == 1);

  NaiveBayes nb(points, labels, 2);
  assert(nb.predict({0.5, 0.5}) == 0);
  assert(nb.predict({9, 9}) == 1);
  // Declaring a third class that no row carries leaves it with zero prior, never predicted.
  NaiveBayes sparse(points, labels, 3);
  assert(isinf(sparse.log_likelihood({0.5, 0.5}, 2)));
  assert(sparse.predict({0.5, 0.5}) == 0 && sparse.predict({9, 9}) == 1);

  vector<double> w = logistic_regression(points, labels, 0.5, 2000);
  assert(logistic_predict(w, {0.5, 0.5}) < 0.5);
  assert(logistic_predict(w, {9, 9}) > 0.5);
  return 0;
}
\end{lstlisting}

\section{\texorpdfstring{Game Theory}{6.8 Game Theory}}
\setcounter{section}{8}
\setcounter{subsection}{0}
\subsection{Nim}
\label{sec:6.8.1}
Solves normal-play nim and related impartial games where each position is the disjoint sum of independent piles. In normal play, the player who makes the last move wins. A nim position is losing exactly when the XOR of all pile sizes is $0$.

This is the base case of the Sprague-Grundy theorem: every finite impartial game position is equivalent to a Nim pile of size equal to its Grundy number, and sums of games combine by XOR.

\begin{compactitem}
  \item \inlinecode{nim\_sum(piles)} returns the XOR of all pile sizes.
  \item \inlinecode{first\_player\_wins(piles)} returns whether the player to move has a winning strategy.
  \item \inlinecode{winning\_move(piles)} returns a move (\inlinecode{pile}, \inlinecode{new\_size}) that moves to XOR $0$, or $(-1, -1)$ if the position is losing.
\end{compactitem}

\textbf{Time Complexity}: $O(n)$ per call, where $n$ is the number of piles.

\textbf{Space Complexity}: $O(1)$ auxiliary.

\begin{lstlisting}[style=cpp]
#include <utility>
#include <vector>

int nim_sum(const std::vector<int> &piles) {
  int x = 0;
  for (int p : piles) {
    x ^= p;
  }
  return x;
}

bool first_player_wins(const std::vector<int> &piles) {
  return nim_sum(piles) != 0;
}

std::pair<int, int> winning_move(const std::vector<int> &piles) {
  int x = nim_sum(piles);
  if (x == 0) {
    return {-1, -1};
  }
  for (int i = 0; i < static_cast<int>(piles.size()); i++) {
    if (int target = piles[i] ^ x; target < piles[i]) {
      return {i, target};
    }
  }
  return {-1, -1};
}
\end{lstlisting}
\Needspace*{4\baselineskip}
\textbf{Example Usage}\par
\begin{lstlisting}[style=cpp]
#include <cassert>
using namespace std;

int main() {
  vector<int> piles{3, 4, 5};
  // 3 ^ 4 ^ 5 = 2, so the position is winning.
  assert(nim_sum(piles) == 2);
  assert(first_player_wins(piles));
  auto [pile_idx, new_size] = winning_move(piles);
  piles[pile_idx] = new_size;
  // A winning move always leaves XOR 0 for the opponent.
  assert(nim_sum(piles) == 0);

  vector<int> losing{1, 2, 3};
  assert(!first_player_wins(losing));
  assert((winning_move(losing) == pair<int, int>{-1, -1}));
  return 0;
}
\end{lstlisting}
\subsection{Nim Product}
\label{sec:6.8.2}
Nimbers are the values of impartial games, added by XOR. Under the additional operation of nim multiplication they form a field: the nonnegative integers below $2^{2^k}$ are closed under both operations, with XOR as addition and nim product as multiplication. Writing $\otimes$ for the nim product and $\oplus$ for XOR, nim multiplication is commutative, associative, and distributive over $\oplus$, and is defined recursively by $a \otimes b = \operatorname{mex}\{(a' \otimes b) \oplus (a \otimes b') \oplus (a' \otimes b') : 0 \le a' < a, \, 0 \le b' < b\}$. Its main use is multiplying the Grundy values of independent ``coin-turning'' games whose move sets combine as a product, such as two-dimensional turning games built from one-dimensional ones.

Rather than evaluate the recursive definition directly, this implementation precomputes the nim product of every pair of single bits $2^i \otimes 2^j$ for $i, j < 64$. Each such product follows from the Fermat-like rule $2^{2^k} \otimes 2^{2^k} = 2^{2^k} \oplus 2^{2^k - 1}$ and the fact that distinct Fermat powers multiply like ordinary powers of two. A full product then XORs together the bit-pair products selected by the set bits of the operands. Because every 64-bit value lies in the field of size $2^{64}$, the operation is total on \inlinecode{uint64\_t} and the nonzero values have multiplicative inverses.

\begin{compactitem}
  \item \inlinecode{nim\_product(x, y)} returns the nim product of \inlinecode{x} and \inlinecode{y}.
  \item \inlinecode{nim\_pow(b, e)} returns \inlinecode{b} raised to the \inlinecode{e}-th power under nim multiplication.
  \item \inlinecode{nim\_inverse(b)} returns the multiplicative inverse of a nonzero \inlinecode{b}, so that \inlinecode{nim\_product(b, nim\_inverse(b)) == 1}.
\end{compactitem}

\Needspace*{3\baselineskip}
\textbf{Time Complexity}:
\begin{compactitem}
  \item $O(1)$ table construction on first use (a fixed 64 by 64 table).
  \item $O(64^{2})$ per call to \inlinecode{nim\_product()}; $O(64^{2} \log e)$ per call to \inlinecode{nim\_pow()} and \inlinecode{nim\_inverse()}.
\end{compactitem}

\textbf{Space Complexity}: $O(1)$ for a fixed 64 by 64 table of 64-bit products.

\begin{lstlisting}[style=cpp]
#include <cassert>
#include <cstdint>

uint64_t bit_prod[64][64];

// Build nim products for pairs of single-bit values. Each entry depends only on entries with a
// smaller first index, so a single ascending pass fills the table.
const bool nim_product_init = [] {
  for (int i = 0; i < 64; i++) {
    for (int j = 0; j < 64; j++) {
      if ((i & j) == 0) {
        bit_prod[i][j] = 1ULL << (i | j);  // Distinct Fermat powers multiply like 2^(i|j).
      } else {
        int a = (i & j) & -(i & j);  // Lowest Fermat power shared by both exponents.
        bit_prod[i][j] = bit_prod[i ^ a][j] ^ bit_prod[(i ^ a) | (a - 1)][(j ^ a) | (i & (a - 1))];
      }
    }
  }
  return true;
}();

uint64_t nim_product(uint64_t x, uint64_t y) {
  uint64_t res = 0;
  for (int i = 0; i < 64 && (x >> i) != 0; i++) {
    if ((x >> i) & 1) {
      for (int j = 0; j < 64 && (y >> j) != 0; j++) {
        if ((y >> j) & 1) {
          res ^= bit_prod[i][j];
        }
      }
    }
  }
  return res;
}

uint64_t nim_pow(uint64_t b, uint64_t e) {
  uint64_t res = 1;
  for (; e > 0; e >>= 1) {
    if (e & 1) {
      res = nim_product(res, b);
    }
    b = nim_product(b, b);
  }
  return res;
}

uint64_t nim_inverse(uint64_t b) {
  assert(b != 0);
  // The multiplicative group of the size-2^64 field has order 2^64 - 1, so b^(2^64 - 2) = b^{-1}.
  return nim_pow(b, static_cast<uint64_t>(-2));
}
\end{lstlisting}
\Needspace*{4\baselineskip}
\textbf{Example Usage}\par
\begin{lstlisting}[style=cpp]
#include <cassert>

int main() {
  // Smallest nontrivial products: {0,1,2,3} form the field GF(4) under XOR and nim product.
  assert(nim_product(2, 2) == 3);
  assert(nim_product(2, 3) == 1);
  assert(nim_product(3, 3) == 2);
  assert(nim_product(1, 12345) == 12345);  // 1 is the multiplicative identity.
  assert(nim_product(0, 12345) == 0);

  // Field axioms on a range of values: commutativity, distributivity over XOR, and inverses.
  for (uint64_t a = 0; a < 64; a++) {
    for (uint64_t b = 0; b < 64; b++) {
      assert(nim_product(a, b) == nim_product(b, a));
      for (uint64_t c = 0; c < 8; c++) {
        assert(nim_product(a, b ^ c) == (nim_product(a, b) ^ nim_product(a, c)));
      }
    }
    if (a != 0) {
      assert(nim_product(a, nim_inverse(a)) == 1);
    }
  }
  assert(nim_product(0x1234567890abcdefULL, nim_inverse(0x1234567890abcdefULL)) == 1);
  return 0;
}
\end{lstlisting}
\subsection{Sprague-Grundy Theorem}
\label{sec:6.8.3}
Computes Grundy numbers for finite impartial games. The Grundy number of a game position is the minimum excluded nonnegative integer (MEX) of the Grundy numbers reachable in one move. A position is losing exactly when its Grundy number is $0$.

For a sum of independent impartial games, the combined Grundy number is the XOR of the component Grundy numbers. This is the Sprague-Grundy theorem.

\begin{compactitem}
  \item \inlinecode{subtraction\_game\_grundy(max\_stones, moves)} computes Grundy numbers for a game where a move subtracts one positive value in \inlinecode{moves} from the pile. It returns a vector of length \inlinecode{max\_stones + 1}, whose element \inlinecode{grundy[stones]} is the Grundy number for a pile of that size. Nonpositive moves are ignored.
  \item \inlinecode{sum\_grundy(grundies)} returns the XOR of component Grundy numbers.
\end{compactitem}

\Needspace*{3\baselineskip}
\textbf{Time Complexity}:
\begin{compactitem}
  \item $O(n \cdot m)$ per call to \inlinecode{subtraction\_game\_grundy(max\_stones, moves)}, where $n$ is \inlinecode{max\_stones} and $m$ is the number of allowed moves.
  \item $O(k)$ per call to \inlinecode{sum\_grundy()}, where $k$ is the input length.
\end{compactitem}

\Needspace*{3\baselineskip}
\textbf{Space Complexity}:
\begin{compactitem}
  \item $O(n)$ output and $O(m)$ auxiliary for \inlinecode{subtraction\_game\_grundy(max\_stones, moves)}.
  \item $O(1)$ auxiliary for \inlinecode{sum\_grundy(grundies)}.
\end{compactitem}

\begin{lstlisting}[style=cpp]
#include <cassert>
#include <vector>

int mex(const std::vector<int> &values) {
  std::vector<char> seen(values.size() + 1);
  for (int v : values) {
    if (0 <= v && v < static_cast<int>(seen.size())) {
      seen[v] = true;
    }
  }
  for (int i = 0; i < static_cast<int>(seen.size()); i++) {
    if (!seen[i]) {
      return i;
    }
  }
  return static_cast<int>(seen.size());
}

std::vector<int> subtraction_game_grundy(int max_stones, const std::vector<int> &moves) {
  assert(max_stones >= 0);
  std::vector<int> grundy(max_stones + 1);
  for (int stones = 1; stones <= max_stones; stones++) {
    std::vector<int> reachable;
    for (int m : moves) {
      if (0 < m && m <= stones) {
        reachable.push_back(grundy[stones - m]);
      }
    }
    grundy[stones] = mex(reachable);
  }
  return grundy;
}

int sum_grundy(const std::vector<int> &grundies) {
  int x = 0;
  for (int g : grundies) {
    x ^= g;
  }
  return x;
}
\end{lstlisting}
\Needspace*{4\baselineskip}
\textbf{Example Usage}\par
\begin{lstlisting}[style=cpp]
#include <cassert>
using namespace std;

int main() {
  vector<int> moves{1, 3, 4};
  auto g = subtraction_game_grundy(10, moves);
  // With moves 1, 3, 4, pile size 2 is losing and pile size 4 has mex {0, 1} = 2.
  assert((g == vector<int>{0, 1, 0, 1, 2, 3, 2, 0, 1, 0, 1}));

  vector<int> heaps{g[5], g[7], g[10]};
  // Independent games combine by xor of their Grundy numbers.
  assert(sum_grundy(heaps) == (g[5] ^ g[7] ^ g[10]));
  return 0;
}
\end{lstlisting}
\subsection{Grundy Numbers on DAG}
\label{sec:6.8.4}
Computes Grundy numbers for impartial games whose positions form a directed acyclic graph. Each node is a game position, and each outgoing edge is a legal move. Terminal nodes have Grundy number $0$; all other nodes take the MEX of their successors' Grundy numbers.

The implementation uses memoized DFS. It assumes the graph is acyclic; if cycles are present, the usual finite impartial-game Grundy recurrence is not directly valid without additional analysis.

\begin{compactitem}
  \item \inlinecode{grundy\_on\_dag(g)} returns the Grundy number of every node in graph \inlinecode{g}, where \inlinecode{g[u]} contains all positions reachable from position \inlinecode{u} in one move and every endpoint is a valid node index.
\end{compactitem}

\textbf{Time Complexity}: $O(n + m)$ per call, where $n$ is the number of nodes and $m$ is the number of edges.

\textbf{Space Complexity}: $O(d)$ auxiliary stack space and $O(n + m)$ auxiliary heap space, where $d$ is DFS depth.

\begin{lstlisting}[style=cpp]
#include <algorithm>
#include <cassert>
#include <vector>

std::vector<int> grundy_on_dag(const std::vector<std::vector<int>> &g) {
  assert(std::all_of(g.begin(), g.end(), [&](const auto &edges) {
    return std::all_of(edges.begin(), edges.end(), [&](int v) {
      return 0 <= v && v < static_cast<int>(g.size());
    });
  }));
  std::vector<int> memo(g.size(), -1);
  auto dfs = [&](auto &&dfs, int u) {
    if (memo[u] != -1) {
      return;
    }
    std::vector<char> seen(g[u].size() + 1);
    for (int v : g[u]) {
      dfs(dfs, v);
      int x = memo[v];
      if (x < static_cast<int>(seen.size())) {
        seen[x] = true;
      }
    }
    int res = 0;
    while (res < static_cast<int>(seen.size()) && seen[res]) {
      res++;
    }
    memo[u] = res;
  };
  for (int u = 0; u < static_cast<int>(g.size()); u++) {
    dfs(dfs, u);
  }
  return memo;
}
\end{lstlisting}
\Needspace*{4\baselineskip}
\textbf{Example Usage}\par
\begin{lstlisting}[style=cpp]
#include <cassert>
using namespace std;

int main() {
  vector<vector<int>> g{{1, 2}, {3}, {3, 4}, {}, {}};
  // Terminal nodes have Grundy 0, their predecessors have mex {0} = 1, and node 0 has mex {1} = 0.
  assert((grundy_on_dag(g) == vector<int>{0, 1, 1, 0, 0}));
  return 0;
}
\end{lstlisting}
\subsection{Minimax and Alpha-Beta Pruning}
\label{sec:6.8.5}
Searches a finite two-player, zero-sum, perfect-information game tree using minimax and alpha-beta pruning. Minimax assumes both players play optimally: the maximizing player chooses the child with largest value, and the minimizing player chooses the child with smallest value.

Alpha-beta pruning computes the same value as minimax, but avoids exploring branches that cannot affect the final answer. It is most effective when good moves are searched first.

The example game is ``take $1$ or $2$ stones'': starting with \inlinecode{stones} stones, players alternate taking $1$ or $2$ stones, and the player who takes the last stone wins. The evaluation returns $1$ when the maximizing player wins and $-1$ when the minimizing player wins.

\begin{compactitem}
  \item \inlinecode{minimax(stones, maximizing)} returns the exact game-tree value, where \inlinecode{maximizing} indicates whose turn it is.
  \item \inlinecode{alpha\_beta(stones, maximizing)} returns the same value using alpha-beta pruning.
  \item \inlinecode{best\_take(stones)} returns an optimal first move using alpha-beta pruning, either $1$ or $2$, for positive \inlinecode{stones}.
\end{compactitem}

\textbf{Time Complexity}: $O(b^{d})$ per call in the worst case, where $b$ is branching factor and $d$ is search depth. Alpha-beta can achieve $O(b^{d/2})$ with optimal move ordering.

\textbf{Space Complexity}: $O(d)$ auxiliary recursion stack space.

\begin{lstlisting}[style=cpp]
#include <algorithm>

int minimax(int stones, bool maximizing) {
  if (stones == 0) {
    return maximizing ? -1 : 1;
  }
  int best = maximizing ? -2 : 2;
  for (int take = 1; take <= 2 && take <= stones; take++) {
    int child = minimax(stones - take, !maximizing);
    best = maximizing ? std::max(best, child) : std::min(best, child);
  }
  return best;
}

int alpha_beta(int stones, bool maximizing, int alpha = -2, int beta = 2) {
  if (stones == 0) {
    return maximizing ? -1 : 1;
  }
  int value = maximizing ? -2 : 2;
  for (int take = 1; take <= 2 && take <= stones; take++) {
    int child = alpha_beta(stones - take, !maximizing, alpha, beta);
    if (maximizing) {
      value = std::max(value, child);
      alpha = std::max(alpha, value);
    } else {
      value = std::min(value, child);
      beta = std::min(beta, value);
    }
    if (alpha >= beta) {
      break;
    }
  }
  return value;
}

int best_take(int stones) {
  int best_move = 1, best_value = -2;
  for (int take = 1; take <= 2 && take <= stones; take++) {
    int value = alpha_beta(stones - take, false);
    if (value > best_value) {
      best_value = value;
      best_move = take;
    }
  }
  return best_move;
}
\end{lstlisting}
\Needspace*{4\baselineskip}
\textbf{Example Usage}\par
\begin{lstlisting}[style=cpp]
#include <cassert>

int main() {
  assert(minimax(1, true) == 1);
  assert(minimax(3, true) == -1);
  for (int stones = 1; stones <= 10; stones++) {
    assert(minimax(stones, true) == alpha_beta(stones, true));
  }
  assert(best_take(4) == 1);  // Move to losing state with 3 stones.
  assert(best_take(5) == 2);  // Move to losing state with 3 stones.
  return 0;
}
\end{lstlisting}

\chapter{\texorpdfstring{Geometry}{7. Geometry}}

\section{\texorpdfstring{Geometry Primitives}{7.1 Geometry Primitives}}
\setcounter{section}{1}
\setcounter{subsection}{0}
\subsection{Point}
\label{sec:7.1.1}
A 2D point class templated on the coordinate type \inlinecode{T}. Algebraic operations preserve \inlinecode{T}, while metric operations use \inlinecode{fp\_t}, which is \inlinecode{T} for floating-point coordinates and \inlinecode{double} otherwise.

Non-floating-point coordinate types such as \inlinecode{Modular} or \inlinecode{Rational} compose for exact operations (arithmetic, \inlinecode{dot}, \inlinecode{cross}, \inlinecode{sqnorm}, comparisons, and \inlinecode{EQ}) using the coordinate type's own operators. Their metric operations convert coordinates explicitly to \inlinecode{double}; for \inlinecode{Modular}, this uses the stored representative and is not finite-field geometry.

\begin{compactitem}
  \item \inlinecode{TPoint\textless{}T\textgreater{}()} constructs the origin, while \inlinecode{TPoint\textless{}T\textgreater{}(x, y)} constructs point $(x, y)$. A pair of values can also be converted explicitly to a point.
  \item Operators \inlinecode{+}, \inlinecode{-}, \inlinecode{*}, \inlinecode{/}, and their compound forms act element-wise on points or apply a scalar. Comparisons are exact and lexicographic, as required by standard algorithms and containers. \inlinecode{EQ(p, q)} instead uses \inlinecode{EPS} for floating-point coordinates and remains exact for other coordinate types.
  \item \inlinecode{sqnorm()}, \inlinecode{dot(p)}, and \inlinecode{cross(p)} return the exact algebraic results in coordinate type \inlinecode{T}.
  \item \inlinecode{norm()}, \inlinecode{arg()}, \inlinecode{proj(p)}, and \inlinecode{normalize()} return metric results in \inlinecode{fp\_t}; division also promotes to \inlinecode{fp\_t} to avoid integer truncation.
  \item \inlinecode{rotate90()}, \inlinecode{rotate180()}, and \inlinecode{rotate270()} rotate by the corresponding counter-clockwise cardinal angle; overloads taking \inlinecode{p} rotate about point \inlinecode{p}.
  \item \inlinecode{rotate\_cw(t)} and \inlinecode{rotate\_ccw(t)} rotate by an arbitrary angle \inlinecode{t} in radians; overloads taking (\inlinecode{p}, \inlinecode{t}) rotate about point \inlinecode{p}.
  \item \inlinecode{reflect(p)} reflects across point \inlinecode{p}, while \inlinecode{reflect(p, q)} reflects across the line through points \inlinecode{p} and \inlinecode{q}.
  \item \inlinecode{to\_double()} and \inlinecode{to\_ldouble()} explicitly convert the coordinate type. Integral points also convert implicitly to floating-point point types, so \inlinecode{PointI} can be passed where \inlinecode{PointD} is expected.
  \item \inlinecode{PointI}, \inlinecode{PointL}, \inlinecode{PointD}, and \inlinecode{PointLD} use \inlinecode{int}, \inlinecode{long long}, \inlinecode{double}, and \inlinecode{long double} coordinates, respectively. \inlinecode{Point} aliases \inlinecode{PointD}.
\end{compactitem}

Overflow warning: the exact products \inlinecode{dot()}, \inlinecode{cross()}, and \inlinecode{sqnorm()} grow like the squared coordinate magnitude. With \inlinecode{PointI} these overflow a 32-bit \inlinecode{int} once coordinates exceed roughly a few tens of thousands, so use \inlinecode{PointL} for larger integer coordinates.

\textbf{Time Complexity}: $O(1)$ per call to the constructor and all other operations.

\Needspace*{3\baselineskip}
\textbf{Space Complexity}:
\begin{compactitem}
  \item $O(1)$ for storage of the point.
  \item $O(1)$ auxiliary for all operations.
\end{compactitem}

\begin{lstlisting}[style=cpp]
#include <cmath>
#include <ostream>
#include <tuple>
#include <type_traits>
#include <utility>

const double EPS = 1e-9;

// Epsilon-aware for floating-point coordinates; exact for int and for coordinate types like Modular
// or Rational, which therefore compose for all of the predicates below.
template<typename T, typename U, typename C = std::common_type_t<T, U>>
bool EQ(T a, U b) {
  if constexpr (std::is_floating_point_v<C>) return C(a) == C(b) || std::fabs(C(a) - C(b)) <= EPS;
  return C(a) == C(b);
}

template<typename T, typename U, typename C = std::common_type_t<T, U>>
bool LT(T a, U b) {
  if constexpr (std::is_floating_point_v<C>) return C(a) < C(b) - EPS;
  return C(a) < C(b);
}

template<typename T>
struct TPoint {
  // Metric operations preserve floating-point coordinates and otherwise promote to double.
  using fp_t = std::conditional_t<std::is_floating_point<T>::value, T, double>;

  T x, y;

  TPoint() : x(0), y(0) {}
  TPoint(T x, T y) : x(x), y(y) {}
  explicit TPoint(const std::pair<T, T> &p) : x(p.first), y(p.second) {}

  // Implicit conversion from integral to floating-point point (optional, can just use to_double()).
  template<
      typename U,
      typename = std::enable_if_t<std::is_integral<U>::value && std::is_floating_point<T>::value>>
  TPoint(const TPoint<U> &p) : x(static_cast<T>(p.x)), y(static_cast<T>(p.y)) {}

  friend bool EQ(const TPoint &a, const TPoint &b) { return EQ(a.x, b.x) && EQ(a.y, b.y); }

  bool operator==(const TPoint &p) const { return x == p.x && y == p.y; }
  bool operator!=(const TPoint &p) const { return !(*this == p); }
  bool operator<(const TPoint &p) const { return std::tie(x, y) < std::tie(p.x, p.y); }
  bool operator>(const TPoint &p) const { return p < *this; }
  bool operator<=(const TPoint &p) const { return !(*this > p); }
  bool operator>=(const TPoint &p) const { return !(*this < p); }
  TPoint operator+(const TPoint &p) const { return {x + p.x, y + p.y}; }
  TPoint operator-(const TPoint &p) const { return {x - p.x, y - p.y}; }
  TPoint operator+(T v) const { return {x + v, y + v}; }
  TPoint operator-(T v) const { return {x - v, y - v}; }
  TPoint operator*(T v) const { return {x * v, y * v}; }

  // Division always promotes to fp_t to avoid integer truncation.
  TPoint<fp_t> operator/(fp_t v) const { return {fp_t(x) / v, fp_t(y) / v}; }

  TPoint &operator+=(const TPoint &p) { x += p.x; y += p.y; return *this; }
  TPoint &operator-=(const TPoint &p) { x -= p.x; y -= p.y; return *this; }
  TPoint &operator+=(T v) { x += v; y += v; return *this; }
  TPoint &operator-=(T v) { x -= v; y -= v; return *this; }
  TPoint &operator*=(T v) { x *= v; y *= v; return *this; }
  friend TPoint operator+(T v, const TPoint &p) { return p + v; }
  friend TPoint operator*(T v, const TPoint &p) { return p * v; }

  // --- Exact operations: return T or TPoint<T>, work for any coordinate type ---

  T sqnorm() const { return x * x + y * y; }                    // Overflow warning.
  T dot(const TPoint &p) const { return x * p.x + y * p.y; }    // Overflow warning.
  T cross(const TPoint &p) const { return x * p.y - y * p.x; }  // Overflow warning.

  // --- Floating-point operations: return fp_t or TPoint<fp_t> ---

  fp_t norm() const { return std::hypot(fp_t(x), fp_t(y)); }
  fp_t arg() const { return std::atan2(fp_t(y), fp_t(x)); }
  fp_t proj(const TPoint &p) const { return (fp_t(x) * p.x + fp_t(y) * p.y) / p.norm(); }

  TPoint<fp_t> normalize() const {
    fp_t n = norm();
    if (n < 1e-30) {              // guard against dividing by a near-zero norm, not a
      return {fp_t{0}, fp_t{0}};  // geometric tolerance (EPS is too coarse here)
    }
    return {fp_t(x) / n, fp_t(y) / n};
  }

  // --- Cardinal rotations: exact for all coordinate types including int ---

  // Returns (x, y) rotated 90/180/270 degrees counter-clockwise about the origin.
  TPoint rotate90() const { return {-y, x}; }
  TPoint rotate180() const { return {-x, -y}; }
  TPoint rotate270() const { return {y, -x}; }

  // Returns (x, y) rotated 90/180/270 degrees counter-clockwise about point p.
  TPoint rotate90(const TPoint &p) const { return (*this - p).rotate90() + p; }
  TPoint rotate180(const TPoint &p) const { return (*this - p).rotate180() + p; }
  TPoint rotate270(const TPoint &p) const { return (*this - p).rotate270() + p; }

  // --- Arbitrary-angle rotations: always return floating-point ---

  // Returns (x, y) rotated t radians clockwise about the origin.
  TPoint<fp_t> rotate_cw(fp_t t) const {
    fp_t fx(x), fy(y);
    return {fx * std::cos(t) + fy * std::sin(t), fy * std::cos(t) - fx * std::sin(t)};
  }

  // Returns (x, y) rotated t radians counter-clockwise about the origin.
  TPoint<fp_t> rotate_ccw(fp_t t) const {
    fp_t fx(x), fy(y);
    return {fx * std::cos(t) - fy * std::sin(t), fx * std::sin(t) + fy * std::cos(t)};
  }

  // Returns (x, y) rotated t radians clockwise about point p.
  TPoint<fp_t> rotate_cw(const TPoint &p, fp_t t) const {
    return TPoint<fp_t>{fp_t(x) - p.x, fp_t(y) - p.y}.rotate_cw(t) +
           TPoint<fp_t>{fp_t(p.x), fp_t(p.y)};
  }

  // Returns (x, y) rotated t radians counter-clockwise about point p.
  TPoint<fp_t> rotate_ccw(const TPoint &p, fp_t t) const {
    return TPoint<fp_t>{fp_t(x) - p.x, fp_t(y) - p.y}.rotate_ccw(t) +
           TPoint<fp_t>{fp_t(p.x), fp_t(p.y)};
  }

  // --- Reflections ---

  // Returns (x, y) reflected across point p. Exact for any coordinate type.
  TPoint reflect(const TPoint &p) const { return {2 * p.x - x, 2 * p.y - y}; }

  // Returns (x, y) reflected across the line containing points p and q.
  // Always returns floating-point coordinates.
  TPoint<fp_t> reflect(const TPoint &p, const TPoint &q) const {
    TPoint<fp_t> fp{fp_t(p.x), fp_t(p.y)};
    if (EQ(p, q)) {
      return TPoint<fp_t>{fp_t(x), fp_t(y)}.reflect(fp);
    }
    TPoint<fp_t> r{fp_t(x) - p.x, fp_t(y) - p.y};
    TPoint<fp_t> s{fp_t(q.x) - p.x, fp_t(q.y) - p.y};
    fp_t ssq = s.sqnorm();
    r = TPoint<fp_t>{(r.x * s.x + r.y * s.y) / ssq, (r.x * s.y - r.y * s.x) / ssq};
    return TPoint<fp_t>{r.x * s.x - r.y * s.y + fp.x, r.x * s.y + r.y * s.x + fp.y};
  }

  // --- Explicit type conversions ---

  TPoint<double> to_double() const { return {static_cast<double>(x), static_cast<double>(y)}; }

  TPoint<long double> to_ldouble() const {
    return {static_cast<long double>(x), static_cast<long double>(y)};
  }

  // --- Friend free-function versions ---

  friend T sqnorm(const TPoint &p) { return p.sqnorm(); }
  friend fp_t norm(const TPoint &p) { return p.norm(); }
  friend fp_t arg(const TPoint &p) { return p.arg(); }
  friend T dot(const TPoint &p, const TPoint &q) { return p.dot(q); }
  friend T cross(const TPoint &p, const TPoint &q) { return p.cross(q); }
  friend fp_t proj(const TPoint &p, const TPoint &q) { return p.proj(q); }
  friend TPoint<fp_t> normalize(const TPoint &p) { return p.normalize(); }
  friend TPoint rotate90(const TPoint &p) { return p.rotate90(); }
  friend TPoint rotate180(const TPoint &p) { return p.rotate180(); }
  friend TPoint rotate270(const TPoint &p) { return p.rotate270(); }
  friend TPoint rotate90(const TPoint &p, const TPoint &q) { return p.rotate90(q); }
  friend TPoint rotate180(const TPoint &p, const TPoint &q) { return p.rotate180(q); }
  friend TPoint rotate270(const TPoint &p, const TPoint &q) { return p.rotate270(q); }
  friend TPoint<fp_t> rotate_cw(const TPoint &p, fp_t t) { return p.rotate_cw(t); }
  friend TPoint<fp_t> rotate_ccw(const TPoint &p, fp_t t) { return p.rotate_ccw(t); }
  friend TPoint<fp_t> rotate_cw(const TPoint &p, const TPoint &q, fp_t t) { return p.rotate_cw(q, t); }
  friend TPoint<fp_t> rotate_ccw(const TPoint &p, const TPoint &q, fp_t t) { return p.rotate_ccw(q, t); }
  friend TPoint reflect(const TPoint &p, const TPoint &q) { return p.reflect(q); }
  friend TPoint<fp_t> reflect(const TPoint &p, const TPoint &a, const TPoint &b) { return p.reflect(a, b); }

  friend std::ostream &operator<<(std::ostream &out, const TPoint &p) {
    if constexpr (std::is_floating_point<T>::value) {
      return out << "(" << (std::fabs(p.x) < EPS ? 0 : p.x) << ","
                 << (std::fabs(p.y) < EPS ? 0 : p.y) << ")";
    }
    return out << "(" << p.x << "," << p.y << ")";
  }
};

using PointI = TPoint<int>;
using PointL = TPoint<long long>;
using PointD = TPoint<double>;
using PointLD = TPoint<long double>;
using Point = PointD;  // Default point type is double.

// Can compose with numerical types from chapter 6:
// using PointB = TPoint<BigInt>;
// using PointR = TPoint<Rational<int64_t>>;
// using PointM = TPoint<Modular<1000000007>>;
\end{lstlisting}
\Needspace*{4\baselineskip}
\textbf{Example Usage}\par
\begin{lstlisting}[style=cpp]
#include <cassert>
using namespace std;

const double PI = acos(-1.0);

int main() {
  Point p(-10, 3);
  assert(EQ(Point(-18, 29), p + Point(-3, 9) * 6.0 / 2.0 - Point(-1, 1)));
  assert(EQ(109, p.sqnorm()));
  assert(EQ(10.44030650891, p.norm()));
  assert(EQ(2.850135859112, p.arg()));
  assert(EQ(0, p.dot(Point(3, 10))));
  assert(EQ(0, p.cross(Point(10, -3))));
  assert(EQ(10, p.proj(Point(-10, 0))));
  assert(EQ(1, p.normalize().norm()));
  assert(EQ(Point(-3, -10), p.rotate90()));
  assert(EQ(Point(10, -3), p.rotate180()));
  assert(EQ(Point(3, 10), p.rotate270()));
  assert(EQ(Point(3, 12), p.rotate_cw(Point(1, 1), PI / 2)));
  assert(EQ(Point(1, -10), p.rotate_ccw(Point(2, 2), PI / 2)));
  assert(EQ(Point(10, -3), p.reflect(Point(0, 0))));
  assert(EQ(Point(-10, -3), p.reflect(Point(-2, 0), Point(5, 0))));

  // Integer point - exact arithmetic, float-only ops return PointD.
  PointI a(3, 4), b(1, 0);
  assert(a + b == PointI(4, 4));
  assert(a - b == PointI(2, 4));
  assert(a * 2 == PointI(6, 8));
  assert(a.sqnorm() == 25);
  assert(a.dot(b) == 3);
  assert(a.cross(b) == -4);
  assert(EQ(a.norm(), 5.0));
  assert(a.rotate90() == PointI(-4, 3));
  assert(a.rotate180() == PointI(-3, -4));
  assert(a.rotate270() == PointI(4, -3));
  PointD anorm = a.normalize();  // returns PointD
  assert(EQ(anorm.norm(), 1.0));
  PointD adiv = a / 2.0;  // returns PointD (division always promotes)
  assert(EQ(adiv.x, 1.5) && EQ(adiv.y, 2.0));
  auto rotated = PointI(1, 0).rotate_ccw(PI / 4);  // returns PointD (arbitrary rotation promotes)
  double root2over2 = sqrt(2.0) / 2.0;
  assert(EQ(rotated.x, root2over2) && EQ(rotated.y, root2over2));

  // Implicit conversion PointI -> PointD.
  PointD promoted = PointI(1, 2);
  assert(promoted == PointD(1, 2));
  return 0;
}
\end{lstlisting}
\subsection{Line}
\label{sec:7.1.2}
A straight line in two dimensions, represented by coefficients of $ax + by + c = 0$ and templated on their type \inlinecode{T}. Coefficients are canonicalized by default so proportional valid coefficient vectors have the same representation, up to floating-point rounding.

The template flag \inlinecode{CANONICALIZE} controls whether construction normalizes the coefficients and defaults to \inlinecode{true}. When enabled, \inlinecode{EXACT = true} divides coefficients by their Euclidean GCD and fixes their sign, while \inlinecode{EXACT = false} scales a nonvertical line to $b = 1$ and a vertical line to $a = 1$. \inlinecode{EXACT} defaults to \inlinecode{true} for built-in integral types; set it explicitly for custom exact types such as \inlinecode{BigInt} and \inlinecode{Rational}. Exact canonicalization requires \inlinecode{\%}, comparisons, and basic arithmetic, and predicates remain exact for integral \inlinecode{T}. Floating-point and \inlinecode{Modular} coefficients use the scaling form.

\begin{compactitem}
  \item \inlinecode{TLine\textless{}T\textgreater{}(a, b, c)} constructs a line from its coefficients, while \inlinecode{TLine\textless{}T\textgreater{}(p, q)} constructs the line through distinct points \inlinecode{p} and \inlinecode{q}. For floating-point \inlinecode{T}, \inlinecode{TLine\textless{}T\textgreater{}(slope, p)} constructs the line of the given slope through \inlinecode{p}. The default constructor produces an invalid line.
  \item \inlinecode{canonicalized()} returns a canonical copy. \inlinecode{operator==} compares stored coefficients exactly, while \inlinecode{EQ(l1, l2)} canonicalizes raw lines if necessary and uses \inlinecode{EPS} for floating-point coefficients.
  \item \inlinecode{valid()}, \inlinecode{horizontal()}, and \inlinecode{vertical()} classify the line.
  \item \inlinecode{slope()}, \inlinecode{x(y)}, and \inlinecode{y(x)} return \inlinecode{fp\_t}, which is \inlinecode{T} for floating-point coefficients and \inlinecode{double} otherwise. An undefined result is returned as \inlinecode{NaN}.
  \item \inlinecode{contains(p)}, \inlinecode{is\_parallel(l)}, and \inlinecode{is\_perpendicular(l)} test geometric relationships.
  \item \inlinecode{parallel(p)} and \inlinecode{perpendicular(p)} return the corresponding line through point \inlinecode{p}.
  \item \inlinecode{operator\textless{}\textless{}} outputs the line in the form \inlinecode{ax+by+c=0}.
  \item \inlinecode{LineI}, \inlinecode{LineL}, \inlinecode{LineD}, and \inlinecode{LineLD} use \inlinecode{int}, \inlinecode{int64\_t}, \inlinecode{double}, and \inlinecode{long double} coefficients, respectively. \inlinecode{Line} aliases \inlinecode{LineD}.
\end{compactitem}

Overflow warning: the exact predicates form products of coefficients, which grow like the squared coordinate magnitude. With \inlinecode{LineI} these may overflow once coordinates reach a few tens of thousands, so use \inlinecode{LineL} for larger integer coordinates.

\Needspace*{3\baselineskip}
\textbf{Time Complexity}:
\begin{compactitem}
  \item $O(\log C)$ per integral canonicalizing constructor or call to \inlinecode{canonicalize()}, where $C$ is the largest absolute coefficient.
  \item $O(1)$ per call to all other operations.
\end{compactitem}

\Needspace*{3\baselineskip}
\textbf{Space Complexity}:
\begin{compactitem}
  \item $O(1)$ for storage of the line.
  \item $O(1)$ auxiliary for all operations.
\end{compactitem}

\begin{lstlisting}[style=cpp]
#include <cmath>
#include <cstdint>
#include <ios>
#include <limits>
#include <ostream>
#include <type_traits>
#include <utility>

const double EPS = 1e-9;
const double M_NAN = std::numeric_limits<double>::quiet_NaN();

// Epsilon-aware for floating-point coordinates; exact for int and for coordinate types like Modular
// or Rational, which therefore compose for all of the predicates below.
template<typename T, typename U, typename C = std::common_type_t<T, U>>
bool EQ(T a, U b) {
  if constexpr (std::is_floating_point_v<C>) return C(a) == C(b) || std::fabs(C(a) - C(b)) <= EPS;
  return C(a) == C(b);
}

// SFINAE guard: valid only when Pt exposes numeric .x/.y members, so templated point constructors
// don't hijack scalar-argument calls.
template<typename Pt>
using if_point = decltype(std::declval<const Pt &>().x, std::declval<const Pt &>().y, void());

template<typename T, bool CANONICALIZE = true, bool EXACT = std::is_integral<T>::value>
struct TLine {
  // Slope and solve operations preserve floating-point coordinates and otherwise promote to double.
  using fp_t = std::conditional_t<std::is_floating_point<T>::value, T, double>;

  T a, b, c;

  void canonicalize() {
    if constexpr (EXACT) {
      auto gcd = [](T x, T y) {
        if (x < T{0}) x = -x;
        if (y < T{0}) y = -y;
        while (y != T{0}) {
          x %= y;
          std::swap(x, y);
        }
        return x;
      };
      T g = gcd(a, gcd(b, c));
      if (g != 0) {
        a /= g;
        b /= g;
        c /= g;
      }
      if (a < 0 || (a == 0 && b < 0)) {
        a = -a;
        b = -b;
        c = -c;
      }
    } else if (!EQ(b, T{0})) {
      a /= b;
      c /= b;
      b = T{1};
    } else if (!EQ(a, T{0})) {
      c /= a;
      a = T{1};
      b = T{0};
    }
  }

  TLine() : a(0), b(0), c(0) {}  // Invalid or uninitialized line.

  TLine(T a, T b, T c) : a(a), b(b), c(c) {
    if constexpr (CANONICALIZE) canonicalize();
  }

  // Line through two points. Coefficients are exact for integral points (overflow warning for c).
  template<typename Pt, typename = if_point<Pt>>
  TLine(const Pt &p, const Pt &q) : a(q.y - p.y), b(p.x - q.x), c(-(a * p.x + b * p.y)) {
    if constexpr (CANONICALIZE) canonicalize();
  }

  // Line of a given slope through a point. Intended for floating-point T (a generic slope
  // produces non-integer coefficients), so this constructor is disabled for integer T.
  template<
      class Pt, typename = if_point<Pt>, typename U = T,
      class = std::enable_if_t<std::is_floating_point<U>::value>>
  TLine(fp_t slope, const Pt &p) : a(-slope), b(1), c(slope * p.x - p.y) {
    if constexpr (CANONICALIZE) canonicalize();
  }

  TLine<T, true, EXACT> canonicalized() const { return TLine<T, true, EXACT>(a, b, c); }

  bool operator==(const TLine &l) const { return a == l.a && b == l.b && c == l.c; }
  bool operator!=(const TLine &l) const { return !(*this == l); }

  friend bool EQ(const TLine &l1, const TLine &l2) {
    if constexpr (CANONICALIZE) {
      return EQ(l1.a, l2.a) && EQ(l1.b, l2.b) && EQ(l1.c, l2.c);
    }
    return EQ(l1.canonicalized(), l2.canonicalized());
  }

  // Returns whether the line is a valid line (not both a and b zero).
  bool valid() const { return !(EQ(a, 0) && EQ(b, 0)); }
  bool horizontal() const { return valid() && EQ(a, 0); }
  bool vertical() const { return valid() && EQ(b, 0); }

  // Slope -a/b, or NaN if the line is vertical or invalid.
  fp_t slope() const { return (!valid() || EQ(b, 0)) ? M_NAN : -fp_t(a) / b; }

  // Solve for x at a given y (NaN if horizontal or invalid).
  fp_t x(fp_t y) const { return (!valid() || EQ(a, 0)) ? M_NAN : -(fp_t(b) * y + c) / a; }

  // Solve for y at a given x (NaN if vertical or invalid).
  fp_t y(fp_t x) const { return (!valid() || EQ(b, 0)) ? M_NAN : -(fp_t(a) * x + c) / b; }

  // Whether point p lies on the line. Exact for integer T and integer p.
  template<typename Pt>
  bool contains(const Pt &p) const {
    return valid() && EQ(a * p.x + b * p.y + c, 0);  // Overflow warning.
  }

  // Parallel iff the normals are parallel (cross == 0); perpendicular iff normals are perpendicular
  // (dot == 0). Both exact for integer T. Overflow warning.
  bool is_parallel(const TLine &l) const { return valid() && l.valid() && EQ(a * l.b, l.a * b); }

  bool is_perpendicular(const TLine &l) const {
    return valid() && l.valid() && EQ(a * l.a, -(b * l.b));
  }

  // Parallel/perpendicular line through point p. Exact for integer T and integer p.
  template<typename Pt, typename = if_point<Pt>>
  TLine parallel(const Pt &p) const {
    return TLine(a, b, -(a * p.x + b * p.y));  // Overflow warning.
  }

  template<typename Pt, typename = if_point<Pt>>
  TLine perpendicular(const Pt &p) const {
    return TLine(-b, a, b * p.x - a * p.y);  // Overflow warning.
  }

  friend std::ostream &operator<<(std::ostream &out, const TLine &l) {
    auto clean = [](T v) {
      if constexpr (std::is_floating_point_v<T>) {
        return std::fabs(v) < EPS ? T{0} : v;
      }
      return v;
    };
    auto flags = out.flags();
    out << clean(l.a) << "x" << std::showpos << clean(l.b) << "y" << clean(l.c) << "=0";
    out.flags(flags);
    return out;
  }
};

using LineI = TLine<int>;
using LineL = TLine<int64_t>;
using LineD = TLine<double>;
using LineLD = TLine<long double>;
using Line = LineD;  // Default line type is double.

// Can compose with numerical types from chapter 6:
// using LineB = TLine<BigInt, true, true>;
// using LineR = TLine<Rational<int64_t>, true, true>;
\end{lstlisting}
\Needspace*{4\baselineskip}
\textbf{Example Usage}\par
\begin{lstlisting}[style=cpp]
#include <cassert>

struct Point {
  double x, y;
  Point(double x = 0, double y = 0) : x(x), y(y) {}
};

struct PointI {
  int x, y;
  PointI(int x = 0, int y = 0) : x(x), y(y) {}
};

int main() {
  Line invalid;
  assert(!invalid.contains(Point(0, 0)) && !invalid.is_parallel(invalid));

  // Floating-point line.
  Line l(2, -5, -8);
  Line para = Line(2, -5, -8).parallel(Point(-6, -2));
  Line perp = Line(2, -5, -8).perpendicular(Point(-6, -2));
  assert(l.is_parallel(para) && l.is_perpendicular(perp));
  assert(l.slope() == 0.4);
  assert(EQ(para, Line(-0.4, 1, -0.4)));
  assert(EQ(perp, Line(2.5, 1, 17)));
  assert(l.contains(Point(1.5, -1)));  // 2(1.5) - 5(-1) - 8 = 0.

  // Integer line through integer points - exact canonical coefficients and predicates.
  LineI il(PointI(0, 0), PointI(2, 1));  // a=1, b=-2, c=0  ->  x - 2y = 0
  assert(il.a == 1 && il.b == -2 && il.c == 0);
  assert(il.contains(PointI(4, 2)));  // exact integer check
  assert(!il.contains(PointI(4, 3)));
  assert(il.is_parallel(LineI(PointI(1, 1), PointI(3, 2))));        // both slope 1/2
  assert(il.is_perpendicular(LineI(PointI(0, 0), PointI(1, -2))));  // dot of normals = 0
  assert(il.slope() == 0.5);                                        // fp_t result

  LineI reduced(2, -4, -6);
  assert(reduced == LineI(1, -2, -3));

  using RawLine = TLine<double, false>;
  RawLine raw(2, -5, -8), scaled(-0.4, 1, 1.6);
  assert(raw != scaled);    // Stored coefficients differ.
  assert(EQ(raw, scaled));  // Canonicalized copies represent the same line.
  RawLine canonical = raw;
  canonical.canonicalize();
  assert(canonical == scaled);
  assert(EQ(raw.canonicalized(), Line(-0.4, 1, 1.6)));
  return 0;
}
\end{lstlisting}
\subsection{Circle}
\label{sec:7.1.3}
A circle in two dimensions with epsilon-aware operations. The circle centered at $(h, k)$ is represented by the relation $(x - h)^2 + (y - k)^2 = r^2$, where the radius $r$ is normalized to a nonnegative number.

This implementation stores and computes circle parameters in \inlinecode{double}. Point-accepting constructors and predicates are templated on the point type \inlinecode{Pt} and only read \inlinecode{.x}/\inlinecode{.y}, so they accept \inlinecode{Point}, \inlinecode{PointD}, or \inlinecode{PointI} from \secref{7.1.1} or any struct with numeric \inlinecode{.x} and \inlinecode{.y} fields.

\begin{compactitem}
  \item \inlinecode{Circle()} constructs the zero-radius circle centered at the origin.
  \item \inlinecode{Circle(r)} constructs a circle of radius \inlinecode{abs(r)} centered at the origin.
  \item \inlinecode{Circle(h, k, r)} constructs a circle of radius \inlinecode{abs(r)} centered at (\inlinecode{h}, \inlinecode{k}).
  \item \inlinecode{Circle(o, r)} constructs a circle of radius \inlinecode{abs(r)} centered at point \inlinecode{o}.
  \item \inlinecode{Circle(a, b)} constructs the circle whose diameter is segment \inlinecode{a}-\inlinecode{b}.
  \item \inlinecode{Circle(a, b, c)} constructs the circumcircle through non-collinear points \inlinecode{a}, \inlinecode{b}, and \inlinecode{c}, throwing if the points are collinear.
  \item \inlinecode{Circle(a, b, r)} constructs one circle of radius \inlinecode{abs(r)} passing through points \inlinecode{a} and \inlinecode{b}, throwing if no finite circle is determined by the inputs.
  \item \inlinecode{operator==} and \inlinecode{operator!=} compare stored centers and radii exactly; \inlinecode{EQ(a, b)} compares them using \inlinecode{EPS}.
  \item \inlinecode{contains(p)} returns whether point \inlinecode{p} lies inside or on the circle.
  \item \inlinecode{on\_boundary(p)} returns whether point \inlinecode{p} lies on the circle boundary.
  \item \inlinecode{incircle(a, b, c)} returns the circle inscribed in triangle \inlinecode{abc}.
  \item \inlinecode{in\_circumcircle(a, b, c, d)} returns $1$ if point \inlinecode{d} lies strictly inside the circle through points \inlinecode{a}, \inlinecode{b}, and \inlinecode{c}, $0$ if it lies on the circle, or $-1$ if it lies outside. This has exact reasoning for integral points by evaluating the determinant directly without constructing a \inlinecode{Circle} or taking square roots. Points \inlinecode{a}, \inlinecode{b}, and \inlinecode{c} must not be collinear. For integer inputs, the sign test is exact provided the chosen intermediate type does not overflow.
\end{compactitem}

Overflow warning: the determinant in \inlinecode{in\_circumcircle()} grows like the fourth power of the coordinate magnitude. For large integer coordinates, use a point or intermediate type wide enough for that determinant.

\textbf{Time Complexity}: $O(1)$ per call to the constructors and all other operations.

\Needspace*{3\baselineskip}
\textbf{Space Complexity}:
\begin{compactitem}
  \item $O(1)$ for storage of the circle.
  \item $O(1)$ auxiliary for all operations.
\end{compactitem}

\begin{lstlisting}[style=cpp]
#include <cmath>
#include <cstdint>
#include <ios>
#include <ostream>
#include <stdexcept>
#include <type_traits>

const double EPS = 1e-9;

template<typename T, typename U, typename C = std::common_type_t<T, U>>
bool EQ(T a, U b) {
  if constexpr (std::is_floating_point_v<C>) return C(a) == C(b) || std::fabs(C(a) - C(b)) <= EPS;
  return C(a) == C(b);
}

template<typename T, typename U, typename C = std::common_type_t<T, U>>
bool LT(T a, U b) {
  if constexpr (std::is_floating_point_v<C>) return C(a) < C(b) - EPS;
  return C(a) < C(b);
}

template<typename T, typename U>
bool LE(T a, U b) {
  return !LT(b, a);
}

// SFINAE guard: valid only when Pt exposes numeric .x/.y members. This keeps the templated point
// constructors from hijacking calls like Circle(double, double, double).
template<typename Pt>
using if_point = decltype(std::declval<const Pt &>().x, std::declval<const Pt &>().y, void());

struct Circle {
  double h, k, r;

  Circle() : h(0), k(0), r(0) {}
  explicit Circle(double r) : h(0), k(0), r(fabs(r)) {}
  Circle(double h, double k, double r) : h(h), k(k), r(fabs(r)) {}

  template<typename Pt, typename = if_point<Pt>>
  Circle(const Pt &o, double r) : h(o.x), k(o.y), r(fabs(r)) {}

  template<typename Pt, typename = if_point<Pt>>
  Circle(const Pt &a, const Pt &b) {
    h = (a.x + b.x) / 2.0;
    k = (a.y + b.y) / 2.0;
    r = std::hypot(a.x - h, a.y - k);
  }

  template<typename Pt, typename = if_point<Pt>>
  Circle(const Pt &a, const Pt &b, const Pt &c) {
    double ax = a.x, ay = a.y, bx = b.x, by = b.y, cx = c.x, cy = c.y;
    double asq = ax * ax + ay * ay, bsq = bx * bx + by * by, csq = cx * cx + cy * cy;
    double d = 2.0 * (ax * (by - cy) + bx * (cy - ay) + cx * (ay - by));
    if (EQ(d, 0)) {
      throw std::runtime_error("No circumcircle from collinear points.");
    }
    h = (asq * (by - cy) + bsq * (cy - ay) + csq * (ay - by)) / d;
    k = (asq * (cx - bx) + bsq * (ax - cx) + csq * (bx - ax)) / d;
    r = std::hypot(ax - h, ay - k);
  }

  template<typename Pt, typename = if_point<Pt>>
  Circle(const Pt &a, const Pt &b, double r) : r(fabs(r)) {
    if (EQ(this->r, 0) && EQ(a.x, b.x) && EQ(a.y, b.y)) {
      h = a.x;
      k = a.y;
      return;
    }
    double d = std::hypot(b.x - a.x, b.y - a.y);
    if (EQ(d, 0)) {
      throw std::runtime_error("Identical points, infinite circles.");
    }
    if (LT(this->r * 2.0, d)) {
      throw std::runtime_error("Points too far away to make Circle.");
    }
    double z = this->r * this->r - d * d / 4.0;
    double v = std::sqrt(z < 0 ? 0 : z) / d;
    double mx = (a.x + b.x) / 2.0, my = (a.y + b.y) / 2.0;
    h = mx + v * (a.y - b.y);
    k = my + v * (b.x - a.x);
    // The other answer is (h, k) = (mx - v*(a.y - b.y), my - v*(b.x - a.x)).
  }

  bool operator==(const Circle &c) const { return h == c.h && k == c.k && r == c.r; }
  bool operator!=(const Circle &c) const { return !(*this == c); }

  friend bool EQ(const Circle &a, const Circle &b) {
    return EQ(a.h, b.h) && EQ(a.k, b.k) && EQ(a.r, b.r);
  }

  template<typename Pt>
  bool contains(const Pt &p) const {
    return LE(std::hypot(p.x - h, p.y - k), r);
  }

  template<typename Pt>
  bool on_boundary(const Pt &p) const {
    return EQ(std::hypot(p.x - h, p.y - k), r);
  }

  friend std::ostream &operator<<(std::ostream &out, const Circle &c) {
    auto flags = out.flags();
    out << std::showpos << "(x" << -(fabs(c.h) < EPS ? 0 : c.h) << ")^2+"
        << "(y" << -(fabs(c.k) < EPS ? 0 : c.k) << ")^2" << std::noshowpos << "="
        << (fabs(c.r) < EPS ? 0 : c.r * c.r);
    out.flags(flags);
    return out;
  }
};

template<typename Pt>
Circle incircle(const Pt &a, const Pt &b, const Pt &c) {
  double al = std::hypot(b.x - c.x, b.y - c.y);
  double bl = std::hypot(a.x - c.x, a.y - c.y);
  double cl = std::hypot(a.x - b.x, a.y - b.y);
  double l = al + bl + cl;
  double px = a.x - c.x, py = a.y - c.y, qx = b.x - c.x, qy = b.y - c.y;
  return EQ(l, 0) ? Circle(a.x, a.y, 0)
                  : Circle(
                        (al * a.x + bl * b.x + cl * c.x) / l, (al * a.y + bl * b.y + cl * c.y) / l,
                        fabs(px * qy - py * qx) / l
                    );
}

template<typename Pt>
int in_circumcircle(const Pt &a, const Pt &b, const Pt &c, const Pt &d) {
  using T = decltype(a.x + a.x);
  using W = std::conditional_t<
      std::is_floating_point<T>::value, long double,
      std::conditional_t<std::is_integral<T>::value, int64_t, T>>;
  W adx = (W)a.x - d.x, ady = (W)a.y - d.y;
  W bdx = (W)b.x - d.x, bdy = (W)b.y - d.y;
  W cdx = (W)c.x - d.x, cdy = (W)c.y - d.y;
  // Overflow warning.
  W det = (adx * adx + ady * ady) * (bdx * cdy - cdx * bdy) +
          (bdx * bdx + bdy * bdy) * (cdx * ady - adx * cdy) +
          (cdx * cdx + cdy * cdy) * (adx * bdy - bdx * ady);
  W orient = ((W)b.x - a.x) * ((W)c.y - a.y) - ((W)b.y - a.y) * ((W)c.x - a.x);
  W val = W{0} < orient ? det : -det;
  return EQ(val, W{0}) ? 0 : (W{0} < val ? 1 : -1);
}
\end{lstlisting}
\Needspace*{4\baselineskip}
\textbf{Example Usage}\par
\begin{lstlisting}[style=cpp]
#include <cassert>

struct Point {
  double x, y;
  Point(double x = 0, double y = 0) : x(x), y(y) {}
};

struct PointI {
  int x, y;
  PointI(int x = 0, int y = 0) : x(x), y(y) {}
};

int main() {
  Circle c(-2, 5, sqrt(10));
  assert(EQ(c, Circle(Point(-2, 5), sqrt(10))));
  assert(EQ(c, Circle(Point(1, 6), Point(-5, 4))));
  assert(EQ(c, Circle(Point(-3, 2), Point(-3, 8), Point(-1, 8))));
  assert(EQ(Circle(Point(0, 0), Point(0, 2), -1), Circle(0, 1, 1)));
  assert(EQ(c, incircle(Point(-12, 5), Point(3, 0), Point(0, 9))));
  assert(c.contains(Point(-2, 8)) && !c.contains(Point(-2, 9)));
  assert(c.on_boundary(Point(-1, 2)) && !c.on_boundary(Point(-1.01, 2)));

  // Integer-coordinate points are accepted; the Circle is computed in double.
  assert(EQ(c, Circle(PointI(1, 6), PointI(-5, 4))));
  assert(EQ(c, Circle(PointI(-3, 2), PointI(-3, 8), PointI(-1, 8))));
  assert(EQ(c, incircle(PointI(-12, 5), PointI(3, 0), PointI(0, 9))));
  assert(c.contains(PointI(-2, 8)) && !c.contains(PointI(-2, 9)));
  assert(c.on_boundary(PointI(-1, 2)));

  // Exact integer in-circle predicate: unit circle through (1,0), (0,1), (-1,0).
  assert(in_circumcircle(PointI(1, 0), PointI(0, 1), PointI(-1, 0), PointI(0, 0)) == 1);   // inside
  assert(in_circumcircle(PointI(1, 0), PointI(0, 1), PointI(-1, 0), PointI(2, 0)) == -1);  // out
  assert(in_circumcircle(PointI(1, 0), PointI(0, 1), PointI(-1, 0), PointI(0, -1)) == 0);  // edge

  // Orientation-independent: reversing a, b, c gives the same answer.
  assert(in_circumcircle(PointI(-1, 0), PointI(0, 1), PointI(1, 0), PointI(0, 0)) == 1);
  return 0;
}
\end{lstlisting}
\subsection{Triangle}
\label{sec:7.1.4}
Common triangle calculations in two dimensions. The functions are templated on the point type \inlinecode{Pt}, which should accept \inlinecode{Point}, \inlinecode{PointD}, or \inlinecode{PointI} from \secref{7.1.1} or any struct with numeric \inlinecode{.x} and \inlinecode{.y} fields. \inlinecode{triangle\_area()} returns \inlinecode{double} regardless of input type (the area is fractional via the \inlinecode{/2}). \inlinecode{same\_side()} and \inlinecode{point\_in\_triangle()} do all arithmetic in the point's own coordinate type and are exact for integer points: they reduce each cross product to a sign before combining, so no precision is lost and cross products are never multiplied together.

\begin{compactitem}
  \item \inlinecode{triangle\_area(a, b, c)} returns the area of the triangle with vertices \inlinecode{a}, \inlinecode{b}, and \inlinecode{c}.
  \item \inlinecode{triangle\_area\_sides(s1, s2, s3)} returns the area of a triangle with side lengths \inlinecode{s1}, \inlinecode{s2}, and \inlinecode{s3}. The given lengths must be nonnegative and form a valid triangle.
  \item \inlinecode{triangle\_area\_medians(m1, m2, m3)} returns the area of a triangle with medians of lengths \inlinecode{m1}, \inlinecode{m2}, and \inlinecode{m3}. The median of a triangle is a line segment joining a vertex to the midpoint of the opposing edge.
  \item \inlinecode{triangle\_area\_altitudes(h1, h2, h3)} returns the area of a triangle with altitudes \inlinecode{h1}, \inlinecode{h2}, and \inlinecode{h3}. An altitude of a triangle is the shortest line between a vertex and the infinite line that is extended from its opposite edge.
  \item \inlinecode{same\_side(p1, p2, a, b, include\_boundary = true)} returns whether points \inlinecode{p1} and \inlinecode{p2} lie on the same side of the line containing points \inlinecode{a} and \inlinecode{b}. If one or both points lie exactly on the line, then the result will depend on \inlinecode{include\_boundary}.
  \item \inlinecode{point\_in\_triangle(p, a, b, c, include\_boundary = true)} returns whether point \inlinecode{p} lies within the triangle with vertices \inlinecode{a}, \inlinecode{b}, and \inlinecode{c}. If the point lies on or close to an edge (by roughly \inlinecode{EPS}), then the result will depend on \inlinecode{include\_boundary}.
\end{compactitem}

Use an integral point type for integer coordinates, and a floating-point type for fractional coordinates.

Overflow warning: each individual cross product is still on the order of the squared coordinate magnitude, so for integer point types use a 64-bit coordinate type (e.g. \inlinecode{PointL} from \secref{7.1.1}) once coordinates reach a few tens of thousands.

\textbf{Time Complexity}: $O(1)$ per operation.

\textbf{Space Complexity}: $O(1)$ auxiliary for all operations.

\begin{lstlisting}[style=cpp]
#include <cmath>
#include <type_traits>

const double EPS = 1e-9;

template<typename T, typename U, typename C = std::common_type_t<T, U>>
bool EQ(T a, U b) {
  if constexpr (std::is_floating_point_v<C>) return C(a) == C(b) || std::fabs(C(a) - C(b)) <= EPS;
  return C(a) == C(b);
}

template<typename T, typename U, typename C = std::common_type_t<T, U>>
bool LT(T a, U b) {
  if constexpr (std::is_floating_point_v<C>) return C(a) < C(b) - EPS;
  return C(a) < C(b);
}

template<typename Pt>
double triangle_area(const Pt &a, const Pt &b, const Pt &c) {
  double acx = static_cast<double>(a.x) - c.x, acy = static_cast<double>(a.y) - c.y;
  double bcx = static_cast<double>(b.x) - c.x, bcy = static_cast<double>(b.y) - c.y;
  return fabs(acx * bcy - acy * bcx) / 2.0;
}

double triangle_area_sides(double s1, double s2, double s3) {
  double s = (s1 + s2 + s3) / 2.0;
  return std::sqrt(s * (s - s1) * (s - s2) * (s - s3));
}

double triangle_area_medians(double m1, double m2, double m3) {
  return 4.0 * triangle_area_sides(m1, m2, m3) / 3.0;
}

double triangle_area_altitudes(double h1, double h2, double h3) {
  if (EQ(h1, 0) || EQ(h2, 0) || EQ(h3, 0)) {
    return 0;
  }
  double x = h1 * h1, y = h2 * h2, z = h3 * h3;
  double v = 2.0 / (x * y) + 2.0 / (x * z) + 2.0 / (y * z);
  return 1.0 / std::sqrt(v - 1.0 / (x * x) - 1.0 / (y * y) - 1.0 / (z * z));
}

template<typename Pt>
bool same_side(const Pt &p1, const Pt &p2, const Pt &a, const Pt &b, bool include_boundary = true) {
  // Cross products in the point's native coordinate type (exact for integer points).
  auto c1 = (b.x - a.x) * (p1.y - a.y) - (b.y - a.y) * (p1.x - a.x);
  auto c2 = (b.x - a.x) * (p2.y - a.y) - (b.y - a.y) * (p2.x - a.x);
  // Compare the signs, not the product c1 * c2, to avoid integer overflow.
  int s1 = LT(0, c1) ? 1 : (LT(c1, 0) ? -1 : 0);
  int s2 = LT(0, c2) ? 1 : (LT(c2, 0) ? -1 : 0);
  return include_boundary ? (s1 * s2 >= 0) : (s1 * s2 > 0);
}

template<typename Pt>
bool point_in_triangle(
    const Pt &p, const Pt &a, const Pt &b, const Pt &c, bool include_boundary = true
) {
  return same_side(p, a, b, c, include_boundary) && same_side(p, b, a, c, include_boundary) &&
         same_side(p, c, a, b, include_boundary);
}
\end{lstlisting}
\Needspace*{4\baselineskip}
\textbf{Example Usage}\par
\begin{lstlisting}[style=cpp]
#include <cassert>

struct Point {
  double x, y;
  Point(double x = 0, double y = 0) : x(x), y(y) {}
};

struct PointI {
  int x, y;
  PointI(int x = 0, int y = 0) : x(x), y(y) {}
};

int main() {
  assert(EQ(6, triangle_area(Point(0, -1), Point(4, -1), Point(0, -4))));
  assert(EQ(6, triangle_area_sides(3, 4, 5)));
  assert(EQ(6, triangle_area_medians(3.605551275, 2.5, 4.272001873)));
  assert(EQ(6, triangle_area_altitudes(3, 4, 2.4)));

  // Fractional coordinates require a floating-point point type.
  assert(point_in_triangle(Point(0, 0), Point(-1, 0), Point(0, -2), Point(4, 0)));
  assert(!point_in_triangle(Point(0, 1), Point(-1, 0), Point(0, -2), Point(4, 0)));
  assert(point_in_triangle(Point(-2.44, 0.82), Point(-1, 0), Point(-3, 1), Point(4, 0)));
  assert(!point_in_triangle(Point(-2.44, 0.7), Point(-1, 0), Point(-3, 1), Point(4, 0)));

  // Integer coordinates: same_side / point_in_triangle are computed exactly in int.
  assert(EQ(8, triangle_area(PointI(0, 0), PointI(4, 0), PointI(0, 4))));
  assert(point_in_triangle(PointI(1, 1), PointI(0, 0), PointI(4, 0), PointI(0, 4)));
  assert(!point_in_triangle(PointI(5, 5), PointI(0, 0), PointI(4, 0), PointI(0, 4)));
  assert(point_in_triangle(PointI(2, 2), PointI(0, 0), PointI(4, 0), PointI(0, 4)));  // on edge
  assert(!point_in_triangle(PointI(2, 2), PointI(0, 0), PointI(4, 0), PointI(0, 4), false));
  return 0;
}
\end{lstlisting}
\subsection{Rectangle}
\label{sec:7.1.5}
Common axis-aligned rectangle calculations in two dimensions. The functions are templated on the point type \inlinecode{Pt}, which should accept \inlinecode{Point}, \inlinecode{PointD}, or \inlinecode{PointI} from \secref{7.1.1}, or any struct with numeric \inlinecode{.x} and \inlinecode{.y} fields.

\begin{compactitem}
  \item \inlinecode{point\_in\_rect(p, v, w, h, include\_boundary = true)} returns whether \inlinecode{p} is inside the axis-aligned rectangle whose opposite corners are \inlinecode{v} and (\inlinecode{v.x + w}, \inlinecode{v.y + h}). Negative widths and heights are supported. The \inlinecode{include\_boundary} flag controls whether boundary points count.
  \item \inlinecode{point\_in\_rect(p, a, b, include\_boundary = true)} returns whether \inlinecode{p} is inside the axis-aligned rectangle with opposite corners \inlinecode{a} and \inlinecode{b}.
  \item \inlinecode{rect\_intersection(a1, b1, a2, b2, \&p, \&q, include\_boundary = true)} finds the intersection of two axis-aligned rectangles, where \inlinecode{a1}/\inlinecode{b1} and \inlinecode{a2}/\inlinecode{b2} are opposite-corner pairs. Returns $-1$ if the rectangles are disjoint, $0$ if they partially intersect, $1$ if the first rectangle is completely inside the second, and $2$ if the second rectangle is completely inside the first. If an intersection exists, its lower-left and upper-right corners are stored in \inlinecode{p} and \inlinecode{q} when those pointers are non-null. The \inlinecode{include\_boundary} flag controls whether boundary-only contact counts as intersecting.
\end{compactitem}

Overflow warning: For integer-coordinate inputs, comparisons are exact as long as intermediate coordinate additions, subtractions, and min/max values do not overflow.

\textbf{Time Complexity}: $O(1)$ per operation.

\textbf{Space Complexity}: $O(1)$ auxiliary for all operations.

\begin{lstlisting}[style=cpp]
#include <algorithm>
#include <cmath>
#include <type_traits>

const double EPS = 1e-9;

template<typename T, typename U, typename C = std::common_type_t<T, U>>
bool EQ(T a, U b) {
  if constexpr (std::is_floating_point_v<C>) return C(a) == C(b) || std::fabs(C(a) - C(b)) <= EPS;
  return C(a) == C(b);
}

template<typename T, typename U, typename C = std::common_type_t<T, U>>
bool LT(T a, U b) {
  if constexpr (std::is_floating_point_v<C>) return C(a) < C(b) - EPS;
  return C(a) < C(b);
}

template<typename T, typename U> bool GT(T a, U b) { return LT(b, a); }
template<typename T, typename U> bool LE(T a, U b) { return !LT(b, a); }
template<typename T, typename U> bool GE(T a, U b) { return !LT(a, b); }

template<typename Pt, typename T>
bool point_in_rect(const Pt &p, const Pt &v, const T &w, const T &h, bool include_boundary = true) {
  if (w < 0) {
    return point_in_rect(p, Pt(v.x + w, v.y), -w, h, include_boundary);
  }
  if (h < 0) {
    return point_in_rect(p, Pt(v.x, v.y + h), w, -h, include_boundary);
  }
  return include_boundary ? (GE(p.x, v.x) && LE(p.x, v.x + w) && GE(p.y, v.y) && LE(p.y, v.y + h))
                          : (GT(p.x, v.x) && LT(p.x, v.x + w) && GT(p.y, v.y) && LT(p.y, v.y + h));
}

template<typename Pt>
bool point_in_rect(const Pt &p, const Pt &a, const Pt &b, bool include_boundary = true) {
  auto xl = std::min(a.x, b.x), yl = std::min(a.y, b.y);
  auto xh = std::max(a.x, b.x), yh = std::max(a.y, b.y);
  return point_in_rect(p, Pt(xl, yl), xh - xl, yh - yl, include_boundary);
}

template<typename Pt>
int rect_intersection(
    const Pt &a1, const Pt &b1, const Pt &a2, const Pt &b2, Pt *p = nullptr, Pt *q = nullptr,
    const bool include_boundary = true
) {
  // Normalize each rectangle to [xl, xh] x [yl, yh] with coordinate-wise min/max only.
  // (Using auto keeps the native coordinate type, so integer rectangles stay exact.)
  auto xl1 = std::min(a1.x, b1.x), xh1 = std::max(a1.x, b1.x);
  auto yl1 = std::min(a1.y, b1.y), yh1 = std::max(a1.y, b1.y);
  auto xl2 = std::min(a2.x, b2.x), xh2 = std::max(a2.x, b2.x);
  auto yl2 = std::min(a2.y, b2.y), yh2 = std::max(a2.y, b2.y);
  // The intersection is the overlap of the two coordinate intervals.
  auto ixl = std::max(xl1, xl2), ixh = std::min(xh1, xh2);
  auto iyl = std::max(yl1, yl2), iyh = std::min(yh1, yh2);
  if (include_boundary ? (GT(ixl, ixh) || GT(iyl, iyh)) : (GE(ixl, ixh) || GE(iyl, iyh))) {
    return -1;  // Completely disjoint.
  }
  // Opposing corners of the overlap: bottom-left in p, top-right in q.
  if (p != nullptr) {
    *p = Pt(ixl, iyl);
  }
  if (q != nullptr) {
    *q = Pt(ixh, iyh);
  }
  // The overlap equals a rectangle exactly when that rectangle is contained in the other.
  if (EQ(ixl, xl1) && EQ(ixh, xh1) && EQ(iyl, yl1) && EQ(iyh, yh1)) {
    return 1;  // Rectangle 1 completely inside 2.
  }
  if (EQ(ixl, xl2) && EQ(ixh, xh2) && EQ(iyl, yl2) && EQ(iyh, yh2)) {
    return 2;  // Rectangle 2 completely inside 1.
  }
  return 0;  // Partial intersection.
}
\end{lstlisting}
\Needspace*{4\baselineskip}
\textbf{Example Usage}\par
\begin{lstlisting}[style=cpp]
#include <cassert>

struct Point {
  double x, y;
  Point(double x = 0, double y = 0) : x(x), y(y) {}
  bool operator==(const Point &p) const { return x == p.x && y == p.y; }
};

int main() {
  assert(point_in_rect(Point(0, -1), Point(0, -3), 3, 2));
  assert(!point_in_rect(Point(0, -1), Point(0, -3), 3, 2, false));
  assert(point_in_rect(Point(2, -2), Point(3, -3), -3, 2));
  assert(!point_in_rect(Point(0, 0), Point(3, -1), -3, -2));
  assert(point_in_rect(Point(2, -2), Point(3, -3), Point(0, -1)));
  assert(!point_in_rect(Point(0, -1), Point(3, -3), Point(0, -1), false));
  assert(!point_in_rect(Point(-1, -2), Point(3, -3), Point(0, -1)));

  Point p, q;
  // Output coordinates come directly from input coordinates via min/max, so exact equality is safe.
  assert(-1 == rect_intersection(Point(0, 0), Point(1, 1), Point(2, 2), Point(3, 3)));
  assert(0 == rect_intersection(Point(1, 1), Point(7, 7), Point(5, 5), Point(0, 0), &p, &q));
  assert(p == Point(1, 1) && q == Point(5, 5));
  assert(1 == rect_intersection(Point(1, 1), Point(0, 0), Point(0, 0), Point(1, 10), &p, &q));
  assert(p == Point(0, 0) && q == Point(1, 1));
  assert(2 == rect_intersection(Point(0, 5), Point(5, 7), Point(1, 6), Point(2, 5), &p, &q));
  assert(p == Point(1, 5) && q == Point(2, 6));
  assert(
      -1 == rect_intersection(Point(0, 0), Point(1, 1), Point(1, 1), Point(2, 2), &p, &q, false)
  );
  return 0;
}
\end{lstlisting}

\section{\texorpdfstring{Angles, Distances, and Intersections}{7.2 Angles, Distances, and Intersections}}
\setcounter{section}{2}
\setcounter{subsection}{0}
\subsection{Angles}
\label{sec:7.2.1}
Angle calculations in two dimensions. All operations are inherently floating-point (\inlinecode{atan2}, \inlinecode{acos}, trigonometry), so these functions use the local double-coordinate \inlinecode{PointD} type. Convert integer points to \inlinecode{PointD} when asking for angles. The constants \inlinecode{DEG} and \inlinecode{RAD} may be used as multipliers to convert between degrees and radians. For example, if \inlinecode{t} is a value in degrees, then \inlinecode{t * DEG} is the equivalent angle in radians; if \inlinecode{t} is in radians, then \inlinecode{t * RAD} is the equivalent angle in degrees.

\begin{compactitem}
  \item \inlinecode{reduce\_deg(t)} takes an angle \inlinecode{t} degrees and returns an equivalent angle in the range $[0, 360)$ degrees (e.g. $-630$ becomes $90$).
  \item \inlinecode{reduce\_rad(t)} takes an angle \inlinecode{t} radians and returns an equivalent angle in the range $[0, 2\pi)$ radians (e.g. $720.5$ becomes $0.5$).
  \item \inlinecode{polar\_x(r, t)} and \inlinecode{polar\_y(r, t)} return the Cartesian coordinates of the point at radius \inlinecode{r} and angle \inlinecode{t} radians in polar coordinates (see \inlinecode{std::polar()}).
  \item \inlinecode{polar\_angle(p)} returns the angle in radians of the line segment from $(0, 0)$ to point \inlinecode{p}, relative counterclockwise to the positive $x$-axis.
  \item \inlinecode{angle(a, o, b)} returns the smallest angle in radians formed by the points \inlinecode{a}, \inlinecode{o}, \inlinecode{b} with vertex at point \inlinecode{o}. The points \inlinecode{a} and \inlinecode{b} must differ from \inlinecode{o}.
  \item \inlinecode{angle\_between(a, b)} returns the directed counter-clockwise angle in $[0, 2\pi)$ from vector \inlinecode{a} to vector \inlinecode{b}, treating both points as vectors from the origin. Both vectors must be nonzero.
  \item \inlinecode{angle\_between(a1, b1, a2, b2)} returns the smaller angle in radians between two lines whose normal vectors are (\inlinecode{a1}, \inlinecode{b1}) and (\inlinecode{a2}, \inlinecode{b2}), limited to $[0, \pi / 2]$. Both coefficient pairs must represent valid lines.
  \item \inlinecode{cross(a, b, o = Pt(0, 0))} returns the signed $z$-component of the three-dimensional cross product between points \inlinecode{a} and \inlinecode{b}, where the $z$-coordinates are implicitly zero and the origin is shifted to point \inlinecode{o}. This is also double the signed area of the triangle from these three points.
  \item \inlinecode{turn(a, o, b)} returns $1$ if the path \inlinecode{a} $\to$ \inlinecode{o} $\to$ \inlinecode{b} forms a left turn on the plane, $0$ if the path forms a straight line segment, or $-1$ if it forms a right turn.
\end{compactitem}

Overflow warning: \inlinecode{cross()} and \inlinecode{turn()} preserve the point's coordinate type, so their products may overflow for integral points. Use a 64-bit coordinate type for large integer coordinates.

\textbf{Time Complexity}: $O(1)$ per operation.

\textbf{Space Complexity}: $O(1)$ auxiliary for all operations.

\begin{lstlisting}[style=cpp]
#include <algorithm>
#include <cmath>
#include <type_traits>

const double EPS = 1e-9;
const double PI = std::acos(-1.0);
const double DEG = PI / 180;
const double RAD = 180 / PI;

template<typename T, typename U, typename C = std::common_type_t<T, U>>
bool EQ(T a, U b) {
  if constexpr (std::is_floating_point_v<C>) return C(a) == C(b) || std::fabs(C(a) - C(b)) <= EPS;
  return C(a) == C(b);
}

template<typename T, typename U, typename C = std::common_type_t<T, U>>
bool LT(T a, U b) {
  if constexpr (std::is_floating_point_v<C>) return C(a) < C(b) - EPS;
  return C(a) < C(b);
}

double reduce_deg(double t) {
  if (t < -360) {
    return reduce_deg(std::fmod(t, 360));
  }
  if (t < 0) {
    return t + 360;
  }
  return (t >= 360) ? std::fmod(t, 360) : t;
}

double reduce_rad(double t) {
  if (t < -2 * PI) {
    return reduce_rad(std::fmod(t, 2 * PI));
  }
  if (t < 0) {
    return t + 2 * PI;
  }
  return (t >= 2 * PI) ? std::fmod(t, 2 * PI) : t;
}

double polar_x(double r, double t) {
  return r * std::cos(t);
}

double polar_y(double r, double t) {
  return r * std::sin(t);
}

template<typename Pt>
double polar_angle(const Pt &p) {
  double t = std::atan2(static_cast<double>(p.y), static_cast<double>(p.x));
  return (t < 0) ? (t + 2 * PI) : t;
}

template<typename Pt>
double angle(const Pt &a, const Pt &o, const Pt &b) {
  double ux = o.x - a.x, uy = o.y - a.y;
  double vx = o.x - b.x, vy = o.y - b.y;
  double cosine = (ux * vx + uy * vy) / (std::hypot(ux, uy) * std::hypot(vx, vy));
  return std::acos(std::clamp(cosine, -1.0, 1.0));
}

template<typename Pt>
double angle_between(const Pt &a, const Pt &b) {
  double t = std::atan2(a.x * b.y - a.y * b.x, a.x * b.x + a.y * b.y);
  return (t < 0) ? (t + 2 * PI) : t;
}

double angle_between(double a1, double b1, double a2, double b2) {
  double t = std::atan2(a1 * b2 - a2 * b1, a1 * a2 + b1 * b2);
  if (t < 0) {
    t += PI;
  }
  return LT(PI / 2, t) ? (PI - t) : t;
}

template<typename Pt>
auto cross(const Pt &a, const Pt &b, const Pt &o = Pt(0, 0)) {
  return (a.x - o.x) * (b.y - o.y) - (a.y - o.y) * (b.x - o.x);
}

template<typename Pt>
int turn(const Pt &a, const Pt &o, const Pt &b) {
  auto c = cross(a, b, o);
  return LT(c, 0) ? 1 : (LT(0, c) ? -1 : 0);
}
\end{lstlisting}
\Needspace*{4\baselineskip}
\textbf{Example Usage}\par
\begin{lstlisting}[style=cpp]
struct Point {
  double x, y;
  Point(double x = 0, double y = 0) : x(x), y(y) {}
  bool operator==(const Point &p) const { return x == p.x && y == p.y; }
  friend bool EQ(const Point &a, const Point &b) { return EQ(a.x, b.x) && EQ(a.y, b.y); }
};

#include <cassert>

int main() {
  assert(EQ(123, reduce_deg(-8 * 360 + 123)));
  assert(EQ(1.2345, reduce_rad(2 * PI * 8 + 1.2345)));
  assert(EQ(-4, polar_x(4, PI)) && EQ(0, polar_y(4, PI)));
  assert(EQ(0, polar_x(4, -PI / 2)) && EQ(-4, polar_y(4, -PI / 2)));
  assert(EQ(45, polar_angle(Point(5, 5)) * RAD));
  assert(EQ(135 * DEG, polar_angle(Point(-4, 4))));
  assert(EQ(90 * DEG, angle(Point(5, 0), Point(0, 5), Point(-5, 0))));
  assert(EQ(225 * DEG, angle_between(Point(0, 5), Point(5, -5))));
  assert(-1 == cross(Point(0, 1), Point(1, 0), Point(0, 0)));
  assert(-1 == turn(Point(0, 1), Point(0, 0), Point(-5, -5)));
  return 0;
}
\end{lstlisting}
\subsection{Distances}
\label{sec:7.2.2}
Distance calculations in two dimensions for points, lines, and line segments. The functions are templated on the point type \inlinecode{Pt}: pass any struct with numeric \inlinecode{.x} and \inlinecode{.y} fields (for example \inlinecode{PointD}/\inlinecode{PointI} from \secref{7.1.1}, or the local example struct). \inlinecode{sqdist} preserves the coordinate type and is exact for integer points. Functions that return an actual distance use \inlinecode{double} for the final division/square root. \inlinecode{closest\_point} returns a point of the same type as its inputs, so use a floating-point \inlinecode{Pt} there for a meaningful (non-truncated) result.

\begin{compactitem}
  \item \inlinecode{dist(a, b)} and \inlinecode{sqdist(a, b)} respectively return the distance and squared distance between points \inlinecode{a} and \inlinecode{b}.
  \item \inlinecode{line\_dist(p, a, b, c)} returns the distance from point \inlinecode{p} to the valid line $ax + by + c = 0$.
  \item \inlinecode{line\_dist(p, a, b)} returns the distance from point \inlinecode{p} to the infinite line containing points \inlinecode{a} and \inlinecode{b}. If \inlinecode{a} = \inlinecode{b}, the distance from \inlinecode{p} to \inlinecode{a} is returned.
  \item \inlinecode{line\_dist(a1, b1, c1, a2, b2, c2)} returns the distance between two valid lines.
  \item \inlinecode{seg\_dist(p, a, b)} returns the distance from point \inlinecode{p} to the line segment \inlinecode{a}-\inlinecode{b}.
  \item \inlinecode{seg\_dist(a, b, c, d)} returns the minimum distance between line segments \inlinecode{a}-\inlinecode{b} and \inlinecode{c}-\inlinecode{d}, or $0$ if the segments touch or intersect.
  \item \inlinecode{closest\_point(a, b, p)} returns the point on segment \inlinecode{a}-\inlinecode{b} closest to point \inlinecode{p}, using the input point type for the result. Use a floating-point point type when the closest point may have fractional coordinates.
\end{compactitem}

Overflow warning: squared distances, dot products, and cross products are formed in the point's coordinate arithmetic type. For integer point types, use a 64-bit coordinate type when coordinates may exceed a few tens of thousands.

\textbf{Time Complexity}: $O(1)$ per operation.

\textbf{Space Complexity}: $O(1)$ auxiliary for all operations.

\begin{lstlisting}[style=cpp]
#include <algorithm>
#include <cmath>
#include <type_traits>

const double EPS = 1e-9;

template<typename T, typename U, typename C = std::common_type_t<T, U>>
bool EQ(T a, U b) {
  if constexpr (std::is_floating_point_v<C>) return C(a) == C(b) || std::fabs(C(a) - C(b)) <= EPS;
  return C(a) == C(b);
}

template<typename T, typename U, typename C = std::common_type_t<T, U>>
bool LT(T a, U b) {
  if constexpr (std::is_floating_point_v<C>) return C(a) < C(b) - EPS;
  return C(a) < C(b);
}

template<typename T, typename U> bool LE(T a, U b) { return !LT(b, a); }
template<typename T, typename U> bool GE(T a, U b) { return !LT(a, b); }

// sqdist returns the coordinate type (exact for integer points).
template<typename Pt>
auto sqdist(const Pt &a, const Pt &b) {
  auto dx = b.x - a.x, dy = b.y - a.y;
  return dx * dx + dy * dy;  // Overflow warning.
}

template<typename Pt>
double dist(const Pt &a, const Pt &b) {
  return std::sqrt(static_cast<double>(sqdist(a, b)));
}

template<typename Pt, typename T>
double line_dist(const Pt &p, const T &a, const T &b, const T &c) {
  return fabs(static_cast<double>(a) * p.x + static_cast<double>(b) * p.y + c) /
         std::sqrt(static_cast<double>(a) * a + static_cast<double>(b) * b);
}

template<typename Pt>
double line_dist(const Pt &p, const Pt &a, const Pt &b) {
  if (EQ(a.x, b.x) && EQ(a.y, b.y)) {
    return dist(p, a);
  }
  auto n = sqdist(a, b);
  auto d = (p.x - a.x) * (b.x - a.x) + (p.y - a.y) * (b.y - a.y);  // Overflow warning.
  double u = static_cast<double>(d) / n;
  double dx = a.x + u * (b.x - a.x) - p.x, dy = a.y + u * (b.y - a.y) - p.y;
  return std::sqrt(dx * dx + dy * dy);
}

template<typename T>
double line_dist(const T &a1, const T &b1, const T &c1, const T &a2, const T &b2, const T &c2) {
  if (EQ(static_cast<double>(a1) * b2, static_cast<double>(a2) * b1)) {
    double factor = EQ(b1, 0) ? (static_cast<double>(a1) / a2) : (static_cast<double>(b1) / b2);
    return EQ(c1, c2 * factor)
               ? 0
               : fabs(c2 * factor - c1) /
                     std::sqrt(static_cast<double>(a1) * a1 + static_cast<double>(b1) * b1);
  }
  return 0;
}

template<typename Pt>
double seg_dist(const Pt &p, const Pt &a, const Pt &b) {
  if (EQ(a.x, b.x) && EQ(a.y, b.y)) {
    return dist(p, a);
  }
  auto n = sqdist(a, b);
  auto d = (p.x - a.x) * (b.x - a.x) + (p.y - a.y) * (b.y - a.y);  // Overflow warning.
  if (LE(d, 0) || EQ(n, 0)) {
    return dist(p, a);
  }
  if (GE(d, n)) {
    return dist(p, b);
  }
  double t = static_cast<double>(d) / n;
  double dx = a.x + t * (b.x - a.x) - p.x, dy = a.y + t * (b.y - a.y) - p.y;
  return std::sqrt(dx * dx + dy * dy);
}

template<typename Pt>
double seg_dist(const Pt &a, const Pt &b, const Pt &c, const Pt &d) {
  if (EQ(a.x, b.x) && EQ(a.y, b.y)) {
    return seg_dist(a, c, d);
  }
  if (EQ(c.x, d.x) && EQ(c.y, d.y)) {
    return seg_dist(c, a, b);
  }
  auto ab_x = b.x - a.x, ab_y = b.y - a.y;
  auto ac_x = c.x - a.x, ac_y = c.y - a.y;
  auto cd_x = d.x - c.x, cd_y = d.y - c.y;
  auto c1 = ab_x * cd_y - ab_y * cd_x;  // Overflow warning.
  auto c2 = ac_x * ab_y - ac_y * ab_x;
  if (EQ(c1, 0) && EQ(c2, 0)) {
    auto less = [](const Pt &u, const Pt &v) { return u.x != v.x ? u.x < v.x : u.y < v.y; };
    auto minp = [&](const Pt &u, const Pt &v) -> const Pt & { return less(v, u) ? v : u; };
    auto maxp = [&](const Pt &u, const Pt &v) -> const Pt & { return less(u, v) ? v : u; };
    const Pt &res1 = maxp(minp(a, b), minp(c, d));
    const Pt &res2 = minp(maxp(a, b), maxp(c, d));
    if (!less(res2, res1)) {
      return 0;
    }
  } else if (!EQ(c1, 0)) {
    auto t_num = ac_x * cd_y - ac_y * cd_x;
    bool c1_pos = c1 > 0;
    bool t_between_01 = c1_pos ? (LE(0, t_num) && LE(t_num, c1)) : (LE(t_num, 0) && LE(c1, t_num));
    bool u_between_01 = c1_pos ? (LE(0, c2) && LE(c2, c1)) : (LE(c2, 0) && LE(c1, c2));
    if (t_between_01 && u_between_01) {
      return 0;
    }
  }
  return std::min(
      std::min(seg_dist(a, c, d), seg_dist(b, c, d)), std::min(seg_dist(c, a, b), seg_dist(d, a, b))
  );
}

template<typename Pt>
Pt closest_point(const Pt &a, const Pt &b, const Pt &p) {
  Pt res{};
  if (EQ(a.x, b.x) && EQ(a.y, b.y)) {
    res.x = a.x;
    res.y = a.y;
    return res;
  }
  auto n = sqdist(a, b);
  auto d = (p.x - a.x) * (b.x - a.x) + (p.y - a.y) * (b.y - a.y);  // Overflow warning.
  double t = static_cast<double>(d) / n;
  if (t <= 0) {
    res.x = a.x;
    res.y = a.y;
  } else if (t >= 1) {
    res.x = b.x;
    res.y = b.y;
  } else {
    res.x = a.x + t * (b.x - a.x);
    res.y = a.y + t * (b.y - a.y);
  }
  return res;
}
\end{lstlisting}
\Needspace*{4\baselineskip}
\textbf{Example Usage}\par
\begin{lstlisting}[style=cpp]
#include <cassert>

struct Point {
  double x, y;
  Point(double x = 0, double y = 0) : x(x), y(y) {}
};

struct PointI {
  int x, y;
  PointI(int x = 0, int y = 0) : x(x), y(y) {}
};

int main() {
  assert(EQ(5, dist(Point(-1, -1), Point(2, 3))));
  assert(EQ(25, sqdist(Point(-1, -1), Point(2, 3))));
  assert(EQ(1.2, line_dist(Point(2, 1), -4, 3, -1)));
  assert(EQ(0.8, line_dist(Point(3, 3), Point(-1, -1), Point(2, 3))));
  assert(EQ(1.2, line_dist(Point(2, 1), Point(-1, -1), Point(2, 3))));
  assert(EQ(0, line_dist(-4, 3, -1, 8, 6, 2)));
  assert(EQ(0.8, line_dist(-4, 3, -1, -8, 6, -10)));
  assert(EQ(1.0, seg_dist(Point(3, 3), Point(-1, -1), Point(2, 3))));
  assert(EQ(1.2, seg_dist(Point(2, 1), Point(-1, -1), Point(2, 3))));
  assert(EQ(0, seg_dist(Point(0, 2), Point(3, 3), Point(-1, -1), Point(2, 3))));
  assert(EQ(0.6, seg_dist(Point(-1, 0), Point(-2, 2), Point(-1, -1), Point(2, 3))));

  // Integer points: sqdist is exact, dist returns double.
  PointI ia{-1, -1}, ib{2, 3};
  assert(sqdist(ia, ib) == 25);   // int result
  assert(EQ(dist(ia, ib), 5.0));  // double result
  return 0;
}
\end{lstlisting}
\subsection{Line Intersection}
\label{sec:7.2.3}
Intersection calculations in two dimensions for straight lines and line segments. The functions are templated on the point type \inlinecode{Pt}: pass any struct with numeric \inlinecode{.x} and \inlinecode{.y} fields (for example \inlinecode{PointD}/\inlinecode{PointI} from \secref{7.1.1}, or the local example struct) for which \inlinecode{operator\textless{}} orders points lexicographically using exact coordinate comparisons. Point output types are deduced from the caller-supplied arguments, so they may differ from the input type, e.g. integer endpoints with a floating-point output point.

\begin{compactitem}
  \item \inlinecode{line\_intersection(a1, b1, c1, a2, b2, c2, \&p)} intersects lines $a_1 x + b_1 y + c_1 = 0$ and $a_2 x + b_2 y + c_2 = 0$, returning $-1$ (parallel), $0$ (one point, stored into \inlinecode{p} if it is not \inlinecode{nullptr}), or $1$ (identical). Both coefficient triples must represent valid lines.
  \item \inlinecode{line\_intersection(p1, p2, p3, p4, \&p)} intersects the infinite lines determined by points \inlinecode{p1}, \inlinecode{p2}, \inlinecode{p3}, and \inlinecode{p4}, returning $-1$ (parallel), $0$ (one point, stored into \inlinecode{p} if it is not \inlinecode{nullptr}), or $1$ (identical). The points in each pair must differ.
  \item \inlinecode{seg\_intersection(a, b, c, d, \&p, \&q, include\_boundary = true)} intersects segments \inlinecode{a}-\inlinecode{b} and \inlinecode{c}-\inlinecode{d}, returning $-1$ (none), $0$ (one point), or $1$ (overlapping segment). The \inlinecode{include\_boundary} flag controls whether segments that meet only at an endpoint count as intersecting. If the intersection is a point (and \inlinecode{p} is not \inlinecode{nullptr}), it will be stored into \inlinecode{p}. If the intersection is an overlapping segment (and \inlinecode{p} and \inlinecode{q} are not \inlinecode{nullptr}) then their endpoints will be stored into pointers \inlinecode{p} and \inlinecode{q}, which must reference a floating-point point type.
  \item \inlinecode{seg\_intersection(a, b, c, d, include\_boundary = true)} is a simplified, detection-only version of the above. Both versions are exact for detection if the input point type is integral.
  \item \inlinecode{closest\_point(a, b, c, p)} returns the closest point on the valid line $ax + by + c = 0$ to point \inlinecode{p}.
\end{compactitem}

Overflow warning: the exact integer paths multiply coordinate differences (cross products and squared lengths), which grow like the squared coordinate magnitude. With 32-bit \inlinecode{int} coordinates these overflow once coordinates exceed roughly a few tens of thousands, so for larger integer inputs use a point type with 64-bit (\inlinecode{int64\_t}) coordinates.

\textbf{Time Complexity}: $O(1)$ per operation.

\textbf{Space Complexity}: $O(1)$ auxiliary for all operations.

\begin{lstlisting}[style=cpp]
#include <algorithm>
#include <cmath>
#include <type_traits>

const double EPS = 1e-9;

template<typename T, typename U, typename C = std::common_type_t<T, U>>
bool EQ(T a, U b) {
  if constexpr (std::is_floating_point_v<C>) return C(a) == C(b) || std::fabs(C(a) - C(b)) <= EPS;
  return C(a) == C(b);
}

template<typename T, typename U, typename C = std::common_type_t<T, U>>
bool LT(T a, U b) {
  if constexpr (std::is_floating_point_v<C>) return C(a) < C(b) - EPS;
  return C(a) < C(b);
}

template<typename T, typename U>
bool LE(T a, U b) {
  return !LT(b, a);
}

// Guard to check outputs are real coordinates, since intersections are generally not integral.
template<typename OutPt>
constexpr bool is_float_pt =
    std::is_floating_point_v<decltype(OutPt::x)> && std::is_floating_point_v<decltype(OutPt::y)>;

template<typename OutPt>
int line_intersection(double a1, double b1, double c1, double a2, double b2, double c2, OutPt *p) {
  static_assert(is_float_pt<OutPt>, "output point coordinates must be floating-point");
  double det = a1 * b2 - a2 * b1;  // Cramer's rule for a1*x + b1*y = -c1, a2*x + b2*y = -c2.
  if (EQ(det, 0)) {                // Parallel (lines need not be given in normalized form).
    // Identical iff the other two 2x2 determinants also vanish.
    return (EQ(a1 * c2 - a2 * c1, 0) && EQ(b1 * c2 - b2 * c1, 0)) ? 1 : -1;
  }
  if (p != nullptr) {
    double x = (b1 * c2 - b2 * c1) / det;  // Local x to ensure y is derived from full-precision x.
    p->x = x;
    p->y = !EQ(b1, 0) ? -(a1 * x + c1) / b1 : -(a2 * x + c2) / b2;
  }
  return 0;
}

template<typename Pt, typename OutPt>
int line_intersection(const Pt &p1, const Pt &p2, const Pt &p3, const Pt &p4, OutPt *p) {
  static_assert(is_float_pt<OutPt>, "output point coordinates must be floating-point");
  auto a1 = p2.y - p1.y, b1 = p1.x - p2.x;
  auto c1 = -(p1.x * p2.y - p2.x * p1.y);
  auto a2 = p4.y - p3.y, b2 = p3.x - p4.x;
  auto c2 = -(p3.x * p4.y - p4.x * p3.y);
  auto x = -(c1 * b2 - c2 * b1), y = -(a1 * c2 - a2 * c1);
  auto det = a1 * b2 - a2 * b1;
  if (EQ(det, 0)) {
    return (EQ(x, 0) && EQ(y, 0)) ? 1 : -1;
  }
  if (p != nullptr) {
    p->x = static_cast<double>(x) / det;
    p->y = static_cast<double>(y) / det;
  }
  return 0;
}

template<typename Pt>
bool point_on_segment(const Pt &p, const Pt &a, const Pt &b) {
  // Overflow risk for integer Pt: these products are ~O(max_coord^2); use int64_t if necessary.
  return EQ((p.x - a.x) * (b.y - a.y) - (p.y - a.y) * (b.x - a.x), 0) &&
         LE(std::min(a.x, b.x), p.x) && LE(p.x, std::max(a.x, b.x)) &&
         LE(std::min(a.y, b.y), p.y) && LE(p.y, std::max(a.y, b.y));
}

// Detection is exact for integer Pt. Intersection point output may use a separate float OutPt.
template<typename Pt, typename OutPt>
int seg_intersection(
    const Pt &a, const Pt &b, const Pt &c, const Pt &d, OutPt *p = nullptr, OutPt *q = nullptr,
    const bool include_boundary = true
) {
  static_assert(is_float_pt<OutPt>, "output point coordinates must be floating-point");
  auto set_out = [](auto *out, auto x, auto y) {
    if (out != nullptr) {
      out->x = x;
      out->y = y;
    }
  };
  auto ab_x = b.x - a.x, ab_y = b.y - a.y;
  auto ac_x = c.x - a.x, ac_y = c.y - a.y;
  auto cd_x = d.x - c.x, cd_y = d.y - c.y;
  auto ab2 = ab_x * ab_x + ab_y * ab_y;
  auto cd2 = cd_x * cd_x + cd_y * cd_y;
  if (EQ(ab2, 0)) {
    if (include_boundary && point_on_segment(a, c, d)) {
      set_out(p, a.x, a.y);
      return 0;  // Degenerate: first segment is a point intersecting the second segment.
    }
    return -1;  // Degenerate: first segment is a point not intersecting the second segment.
  }
  if (EQ(cd2, 0)) {
    if (include_boundary && point_on_segment(c, a, b)) {
      set_out(p, c.x, c.y);
      return 0;  // Degenerate: second segment is a point intersecting the first segment.
    }
    return -1;  // Degenerate: second segment is a point not intersecting the first segment.
  }
  auto c1 = ab_x * cd_y - ab_y * cd_x;
  auto c2 = ac_x * ab_y - ac_y * ab_x;
  if (EQ(c1, 0) && EQ(c2, 0)) {  // Collinear.
    Pt res1 = std::max(std::min(a, b), std::min(c, d));
    Pt res2 = std::min(std::max(a, b), std::max(c, d));
    if (include_boundary ? !(res2 < res1) : (res1 < res2)) {
      if (res1 == res2) {
        set_out(p, res1.x, res1.y);
        return 0;  // Collinear and overlapping: touch at one endpoint.
      }
      if (p != nullptr && q != nullptr) {  // Both ends are reported, or neither is.
        set_out(p, res1.x, res1.y);
        set_out(q, res2.x, res2.y);
      }
      return 1;  // Collinear and overlapping: overlap is a segment.
    }
    return -1;  // Collinear and disjoint.
  }
  if (EQ(c1, 0)) {
    return -1;  // Parallel and disjoint.
  }
  auto t_num = ac_x * cd_y - ac_y * cd_x;
  bool c1_pos = c1 > 0;
  bool t_ok =
      c1_pos
          ? (include_boundary ? (LE(0, t_num) && LE(t_num, c1)) : (LT(0, t_num) && LT(t_num, c1)))
          : (include_boundary ? (LE(c1, t_num) && LE(t_num, 0)) : (LT(c1, t_num) && LT(t_num, 0)));
  bool u_ok = c1_pos ? (include_boundary ? (LE(0, c2) && LE(c2, c1)) : (LT(0, c2) && LT(c2, c1)))
                     : (include_boundary ? (LE(c1, c2) && LE(c2, 0)) : (LT(c1, c2) && LT(c2, 0)));
  if (t_ok && u_ok) {
    double t = static_cast<double>(t_num) / c1;
    set_out(p, a.x + t * ab_x, a.y + t * ab_y);
    return 0;  // Non-parallel with one intersection.
  }
  return -1;  // Non-parallel with no intersection.
}

// Simplified detection-only version (no output points needed); exact for integer Pt.
// Alternatively, just call the version above and pass static_cast<Pt *>(nullptr) for p and q.
template<typename Pt>
int seg_intersection(
    const Pt &a, const Pt &b, const Pt &c, const Pt &d, bool include_boundary = true
) {
  auto ab_x = b.x - a.x, ab_y = b.y - a.y;
  auto ac_x = c.x - a.x, ac_y = c.y - a.y;
  auto cd_x = d.x - c.x, cd_y = d.y - c.y;
  auto ab2 = ab_x * ab_x + ab_y * ab_y;
  auto cd2 = cd_x * cd_x + cd_y * cd_y;
  if (EQ(ab2, 0)) {
    return (include_boundary && point_on_segment(a, c, d)) ? 0 : -1;
  }
  if (EQ(cd2, 0)) {
    return (include_boundary && point_on_segment(c, a, b)) ? 0 : -1;
  }
  auto c1 = ab_x * cd_y - ab_y * cd_x;
  auto c2 = ac_x * ab_y - ac_y * ab_x;
  if (EQ(c1, 0) && EQ(c2, 0)) {  // Collinear.
    Pt res1 = std::max(std::min(a, b), std::min(c, d));
    Pt res2 = std::min(std::max(a, b), std::max(c, d));
    if (include_boundary ? !(res2 < res1) : (res1 < res2)) {
      return (res1 == res2) ? 0 : 1;
    }
    return -1;
  }
  if (EQ(c1, 0)) {
    return -1;
  }
  auto t_num = ac_x * cd_y - ac_y * cd_x;
  bool c1_pos = c1 > 0;
  bool t_ok =
      c1_pos
          ? (include_boundary ? (LE(0, t_num) && LE(t_num, c1)) : (LT(0, t_num) && LT(t_num, c1)))
          : (include_boundary ? (LE(c1, t_num) && LE(t_num, 0)) : (LT(c1, t_num) && LT(t_num, 0)));
  bool u_ok = c1_pos ? (include_boundary ? (LE(0, c2) && LE(c2, c1)) : (LT(0, c2) && LT(c2, c1)))
                     : (include_boundary ? (LE(c1, c2) && LE(c2, 0)) : (LT(c1, c2) && LT(c2, 0)));
  return (t_ok && u_ok) ? 0 : -1;
}

template<typename Pt>
Pt closest_point(double a, double b, double c, const Pt &p) {
  Pt res{};
  if (EQ(a, 0)) {
    res.x = p.x;
    res.y = -c / b;
  } else if (EQ(b, 0)) {
    res.x = -c / a;
    res.y = p.y;
  } else {
    line_intersection(a, b, c, -b, a, b * p.x - a * p.y, &res);
  }
  return res;
}
\end{lstlisting}
\Needspace*{4\baselineskip}
\textbf{Example Usage}\par
\begin{lstlisting}[style=cpp]
#include <cassert>

struct Point {
  double x, y;
  Point(double x = 0, double y = 0) : x(x), y(y) {}
  bool operator==(const Point &p) const { return x == p.x && y == p.y; }
  friend bool EQ(const Point &a, const Point &b) { return EQ(a.x, b.x) && EQ(a.y, b.y); }
  bool operator!=(const Point &p) const { return !(*this == p); }
  bool operator<(const Point &p) const { return x != p.x ? x < p.x : y < p.y; }
  bool operator>(const Point &p) const { return p < *this; }
};

struct PointI {
  int x, y;
  PointI(int x = 0, int y = 0) : x(x), y(y) {}
  bool operator==(const PointI &p) const { return x == p.x && y == p.y; }
  bool operator!=(const PointI &p) const { return !(*this == p); }
  bool operator<(const PointI &p) const { return x != p.x ? x < p.x : y < p.y; }
  bool operator>(const PointI &p) const { return p < *this; }
};

int main() {
  Point p, q;
  assert(line_intersection(-1.0, 1.0, 0.0, 1.0, 1.0, -3.0, &p) == 0);
  assert(EQ(p, Point(1.5, 1.5)));
  assert(line_intersection(Point(0, 0), Point(1, 1), Point(0, 4), Point(4, 0), &p) == 0);
  assert(EQ(p, Point(2, 2)));

  auto test = [&](double a, double b, double c, double d, double e, double f, double g, double h) {
    return seg_intersection(Point(a, b), Point(c, d), Point(e, f), Point(g, h), &p, &q);
  };
  assert(0 == test(-4, 0, 4, 0, 0, -4, 0, 4) && EQ(p, Point(0, 0)));
  assert(0 == test(0, 0, 10, 10, 2, 2, 16, 4) && EQ(p, Point(2, 2)));
  assert(0 == test(-2, 2, -2, -2, -2, 0, 0, 0) && EQ(p, Point(-2, 0)));
  assert(0 == test(0, 4, 4, 4, 4, 0, 4, 8) && EQ(p, Point(4, 4)));
  assert(1 == test(10, 10, 0, 0, 2, 2, 6, 6));
  assert(EQ(p, Point(2, 2)) && EQ(q, Point(6, 6)));
  assert(-1 == test(6, 8, 8, 10, 12, 12, 4, 4));
  assert(-1 == test(-4, 2, -8, 8, 0, 0, -4, 6));

  assert(EQ(Point(2.5, 2.5), closest_point(-1.0, -1.0, 5.0, Point(0, 0))));
  assert(EQ(Point(3, 0), closest_point(1.0, 0.0, -3.0, Point(0, 0))));

  // Detection-only with integer coordinates (exact, no float needed).
  assert(seg_intersection(PointI(0, 0), PointI(4, 4), PointI(0, 4), PointI(4, 0)) == 0);
  assert(seg_intersection(PointI(0, 0), PointI(1, 0), PointI(2, 0), PointI(3, 0)) == -1);
  assert(seg_intersection(PointI(0, 0), PointI(4, 4), PointI(0, 4), PointI(4, 0), &p) == 0);
  assert(EQ(p, Point(2, 2)));

  // include_boundary = false treats endpoint-only contact as non-intersecting. The T-junction
  // and shared-endpoint cases below intersect by default but are rejected when the flag is off.
  assert(0 == seg_intersection(Point(0, 4), Point(4, 4), Point(4, 0), Point(4, 8), &p, &q));
  assert(-1 == seg_intersection(Point(0, 4), Point(4, 4), Point(4, 0), Point(4, 8), &p, &q, false));
  assert(-1 == seg_intersection(PointI(0, 0), PointI(2, 2), PointI(2, 2), PointI(6, 6), false));
  return 0;
}
\end{lstlisting}
\subsection{Circle Intersection}
\label{sec:7.2.4}
Circle tangent and intersection calculations in two dimensions. Such calculations are floating-point by nature, so this section mostly computes in \inlinecode{double}. The input point \inlinecode{p} to \inlinecode{tangent()} can be any type with numeric \inlinecode{.x} and \inlinecode{.y} fields, and intersection point outputs are written through a caller-supplied point type.

Circle radii are normalized to be nonnegative.

\begin{compactitem}
  \item \inlinecode{tangent(c, p, \&l1, \&l2)} determines the line(s) tangent to circle \inlinecode{c} that pass through point \inlinecode{p}, returning $-1$ if there is no tangent line because \inlinecode{p} is strictly inside \inlinecode{c}, $0$ if there is exactly one tangent line because \inlinecode{p} is on the boundary of \inlinecode{c} (in which case the line will be stored into pointer \inlinecode{l1} if it's not \inlinecode{nullptr}), or $1$ if there are two tangent lines because \inlinecode{p} is strictly outside of \inlinecode{c} (in which case the lines will be stored into pointers \inlinecode{l1} and \inlinecode{l2} if they are not \inlinecode{nullptr}). The circle must have positive radius.
  \item \inlinecode{intersection(c, l, \&p, \&q)} determines the intersection between the circle \inlinecode{c} and line \inlinecode{l}, returning $-1$ if there is no intersection, $0$ if the line has one intersection point because the line is tangent (in which case it will be stored into pointer \inlinecode{p} if it's not \inlinecode{nullptr}), or $1$ if there are two intersection points because the line crosses through the circle (stored independently into \inlinecode{p} and \inlinecode{q} when those pointers are non-null). The line must be valid.
  \item \inlinecode{intersection(c1, c2, \&p, \&q)} determines the intersection points between two circles \inlinecode{c1} and \inlinecode{c2}, returning $-2$ if circle \inlinecode{c2} completely encloses circle \inlinecode{c1}, $-1$ if circle \inlinecode{c1} completely encloses circle \inlinecode{c2}, $0$ if the circles are completely disjoint, $1$ if the circles are tangent with one intersection (stored in \inlinecode{p} if it's not \inlinecode{nullptr}), $2$ if the circles intersect at two points (stored independently in \inlinecode{p} and \inlinecode{q} when those pointers are non-null), or $3$ if the circles are equal and intersect at infinite points.
  \item \inlinecode{intersection\_area(c1, c2)} returns the intersection area of circles \inlinecode{c1} and \inlinecode{c2}.
\end{compactitem}

\textbf{Time Complexity}: $O(1)$ per operation.

\textbf{Space Complexity}: $O(1)$ auxiliary for all operations.

\begin{lstlisting}[style=cpp]
#include <algorithm>
#include <cmath>
#include <type_traits>

const double EPS = 1e-9, PI = std::acos(-1.0);

template<typename T, typename U, typename C = std::common_type_t<T, U>>
bool EQ(T a, U b) {
  if constexpr (std::is_floating_point_v<C>) return C(a) == C(b) || std::fabs(C(a) - C(b)) <= EPS;
  return C(a) == C(b);
}

template<typename T, typename U, typename C = std::common_type_t<T, U>>
bool LT(T a, U b) {
  if constexpr (std::is_floating_point_v<C>) return C(a) < C(b) - EPS;
  return C(a) < C(b);
}

template<typename T, typename U>
bool LE(T a, U b) {
  return !LT(b, a);
}

struct Circle {
  double h, k, r;

  Circle(double h, double k, double r) : h(h), k(k), r(fabs(r)) {}
};

struct Line {
  double a, b, c;

  Line() : a(0), b(0), c(0) {}

  Line(double a, double b, double c) {
    if (!EQ(b, 0)) {
      this->a = a / b;
      this->c = c / b;
      this->b = 1;
    } else {
      this->c = c / a;
      this->a = 1;
      this->b = 0;
    }
  }

  Line(double px, double py, double qx, double qy) : a(0), b(0), c(0) {
    if (EQ(px, qx)) {
      if (!EQ(py, qy)) {  // Vertical line.
        a = 1;
        b = 0;
        c = -px;
      }  // Else, invalid line.
    } else {
      a = -static_cast<double>(py - qy) / (px - qx);
      b = 1;
      c = -(a * px) - (b * py);
    }
  }

  bool operator==(const Line &l) const { return a == l.a && b == l.b && c == l.c; }

  friend bool EQ(const Line &l1, const Line &l2) {
    return EQ(l1.a, l2.a) && EQ(l1.b, l2.b) && EQ(l1.c, l2.c);
  }
};

template<typename Pt>
int tangent(const Circle &c, const Pt &p, Line *l1 = nullptr, Line *l2 = nullptr) {
  auto sqnorm = [](double x, double y) { return x * x + y * y; };
  double vopx = p.x - c.h, vopy = p.y - c.k;
  if (EQ(sqnorm(vopx, vopy), c.r * c.r)) {  // Point on an edge.
    if (l1 != nullptr) {                    // Tangent through p is perpendicular to the radius cp.
      *l1 = Line(vopx, vopy, -(vopx * p.x + vopy * p.y));
    }
    return 0;
  }
  if (LE(sqnorm(vopx, vopy), c.r * c.r)) {
    return -1;  // Point inside Circle, no intersection.
  }
  double qx = vopx / c.r, qy = vopy / c.r;
  double n = sqnorm(qx, qy), d = qy * std::sqrt(n - 1.0);
  double t1x = (qx - d) / n, t1y = c.k;
  double t2x = (qx + d) / n, t2y = c.k;
  if (!EQ(qy, 0)) {  // Common case.
    t1y += c.r * (1.0 - t1x * qx) / qy;
    t2y += c.r * (1.0 - t2x * qx) / qy;
  } else {  // Point at center horizontal, y = 0.
    d = c.r * std::sqrt(1.0 - t1x * t1x);
    t1y += d;
    t2y -= d;
  }
  t1x = t1x * c.r + c.h;
  t2x = t2x * c.r + c.h;
  // Note: here, t1 and t2 are the two points of tangencies
  if (l1 != nullptr) {
    *l1 = Line(p.x, p.y, t1x, t1y);
  }
  if (l2 != nullptr) {
    *l2 = Line(p.x, p.y, t2x, t2y);
  }
  return 1;
}

// Guard to check outputs are real coordinates, since intersections are generally not integral.
template<typename OutPt>
constexpr bool is_float_pt =
    std::is_floating_point_v<decltype(OutPt::x)> && std::is_floating_point_v<decltype(OutPt::y)>;

template<typename OutPt>
int intersection(const Circle &c, const Line &l, OutPt *p = nullptr, OutPt *q = nullptr) {
  static_assert(is_float_pt<OutPt>, "output point coordinates must be floating-point");
  auto set_out = [](auto *out, auto x, auto y) {
    if (out != nullptr) {
      out->x = x;
      out->y = y;
    }
  };
  double v = c.h * l.a + c.k * l.b + l.c;
  double aabb = l.a * l.a + l.b * l.b;
  double disc = v * v / aabb - c.r * c.r;
  if (disc > EPS) {
    return -1;
  }
  double x0 = -l.a * v / aabb, y0 = -l.b * v / aabb;
  if (disc > -EPS) {
    set_out(p, x0 + c.h, y0 + c.k);
    return 0;
  }
  double k = std::sqrt(std::max(0.0, disc / -aabb));
  set_out(p, x0 + k * l.b + c.h, y0 - k * l.a + c.k);
  set_out(q, x0 - k * l.b + c.h, y0 + k * l.a + c.k);
  return 1;
}

template<typename OutPt>
int intersection(const Circle &c1, const Circle &c2, OutPt *p = nullptr, OutPt *q = nullptr) {
  static_assert(is_float_pt<OutPt>, "output point coordinates must be floating-point");
  auto set_out = [](auto *out, auto x, auto y) {
    if (out != nullptr) {
      out->x = x;
      out->y = y;
    }
  };
  if (EQ(c1.h, c2.h) && EQ(c1.k, c2.k)) {
    return EQ(c1.r, c2.r) ? 3 : (c1.r > c2.r ? -1 : -2);
  }
  double d12x = c2.h - c1.h, d12y = c2.k - c1.k;
  double d = std::hypot(d12x, d12y);
  if (LT(c1.r + c2.r, d)) {
    return 0;
  }
  if (LT(d, fabs(c1.r - c2.r))) {
    return c1.r > c2.r ? -1 : -2;
  }
  double a = (c1.r * c1.r - c2.r * c2.r + d * d) / (2 * d);
  double x0 = c1.h + (d12x * a / d), y0 = c1.k + (d12y * a / d);
  double s = std::sqrt(std::max(0.0, c1.r * c1.r - a * a));
  double rx = -d12y * s / d, ry = d12x * s / d;
  if (EQ(rx, 0) && EQ(ry, 0)) {
    set_out(p, x0, y0);
    return 1;
  }
  set_out(p, x0 - rx, y0 - ry);
  set_out(q, x0 + rx, y0 + ry);
  return 2;
}

double intersection_area(const Circle &c1, const Circle &c2) {
  double r = std::min(c1.r, c2.r), R = std::max(c1.r, c2.r);
  double d = std::hypot(c2.h - c1.h, c2.k - c1.k);
  if (LE(d, R - r)) {
    return PI * r * r;
  }
  if (LE(R + r, d)) {
    return 0;
  }
  double alpha = std::clamp((d * d + r * r - R * R) / 2 / d / r, -1.0, 1.0);
  double beta = std::clamp((d * d + R * R - r * r) / 2 / d / R, -1.0, 1.0);
  return r * r * std::acos(alpha) + R * R * std::acos(beta) -
         0.5 * std::sqrt((-d + r + R) * (d + r - R) * (d - r + R) * (d + r + R));
}
\end{lstlisting}
\Needspace*{4\baselineskip}
\textbf{Example Usage}\par
\begin{lstlisting}[style=cpp]
#include <cassert>

struct Point {
  double x, y;
  Point(double x = 0, double y = 0) : x(x), y(y) {}
  bool operator==(const Point &p) const { return x == p.x && y == p.y; }
  friend bool EQ(const Point &a, const Point &b) { return EQ(a.x, b.x) && EQ(a.y, b.y); }
};

int main() {
  Line l1, l2;
  assert(-1 == tangent(Circle(0, 0, 4), Point(1, 1), &l1, &l2));
  assert(0 == tangent(Circle(0, 0, sqrt(2)), Point(1, 1), &l1, &l2));
  assert(EQ(l1, Line(-1, -1, 2)));
  assert(1 == tangent(Circle(0, 0, 2), Point(2, 2), &l1, &l2));
  assert(EQ(l1, Line(0, -2, 4)));
  assert(EQ(l2, Line(2, 0, -4)));

  Point p, q;
  assert(-1 == intersection(Circle(1, 1, 3), Line(5, 3, -30), &p, &q));
  assert(0 == intersection(Circle(1, 1, 3), Line(0, 1, -4), &p, &q));
  assert(EQ(p, Point(1, 4)));
  assert(1 == intersection(Circle(1, 1, 3), Line(0, 1, -1), &p, &q));
  assert(EQ(p, Point(4, 1)));
  assert(EQ(q, Point(-2, 1)));
  assert(1 == intersection(Circle(1, 1, 3), Line(0, 1, -1), &p));
  assert(EQ(p, Point(4, 1)));

  assert(-2 == intersection(Circle(1, 1, 1), Circle(0, 0, 3), &p, &q));
  assert(-1 == intersection(Circle(0, 0, 3), Circle(1, 1, 1), &p, &q));
  assert(0 == intersection(Circle(5, 0, 4), Circle(-5, 0, 4), &p, &q));
  assert(1 == intersection(Circle(-5, 0, 5), Circle(5, 0, 5), &p, &q));
  assert(EQ(p, Point(0, 0)));
  assert(2 == intersection(Circle(-0.5, 0, 1), Circle(0.5, 0, 1), &p, &q));
  assert(EQ(p, Point(0, -sqrt(3) / 2)));
  assert(EQ(q, Point(0, sqrt(3) / 2)));
  assert(
      2 == intersection(Circle(-0.5, 0, 1), Circle(0.5, 0, 1), static_cast<Point *>(nullptr), &q)
  );
  assert(EQ(q, Point(0, sqrt(3) / 2)));

  // Each circle passes through the other's center.
  double r = 3, a = intersection_area(Circle(-r / 2, 0, r), Circle(r / 2, 0, r));
  assert(EQ(a, r * r * (2 * PI / 3 - sqrt(3) / 2)));
  return 0;
}
\end{lstlisting}
\subsection{Circle Tangents}
\label{sec:7.2.5}
Given two circles, compute their common tangent lines by returning the two points of tangency on each circle. External tangents touch both circles on the same side of the line joining their centers; internal tangents cross between the circles. A point can be treated as a zero-radius circle, so this also covers tangents from a point to a circle.

\begin{compactitem}
  \item \inlinecode{circle\_tangents(c1, r1, c2, r2, internal = false)} returns $0$, $1$, or $2$ tangent point pairs. Each pair $(p, q)$ means the tangent line touches the first circle at $p$ and the second circle at $q$. Set \inlinecode{internal = true} to get internal tangents; otherwise external tangents are returned.
  \item \inlinecode{all\_circle\_tangents(c1, r1, c2, r2)} returns all external and internal tangents.
\end{compactitem}

If the two circles are identical, there are infinitely many tangents and the functions return an empty vector. If one circle strictly contains the other, there are no external tangents; if the circles overlap, there are no internal tangents.

\textbf{Time Complexity}: $O(1)$ per call.

\textbf{Space Complexity}: $O(1)$ auxiliary, excluding the returned vector.

\begin{lstlisting}[style=cpp]
#include <algorithm>
#include <cmath>
#include <type_traits>
#include <utility>
#include <vector>

const double EPS = 1e-9;

template<typename T, typename U, typename C = std::common_type_t<T, U>>
bool EQ(T a, U b) {
  if constexpr (std::is_floating_point_v<C>) return C(a) == C(b) || std::fabs(C(a) - C(b)) <= EPS;
  return C(a) == C(b);
}

template<typename T, typename U, typename C = std::common_type_t<T, U>>
bool LT(T a, U b) {
  if constexpr (std::is_floating_point_v<C>) return C(a) < C(b) - EPS;
  return C(a) < C(b);
}

struct Point {
  double x, y;
  Point(double x = 0, double y = 0) : x(x), y(y) {}
  bool operator==(const Point &p) const { return x == p.x && y == p.y; }
  friend bool EQ(const Point &a, const Point &b) { return EQ(a.x, b.x) && EQ(a.y, b.y); }
  Point operator+(const Point &p) const { return {x + p.x, y + p.y}; }
  Point operator-(const Point &p) const { return {x - p.x, y - p.y}; }
  Point operator*(double k) const { return {x * k, y * k}; }
  Point operator/(double k) const { return {x / k, y / k}; }
  double sqnorm() const { return x * x + y * y; }
  Point rotate90() const { return {-y, x}; }
};

template<typename PtA, typename PtB>
std::vector<std::pair<Point, Point>> circle_tangents(
    const PtA &c1, double r1, const PtB &c2, double r2, bool internal = false
) {
  r1 = fabs(r1);
  r2 = fabs(r2);
  if (internal) {
    r2 = -r2;
  }
  Point d(c2.x - c1.x, c2.y - c1.y);
  double dr = r1 - r2, d2 = d.sqnorm(), h2 = d2 - dr * dr;
  if (EQ(d2, 0) || LT(h2, 0)) {
    return {};
  }
  std::vector<std::pair<Point, Point>> res;
  for (double s : {-1.0, 1.0}) {
    Point v = (d * dr + d.rotate90() * (std::sqrt(std::max(0.0, h2)) * s)) / d2;
    res.emplace_back(Point(c1.x, c1.y) + v * r1, Point(c2.x, c2.y) + v * r2);
  }
  if (EQ(h2, 0)) {
    res.pop_back();
  }
  return res;
}

template<typename PtA, typename PtB>
std::vector<std::pair<Point, Point>> all_circle_tangents(
    const PtA &c1, double r1, const PtB &c2, double r2
) {
  auto res = circle_tangents(c1, r1, c2, r2);
  auto internal = circle_tangents(c1, r1, c2, r2, true);
  res.insert(res.end(), internal.begin(), internal.end());
  return res;
}
\end{lstlisting}
\Needspace*{4\baselineskip}
\textbf{Example Usage}\par
\begin{lstlisting}[style=cpp]
#include <cassert>
using namespace std;

int main() {
  Point a(0, 0), b(4, 0);
  auto ext = circle_tangents(a, 1, b, 1);
  assert(ext.size() == 2);
  assert(EQ(ext[0].first, Point(0, -1)) && EQ(ext[0].second, Point(4, -1)));
  assert(EQ(ext[1].first, Point(0, 1)) && EQ(ext[1].second, Point(4, 1)));

  auto in = circle_tangents(a, 1, b, 1, true);
  assert(in.size() == 2);
  assert(
      EQ(in[0].first, Point(0.5, -sqrt(3.0) / 2)) && EQ(in[0].second, Point(3.5, sqrt(3.0) / 2))
  );
  assert(
      EQ(in[1].first, Point(0.5, sqrt(3.0) / 2)) && EQ(in[1].second, Point(3.5, -sqrt(3.0) / 2))
  );
  auto all = ext;
  all.insert(all.end(), in.begin(), in.end());
  assert(all_circle_tangents(a, 1, b, 1) == all);
  auto point_to_circle = circle_tangents(Point(0, 0), 0, Point(2, 0), 1);
  assert(point_to_circle.size() == 2);
  assert(
      EQ(point_to_circle[0].first, Point(0, 0)) &&
      EQ(point_to_circle[0].second, Point(1.5, -sqrt(3.0) / 2))
  );
  assert(
      EQ(point_to_circle[1].first, Point(0, 0)) &&
      EQ(point_to_circle[1].second, Point(1.5, sqrt(3.0) / 2))
  );
  assert(circle_tangents(Point(0, 0), 5, Point(1, 0), 1).empty());        // contained
  assert(circle_tangents(Point(0, 0), 2, Point(3, 0), 2, true).empty());  // overlapping
  return 0;
}
\end{lstlisting}

\section{\texorpdfstring{Polygons and Point Sets}{7.3 Polygons and Point Sets}}
\setcounter{section}{3}
\setcounter{subsection}{0}
\subsection{Polygon Sorting and Area}
\label{sec:7.3.1}
Given a list of distinct points in two dimensions, order them into a valid polygon and determine the area. The sorting comparator orders two points by the sign of their cross product around a chosen center, that is, by angle. To form a simple polygon, choose a center inside the points' convex hull that differs from every input point; the arithmetic mean used in the example is usually convenient. The area functions use the shoelace formula, summing the cross products of consecutive vertex pairs. All functions accept \inlinecode{Point}, \inlinecode{PointI}, or \inlinecode{PointL} from \secref{7.1.1}, or any compatible point type with numeric \inlinecode{.x} and \inlinecode{.y} fields and the comparison operators used by the requested operation.

\begin{compactitem}
  \item \inlinecode{cw\_comp(a, b, c)} returns whether point \inlinecode{a} compares clockwise before \inlinecode{b} about \inlinecode{c}.
  \item \inlinecode{polygon\_area\_2x(lo, hi)} returns exactly double the area of the polygon with vertices specified by the range $[{\inlinecodemath{lo}}, {\inlinecodemath{hi}})$ of points in either clockwise or counter-clockwise order. The return value is integral or floating-point, depending on the input point type. For integer vertices, divide by $2$ in the caller if the exact area is needed.
  \item \inlinecode{polygon\_area(lo, hi)} returns the area as \inlinecode{double}.
  \item \inlinecode{polygon\_centroid(lo, hi)} returns the centroid (center of mass) of a non-degenerate simple polygon with vertices in boundary order as an \inlinecode{std::pair\textless{}double, double\textgreater{}}. It uses signed shoelace weights, so either clockwise or counter-clockwise inputs will work.
\end{compactitem}

Overflow warning: \inlinecode{cw\_comp} and \inlinecode{polygon\_area\_2x} form cross products that grow like the squared coordinate magnitude (and the shoelace sum accumulates over all vertices). For integer point types use a 64-bit coordinate type (e.g. \inlinecode{PointL} from \secref{7.1.1}) for large or numerous coordinates.

\Needspace*{3\baselineskip}
\textbf{Time Complexity}:
\begin{compactitem}
  \item $O(n)$ per call to \inlinecode{polygon\_area\_2x()}, \inlinecode{polygon\_area()}, and \inlinecode{polygon\_centroid()}, where $n$ is the number of points.
  \item $O(1)$ per call to \inlinecode{cw\_comp()} and the comparators.
\end{compactitem}

\textbf{Space Complexity}: $O(1)$ auxiliary for all operations.

\begin{lstlisting}[style=cpp]
#include <cassert>
#include <cmath>
#include <stdexcept>
#include <type_traits>
#include <utility>

const double EPS = 1e-9;

template<typename T, typename U, typename C = std::common_type_t<T, U>>
bool EQ(T a, U b) {
  if constexpr (std::is_floating_point_v<C>) return C(a) == C(b) || std::fabs(C(a) - C(b)) <= EPS;
  return C(a) == C(b);
}

// Comparator (clockwise angular order about c). Exact: comparisons are pure sign tests on
// coordinate differences and the cross product, so it is a valid strict weak ordering and is
// exact for integer-coordinate points.
template<typename Pt>
bool cw_comp(const Pt &a, const Pt &b, const Pt &c) {
  if (a.x - c.x >= 0 && b.x - c.x < 0) {
    return true;
  }
  if (a.x - c.x < 0 && b.x - c.x >= 0) {
    return false;
  }
  if (a.x - c.x == 0 && b.x - c.x == 0) {
    if (a.y - c.y >= 0 || b.y - c.y >= 0) {
      return a.y > b.y;
    }
    return b.y > a.y;
  }
  auto acx = a.x - c.x, acy = a.y - c.y;
  auto bcx = b.x - c.x, bcy = b.y - c.y;
  auto det = acx * bcy - acy * bcx;  // Overflow warning.
  if (det == 0) {
    auto acnorm = acx * acx + acy * acy;  // Overflow warning.
    auto bcnorm = bcx * bcx + bcy * bcy;
    return acnorm > bcnorm;
  }
  return det < 0;
}

// Returns 2 * area. Result is exact (integer) for integer-coordinate points.
template<typename It>
auto polygon_area_2x(It lo, It hi) {
  using T = decltype(lo->x * lo->y);
  if (lo == hi) {
    return T{0};
  }
  T area = 0;
  for (It i = lo, j = hi - 1; i != hi; j = i++) {
    area += static_cast<T>(j->x - i->x) * static_cast<T>(j->y + i->y);  // Overflow warning.
  }
  return area < T{0} ? -area : area;
}

template<typename It>
double polygon_area(It lo, It hi) {
  return static_cast<double>(polygon_area_2x(lo, hi)) / 2.0;
}

template<typename It>
std::pair<double, double> polygon_centroid(It lo, It hi) {
  assert(hi - lo >= 3);
  double cx = 0, cy = 0, area2 = 0;
  for (It i = lo, j = hi - 1; i != hi; j = i++) {
    double cross = static_cast<double>(j->x) * i->y - static_cast<double>(i->x) * j->y;
    cx += (static_cast<double>(j->x) + i->x) * cross;
    cy += (static_cast<double>(j->y) + i->y) * cross;
    area2 += cross;
  }
  if (EQ(area2, 0)) {
    throw std::runtime_error("Cannot compute centroid of zero-area polygon.");
  }
  return {cx / (3 * area2), cy / (3 * area2)};
}
\end{lstlisting}
\Needspace*{4\baselineskip}
\textbf{Example Usage}\par
\begin{lstlisting}[style=cpp]
#include <algorithm>
#include <cassert>
#include <random>
#include <vector>
using namespace std;

struct Point {
  double x, y;
  Point(double x = 0, double y = 0) : x(x), y(y) {}
  bool operator==(const Point &p) const { return x == p.x && y == p.y; }
  friend bool EQ(const Point &a, const Point &b) { return EQ(a.x, b.x) && EQ(a.y, b.y); }
  bool operator!=(const Point &p) const { return !(*this == p); }
  bool operator<(const Point &p) const { return x != p.x ? x < p.x : y < p.y; }
  bool operator>(const Point &p) const { return p < *this; }
};

struct PointI {
  int x, y;
  PointI(int x = 0, int y = 0) : x(x), y(y) {}
  bool operator==(const PointI &p) const { return x == p.x && y == p.y; }
  bool operator!=(const PointI &p) const { return !(*this == p); }
  bool operator<(const PointI &p) const { return x != p.x ? x < p.x : y < p.y; }
  bool operator>(const PointI &p) const { return p < *this; }
};

template<typename It>
Point mean_center(It lo, It hi) {
  assert(lo != hi);
  double x_sum = 0, y_sum = 0, n = hi - lo;
  for (It it = lo; it != hi; ++it) {
    x_sum += it->x;
    y_sum += it->y;
  }
  return Point(x_sum / n, y_sum / n);
}

int main() {
  vector<Point> points{Point(1, 3), Point(1, 2), Point(2, 1), Point(0, 0), Point(-1, 3)};
  vector<Point> v(points);
  mt19937 rng(1234567);  // Fixed seed for reproducibility.
  shuffle(v.begin(), v.end(), rng);
  Point c = mean_center(v.begin(), v.end());
  assert(EQ(c, Point(0.6, 1.8)));
  sort(v.begin(), v.end(), [c](const Point &a, const Point &b) { return cw_comp(a, b, c); });
  assert((v == vector<Point>{{1, 3}, {1, 2}, {2, 1}, {0, 0}, {-1, 3}}));
  assert(EQ(polygon_area(v.begin(), v.end()), 5));

  sort(v.begin(), v.end(), [c](const Point &a, const Point &b) { return cw_comp(b, a, c); });
  assert((v == vector<Point>{{-1, 3}, {0, 0}, {2, 1}, {1, 2}, {1, 3}}));
  assert(EQ(polygon_area(v.begin(), v.end()), 5));

  // Integer points: polygon_area_2x is exact (no float arithmetic).
  vector<PointI> iv{{0, 0}, {4, 0}, {0, 3}};            // right triangle, area = 6
  assert(polygon_area_2x(iv.begin(), iv.end()) == 12);  // exact int
  assert(EQ(polygon_area(iv.begin(), iv.end()), 6.0));
  auto centroid = polygon_centroid(iv.begin(), iv.end());
  assert(EQ(centroid.first, 4.0 / 3.0) && EQ(centroid.second, 1.0));
  return 0;
}
\end{lstlisting}
\subsection{Point-in-Polygon}
\label{sec:7.3.2}
Given a point \inlinecode{p} and a polygon, determines whether \inlinecode{p} lies inside using ray casting. A ray cast from \inlinecode{p} must cross the polygon's boundary an odd number of times exactly when \inlinecode{p} is inside, so the function counts crossings between a horizontal ray and each polygon edge. The function is templated on the point type \inlinecode{Pt} and accepts \inlinecode{Point}, \inlinecode{PointI}, or \inlinecode{PointL} from \secref{7.1.1}, or any compatible type with numeric \inlinecode{.x} and \inlinecode{.y} fields. All comparisons are exact (no epsilon): the ray-crossing parity needs a consistent comparison to be correct, and the result is exact for integer-coordinate points.

\begin{compactitem}
  \item \inlinecode{point\_in\_polygon(p, lo, hi, include\_boundary = true)} returns whether \inlinecode{p} lies within the polygon with vertices specified by the range $[{\inlinecodemath{lo}}, {\inlinecodemath{hi}})$ of points in either clockwise or counter-clockwise order. The \inlinecode{include\_boundary} flag controls whether points exactly on an edge or vertex count as inside.
\end{compactitem}

Overflow warning: the edge orientation test forms a cross product that grows like the squared coordinate magnitude. For integer point types use a 64-bit coordinate type (e.g. \inlinecode{PointL} from \secref{7.1.1}) once coordinates exceed roughly a few tens of thousands.

\textbf{Time Complexity}: $O(n)$ per call, where $n$ is the distance between \inlinecode{lo} and \inlinecode{hi}.

\textbf{Space Complexity}: $O(1)$ auxiliary.

\begin{lstlisting}[style=cpp]
template<typename Pt>
auto cross(const Pt &a, const Pt &b, const Pt &o) {
  // Overflow risk for integer Pt: these products are ~O(max_coord^2); use int64_t if necessary.
  return (a.x - o.x) * (b.y - o.y) - (a.y - o.y) * (b.x - o.x);
}

// Detection is exact for integer-coordinate Pt.
template<typename Pt, typename It>
bool point_in_polygon(const Pt &p, It lo, It hi, bool include_boundary = true) {
  if (lo == hi) {
    return false;
  }
  bool res = false;
  for (It i = lo, j = hi - 1; i != hi; j = i++) {
    if (i->y == p.y && (i->x == p.x || (j->y == p.y && (i->x <= p.x || j->x <= p.x)))) {
      return include_boundary;
    }
    if ((p.y < i->y) != (p.y < j->y)) {
      auto det = cross(*i, *j, p);
      if (det == 0) {
        return include_boundary;
      }
      if ((0 < det) != (i->y < j->y)) {
        res = !res;
      }
    }
  }
  return res;
}
\end{lstlisting}
\Needspace*{4\baselineskip}
\textbf{Example Usage}\par
\begin{lstlisting}[style=cpp]
#include <cassert>
#include <vector>
using namespace std;

struct Point {
  double x, y;
  Point(double x = 0, double y = 0) : x(x), y(y) {}
};

struct PointI {
  int x, y;
  PointI(int x = 0, int y = 0) : x(x), y(y) {}
};

int main() {
  vector<Point> poly{Point(-1, 3), Point(1, 3), Point(2, 1), Point(0, 0)};
  assert(point_in_polygon(Point(1, 2), poly.begin(), poly.end()));
  assert(point_in_polygon(Point(0, 3), poly.begin(), poly.end()));
  assert(!point_in_polygon(Point(0, 3), poly.begin(), poly.end(), false));
  assert(!point_in_polygon(Point(0, 3.01), poly.begin(), poly.end()));
  assert(!point_in_polygon(Point(2, 2), poly.begin(), poly.end()));

  // Integer-coordinate points: exact detection.
  vector<PointI> ipoly{{0, 0}, {4, 0}, {4, 4}, {0, 4}};
  assert(point_in_polygon(PointI(2, 2), ipoly.begin(), ipoly.end()));
  assert(!point_in_polygon(PointI(4, 2), ipoly.begin(), ipoly.end(), false));
  assert(!point_in_polygon(PointI(5, 2), ipoly.begin(), ipoly.end()));
  return 0;
}
\end{lstlisting}
\subsection{Convex Hull and Diametral Pair}
\label{sec:7.3.3}
The convex hull of points in two dimensions is the smallest convex set containing them. A diametral pair is a pair of input points at maximum distance. Monotone chain computes the hull by sorting the points lexicographically and building the lower and upper boundaries in one pass each, popping any point that would fail to create a counter-clockwise turn. Rotating calipers then walks two antipodal pointers around the hull, advancing whichever increases the separation and visiting every candidate diametral pair in linear time. Both functions accept either floating-point or integral coordinates, and use exact comparisons for integral points.

\begin{compactitem}
  \item \inlinecode{convex\_hull(p)} returns the convex hull of points \inlinecode{p} in counter-clockwise order. Duplicate input points and collinear points along hull edges are omitted. To instead return the hull points in clockwise order, replace the cross product comparisons \inlinecode{\textless{}= 0} with \inlinecode{\textgreater{}= 0}.
  \item \inlinecode{diametral\_pair(p)} returns the maximum-distance pair of points in \inlinecode{p}.
\end{compactitem}

Overflow warning: \inlinecode{cross()} and the squared distance in \inlinecode{diametral\_pair()} grow like the squared coordinate magnitude. With 32-bit \inlinecode{int} coordinates they overflow once coordinates exceed a few tens of thousands; use a 64-bit (\inlinecode{int64\_t}) coordinate type for larger integer inputs.

\textbf{Time Complexity}: $O(n \log n)$ per call, where $n$ is the number of points.

\textbf{Space Complexity}: $O(n)$ auxiliary for storage of the convex hull.

\begin{lstlisting}[style=cpp]
#include <algorithm>
#include <cmath>
#include <utility>
#include <vector>

template<typename Pt>
auto cross(const Pt &a, const Pt &b, const Pt &o) {
  // Overflow risk for integer Pt: these products are ~O(max_coord^2); use int64_t if necessary.
  return (a.x - o.x) * (b.y - o.y) - (a.y - o.y) * (b.x - o.x);
}

// Convex hull: exact for integer-coordinate points.
template<typename Pt>
std::vector<Pt> convex_hull(std::vector<Pt> p) {
  std::sort(p.begin(), p.end());
  auto same_point = [](const Pt &a, const Pt &b) { return !(a < b) && !(b < a); };
  p.erase(std::unique(p.begin(), p.end(), same_point), p.end());
  int n = static_cast<int>(p.size());
  if (n <= 1) {
    return p;
  }
  int k = 0;
  std::vector<Pt> res(2 * n);
  for (const Pt &q : p) {
    while (k >= 2 && cross(res[k - 1], q, res[k - 2]) <= 0) {
      k--;
    }
    res[k++] = q;
  }
  int t = k + 1;
  for (int i = n - 2; i >= 0; i--) {
    while (k >= t && cross(res[k - 1], p[i], res[k - 2]) <= 0) {
      k--;
    }
    res[k++] = p[i];
  }
  res.resize(k - 1);
  return res;
}

// Diametral pair: squared-distance comparisons are exact for integer-coordinate points.
template<typename Pt>
std::pair<Pt, Pt> diametral_pair(const std::vector<Pt> &p) {
  auto h = convex_hull(p);
  int m = static_cast<int>(h.size());
  if (m == 0) {
    return std::pair<Pt, Pt>{};
  }
  if (m == 1) {
    return std::pair<Pt, Pt>{h[0], h[0]};
  }
  if (m == 2) {
    return std::pair<Pt, Pt>{h[0], h[1]};
  }
  int k = 1;
  while (std::abs(cross(h[0], h[(k + 1) % m], h[m - 1])) > std::abs(cross(h[0], h[k], h[m - 1]))) {
    k++;
  }
  auto sqdist = [](const Pt &a, const Pt &b) {
    // Overflow risk for integer Pt: ~O(max_coord^2); use int64_t if necessary.
    auto dx = a.x - b.x, dy = a.y - b.y;
    return dx * dx + dy * dy;
  };
  auto maxsq = sqdist(h[0], h[0]);
  std::pair<Pt, Pt> res{h[0], h[0]};
  for (int i = 0, j = k; i <= k && j < m; i++) {
    auto d = sqdist(h[i], h[j]);
    if (d > maxsq) {
      maxsq = d;
      res = {h[i], h[j]};
    }
    while (j < m && std::abs(cross(h[(i + 1) % m], h[(j + 1) % m], h[i])) >
                        std::abs(cross(h[(i + 1) % m], h[j], h[i]))) {
      d = sqdist(h[i], h[(j + 1) % m]);
      if (d > maxsq) {
        maxsq = d;
        res = {h[i], h[(j + 1) % m]};
      }
      j++;
    }
  }
  return res;
}
\end{lstlisting}
\Needspace*{4\baselineskip}
\textbf{Example Usage}\par
\begin{lstlisting}[style=cpp]
#include <cassert>
#include <random>
#include <vector>
using namespace std;

struct PointI {
  int x, y;
  PointI(int x = 0, int y = 0) : x(x), y(y) {}
  bool operator==(const PointI &o) const { return x == o.x && y == o.y; }
  bool operator!=(const PointI &o) const { return !(*this == o); }
  bool operator<(const PointI &o) const { return x != o.x ? x < o.x : y < o.y; }
  bool operator>(const PointI &o) const { return o < *this; }
};

int main() {
  {
    vector<PointI> v{{1, 3}, {1, 2}, {2, 1}, {0, 0}, {-1, 3}};
    mt19937 rng(1234567);  // Fixed seed for reproducibility.
    shuffle(v.begin(), v.end(), rng);
    vector<PointI> h{{-1, 3}, {0, 0}, {2, 1}, {1, 3}};
    assert(convex_hull(v) == h);
  }
  {
    auto [p1, p2] = diametral_pair(vector<PointI>{{0, 0}, {3, 0}, {0, 3}, {1, 1}, {4, 4}});
    assert(p1 == PointI(0, 0) && p2 == PointI(4, 4));
  }
  {
    vector<PointI> v{{0, 0}, {4, 0}, {4, 4}, {0, 4}, {2, 2}};
    auto h = convex_hull(v);
    assert(h.size() == 4);  // interior point (2,2) excluded
    auto [p1, p2] = diametral_pair(v);
    // diametral pair is exact
    assert(
        (p1 == PointI(0, 0) && p2 == PointI(4, 4)) || (p1 == PointI(4, 4) && p2 == PointI(0, 0)) ||
        (p1 == PointI(4, 0) && p2 == PointI(0, 4)) || (p1 == PointI(0, 4) && p2 == PointI(4, 0))
    );
  }
  {
    assert((convex_hull(vector<PointI>{{2, 3}, {2, 3}, {2, 3}}) == vector<PointI>{{2, 3}}));
  }
  return 0;
}
\end{lstlisting}
\subsection{Convex Polygon Queries}
\label{sec:7.3.4}
Given a convex polygon in counter-clockwise order, answer containment and tangent/intersection queries. Point containment is logarithmic by splitting the polygon into triangles sharing \inlinecode{poly[0]}. The line and tangent helpers below are written as simple linear scans; this is shorter and less assumption-heavy than the classic logarithmic versions, while still documenting the exact query semantics in one place.

\begin{compactitem}
  \item \inlinecode{point\_in\_convex\_polygon(poly, p, include\_boundary = true)} returns whether point \inlinecode{p} lies inside the convex polygon \inlinecode{poly}. The polygon must be in counter-clockwise order, with no repeated final vertex. Points on the boundary count as inside unless \inlinecode{include\_boundary} is set to \inlinecode{false}.
  \item \inlinecode{line\_convex\_polygon\_intersection(poly, a, b)} describes where the infinite directed line through \inlinecode{a} $\to$ \inlinecode{b} hits the polygon: $(-1, -1)$ for no intersection, $(i, -1)$ for touching vertex $i$, $(i, i)$ for containing side $i \to i+1$, or $(i, j)$ for crossing sides $i \to i+1$ and $j \to j+1$, with side indices taken modulo the number of vertices. The polygon must have at least three vertices, and \inlinecode{a} and \inlinecode{b} must differ. When the line crosses two sides, their order is determined using a floating-point line parameter.
  \item \inlinecode{convex\_polygon\_tangents(poly, p)} returns the two tangent vertex indices from an outside point \inlinecode{p} as (\inlinecode{left}, \inlinecode{right}). The left tangent has all polygon vertices on or to the left of the directed line \inlinecode{p} $\to$ \inlinecode{poly[left]}; the right tangent is analogous for the right side.
\end{compactitem}

Overflow warning: For integer-coordinate inputs, all orientation tests are exact provided the cross products do not overflow.

\Needspace*{3\baselineskip}
\textbf{Time Complexity}:
\begin{compactitem}
  \item $O(\log n)$ per call to \inlinecode{point\_in\_convex\_polygon()}, where $n$ is the number of polygon vertices.
  \item $O(n)$ per call to \inlinecode{line\_convex\_polygon\_intersection} and \inlinecode{convex\_polygon\_tangents}.
\end{compactitem}

\textbf{Space Complexity}: $O(1)$ auxiliary.

\begin{lstlisting}[style=cpp]
#include <algorithm>
#include <cassert>
#include <utility>
#include <vector>

template<typename Pt>
auto cross(const Pt &a, const Pt &b, const Pt &o) {
  // Overflow risk for integer Pt: these products are ~O(max_coord^2); use int64_t if necessary.
  return (a.x - o.x) * (b.y - o.y) - (a.y - o.y) * (b.x - o.x);
}

template<typename T>
int sgn(const T &x) {
  return (T{0} < x) - (x < T{0});
}

template<typename Pt>
bool point_on_segment(const Pt &p, const Pt &a, const Pt &b) {
  return cross(a, b, p) == 0 && std::min(a.x, b.x) <= p.x && p.x <= std::max(a.x, b.x) &&
         std::min(a.y, b.y) <= p.y && p.y <= std::max(a.y, b.y);
}

template<typename Pt>
bool point_in_convex_polygon(
    const std::vector<Pt> &poly, const Pt &p, bool include_boundary = true
) {
  int n = static_cast<int>(poly.size());
  if (n == 0) {
    return false;
  }
  if (n == 1) {
    return include_boundary && p.x == poly[0].x && p.y == poly[0].y;
  }
  if (n == 2) {
    return include_boundary && point_on_segment(p, poly[0], poly[1]);
  }
  int left = sgn(cross(poly[1], p, poly[0]));
  int right = sgn(cross(poly[n - 1], p, poly[0]));
  if (left < 0 || right > 0) {
    return false;
  }
  if (left == 0) {
    return include_boundary && point_on_segment(p, poly[0], poly[1]);
  }
  if (right == 0) {
    return include_boundary && point_on_segment(p, poly[0], poly[n - 1]);
  }
  int lo = 1, hi = n - 1;
  while (hi - lo > 1) {
    int mid = lo + (hi - lo) / 2;
    if (cross(poly[mid], p, poly[0]) >= 0) {
      lo = mid;
    } else {
      hi = mid;
    }
  }
  int s = sgn(cross(poly[lo + 1], p, poly[lo]));
  return include_boundary ? s >= 0 : s > 0;
}

template<typename Pt>
std::pair<int, int> line_convex_polygon_intersection(
    const std::vector<Pt> &poly, const Pt &a, const Pt &b
) {
  int n = static_cast<int>(poly.size());
  assert(n >= 3 && (a.x != b.x || a.y != b.y));
  std::vector<int> on, cross_edges;
  bool has_pos = false, has_neg = false;
  for (int i = 0; i < n; i++) {
    int s = sgn(cross(b, poly[i], a));
    has_pos |= s > 0;
    has_neg |= s < 0;
    if (s == 0) {
      on.push_back(i);
    }
  }
  if (!has_pos && !has_neg) {
    return n >= 2 ? std::make_pair(0, 0) : std::make_pair(-1, -1);
  }
  if (!(has_pos && has_neg)) {
    for (int i = 0; i < static_cast<int>(on.size()); i++) {
      int j = (i + 1) % static_cast<int>(on.size());
      if ((on[i] + 1) % n == on[j]) {
        return {on[i], on[i]};
      }
    }
    return on.empty() ? std::make_pair(-1, -1) : std::make_pair(on[0], -1);
  }
  for (int i = 0; i < n; i++) {
    int j = (i + 1) % n;
    int si = sgn(cross(b, poly[i], a)), sj = sgn(cross(b, poly[j], a));
    if (si == 0 || si * sj < 0) {
      cross_edges.push_back(i);
    }
  }
  if (cross_edges.size() == 1) {
    return {cross_edges[0], -1};
  }
  auto line_parameter = [&](int i) {
    int j = (i + 1) % n;
    auto ex = poly[j].x - poly[i].x, ey = poly[j].y - poly[i].y;
    auto ax = a.x, ay = a.y, dx = b.x - a.x, dy = b.y - a.y;
    auto det = dx * ey - dy * ex;
    return static_cast<double>((poly[i].x - ax) * ey - (poly[i].y - ay) * ex) / det;
  };
  if (line_parameter(cross_edges[1]) < line_parameter(cross_edges[0])) {
    std::swap(cross_edges[0], cross_edges[1]);
  }
  return {cross_edges[0], cross_edges[1]};
}

template<typename Pt>
std::pair<int, int> convex_polygon_tangents(const std::vector<Pt> &poly, const Pt &p) {
  int n = static_cast<int>(poly.size());
  std::pair<int, int> res{-1, -1};
  for (int i = 0; i < n; i++) {
    int prev = (i + n - 1) % n, next = (i + 1) % n;
    int s1 = sgn(cross(poly[i], poly[prev], p));
    int s2 = sgn(cross(poly[i], poly[next], p));
    bool left = s1 >= 0 && s2 >= 0;
    bool right = s1 <= 0 && s2 <= 0;
    if (left) {
      res.first = i;
    }
    if (right) {
      res.second = i;
    }
  }
  return res;
}
\end{lstlisting}
\Needspace*{4\baselineskip}
\textbf{Example Usage}\par
\begin{lstlisting}[style=cpp]
#include <algorithm>
#include <cstdint>
using namespace std;

struct PointI {
  int64_t x, y;
  PointI(int64_t x = 0, int64_t y = 0) : x(x), y(y) {}
};

int main() {
  vector<PointI> square{{0, 0}, {4, 0}, {4, 4}, {0, 4}};
  assert(point_in_convex_polygon(square, PointI(2, 2)));
  assert(point_in_convex_polygon(square, PointI(4, 2)));
  assert(!point_in_convex_polygon(square, PointI(4, 2), false));
  assert(!point_in_convex_polygon(square, PointI(5, 2)));

  vector<PointI> pentagon{{0, 0}, {3, 0}, {5, 2}, {2, 5}, {-1, 3}};
  assert(point_in_convex_polygon(pentagon, PointI(2, 3)));
  assert(point_in_convex_polygon(pentagon, PointI(0, 0)));
  assert(!point_in_convex_polygon(pentagon, PointI(0, 0), false));
  assert(!point_in_convex_polygon(pentagon, PointI(5, 4)));

  assert(line_convex_polygon_intersection(square, PointI(-1, 2), PointI(5, 2)) == make_pair(3, 1));
  assert(line_convex_polygon_intersection(square, PointI(0, 0), PointI(4, 0)) == make_pair(0, 0));
  assert(line_convex_polygon_intersection(square, PointI(5, 0), PointI(5, 4)) == make_pair(-1, -1));

  // From the point to the right of the square, the upper-right and lower-right vertices are
  // tangent.
  assert(convex_polygon_tangents(square, PointI(6, 2)) == make_pair(2, 1));
  return 0;
}
\end{lstlisting}
\subsection{Minimum Enclosing Circle}
\label{sec:7.3.5}
Given a set of points in two dimensions, finds the unique circle of minimum radius (equivalently, minimum area) containing all points.

This implementation uses Welzl's randomized incremental algorithm. The points are shuffled and processed one at a time while maintaining the current minimum enclosing circle. Whenever a point lies outside the current circle, the circle is rebuilt with that point constrained to lie on the boundary. The final circle is determined by at most three boundary points.

\begin{compactitem}
  \item \inlinecode{min\_enclosing\_circle(p)} returns the minimum enclosing circle for the points in \inlinecode{p}. The point type may be any type exposing numeric \inlinecode{.x} and \inlinecode{.y} members. The returned circle always uses \inlinecode{double} coordinates, so integer-coordinate inputs are accepted.
\end{compactitem}

\textbf{Time Complexity}: $O(n)$ expected time per call, where $n$ is the number of points, or $O(n^{3})$ in the worst case for a particular shuffle order.

\textbf{Space Complexity}: $O(n)$ auxiliary for the shuffled copy of the points.

\begin{lstlisting}[style=cpp]
#include <algorithm>
#include <chrono>
#include <cmath>
#include <random>
#include <stdexcept>
#include <type_traits>
#include <vector>

const double EPS = 1e-9;

template<typename T, typename U, typename C = std::common_type_t<T, U>>
bool EQ(T a, U b) {
  if constexpr (std::is_floating_point_v<C>) return C(a) == C(b) || std::fabs(C(a) - C(b)) <= EPS;
  return C(a) == C(b);
}

template<typename T, typename U, typename C = std::common_type_t<T, U>>
bool LT(T a, U b) {
  if constexpr (std::is_floating_point_v<C>) return C(a) < C(b) - EPS;
  return C(a) < C(b);
}

template<typename T, typename U>
bool LE(T a, U b) {
  return !LT(b, a);
}

double sqnorm(double x, double y) {
  return x * x + y * y;
}

// SFINAE guard: valid only when Pt exposes numeric .x/.y members. This keeps the templated point
// constructors from hijacking calls like Circle(double, double, double).
template<typename Pt>
using if_point = decltype(std::declval<const Pt &>().x, std::declval<const Pt &>().y, void());

struct Circle {
  double h, k, r;

  Circle() : h(0), k(0), r(0) {}
  Circle(double h, double k, double r) : h(h), k(k), r(fabs(r)) {}

  template<typename Pt, typename = if_point<Pt>>
  Circle(const Pt &a, const Pt &b) {
    h = (a.x + b.x) / 2.0;
    k = (a.y + b.y) / 2.0;
    r = std::hypot(a.x - h, a.y - k);
  }

  template<typename Pt, typename = if_point<Pt>>
  Circle(const Pt &a, const Pt &b, const Pt &c) {
    double an = sqnorm(b.x - c.x, b.y - c.y);
    double bn = sqnorm(a.x - c.x, a.y - c.y);
    double cn = sqnorm(a.x - b.x, a.y - b.y);
    double wa = an * (bn + cn - an), wb = bn * (an + cn - bn), wc = cn * (an + bn - cn);
    double w = wa + wb + wc;
    if (EQ(w, 0)) {
      throw std::runtime_error("No circumcircle from collinear points.");
    }
    h = (wa * a.x + wb * b.x + wc * c.x) / w;
    k = (wa * a.y + wb * b.y + wc * c.y) / w;
    r = std::hypot(a.x - h, a.y - k);
  }

  template<typename Pt, typename = if_point<Pt>>
  bool contains(const Pt &p) const {
    return LE(sqnorm(p.x - h, p.y - k), r * r);
  }
};

// Input points can be any type with .x and .y fields; Circle output is always double.
template<typename Pt>
Circle min_enclosing_circle(std::vector<Pt> p) {
  if (p.empty()) {
    return Circle(0, 0, 0);
  }
  if (p.size() == 1) {
    return Circle(p[0].x, p[0].y, 0);
  }
  static std::mt19937 rng(std::chrono::steady_clock::now().time_since_epoch().count());
  std::shuffle(p.begin(), p.end(), rng);
  Circle res(p[0], p[1]);
  for (int i = 2; i < static_cast<int>(p.size()); i++) {
    if (res.contains(p[i])) {
      continue;
    }
    res = Circle(p[0], p[i]);
    for (int j = 1; j < i; j++) {
      if (res.contains(p[j])) {
        continue;
      }
      res = Circle(p[i], p[j]);
      for (int k = 0; k < j; k++) {
        if (!res.contains(p[k])) {
          res = Circle(p[i], p[j], p[k]);
        }
      }
    }
  }
  return res;
}
\end{lstlisting}
\Needspace*{4\baselineskip}
\textbf{Example Usage}\par
\begin{lstlisting}[style=cpp]
#include <cassert>
using namespace std;

struct Point {
  double x, y;
  Point(double x = 0, double y = 0) : x(x), y(y) {}
  bool operator==(const Point &p) const { return x == p.x && y == p.y; }
};

struct PointI {
  int x, y;
  PointI(int x = 0, int y = 0) : x(x), y(y) {}
};

int main() {
  vector<Point> empty;
  Circle none = min_enclosing_circle(empty);
  assert(EQ(none.h, 0) && EQ(none.k, 0) && EQ(none.r, 0));

  vector<Point> one{{2, 3}};
  Circle single = min_enclosing_circle(one);
  assert(EQ(single.h, 2) && EQ(single.k, 3) && EQ(single.r, 0));

  vector<Point> two{{0, 0}, {4, 0}};
  Circle diameter = min_enclosing_circle(two);
  assert(EQ(diameter.h, 2) && EQ(diameter.k, 0) && EQ(diameter.r, 2));

  vector<Point> v{{0, 0}, {0, 1}, {1, 0}, {1, 1}};
  Circle res = min_enclosing_circle(v);
  assert(EQ(res.h, 0.5) && EQ(res.k, 0.5) && EQ(res.r, 1 / sqrt(2)));

  // Integer-coordinate input: Circle output is always double.
  vector<PointI> iv{{0, 0}, {0, 2}, {2, 0}, {2, 2}};
  Circle ic = min_enclosing_circle(iv);
  assert(EQ(ic.h, 1.0) && EQ(ic.k, 1.0));
  return 0;
}
\end{lstlisting}
\subsection{Closest Pair}
\label{sec:7.3.6}
Given a list of points in two dimensions, finds the closest pair using a divide-and-conquer algorithm. The points are split in half by $x$-coordinate and each half is solved recursively; the combining step then only needs to examine points within the best distance so far of the dividing line, where each point in this strip is compared to a constant number of $y$-ordered neighbors.

\begin{compactitem}
  \item \inlinecode{closest\_pair(p, \&res)} returns the minimum squared distance between any two points in \inlinecode{p}. If \inlinecode{res} is non-null, one closest pair is stored there in lexicographic order. With fewer than two points, the maximum value of the squared-distance type is returned and \inlinecode{res} is unchanged. The point type may be any type exposing numeric \inlinecode{.x} and \inlinecode{.y} members and a lexicographic \inlinecode{operator\textless{}}. The distance preserves the coordinate arithmetic type, and the returned pair preserves the point type. For integer coordinates, the distance is exact provided intermediate products do not overflow.
\end{compactitem}

Overflow warning: squared distances grow like the square of the coordinate magnitude. For integer point types, use a 64-bit coordinate type (e.g. \inlinecode{PointL} from \secref{7.1.1}) when coordinates may exceed a few tens of thousands.

\textbf{Time Complexity}: $O(n \log n)$ per call, where $n$ is the number of points.

\Needspace*{3\baselineskip}
\textbf{Space Complexity}:
\begin{compactitem}
  \item $O(n)$ auxiliary heap space.
  \item $O(\log n)$ auxiliary stack space.
\end{compactitem}

\begin{lstlisting}[style=cpp]
#include <algorithm>
#include <iterator>
#include <limits>
#include <utility>
#include <vector>

template<typename Pt>
auto sqdist(const Pt &a, const Pt &b) {
  auto dx = b.x - a.x, dy = b.y - a.y;
  return dx * dx + dy * dy;  // Overflow warning.
}

template<typename It, typename T, typename Pt = typename std::iterator_traits<It>::value_type>
void closest_pair_rec(It lo, It hi, std::vector<Pt> &tmp, T &best, std::pair<Pt, Pt> *res) {
  int n = hi - lo;
  auto by_y = [](const Pt &a, const Pt &b) { return a.y != b.y ? a.y < b.y : a.x < b.x; };
  if (n <= 3) {
    for (It i = lo; i != hi; ++i) {
      for (It j = i + 1; j != hi; ++j) {
        T d = sqdist(*i, *j);
        if (d < best) {
          best = d;
          if (res != nullptr) {
            *res = std::minmax(*i, *j);
          }
        }
      }
    }
    std::sort(lo, hi, by_y);
    return;
  }
  It mid = lo + n / 2;
  auto midx = mid->x;
  closest_pair_rec(lo, mid, tmp, best, res);
  closest_pair_rec(mid, hi, tmp, best, res);
  // Each half is now y-sorted, so merge them before examining the center strip.
  tmp.clear();
  std::merge(lo, mid, mid, hi, std::back_inserter(tmp), by_y);
  std::move(tmp.begin(), tmp.end(), lo);
  tmp.clear();
  for (It it = lo; it != hi; ++it) {
    auto dx = it->x - midx;
    if (dx * dx < best) {
      tmp.push_back(*it);
    }
  }
  for (int i = 0; i < static_cast<int>(tmp.size()); i++) {
    for (int j = i + 1; j < static_cast<int>(tmp.size()); j++) {
      auto dy = tmp[j].y - tmp[i].y;
      if (dy * dy >= best) {
        break;
      }
      T d = sqdist(tmp[i], tmp[j]);
      if (d < best) {
        best = d;
        if (res != nullptr) {
          *res = std::minmax(tmp[i], tmp[j]);
        }
      }
    }
  }
  tmp.clear();
}

template<typename Pt>
auto closest_pair(std::vector<Pt> p, std::pair<Pt, Pt> *res = nullptr) {
  using T = decltype(sqdist(p[0], p[0]));
  T best = std::numeric_limits<T>::max();
  std::sort(p.begin(), p.end(), [](const Pt &a, const Pt &b) {
    return a.x != b.x ? a.x < b.x : a.y < b.y;
  });
  std::vector<Pt> tmp;
  tmp.reserve(p.size());
  closest_pair_rec(p.begin(), p.end(), tmp, best, res);
  return best;
}
\end{lstlisting}
\Needspace*{4\baselineskip}
\textbf{Example Usage}\par
\begin{lstlisting}[style=cpp]
#include <cassert>
#include <cmath>
#include <vector>
using namespace std;

bool EQ(double a, double b) {
  return fabs(a - b) < 1e-9;
}

struct Point {
  double x, y;
  Point(double x = 0, double y = 0) : x(x), y(y) {}
  bool operator==(const Point &p) const { return x == p.x && y == p.y; }
  bool operator<(const Point &p) const { return x != p.x ? x < p.x : y < p.y; }
};

struct PointI {
  int x, y;
  PointI(int x = 0, int y = 0) : x(x), y(y) {}
  bool operator==(const PointI &p) const { return x == p.x && y == p.y; }
  bool operator<(const PointI &o) const { return x != o.x ? x < o.x : y < o.y; }
};

int main() {
  vector<Point> v{{2, 3}, {12, 30}, {40, 50}, {5, 1}, {12, 10}, {3, 4}};
  pair<Point, Point> res;
  assert(EQ(closest_pair(v, &res), 2));
  auto [p1, p2] = res;
  assert(p1 == Point(2, 3) && p2 == Point(3, 4));

  // Integer-coordinate input: exact pair selection, int squared distance returned.
  vector<PointI> iv{{0, 0}, {10, 10}, {3, 4}};
  pair<PointI, PointI> ires;
  assert(closest_pair(iv, &ires) == 25);
  auto [i1, i2] = ires;
  assert(i1 == PointI(0, 0) && i2 == PointI(3, 4));
  return 0;
}
\end{lstlisting}
\subsection{Maximum Collinear Points}
\label{sec:7.3.7}
Given a set of points with integer coordinates, compute the maximum number of points that lie on one line. For every anchor point, normalize the direction vector to every other point by dividing by \inlinecode{gcd(abs(dx), abs(dy))} and choosing a canonical sign; equal normalized directions from the same anchor are collinear with it.

\begin{compactitem}
  \item \inlinecode{max\_collinear\_points(p)} returns the largest number of collinear points. Duplicate points are counted as separate input points.
\end{compactitem}

Overflow warning: coordinate differences and their absolute values must fit in \inlinecode{int64\_t}.

\textbf{Time Complexity}: $O(n^{2} \log n)$ per call, where $n$ is the number of points.

\textbf{Space Complexity}: $O(n)$ auxiliary.

\begin{lstlisting}[style=cpp]
#include <algorithm>
#include <cmath>
#include <cstdint>
#include <map>
#include <numeric>
#include <utility>
#include <vector>

struct Point {
  int64_t x, y;
  Point(int64_t x = 0, int64_t y = 0) : x(x), y(y) {}
  bool operator==(const Point &p) const { return x == p.x && y == p.y; }
};

int max_collinear_points(const std::vector<Point> &p) {
  int n = static_cast<int>(p.size()), ans = std::min(n, 1);
  for (int i = 0; i < n; i++) {
    std::map<std::pair<int64_t, int64_t>, int> cnt;
    int duplicates = 0, best = 0;
    for (int j = 0; j < n; j++) {
      if (i == j) {
        continue;
      }
      int64_t dx = p[j].x - p[i].x, dy = p[j].y - p[i].y;  // Overflow warning.
      if (dx == 0 && dy == 0) {
        duplicates++;
        continue;
      }
      int64_t g = std::gcd(std::abs(dx), std::abs(dy));
      dx /= g;
      dy /= g;
      if (dx < 0 || (dx == 0 && dy < 0)) {
        dx = -dx;
        dy = -dy;
      }
      best = std::max(best, ++cnt[{dx, dy}]);
    }
    ans = std::max(ans, best + duplicates + 1);
  }
  return ans;
}
\end{lstlisting}
\Needspace*{4\baselineskip}
\textbf{Example Usage}\par
\begin{lstlisting}[style=cpp]
#include <cassert>
using namespace std;

int main() {
  vector<Point> p{{0, 0}, {1, 1}, {2, 2}, {3, 5}, {4, 6}, {3, 3}};
  assert(max_collinear_points(p) == 4);
  p.emplace_back(1, 1);
  assert(max_collinear_points(p) == 5);
  assert(max_collinear_points({}) == 0);
  assert(max_collinear_points({Point(2, 3)}) == 1);
  return 0;
}
\end{lstlisting}
\subsection{Rectangle Coverage}
\label{sec:7.3.8}
Given weighted axis-aligned rectangles, iterate over every elementary cell induced by their coordinate-compressed boundaries. Within each such cell, the accumulated weight is constant, so the callback can compute rectangle union area, exactly-$k$ coverage area, threshold area, maximum overlap, weighted integrals, coverage histograms, or other offline coverage statistics.

Rectangles are half-open, $[x_1, x_2) \times [y_1, y_2)$, so adjacent rectangles sharing only an edge do not double-count any area. This also matches the difference-array update: add at the lower boundaries and subtract at the first coordinate after the rectangle.

\begin{compactitem}
  \item \inlinecode{for\_each\_rect\_cell(rects, visit)} calls \inlinecode{visit(x1, y1, x2, y2, weight)} once for each positive-area compressed cell inside the bounding box of the rectangles, where \inlinecode{weight} is the sum of weights covering that cell. Cells are visited in increasing \inlinecode{x1}, then increasing \inlinecode{y1}. Each rectangle must satisfy ${\inlinecodemath{x1}} < {\inlinecodemath{x2}}$ and ${\inlinecodemath{y1}} < {\inlinecodemath{y2}}$.
\end{compactitem}

Enumerating cells costs $O(n^{2})$ time and memory, which is the price of answering any statistic at all rather than one chosen in advance. When only the union area or perimeter is wanted, the sweep of section \secref{7.3.9} computes it in $O(n \log n)$ time and $O(n)$ memory, which is the only option once the rectangle count reaches the thousands.

Overflow warning: Coordinate differences, cell areas, and accumulated weights must fit in \inlinecode{int64\_t}.

\textbf{Time Complexity}: $O(n^{2} + n \log n)$ per call, where $n$ is the number of rectangles.

\textbf{Space Complexity}: $O(n^{2})$ auxiliary.

\begin{lstlisting}[style=cpp]
#include <algorithm>
#include <cassert>
#include <cstdint>
#include <vector>

struct Rect {
  int64_t x1, y1, x2, y2, w = 1;
};

template<typename Fn>
void for_each_rect_cell(const std::vector<Rect> &rects, Fn visit) {
  assert(std::all_of(rects.begin(), rects.end(), [](const Rect &r) {
    return r.x1 < r.x2 && r.y1 < r.y2;
  }));
  std::vector<int64_t> xs, ys;
  for (const Rect &r : rects) {
    xs.push_back(r.x1);
    xs.push_back(r.x2);
    ys.push_back(r.y1);
    ys.push_back(r.y2);
  }
  std::sort(xs.begin(), xs.end());
  std::sort(ys.begin(), ys.end());
  xs.erase(std::unique(xs.begin(), xs.end()), xs.end());
  ys.erase(std::unique(ys.begin(), ys.end()), ys.end());
  int X = static_cast<int>(xs.size()), Y = static_cast<int>(ys.size());
  std::vector<std::vector<int64_t>> diff(X + 1, std::vector<int64_t>(Y + 1));
  auto index = [](const std::vector<int64_t> &v, int64_t x) {
    return static_cast<int>(std::lower_bound(v.begin(), v.end(), x) - v.begin());
  };
  for (const Rect &r : rects) {
    int x1 = index(xs, r.x1), x2 = index(xs, r.x2);
    int y1 = index(ys, r.y1), y2 = index(ys, r.y2);
    diff[x1][y1] += r.w;  // Overflow warning.
    diff[x2][y1] -= r.w;
    diff[x1][y2] -= r.w;
    diff[x2][y2] += r.w;
  }
  for (int x = 0; x < X; x++) {
    for (int y = 0; y < Y; y++) {
      diff[x][y] += (x > 0) ? diff[x - 1][y] : 0;
      diff[x][y] += (y > 0) ? diff[x][y - 1] : 0;
      diff[x][y] -= (x > 0 && y > 0) ? diff[x - 1][y - 1] : 0;
      if (x + 1 < X && y + 1 < Y) {
        visit(xs[x], ys[y], xs[x + 1], ys[y + 1], diff[x][y]);
      }
    }
  }
}
\end{lstlisting}
\Needspace*{4\baselineskip}
\textbf{Example Usage}\par
\begin{lstlisting}[style=cpp]
#include <cassert>
using namespace std;

int main() {
  auto rect_area_if = [](const vector<Rect> &rects, auto pred) {
    int64_t area = 0;
    for_each_rect_cell(rects, [&](int64_t x1, int64_t y1, int64_t x2, int64_t y2, int64_t w) {
      if (pred(w)) {
        area += (x2 - x1) * (y2 - y1);  // Overflow warning.
      }
    });
    return area;
  };

  // Rectangles A=[0,2)x[0,2), B=[1,3)x[1,3), and C=[2,4)x[0,1) all with default weight 1.
  //
  // y=3 +---+---+---+---+
  //     |   | B | B |   |
  // y=2 +---+---+---+---+
  //     | A |A/B| B |   |
  // y=1 +---+---+---+---+
  //     | A | A | C | C |
  // y=0 +---+---+---+---+
  //   x=0   1   2   3   4
  vector<Rect> rects{{0, 0, 2, 2}, {1, 1, 3, 3}, {2, 0, 4, 1}};
  assert(rect_area_if(rects, [](int64_t w) { return w > 0; }) == 9);
  assert(rect_area_if(rects, [](int64_t w) { return w == 2; }) == 1);

  // Rectangles A=[0,4)x[0,3) of w=2, B=[2,6)x[1,4) of w=3, and C=[1,5)x[2,5) of w=1.
  //
  // y=5 +---+---+---+---+---+---+
  //     |   | 1 | 1 | 1 | 1 |   |
  // y=4 +---+---+---+---+---+---+
  //     |   | 1 | 4 | 4 | 4 | 3 |
  // y=3 +---+---+---+---+---+---+
  //     | 2 | 3 | 6 | 6 | 4 | 3 |
  // y=2 +---+---+---+---+---+---+
  //     | 2 | 2 | 5 | 5 | 3 | 3 |
  // y=1 +---+---+---+---+---+---+
  //     | 2 | 2 | 2 | 2 |   |   |
  // y=0 +---+---+---+---+---+---+
  //   x=0   1   2   3   4   5   6
  vector<Rect> weighted{{0, 0, 4, 3, 2}, {2, 1, 6, 4, 3}, {1, 2, 5, 5, 1}};
  assert(rect_area_if(weighted, [](int64_t w) { return w >= 3; }) == 13);
  assert(rect_area_if(weighted, [](int64_t w) { return w >= 5; }) == 4);

  int64_t max_weight = 0, integral = 0;
  for_each_rect_cell(weighted, [&](int64_t x1, int64_t y1, int64_t x2, int64_t y2, int64_t w) {
    max_weight = max(max_weight, w);
    integral += w * (x2 - x1) * (y2 - y1);  // Overflow warning.
  });
  assert(max_weight == 6);
  assert(integral == 12 * 2 + 12 * 3 + 12);
  return 0;
}
\end{lstlisting}
\subsection{Rectangle Union}
\label{sec:7.3.9}
Given axis-aligned rectangles that may overlap arbitrarily, measure their union. Klee's algorithm sweeps a vertical line left to right, maintaining how much of that line the rectangles currently cover. That covered length is constant between consecutive events, so each strip contributes the length times its width, and every other measure below is another aggregate of the same cross section.

A segment tree over the compressed $y$ coordinates holds the state, storing at each node the number of intervals covering its whole span and the length of that span which is covered. It never pushes anything down: a node's count refers only to intervals covering it entirely, so a node with a positive count reports its full span while one with a zero count defers to its children. Every addition is matched by a later removal, so counts never go negative and no lazy tag is needed.

Three further aggregates ride along on the same nodes. Doubly covered length shifts the recurrence by one, a node covered twice reporting its whole span and a node covered once reporting whatever its children cover at all. The weighted maximum is the largest root-to-leaf sum of weights, so each node reports its own weight plus the larger of its children's; left at the default weight of one per rectangle, that is simply the greatest number of rectangles covering a point. The perimeter needs the number of maximal covered runs, plus whether a node's first and last units are covered so that a run crossing a boundary is not counted twice: vertical edges are then the change in covered length at each event, and horizontal edges are twice the run count times the strip width.

Rectangles are half-open, $[x_1, x_2) \times [y_1, y_2)$, so adjacent rectangles sharing only an edge neither double-count any area nor leave a seam.

\begin{compactitem}
  \item \inlinecode{rect\_union(rects)} returns a \inlinecode{UnionMeasure} holding the \inlinecode{area} covered at least once, the \inlinecode{overlap\_area} covered at least twice, the \inlinecode{perimeter} of the outline including any enclosed hole, and \inlinecode{max\_weight}, the greatest weighted sum at any point. Subtracting the second from the first leaves the area covered exactly once. Each rectangle must satisfy ${\inlinecodemath{x1}} < {\inlinecodemath{x2}}$ and ${\inlinecodemath{y1}} < {\inlinecodemath{y2}}$.
\end{compactitem}

Weights change only the maximum, since the three geometric measures count coverage rather than payload, and weighted area needs no sweep at all: it is $\sum w_i$ times area, one loop over the input. Area covered at least $k$ times instead needs a vector of $k$ lengths per node; for that or any aggregate that does not decompose this way, the compression of section \secref{7.3.8} enumerates every elementary cell at $O(n^{2})$.

Overflow warning: Coordinate differences, areas, perimeters, and weight sums must fit in \inlinecode{int64\_t}.

\textbf{Time Complexity}: $O(n \log n)$ per call, where $n$ is the number of rectangles.

\textbf{Space Complexity}: $O(n)$ auxiliary.

\begin{lstlisting}[style=cpp]
#include <algorithm>
#include <cassert>
#include <cstdint>
#include <cstdlib>
#include <tuple>
#include <utility>
#include <vector>

struct Rect {
  int64_t x1, y1, x2, y2, w = 1;
};

struct UnionMeasure {
  int64_t area = 0, overlap_area = 0, perimeter = 0, max_weight = 0;
};

// The cross section of the sweep: over the currently covered y intervals, the length covered at
// least once and at least twice, the number of maximal covered runs, and the greatest weight at any
// point. Each is its own recurrence, so each is pulled up separately.
class CoverTree {
  static int checked_len(const std::vector<int64_t> &ys) {
    assert(ys.size() >= 2);
    return static_cast<int>(ys.size()) - 1;
  }

  std::vector<int64_t> ys;  // Compressed coordinates; leaf i spans [ys[i], ys[i + 1]).
  int len;
  std::vector<int> count, runs;
  std::vector<int64_t> covered, covered_twice, weight, heaviest;
  std::vector<char> covers_lo, covers_hi;

  // Length covered at least once and at least twice. The second reads the first one level down, so
  // covering at least k times extends this chain by one array per level.
  void pull_coverage(int i, int lo, int hi) {
    int left = 2 * i + 1, right = left + 1;
    if (count[i] > 0) {
      covered[i] = ys[hi + 1] - ys[lo];
      // One cover here turns anything already covered below into a double cover.
      covered_twice[i] =
          count[i] > 1 ? covered[i] : (lo == hi ? 0 : covered[left] + covered[right]);
    } else if (lo == hi) {
      covered[i] = covered_twice[i] = 0;
    } else {
      covered[i] = covered[left] + covered[right];
      covered_twice[i] = covered_twice[left] + covered_twice[right];
    }
  }

  // Number of maximal covered runs, with the flags that keep a run spanning a boundary from being
  // counted in both children.
  void pull_runs(int i, int lo, int hi) {
    int left = 2 * i + 1, right = left + 1;
    if (count[i] > 0) {
      runs[i] = covers_lo[i] = covers_hi[i] = 1;
    } else if (lo == hi) {
      runs[i] = covers_lo[i] = covers_hi[i] = 0;
    } else {
      runs[i] = runs[left] + runs[right] - (covers_hi[left] && covers_lo[right] ? 1 : 0);
      covers_lo[i] = covers_lo[left];
      covers_hi[i] = covers_hi[right];
    }
  }

  // The weight over an interval is the sum of the weights above it, so the heaviest point is the
  // largest root-to-leaf sum. This one ignores the cover count entirely.
  void pull_weight(int i, int lo, int hi) {
    heaviest[i] = weight[i] + (lo == hi ? 0 : std::max(heaviest[2 * i + 1], heaviest[2 * i + 2]));
  }

 public:
  explicit CoverTree(std::vector<int64_t> coords)
      : ys(std::move(coords)),
        len(checked_len(ys)),
        count(4 * len),
        runs(4 * len),
        covered(4 * len),
        covered_twice(4 * len),
        weight(4 * len),
        heaviest(4 * len),
        covers_lo(4 * len),
        covers_hi(4 * len) {}

  // Adds delta covers of weight w to every elementary interval in [lo, hi]; a removal negates both.
  void update(int lo, int hi, int delta, int64_t w) {
    auto rec = [&](auto &&rec, int i, int l, int h) {
      if (hi < l || h < lo) {
        return;
      }
      if (lo <= l && h <= hi) {
        count[i] += delta;
        weight[i] += delta * w;
      } else {
        int mid = l + (h - l) / 2;
        rec(rec, 2 * i + 1, l, mid);
        rec(rec, 2 * i + 2, mid + 1, h);
      }
      pull_coverage(i, l, h);
      pull_runs(i, l, h);
      pull_weight(i, l, h);
    };
    rec(rec, 0, 0, len - 1);
  }

  int64_t covered_length() const { return covered[0]; }
  int64_t twice_covered_length() const { return covered_twice[0]; }
  int run_count() const { return runs[0]; }
  int64_t heaviest_weight() const { return heaviest[0]; }
};

UnionMeasure rect_union(const std::vector<Rect> &rects) {
  if (rects.empty()) {
    return {};
  }
  assert(std::all_of(rects.begin(), rects.end(), [](const Rect &r) {
    return r.x1 < r.x2 && r.y1 < r.y2;
  }));
  std::vector<int64_t> ys;
  for (const Rect &r : rects) {
    ys.push_back(r.y1);
    ys.push_back(r.y2);
  }
  std::sort(ys.begin(), ys.end());
  ys.erase(std::unique(ys.begin(), ys.end()), ys.end());
  auto index = [&](int64_t y) {
    return static_cast<int>(std::lower_bound(ys.begin(), ys.end(), y) - ys.begin());
  };
  std::vector<std::tuple<int64_t, int, int, int, int64_t>> events;  // (x, lo, hi, delta, weight)
  for (const Rect &r : rects) {
    events.emplace_back(r.x1, index(r.y1), index(r.y2) - 1, 1, r.w);
    events.emplace_back(r.x2, index(r.y1), index(r.y2) - 1, -1, r.w);
  }
  std::sort(events.begin(), events.end());
  CoverTree tree(std::move(ys));
  UnionMeasure res;
  int64_t width = 0, doubled = 0, strips = 0, prev_x = 0;
  for (int i = 0, nevents = static_cast<int>(events.size()); i < nevents;) {
    int64_t x = std::get<0>(events[i]);
    if (i > 0) {  // The strip since the previous event has a constant cross-section.
      res.area += width * (x - prev_x);
      res.overlap_area += doubled * (x - prev_x);
      res.perimeter += 2 * strips * (x - prev_x);
    }
    int end = i;
    while (end < nevents && std::get<0>(events[end]) == x) {
      end++;
    }
    // Add before removing at the same x. The sum of the resulting length changes is the symmetric
    // difference of the old and new cross-sections: shared edges vanish, while disjoint entering
    // and leaving intervals are both counted.
    for (int wanted_delta : {1, -1}) {
      int64_t before = tree.covered_length();
      for (int j = i; j < end; j++) {
        auto [_, tgt_lo, tgt_hi, delta, w] = events[j];
        if (delta == wanted_delta) {
          tree.update(tgt_lo, tgt_hi, delta, w);
        }
      }
      res.perimeter += std::llabs(tree.covered_length() - before);
    }
    i = end;
    res.max_weight = std::max(res.max_weight, tree.heaviest_weight());
    width = tree.covered_length();
    doubled = tree.twice_covered_length();
    strips = tree.run_count();
    prev_x = x;
  }
  return res;
}
\end{lstlisting}
\Needspace*{4\baselineskip}
\textbf{Example Usage}\par
\begin{lstlisting}[style=cpp]
#include <cassert>
using namespace std;

int main() {
  assert(rect_union({}).area == 0 && rect_union({}).perimeter == 0);

  UnionMeasure one = rect_union({{0, 0, 2, 3}});
  assert(one.area == 6 && one.perimeter == 10);
  assert(one.overlap_area == 0 && one.max_weight == 1);

  // Disjoint rectangles simply add up, and nothing is covered twice.
  UnionMeasure apart = rect_union({{0, 0, 1, 1}, {2, 0, 3, 1}});
  assert(apart.area == 2 && apart.perimeter == 8);
  assert(apart.overlap_area == 0 && apart.max_weight == 1);

  // Disjoint intervals can leave and enter at the same x; both vertical edges still count.
  UnionMeasure crossing = rect_union({{0, 0, 1, 1}, {1, 2, 2, 3}});
  assert(crossing.area == 2 && crossing.perimeter == 8);

  // Two squares meeting in a unit square, whose union outlines an L-shaped hexagon.
  UnionMeasure pair = rect_union({{0, 0, 2, 2}, {1, 1, 3, 3}});
  assert(pair.area == 7 && pair.perimeter == 12);
  assert(pair.overlap_area == 1 && pair.max_weight == 2);
  assert(pair.area - pair.overlap_area == 6);  // Area belonging to exactly one square.

  // A rectangle nested inside another adds no area, but doubles the cover where it sits.
  UnionMeasure nested = rect_union({{0, 0, 10, 10}, {2, 2, 4, 4}});
  assert(nested.area == 100 && nested.perimeter == 40);
  assert(nested.overlap_area == 4 && nested.max_weight == 2);

  // Four rectangles enclosing a hole, overlapping only at the four corners.
  UnionMeasure ring = rect_union({{0, 0, 3, 1}, {0, 2, 3, 3}, {0, 0, 1, 3}, {2, 0, 3, 3}});
  assert(ring.area == 8);
  assert(ring.perimeter == 16);  // 12 around the outside and 4 around the unit hole.
  assert(ring.overlap_area == 4 && ring.max_weight == 2);

  // Three rectangles stacked over one cell.
  UnionMeasure stack = rect_union({{0, 0, 2, 2}, {1, 0, 3, 2}, {1, 1, 3, 3}});
  assert(stack.max_weight == 3 && stack.overlap_area == 3);
  return 0;
}
\end{lstlisting}
\subsection{Multiple Segment Intersection}
\label{sec:7.3.10}
Given a list of line segments in two dimensions, determine whether any pair of segments intersect using a sweep line algorithm. Endpoint events are processed from left to right while an ordered set maintains the segments currently crossing the sweep line, sorted by $y$; only segments that become adjacent in this set need to be tested, since any leftmost intersection must involve a pair that is adjacent just before it occurs. The input type \inlinecode{Segment\textless{}Pt\textgreater{}} is templated on the point type, so endpoints may be integer (\inlinecode{PointI}) or floating-point (\inlinecode{Point} or \inlinecode{PointD}), but must support \inlinecode{operator\textless{}} which orders points lexicographically. The cross-product sign tests in \inlinecode{seg\_intersection} are exact for integer endpoints, so intersection detection is exact.

This implementation counts endpoint contacts as intersections. See \secref{7.2.3} for pairwise intersection predicates when that distinction matters.

\begin{compactitem}
  \item \inlinecode{find\_intersecting\_pair(segs)} returns the indices of one pair of intersecting segments, or \inlinecode{std::nullopt} if none exist. Constructing a \inlinecode{Segment} places its lexicographically smaller endpoint first, as required by the sweep.
\end{compactitem}

Overflow warning: \inlinecode{seg\_intersection} forms the usual quadratic cross products, but the $y$-ordering cross-multiplication (\inlinecode{ay * bdx}) is cubic in the coordinate magnitude. Use \inlinecode{int64\_t} coordinates for modest integer inputs and a wider intermediate type when these cubic products may exceed it.

\textbf{Time Complexity}: $O(n \log n)$ per call, where $n$ is the number of segments.

\textbf{Space Complexity}: $O(n)$ auxiliary, where $n$ is the number of segments.

\begin{lstlisting}[style=cpp]
#include <algorithm>
#include <optional>
#include <set>
#include <tuple>
#include <utility>
#include <vector>

// Simplified detection-only version of seg_intersection() from 7.2.3 (exact for integral Pt).
template<typename Pt>
bool seg_intersection(const Pt &a, const Pt &b, const Pt &c, const Pt &d) {
  auto cross = [](const Pt &p, const Pt &q, const Pt &r) {
    return (q.x - p.x) * (r.y - p.y) - (q.y - p.y) * (r.x - p.x);
  };
  auto c1 = cross(a, b, c), c2 = cross(a, b, d), c3 = cross(c, d, a), c4 = cross(c, d, b);
  if (c1 == 0 && c2 == 0 && c3 == 0 && c4 == 0) {
    Pt p = std::max(std::min(a, b), std::min(c, d));
    Pt q = std::min(std::max(a, b), std::max(c, d));
    return !(q < p);
  }
  auto straddles = [](auto x, auto y) { return (x <= 0 && 0 <= y) || (y <= 0 && 0 <= x); };
  return straddles(c1, c2) && straddles(c3, c4);
}

// The point type Pt must support operator< (used to canonicalize endpoint order).
template<typename Pt>
struct Segment {
  Pt p, q;
  Segment(const Pt &p, const Pt &q) : p(std::min(p, q)), q(std::max(p, q)) {}
};

template<typename Pt>
bool intersects(const Segment<Pt> &s1, const Segment<Pt> &s2) {
  return seg_intersection(s1.p, s1.q, s2.p, s2.q);
}

template<typename Pt>
std::optional<std::pair<int, int>> find_intersecting_pair(const std::vector<Segment<Pt>> &segs) {
  struct Event {
    Pt p;
    bool is_end;
    int seg;
  };
  std::vector<Event> e;
  for (int i = 0; i < static_cast<int>(segs.size()); i++) {
    e.push_back(Event{segs[i].p, false, i});
    e.push_back(Event{segs[i].q, true, i});
  }
  std::sort(e.begin(), e.end(), [](const Event &a, const Event &b) {
    return std::tie(a.p.x, a.is_end, a.p.y, a.seg) < std::tie(b.p.x, b.is_end, b.p.y, b.seg);
  });
  // Compare y-values at sweep coordinate x without division: y = (p.y*dx + dy*(x - p.x)) / dx.
  // Vertical segments use their lower endpoint.
  auto ycmp = [](const auto &a, const auto &b, auto x) {
    // Overflow risk for integer Pt: the final cross-multiply is ~O(max_coord^3).
    auto adx = a.q.x - a.p.x, bdx = b.q.x - b.p.x;
    auto ay = a.p.y * adx + (a.q.y - a.p.y) * (x - a.p.x);
    auto by = b.p.y * bdx + (b.q.y - b.p.y) * (x - b.p.x);
    if (adx == 0) {
      ay = a.p.y;
      adx = 1;
    }
    if (bdx == 0) {
      by = b.p.y;
      bdx = 1;
    }
    auto lhs = ay * bdx, rhs = by * adx;
    return (lhs < rhs) ? -1 : ((rhs < lhs) ? 1 : 0);
  };
  auto cmp = [&segs, ycmp](int a, int b) {
    if (a == b) {
      return false;
    }
    auto x = std::max(segs[a].p.x, segs[b].p.x);
    int order = ycmp(segs[a], segs[b], x);
    return (order == 0) ? (a < b) : (order < 0);
  };
  using ActiveSet = std::set<int, decltype(cmp)>;
  ActiveSet s(cmp);
  std::vector<typename ActiveSet::iterator> position(segs.size());
  auto result = [](int i, int j) { return std::pair{std::min(i, j), std::max(i, j)}; };
  for (const auto &ev : e) {
    int seg = ev.seg;
    if (!ev.is_end) {
      auto it = s.insert(seg).first;
      position[seg] = it;
      auto next = it;
      if (++next != s.end() && intersects(segs[*next], segs[seg])) {
        return result(*next, seg);
      }
      if (it != s.begin() && intersects(segs[*--it], segs[seg])) {
        return result(*it, seg);
      }
    } else {
      auto it = position[seg];
      auto next = it;
      if (++next != s.end() && it != s.begin()) {
        auto prev = it;
        if (intersects(segs[*next], segs[*--prev])) {
          return result(*next, *prev);
        }
      }
      s.erase(it);
    }
  }
  return std::nullopt;
}
\end{lstlisting}
\Needspace*{4\baselineskip}
\textbf{Example Usage}\par
\begin{lstlisting}[style=cpp]
#include <cassert>
#include <vector>
using namespace std;

struct Point {
  double x, y;
  Point(double x = 0, double y = 0) : x(x), y(y) {}
  bool operator<(const Point &p) const { return x != p.x ? x < p.x : y < p.y; }
};

struct PointI {
  int x, y;
  PointI(int x = 0, int y = 0) : x(x), y(y) {}
  bool operator<(const PointI &p) const { return x != p.x ? x < p.x : y < p.y; }
};

int main() {
  {  // Double-coordinate segments.
    vector<Segment<Point>> v{
        Segment<Point>(Point(0, 0), Point(2, 2)), Segment<Point>(Point(3, 0), Point(0, -1)),
        Segment<Point>(Point(0, 2), Point(2, -2)), Segment<Point>(Point(0, 3), Point(9, 0))
    };
    assert((find_intersecting_pair(v) == pair{0, 2}));
  }
  {  // Integer-coordinate segments: detection is exact.
    vector<Segment<PointI>> v{
        Segment<PointI>({0, 0}, {2, 2}), Segment<PointI>({3, 0}, {0, -1}),
        Segment<PointI>({0, 2}, {2, -2}), Segment<PointI>({0, 3}, {9, 0})
    };
    assert((find_intersecting_pair(v) == pair{0, 2}));

    const vector<Segment<PointI>> disjoint{
        Segment<PointI>({0, 0}, {1, 0}), Segment<PointI>({0, 5}, {1, 5})
    };
    assert(!find_intersecting_pair(disjoint));
    const vector<Segment<PointI>> duplicate{
        Segment<PointI>({0, 0}, {2, 2}), Segment<PointI>({0, 0}, {2, 2})
    };
    assert((find_intersecting_pair(duplicate) == pair{0, 1}));
  }
  {  // Shared endpoints count as intersections.
    vector<Segment<PointI>> shared{
        Segment<PointI>({0, 0}, {2, 2}), Segment<PointI>({2, 2}, {4, 0})
    };
    assert((find_intersecting_pair(shared) == pair{0, 1}));
  }
  return 0;
}
\end{lstlisting}
\subsection{Manhattan MST}
\label{sec:7.3.11}
Given points in the plane, produce a small set of candidate edges guaranteed to contain a minimum spanning tree under Manhattan distance $|x_1 - x_2| + |y_1 - y_2|$. The plane is swept in four rotations/reflections; in each sweep, dominance by $x + y$ identifies the only nearby candidates that can matter. Run Kruskal on the returned edges to obtain the MST.

\begin{compactitem}
  \item \inlinecode{manhattan\_mst\_candidates(p)} returns $O(n)$ candidate edges as tuples (\inlinecode{weight}, \inlinecode{u}, \inlinecode{v}), where $n = {\inlinecodemath{p.size()}}$ and \inlinecode{u} and \inlinecode{v} are indices into \inlinecode{p}.
  \item \inlinecode{manhattan\_mst\_weight(p)} returns the MST weight by running Kruskal's algorithm on those candidate edges.
\end{compactitem}

Overflow warning: the sweeps form coordinate sums, negations, and differences such as \inlinecode{x + y} and \inlinecode{dx + dy}. With \inlinecode{int64\_t} coordinates, those intermediate values must fit in \inlinecode{int64\_t}.

\textbf{Time Complexity}: $O(n \log n)$ per call, where $n$ is the number of points.

\textbf{Space Complexity}: $O(n)$ auxiliary.

\begin{lstlisting}[style=cpp]
#include <algorithm>
#include <cstdint>
#include <cstdlib>
#include <map>
#include <numeric>
#include <tuple>
#include <utility>
#include <vector>

struct Point {
  int64_t x, y;
  Point(int64_t x = 0, int64_t y = 0) : x(x), y(y) {}
};

std::vector<std::tuple<int64_t, int, int>> manhattan_mst_candidates(std::vector<Point> p) {
  int n = static_cast<int>(p.size());
  std::vector<int> id(n);
  std::iota(id.begin(), id.end(), 0);
  std::vector<std::tuple<int64_t, int, int>> edges;
  for (int rot = 0; rot < 4; rot++) {
    std::sort(id.begin(), id.end(), [&](int i, int j) {
      return p[i].x + p[i].y < p[j].x + p[j].y;
    });
    std::map<int64_t, int> sweep;
    for (int i : id) {
      for (auto it = sweep.lower_bound(-p[i].y); it != sweep.end();) {
        int j = it->second;
        int64_t dx = p[i].x - p[j].x, dy = p[i].y - p[j].y;
        if (dy > dx) {
          break;
        }
        edges.emplace_back(dx + dy, i, j);
        it = sweep.erase(it);
      }
      sweep[-p[i].y] = i;
    }
    for (Point &q : p) {
      if (rot & 1) {
        q.x = -q.x;
      } else {
        std::swap(q.x, q.y);
      }
    }
  }
  return edges;
}

struct DSU {
  std::vector<int> root, rank;

  explicit DSU(int n) : root(n), rank(n, 0) { std::iota(root.begin(), root.end(), 0); }
  int find(int u) { return root[u] == u ? u : root[u] = find(root[u]); }

  bool unite(int u, int v) {
    u = find(u);
    v = find(v);
    if (u == v) {
      return false;
    }
    if (rank[u] < rank[v]) {
      std::swap(u, v);
    }
    root[v] = u;
    if (rank[u] == rank[v]) {
      rank[u]++;
    }
    return true;
  }
};

int64_t manhattan_mst_weight(const std::vector<Point> &p) {
  auto edges = manhattan_mst_candidates(p);
  std::sort(edges.begin(), edges.end());
  DSU dsu(static_cast<int>(p.size()));
  int64_t res = 0;
  for (auto [w, u, v] : edges) {
    if (dsu.unite(u, v)) {
      res += w;
    }
  }
  return res;
}
\end{lstlisting}
\Needspace*{4\baselineskip}
\textbf{Example Usage}\par
\begin{lstlisting}[style=cpp]
#include <cassert>
using namespace std;

int main() {
  assert(manhattan_mst_weight({Point(5, -7)}) == 0);
  assert(manhattan_mst_weight({Point(0, 0), Point(3, -4)}) == 7);

  vector<Point> p{{0, 0}, {2, 1}, {3, 4}, {-1, 2}, {5, 0}};
  assert(manhattan_mst_weight(p) == 14);
  auto edges = manhattan_mst_candidates(p);
  assert(edges.size() <= 4 * p.size());
  for (auto [w, u, v] : edges) {
    assert(0 <= u && u < static_cast<int>(p.size()));
    assert(0 <= v && v < static_cast<int>(p.size()));
    assert(w == abs(p[u].x - p[v].x) + abs(p[u].y - p[v].y));
  }
  return 0;
}
\end{lstlisting}

\section{\texorpdfstring{Advanced Planar Geometry}{7.4 Advanced Planar Geometry}}
\setcounter{section}{4}
\setcounter{subsection}{0}
\subsection{Convex Polygon Cut}
\label{sec:7.4.1}
Given a convex polygon and a directed line from \inlinecode{p} to \inlinecode{q}, clips the polygon to the closed left half-plane of that line. The polygon is walked edge by edge: vertices on the left side of the line are kept, and whenever an edge crosses the line, the intersection point is appended to the output.

\begin{compactitem}
  \item \inlinecode{convex\_cut(lo, hi, p, q)} returns the portion of the polygon lying on or to the left of the directed line from \inlinecode{p} to \inlinecode{q}. The input range $[{\inlinecodemath{lo}}, {\inlinecodemath{hi}})$ must contain the vertices of a convex polygon in boundary order, either clockwise or counterclockwise. The returned polygon preserves that boundary order and uses \inlinecode{Point} with \inlinecode{double} coordinates because edge-line intersections are computed in floating point. The points \inlinecode{p} and \inlinecode{q} must differ.
\end{compactitem}

Overflow warning: For integer-coordinate inputs, classification is exact only if the intermediate products do not overflow.

\textbf{Time Complexity}: $O(n)$ per call, where $n$ is the distance between \inlinecode{lo} and \inlinecode{hi}.

\textbf{Space Complexity}: $O(n)$ for the returned polygon and $O(1)$ auxiliary.

\begin{lstlisting}[style=cpp]
#include <cassert>
#include <cmath>
#include <type_traits>
#include <vector>

const double EPS = 1e-9;

template<typename T, typename U, typename C = std::common_type_t<T, U>>
bool EQ(T a, U b) {
  if constexpr (std::is_floating_point_v<C>) return C(a) == C(b) || std::fabs(C(a) - C(b)) <= EPS;
  return C(a) == C(b);
}

template<typename T, typename U, typename C = std::common_type_t<T, U>>
bool LT(T a, U b) {
  if constexpr (std::is_floating_point_v<C>) return C(a) < C(b) - EPS;
  return C(a) < C(b);
}

struct Point {
  double x, y;
  Point(double x = 0, double y = 0) : x(x), y(y) {}
  bool operator==(const Point &p) const { return x == p.x && y == p.y; }
  friend bool EQ(const Point &a, const Point &b) { return EQ(a.x, b.x) && EQ(a.y, b.y); }
};

template<typename PtA, typename PtB, typename PtO>
auto cross(const PtA &a, const PtB &b, const PtO &o) {
  return (a.x - o.x) * (b.y - o.y) - (a.y - o.y) * (b.x - o.x);
}

template<typename PtA, typename PtO, typename PtB>
int turn(const PtA &a, const PtO &o, const PtB &b) {
  auto c = cross(a, b, o);
  return LT(c, 0) ? 1 : (LT(0, c) ? -1 : 0);
}

template<typename PtA, typename PtB>
int line_intersection(
    const PtA &p1, const PtA &p2, const PtB &p3, const PtB &p4, Point *p = nullptr
) {
  double a1 = static_cast<double>(p2.y) - p1.y, b1 = static_cast<double>(p1.x) - p2.x;
  double c1 = -(static_cast<double>(p1.x) * p2.y - static_cast<double>(p2.x) * p1.y);
  double a2 = static_cast<double>(p4.y) - p3.y, b2 = static_cast<double>(p3.x) - p4.x;
  double c2 = -(static_cast<double>(p3.x) * p4.y - static_cast<double>(p4.x) * p3.y);
  double x = -(c1 * b2 - c2 * b1), y = -(a1 * c2 - a2 * c1);
  double det = a1 * b2 - a2 * b1;
  if (EQ(det, 0)) {
    return (EQ(x, 0) && EQ(y, 0)) ? 1 : -1;
  }
  if (p != nullptr) {
    *p = Point(x / det, y / det);
  }
  return 0;
}

template<typename It, typename Pt>
std::vector<Point> convex_cut(It lo, It hi, const Pt &p, const Pt &q) {
  assert(!EQ(p.x, q.x) || !EQ(p.y, q.y));
  if (lo == hi) {
    return {};
  }
  std::vector<Point> res;
  for (It i = lo, j = hi - 1; i != hi; j = i++) {
    Point pj(static_cast<double>(j->x), static_cast<double>(j->y));
    Point pi(static_cast<double>(i->x), static_cast<double>(i->y));
    int d1 = turn(q, p, *j), d2 = turn(q, p, *i);
    if (d1 <= 0) {
      res.push_back(pj);
    }
    if (d1 * d2 < 0) {
      Point r;
      line_intersection(p, q, *j, *i, &r);
      res.push_back(r);
    }
  }
  return res;
}
\end{lstlisting}
\Needspace*{4\baselineskip}
\textbf{Example Usage}\par
\begin{lstlisting}[style=cpp]
using namespace std;

struct PointI {
  int x, y;
  PointI(int x = 0, int y = 0) : x(x), y(y) {}
};

bool EQ(const vector<Point> &a, const vector<Point> &b) {
  return a.size() == b.size() &&
         equal(a.begin(), a.end(), b.begin(), [](const Point &p, const Point &q) {
           return EQ(p, q);
         });
}

int main() {
  {
    vector<Point> v{{1, 3}, {2, 2}, {2, 1}, {0, 0}, {-1, 3}};
    // Cut using the vertical line through (0, 0).
    vector<Point> c{{-1, 3}, {0, 3}, {0, 0}};
    assert(EQ(convex_cut(v.begin(), v.end(), Point(0, 0), Point(0, 1)), c));
  }
  {  // On a non-convex input, the result may be multiple disjoint polygons!
    vector<Point> v{{0, 0}, {2, 2}, {0, 4}, {3, 4}, {3, 0}};
    vector<Point> c{{1, 0}, {0, 0}, {1, 1}, {1, 3}, {0, 4}, {1, 4}};
    assert(EQ(convex_cut(v.begin(), v.end(), Point(1, 0), Point(1, 4)), c));
  }
  {
    vector<PointI> v{{0, 0}, {4, 0}, {4, 4}, {0, 4}};
    vector<Point> c{{0, 4}, {0, 0}, {2, 0}, {2, 4}};
    assert(EQ(convex_cut(v.begin(), v.end(), PointI(2, 0), PointI(2, 4)), c));
  }
  return 0;
}
\end{lstlisting}
\subsection{Half-Plane Intersection}
\label{sec:7.4.2}
Given a list of directed lines, compute the convex polygon contained in every left half-plane. Half-plane intersection sorts the lines by direction, keeps only the tightest representative among parallel lines, and maintains a deque of lines whose pairwise intersections form the current boundary. Directions are ordered by half-plane and cross-product signs, without epsilon comparisons, to satisfy \inlinecode{std::sort}'s strict ordering requirement. \inlinecode{EPS} is then used outside the comparator to group parallel directions and retain only the tightest half-plane in each group.

\begin{compactitem}
  \item \inlinecode{half\_plane\_intersection(planes)} returns the polygon cut out by the half-planes in counter-clockwise order. Each half-plane is represented by \inlinecode{HalfPlane(p, q)}, meaning the closed region to the left of the directed line \inlinecode{p} $\to$ \inlinecode{q}. This implementation is intended for intersections that form a bounded polygon with positive area; empty, unbounded, or degenerate inputs may return an empty vector. Add explicit bounding-box half-planes when a bounded polygon is required. The two points defining each half-plane must differ.
  \item \inlinecode{out(q)} returns whether point \inlinecode{q} lies strictly outside this half-plane, that is, to the right of its directed line.
  \item \inlinecode{intersect(h)} returns the point where the boundary lines of this half-plane and \inlinecode{h} meet, which must not be parallel.
\end{compactitem}

\textbf{Time Complexity}: $O(n \log n)$ per call, where $n$ is the number of half-planes.

\textbf{Space Complexity}: $O(n)$ auxiliary.

\begin{lstlisting}[style=cpp]
#include <algorithm>
#include <cmath>
#include <deque>
#include <type_traits>
#include <vector>

const double EPS = 1e-9;

template<typename T, typename U, typename C = std::common_type_t<T, U>>
bool EQ(T a, U b) {
  if constexpr (std::is_floating_point_v<C>) return C(a) == C(b) || std::fabs(C(a) - C(b)) <= EPS;
  return C(a) == C(b);
}

template<typename T, typename U, typename C = std::common_type_t<T, U>>
bool LT(T a, U b) {
  if constexpr (std::is_floating_point_v<C>) return C(a) < C(b) - EPS;
  return C(a) < C(b);
}

struct Point {
  double x, y;
  Point(double x = 0, double y = 0) : x(x), y(y) {}
  bool operator==(const Point &p) const { return x == p.x && y == p.y; }
  friend bool EQ(const Point &a, const Point &b) { return EQ(a.x, b.x) && EQ(a.y, b.y); }
  Point operator+(const Point &p) const { return {x + p.x, y + p.y}; }
  Point operator-(const Point &p) const { return {x - p.x, y - p.y}; }
  Point operator*(double k) const { return {x * k, y * k}; }
  double cross(const Point &p) const { return x * p.y - y * p.x; }
};

struct HalfPlane {
  Point p, dir;

  HalfPlane() = default;
  HalfPlane(Point p, Point q) : p(p), dir(q - p) {}

  bool out(const Point &q) const { return LT(dir.cross(q - p), 0); }

  Point intersect(const HalfPlane &h) const {
    double t = h.dir.cross(p - h.p) / dir.cross(h.dir);
    return p + dir * t;
  }

  bool operator<(const HalfPlane &h) const {
    auto half = [](const Point &d) { return d.y < 0 || (d.y == 0 && d.x < 0); };
    if (half(dir) != half(h.dir)) {
      return half(dir) < half(h.dir);
    }
    return dir.cross(h.dir) > 0;
  }
};

std::vector<Point> half_plane_intersection(std::vector<HalfPlane> planes) {
  std::sort(planes.begin(), planes.end());
  std::vector<HalfPlane> unique;
  for (const HalfPlane &h : planes) {
    if (!unique.empty() && EQ(unique.back().dir.cross(h.dir), 0)) {
      if (unique.back().dir.cross(h.p - unique.back().p) > 0) {
        unique.back() = h;
      }
      continue;
    }
    unique.push_back(h);
  }
  std::deque<HalfPlane> dq;
  auto bad_back = [&](const HalfPlane &h) {
    return dq.size() >= 2 && h.out(dq[dq.size() - 2].intersect(dq.back()));
  };
  auto bad_front = [&](const HalfPlane &h) {
    return dq.size() >= 2 && h.out(dq[0].intersect(dq[1]));
  };
  for (const HalfPlane &h : unique) {
    while (bad_back(h)) {
      dq.pop_back();
    }
    while (bad_front(h)) {
      dq.pop_front();
    }
    dq.push_back(h);
  }
  while (dq.size() >= 3 && dq.front().out(dq[dq.size() - 2].intersect(dq.back()))) {
    dq.pop_back();
  }
  while (dq.size() >= 3 && dq.back().out(dq[0].intersect(dq[1]))) {
    dq.pop_front();
  }
  if (dq.size() < 3) {
    return {};
  }
  std::vector<Point> res;
  for (int i = 0; i < static_cast<int>(dq.size()); i++) {
    if (EQ(dq[i].dir.cross(dq[(i + 1) % dq.size()].dir), 0)) {
      return {};
    }
    res.push_back(dq[i].intersect(dq[(i + 1) % dq.size()]));
  }
  return res;
}
\end{lstlisting}
\Needspace*{4\baselineskip}
\textbf{Example Usage}\par
\begin{lstlisting}[style=cpp]
#include <cassert>
using namespace std;

double polygon_area(const vector<Point> &p) {
  double area = 0;
  for (int i = 0, j = static_cast<int>(p.size()) - 1; i < static_cast<int>(p.size()); j = i++) {
    area += p[j].cross(p[i]);
  }
  return fabs(area) / 2.0;
}

int main() {
  vector<HalfPlane> box{
      HalfPlane(Point(0, 0), Point(4, 0)),
      HalfPlane(Point(4, 0), Point(4, 4)),
      HalfPlane(Point(4, 4), Point(0, 4)),
      HalfPlane(Point(0, 4), Point(0, 0)),
  };
  auto square = half_plane_intersection(box);
  assert(square.size() == 4 && EQ(polygon_area(square), 16));

  auto tighter = box;
  tighter.emplace_back(Point(0, 1), Point(4, 1));  // y >= 1
  assert(EQ(polygon_area(half_plane_intersection(tighter)), 12));

  box.emplace_back(Point(2, 0), Point(2, 4));  // x <= 2
  auto rect = half_plane_intersection(box);
  assert(rect.size() == 4 && EQ(polygon_area(rect), 8));

  vector<HalfPlane> empty{
      HalfPlane(Point(0, 0), Point(1, 0)),  // y >= 0
      HalfPlane(Point(1, 1), Point(0, 1)),  // y <= 1
      HalfPlane(Point(0, 2), Point(1, 2)),  // y >= 2
  };
  assert(half_plane_intersection(empty).empty());
  return 0;
}
\end{lstlisting}
\subsection{Minkowski Sum}
\label{sec:7.4.3}
Given two convex polygons, compute their Minkowski sum: the set of all points $a + b$ where $a$ comes from the first polygon and $b$ comes from the second. After rotating both polygons to start at their lowest-leftmost vertex, the boundary edges of the sum are obtained by merging the two cyclic edge-angle lists. Minkowski sums are useful for collision/translation problems, convex polygon distance checks, and diameter computations via $P + (-P)$.

\begin{compactitem}
  \item \inlinecode{minkowski\_sum(a, b)} returns the convex polygon representing the sum. Inputs must be convex polygons in boundary order, clockwise or counter-clockwise, with no repeated final vertex. Collinear consecutive output vertices are removed.
\end{compactitem}

\textbf{Time Complexity}: $O(n + m)$ per call, where $n$ and $m$ are the polygon sizes.

\textbf{Space Complexity}: $O(n + m)$ auxiliary.

\begin{lstlisting}[style=cpp]
#include <algorithm>
#include <cmath>
#include <type_traits>
#include <utility>
#include <vector>

const double EPS = 1e-9;

template<typename T, typename U, typename C = std::common_type_t<T, U>>
bool EQ(T a, U b) {
  if constexpr (std::is_floating_point_v<C>) return C(a) == C(b) || std::fabs(C(a) - C(b)) <= EPS;
  return C(a) == C(b);
}

struct Point {
  double x, y;
  Point(double x = 0, double y = 0) : x(x), y(y) {}
  bool operator==(const Point &p) const { return x == p.x && y == p.y; }
  friend bool EQ(const Point &a, const Point &b) { return EQ(a.x, b.x) && EQ(a.y, b.y); }
  Point operator+(const Point &p) const { return {x + p.x, y + p.y}; }
  Point operator-(const Point &p) const { return {x - p.x, y - p.y}; }
  double cross(const Point &p) const { return x * p.y - y * p.x; }
};

std::vector<Point> normalize(std::vector<Point> p) {
  while (p.size() > 1 && EQ(p.front(), p.back())) {
    p.pop_back();
  }
  if (p.size() <= 1) {
    return p;
  }
  double area = 0;
  for (int i = 0, j = static_cast<int>(p.size()) - 1; i < static_cast<int>(p.size()); j = i++) {
    area += p[j].cross(p[i]);
  }
  if (area < 0) {
    std::reverse(p.begin(), p.end());
  }
  int start = 0;
  for (int i = 1; i < static_cast<int>(p.size()); i++) {
    if (p[i].y < p[start].y || (EQ(p[i].y, p[start].y) && p[i].x < p[start].x)) {
      start = i;
    }
  }
  std::rotate(p.begin(), p.begin() + start, p.end());
  return p;
}

std::vector<Point> remove_collinear(const std::vector<Point> &p) {
  std::vector<Point> res;
  for (const Point &q : p) {
    while (res.size() >= 2 && EQ((res.back() - res[res.size() - 2]).cross(q - res.back()), 0)) {
      res.pop_back();
    }
    if (res.empty() || !EQ(q, res.back())) {
      res.push_back(q);
    }
  }
  while (res.size() >= 3 && EQ((res.back() - res[res.size() - 2]).cross(res[0] - res.back()), 0)) {
    res.pop_back();
  }
  while (res.size() >= 3 && EQ((res[0] - res.back()).cross(res[1] - res[0]), 0)) {
    res.erase(res.begin());
  }
  return res;
}

std::vector<Point> minkowski_sum(std::vector<Point> a, std::vector<Point> b) {
  a = normalize(std::move(a));
  b = normalize(std::move(b));
  if (a.empty() || b.empty()) {
    return {};
  }
  if (a.size() == 1) {
    for (Point &p : b) {
      p = p + a[0];
    }
    return b;
  }
  if (b.size() == 1) {
    for (Point &p : a) {
      p = p + b[0];
    }
    return a;
  }
  int n = static_cast<int>(a.size()), m = static_cast<int>(b.size());
  a.push_back(a[0]);
  a.push_back(a[1]);
  b.push_back(b[0]);
  b.push_back(b[1]);
  std::vector<Point> res;
  int i = 0, j = 0;
  while (i < n || j < m) {
    res.push_back(a[i] + b[j]);
    if (i == n) {
      j++;
      continue;
    }
    if (j == m) {
      i++;
      continue;
    }
    double cr = (a[i + 1] - a[i]).cross(b[j + 1] - b[j]);
    if (cr > -EPS) {
      i++;
    }
    if (cr < EPS) {
      j++;
    }
  }
  return remove_collinear(res);
}
\end{lstlisting}
\Needspace*{4\baselineskip}
\textbf{Example Usage}\par
\begin{lstlisting}[style=cpp]
#include <cassert>
using namespace std;

int main() {
  vector<Point> empty;
  assert(minkowski_sum(empty, {Point(1, 2)}).empty());

  vector<Point> square{{0, 0}, {1, 0}, {1, 1}, {0, 1}};
  vector<Point> tri{{0, 0}, {2, 0}, {0, 2}};
  // The square thickens the triangle by one unit in the positive x/y directions.
  assert((minkowski_sum(square, tri) == vector<Point>{{0, 0}, {3, 0}, {3, 1}, {1, 3}, {0, 3}}));

  // Adding a single point translates the polygon.
  assert((minkowski_sum(square, {Point(2, 3)}) == vector<Point>{{2, 3}, {3, 3}, {3, 4}, {2, 4}}));
  assert((minkowski_sum({Point(1, 2)}, {Point(3, 4)}) == vector<Point>{{4, 6}}));
  return 0;
}
\end{lstlisting}
\subsection{Polygon Intersection and Union}
\label{sec:7.4.4}
Given two simple (non-self-intersecting) polygons, determine the areas of their intersection and union using a sweep line algorithm and the inclusion-exclusion principle. Every vertex and every pairwise edge intersection contributes its $x$-coordinate as a sweep boundary, so within each vertical slab between consecutive boundaries the two borders cannot cross, and the slab's overlap is accumulated exactly from trapezoid areas. The union then follows by inclusion-exclusion as the sum of the individual areas minus the intersection.

\begin{compactitem}
  \item \inlinecode{intersection\_area(a, b)} returns the intersection area of polygons \inlinecode{a} and \inlinecode{b}, whose vertices are listed in boundary order.
  \item \inlinecode{union\_area(a, b)} returns the union area of polygons \inlinecode{a} and \inlinecode{b}.
\end{compactitem}

Overflow warning: the segment-intersection tests form cross products in the input point's coordinate type, which grow like the squared coordinate magnitude. For integral point types, use 64-bit coordinates (e.g. \inlinecode{PointL} from \secref{7.1.1}) if necessary.

\textbf{Time Complexity}: $O(n \cdot m \cdot (n + m) \cdot \log\left(n + m\right))$ per call to \inlinecode{intersection\_area()} and \inlinecode{union\_area()}, where $n$ is the number of vertices in the first polygon and $m$ is the number of vertices in the second polygon (or $O(N^{3} \log N)$ for $N = n + m$).

\textbf{Space Complexity}: $O(n \cdot m)$ auxiliary for \inlinecode{intersection\_area()} and \inlinecode{union\_area()}, where $n$ and $m$ are the respective polygon sizes.

\begin{lstlisting}[style=cpp]
#include <algorithm>
#include <cmath>
#include <type_traits>
#include <utility>
#include <vector>

const double EPS = 1e-9;

template<typename T, typename U, typename C = std::common_type_t<T, U>>
bool EQ(T a, U b) {
  if constexpr (std::is_floating_point_v<C>) return C(a) == C(b) || std::fabs(C(a) - C(b)) <= EPS;
  return C(a) == C(b);
}

template<typename T, typename U, typename C = std::common_type_t<T, U>>
bool LT(T a, U b) {
  if constexpr (std::is_floating_point_v<C>) return C(a) < C(b) - EPS;
  return C(a) < C(b);
}

template<typename T, typename U>
bool LE(T a, U b) {
  return !LT(b, a);
}

template<typename Pt>
bool point_on_segment(const Pt &p, const Pt &a, const Pt &b) {
  // Overflow risk for integer Pt: these products are ~O(max_coord^2); use int64_t if necessary.
  return EQ((p.x - a.x) * (b.y - a.y) - (p.y - a.y) * (b.x - a.x), 0) &&
         LE(std::min(a.x, b.x), p.x) && LE(p.x, std::max(a.x, b.x)) &&
         LE(std::min(a.y, b.y), p.y) && LE(p.y, std::max(a.y, b.y));
}

// Specialized version of seg_intersection() from 7.2.3, simplified for include_boundary = true,
// returning -1 for no intersection, 0 for one intersection point, 1 for positive length overlap.
template<typename Pt>
int seg_intersection1(
    const Pt &a, const Pt &b, const Pt &c, const Pt &d, double *outx, double *outy
) {
  auto ab_x = b.x - a.x, ab_y = b.y - a.y;
  auto ac_x = c.x - a.x, ac_y = c.y - a.y;
  auto cd_x = d.x - c.x, cd_y = d.y - c.y;
  auto ab2 = ab_x * ab_x + ab_y * ab_y;
  auto cd2 = cd_x * cd_x + cd_y * cd_y;
  if (EQ(ab2, 0)) {
    if (point_on_segment(a, c, d)) {
      *outx = a.x;
      *outy = a.y;
      return 0;
    }
    return -1;
  }
  if (EQ(cd2, 0)) {
    if (point_on_segment(c, a, b)) {
      *outx = c.x;
      *outy = c.y;
      return 0;
    }
    return -1;
  }
  auto c1 = ab_x * cd_y - ab_y * cd_x;
  auto c2 = ac_x * ab_y - ac_y * ab_x;
  if (EQ(c1, 0) && EQ(c2, 0)) {  // Collinear.
    Pt res1 = std::max(std::min(a, b), std::min(c, d));
    Pt res2 = std::min(std::max(a, b), std::max(c, d));
    if (!(res2 < res1)) {
      if (res1 == res2) {
        *outx = static_cast<double>(res1.x);
        *outy = static_cast<double>(res1.y);
        return 0;
      }
      return 1;
    }
    return -1;
  }
  if (EQ(c1, 0)) {
    return -1;
  }
  auto t_num = ac_x * cd_y - ac_y * cd_x;
  bool c1_pos = c1 > 0;
  bool t_ok = c1_pos ? (LE(0, t_num) && LE(t_num, c1)) : (LE(c1, t_num) && LE(t_num, 0));
  bool u_ok = c1_pos ? (LE(0, c2) && LE(c2, c1)) : (LE(c1, c2) && LE(c2, 0));
  if (t_ok && u_ok) {
    double t = static_cast<double>(t_num) / c1;
    *outx = static_cast<double>(a.x + t * ab_x);
    *outy = static_cast<double>(a.y + t * ab_y);
    return 0;
  }
  return -1;
}

// The two segments may use different point types (e.g. polygon edge vs. sweep line).
template<typename PtA, typename PtB>
int line_intersection1(
    const PtA &p1, const PtA &p2, const PtB &p3, const PtB &p4, double *outx, double *outy
) {
  double a1 = static_cast<double>(p2.y) - p1.y, b1 = static_cast<double>(p1.x) - p2.x;
  double c1 = -(a1 * p1.x + b1 * p1.y);
  double a2 = static_cast<double>(p4.y) - p3.y, b2 = static_cast<double>(p3.x) - p4.x;
  double c2 = -(a2 * p3.x + b2 * p3.y);
  double x = -(c1 * b2 - c2 * b1), y = -(a1 * c2 - a2 * c1);
  double det = a1 * b2 - a2 * b1;
  if (EQ(det, 0)) {
    return (EQ(x, 0) && EQ(y, 0)) ? 1 : -1;
  }
  if (outx != nullptr && outy != nullptr) {
    *outx = x / det;
    *outy = y / det;
  }
  return 0;
}

template<typename Pt>
double intersection_area(const std::vector<Pt> &a, const std::vector<Pt> &b) {
  if (a.empty() || b.empty()) {
    return 0;
  }
  struct SweepPoint {
    double x, y;
  };  // For line intersection with the sweep line.
  const std::vector<Pt> *polys[2] = {&a, &b};
  int orientation[2];
  for (int id = 0; id < 2; id++) {
    const auto &p = *polys[id];
    int n = static_cast<int>(p.size());
    double area = 0;
    for (int i = 0, j = n - 1; i < n; j = i++) {
      area += (static_cast<double>(p[j].x) - p[i].x) * (static_cast<double>(p[j].y) + p[i].y);
    }
    orientation[id] = (area < 0 ? 1 : (area > 0 ? -1 : 0));
  }
  std::vector<double> x_coords;
  x_coords.reserve(a.size() + b.size());
  for (const Pt &p : a) {
    x_coords.push_back(p.x);
  }
  for (const Pt &p : b) {
    x_coords.push_back(p.x);
  }
  int n = static_cast<int>(a.size()), m = static_cast<int>(b.size());
  for (int i1 = 0, j1 = n - 1; i1 < n; j1 = i1++) {
    for (int i2 = 0, j2 = m - 1; i2 < m; j2 = i2++) {
      double outx, outy;
      if (seg_intersection1(a[i1], a[j1], b[i2], b[j2], &outx, &outy) == 0) {
        x_coords.push_back(outx);
      }
    }
  }
  std::sort(x_coords.begin(), x_coords.end());
  x_coords.erase(
      std::unique(x_coords.begin(), x_coords.end(), [](double a, double b) { return EQ(a, b); }),
      x_coords.end()
  );
  double res = 0;
  for (int k = 0; k + 1 < static_cast<int>(x_coords.size()); k++) {
    double x = (x_coords[k] + x_coords[k + 1]) / 2;
    SweepPoint sweep0{x, 0}, sweep1{x, 1};
    std::vector<std::pair<double, int>> events;  // (y, mask delta)
    for (int id = 0; id < 2; id++) {
      const auto &p = *polys[id];
      int size = static_cast<int>(p.size());
      for (int j = 0, i = size - 1; j < size; i = j++) {
        double px, py;
        if (line_intersection1(p[j], p[i], sweep0, sweep1, &px, &py) == 0) {
          double y = py, x0 = p[i].x, x1 = p[j].x;
          if (x0 < x && x1 > x) {
            events.emplace_back(y, orientation[id] * (1 << id));
          } else if (x0 > x && x1 < x) {
            events.emplace_back(y, -orientation[id] * (1 << id));
          }
        }
      }
    }
    std::sort(events.begin(), events.end());
    double height = 0;
    int coverage[2] = {0, 0};
    for (int j = 0; j < static_cast<int>(events.size()); j++) {
      auto [y, mask_delta] = events[j];
      if (coverage[0] != 0 && coverage[1] != 0) {
        height += y - events[j - 1].first;
      }
      int bit = std::abs(mask_delta);
      int poly = (bit == 1 ? 0 : 1);
      coverage[poly] += mask_delta / bit;
    }
    res += height * (x_coords[k + 1] - x_coords[k]);
  }
  return res;
}

template<typename Pt>
double polygon_area(const std::vector<Pt> &p) {
  if (p.empty()) {
    return 0;
  }
  int n = static_cast<int>(p.size());
  double area = 0;
  for (int i = 0, j = n - 1; i < n; j = i++) {
    area += (static_cast<double>(p[j].x) - p[i].x) * (static_cast<double>(p[j].y) + p[i].y);
  }
  return fabs(area / 2.0);
}

template<typename Pt>
double union_area(const std::vector<Pt> &a, const std::vector<Pt> &b) {
  return polygon_area(a) + polygon_area(b) - intersection_area(a, b);
}
\end{lstlisting}
\Needspace*{4\baselineskip}
\textbf{Example Usage}\par
\begin{lstlisting}[style=cpp]
#include <cassert>
using namespace std;

struct Point {
  double x, y;
  Point(double x = 0, double y = 0) : x(x), y(y) {}
  bool operator==(const Point &p) const { return x == p.x && y == p.y; }
  bool operator!=(const Point &p) const { return !(*this == p); }
  bool operator<(const Point &p) const { return x != p.x ? x < p.x : y < p.y; }
};

struct PointI {
  int x, y;
  PointI(int x = 0, int y = 0) : x(x), y(y) {}
  bool operator==(const PointI &o) const { return x == o.x && y == o.y; }
  bool operator!=(const PointI &p) const { return !(*this == p); }
  bool operator<(const PointI &p) const { return x != p.x ? x < p.x : y < p.y; }
};

int main() {
  vector<Point> p, s;
  // Irregular pentagon overlapping the square below in a triangle of area 1.5.
  p.emplace_back(1, 3);
  p.emplace_back(1, 2);
  p.emplace_back(2, 1);
  p.emplace_back(0, 0);
  p.emplace_back(-1, 3);
  // Square of area 12.5 in quadrant 2.
  s.emplace_back(0, 0);
  s.emplace_back(0, 3);
  s.emplace_back(-3, 3);
  s.emplace_back(-3, 0);
  assert(EQ(1.5, intersection_area(p, s)));
  assert(EQ(12.5, union_area(p, s)));

  // Integer-coordinate polygons are accepted (computation proceeds in double).
  vector<PointI> ip{{1, 3}, {1, 2}, {2, 1}, {0, 0}, {-1, 3}};
  vector<PointI> is{{0, 0}, {0, 3}, {-3, 3}, {-3, 0}};
  assert(EQ(1.5, intersection_area(ip, is)));
  assert(EQ(12.5, union_area(ip, is)));
  return 0;
}
\end{lstlisting}
\subsection{Delaunay Triangulation (Simple)}
\label{sec:7.4.5}
Given a set $P$ of distinct two-dimensional points, a Delaunay triangulation of $P$ is a triangulation of the convex hull of $P$ such that no point of $P$ lies strictly inside the circumcircle of any triangle in the triangulation. The Delaunay triangulation is not necessarily unique when four or more points are cocircular.

This implementation produces one valid triangulation using a simple brute-force algorithm. Each candidate triangle is tested against the empty-circumcircle condition, and candidates whose edges would properly cross already accepted triangles are rejected.

All arithmetic uses the point's own coordinate type, so integer inputs yield an exact triangulation.

\begin{compactitem}
  \item \inlinecode{delaunay\_triangulation(p)} returns a Delaunay triangulation of the distinct points in \inlinecode{p} as a vector of clockwise-oriented vertex triples, or an empty vector if no triangulation exists (e.g. fewer than three points, or all points collinear).
\end{compactitem}

Overflow warning: the lifted-paraboloid test grows like the fourth power of the coordinate magnitude. With 64-bit integer coordinates, it overflows once coordinates exceed roughly $17600$. For floating-point coordinates the test is subject to the usual rounding error.

\textbf{Time Complexity}: $O(n^{6})$ per call, where $n$ is the number of points. There are $O(n^{3})$ candidate triangles; each scans all points and, conservatively, up to $O(n^{3})$ previously accepted candidates.

\textbf{Space Complexity}: $O(n)$ auxiliary, excluding the returned triangulation.

\begin{lstlisting}[style=cpp]
#include <cmath>
#include <tuple>
#include <type_traits>
#include <vector>

const double EPS = 1e-9;

template<typename T, typename U, typename C = std::common_type_t<T, U>>
bool EQ(T a, U b) {
  if constexpr (std::is_floating_point_v<C>) return C(a) == C(b) || std::fabs(C(a) - C(b)) <= EPS;
  return C(a) == C(b);
}

template<typename T, typename U, typename C = std::common_type_t<T, U>>
bool LT(T a, U b) {
  if constexpr (std::is_floating_point_v<C>) return C(a) < C(b) - EPS;
  return C(a) < C(b);
}

template<typename T, typename U>
bool LE(T a, U b) {
  return !LT(b, a);
}

// Specialized version of seg_intersection() from 7.2.3, simplified for include_boundary = false,
// i.e. we're detecting proper crossings of two non-parallel segments whose interiors intersect.
template<typename Pt>
bool proper_seg_intersection(const Pt &a, const Pt &b, const Pt &c, const Pt &d) {
  auto cross = [](const Pt &o, const Pt &p, const Pt &q) {
    return (p.x - o.x) * (q.y - o.y) - (p.y - o.y) * (q.x - o.x);  // Overflow warning.
  };
  auto c1 = cross(a, b, c), c2 = cross(a, b, d), c3 = cross(c, d, a), c4 = cross(c, d, b);
  return ((LT(c1, 0) && LT(0, c2)) || (LT(c2, 0) && LT(0, c1))) &&
         ((LT(c3, 0) && LT(0, c4)) || (LT(c4, 0) && LT(0, c3)));
}

template<typename Pt>
std::vector<std::tuple<Pt, Pt, Pt>> delaunay_triangulation(const std::vector<Pt> &p) {
  int n = static_cast<int>(p.size());
  std::vector<decltype(Pt::x)> z;
  z.reserve(n);
  for (const auto &[px, py] : p) {
    z.emplace_back(px * px + py * py);  // Overflow warning.
  }
  std::vector<std::tuple<Pt, Pt, Pt>> res;
  std::vector<std::tuple<int, int, int>> res_idx;
  for (int i = 0; i < n - 2; i++) {
    for (int j = i + 1; j < n; j++) {
      for (int k = i + 1; k < n; k++) {
        if (j == k) {
          continue;
        }
        // Overflow warning.
        auto nx = (p[j].y - p[i].y) * (z[k] - z[i]) - (p[k].y - p[i].y) * (z[j] - z[i]);
        auto ny = (p[k].x - p[i].x) * (z[j] - z[i]) - (p[j].x - p[i].x) * (z[k] - z[i]);
        auto nz = (p[j].x - p[i].x) * (p[k].y - p[i].y) - (p[k].x - p[i].x) * (p[j].y - p[i].y);
        if (LE(0, nz)) {
          continue;
        }
        std::vector<Pt> s1{p[i], p[j], p[k], p[i]};
        for (int m = 0; m < n; m++) {
          if (LT(0, nx * (p[m].x - p[i].x) + ny * (p[m].y - p[i].y) + nz * (z[m] - z[i]))) {
            goto skip;
          }
        }
        // Reject candidates whose edges properly cross already accepted triangles.
        for (auto [t0, t1, t2] : res_idx) {
          std::vector<Pt> s2{p[t0], p[t1], p[t2], p[t0]};
          for (int u = 0; u < 3; u++) {
            for (int v = 0; v < 3; v++) {
              if (proper_seg_intersection(s1[u], s1[u + 1], s2[v], s2[v + 1])) {
                goto skip;
              }
            }
          }
        }
        res.emplace_back(p[i], p[j], p[k]);
        res_idx.emplace_back(i, j, k);
      skip:;
      }
    }
  }
  return res;
}
\end{lstlisting}
\Needspace*{4\baselineskip}
\textbf{Example Usage}\par
\begin{lstlisting}[style=cpp]
#include <cassert>
using namespace std;

struct Point {
  double x, y;
  Point(double x = 0, double y = 0) : x(x), y(y) {}
  bool operator==(const Point &p) const { return x == p.x && y == p.y; }
  bool operator<(const Point &p) const { return x != p.x ? x < p.x : y < p.y; }
};

struct PointI {
  int x, y;
  PointI(int x = 0, int y = 0) : x(x), y(y) {}
  bool operator==(const PointI &p) const { return x == p.x && y == p.y; }
  bool operator<(const PointI &p) const { return x != p.x ? x < p.x : y < p.y; }
};

int main() {
  vector<Point> v{{1, 3}, {1, 2}, {2, 1}, {0, 0}, {-1, 3}};
  vector<tuple<Point, Point, Point>> t{
      {Point{1, 3}, Point{1, 2}, Point{-1, 3}},
      {Point{1, 3}, Point{2, 1}, Point{1, 2}},
      {Point{1, 2}, Point{2, 1}, Point{0, 0}},
      {Point{1, 2}, Point{0, 0}, Point{-1, 3}}
  };
  assert(delaunay_triangulation(v) == t);

  // Integer-coordinate inputs are triangulated using exact integer arithmetic.
  vector<PointI> iv{{1, 3}, {1, 2}, {2, 1}, {0, 0}, {-1, 3}};
  vector<tuple<PointI, PointI, PointI>> ti{
      {PointI{1, 3}, PointI{1, 2}, PointI{-1, 3}},
      {PointI{1, 3}, PointI{2, 1}, PointI{1, 2}},
      {PointI{1, 2}, PointI{2, 1}, PointI{0, 0}},
      {PointI{1, 2}, PointI{0, 0}, PointI{-1, 3}}
  };
  assert(delaunay_triangulation(iv) == ti);
  return 0;
}
\end{lstlisting}
\subsection{Delaunay Triangulation (Guibas-Stolfi)}
\label{sec:7.4.6}
Given a set $P$ of distinct two-dimensional points, a Delaunay triangulation of $P$ is a triangulation of the convex hull of $P$ such that no point of $P$ lies strictly inside the circumcircle of any triangle in the triangulation. The Delaunay triangulation is not necessarily unique when four or more points are cocircular.

This implementation produces one valid triangulation using the Guibas-Stolfi divide-and-conquer algorithm with a quad-edge data structure. The point set is split in half, each half is triangulated recursively, and the halves are stitched together from the bottom up by a sequence of cross edges, with circumcircle tests deciding each connecting edge and deleting invalidated ones.

The point type must provide exact lexicographic \inlinecode{operator\textless{}}, since sorting and duplicate removal require a strict ordering rather than epsilon equality. All arithmetic uses the point's coordinate type, so integer inputs yield an exact triangulation.

\begin{compactitem}
  \item \inlinecode{delaunay\_triangulation(p)} returns the triangles of one Delaunay triangulation as a vector of counterclockwise-oriented \inlinecode{Point\textless{}T\textgreater{}} triples (\inlinecode{std::tuple}). Duplicate points are ignored. If fewer than three non-collinear unique points remain, the result is empty.
\end{compactitem}

Overflow warning: the in-circle test grows like the fourth power of the coordinate magnitude and dominates the overflow budget. With 64-bit integer coordinates, it overflows once coordinates exceed roughly $14000$. For floating-point coordinates, the predicates are subject to the usual rounding error.

\textbf{Time Complexity}: $O(n \log n)$ per call, where $n$ is the number of input points.

\Needspace*{3\baselineskip}
\textbf{Space Complexity}:
\begin{compactitem}
  \item $O(n \log n)$ for allocated quad-edges in the worst case because deleted edges remain in the pool.
  \item $O(n)$ for the returned triangulation.
\end{compactitem}

\begin{lstlisting}[style=cpp]
#include <algorithm>
#include <cmath>
#include <tuple>
#include <type_traits>
#include <utility>
#include <vector>

const double EPS = 1e-9;

template<typename T, typename U, typename C = std::common_type_t<T, U>>
bool EQ(T a, U b) {
  if constexpr (std::is_floating_point_v<C>) return C(a) == C(b) || std::fabs(C(a) - C(b)) <= EPS;
  return C(a) == C(b);
}

template<typename T, typename U, typename C = std::common_type_t<T, U>>
bool LT(T a, U b) {
  if constexpr (std::is_floating_point_v<C>) return C(a) < C(b) - EPS;
  return C(a) < C(b);
}

template<typename T>
struct Point {
  T x, y;
  Point(T x = 0, T y = 0) : x(x), y(y) {}
  bool operator==(const Point &p) const { return x == p.x && y == p.y; }
  friend bool EQ(const Point &a, const Point &b) { return EQ(a.x, b.x) && EQ(a.y, b.y); }
  bool operator<(const Point &p) const { return std::tie(x, y) < std::tie(p.x, p.y); }
};

template<typename T>
T cross(const Point<T> &a, const Point<T> &b, const Point<T> &o) {
  return (a.x - o.x) * (b.y - o.y) - (a.y - o.y) * (b.x - o.x);  // Overflow warning.
}

template<typename T>
struct Edge {
  Point<T> origin;
  Edge *rot, *onext;
  bool primal, deleted, used;

  Edge *rev() const { return rot->rot; }
  Edge *lnext() const { return rot->rev()->onext->rot; }
  Edge *oprev() const { return rot->onext->rot; }
  Point<T> dest() const { return rev()->origin; }
};

template<typename T>
struct QuadEdgePool {
  std::vector<Edge<T> *> edges;

  QuadEdgePool() = default;
  QuadEdgePool(const QuadEdgePool &) = delete;
  QuadEdgePool &operator=(const QuadEdgePool &) = delete;

  ~QuadEdgePool() {
    for (Edge<T> *e : edges) {
      delete e;
    }
  }

  Edge<T> *make_edge(const Point<T> &u, const Point<T> &v) {
    Edge<T> *e1 = new Edge<T>, *e2 = new Edge<T>, *e3 = new Edge<T>, *e4 = new Edge<T>;
    edges.push_back(e1);
    edges.push_back(e2);
    edges.push_back(e3);
    edges.push_back(e4);
    e1->origin = u;
    e2->origin = v;
    e1->rot = e3;
    e2->rot = e4;
    e3->rot = e2;
    e4->rot = e1;
    e1->onext = e1;
    e2->onext = e2;
    e3->onext = e4;
    e4->onext = e3;
    e1->primal = e2->primal = true;
    e3->primal = e4->primal = false;
    e1->deleted = e2->deleted = e3->deleted = e4->deleted = false;
    e1->used = e2->used = e3->used = e4->used = false;
    return e1;
  }

  void delete_edge(Edge<T> *e) {
    splice(e, e->oprev());
    splice(e->rev(), e->rev()->oprev());
    e->deleted = e->rot->deleted = e->rev()->deleted = e->rev()->rot->deleted = true;
  }

  Edge<T> *connect(Edge<T> *a, Edge<T> *b) {
    Edge<T> *e = make_edge(a->dest(), b->origin);
    splice(e, a->lnext());
    splice(e->rev(), b);
    return e;
  }

  static void splice(Edge<T> *a, Edge<T> *b) {
    std::swap(a->onext->rot->onext, b->onext->rot->onext);
    std::swap(a->onext, b->onext);
  }
};

template<typename T>
bool left_of(const Point<T> &p, Edge<T> *e) {
  return LT(0, cross(e->origin, e->dest(), p));
}

template<typename T>
bool right_of(const Point<T> &p, Edge<T> *e) {
  return LT(cross(e->origin, e->dest(), p), 0);
}

template<typename T>
bool in_circle(const Point<T> &a, const Point<T> &b, const Point<T> &c, const Point<T> &d) {
  T adx = a.x - d.x, ady = a.y - d.y;
  T bdx = b.x - d.x, bdy = b.y - d.y;
  T cdx = c.x - d.x, cdy = c.y - d.y;
  T abdet = adx * bdy - bdx * ady;  // Overflow warning.
  T bcdet = bdx * cdy - cdx * bdy;
  T cadet = cdx * ady - adx * cdy;
  T alift = adx * adx + ady * ady;
  T blift = bdx * bdx + bdy * bdy;
  T clift = cdx * cdx + cdy * cdy;
  return LT(0, alift * bcdet + blift * cadet + clift * abdet);
}

template<typename T>
std::pair<Edge<T> *, Edge<T> *> build_triangulation(
    const std::vector<Point<T>> &p, int lo, int hi, QuadEdgePool<T> &pool
) {
  if (hi - lo == 1) {
    Edge<T> *a = pool.make_edge(p[lo], p[hi]);
    return {a, a->rev()};
  }
  if (hi - lo == 2) {
    Edge<T> *a = pool.make_edge(p[lo], p[lo + 1]);
    Edge<T> *b = pool.make_edge(p[lo + 1], p[hi]);
    pool.splice(a->rev(), b);
    int side = LT(0, cross(p[lo], p[lo + 1], p[hi]))
                   ? 1
                   : (LT(cross(p[lo], p[lo + 1], p[hi]), 0) ? -1 : 0);
    if (side == 0) {
      return {a, b->rev()};
    }
    Edge<T> *c = pool.connect(b, a);
    if (side == 1) {
      return {a, b->rev()};
    }
    return {c->rev(), c};
  }
  int mid = lo + (hi - lo) / 2;
  auto [ldo, ldi] = build_triangulation(p, lo, mid, pool);
  auto [rdi, rdo] = build_triangulation(p, mid + 1, hi, pool);
  while (true) {
    if (left_of(rdi->origin, ldi)) {
      ldi = ldi->lnext();
    } else if (right_of(ldi->origin, rdi)) {
      rdi = rdi->rev()->onext;
    } else {
      break;
    }
  }
  Edge<T> *base = pool.connect(rdi->rev(), ldi);
  if (ldi->origin == ldo->origin) {
    ldo = base->rev();
  }
  if (rdi->origin == rdo->origin) {
    rdo = base;
  }
  while (true) {
    Edge<T> *lcand = base->rev()->onext;
    if (right_of(lcand->dest(), base)) {
      while (in_circle(base->dest(), base->origin, lcand->dest(), lcand->onext->dest())) {
        Edge<T> *next = lcand->onext;
        pool.delete_edge(lcand);
        lcand = next;
      }
    }
    Edge<T> *rcand = base->oprev();
    if (right_of(rcand->dest(), base)) {
      while (in_circle(base->dest(), base->origin, rcand->dest(), rcand->oprev()->dest())) {
        Edge<T> *prev = rcand->oprev();
        pool.delete_edge(rcand);
        rcand = prev;
      }
    }
    bool lvalid = right_of(lcand->dest(), base), rvalid = right_of(rcand->dest(), base);
    if (!lvalid && !rvalid) {
      break;
    }
    if (!lvalid ||
        (rvalid && in_circle(lcand->dest(), lcand->origin, rcand->origin, rcand->dest()))) {
      base = pool.connect(rcand, base->rev());
    } else {
      base = pool.connect(base->rev(), lcand->rev());
    }
  }
  return {ldo, rdo};
}

template<typename T>
auto delaunay_triangulation(std::vector<Point<T>> p) {
  using Pt = Point<T>;
  using Triangle = std::tuple<Pt, Pt, Pt>;
  // Rotates triangle vertices so that the lexicographically smallest comes first, preserving the
  // cyclic (counterclockwise) order, to give each triangle a canonical, comparable representation.
  auto canonical = [](const Pt &a, const Pt &b, const Pt &c) -> Triangle {
    if (b < a && b < c) {
      return {b, c, a};
    }
    if (c < a && c < b) {
      return {c, a, b};
    }
    return {a, b, c};
  };
  std::sort(p.begin(), p.end());
  auto same_point = [](const Pt &a, const Pt &b) { return !(a < b) && !(b < a); };
  p.erase(std::unique(p.begin(), p.end(), same_point), p.end());
  std::vector<Triangle> res;
  if (p.size() < 3) {
    return res;
  }
  QuadEdgePool<T> pool;
  build_triangulation(p, 0, static_cast<int>(p.size()) - 1, pool);
  for (Edge<T> *start : pool.edges) {
    if (!start->primal || start->deleted || start->used) {
      continue;
    }
    std::vector<Pt> face;
    Edge<T> *curr = start;
    do {
      curr->used = true;
      face.push_back(curr->origin);
      curr = curr->lnext();
    } while (curr != start);
    if (face.size() == 3 && LT(0, cross(face[0], face[1], face[2]))) {
      res.push_back(canonical(face[0], face[1], face[2]));
    }
  }
  std::sort(res.begin(), res.end());
  return res;
}
\end{lstlisting}
\Needspace*{4\baselineskip}
\textbf{Example Usage}\par
\begin{lstlisting}[style=cpp]
#include <cassert>
#include <cstdint>
using namespace std;

using PointL = Point<int64_t>;
using PointD = Point<double>;

int main() {
  vector<PointD> v{{1, 3}, {1, 2}, {2, 1}, {0, 0}, {-1, 3}};
  vector<tuple<PointD, PointD, PointD>> t{
      {PointD{-1, 3}, PointD{0, 0}, PointD{1, 2}},
      {PointD{-1, 3}, PointD{1, 2}, PointD{1, 3}},
      {PointD{0, 0}, PointD{2, 1}, PointD{1, 2}},
      {PointD{1, 2}, PointD{2, 1}, PointD{1, 3}}
  };
  assert(delaunay_triangulation(v) == t);

  // Integer-coordinate inputs are triangulated using exact integer arithmetic.
  vector<PointL> iv{{1, 3}, {1, 2}, {2, 1}, {0, 0}, {-1, 3}};
  vector<tuple<PointL, PointL, PointL>> ti{
      {PointL{-1, 3}, PointL{0, 0}, PointL{1, 2}},
      {PointL{-1, 3}, PointL{1, 2}, PointL{1, 3}},
      {PointL{0, 0}, PointL{2, 1}, PointL{1, 2}},
      {PointL{1, 2}, PointL{2, 1}, PointL{1, 3}}
  };
  assert(delaunay_triangulation(iv) == ti);
  return 0;
}
\end{lstlisting}

\section{\texorpdfstring{3D Geometry}{7.5 3D Geometry}}
\setcounter{section}{5}
\setcounter{subsection}{0}
\subsection{3D Point}
\label{sec:7.5.1}
A 3D point class templated on the coordinate type \inlinecode{T}. Algebraic operations preserve \inlinecode{T}, while metric operations use \inlinecode{fp\_t}, which is \inlinecode{T} for floating-point coordinates and \inlinecode{double} otherwise.

\begin{compactitem}
  \item \inlinecode{TPoint3\textless{}T\textgreater{}(x = 0, y = 0, z = 0)} constructs point (\inlinecode{x}, \inlinecode{y}, \inlinecode{z}).
  \item Operators \inlinecode{+}, \inlinecode{-}, \inlinecode{*}, and \inlinecode{/} preserve the coordinate type except that division promotes to \inlinecode{fp\_t}. Comparisons are exact and lexicographic, while \inlinecode{EQ(p, q)} uses \inlinecode{EPS} for floating-point coordinates and remains exact for other types.
  \item \inlinecode{sqnorm()}, \inlinecode{dot(p)}, and \inlinecode{cross(p)} return exact algebraic results in coordinate type \inlinecode{T}.
  \item \inlinecode{norm()}, \inlinecode{phi()}, and \inlinecode{theta()} return the length, azimuth, and polar angle in \inlinecode{fp\_t}.
  \item \inlinecode{unit()} returns the unit vector and requires a nonzero point; \inlinecode{normal(p)} returns a unit normal and requires nonparallel vectors.
  \item \inlinecode{rotate(angle, axis)} rotates the point about a nonzero \inlinecode{axis} using Rodrigues' formula.
  \item \inlinecode{Point3I}, \inlinecode{Point3L}, \inlinecode{Point3D}, and \inlinecode{Point3LD} use \inlinecode{int}, \inlinecode{int64\_t}, \inlinecode{double}, and \inlinecode{long double} coordinates, respectively.
\end{compactitem}

Overflow warning: the exact products \inlinecode{dot()}, \inlinecode{cross()}, and \inlinecode{sqnorm()} grow like the squared coordinate magnitude. With \inlinecode{TPoint3\textless{}int\textgreater{}} these overflow a 32-bit \inlinecode{int} once coordinates exceed a few tens of thousands, so use \inlinecode{Point3L} (\inlinecode{TPoint3\textless{}int64\_t\textgreater{}}) for larger integer coordinates.

\textbf{Time Complexity}: $O(1)$ per operation.

\textbf{Space Complexity}: $O(1)$ for storage and auxiliary.

\begin{lstlisting}[style=cpp]
#include <cmath>
#include <cstdint>
#include <tuple>
#include <type_traits>

const double EPS = 1e-9;

template<typename T, typename U, typename C = std::common_type_t<T, U>>
bool EQ(T a, U b) {
  if constexpr (std::is_floating_point_v<C>) return C(a) == C(b) || std::fabs(C(a) - C(b)) <= EPS;
  return C(a) == C(b);
}

template<typename T>
struct TPoint3 {
  using fp_t = std::conditional_t<std::is_floating_point<T>::value, T, double>;

  T x, y, z;

  TPoint3(T x = 0, T y = 0, T z = 0) : x(x), y(y), z(z) {}
  bool operator==(const TPoint3 &p) const { return std::tie(x, y, z) == std::tie(p.x, p.y, p.z); }

  friend bool EQ(const TPoint3 &a, const TPoint3 &b) {
    return EQ(a.x, b.x) && EQ(a.y, b.y) && EQ(a.z, b.z);
  }

  bool operator!=(const TPoint3 &p) const { return !(*this == p); }
  bool operator<(const TPoint3 &p) const { return std::tie(x, y, z) < std::tie(p.x, p.y, p.z); }
  TPoint3 operator+(const TPoint3 &p) const { return {x + p.x, y + p.y, z + p.z}; }
  TPoint3 operator-(const TPoint3 &p) const { return {x - p.x, y - p.y, z - p.z}; }
  TPoint3 operator*(T k) const { return {x * k, y * k, z * k}; }
  TPoint3<fp_t> operator/(fp_t k) const { return {fp_t(x) / k, fp_t(y) / k, fp_t(z) / k}; }
  T dot(const TPoint3 &p) const { return x * p.x + y * p.y + z * p.z; }  // Overflow warning.
  T sqnorm() const { return x * x + y * y + z * z; }                     // Overflow warning.
  fp_t norm() const { return std::hypot(std::hypot(fp_t(x), fp_t(y)), fp_t(z)); }
  fp_t phi() const { return std::atan2(fp_t(y), fp_t(x)); }
  fp_t theta() const { return std::atan2(std::hypot(fp_t(x), fp_t(y)), fp_t(z)); }
  TPoint3<fp_t> unit() const { return *this / norm(); }

  TPoint3 cross(const TPoint3 &p) const {
    return {y * p.z - z * p.y, z * p.x - x * p.z, x * p.y - y * p.x};  // Overflow warning.
  }

  TPoint3<fp_t> normal(const TPoint3 &p) const { return cross(p).unit(); }

  TPoint3<fp_t> rotate(fp_t angle, const TPoint3 &axis) const {
    fp_t s = std::sin(angle), c = std::cos(angle);
    TPoint3<fp_t> u = axis.unit(), v{fp_t(x), fp_t(y), fp_t(z)};
    return u * v.dot(u) * (1 - c) + v * c - v.cross(u) * s;
  }
};

using Point3I = TPoint3<int>;
using Point3L = TPoint3<int64_t>;
using Point3D = TPoint3<double>;
using Point3LD = TPoint3<long double>;
\end{lstlisting}
\Needspace*{4\baselineskip}
\textbf{Example Usage}\par
\begin{lstlisting}[style=cpp]
#include <cassert>
using namespace std;

int main() {
  Point3I a(1, 0, 0), b(0, 1, 0);
  assert(a.dot(b) == 0);
  assert(a.cross(b) == Point3I(0, 0, 1));
  assert(a.sqnorm() == 1);

  Point3D p(1, 0, 0);
  assert(EQ(p.rotate(acos(-1.0) / 2, Point3D(0, 0, 1)), Point3D(0, 1, 0)));
  return 0;
}
\end{lstlisting}
\subsection{Polyhedron Volume}
\label{sec:7.5.2}
Given an oriented triangular mesh of a closed polyhedron, compute its signed volume. Each triangle contributes the signed volume of the tetrahedron formed by the triangle $abc$ and the origin $((a \times b) \cdot c) / 6$. If faces are oriented outward, the result is positive for the orientation used by \inlinecode{convex\_hull\_3d()} in \secref{7.5.3}; take \inlinecode{abs()} if only the magnitude is needed.

\begin{compactitem}
  \item \inlinecode{signed\_polyhedron\_volume(p, faces)} returns the signed volume. Each face triplet (\inlinecode{a}, \inlinecode{b}, \inlinecode{c}) represents the triangle with vertices \inlinecode{p[a]}, \inlinecode{p[b]}, and \inlinecode{p[c]}, and its indices must be valid.
\end{compactitem}

Overflow warning: for integral point types, each scalar triple product must fit in the point's coordinate arithmetic type before it is converted to \inlinecode{double}.

\textbf{Time Complexity}: $O(f)$ per call, where $f$ is the number of triangular faces.

\textbf{Space Complexity}: $O(1)$ auxiliary.

\begin{lstlisting}[style=cpp]
#include <algorithm>
#include <cassert>
#include <vector>

struct Point3D {
  double x, y, z;

  Point3D(double x = 0, double y = 0, double z = 0) : x(x), y(y), z(z) {}
  double dot(const Point3D &p) const { return x * p.x + y * p.y + z * p.z; }

  Point3D cross(const Point3D &p) const {
    return {y * p.z - z * p.y, z * p.x - x * p.z, x * p.y - y * p.x};
  }
};

template<typename Pt, typename F>
double signed_polyhedron_volume(const std::vector<Pt> &p, const std::vector<F> &faces) {
  assert(std::all_of(faces.begin(), faces.end(), [&](const F &face) {
    auto [a, b, c] = face;
    return 0 <= a && a < static_cast<int>(p.size()) && 0 <= b && b < static_cast<int>(p.size()) &&
           0 <= c && c < static_cast<int>(p.size());
  }));
  double volume = 0;
  for (const auto &[a, b, c] : faces) {
    volume += p[a].cross(p[b]).dot(p[c]);  // Overflow warning.
  }
  return volume / 6.0;
}
\end{lstlisting}
\Needspace*{4\baselineskip}
\textbf{Example Usage}\par
\begin{lstlisting}[style=cpp]
#include <cassert>
#include <cmath>
#include <tuple>
using namespace std;

bool EQ(double a, double b) {
  return fabs(a - b) < 1e-9;
}

int main() {
  using Face = tuple<int, int, int>;  // (a, b, c)
  vector<Point3D> p{{0, 0, 0}, {1, 0, 0}, {0, 1, 0}, {0, 0, 1}};
  vector<Face> faces{
      Face(0, 2, 1),
      Face(0, 1, 3),
      Face(0, 3, 2),
      Face(1, 2, 3),
  };
  assert(EQ(signed_polyhedron_volume(p, faces), 1.0 / 6.0));
  return 0;
}
\end{lstlisting}
\subsection{3D Convex Hull}
\label{sec:7.5.3}
The convex hull of points in 3D is the smallest convex polyhedron containing them. This incremental algorithm computes its triangular faces: start with a tetrahedron, delete every visible face for each new point, and stitch the horizon edges to that point. Returned faces are oriented outward.

\begin{compactitem}
  \item \inlinecode{convex\_hull\_3d(p)} returns triangular faces as \inlinecode{Face\{a, b, c\}} triples, where each face has vertices \inlinecode{p[a]}, \inlinecode{p[b]}, and \inlinecode{p[c]}. The input must contain at least four points, with no four coplanar. For coplanar or highly degenerate inputs, perturb the points slightly or use a more robust library.
\end{compactitem}

\textbf{Time Complexity}: $O(n^{2})$ per call, where $n$ is the number of points.

\textbf{Space Complexity}: $O(n^{2})$ auxiliary.

\begin{lstlisting}[style=cpp]
#include <array>
#include <cassert>
#include <utility>
#include <vector>

const double EPS = 1e-9;

struct Point3D {
  double x, y, z;

  Point3D(double x = 0, double y = 0, double z = 0) : x(x), y(y), z(z) {}
  Point3D operator-(const Point3D &p) const { return {x - p.x, y - p.y, z - p.z}; }
  Point3D operator*(double k) const { return {x * k, y * k, z * k}; }
  double dot(const Point3D &p) const { return x * p.x + y * p.y + z * p.z; }

  Point3D cross(const Point3D &p) const {
    return {y * p.z - z * p.y, z * p.x - x * p.z, x * p.y - y * p.x};
  }
};

struct Face {
  int a, b, c;
  Face(int a = 0, int b = 0, int c = 0) : a(a), b(b), c(c) {}
};

std::vector<Face> convex_hull_3d(const std::vector<Point3D> &p) {
  int n = static_cast<int>(p.size());
  assert(n >= 4);
  std::vector<std::vector<std::array<int, 2>>> edge(
      n, std::vector<std::array<int, 2>>(n, {-1, -1})
  );
  std::vector<Face> faces;
  auto normal = [&](const Face &f) { return (p[f.b] - p[f.a]).cross(p[f.c] - p[f.a]); };
  auto visible = [&](const Face &f, int i) { return normal(f).dot(p[i] - p[f.a]) > EPS; };
  auto add_edge = [&](int a, int b, int c) {
    if (a > b) {
      std::swap(a, b);
    }
    auto &e = edge[a][b];
    (e[0] == -1 ? e[0] : e[1]) = c;
  };
  auto remove_edge = [&](int a, int b, int c) {
    if (a > b) {
      std::swap(a, b);
    }
    auto &e = edge[a][b];
    (e[0] == c ? e[0] : e[1]) = -1;
  };
  auto edge_count = [&](int a, int b) {
    if (a > b) {
      std::swap(a, b);
    }
    return (edge[a][b][0] != -1) + (edge[a][b][1] != -1);
  };
  auto make_face = [&](int a, int b, int c, int inside) {
    Face f(a, b, c);
    if (normal(f).dot(p[inside] - p[a]) > 0) {
      std::swap(f.b, f.c);
    }
    add_edge(f.a, f.b, f.c);
    add_edge(f.b, f.c, f.a);
    add_edge(f.c, f.a, f.b);
    faces.push_back(f);
  };
  for (int i = 0; i < 4; i++) {
    for (int j = i + 1; j < 4; j++) {
      for (int k = j + 1; k < 4; k++) {
        make_face(i, j, k, 6 - i - j - k);
      }
    }
  }
  for (int i = 4; i < n; i++) {
    for (int j = 0; j < static_cast<int>(faces.size()); j++) {
      Face f = faces[j];
      if (visible(f, i)) {
        remove_edge(f.a, f.b, f.c);
        remove_edge(f.b, f.c, f.a);
        remove_edge(f.c, f.a, f.b);
        std::swap(faces[j--], faces.back());
        faces.pop_back();
      }
    }
    int old_faces = static_cast<int>(faces.size());
    for (int j = 0; j < old_faces; j++) {
      Face f = faces[j];
      if (edge_count(f.a, f.b) != 2) {
        make_face(f.a, f.b, i, f.c);
      }
      if (edge_count(f.b, f.c) != 2) {
        make_face(f.b, f.c, i, f.a);
      }
      if (edge_count(f.c, f.a) != 2) {
        make_face(f.c, f.a, i, f.b);
      }
    }
  }
  return faces;
}
\end{lstlisting}
\Needspace*{4\baselineskip}
\textbf{Example Usage}\par
\begin{lstlisting}[style=cpp]
using namespace std;

void assert_outward(const vector<Point3D> &p, const vector<Face> &faces) {
  for (const Face &f : faces) {
    assert(0 <= f.a && f.a < static_cast<int>(p.size()));
    assert(0 <= f.b && f.b < static_cast<int>(p.size()));
    assert(0 <= f.c && f.c < static_cast<int>(p.size()));
    assert(f.a != f.b && f.b != f.c && f.c != f.a);
    Point3D normal = (p[f.b] - p[f.a]).cross(p[f.c] - p[f.a]);
    for (const Point3D &q : p) {
      assert(normal.dot(q - p[f.a]) <= EPS);
    }
  }
}

int main() {
  vector<Point3D> tetra{{0, 0, 0}, {1, 0, 0}, {0, 1, 0}, {0, 0, 1}};
  auto faces = convex_hull_3d(tetra);
  // A tetrahedron has four triangular faces.
  assert(faces.size() == 4);
  assert_outward(tetra, faces);

  tetra.emplace_back(1, 1, 1);
  faces = convex_hull_3d(tetra);
  // Adding the opposite point makes a triangular bipyramid with six triangular faces.
  assert(faces.size() == 6);
  assert_outward(tetra, faces);
  return 0;
}
\end{lstlisting}

\chapter{\texorpdfstring{Miscellany}{8. Miscellany}}

\section{\texorpdfstring{Contest Utilities}{8.1 Contest Utilities}}
\setcounter{section}{1}
\label{sec:8.1}
Small contest convenience helpers that are useful across many algorithms. These snippets avoid algorithm-specific policy and are meant to be pasted near the top of a solution file.

A personal template may also collect shortened versions of utilities elsewhere in the anthology, such as fast I/O and debugging, custom hashing, coordinate compression, bit operations, modular arithmetic, binary search, and GNU policy-based data structures. They remain in their dedicated sections here to avoid duplication and let each template be assembled to taste.

\begin{compactitem}
  \item \inlinecode{Rep(i,N)}, \inlinecode{For(i,L,H)}, \inlinecode{Rev(i,N)}, \inlinecode{Dwn(i,H,L)}, and \inlinecode{Each(x,C)} are compact loop macros.
  \item \inlinecode{All(c)} and \inlinecode{Rall(c)} expand to the forward or reverse iterator pair of a container or array.
  \item \inlinecode{sz(c)} returns the signed size of a container.
  \item \inlinecode{ckmin(a, b)} assigns \inlinecode{a = b} only when \inlinecode{b} \textless{} \inlinecode{a}, and returns whether the assignment happened.
  \item \inlinecode{ckmax(a, b)} assigns \inlinecode{a = b} only when \inlinecode{a} \textless{} \inlinecode{b}, and returns whether the assignment happened.
  \item \inlinecode{floor\_div(a, b)} and \inlinecode{ceil\_div(a, b)} divide signed integers with mathematical rounding toward negative or positive infinity. Requires nonzero \inlinecode{b} and the resulting quotient to be representable by the result type; in particular, the minimum value cannot be divided by $-1$.
  \item \inlinecode{make\_vec(n1, ..., nk, v)} returns a \inlinecode{k}-dimensional nested vector with the given sizes and every element initialized to \inlinecode{v}, whose type is deduced from \inlinecode{v}. The last argument is always the fill value, so \inlinecode{make\_vec(3, 5)} is a vector of three fives rather than a $3$ by $5$ grid.
  \item \inlinecode{indices(n)} returns the vector $[0, 1, \ldots, {\inlinecodemath{n}} - 1]$.
  \item \inlinecode{lb(a, x, comp = std::less\textless{}\textgreater{})} and \inlinecode{ub(a, x, comp = std::less\textless{}\textgreater{})} return lower and upper bound indices in random-access sequence \inlinecode{a}, which must be sorted according to \inlinecode{comp}.
  \item \inlinecode{sort\_unique(v, comp = std::less\textless{}\textgreater{})} sorts a vector and removes duplicates, which are the elements that \inlinecode{comp} deems equivalent.
  \item \inlinecode{argsort(a, comp = std::less\textless{}\textgreater{})} returns the indices of \inlinecode{a} ordered by their values according to \inlinecode{comp}; the order of indices whose values compare equal is unspecified.
  \item \inlinecode{find\_ptr(m, k)} returns a pointer to the value that key \inlinecode{k} maps to in the associative container \inlinecode{m}, or \inlinecode{nullptr} if the key is absent. The pointer is const only when \inlinecode{m} is, so \inlinecode{if (auto *p = find\_ptr(m, k))} both tests for the key and updates it in a single lookup.
  \item \inlinecode{get\_or(m, k, def = \{\})} returns the value that key \inlinecode{k} maps to in the associative container \inlinecode{m}, or \inlinecode{def} if the key is absent. Unlike \inlinecode{operator[]}, it never inserts and works on a const \inlinecode{m}.
  \item \inlinecode{erase\_one(c, x)} erases one copy of \inlinecode{x} from associative container \inlinecode{c}, asserting that it is present.
  \item \inlinecode{min\_heap\textless{}T\textgreater{}} is a min-heap alias.
  \item \inlinecode{y\_combinator(f)} wraps a recursive lambda so that it can call itself as its first argument.
\end{compactitem}

\begin{lstlisting}[style=cpp]
#include <algorithm>
#include <cassert>
#include <functional>
#include <iterator>
#include <numeric>
#include <queue>
#include <type_traits>
#include <utility>
#include <vector>

#define Rep(i, N)    for (int i = 0, _##i = (N); i < _##i; ++i)     //   0 to N-1
#define For(i, L, H) for (int i = (L), _##i = (H); i <= _##i; ++i)  //   L to H
#define Rev(i, N)    for (int i = (N); --i >= 0;)                   // N-1 to 0
#define Dwn(i, H, L) for (int i = (H), _##i = (L); i >= _##i; --i)  //   H to L
#define Each(x, C)   for (auto &x : (C))
#define All(C)       std::begin(C), std::end(C)
#define Rall(C)      std::rbegin(C), std::rend(C)

template<typename C> int sz(const C &c) { return static_cast<int>(c.size()); }
template<typename T, typename U> bool ckmin(T &a, const U &b) { return b < a && (a = b, true); }
template<typename T, typename U> bool ckmax(T &a, const U &b) { return a < b && (a = b, true); }

template<typename T>
T floor_div(T a, T b) {
  assert(b != 0);
  T q = a / b, r = a % b;
  return q - (r != 0 && ((r < 0) != (b < 0)));
}

template<typename T>
T ceil_div(T a, T b) {
  assert(b != 0);
  T q = a / b, r = a % b;
  return q + (r != 0 && ((r < 0) == (b < 0)));
}

template<typename T>
std::vector<T> make_vec(int n, const T &v) {
  return std::vector<T>(n, v);
}

template<typename... Args>
auto make_vec(int n, Args... args) {
  auto inner = make_vec(args...);
  return std::vector<decltype(inner)>(n, inner);
}

std::vector<int> indices(int n) {
  std::vector<int> v(n);
  std::iota(v.begin(), v.end(), 0);
  return v;
}

template<typename Seq, typename T, typename Compare = std::less<>>
int lb(const Seq &a, const T &x, Compare comp = Compare{}) {
  return static_cast<int>(std::lower_bound(a.begin(), a.end(), x, comp) - a.begin());
}

template<typename Seq, typename T, typename Compare = std::less<>>
int ub(const Seq &a, const T &x, Compare comp = Compare{}) {
  return static_cast<int>(std::upper_bound(a.begin(), a.end(), x, comp) - a.begin());
}

template<typename T, typename Compare = std::less<>>
void sort_unique(std::vector<T> &v, Compare comp = Compare{}) {
  std::sort(v.begin(), v.end(), comp);
  auto same = [&](const T &a, const T &b) { return !comp(a, b) && !comp(b, a); };
  v.erase(std::unique(v.begin(), v.end(), same), v.end());
}

template<typename Seq, typename Compare = std::less<>>
std::vector<int> argsort(const Seq &a, Compare comp = Compare{}) {
  std::vector<int> p = indices(static_cast<int>(a.size()));
  std::sort(p.begin(), p.end(), [&](int i, int j) { return comp(a[i], a[j]); });
  return p;
}

template<typename M>
auto find_ptr(M &m, const typename M::key_type &k) {
  auto it = m.find(k);
  return it == m.end() ? nullptr : &it->second;
}

template<typename M>
typename M::mapped_type get_or(
    const M &m, const typename M::key_type &k, const typename M::mapped_type &def = {}
) {
  auto it = m.find(k);
  return it == m.end() ? def : it->second;
}

template<typename C, typename T>
void erase_one(C &c, const T &x) {
  auto it = c.find(x);
  assert(it != c.end());
  c.erase(it);
}

template<typename T>
using min_heap = std::priority_queue<T, std::vector<T>, std::greater<T>>;

template<typename Fn>
class y_combinator_result {
  Fn fn;

 public:
  template<typename T>
  explicit y_combinator_result(T &&f) : fn(std::forward<T>(f)) {}

  template<typename... Args>
  decltype(auto) operator()(Args &&...args) {
    return fn(std::ref(*this), std::forward<Args>(args)...);
  }
};

template<typename Fn>
decltype(auto) y_combinator(Fn &&f) {
  return y_combinator_result<std::decay_t<Fn>>(std::forward<Fn>(f));
}
\end{lstlisting}
\Needspace*{4\baselineskip}
\textbf{Example Usage}\par
\begin{lstlisting}[style=cpp]
#include <map>
#include <set>
using namespace std;

int main() {
  int tot = 0;
  Rep(i, 4) tot += i;     // 0 + 1 + 2 + 3 = 6
  For(i, 2, 5) tot += i;  // 6 + (2 + 3 + 4 + 5) = 20
  Rev(i, 3) tot += i;     // 20 + (2 + 1 + 0) = 23
  Dwn(i, 8, 6) tot += i;  // 23 + (8 + 7 + 6) = 44
  assert(tot == 44);

  int x = 10;
  assert(ckmin(x, 7) && x == 7);
  assert(!ckmin(x, 9) && x == 7);
  assert(ckmax(x, 11) && x == 11);

  vector<int> v{1, 2, 3};
  assert(sz(v) == 3);
  sort(Rall(v));
  assert((v == vector<int>{3, 2, 1}));
  assert(lb(v, 2, greater<>{}) == 1 && ub(v, 2, greater<>{}) == 2);
  sort(All(v));
  assert(lb(v, 2) == 1 && ub(v, 2) == 2);
  assert(floor_div(7, -3) == -3);
  assert(floor_div(-7, 3) == -3);
  assert(ceil_div(7, -3) == -2);
  assert(ceil_div(-7, 3) == -2);

  auto grid = make_vec(2, 3, -1);
  grid[0][0] = 7;
  assert(sz(grid) == 2 && sz(grid[1]) == 3 && grid[1][0] == -1);
  v = {3, 1, 3, 2, 1};
  sort_unique(v);
  assert((v == vector<int>{1, 2, 3}));
  v = {3, 1, 3, 2, 1};
  sort_unique(v, greater<>{});
  assert((v == vector<int>{3, 2, 1}));
  assert((indices(4) == vector<int>{0, 1, 2, 3}));
  assert((argsort(vector<int>{30, 10, 20}) == vector<int>{1, 2, 0}));
  assert((argsort(vector<int>{30, 10, 20}, greater<>{}) == vector<int>{0, 2, 1}));

  map<int, int> cnt{{2, 5}};
  if (int *p = find_ptr(cnt, 2)) {
    *p += 10;  // A single lookup both tests for the key and updates it.
  }
  const map<int, int> &view = cnt;  // A const container yields a pointer to const.
  assert(*find_ptr(view, 2) == 15 && find_ptr(view, 7) == nullptr);
  assert(get_or(cnt, 2) == 15 && get_or(cnt, 7, -1) == -1 && sz(cnt) == 1);
  min_heap<int> pq;
  pq.push(3);
  pq.push(1);
  assert(pq.top() == 1);

  multiset<int> ms{1, 2, 2, 3};
  erase_one(ms, 2);
  assert(ms.count(2) == 1);

  // Recursive lambda with y_combinator.
  auto fib =
      y_combinator([](auto fib, int n) -> int { return n < 2 ? n : fib(n - 1) + fib(n - 2); });
  assert(fib(10) == 55);

  // Recursive lambda without y_combinator: pass the lambda to itself as the first argument.
  vector<vector<int>> adj{{1, 2}, {0}, {0}};
  int seen = 0;
  auto dfs = [&](auto &&dfs, int u, int p) -> void {
    seen++;
    Each(v, adj[u]) if (v != p) dfs(dfs, v, u);
  };
  dfs(dfs, 0, -1);
  assert(seen == 3);

  // Recursive lambdas are automatically supported in C++23.
#if __cplusplus >= 202302L
  auto fact = [&](this auto fact, int n) -> int { return n ? n * fact(n - 1) : 1; };
  assert(fact(5) == 120);
#endif
  return 0;
}
\end{lstlisting}

\section{\texorpdfstring{Fast IO}{8.2 Fast IO}}
\setcounter{section}{2}
\label{sec:8.2}
Fast input and output helpers. For most cases, it is enough to speed up standard I/O simply with \inlinecode{ios::sync\_with\_stdio(false); cin.tie(nullptr);}. The following classes only come into play when input is genuinely huge or \inlinecode{iostream} overhead is measurable.

\begin{compactitem}
  \item \inlinecode{set\_in(name)}, \inlinecode{set\_out(name)}, and \inlinecode{set\_io(iname, oname)} redirect standard input/output to files. For example, \inlinecode{set\_io("task.in", "task.out")} is convenient for USACO-style problems. An empty name leaves that stream untouched, so \inlinecode{set\_io("", "")} is a no-op that can stay in place when the same solution is submitted to a judge that pipes its input.
\end{compactitem}

File-redirection failures throw \inlinecode{std::runtime\_error}.

\begin{lstlisting}[style=cpp]
#include <cctype>
#include <cstddef>
#include <cstdio>
#include <cstdlib>
#include <limits>
#include <stdexcept>
#include <string>
#include <type_traits>

void set_in(const std::string &name) {
  if (!name.empty() && !freopen(name.c_str(), "r", stdin)) {
    throw std::runtime_error("Failed to open input file: " + name);
  }
}

void set_out(const std::string &name) {
  if (!name.empty() && !freopen(name.c_str(), "w", stdout)) {
    throw std::runtime_error("Failed to open output file: " + name);
  }
}

void set_io(const std::string &iname, const std::string &oname) {
  set_in(iname);
  set_out(oname);
}
\end{lstlisting}
When only integers are read, one loop over a character reader is the whole technique and is short enough to copy or retype: skip anything that starts no number, then fold digits in with \inlinecode{x * 10 + d}. Which character reader it uses decides everything, since the per-character mutex that the locked \inlinecode{getchar()} takes costs more than the parsing does. The example below times the alternatives against each other on the same input.

\begin{compactitem}
  \item \inlinecode{read\_int(x)} reads one integer from \inlinecode{stdin} into \inlinecode{x}, skipping any leading character that cannot start a number. A leading \inlinecode{-} is honored, so \inlinecode{Int} must be signed for negative input. Behavior is undefined at end of file.
\end{compactitem}

\begin{lstlisting}[style=cpp]
// The unlocked reader is where nearly all of the speed comes from. It is POSIX, not standard C++.
inline int read_char() {
#if defined(_WIN32)
  return _getchar_nolock();
#elif defined(__unix__) || defined(__APPLE__)
  return getchar_unlocked();
#else
  return getchar();
#endif
}

template<typename Int>
void read_int(Int &x) {
  int c, s = 1;
  while ((c = read_char()) != '-' && (c < '0' || c > '9'));
  if (c == '-') s = -1, c = read_char();
  for (x = 0; '0' <= c && c <= '9'; c = read_char()) x = x * 10 + c - '0';
  x *= s;
}
\end{lstlisting}
\inlinecode{FastInput} reads large byte blocks with \inlinecode{fread()} and parses tokens from an internal buffer.

\begin{compactitem}
  \item \inlinecode{FastInput in(file = stdin)} constructs a faster reader from a \inlinecode{FILE*}.
  \item \inlinecode{in \textgreater{}\textgreater{} x} reads a non-whitespace token into \inlinecode{char}, \inlinecode{std::string}, integral types, or floating-point types.
\end{compactitem}

Detected read failures throw \inlinecode{std::runtime\_error}. End-of-file is not an error, and the parser otherwise assumes valid input. Floating-point input is provided for convenience, not as the main performance path.

\begin{lstlisting}[style=cpp]
class FastInput {
  static constexpr int BUF_SIZE = 1 << 20;
  FILE *file;
  char buf[BUF_SIZE];
  int pos = 0, len = 0;

  char get_char() {
    if (pos == len) {
      len = static_cast<int>(fread(buf, 1, BUF_SIZE, file));
      pos = 0;
      if (len == 0) {
        if (ferror(file)) {
          throw std::runtime_error("Failed to read input.");
        }
        return 0;
      }
    }
    return buf[pos++];
  }

  void skip_blanks() {
    while (true) {
      char c = get_char();
      if (!c) {
        return;
      }
      if (!std::isspace(static_cast<unsigned char>(c))) {
        --pos;
        return;
      }
    }
  }

 public:
  explicit FastInput(FILE *file_ = stdin) : file(file_) {}

  FastInput(const FastInput &) = delete;
  FastInput &operator=(const FastInput &) = delete;

  FastInput &operator>>(char &c) {
    skip_blanks();
    c = get_char();
    return *this;
  }

  FastInput &operator>>(std::string &s) {
    skip_blanks();
    s.clear();
    while (true) {
      char c = get_char();
      if (!c || std::isspace(static_cast<unsigned char>(c))) {
        break;
      }
      s += c;
    }
    return *this;
  }

  FastInput &operator>>(bool &x) {
    int val;
    *this >> val;
    x = val != 0;
    return *this;
  }

  template<typename T>
  typename std::enable_if<
      std::is_integral<T>::value && !std::is_same<T, bool>::value, FastInput &>::type
  operator>>(T &x) {
    skip_blanks();
    char c = get_char();
    bool neg = false;
    if (c == '-') {
      neg = true;
      c = get_char();
    }
    using U = typename std::make_unsigned<T>::type;
    U val = 0;
    while (c && !std::isspace(static_cast<unsigned char>(c))) {
      val = U{10} * val + U(c - '0');
      c = get_char();
    }
    if constexpr (std::is_signed<T>::value) {
      x = neg ? T(U{0} - val) : T(val);
    } else {
      x = T(val);
    }
    return *this;
  }

  template<typename T>
  typename std::enable_if<std::is_floating_point<T>::value, FastInput &>::type operator>>(T &x) {
    std::string s;
    *this >> s;
    x = static_cast<T>(std::strtold(s.c_str(), nullptr));
    return *this;
  }
};
\end{lstlisting}
\inlinecode{FastOutput} formats values into an internal buffer and writes them in large blocks with \inlinecode{fwrite()}.

\begin{compactitem}
  \item \inlinecode{FastOutput out(file = stdout)} constructs a faster writer to a \inlinecode{FILE*}.
  \item \inlinecode{out \textless{}\textless{} x} writes \inlinecode{char}, C strings, \inlinecode{std::string}, integral types, or floating-point types.
  \item \inlinecode{out.flush()} writes any buffered output immediately.
\end{compactitem}

Detected write failures throw \inlinecode{std::runtime\_error}. Destruction flushes on a best-effort basis; call \inlinecode{flush()} explicitly to detect a final write failure. Floating-point output is provided for convenience, not as the main performance path.

\begin{lstlisting}[style=cpp]
class FastOutput {
  static constexpr int BUF_SIZE = 1 << 20;
  FILE *file;
  char buf[BUF_SIZE];
  int pos = 0;

  bool flush_buf() {
    int n = pos;
    pos = 0;
    return fwrite(buf, 1, n, file) == static_cast<std::size_t>(n);
  }

  void put_char(char c) {
    if (pos == BUF_SIZE) {
      flush();
    }
    buf[pos++] = c;
  }

 public:
  explicit FastOutput(FILE *file_ = stdout) : file(file_) {}

  ~FastOutput() { flush_buf(); }
  FastOutput(const FastOutput &) = delete;
  FastOutput &operator=(const FastOutput &) = delete;

  void flush() {
    if (!flush_buf()) {
      throw std::runtime_error("Failed to write output.");
    }
  }

  FastOutput &operator<<(char c) {
    put_char(c);
    return *this;
  }

  FastOutput &operator<<(const char *s) {
    while (*s) {
      put_char(*s++);
    }
    return *this;
  }

  FastOutput &operator<<(const std::string &s) {
    for (char c : s) {
      put_char(c);
    }
    return *this;
  }

  FastOutput &operator<<(bool x) {
    put_char(x ? '1' : '0');
    return *this;
  }

  template<typename T>
  typename std::enable_if<
      std::is_integral<T>::value && !std::is_same<T, bool>::value, FastOutput &>::type
  operator<<(T x) {
    using U = typename std::make_unsigned<T>::type;
    U val = U(x);
    if constexpr (std::is_signed<T>::value) {
      if (x < 0) {
        put_char('-');
        val = U{0} - val;
      }
    }
    char s[std::numeric_limits<U>::digits10 + 1];
    int n = 0;
    do {
      s[n++] = char('0' + val % 10);
      val /= 10;
    } while (val);
    while (n--) {
      put_char(s[n]);
    }
    return *this;
  }

  template<typename T>
  typename std::enable_if<std::is_floating_point<T>::value, FastOutput &>::type operator<<(T x) {
    char s[64];
    std::snprintf(
        s, sizeof(s), "%.*Lg", std::numeric_limits<T>::max_digits10, static_cast<long double>(x)
    );
    return *this << s;
  }
};
\end{lstlisting}
\Needspace*{4\baselineskip}
\textbf{Example Usage}\par
\begin{lstlisting}[style=cpp]
#include <cassert>
#include <ctime>
#include <filesystem>
#include <fstream>
#include <iostream>
using namespace std;

int main() {
  // set_io("file.in", "file.out");

  FILE *input = tmpfile();
  fputs("42 hello 3.5 -2147483648 1 Z 1234567890123\n", input);
  rewind(input);

  FastInput in(input);
  int x;
  string s;
  double y;
  int z;
  bool flag;
  char ch;
  long long big;
  in >> x >> s >> y >> z >> flag >> ch >> big;
  assert(x == 42 && s == "hello" && y == 3.5 && z == -2147483648);
  assert(flag && ch == 'Z' && big == 1234567890123LL);
  fclose(input);

  FILE *output = tmpfile();
  {
    FastOutput out(output);
    out << x << ' ' << s << ' ' << y << ' ' << z << ' ' << flag << ' ' << ch << ' ' << big << '\n';
    out.flush();
  }
  rewind(output);
  char buf[96] = {};
  assert(fgets(buf, sizeof(buf), output));
  assert(string(buf) == "42 hello 3.5 -2147483648 1 Z 1234567890123\n");
  fclose(output);

  // Benchmark various ways of reading the same million integers from a scratch file.
  const int n = 1000000;
  string text;
  for (int i = 0; i < n; i++) {
    text += to_string(i) + ' ';
  }
  string path = (filesystem::temp_directory_path() / "fast_io.tmp").string();
  ofstream(path) << text;
  auto time_reads = [&](const char *name, auto read_one) {
    set_in(path);
    clock_t start = clock();
    for (int i = 0, v = 0; i < n; i++) read_one(v);
    printf("%-22s %6.3f s\n", name, double(clock() - start) / CLOCKS_PER_SEC);
  };
  printf("Reading %d integers:\n", n);
  time_reads("cin", [](int &v) { cin >> v; });
  time_reads("scanf", [](int &v) { scanf("%d", &v); });
  time_reads("read_int (unlocked)", [](int &v) { read_int(v); });
  time_reads("FastInput", [](int &v) {
    static FastInput in;
    in >> v;
  });

  auto time_writes = [&](const char *name, auto write_all) {
    clock_t start = clock();
    write_all(n);
    printf("%-22s %6.3f s\n", name, double(clock() - start) / CLOCKS_PER_SEC);
  };
  printf("\nWriting %d integers:\n", n);
  time_writes("ofstream", [&](int k) {
    ofstream out(path);
    for (int i = 0; i < k; i++) out << i << ' ';
  });
  time_writes("fprintf", [&](int k) {
    FILE *f = fopen(path.c_str(), "w");
    for (int i = 0; i < k; i++) fprintf(f, "%d ", i);
    fclose(f);
  });
  time_writes("FastOutput", [&](int k) {
    FILE *f = fopen(path.c_str(), "w");
    {
      FastOutput out(f);
      for (int i = 0; i < k; i++) out << i << ' ';
    }
    fclose(f);
  });
  remove(path.c_str());
  return 0;
}
\end{lstlisting}
\begin{exampleoutput}
Reading 1000000 integers:
cin                     0.959 s
scanf                   0.061 s
read_int (unlocked)     0.010 s
FastInput               0.021 s

Writing 1000000 integers:
ofstream                0.037 s
fprintf                 0.071 s
FastOutput              0.009 s
\end{exampleoutput}

\section{\texorpdfstring{Debugging}{8.3 Debugging}}
\setcounter{section}{3}
\label{sec:8.3}
Local-only debug printing for contest code. \inlinecode{dbg}, \inlinecode{pr}, \inlinecode{dbg\_rows}, and \inlinecode{pr\_rows} expand to no-ops unless \inlinecode{LOCAL} is defined (e.g. via the compiler flag \inlinecode{-DLOCAL}), so they can remain in submitted code without producing output.

\begin{compactitem}
  \item \inlinecode{dbg\_repr(x)} returns a string representation of arithmetic types, strings, pairs, tuples, iterable containers, and custom types that support insertion into \inlinecode{std::ostream}.
  \item \inlinecode{dbg(a, b, c)} prints the function name, line number, argument names, and values to \inlinecode{std::cerr}.
  \item \inlinecode{pr(a, b, c)} prints only the values to \inlinecode{std::cerr}; \inlinecode{pr()} prints a blank line.
  \item \inlinecode{dbg\_rows(rows, col\_width, limit = -1)} prints the function name, line number, argument name, and each iterable row on a separate indexed line. A positive \inlinecode{col\_width} right-aligns each element in a field of that width. Pass $0$ for ordinary container formatting. A nonnegative \inlinecode{limit} prints at most that many rows and elements per row, marking omitted entries with \inlinecode{...}.
  \item \inlinecode{pr\_rows(rows, col\_width, limit = -1)} prints only the indexed rows with the same formatting.
\end{compactitem}

Avoid naming these helpers \inlinecode{to\_string()} unless you want them mixed into overload resolution with \inlinecode{std::to\_string()} and other user-defined \inlinecode{to\_string()} functions. A separate name keeps debug formatting local to this snippet.

\begin{lstlisting}[style=cpp]
#include <cstddef>
#include <iomanip>
#include <iostream>
#include <sstream>
#include <string>
#include <tuple>
#include <type_traits>
#include <utility>

template<typename T, typename = void>
struct is_iterable : std::false_type {};

template<typename T>
struct is_iterable<
    T, std::void_t<decltype(std::begin(std::declval<T>())), decltype(std::end(std::declval<T>()))>>
    : std::true_type {};

template<typename T, typename = void>
struct is_streamable : std::false_type {};

template<typename T>
struct is_streamable<
    T, std::void_t<decltype(std::declval<std::ostream &>() << std::declval<const T &>())>>
    : std::true_type {};

std::string dbg_repr(const std::string &s) { return '"' + s + '"'; }
std::string dbg_repr(const char *s) { return dbg_repr(std::string(s)); }
std::string dbg_repr(char c) { return std::string("'") + c + "'"; }
std::string dbg_repr(bool b) { return b ? "true" : "false"; }

template<typename T>
typename std::enable_if<
    std::is_arithmetic<T>::value && !std::is_same<T, char>::value &&
        !std::is_same<T, bool>::value && !std::is_floating_point<T>::value,
    std::string>::type
dbg_repr(T x) {
  std::ostringstream out;
  out << x;
  return out.str();
}

template<typename T>
typename std::enable_if<std::is_floating_point<T>::value, std::string>::type dbg_repr(T x) {
  std::ostringstream out;
  out << std::setprecision(10) << x;  // Set floating point dbg precision here.
  return out.str();
}

template<typename T>
typename std::enable_if<
    is_streamable<T>::value && !is_iterable<T>::value && !std::is_arithmetic<T>::value &&
        !std::is_convertible<T, std::string>::value,
    std::string>::type
dbg_repr(const T &x) {
  std::ostringstream out;
  out << x;
  return out.str();
}

template<typename A, typename B>
std::string dbg_repr(const std::pair<A, B> &p) {
  return "(" + dbg_repr(p.first) + ", " + dbg_repr(p.second) + ")";
}

template<typename Tuple, std::size_t... Is>
std::string dbg_tuple_repr(const Tuple &t, std::index_sequence<Is...>) {
  std::string res = "(";
  bool first = true;
  ((res += (first ? (first = false, "") : ", ") + dbg_repr(std::get<Is>(t))), ...);
  return res + ")";
}

template<typename... Ts>
std::string dbg_repr(const std::tuple<Ts...> &t) {
  return dbg_tuple_repr(t, std::index_sequence_for<Ts...>{});
}

template<typename T>
typename std::enable_if<
    is_iterable<T>::value && !std::is_convertible<T, std::string>::value, std::string>::type
dbg_repr(const T &v) {
  std::string res = "{";
  bool first = true;
  for (const auto &x : v) {
    if (!first) {
      res += ", ";
    }
    first = false;
    res += dbg_repr(x);
  }
  return res + "}";
}

template<typename Head, typename... Tail>
void dbg_out(bool leading_space, const Head &head, const Tail &...tail) {
  if (leading_space) {
    std::cerr << ' ';
  }
  std::cerr << dbg_repr(head);
  ((std::cerr << ' ' << dbg_repr(tail)), ...);
  std::cerr << '\n';
}

void pr_out() {
  std::cerr << '\n';
}

template<typename Head, typename... Tail>
void pr_out(const Head &head, const Tail &...tail) {
  dbg_out(false, head, tail...);
}

template<typename Rows>
void dbg_rows_out(const Rows &rows, int col_width = 0, int limit = -1) {
  int i = 0;
  for (const auto &row : rows) {
    if (i == limit) {
      std::cerr << "...\n";
      break;
    }
    std::cerr << i++ << ':' << (col_width > 0 ? "" : " {");
    int j = 0;
    for (const auto &x : row) {
      if (j == limit) {
        std::cerr << (col_width > 0 ? " ..." : (j == 0 ? "..." : ", ..."));
        break;
      }
      std::cerr << (col_width > 0 ? " " : (j == 0 ? "" : ", ")) << std::setw(col_width)
                << dbg_repr(x);
      j++;
    }
    std::cerr << (col_width > 0 ? "\n" : "}\n");
  }
}

#ifdef LOCAL
#define dbg(...)                                                                  \
  std::cerr << "[" << __func__ << ":" << __LINE__ << "] " << #__VA_ARGS__ << ":", \
      dbg_out(true, __VA_ARGS__)
#define pr(...) pr_out(__VA_ARGS__)
#define dbg_rows(rows, ...)                                                  \
  std::cerr << "[" << __func__ << ":" << __LINE__ << "] " << #rows << ":\n", \
      dbg_rows_out(rows, __VA_ARGS__)
#define pr_rows(rows, ...) dbg_rows_out(rows, __VA_ARGS__)
#else
#define dbg(...) ((void)0)
#define pr(...) ((void)0)
#define dbg_rows(...) ((void)0)
#define pr_rows(...) ((void)0)
#endif
\end{lstlisting}
\Needspace*{4\baselineskip}
\textbf{Example Usage}\par
\begin{lstlisting}[style=cpp]
#include <cassert>
#include <map>
#include <vector>
using namespace std;

struct Point {
  int x, y;
};

// Can overload dbg_repr to dbg your custom structs.
string dbg_repr(const Point &p) {
  return "Point(" + dbg_repr(p.x) + ", " + dbg_repr(p.y) + ")";
}

struct Edge {
  int u, v, w;
};

string dbg_repr(const Edge &e) {
  return "Edge{" + dbg_repr(e.u) + "->" + dbg_repr(e.v) + ", w=" + dbg_repr(e.w) + "}";
}

struct Streamable {
  int value;

  friend ostream &operator<<(ostream &out, const Streamable &x) {
    return out << "Streamable(" << x.value << ")";
  }
};

int main() {
  vector<int> v{1, 2, 3};
  auto p = make_pair(4, string("x"));
  auto t = make_tuple(1, true, string("ok"));
  map<string, int> mp{{"a", 1}, {"b", 2}};
  Point pt{2, 3};
  Edge e{0, 1, 5};
  double d = 3.141592653589793;

  assert(dbg_repr(v) == "{1, 2, 3}");
  assert(dbg_repr(p) == "(4, \"x\")");
  assert(dbg_repr(t) == "(1, true, \"ok\")");
  assert(dbg_repr(mp) == "{(\"a\", 1), (\"b\", 2)}");
  assert(dbg_repr(pt) == "Point(2, 3)");
  assert(dbg_repr(e) == "Edge{0->1, w=5}");
  assert(dbg_repr(Streamable{6}) == "Streamable(6)");
  assert(dbg_repr(vector<Streamable>{{7}, {8}}) == "{Streamable(7), Streamable(8)}");
  assert(dbg_repr(1.0 / 3.0) == "0.3333333333");
  assert(dbg_repr(d) == "3.141592654");

  dbg(v);
  dbg(p, t, mp);
  dbg(pt, e);
  pr("plain output", d);

  vector<vector<int>> adj{{1, 2}, {0}, {0}};
  pr();
  dbg_rows(adj, 0);

  vector<vector<int>> grid{{1, 20, -3}, {400, 5, 6}, {7, 8, 9}};
  pr();
  dbg_rows(grid, 3, 2);
  return 0;
}
\end{lstlisting}
\begin{exampleoutput}
[main:250] v: {1, 2, 3}
[main:251] p, t, mp: (4, "x") (1, true, "ok") {("a", 1), ("b", 2)}
[main:252] pt, e: Point(2, 3) Edge{0->1, w=5}
"plain output" 3.141592654

[main:257] adj:
0: {1, 2}
1: {0}
2: {0}

[main:261] grid:
0:   1  20 ...
1: 400   5 ...
...
\end{exampleoutput}

\section{\texorpdfstring{Stress-Testing}{8.4 Stress-Testing}}
\setcounter{section}{4}
\label{sec:8.4}
Provides random input generators and a harness that compares an optimized solution against a trusted brute force implementation. Small tests keep brute force practical and produce manageable counterexamples. The generators create common contest inputs such as distinct samples, permutations, graphs, trees, and integer compositions for stress tests, randomized tests, and benchmarks.

\begin{compactitem}
  \item \inlinecode{rng} is a global 64-bit Mersenne Twister. Call \inlinecode{rng.seed(seed)} before a stress test to make it reproducible. Use an explicit seed so generated failures are reproducible.
  \item \inlinecode{rand\_int(lo, hi)} returns a uniformly random integer in range $[{\inlinecodemath{lo}}, {\inlinecodemath{hi}}]$.
  \item \inlinecode{rand\_biased(lo, hi, bias)} returns an integer in $[{\inlinecodemath{lo}}, {\inlinecodemath{hi}}]$. Positive \inlinecode{bias} favors larger values by taking the maximum of \inlinecode{bias + 1} samples, negative \inlinecode{bias} similarly favors smaller values, and zero is uniform.
  \item \inlinecode{rand\_real(lo = 0.0, hi = 1.0)} returns a uniformly random real number in the half-open range $[{\inlinecodemath{lo}}, {\inlinecodemath{hi}})$.
  \item \inlinecode{rand\_choice(values)} returns a uniformly random element of the nonempty vector \inlinecode{values}.
  \item \inlinecode{rand\_weighted(weights)} returns a sampled 0-based index with probability proportional to its nonnegative weight.
  \item \inlinecode{rand\_vec(n, lo, hi)} returns \inlinecode{n} independent uniformly random integers in $[{\inlinecodemath{lo}}, {\inlinecodemath{hi}}]$.
  \item \inlinecode{rand\_str(n, alphabet)} returns a length-\inlinecode{n} string whose characters are sampled independently and uniformly from the nonempty string \inlinecode{alphabet}.
  \item \inlinecode{rand\_perm(n, first = 0)} returns a random permutation of the \inlinecode{n} consecutive integers beginning at \inlinecode{first}.
  \item \inlinecode{rand\_distinct(k, lo, hi)} returns \inlinecode{k} distinct integers sampled uniformly from $[{\inlinecodemath{lo}}, {\inlinecodemath{hi}}]$.
  \item \inlinecode{rand\_tree(n)} returns a uniformly random labeled tree on nodes $[0, {\inlinecodemath{n}})$ as an edge list using a Prüfer code.
  \item \inlinecode{rand\_graph(n, m, directed = false, connected = false)} returns a random simple graph with \inlinecode{n} nodes numbered $[0, {\inlinecodemath{n}})$ and exactly \inlinecode{m} edges. Without connectivity it is uniform; with connectivity it starts from \inlinecode{rand\_tree(n)} and is not uniform over connected graphs. A connected directed graph is weakly connected.
  \item \inlinecode{rand\_composition(parts, sum, min\_part)} uniformly splits \inlinecode{sum} into \inlinecode{parts} ordered integers, each at least \inlinecode{min\_part}.
\end{compactitem}

\Needspace*{3\baselineskip}
\textbf{Time Complexity}:
\begin{compactitem}
  \item $O(1)$ per call to \inlinecode{rand\_int()}, \inlinecode{rand\_real()}, and \inlinecode{rand\_choice()}.
  \item $O(|b| + 1)$ per call to \inlinecode{rand\_biased()} where $b$ is the bias.
  \item $O(n)$ per call to \inlinecode{rand\_weighted()}, \inlinecode{rand\_vec()}, \inlinecode{rand\_str()}, and \inlinecode{rand\_perm()}, where $n$ is the number of weights or generated elements.
  \item $O(k)$ expected per call to \inlinecode{rand\_distinct()}.
  \item $O(n \log n)$ per call to \inlinecode{rand\_tree()}.
  \item $O(m)$ expected per call to \inlinecode{rand\_graph()} without connectivity, or $O(n \log n + m)$ with it.
  \item $O(k \log k)$ per call to \inlinecode{rand\_composition()}, where $k$ is \inlinecode{parts}.
\end{compactitem}

\Needspace*{3\baselineskip}
\textbf{Space Complexity}:
\begin{compactitem}
  \item $O(1)$ auxiliary for the scalar generators and \inlinecode{rand\_choice()}.
  \item $O(n)$ auxiliary for \inlinecode{rand\_weighted()}.
  \item $O(n)$ for the values returned by \inlinecode{rand\_vec()}, \inlinecode{rand\_str()}, \inlinecode{rand\_perm()}, and \inlinecode{rand\_tree()}.
  \item $O(k)$ auxiliary and output space for \inlinecode{rand\_distinct()} and \inlinecode{rand\_composition()}.
  \item $O(m)$ auxiliary and output space for \inlinecode{rand\_graph()}.
\end{compactitem}

\begin{lstlisting}[style=cpp]
#include <algorithm>
#include <cassert>
#include <climits>
#include <cmath>
#include <cstdint>
#include <functional>
#include <initializer_list>
#include <numeric>
#include <optional>
#include <queue>
#include <random>
#include <string>
#include <type_traits>
#include <unordered_set>
#include <utility>
#include <vector>

std::mt19937_64 rng;

template<typename Int>
Int rand_int(Int lo, Int hi) {
  static_assert(std::is_integral<Int>::value, "rand_int() requires an integral type");
  assert(lo <= hi);
  return std::uniform_int_distribution<Int>(lo, hi)(rng);
}

template<typename Int>
Int rand_biased(Int lo, Int hi, int bias) {
  Int value = rand_int(lo, hi);
  while (bias > 0) {
    value = std::max(value, rand_int(lo, hi));
    bias--;
  }
  while (bias < 0) {
    value = std::min(value, rand_int(lo, hi));
    bias++;
  }
  return value;
}

double rand_real(double lo = 0.0, double hi = 1.0) {
  assert(lo < hi);
  return std::uniform_real_distribution<double>(lo, hi)(rng);
}

template<typename T>
T rand_choice(const std::vector<T> &values) {
  assert(!values.empty() && values.size() <= INT_MAX);
  return values[rand_int(0, static_cast<int>(values.size()) - 1)];
}

template<typename W>
int rand_weighted(const std::vector<W> &weights) {
  static_assert(std::is_arithmetic<W>::value, "weights must be numeric");
  assert(!weights.empty() && weights.size() <= INT_MAX);
  std::vector<double> probabilities;
  probabilities.reserve(weights.size());
  double total = 0;
  for (const W &value : weights) {
    double weight = static_cast<double>(value);
    assert(std::isfinite(weight) && weight >= 0);
    probabilities.push_back(weight);
    total += weight;
  }
  assert(std::isfinite(total) && total > 0);
  return std::discrete_distribution<int>(probabilities.begin(), probabilities.end())(rng);
}

template<typename W>
int rand_weighted(std::initializer_list<W> weights) {
  return rand_weighted(std::vector<W>(weights));
}

template<typename Int>
std::vector<Int> rand_vec(int n, Int lo, Int hi) {
  assert(n >= 0 && lo <= hi);
  std::vector<Int> values(n);
  for (Int &x : values) {
    x = rand_int(lo, hi);
  }
  return values;
}

std::string rand_str(int n, const std::string &alphabet) {
  assert(n >= 0 && !alphabet.empty() && alphabet.size() <= INT_MAX);
  std::string s(n, ' ');
  for (char &c : s) {
    c = alphabet[rand_int(0, static_cast<int>(alphabet.size()) - 1)];
  }
  return s;
}

std::vector<int> rand_perm(int n, int first = 0) {
  assert(n >= 0 && (n == 0 || static_cast<int64_t>(first) + n - 1 <= INT_MAX));
  std::vector<int> p(n);
  std::iota(p.begin(), p.end(), first);
  std::shuffle(p.begin(), p.end(), rng);
  return p;
}

std::vector<int> rand_distinct(int k, int lo, int hi) {
  int64_t count = static_cast<int64_t>(hi) - lo + 1;
  assert(lo <= hi && 0 <= k && k <= count);
  std::unordered_set<int> used;
  used.reserve(k);
  std::vector<int> values;
  values.reserve(k);
  // Floyd's algorithm avoids materializing the entire interval.
  for (int64_t j = static_cast<int64_t>(hi) - k + 1; j <= hi; ++j) {
    int x = rand_int(lo, static_cast<int>(j));
    if (!used.insert(x).second) {
      x = static_cast<int>(j);
      used.insert(x);
    }
    values.push_back(x);
  }
  std::shuffle(values.begin(), values.end(), rng);
  return values;
}

std::vector<std::pair<int, int>> rand_tree(int n) {
  assert(n >= 1);
  if (n == 1) {
    return {};
  }
  std::vector<int> degree(n, 1), code(n - 2);
  for (int &u : code) {
    degree[u = rand_int(0, n - 1)]++;
  }
  std::priority_queue<int, std::vector<int>, std::greater<int>> leaves;
  for (int u = 0; u < n; ++u) {
    if (degree[u] == 1) {
      leaves.push(u);
    }
  }
  std::vector<std::pair<int, int>> edges;
  for (int u : code) {
    int leaf = leaves.top();
    leaves.pop();
    edges.emplace_back(leaf, u);
    if (--degree[u] == 1) {
      leaves.push(u);
    }
  }
  int u = leaves.top();
  leaves.pop();
  edges.emplace_back(u, leaves.top());
  std::shuffle(edges.begin(), edges.end(), rng);
  return edges;
}

std::vector<std::pair<int, int>> rand_graph(
    int n, int m, bool directed = false, bool connected = false
) {
  assert(n >= 0 && m >= 0 && (!connected || n >= 1));
  int64_t total =
      directed ? static_cast<int64_t>(n) * (n - 1) : static_cast<int64_t>(n) * (n - 1) / 2;
  int required = connected ? n - 1 : 0;
  assert(required <= m && m <= total);
  std::vector<std::pair<int, int>> edges;
  if (connected) {
    edges = rand_tree(n);
  }
  std::unordered_set<int64_t> required_edges;
  required_edges.reserve(required);
  for (auto &[u, v] : edges) {
    if ((directed && rand_int(0, 1)) || (!directed && u > v)) {
      std::swap(u, v);
    }
    required_edges.insert(static_cast<int64_t>(u) * n + v);
  }
  int64_t available = total - required;
  int extra = m - required;
  bool complement = extra > available / 2;
  int sample = complement ? static_cast<int>(total - m) : extra;
  // For dense graphs, sample the omitted non-tree edges instead of the included ones.
  std::unordered_set<int64_t> sampled;
  sampled.reserve(sample);
  while (static_cast<int>(sampled.size()) < sample) {
    int u = rand_int(0, n - 1), v = rand_int(0, n - 2);
    if (v >= u) {
      v++;
    }
    if (!directed && u > v) {
      std::swap(u, v);
    }
    int64_t key = static_cast<int64_t>(u) * n + v;
    if (!required_edges.count(key)) {
      sampled.insert(key);
    }
  }
  edges.reserve(m);
  if (!complement) {
    for (int64_t key : sampled) {
      edges.emplace_back(static_cast<int>(key / n), static_cast<int>(key % n));
    }
  } else {
    edges.clear();
    for (int u = 0; u < n; u++) {
      for (int v = directed ? 0 : u + 1; v < n; v++) {
        if (u != v && !sampled.count(static_cast<int64_t>(u) * n + v)) {
          edges.emplace_back(u, v);
        }
      }
    }
  }
  std::shuffle(edges.begin(), edges.end(), rng);
  return edges;
}

std::vector<int> rand_composition(int parts, int sum, int min_part) {
  int64_t remaining = static_cast<int64_t>(sum) - static_cast<int64_t>(parts) * min_part;
  assert(parts >= 1 && min_part >= 0 && remaining >= 0 && remaining + parts - 2 <= INT_MAX);
  if (parts == 1) {
    return {sum};
  }
  auto cuts = rand_distinct(parts - 1, 0, static_cast<int>(remaining) + parts - 2);
  std::sort(cuts.begin(), cuts.end());
  std::vector<int> values;
  int prev = -1;
  for (int cut : cuts) {
    values.push_back(min_part + cut - prev - 1);
    prev = cut;
  }
  values.push_back(min_part + static_cast<int>(remaining) + parts - 2 - prev);
  return values;
}
\end{lstlisting}
A stress test compares an optimized implementation against a trusted brute force implementation on many generated inputs. Small tests make brute force practical and produce manageable counterexamples.

\begin{compactitem}
  \item \inlinecode{stress(trials, gen, solve, brute, equal = std::equal\_to\textless{}\textgreater{}())} runs \inlinecode{trials} tests, calling \inlinecode{gen(trial)} with each 1-based trial number and comparing \inlinecode{solve(test)} with \inlinecode{brute(test)} using \inlinecode{equal}. It returns a \inlinecode{StressFailure} containing the number, first failing test, and actual and expected results. The generated test type must be copyable so each solution receives an independent input, and \inlinecode{equal} must accept the respective results of \inlinecode{solve(test)} and \inlinecode{brute(test)}. If all generated tests pass, it returns \inlinecode{std::nullopt}.
\end{compactitem}

\textbf{Time Complexity}: $O(t \cdot (G + X + S + B + C))$ per call in the worst case, where $t$ is \inlinecode{trials}, $X$ is the cost of copying a generated test, and $G$, $S$, $B$, and $C$ are the costs of \inlinecode{gen}, \inlinecode{solve}, \inlinecode{brute}, and \inlinecode{equal}, respectively.

\textbf{Space Complexity}: $O(x + a + e)$ for the test copies and results, where $x$, $a$, and $e$ are their respective sizes, in addition to the working space used by the callbacks.

\begin{lstlisting}[style=cpp]
template<typename Test, typename Actual, typename Expected>
struct StressFailure {
  int trial;
  Test test;
  Actual actual;
  Expected expected;
};

template<typename Generate, typename Solve, typename Brute, typename Equal = std::equal_to<>>
auto stress(int trials, Generate gen, Solve solve, Brute brute, Equal equal = {}) {
  assert(trials >= 0);
  using Test = std::decay_t<decltype(gen(1))>;
  using Actual = std::decay_t<decltype(solve(std::declval<Test &>()))>;
  using Expected = std::decay_t<decltype(brute(std::declval<Test &>()))>;
  using Failure = StressFailure<Test, Actual, Expected>;
  for (int trial = 1; trial <= trials; ++trial) {
    auto test = gen(trial);
    auto solve_test = test, brute_test = test;
    auto actual = solve(solve_test);
    auto expected = brute(brute_test);
    if (!equal(actual, expected)) {
      return std::optional<Failure>{
          Failure{trial, std::move(test), std::move(actual), std::move(expected)}
      };
    }
  }
  return std::optional<Failure>{};
}
\end{lstlisting}
\Needspace*{4\baselineskip}
\textbf{Example Usage}\par
\begin{lstlisting}[style=cpp]
using namespace std;

bool EQ(double a, double b) {
  return a == b || fabs(a - b) <= 1e-9;
}

int main() {
  rng.seed(1234567);  // Fixed seed for reproducibility.

  double real = rand_real(-1.0, 1.0);
  assert(-1.0 <= real && real < 1.0);
  int biased = rand_biased(1, 100000, 3);
  assert(1 <= biased && biased <= 100000);
  assert(rand_weighted({20, 20, 30}) < 3);
  auto values = rand_vec(8, -5, 5);
  assert(values.size() == 8 && all_of(values.begin(), values.end(), [](int x) {
           return -5 <= x && x <= 5;
         }));
  assert(rand_str(12, "abc").size() == 12);
  auto graph = rand_graph(5, 7);
  assert(graph.size() == 7 && all_of(graph.begin(), graph.end(), [](auto edge) {
           return edge.first < edge.second;
         }));
  assert(rand_tree(6).size() == 5);
  assert(rand_graph(6, 8, false, true).size() == 8);
  auto composition = rand_composition(4, 20, 2);
  assert(accumulate(composition.begin(), composition.end(), 0) == 20);
  assert(*min_element(composition.begin(), composition.end()) >= 2);

  auto generate = [](int trial) {
    int n = rand_int(1, 8), kind = trial <= 3 ? trial - 1 : rand_int(0, 2);
    switch (kind) {
      case 0:
        return rand_perm(n, 1);
      case 1:
        return rand_distinct(n, -5, 5);
      default:
        break;
    }
    static const vector<int> boundary_values{-5, 0, 5};
    bernoulli_distribution use_boundary(0.25);
    vector<int> a(n);
    for (int &x : a) {
      x = use_boundary(rng) ? rand_choice(boundary_values) : rand_int(-5, 5);
    }
    return a;
  };
  auto solve = [](vector<int> a) {
    sort(a.begin(), a.end());
    return a;
  };
  auto brute = [](vector<int> a) {
    do {
      if (is_sorted(a.begin(), a.end())) {
        return a;
      }
    } while (next_permutation(a.begin(), a.end()));
    sort(a.begin(), a.end());
    return a;
  };

  assert(!stress(100, generate, solve, brute));
  assert(!stress(
      1, [](int) { return 0; }, [](int) { return 1.0; }, [](int) { return 1.0 + 1e-12; }, EQ
  ));

  // Mutating solvers still receive independent copies of the generated test.
  auto mutate = [](vector<int> &a) { return ++a[0]; };
  auto unchanged = [](const vector<int> &a) { return a[0]; };
  auto failure = stress(1, [](int) { return vector<int>{1}; }, mutate, unchanged);
  assert(failure && failure->trial == 1 && failure->test == vector<int>{1});
  assert(failure->actual == 2 && failure->expected == 1);
  return 0;
}
\end{lstlisting}

\section{\texorpdfstring{Benchmarking}{8.5 Benchmarking}}
\setcounter{section}{5}
\label{sec:8.5}
Small timing helpers for local profiling and codebook examples. Use \inlinecode{std::chrono::steady\_clock} instead of wall-clock time so elapsed durations are monotonic.

\begin{compactitem}
  \item \inlinecode{Timer()} starts a timer immediately.
  \item \inlinecode{timer.reset()} restarts the timer.
  \item \inlinecode{timer.elapsed()} returns the elapsed time in milliseconds as a double.
  \item \inlinecode{benchmark(iterations, f)} runs \inlinecode{f()} \inlinecode{iterations} times and returns the average elapsed milliseconds per call. \inlinecode{iterations} must be positive.
\end{compactitem}

Benchmarks should be treated as local measurements, not proofs of asymptotic performance. Warm-up, CPU frequency scaling, compiler optimization, and input choice can dominate tiny timings.

\begin{lstlisting}[style=cpp]
#include <cassert>
#include <chrono>

class Timer {
  using Clock = std::chrono::steady_clock;
  Clock::time_point start_time;

 public:
  Timer() { reset(); }

  void reset() { start_time = Clock::now(); }

  double elapsed() const {
    return std::chrono::duration<double, std::milli>(Clock::now() - start_time).count();
  }
};

template<typename Fn>
double benchmark(int iterations, Fn f) {
  assert(iterations > 0);
  Timer timer;
  for (int i = 0; i < iterations; ++i) {
    f();
  }
  return timer.elapsed() / iterations;
}
\end{lstlisting}
\Needspace*{4\baselineskip}
\textbf{Example Usage}\par
\begin{lstlisting}[style=cpp]
#include <cassert>
#include <numeric>
#include <vector>
using namespace std;

int main() {
  Timer timer;
  vector<int> v(1000);
  iota(v.begin(), v.end(), 0);
  assert(timer.elapsed() >= 0.0);

  volatile long long sink = 0;
  double avg_time = benchmark(5, [&] { sink += accumulate(v.begin(), v.end(), 0LL); });
  assert(avg_time >= 0.0);
  assert(sink > 0);
  return 0;
}
\end{lstlisting}

\section{\texorpdfstring{GNU Policy-Based Data Structures}{8.6 GNU Policy-Based Data Structures}}
\setcounter{section}{6}
\label{sec:8.6}
GNU policy-based data structures (PBDS) provide non-standard associative containers and priority queues, including order-statistic trees and fast hash tables. They are not part of the C++ standard library, but they are available on many GNU C++ contest judges.

\begin{compactitem}
  \item \inlinecode{HashMap\textless{}K, V\textgreater{}} and \inlinecode{HashSet\textless{}K\textgreater{}} behave like faster, non-standard alternatives to \inlinecode{std::unordered\_map\textless{}K, V\textgreater{}} and \inlinecode{std::unordered\_set\textless{}K\textgreater{}}, using \inlinecode{gp\_hash\_table}. The default \inlinecode{SplitMix64Hasher} is for integer-like keys; pass a custom hash for other key types. It randomizes the hash seed to avoid weak adversarial integer hashes.
  \item \inlinecode{OrderedSet\textless{}T\textgreater{}} and \inlinecode{OrderedMap\textless{}K, V\textgreater{}} behave like \inlinecode{std::set\textless{}T\textgreater{}} and \inlinecode{std::map\textless{}K, V\textgreater{}} but also support \inlinecode{find\_by\_order(k)} and \inlinecode{order\_of\_key(x)}.
  \item \inlinecode{OrderedMultiset\textless{}T\textgreater{}} behaves like \inlinecode{std::multiset\textless{}T\textgreater{}} by storing duplicates as (value, unique ID) pairs.
  \item \inlinecode{OrderedMultimap\textless{}K, V\textgreater{}} behaves like \inlinecode{std::multimap\textless{}K, V\textgreater{}} by storing duplicate keys as (key, unique ID) pairs.
  \item \inlinecode{PairingHeap\textless{}T\textgreater{}} is a meldable priority queue; \inlinecode{push()} returns an iterator that can be passed to \inlinecode{modify()} or \inlinecode{erase()}.
  \item \inlinecode{BinaryHeap\textless{}T\textgreater{}}, \inlinecode{BinomialHeap\textless{}T\textgreater{}}, and \inlinecode{RcBinomialHeap\textless{}T\textgreater{}} are alternative GNU priority queue tags with the same interface.
\end{compactitem}

The order-statistic operations use GNU PBDS conventions: \inlinecode{find\_by\_order(k)} is 0-based and \inlinecode{order\_of\_key(x)} returns the number of keys strictly less than \inlinecode{x}, even if \inlinecode{x} is absent. Raw PBDS ordered trees return iterators from \inlinecode{find\_by\_order(k)}; the multiset and multimap wrappers below unwrap that iterator into a value. The priority queue \inlinecode{join()} operation melds another heap into the current one. Since this depends on \inlinecode{\_\_gnu\_pbds}, keep a standard library fallback in mind when portability matters.

\Needspace*{3\baselineskip}
\textbf{Time Complexity}:
\begin{compactitem}
  \item \inlinecode{HashMap}/\inlinecode{HashSet}: same expected bounds as \inlinecode{std::unordered\_map}/\inlinecode{std::unordered\_set} (expected $O(1)$, worst-case $O(n)$).
  \item \inlinecode{OrderedSet}/\inlinecode{OrderedMap}: same as \inlinecode{std::set}/\inlinecode{std::map}, with $O(\log n)$ for \inlinecode{find\_by\_order} and \inlinecode{order\_of\_key}.
  \item \inlinecode{OrderedMultiset}/\inlinecode{OrderedMultimap}: same as \inlinecode{std::multiset}/\inlinecode{std::multimap}, with $O(\log n)$ for \inlinecode{find\_by\_order} and \inlinecode{order\_of\_key}.
  \item $O(1)$ amortized per call to \inlinecode{push()} and \inlinecode{join()} for \inlinecode{PairingHeap}.
  \item $O(\log n)$ amortized per call to \inlinecode{pop()}, \inlinecode{modify()}, and \inlinecode{erase()} for \inlinecode{PairingHeap}.
\end{compactitem}

\textbf{Space Complexity}: $O(n)$ for all containers.

\begin{lstlisting}[style=cpp]
#include <cassert>
#include <chrono>
#include <cstddef>
#include <cstdint>
#include <ext/pb_ds/assoc_container.hpp>
#include <ext/pb_ds/priority_queue.hpp>
#include <ext/pb_ds/tree_policy.hpp>
#include <functional>
#include <utility>

struct SplitMix64Hasher {
  static uint64_t splitmix64(uint64_t x) {
    x += 0x9e3779b97f4a7c15ULL;
    x = (x ^ (x >> 30)) * 0xbf58476d1ce4e5b9ULL;
    x = (x ^ (x >> 27)) * 0x94d049bb133111ebULL;
    return x ^ (x >> 31);
  }

  std::size_t operator()(uint64_t x) const {
    static const uint64_t RAND_SEED = std::chrono::steady_clock::now().time_since_epoch().count();
    return splitmix64(x + RAND_SEED);
  }
};

template<typename K, typename Hash = SplitMix64Hasher>
using HashSet = __gnu_pbds::gp_hash_table<K, __gnu_pbds::null_type, Hash>;

template<typename K, typename V, typename Hash = SplitMix64Hasher>
using HashMap = __gnu_pbds::gp_hash_table<K, V, Hash>;

template<typename K, typename V, typename Compare = std::less<K>>
using OrderedMap = __gnu_pbds::tree<
    K, V, Compare, __gnu_pbds::rb_tree_tag, __gnu_pbds::tree_order_statistics_node_update>;

template<typename T, typename Compare = std::less<T>>
using OrderedSet = OrderedMap<T, __gnu_pbds::null_type, Compare>;

template<typename T>
class OrderedMultiset {
  OrderedSet<std::pair<T, int>> tree;
  int timer = 0;

 public:
  void insert(const T &x) { tree.insert({x, timer++}); }

  bool erase_one(const T &x) {
    auto it = tree.lower_bound({x, -1});
    if (it == tree.end() || it->first != x) {
      return false;
    }
    tree.erase(it);
    return true;
  }

  T find_by_order(int k) const {
    assert(0 <= k && k < size());
    return tree.find_by_order(k)->first;
  }

  int order_of_key(const T &x) const { return tree.order_of_key({x, -1}); }
  int size() const { return static_cast<int>(tree.size()); }
};

template<typename K, typename V>
class OrderedMultimap {
  OrderedMap<std::pair<K, int>, V> tree;
  int timer = 0;

 public:
  void insert(const K &key, const V &value) { tree.insert({{key, timer++}, value}); }

  bool erase_one(const K &key) {
    auto it = tree.lower_bound({key, -1});
    if (it == tree.end() || it->first.first != key) {
      return false;
    }
    tree.erase(it);
    return true;
  }

  std::pair<K, V> find_by_order(int k) const {
    assert(0 <= k && k < size());
    auto it = tree.find_by_order(k);
    return {it->first.first, it->second};
  }

  int order_of_key(const K &x) const { return tree.order_of_key({x, -1}); }
  int size() const { return static_cast<int>(tree.size()); }
};

template<typename T, typename Compare = std::less<T>>
using PairingHeap = __gnu_pbds::priority_queue<T, Compare, __gnu_pbds::pairing_heap_tag>;

template<typename T, typename Compare = std::less<T>>
using BinaryHeap = __gnu_pbds::priority_queue<T, Compare, __gnu_pbds::binary_heap_tag>;

template<typename T, typename Compare = std::less<T>>
using BinomialHeap = __gnu_pbds::priority_queue<T, Compare, __gnu_pbds::binomial_heap_tag>;

template<typename T, typename Compare = std::less<T>>
using RcBinomialHeap = __gnu_pbds::priority_queue<T, Compare, __gnu_pbds::rc_binomial_heap_tag>;
\end{lstlisting}
\Needspace*{4\baselineskip}
\textbf{Example Usage}\par
\begin{lstlisting}[style=cpp]
int main() {
  OrderedSet<int> s;
  s.insert(10);
  s.insert(20);
  s.insert(30);
  s.insert(20);  // Duplicate ignored, like std::set.
  // Net new over std::set: order-statistic queries.
  assert(s.size() == 3);
  assert(*s.find_by_order(0) == 10);
  assert(*s.find_by_order(1) == 20);
  assert(s.order_of_key(25) == 2);
  assert(s.order_of_key(100) == 3);

  HashMap<int, int> hm;
  hm[10] = 1;
  hm[20] = 2;
  assert(hm[10] == 1);
  HashSet<int> hs;
  hs.insert(7);
  assert(hs.find(7) != hs.end());

  OrderedMultiset<int> ms;
  ms.insert(5);
  ms.insert(5);
  ms.insert(7);
  assert(ms.size() == 3);
  // Net new over std::multiset: order statistics with duplicates.
  assert(ms.order_of_key(7) == 2);
  assert(ms.find_by_order(1) == 5);
  assert(ms.erase_one(5));
  assert(ms.order_of_key(7) == 1);
  assert(ms.erase_one(5));
  assert(!ms.erase_one(5));
  assert(ms.find_by_order(0) == 7);

  OrderedMultimap<int, char> mp;
  mp.insert(5, 'a');
  mp.insert(5, 'b');
  mp.insert(7, 'c');
  // Net new over std::multimap: order statistics with duplicate keys.
  assert(mp.order_of_key(7) == 2);
  assert(mp.find_by_order(1).first == 5);
  assert(mp.erase_one(5));
  assert(mp.order_of_key(7) == 1);
  assert(mp.erase_one(5));
  assert(!mp.erase_one(5));
  assert(mp.find_by_order(0).first == 7);

  PairingHeap<int> h1, h2;
  auto it = h1.push(10);
  auto erased = h1.push(5);
  h1.modify(it, 20);  // Net new over std::priority_queue: update through a handle.
  h1.erase(erased);   // Also net new: remove an arbitrary heap item through a handle.
  h2.push(7);
  h1.join(h2);  // Also net new: meld another heap into this one.
  assert(h2.empty());
  assert(h1.top() == 20);
  h1.pop();
  assert(h1.top() == 7);
  return 0;
}
\end{lstlisting}

\section{\texorpdfstring{Multithreading}{8.7 Multithreading}}
\setcounter{section}{7}
\label{sec:8.7}
Basic multithreading can speed up independent local work such as stress tests, batch experiments, or offline preprocessing. It is rarely useful in standard contest submissions, and many judges restrict threads, so use it only when the environment allows it.

\begin{compactitem}
  \item \inlinecode{parallel\_for(n, threads, f)} runs \inlinecode{f(i)} once for each $0 \leq {\inlinecodemath{i}} < {\inlinecodemath{n}}$, splitting the indices among at most \inlinecode{threads} worker threads.
  \item \inlinecode{parallel\_cases(cases, threads, solve)} solves independent test cases in parallel and returns the outputs in the original case order, using at most \inlinecode{threads} workers and calling \inlinecode{solve(case, output)} for each case.
\end{compactitem}

Each call to \inlinecode{f(i)} must be safe to run concurrently with the others. If it writes shared state, use separate output slots or synchronization.

\begin{lstlisting}[style=cpp]
#include <algorithm>
#include <sstream>
#include <string>
#include <thread>
#include <vector>

template<typename Fn>
void parallel_for(int n, int threads, Fn f) {
  threads = std::max(1, std::min(threads, n));
  std::vector<std::thread> workers;
  workers.reserve(threads);
  for (int id = 0; id < threads; ++id) {
    workers.emplace_back([=, &f] {
      for (int i = id; i < n; i += threads) {
        f(i);
      }
    });
  }
  for (auto &worker : workers) {
    worker.join();
  }
}

template<typename Case, typename Solve>
std::vector<std::string> parallel_cases(const std::vector<Case> &cases, int threads, Solve solve) {
  std::vector<std::string> out(cases.size());
  parallel_for(static_cast<int>(cases.size()), threads, [&](int i) {
    std::ostringstream os;
    solve(cases[i], os);
    out[i] = os.str();
  });
  return out;
}
\end{lstlisting}
\Needspace*{4\baselineskip}
\textbf{Example Usage}\par
\begin{lstlisting}[style=cpp]
#include <cassert>
#include <numeric>
using namespace std;

int main() {
  vector<long long> squares(1000);
  parallel_for(static_cast<int>(squares.size()), 4, [&](int i) { squares[i] = 1LL * i * i; });
  assert(squares[17] == 289);
  assert(accumulate(squares.begin(), squares.end(), 0LL) == 332833500);

  vector<int> cases{3, 1, 4, 2};
  auto out =
      parallel_cases(cases, 2, [](int x, ostringstream &os) { os << x << "^2 = " << x * x; });
  assert((out == vector<string>{"3^2 = 9", "1^2 = 1", "4^2 = 16", "2^2 = 4"}));
  return 0;
}
\end{lstlisting}

\section{\texorpdfstring{Bump Allocation}{8.8 Bump Allocation}}
\setcounter{section}{8}
\label{sec:8.8}
A typed bump pool replaces many small heap allocations with consecutive slots in one fixed vector. This simple form is useful for pointer-based trees, tries, and linked structures whose nodes all live until the algorithm finishes. Allocation advances one index, and individual nodes are never freed.

\begin{compactitem}
  \item \inlinecode{make\_node(value, next = nullptr)} returns a pointer to a new node initialized with \inlinecode{value} and \inlinecode{next}.
  \item \inlinecode{BumpArena\textless{}N\textgreater{}} owns an inline $N$-byte monotonic buffer. Pass \inlinecode{arena.resource()} to a \inlinecode{std::pmr} container such as \inlinecode{std::pmr::vector\textless{}T\textgreater{}} or \inlinecode{std::pmr::map\textless{}K, V\textgreater{}}.
\end{compactitem}

The vector is created at its final size and never resized, so returned pointers remain valid. Adjust \inlinecode{MAX\_NODES} and the fields of \inlinecode{Node} for the problem at hand.

For standard containers, the C++17 polymorphic allocator interface avoids both a global \inlinecode{operator new} override and the full allocator protocol. An arena must outlive every container using it; declare a large arena statically if it would exceed the stack limit. The wrapper uses \inlinecode{std::pmr::null\_memory\_resource()} as its upstream resource, so exhausting the buffer throws \inlinecode{std::bad\_alloc} instead of silently allocating from the heap.

\Needspace*{3\baselineskip}
\textbf{Time Complexity}:
\begin{compactitem}
  \item $O({\inlinecodemath{MAX\_NODES}})$ for construction of the pool.
  \item $O(1)$ per call to \inlinecode{make\_node()}.
\end{compactitem}

\Needspace*{3\baselineskip}
\textbf{Space Complexity}:
\begin{compactitem}
  \item $O({\inlinecodemath{MAX\_NODES}})$ for storage of the pool.
  \item $O(N)$ for storage of each \inlinecode{BumpArena\textless{}N\textgreater{}}.
\end{compactitem}

\begin{lstlisting}[style=cpp]
#include <cassert>
#include <cstddef>
#include <memory_resource>
#include <vector>

struct Node {
  int value;
  Node *next;
};

static constexpr int MAX_NODES = 100000;
static std::vector<Node> node_pool(MAX_NODES);
static int node_count = 0;

Node *make_node(int value, Node *next = nullptr) {
  assert(node_count < MAX_NODES);
  node_pool[node_count] = {value, next};
  return &node_pool[node_count++];
}

template<std::size_t N>
class BumpArena {
  alignas(std::max_align_t) std::byte buffer[N];
  std::pmr::monotonic_buffer_resource arena{buffer, N, std::pmr::null_memory_resource()};

 public:
  std::pmr::memory_resource *resource() { return &arena; }
};
\end{lstlisting}
\Needspace*{4\baselineskip}
\textbf{Example Usage}\par
\begin{lstlisting}[style=cpp]
#include <map>
using namespace std;

int main() {
  Node *tail = make_node(2);
  Node *head = make_node(1, tail);
  assert(head->value == 1 && head->next->value == 2);

  BumpArena<1024> arena;
  pmr::vector<int> values(arena.resource());
  values.reserve(4);
  values.insert(values.end(), {1, 2, 3, 4});
  assert(values.size() == 4 && values.back() == 4);

  pmr::map<int, int> counts(arena.resource());
  counts.emplace(7, 2);
  assert(counts[7] == 2);
  return 0;
}
\end{lstlisting}


\backmatter

\printindex

\end{document}
